\beginnumbering
\eledsection*{Adhyāya 1}
\pstart \ledsidenote{\passref{3.1.1}}\skt{अथातः }\edtext{\skt{सर्व्वभूतचिन्तां}}{\lemma{\skt{सर्व्वभूतचिन्तां}}\Afootnote{K H N; \skt{सर्वभूतचिन्ता}\ A; \skt{सर्वभूतञ्चिन्ता}\ L}}\skt{ शारीरं }\edtext{\skt{व्याख्यास्यामः}}{\lemma{\skt{व्याख्यास्यामः}}\Afootnote{H A L N; \emph{om.}\ K}}\skt{~॥}\footnoteA{Witness K breaks off here, resuming only in the second chapter of the Cikitsāsthāna.}\pend

\pstart \ledsidenote{\passref{3.1.3}}\edtext{\edtext{\skt{सर्व्व\-भू\-ता\-ना\-ङ्का\-र\-ण\-म\-का\-रणं}}{\lemma{\skt{सर्व्वभूतानाङ्}}\Afootnote{A N H; \skt{अथ खल्ववर्म स सुश्रुत सर्वभूतानां}\ L}}}{\lemma{\skt{कारणम्}}\Afootnote{A N H; \skt{कारिणम्}\ L}}\skt{ }\edtext{\edtext{\edtext{\skt{सत्त्व\-र\-ज\-स्त\-मो\-ल\-क्ष\-ण\-म\-ष्ट\-रू\-प\-म\-खि\-ल\-स्यास्य}}{\lemma{\skt{सत्त्व\appbrk{}र\appbrk{}ज\appbrk{}स्त\appbrk{}मो\appbrk{}ल\appbrk{}क्ष\appbrk{}णम्}}\Afootnote{A N\textsuperscript{pc} H L; \skt{सत्व\appbrk{}र\appbrk{}ज\appbrk{}स्त\appbrk{}मो\appbrk{}क्ष\appbrk{}ल्ल\appbrk{}क्ष\appbrk{}णम्}\ N\textsuperscript{ac}}}}{\lemma{\skt{अष्टरूपम्}}\Afootnote{A N\textsuperscript{pc} H; \skt{अष्टरूपाम्}\ N\textsuperscript{ac}; \skt{ष्टरूपम्}\ L}}}{\lemma{\skt{अखिलस्यास्य}}\Afootnote{N H; \skt{अखिलस्य}\ A L}}\skt{ जगतः }\edtext{\skt{संभ\-व\-हे\-तु\-र\-व्य\-क्त\-न्नाम}}{\lemma{\skt{संभवहेतुर्}}\Afootnote{A N H; \skt{संभविहेतुर्}\ L}}\skt{~। }\edtext{\skt{तदे\-क\-म\-ने\-केषां}}{\lemma{\skt{अनेकेषां}}\Afootnote{N H L; \skt{बहूनां}\ A}}\skt{ }\edtext{\skt{क्षेत्र\-ज्ञा\-ना\-म\-धि\-ष्ठानं}}{\lemma{\skt{अधिष्ठानं}}\Afootnote{A N L; \skt{अधिष्ठानां}\ H}}\skt{ समुद्र }\edtext{\skt{इवोदकौकसां}}{\lemma{\skt{इवोदकौकसां}}\Afootnote{\skt{इवौदकानां}\ A L; \skt{इवा\appbrk{}द\appbrk{}कौ\appbrk{}ज\appbrk{}साम्}\ N; \skt{इवो\appbrk{}द\appbrk{}कौ\appbrk{}ज\appbrk{}साम्}\ H}}\skt{ }\edtext{\skt{भावानाम्}}{\lemma{\skt{भावानाम्}}\Afootnote{A H L; \skt{भावानान्}\ N}}\skt{~॥}\footnoteA{Witnesses N and H read \emph{samudra ivodakaujasāṃ bhāvānām}, which is linguistically and semantically implausible. It might mean "like the ocean is a support for creatures that take their strength (\emph{ojas}) from water." But the simile and the language remain obscure and unlikely. Philipp Maas has suggested the minimal emendation to \emph{ivodakaukasāṃ} which is compelling semantically and palaeographically. It relies only on the confusion of the Nepalese characters \emph{𑐖} and \emph{𑐎}.}\pend

\pstart \ledsidenote{\passref{3.1.4}}\edtext{\skt{तस्मा\-द\-व्य\-क्ता\-न्म\-हा\-नु\-त्प\-द्यते}}{\lemma{\skt{तस्मादव्यक्तान्}}\Afootnote{A H L; \skt{तस्मादव्यक्ताम्}\ N}}\skt{ तल्लिङ्ग }\edtext{\skt{एव}}{\lemma{\skt{एव}}\Afootnote{A N H; \skt{ए}\ L}}\skt{~। }\edtext{\skt{तल्लिङ्गाच्च}}{\lemma{\skt{तल्लिङ्गाच्च}}\Afootnote{A N H; \skt{तल्लिंगाच्चं}\ L}}\skt{ }\edtext{\skt{मह\-त\-स्त\-ल्ल\-क्षण}}{\lemma{\skt{महतस्}}\Afootnote{A N\textsuperscript{pc} H L; \skt{महातस्}\ N\textsuperscript{ac}}}\skt{ एवाहङ्कार उत्पद्यते~। स }\edtext{\skt{त्रिविधो}}{\lemma{\skt{त्रिविधो}}\Afootnote{A N\textsuperscript{pc} H L; \skt{त्रिविधौ}\ N\textsuperscript{ac}}}\skt{ }\edtext{\edtext{\skt{वैका\-रि\-क\-स्तै\-जसो}}{\lemma{\skt{वैकारिकस्}}\Afootnote{A N H; \skt{वैकारिक}\ L}}}{\lemma{\skt{तैजसो}}\Afootnote{A H; \skt{तेजसो}\ N; \skt{स्तैजसो}\ L}}\skt{ भूतादिरिति~। तत्र }\edtext{\edtext{\skt{वैका\-रि\-का\-द\-हं\-का\-रा\-त्त\-ल्ल\-क्ष\-णा\-न्ये\-वै\-का\-द\-शे\-न्द्रि\-या\-ण्यु\-त्प\-द्यन्ते}}{\lemma{\skt{वैका\appbrk{}रि\appbrk{}का\appbrk{}द\appbrk{}हं\appbrk{}का\appbrk{}रात्}}\Afootnote{N H L; \skt{वैका\appbrk{}रि\appbrk{}का\appbrk{}द\appbrk{}ह\appbrk{}ङ्का\appbrk{}रा\appbrk{}त्तै\appbrk{}ज\appbrk{}स\appbrk{}स\appbrk{}हा\appbrk{}यात्}\ A}}}{\lemma{\skt{एवै\appbrk{}का\appbrk{}द\appbrk{}शे\appbrk{}न्द्रि\appbrk{}याण्य्}}\Afootnote{A N H; \skt{एवै\appbrk{}का\appbrk{}दे\appbrk{}शे\appbrk{}न्द्रि\appbrk{}याण्य्}\ L}}\skt{~॥ }\edtext{\skt{तद्यथा}}{\lemma{\skt{तद्}}\Afootnote{A N H L\textsuperscript{pc}; \skt{तथ्यद्}\ L\textsuperscript{ac}}}\skt{~। }\edtext{\skt{श्रोत्रं}}{\lemma{\skt{श्रोत्रं}}\Afootnote{H; \skt{श्रोत्र}\ A N L}}\skt{ }\edtext{\edtext{\skt{त्वक्चक्षुर्ज्जिह्वा}}{\lemma{\skt{त्वक्}}\Afootnote{A N H; \skt{त्व}\ L}}}{\lemma{\skt{चक्षुर्ज्जिह्वा}}\Afootnote{H; \skt{चक्षुर्जिह्वा}\ A; \skt{चक्षु जिह्वा}\ N L}}\skt{ }\edtext{\skt{घ्राणो}}{\lemma{\skt{घ्राणो}}\Afootnote{N H; \skt{घ्राण}\ A L}}\skt{ }\edtext{\skt{वाग्घ\-स्तो\-प\-स्थ\-पा\-यु\-पा\-द\-म\-ना\-न्सीति}}{\lemma{\skt{वाग्घ\appbrk{}स्तो\appbrk{}प\appbrk{}स्थ\appbrk{}पा\appbrk{}यु\appbrk{}पा\appbrk{}द\appbrk{}म\appbrk{}ना\appbrk{}न्सीति}}\Afootnote{A H; \skt{वाग्म\appbrk{}स्तो\appbrk{}प\appbrk{}स्थ\appbrk{}पा\appbrk{}यु\appbrk{}पा\appbrk{}द\appbrk{}म\appbrk{}नां\appbrk{}सीति}\ N; \skt{वाघ्घ\appbrk{}स्तो\appbrk{}प\appbrk{}स्थ\appbrk{}पा\appbrk{}यु\appbrk{}पा\appbrk{}द\appbrk{}म\appbrk{}नांसि}\ L}}\skt{~। तत्र पूर्व्वाणि पञ्च बुद्धीन्द्रियाणि~। इतराणि पञ्च }\edtext{\skt{कर्म्मेन्द्रियाणि}}{\lemma{\skt{कर्म्मेन्द्रियाणि}}\Afootnote{A N H; \skt{कर्म्मोन्द्रियाणि}\ L}}\skt{~। उभ\-या\-त्म\-क\-म्मनः~। }\edtext{\edtext{\skt{भूता\-दे\-स्त\-ल्ल\-क्ष\-णा\-न्येव}}{\lemma{\skt{तल्लक्षणान्य्}}\Afootnote{N H L; \skt{अपि तैज\appbrk{}स\appbrk{}स\appbrk{}हा\appbrk{}या\appbrk{}त्त\appbrk{}ल्ल\appbrk{}क्ष\appbrk{}णान्य्}\ A}}}{\lemma{\skt{एव}}\Afootnote{A N H L\textsuperscript{pc}; \skt{एवआ}\ L\textsuperscript{ac}}}\skt{ पञ्च\-त\-न्मा\-त्रा\-ण्यु\-त्प\-द्यन्ते~॥ }\edtext{\edtext{\skt{तद्यथा}}{\lemma{\skt{तद्}}\Afootnote{N H L; \emph{om.}\ A}}}{\lemma{\skt{यथा}}\Afootnote{N H L; \emph{om.}\ A}}\skt{~॥ शब्दतन्मात्रं~। }\edtext{\skt{स्पर्शतन्मात्रं}}{\lemma{\skt{स्पर्शतन्मात्रं}}\Afootnote{A H; \skt{स्पर्शतन्मा}\ N; \skt{स्पर्श्शतन्मात्रम्}\ L}}\skt{~। रूपतन्मात्रं~। रसतन्मात्रं~। गन्ध\-त\-न्मा\-त्र\-मिति~। }\edtext{\skt{तेभ्यो}}{\lemma{\skt{तेभ्यो}}\Afootnote{N H; \skt{ते}\ A; \emph{om.}\ L}}\skt{ }\edtext{\skt{भूतानि}}{\lemma{\skt{भूतानि}}\Afootnote{N H; \emph{om.}\ A; \skt{भूता}\ L}}\skt{ }\edtext{\edtext{\edtext{\skt{व्योमा\-नि\-ला\-न\-ल\-ज\-लो\-र्व्व्य\-स्ते\-भ्य\-स्त\-द्वि\-शेषाः}}{\lemma{\skt{व्योमा\appbrk{}नि\appbrk{}ला\appbrk{}न\appbrk{}ल\appbrk{}ज\appbrk{}लो\appbrk{}र्व्व्यस्}}\Afootnote{H; \skt{षां}\ A; \skt{व्योम\appbrk{}नि\appbrk{}ला\appbrk{}न\appbrk{}ल\appbrk{}ज\appbrk{}लो\appbrk{}र्व्यस्}\ N; \skt{नि}\ L}}}{\lemma{\skt{तेभ्यस्}}\Afootnote{N H; \emph{om.}\ A; \skt{ते सस्}\ L}}}{\lemma{\skt{तद्विशेषाः}}\Afootnote{N H L; \skt{विशेषाः}\ A}}\skt{ }\edtext{\skt{शब्द\-स्प\-र्श\-रू\-प\-र\-स\-गन्धाः}}{\lemma{\skt{शब्द\appbrk{}स्प\appbrk{}र्श\appbrk{}रू\appbrk{}प\appbrk{}र\appbrk{}स\appbrk{}गन्धाः}}\Afootnote{N H L; \skt{शब्द\appbrk{}स्प\appbrk{}र्श\appbrk{}रू\appbrk{}प\appbrk{}र\appbrk{}स\appbrk{}गन्धाः}\ A}}\skt{~। }\edtext{\edtext{\skt{एवमेतानि}}{\lemma{\skt{एवम्}}\Afootnote{N H; \skt{तेभ्यो भूतानि -- व्योमानि}\ A; \skt{एवव्मेतानि}\ L\textsuperscript{ac}; \skt{एवमेतानि}\ L\textsuperscript{pc}}}}{\lemma{\skt{एतानि}}\Afootnote{N H; \skt{लानल}\ A; \emph{om.}\ L}}\skt{ }\edtext{\skt{चतुर्व्विंशति}}{\lemma{\skt{चतुर्व्विंशति}}\Afootnote{N H; \skt{जलोर्व्यः; एवमेषा}\ A; \emph{om.}\ L}}\skt{ }\edtext{\skt{तत्त्वानि}}{\lemma{\skt{तत्त्वानि}}\Afootnote{N H L; \skt{तत्त्व}\ A}}\skt{ }\edtext{\skt{व्याख्यातानि}}{\lemma{\skt{व्याख्यातानि}}\Afootnote{N H; \skt{चतु\appbrk{}र्विं\appbrk{}श\appbrk{}ति\appbrk{}र्व्या\appbrk{}ख्याता}\ A; \skt{चतु\appbrk{}र्विं\appbrk{}श\appbrk{}ति\appbrk{}र्व्या\appbrk{}ख्या\appbrk{}तानि}\ L}}\skt{ }\edtext{\skt{भवन्ति}}{\lemma{\skt{भवन्ति}}\Afootnote{N H; \emph{om.}\ A L}}\skt{~।}\pend

\pstart \ledsidenote{\passref{3.1.5}}\skt{तत्र बुद्धीन्द्रियाणां शब्दादयो }\edtext{\skt{विषयाः}}{\lemma{\skt{विषयाः}}\Afootnote{A N H; \skt{विषयोः}\ L}}\skt{~। कर्म्मेन्द्रियाणां }\edtext{\skt{यथासंख्यं}}{\lemma{\skt{यथासंख्यं}}\Afootnote{A L; \skt{यथासंख्य}\ N; \skt{यथासंख्या}\ H}}\skt{ }\edtext{\skt{वच\-ना\-दा\-ना\-न\-न्द\-वि\-स\-र्ग्ग\-वि\-ह\-र\-णानि}}{\lemma{\skt{वच\appbrk{}ना\appbrk{}दा\appbrk{}ना\appbrk{}न\appbrk{}न्द\appbrk{}वि\appbrk{}स\appbrk{}र्ग्ग\appbrk{}वि\appbrk{}ह\appbrk{}र\appbrk{}णानि}}\Afootnote{A N H; \skt{वच\appbrk{}ना\appbrk{}दा\appbrk{}ना\appbrk{}न\appbrk{}न्द\appbrk{}वि\appbrk{}स\appbrk{}र्ग\appbrk{}वि\appbrk{}चा\appbrk{}र\appbrk{}णानि}\ L}}\skt{~।}\pend

\pstart \ledsidenote{\passref{3.1.6}}\edtext{\edtext{\skt{अव्य\-क्त\-म्म\-हा\-न\-हं\-कारः}}{\lemma{\skt{अव्यक्तम्}}\Afootnote{A N H; \skt{अर्वव्यक्त}\ L\textsuperscript{ac}; \skt{अव्यक्त}\ L\textsuperscript{pc}}}}{\lemma{\skt{महानहंकारः}}\Afootnote{A N H; \skt{महानहंकाराः}\ L}}\skt{ पञ्च तन्मात्राणि }\edtext{\skt{चेत्यष्टौ}}{\lemma{\skt{चेत्य्}}\Afootnote{A N\textsuperscript{pc} H L; \skt{चे>ष्टत्य्}\ N\textsuperscript{ac}}}\skt{ प्रकृतयः~। शेषाः षोडश विकाराः~।}\pend

\pstart \ledsidenote{\passref{3.1.7}}\edtext{\skt{स्वः}}{\lemma{\skt{स्वः}}\Afootnote{A; \skt{स्वं}\ N H}}\skt{ }\edtext{\edtext{\skt{स्वश्चैषां}}{\lemma{\skt{स्वश्}}\Afootnote{A; \skt{स्वंश्}\ N; \skt{स्वं}\ H}}}{\lemma{\skt{चैषां}}\Afootnote{A N; \skt{स्वैषां}\ H}}\skt{ }\edtext{\edtext{\skt{विष\-योऽ\-धि\-भू\-तम्}}{\lemma{\skt{विषयो}}\Afootnote{A; \skt{विषये}\ N H}}}{\lemma{\skt{ऽधिभूतम्}}\Afootnote{A; \skt{विभुत्वं}\ N H}}\skt{~।}\footnoteA{We emend to the vulgate text, especially in the light of the other very rare instances of these parallels in Sanskrit literature. It is hard to make grammatical or semantic sense of the Nepalese readings \emph{svaṃ svaṃ svaiṣām} (or \emph{ścaiṣām} or \emph{scaiṣām}). Perhaps the original reading \emph{dhibhūtam} was first mistaken to be \emph{vibhūtam}. This is a plausible error since \emph{dhi} and \emph{vi} are similar in early Nepalese orthography. Subsequently, a scribe might have normalized \emph{vibhūtam} to \emph{vibhutvam} in order to make sense in the light of the following list of deities. If these conjectures are correct, these changes must have happened very early in the transmission, in an ancestor manuscript of both witnesses K and H. The word \emph{vibhutva} occurs in Sāṅkhyakārikā 42, although in a quite different context.}\skt{ }\edtext{\skt{स्वयमध्यात्मम्}}{\lemma{\skt{अध्यात्मम्}}\Afootnote{A H; \skt{अध्यात्म}\ N}}\skt{~। }\edtext{\skt{अधिदेवन्तु}}{\lemma{\skt{अधिदेवन्तु}}\Afootnote{H; \skt{अधिदैवतम्}\ A; \skt{धिदेवन्तु}\ N}}\skt{~।}\footnoteA{The \emph{tu} is awkward. Witnesses N and H have daṇḍas after the \emph{tu}.}\skt{ अथ }\edtext{\edtext{\skt{बुद्धेर्ब्रह्मा}}{\lemma{\skt{बुद्धेर्}}\Afootnote{A H; \skt{बुद्धे}\ N}}}{\lemma{\skt{ब्रह्मा}}\Afootnote{A H; \skt{व्रह्मी}\ N}}\skt{~। अहं\-का\-र\-स्ये\-श्वरः~। मनसश्चन्द्रमाः~। दिशः श्रोत्रस्य~। त्वचो वायुः~। सूर्यश्चक्षुषः~। }\edtext{\skt{रसनस्यापः}}{\lemma{\skt{रसनस्यापः}}\Afootnote{A N; \skt{रसना…पः}\ H}}\skt{~। पृथिवी घ्राणस्य~। }\edtext{\skt{वाचोऽग्निः}}{\lemma{\skt{ऽग्निः}}\Afootnote{A H; \skt{ग्निः}\ N}}\skt{~। }\edtext{\skt{हस्तयोरिन्द्रः}}{\lemma{\skt{इन्द्रः}}\Afootnote{A N; \skt{इन्द्र}\ H}}\skt{~। पादयोर्व्विष्णुः~। पायोर्म्मित्रः~। प्रजा\-प\-ति\-रु\-प\-स्थ\-स्येति~।}\footnoteA{Parallel to the homologies described in Mahābhārata 14.42.27–40, which relate ontological triples.}\pend

\pstart \ledsidenote{\passref{3.1.8}}\skt{तत्र सर्व्व एवाचेतनः }\edtext{\skt{एष}}{\lemma{\skt{एष}}\Afootnote{A N H; \emph{om.}\ L}}\skt{ वर्ग्गः~।}\footnoteA{It is tempting to emend \emph{varga} to \emph{sarga} {\symfont `}{\symfont `}the emitted creation,'' given the use of the term in Sāṅkhyakārikā 24 et passim. However, the vulgate text and Ḍalhaṇa's commentary unanimously read \emph{varga}.}\skt{ }\edtext{\skt{पुरुषश्च}}{\lemma{\skt{पुरुषश्च}}\Afootnote{N H; \skt{पुरुषः}\ A L}}\skt{ }\edtext{\edtext{\skt{पञ्च\-विं\-श\-ति\-त\-म\-श्चे\-त\-यिता}}{\lemma{\skt{पञ्च\appbrk{}विं\appbrk{}श\appbrk{}ति\appbrk{}त\appbrk{}मश्}}\Afootnote{A H; \skt{पञ्चविङ्शतिश्}\ N; \skt{पञ्च\appbrk{}विं\appbrk{}श\appbrk{}ति\appbrk{}कञ्}\ L}}\skt{~। स तु}}{\lemma{\skt{चेतयिता~। स तु}}\Afootnote{N H; \emph{om.}\ A; \skt{चेतयिता सो वि}\ L}}\skt{ }\edtext{\skt{कर\-ण\-का\-र्य\-सं\-प्र\-युक्तः~। }\edtext{\skt{सत्यप्यचैतन्ये}}{\lemma{\skt{अचैतन्ये}}\Afootnote{A H L; \skt{आचैतन्ये}\ N}}}{\lemma{\skt{कर\appbrk{}ण\appbrk{}का\appbrk{}र्य\appbrk{}सं\appbrk{}प्र\appbrk{}युक्तः~। सत्य्}}\Afootnote{H; \skt{कार्य\appbrk{}का\appbrk{}र\appbrk{}ण\appbrk{}सं\appbrk{}यु\appbrk{}क्त\appbrk{}श्चे\appbrk{}त\appbrk{}यिता भवति~। सत्य्}\ A; \skt{कर\appbrk{}ण\appbrk{}का\appbrk{}र्य\appbrk{}स\appbrk{}म्प्र\appbrk{}युक्ताः~। सत्य्}\ N; \skt{कार्य\appbrk{}का\appbrk{}र\appbrk{}ण\appbrk{}सं\appbrk{}यु\appbrk{}क्त\appbrk{}स्सत्य्}\ L}}\skt{ }\edtext{\skt{प्रधानस्य}}{\lemma{\skt{प्रधानस्य}}\Afootnote{A N H; \skt{प्रधाने}\ L}}\skt{ }\edtext{\edtext{\skt{पुरु\-ष\-कै\-व\-ल्या\-र्था\-म्प्र\-वृ\-त्ति\-मु\-प\-दि\-शन्ति}}{\lemma{\skt{पुरु\appbrk{}ष\appbrk{}कै\appbrk{}व\appbrk{}ल्या\appbrk{}र्थाम्}}\Afootnote{N; \skt{पुरु\appbrk{}ष\appbrk{}कै\appbrk{}व\appbrk{}ल्यार्थं}\ A; \skt{पुरु\appbrk{}ष\appbrk{}कै\appbrk{}व\appbrk{}ल्या\appbrk{}र्थान्}\ H; \skt{पुरुषस्य कैवल्यार्त्थां}\ L}}}{\lemma{\skt{प्रवृत्तिम्}}\Afootnote{A H L; \skt{प्रवृत्तम्}\ N}}\skt{~। }\edtext{\skt{क्षीरादींश्चात्र}}{\lemma{\skt{क्षीरादींश्}}\Afootnote{A N H; \skt{क्षीरादीञ्}\ L}}\skt{ हेतू\-नु\-दा\-ह\-रन्ति~॥}\pend

\pstart \ledsidenote{\passref{3.1.9}}\skt{अत ऊर्ध्व\-म्प्र\-कृ\-ति\-पु\-रु\-षयोः साधर्म्म्यवैधर्म्म्ये }\edtext{\skt{व्याख्यास्यामः}}{\lemma{\skt{व्याख्यास्यामः}}\Afootnote{N H; \skt{व्याख्यास्यामः~। तद्यथा}\ A L}}\skt{~॥ उभावप्यनादी }\edtext{\skt{उभावप्यनन्तौ}}{\lemma{\skt{अनन्तौ}}\Afootnote{A N H; \skt{अलिंगौ}\ L}}\skt{~। उभावपि }\edtext{\skt{नित्यौ}}{\lemma{\skt{नित्यौ}}\Afootnote{\skt{अलिङ्गौ, उभावपि नित्यौ}\ A; \skt{अनित्यौ}\ N H; \skt{अनन्तौ~। उभावपि नित्यौ}\ L}}\skt{~।}\footnoteA{It is unexpected to assert, as the Nepalese witnesses do, that Puruṣa and Prakṛti are impermanent. Perhaps at an earlier point in the text's transmission, a scribe improperly repeated the initial alpha-privative that is used in every other qualifier in the list. This emendation is influenced by the very similar wording of Mahābhārata Śāntiparvan 210.8ab: \emph{ubhau nityau ....} (see note on testimonia). Note also that the sequence of qualifiers in the Nepalese manuscripts is different than the vulgate: \emph{anitya-apara-aliṅga} rather than \emph{aliṅga-nitya-anapara}.}\skt{ }\edtext{\edtext{\edtext{\skt{उभावप्यपरौ}}{\lemma{\skt{उभाव्}}\Afootnote{A N H; \emph{om.}\ L}}}{\lemma{\skt{अप्य्}}\Afootnote{A N H; \emph{om.}\ L}}}{\lemma{\skt{अपरौ}}\Afootnote{N H; \skt{अन}\ A; \emph{om.}\ L}}\skt{ }\edtext{\edtext{\edtext{\skt{उभावप्यलिङ्गौ}}{\lemma{\skt{उभाव्}}\Afootnote{N H; \emph{om.}\ A L}}}{\lemma{\skt{अप्य्}}\Afootnote{N H; \emph{om.}\ A L}}}{\lemma{\skt{अलिङ्गौ}}\Afootnote{H; \skt{परौ}\ A; \skt{अलिङ्गे}\ N; \emph{om.}\ L}}\skt{ उभौ च सर्व्वगताविति~॥}\footnoteA{Cf. Mahābhārata Śāntiparvan 210.7cd-8ab: \emph{ anādyantāv ubhāv etāv aliṅgau cāpy ubhāv api // 7 ubhau nityau sūkṣmatarau mahadbhyaś ca mahattarau /}.}\skt{ एका तु प्रकृ\-ति\-र\-चे\-तना }\edtext{\skt{त्रिगुणा}}{\lemma{\skt{त्रिगुणा}}\Afootnote{A H L; \skt{तिगुणा}\ N}}\skt{ }\edtext{\skt{बीजधर्म्मिणी}}{\lemma{\skt{बीजधर्म्मिणी}}\Afootnote{A N H; \skt{बीजधर्मणी}\ L}}\skt{ }\edtext{\skt{प्रसवधर्म्मिनी}}{\lemma{\skt{प्रसवधर्म्मिनी}}\Afootnote{N H; \skt{प्रसवधर्मिण्य्}\ A; \emph{om.}\ L}}\skt{ }\edtext{\skt{अम\-ध्य\-स्थ\-ध\-र्म्मिणी}}{\lemma{\skt{अम\appbrk{}ध्य\appbrk{}स्थ\appbrk{}ध\appbrk{}र्म्मिणी}}\Afootnote{A N; \skt{मध्यधर्म्मिणी}\ H; \skt{अम\appbrk{}द्ध्य\appbrk{}स्थ\appbrk{}ध\appbrk{}र्म्मणी}\ L}}\footnoteA{Note the sandhi hiatus, \emph{-dharmiṇī - amadhya-}}\skt{ चेति~। बहवस्तु }\edtext{\skt{पुरु\-षा\-श्चे\-त\-ना\-वन्तः}}{\lemma{\skt{पुरुषाश्}}\Afootnote{A H L; \skt{पुरुषा}\ N}}\skt{~।}\footnoteA{The daṇḍa here breaks the sentence but disambiguates the sandhi.}\skt{ अगुणा }\edtext{\edtext{\skt{अबी\-ज\-ध\-र्म्मि\-णोऽ\-प्र\-स\-व\-ध\-र्म्मिणो}}{\lemma{\skt{अबीजधर्म्मिणो}}\Afootnote{H; \skt{अबीजधर्माणो}\ A; \skt{अबीज}\ N; \skt{अबीजधर्म्मिणीः}\ L}}}{\lemma{\skt{ऽप्र\appbrk{}स\appbrk{}व\appbrk{}ध\appbrk{}र्म्मिणो}}\Afootnote{H L; \skt{ऽप्र\appbrk{}स\appbrk{}व\appbrk{}ध\appbrk{}र्माणो}\ A; \skt{धर्मिणो}\ N}}\skt{ }\edtext{\skt{मध्य\-स्थ\-ध\-र्म्मि\-ण\-श्चेति}}{\lemma{\skt{मध्य\appbrk{}स्थ\appbrk{}ध\appbrk{}र्म्मि\appbrk{}णश्}}\Afootnote{N; \skt{मध्य\appbrk{}स्थ\appbrk{}ध\appbrk{}र्मा\appbrk{}णश्}\ A; \skt{ऽम\appbrk{}ध्य\appbrk{}स्थ\appbrk{}ध\appbrk{}र्म्मि\appbrk{}णश्}\ H; \skt{मद्ध्यस्थाश्}\ L}}\skt{~।}\pend

\pstart \ledsidenote{\passref{3.1.10}}\skt{तत्र कारणानुरूपं कार्यमिति कृत्वा सर्व्व }\edtext{\skt{एवैते च}}{\lemma{\skt{एवैते च}}\Afootnote{H; \skt{एवैते}\ A L N}}\skt{ }\edtext{\skt{विशेषाः}}{\lemma{\skt{विशेषाः}}\Afootnote{A N\textsuperscript{pc} H L; \skt{विशोषाः}\ N\textsuperscript{ac}}}\skt{ सत्त्व\-र\-ज\-स्त\-मो\-मया }\edtext{\skt{भवन्ति}}{\lemma{\skt{भवन्ति}}\Afootnote{A N H; \skt{भवन्ती}\ L}}\skt{~। }\edtext{\edtext{\skt{तद\-ञ्ज\-न\-त्वा\-त्त\-न्म\-य\-त्वाच्च}}{\lemma{\skt{तदञ्जनत्वात्}}\Afootnote{A N H; \skt{तद्व्यञ्जनत्वात्}\ L}}}{\lemma{\skt{तन्मयत्वाच्च}}\Afootnote{A N H; \skt{तज्जन्मत्वाच्च}\ L}}\skt{ }\edtext{\skt{तद्गुण}}{\lemma{\skt{तद्गुण}}\Afootnote{N H; \skt{तद्गुणा}\ A; \skt{तत्गुणा}\ L}}\skt{ एव }\edtext{\skt{पुरुषो}}{\lemma{\skt{पुरुषो}}\Afootnote{N H; \skt{पुरुषा}\ A L}}\skt{ }\edtext{\skt{भवतीत्येके}}{\lemma{\skt{भवतीत्य्}}\Afootnote{N H; \skt{भवन्तीत्य्}\ A L}}\skt{ भाषन्ते~।}\pend

\pstart\versepar \ledsidenote{\passref{3.1.11}}\edtext{\skt{वैदिके तु}}{\lemma{\skt{वैदिके तु}}\Afootnote{N L\textsuperscript{pc}; \skt{वैद्यके तु}\ A H; \skt{वैतदिके तु}\ L\textsuperscript{ac}}}\skt{~॥}\footnoteA{We read \emph{vaidike} with witness N, against H and the vulgate, because \emph{lectio dificilior potior}; see further discussion in the translation.}\skt{ }\pend
\pstart\versepar \skt{}\edtext{\skt{स्वभावो}}{\lemma{\skt{स्वभावो}}\Afootnote{N H; \skt{स्वभावम्}\ A L}}\skt{ }\edtext{\skt{नियतिः}}{\lemma{\skt{नियतिः}}\Afootnote{N H; \skt{ईश्वरं कालं यदृच्छां नियतिं}\ A; \skt{ईश्वरः कालो यद्रच्छा नियतिस्}\ L}}\skt{ }\edtext{\skt{कालः}}{\lemma{\skt{कालः}}\Afootnote{N H; \skt{तथा}\ A L}}\skt{ }\edtext{\edtext{\skt{परि\-णा\-म\-स्त\-थे\-श्वरः}}{\lemma{\skt{परिणामस्}}\Afootnote{N H; \skt{परिणामं}\ A; \skt{परिणां}\ L}}}{\lemma{\skt{तथेश्वरः}}\Afootnote{N H; \emph{om.}\ A L}}\skt{~। }\pend
\pstart\versepar \skt{}\edtext{\skt{यदृच्छेति च}}{\lemma{\skt{यदृच्छेति च}}\Afootnote{N; \skt{च}\ A; \skt{यदृच्छति च}\ H; \skt{इति प्}\ L}}\skt{ }\edtext{\skt{मन्यन्ते}}{\lemma{\skt{मन्यन्ते}}\Afootnote{A N H; \skt{राहुः}\ L}}\skt{ }\edtext{\skt{प्रकृ\-ति\-म्पृ\-थु\-द\-र्शिणः}}{\lemma{\skt{पृथुदर्शिणः}}\Afootnote{N H; \skt{पृथुदर्शिनः}\ A L}}\skt{~॥}\pend

\pstart\versepar \ledsidenote{\passref{3.1.12}}\edtext{\skt{ततो}}{\lemma{\skt{ततो}}\Afootnote{N H; \skt{तन्}\ A; \skt{तमो}\ L}}\skt{ }\edtext{\skt{जातानि}}{\lemma{\skt{जातानि}}\Afootnote{N H; \skt{मयान्य्}\ A L}}\skt{ }\edtext{\skt{भूतानि}}{\lemma{\skt{भूतानि}}\Afootnote{N H L; \skt{एव भूतानि}\ A}}\skt{ }\edtext{\skt{तद्गुणान्येव}}{\lemma{\skt{तद्गुणान्य्}}\Afootnote{A N H; \skt{तत्\_गुणान्य्}\ L}}\skt{ }\edtext{\skt{निर्द्दिशेत्}}{\lemma{\skt{निर्द्दिशेत्}}\Afootnote{N H; \skt{चादिशेत्}\ A; \skt{लक्षयेत्}\ L}}\skt{~।}\footnoteA{According to Ḍalhaṇa (1938: 341), Gayadāsa also read \emph{tato jātāni}, with the Nepalese version and against the vulgate.}\skt{ }\pend
\pstart\versepar \skt{}\edtext{\edtext{\skt{तेभ्यस्तल्लक्षणः}}{\lemma{\skt{तेभ्यस्}}\Afootnote{N H; \skt{तैश्च}\ A L}}}{\lemma{\skt{तल्लक्षणः}}\Afootnote{A H; \skt{तल्लक्षणोः}\ N; \skt{तल्लक्षणाः}\ L}}\skt{ }\edtext{\skt{कृत्स्नो}}{\lemma{\skt{कृत्स्नो}}\Afootnote{A N H; \skt{कृस्ना}\ L}}\skt{ }\edtext{\skt{भूतग्रामः}}{\lemma{\skt{भूतग्रामः}}\Afootnote{A H; \skt{भूताग्राम}\ N\textsuperscript{ac}; \skt{भूतग्राम}\ N\textsuperscript{pc}; \skt{भूतग्रामाः}\ L}}\skt{ }\edtext{\skt{प्रजायते}}{\lemma{\skt{प्रजायते}}\Afootnote{N H; \skt{व्यजन्यत}\ A; \skt{प्रजज्ञिरे}\ L}}\skt{~॥}\pend

\pstart\versepar \ledsidenote{\passref{3.1.13}}\edtext{\skt{तस्यो\-प\-यो\-गोऽ\-भि\-हि\-त\-श्चि\-कि\-त्सा\-म्प्रति}}{\lemma{\skt{ऽभिहितश्}}\Afootnote{A; \skt{भिहितञ्}\ N; \skt{भिहितश्}\ H L}}\skt{ सर्व्वदा~। }\pend
\pstart\versepar \skt{}\edtext{\skt{भूतेभ्योऽपि}}{\lemma{\skt{ऽपि}}\Afootnote{N H; \skt{हि}\ A; \skt{हि प}\ L}}\skt{ }\edtext{\skt{परन्तस्मान्नास्ति}}{\lemma{\skt{परन्तस्मान्}}\Afootnote{H; \skt{परं यस्मान्}\ A; \skt{परेत्तस्मान्}\ N; \skt{रा चिन्ता यस्मात्\_}\ L}}\skt{ }\edtext{\skt{चिन्ता}}{\lemma{\skt{चिन्ता}}\Afootnote{A N H; \emph{om.}\ L}}\skt{ चिकित्सिते~॥}\pend

\pstart \ledsidenote{\passref{3.1.14}}\edtext{\skt{यतोऽभिहितं}}{\lemma{\skt{ऽभिहितं}}\Afootnote{A H; \skt{भिहितन्}\ N; \skt{हितन्}\ L}}\skt{ तत्स\-म्भ\-व\-द्र\-व्य\-स\-मूहो }\edtext{\edtext{\skt{भूतादिरुक्त}}{\lemma{\skt{भूतादिर्}}\Afootnote{A N H L\textsuperscript{pc}; \skt{भूतार्}\ L\textsuperscript{ac}}}}{\lemma{\skt{उक्त}}\Afootnote{N H; \skt{उक्तः”}\ A; \emph{om.}\ L}}\skt{ }\edtext{\skt{इति}}{\lemma{\skt{इति}}\Afootnote{N H L; \emph{om.}\ A}}\skt{~।}\footnoteA{The passage referred to is Suśrutasaṃhitā, sūtrasthāna 1.38.}\skt{ }\edtext{\skt{भौतिकानि}}{\lemma{\skt{भौतिकानि}}\Afootnote{N H; \skt{(सू.१); भौतिकानि}\ A; \skt{इमानि भौतिकानि}\ L}}\skt{ }\edtext{\skt{चेन्द्रि\-या\-ण्या\-यु\-र्व्वेदे}}{\lemma{\skt{आयुर्व्वेदे}}\Afootnote{A N H; \skt{आयुर्वेदद्विदो}\ L}}\skt{ }\edtext{\skt{वर्ण्ण्यन्ते}}{\lemma{\skt{वर्ण्ण्यन्ते}}\Afootnote{A N H; \skt{उबदन्ति}\ L}}\skt{ }\edtext{\skt{तथैवेन्द्रियार्थाः}}{\lemma{\skt{तथैवेन्द्रियार्थाः}}\Afootnote{N H; \skt{तथेन्द्रियार्थाः}\ A; \skt{तथा~। इन्द्रियार्त्थांश्च}\ L}}\skt{~॥}\pend

\pstart\versepar \ledsidenote{\passref{3.1.15}}\edtext{\skt{भवति}}{\lemma{\skt{भवति}}\Afootnote{A N H; \emph{om.}\ L}}\skt{ }\edtext{\skt{चात्र}}{\lemma{\skt{चात्र}}\Afootnote{A N H; \skt{भ}\ L}}\skt{~॥ }\pend
\pstart\versepar \skt{}\edtext{\edtext{\skt{इन्द्रि\-यै\-रि\-न्द्रि\-यार्थं}}{\lemma{\skt{इन्द्रियैर्}}\Afootnote{N H; \skt{इन्द्रि}\ A; \skt{क्~। इन्द्रि}\ L}}\skt{ तु}}{\lemma{\skt{इन्द्रियार्थं तु}}\Afootnote{N\textsuperscript{pc} H; \skt{येणेन्दियार्थं तु}\ A; \skt{इयन्द्रियार्थन्तु}\ N\textsuperscript{ac}; \skt{येणेन्दियर्त्थं तु}\ L}}\skt{ स्वं स्वं गृह्णाति मानवः~। }\pend
\pstart\versepar \skt{नियतं }\edtext{\edtext{\skt{तुल्य\-यो\-नि\-त्वा\-न्ना\-न्ये\-ना\-न्य\-मिति}}{\lemma{\skt{तुल्ययोनित्वान्}}\Afootnote{A N H; \skt{तुल्ययोनित्वात्\_}\ L}}}{\lemma{\skt{नान्येनान्यम्}}\Afootnote{A N L; \skt{नाम्येनान्यम्}\ H}}\skt{ स्थितिः~॥}\pend

\pstart \ledsidenote{\passref{3.1.16}}\edtext{\skt{न चायु\-र्व्वे\-द\-शा\-स्त्रे\-षू\-प\-दि\-श्यन्ते}}{\lemma{\skt{न चायु\appbrk{}र्व्वे\appbrk{}द\appbrk{}शा\appbrk{}स्त्रे\appbrk{}षू\appbrk{}प\appbrk{}दि\appbrk{}श्यन्ते}}\Afootnote{A N H L\textsuperscript{pc}; \skt{न चायु\appbrk{}र्वे\appbrk{}द\appbrk{}शा\appbrk{}स्त्रे\appbrk{}षू\appbrk{}प\appbrk{}दि\appbrk{}श्यत्त्न्ते}\ L\textsuperscript{ac}}}\skt{ }\edtext{\skt{सर्व्वगताः}}{\lemma{\skt{सर्व्वगताः}}\Afootnote{A L; \skt{सर्वगता}\ N H}}\skt{ }\edtext{\skt{क्षेत्रज्ञा}}{\lemma{\skt{क्षेत्रज्ञा}}\Afootnote{A L; \skt{क्षेत्रज्ञ}\ N H}}\skt{ नित्याश्च~। }\edtext{\skt{असर्व्वगतेषु च}}{\lemma{\skt{असर्व्वगतेषु च}}\Afootnote{A N H; \skt{असर्वगतेषु}\ L}}\skt{ क्षेत्रज्ञेषु }\edtext{\skt{नित्येषु च}}{\lemma{\skt{नित्येषु च}}\Afootnote{\skt{नित्य}\ A; \skt{अनित्येषु च}\ N H; \skt{अनित्येषु}\ L}}\skt{ }\edtext{\edtext{\skt{पुरु\-षा\-ख्या\-प\-का\-न्हे\-तू\-नु\-दा\-ह\-रन्ति}}{\lemma{\skt{पुरु\appbrk{}षा\appbrk{}ख्या\appbrk{}प\appbrk{}कान्}}\Afootnote{H; \skt{पुरु\appbrk{}ष\appbrk{}ख्या\appbrk{}प\appbrk{}कान्}\ A; \skt{पुरुषख्यापकां}\ N; \skt{पुरु\appbrk{}ष\appbrk{}सं\appbrk{}ख्या\appbrk{}ति\appbrk{}प्र\appbrk{}का\appbrk{}रान्}\ L}}}{\lemma{\skt{हेतू\appbrk{}नु\appbrk{}दा\appbrk{}ह\appbrk{}रन्ति}}\Afootnote{A N H; \skt{हेतू\appbrk{}नु\appbrk{}दा\appbrk{}ह\appbrk{}रन्ती}\ L}}\skt{ }\edtext{\skt{आयु\-र्व्वे\-द\-सि\-द्धा\-न्तात्।}\footnoteA{We emend against the reading \emph{kṣetrajñeṣv anityeṣu} of witnesses N and H, for the logic of the passage and following the vulgate. The character \emph{ṣva} is clear in both N and H, but a copying error from \emph{-ṣu } is not inconceivable.}\skt{ असर्व्वगताः}}{\lemma{\skt{आयु\appbrk{}र्व्वे\appbrk{}द\appbrk{}सि\appbrk{}द्धा\appbrk{}न्तात्। असर्व्वगताः}}\Afootnote{\skt{आयु\appbrk{}र्वे\appbrk{}द\appbrk{}शा\appbrk{}स्त्रे\appbrk{}ष्व\appbrk{}स\appbrk{}र्व\appbrk{}गताः}\ A; \skt{आयु\appbrk{}र्वे\appbrk{}द\appbrk{}सि\appbrk{}द्धा\appbrk{}न्ता\appbrk{}द\appbrk{}स\appbrk{}र्व\appbrk{}गता}\ N H; \skt{आयु\appbrk{}र्वे\appbrk{}द\appbrk{}सि\appbrk{}द्धान्ताः असर्वगता}\ L}}\skt{ क्षेत्रज्ञा }\edtext{\skt{नित्याश्च}}{\lemma{\skt{नित्याश्च}}\Afootnote{A H; \skt{ऽनित्याश्च}\ N L}}\skt{ }\edtext{\skt{तिर्य\-ग्यो\-नि\-मा\-नु\-ष\-दे\-वेषु च}}{\lemma{\skt{तिर्य\appbrk{}ग्यो\appbrk{}नि\appbrk{}मा\appbrk{}नु\appbrk{}ष\appbrk{}दे\appbrk{}वेषु च}}\Afootnote{N H; \skt{तिर्य\appbrk{}ग्यो\appbrk{}नि\appbrk{}मा\appbrk{}नु\appbrk{}ष\appbrk{}दे\appbrk{}वेषु}\ A L}}\skt{ सञ्चरन्ति }\edtext{\skt{धर्म्मा\-ध\-र्म्म\-नि\-मित्तं~।}\footnoteA{The \emph{ca} after \emph{tiryagyonimānuṣadeveṣu} unnecessarily conjoins the members of the dvandva. While rare, this usage has been observed: Wackernagel 1896--1964, II.1, para 70.}\skt{ त}}{\lemma{\skt{धर्म्मा\appbrk{}ध\appbrk{}र्म्म\appbrk{}नि\appbrk{}मित्तं~। त}}\Afootnote{A N; \skt{धर्म्मा\appbrk{}ध\appbrk{}र्म्म\appbrk{}नि\appbrk{}मित्तं तु}\ H; \skt{धर्म्मा\appbrk{}ध\appbrk{}र्म\appbrk{}नि\appbrk{}मि\appbrk{}त्तम्\_}\ L}}\skt{ }\edtext{\edtext{\skt{एतेऽ\-नु\-मा\-न\-ग्रा\-ह्या\-श्चराः}}{\lemma{\skt{ऽनु\appbrk{}मा\appbrk{}न\appbrk{}ग्रा\appbrk{}ह्याश्}}\Afootnote{A; \skt{नुमानग्राह्याश्}\ N H L}}}{\lemma{\skt{चराः}}\Afootnote{N H; \emph{om.}\ A L}}\skt{ }\edtext{\edtext{\skt{पर\-म\-सू\-क्ष्मा\-श्चे\-त\-ना\-वन्तः}}{\lemma{\skt{परमसूक्ष्माश्}}\Afootnote{N H L; \skt{परमसूक्ष्मा}\ A}}}{\lemma{\skt{चेतनावन्तः}}\Afootnote{N\textsuperscript{pc} H L; \skt{श्चेतनावन्तः}\ A; \skt{चेति नावन्तः}\ N\textsuperscript{ac}}}\skt{ }\edtext{\skt{शाश्वताः}}{\lemma{\skt{शाश्वताः}}\Afootnote{N H; \skt{शाश्वता}\ A; \skt{च}\ L}}\skt{ }\edtext{\skt{शुक्रार्त्तवयोः}}{\lemma{\skt{शुक्रार्त्तवयोः}}\Afootnote{H; \skt{लोहितरेतसोः}\ A; \skt{शुक्रार्त्तवयोस्स}\ N; \skt{शोणितनेतसां}\ L}}\skt{ }\edtext{\edtext{\skt{सन्नि\-पा\-ते\-ष्व\-भि\-व्य\-ज्यन्ते}}{\lemma{\skt{सन्निपातेष्व्}}\Afootnote{A N H; \skt{सन्निपादेषु}\ L}}}{\lemma{\skt{अभिव्यज्यन्ते}}\Afootnote{A N H; \skt{अभिव्यञ्ज्यन्ते}\ L}}\skt{~। }\edtext{\edtext{\skt{यतोऽ\-भि\-हि\-त\-म्प\-ञ्च\-म\-हा\-भू\-त\-श\-री\-रि\-स\-म\-वायः}}{\lemma{\skt{ऽभिहितम्}}\Afootnote{A H; \skt{भिहितम्}\ N L}}}{\lemma{\skt{पञ्च\appbrk{}म\appbrk{}हा\appbrk{}भू\appbrk{}त\appbrk{}श\appbrk{}री\appbrk{}रि\appbrk{}स\appbrk{}म\appbrk{}वायः}}\Afootnote{A; \skt{पञ्चमहाभूतः शरी\appbrk{}रि\appbrk{}स\appbrk{}म\appbrk{}वायः}\ N; \skt{पञ्चमहाभूतः शरी\appbrk{}र\appbrk{}स\appbrk{}म\appbrk{}वायः}\ H; \skt{पञ्च\appbrk{}म\appbrk{}हा\appbrk{}भू\appbrk{}त\appbrk{}श\appbrk{}री\appbrk{}र\appbrk{}स\appbrk{}म\appbrk{}वायः}\ L}}\skt{ }\edtext{\skt{पुरुष}}{\lemma{\skt{पुरुष}}\Afootnote{N H L; \skt{पुरुषः’}\ A}}\skt{ }\edtext{\skt{इति~।}\footnoteA{We emend to \emph{-śarīri-} against N and H, for sense and because this is a quotation from Suśrutasaṃhitā 1.1.22.}\skt{ स}}{\lemma{\skt{इति~। स}}\Afootnote{N H L; \skt{(सू.अ.१) इति; स}\ A}}\skt{ }\edtext{\skt{खल्वेष}}{\lemma{\skt{खल्वेष}}\Afootnote{N H L; \skt{एष}\ A}}\skt{ }\edtext{\edtext{\edtext{\skt{कर्म्म\-पु\-रु\-ष\-श्चि\-कि\-त्सा\-या\-म\-धि\-कृतः}}{\lemma{\skt{कर्म्मपुरुषश्}}\Afootnote{A H L; \skt{कर्मपुरुष}\ N}}}{\lemma{\skt{चिकित्सायाम्}}\Afootnote{N H L; \skt{चिकित्सा}\ A}}}{\lemma{\skt{अधिकृतः}}\Afootnote{N H L; \skt{धिकृतः}\ A}}\skt{~।}\pend

\pstart \ledsidenote{\passref{3.1.17}}\skt{तस्य }\edtext{\skt{खलु}}{\lemma{\skt{खलु}}\Afootnote{N H; \emph{om.}\ A L}}\skt{ }\edtext{\skt{सुख\-दुः\-खे\-च्छा\-द्वे\-ष\-प्र\-यत्नः}}{\lemma{\skt{सुख\appbrk{}दुः\appbrk{}खे\appbrk{}च्छा\appbrk{}द्वे\appbrk{}ष\appbrk{}प्र\appbrk{}यत्नः}}\Afootnote{N H; \skt{सुखदुःखे इच्छाद्वेषौ प्रयत्नः}\ A; \skt{सुख\appbrk{}दुः\appbrk{}ख\appbrk{}मि\appbrk{}च्छा\appbrk{}द्वेषौ प्रयत्नः}\ L}}\skt{ प्राणापानौ }\edtext{\skt{निमेषोन्मेषौ}}{\lemma{\skt{निमेषोन्मेषौ}}\Afootnote{N H; \skt{उन्मेषनिमेषौ}\ A L}}\skt{ बुद्धिर्म्मनः संकल्पो }\edtext{\edtext{\skt{विचा\-र\-णा\-नु\-स्मृ\-ति\-ज्ञा\-न\-म\-ध्य\-व\-सायो}}{\lemma{\skt{विचा\appbrk{}र\appbrk{}णा\appbrk{}नु\appbrk{}स्मृ\appbrk{}ति\appbrk{}ज्ञा\appbrk{}नम्}}\Afootnote{N\textsuperscript{pc} H; \skt{विचारणा स्मृतिर्विज्ञानम्}\ A L; \skt{विचा\appbrk{}रि\appbrk{}णा\appbrk{}नु\appbrk{}स्मृ\appbrk{}ति\appbrk{}ज्ञा\appbrk{}नम्}\ N\textsuperscript{ac}}}}{\lemma{\skt{अध्यवसायो}}\Afootnote{A N H; \skt{द्ध्यवसायो}\ L}}\skt{ विष\-यो\-प\-ल\-ब्धिश्च }\edtext{\skt{गुणाः}}{\lemma{\skt{गुणाः}}\Afootnote{A N H; \skt{तु ध गुणाः}\ L\textsuperscript{ac}; \skt{तु गुणाः}\ L\textsuperscript{pc}}}\skt{~।}\pend

\pstart \ledsidenote{\passref{3.1.18}}\skt{सात्त्विकास्तु आनृशंस्यं सम्वि\-भा\-ग\-रु\-चिता }\edtext{\skt{तितिक्षा}}{\lemma{\skt{तितिक्षा}}\Afootnote{A H L; \skt{तिक्षा}\ N}}\skt{ }\edtext{\edtext{\skt{सत्यन्धर्म}}{\lemma{\skt{सत्यन्}}\Afootnote{A N H; \skt{सत्य}\ L}}}{\lemma{\skt{धर्म}}\Afootnote{A N H; \skt{धर्माव्}\ L}}\skt{ आस्तिक्यं ज्ञानं }\edtext{\edtext{\skt{बुद्धिर्मेधा}}{\lemma{\skt{बुद्धिर्}}\Afootnote{A N H; \skt{विज्ञानं बुद्धिर्}\ L}}}{\lemma{\skt{मेधा}}\Afootnote{A N H; \skt{म्मेवा}\ L}}\skt{ }\edtext{\edtext{\skt{धृति\-र\-न\-भि\-ष्व\-ङ्गता}}{\lemma{\skt{धृतिर्}}\Afootnote{N; \skt{स्मृतिर्धृतिर्}\ A; \skt{स्मृतिर्}\ H; \skt{धृति}\ L}}\skt{ च}}{\lemma{\skt{अनभिष्वङ्गता च}}\Afootnote{H; \skt{अनभिषङ्गश्च}\ A; \skt{अनभिष्ठङ्गता च}\ N; \skt{नभिषगश्चेति}\ L}}\skt{~॥ राजसास्तु }\edtext{\skt{दुःख\-ब\-हु\-ल\-ता\-प्र\-दानं}}{\lemma{\skt{दुःख\appbrk{}ब\appbrk{}हु\appbrk{}ल\appbrk{}ता\appbrk{}प्र\appbrk{}दानं}}\Afootnote{N H; \skt{दुःख\appbrk{}ब\appbrk{}हु\appbrk{}ल\appbrk{}ता\appbrk{}ट\appbrk{}न\appbrk{}शी\appbrk{}ल\appbrk{}ता\appbrk{}धृ\appbrk{}तिर्}\ A; \skt{दुःख\appbrk{}ब\appbrk{}हु\appbrk{}ल\appbrk{}ता\appbrk{}ट\appbrk{}न\appbrk{}त\appbrk{}र\appbrk{}दुश्}\ L}}\skt{ }\edtext{\edtext{\edtext{\skt{दुःशी\-ल\-त्व\-म\-का\-रु\-ण्य\-मा\-नृ\-ति\-क\-त्व\-म\-हं\-कारो}}{\lemma{\skt{दुःशीलत्वम्}}\Afootnote{N H; \skt{अहङ्}\ A; \skt{शीलत्वम्}\ L}}}{\lemma{\skt{अकारुण्यम्}}\Afootnote{N H L; \skt{कार}\ A}}}{\lemma{\skt{अहंकारो}}\Afootnote{H; \skt{अकारुण्यं}\ A; \skt{अहंकरो}\ N; \skt{अहंकारो}\ L}}\skt{ }\edtext{\skt{दम्भो}}{\lemma{\skt{दम्भो}}\Afootnote{A N H; \skt{सम्\_भ}\ L}}\skt{ }\edtext{\skt{मानो}}{\lemma{\skt{मानो}}\Afootnote{A N H; \skt{मान}\ L}}\skt{ }\edtext{\skt{हर्षः}}{\lemma{\skt{हर्षः}}\Afootnote{A N H; \skt{हर्ष}\ L}}\skt{ }\edtext{\edtext{\skt{कामक्रोधश्चेति}}{\lemma{\skt{कामक्रोधश्}}\Afootnote{N H; \skt{कामः क्रोधश्}\ A; \skt{क्रोध}\ L}}}{\lemma{\skt{चेति}}\Afootnote{H; \skt{च}\ A; \skt{चेति~॥ ना}\ N; \skt{कामाश्च भवन्ति}\ L}}\skt{~॥ }\edtext{\skt{तामसास्तु}}{\lemma{\skt{तामसास्तु}}\Afootnote{A H L; \skt{मसास्तु}\ N}}\skt{ }\edtext{\skt{विषादित्वं}}{\lemma{\skt{विषादित्वं}}\Afootnote{A H L; \skt{विशादिन्वन्}\ N}}\skt{ }\edtext{\skt{नास्ति\-क्य\-म\-ध\-र्म्म\-शी\-लता}}{\lemma{\skt{अधर्म्मशीलता}}\Afootnote{A N H; \skt{धर्म्मशीलतत्}\ L}}\skt{ }\edtext{\edtext{\skt{बुद्धि\-नि\-रो\-धोऽ\-ज्ञानं}}{\lemma{\skt{बुद्धिनिरोधो}}\Afootnote{N H L; \skt{बुद्धेर्निरोधो}\ A}}}{\lemma{\skt{ऽज्ञानं}}\Afootnote{A H; \skt{ज्ञानं}\ N L}}\skt{ }\edtext{\skt{दुर्मे\-ध\-स्त्व\-म\-क\-र्म्म\-शी\-लता}}{\lemma{\skt{दुर्मेधस्त्वम्}}\Afootnote{A N H; \skt{दुर्म्मेधत्वम्}\ L}}\skt{ निद्रालुत्वञ्चेति~॥}\pend

\pstart \ledsidenote{\passref{3.1.19}}\edtext{\edtext{\skt{आन्तरिक्षास्तु}}{\lemma{\skt{आन्तरिक्षास्}}\Afootnote{A; \skt{अन्तरिक्षास्}\ N H; \skt{आन्तरिक्षा}\ L}}}{\lemma{\skt{तु}}\Afootnote{N H; \emph{om.}\ A; \skt{स्तु}\ L}}\skt{ }\edtext{\skt{शब्दः}}{\lemma{\skt{शब्दः}}\Afootnote{A N L; \skt{शब्द}\ H}}\skt{ }\edtext{\skt{शब्देन्द्रियं}}{\lemma{\skt{शब्देन्द्रियं}}\Afootnote{A H L; \skt{सर्वेन्द्रियं}\ N}}\skt{ }\edtext{\skt{सर्व्व\-च्छि\-द्र\-स\-मूहो}}{\lemma{\skt{सर्व्व\appbrk{}च्छि\appbrk{}द्र\appbrk{}स\appbrk{}मूहो}}\Afootnote{A N H; \skt{सर्व\appbrk{}जि\appbrk{}द्र\appbrk{}स\appbrk{}मूहे र}\ L}}\skt{ विविक्तता च~।}\footnoteA{We emend \emph{antar-} to \emph{āntar-}, following the vulgate reading for sense. And we separate \emph{śabdaḥ} for the same reason.}\skt{ वायव्यास्तु }\edtext{\skt{स्पर्शः}}{\lemma{\skt{स्पर्शः}}\Afootnote{A N\textsuperscript{pc} H; \skt{स्पषर्शः}\ N\textsuperscript{ac}; \skt{स्पर्श}\ L}}\skt{ }\edtext{\skt{स्पर्शेन्द्रियं}}{\lemma{\skt{स्पर्शेन्द्रियं}}\Afootnote{A H L; \skt{स्पर्शेन्द्रिय}\ N}}\skt{ सर्व्व\-चे\-ष्टा\-स\-मूहः }\edtext{\skt{सर्व्व\-श\-री\-र\-स्प\-न्दनं}}{\lemma{\skt{सर्व्व\appbrk{}श\appbrk{}री\appbrk{}र\appbrk{}स्प\appbrk{}न्दनं}}\Afootnote{A H L; \skt{सर्व\appbrk{}श\appbrk{}री\appbrk{}र\appbrk{}स्प\appbrk{}न्दन}\ N}}\skt{ लघुता च~। }\edtext{\skt{तैजसास्तु}}{\lemma{\skt{तैजसास्तु}}\Afootnote{A H L; \skt{तेजसास्तु}\ N}}\skt{ रूपं }\edtext{\skt{रूपेन्द्रियं}}{\lemma{\skt{रूपेन्द्रियं}}\Afootnote{A N H; \skt{रूपेन्द्रियम्\_ सं}\ L}}\skt{ }\edtext{\skt{वर्ण्णः}}{\lemma{\skt{वर्ण्णः}}\Afootnote{A N; \skt{वर्ण्ण}\ H L}}\skt{ }\edtext{\skt{सन्तापो}}{\lemma{\skt{सन्तापो}}\Afootnote{A N H; \skt{समूहो}\ L}}\skt{ }\edtext{\skt{भ्राजिष्णुता}}{\lemma{\skt{भ्राजिष्णुता}}\Afootnote{A H; \skt{भ्राजिष्ठुता}\ N; \skt{द्राजिष्णुता}\ L}}\skt{ }\edtext{\edtext{\skt{पक्ति\-र\-म\-र्ष\-स्तै\-क्ष्ण्यञ्च}}{\lemma{\skt{अमर्षस्}}\Afootnote{A N\textsuperscript{pc} H L; \skt{अमर्ष}\ N\textsuperscript{ac}}}}{\lemma{\skt{तैक्ष्ण्यञ्च}}\Afootnote{N\textsuperscript{pc} H; \skt{तैक्ष्ण्यं शौर्यं च}\ A; \skt{स्तैस्तैक्ष्ण्यैक्ष्ण्यञ्च}\ N\textsuperscript{ac}; \skt{तैक्ष्ण्यम्\_ पीव्यञ्च}\ L}}\skt{~। आप्यास्तु }\edtext{\skt{रसो}}{\lemma{\skt{रसो}}\Afootnote{A H L; \skt{र}\ N}}\skt{ }\edtext{\skt{च रसेन्द्रियं}}{\lemma{\skt{च रसेन्द्रियं}}\Afootnote{H; \skt{रसनेन्द्रियं}\ A L; \skt{सारसेन्द्रियं}\ N}}\skt{ सर्व्व\-द्र\-व\-स\-मूहो }\edtext{\skt{गुरुता}}{\lemma{\skt{गुरुता}}\Afootnote{A N H; \emph{om.}\ L}}\skt{ शैत्यं }\edtext{\skt{स्नेहो}}{\lemma{\skt{स्नेहो}}\Afootnote{A N H L\textsuperscript{pc}; \skt{स्नेहो}\ L\textsuperscript{ac}}}\skt{ }\edtext{\skt{रेतश्च}}{\lemma{\skt{रेतश्च}}\Afootnote{A N H L\textsuperscript{pc}; \skt{त्रेतश्च}\ L\textsuperscript{ac}}}\skt{~। पार्थिवास्तु गन्धो गन्धेन्द्रियं }\edtext{\skt{सर्व्व\-मू\-र्त्ति\-स\-मूहो}}{\lemma{\skt{सर्व्व\appbrk{}मू\appbrk{}र्त्ति\appbrk{}स\appbrk{}मूहो}}\Afootnote{H L; \skt{सर्व\appbrk{}मू\appbrk{}र्त\appbrk{}स\appbrk{}मूहो}\ A; \skt{सर्व\appbrk{}मू\appbrk{}र्र्ति\appbrk{}स\appbrk{}मूहो}\ N}}\skt{ }\edtext{\skt{गुरुता च}}{\lemma{\skt{गुरुता च}}\Afootnote{N H; \skt{गुरुता चेति}\ A L}}\skt{~।}\pend

\pstart \ledsidenote{\passref{3.1.20}}\skt{तत्र }\edtext{\skt{सत्त्व\-ब\-हु\-ल\-मा\-काशं}}{\lemma{\skt{आकाशं}}\Afootnote{A H L; \skt{आकाशां}\ N}}\skt{~। रजोबहुलो वायुः~। }\edtext{\edtext{\skt{सत्त्व\-र\-जो\-ब\-हु\-लोऽग्निः}}{\lemma{\skt{सत्त्व\appbrk{}र\appbrk{}जो\appbrk{}ब\appbrk{}हुलो}}\Afootnote{A H; \skt{सत्व\appbrk{}र\appbrk{}जो\appbrk{}ब\appbrk{}हुलो}\ N; \skt{सत्व\appbrk{}र\appbrk{}जो\appbrk{}ब\appbrk{}हुला}\ L}}}{\lemma{\skt{ऽग्निः}}\Afootnote{A H; \skt{ग्निः}\ N L}}\skt{~। }\edtext{\skt{सत्त्व\-त\-मो\-ब\-हुला}}{\lemma{\skt{सत्त्व\appbrk{}त\appbrk{}मो\appbrk{}ब\appbrk{}हुला}}\Afootnote{A L; \skt{सत्व\appbrk{}त\appbrk{}मो\appbrk{}ब\appbrk{}हुलो}\ N; \skt{सत्व\appbrk{}त\appbrk{}मो\appbrk{}व\appbrk{}हुल}\ H}}\skt{ आपः~। तमोबहुला }\edtext{\skt{पृथिवी}}{\lemma{\skt{पृथिवी}}\Afootnote{N H; \skt{पृथिवीति}\ A; \skt{पृथिविति}\ L}}\skt{~॥}\pend

\pstart\versepar \ledsidenote{\passref{3.1.21}}\edtext{\skt{भवति}}{\lemma{\skt{भवति}}\Afootnote{H L; \skt{श्लोकौ}\ A; \emph{om.}\ N}}\skt{ }\edtext{\skt{चात्र}}{\lemma{\skt{चात्र}}\Afootnote{A H L; \emph{om.}\ N}}\skt{~॥ }\pend
\pstart\versepar \skt{}\edtext{\skt{अन्योन्यानि प्रविष्टानि}}{\lemma{\skt{अन्योन्यानि प्रविष्टानि}}\Afootnote{H; \skt{भवतः -- अन्यो\appbrk{}न्या\appbrk{}नु\appbrk{}प्र\appbrk{}वि\appbrk{}ष्टानि}\ A; \skt{भवति चात्र~॥ अन्यो न्यानुप्रविष्टानि}\ N\textsuperscript{ac}; \skt{भवति चात्र~॥ अन्यो न्याप्रविष्टानि}\ N\textsuperscript{pc}; \skt{श्ल्लोकः}\ L}}\skt{ }\edtext{\edtext{\skt{सर्व्वाण्येतानि}}{\lemma{\skt{सर्व्वाण्य्}}\Afootnote{A N H; \emph{om.}\ L}}}{\lemma{\skt{एतानि}}\Afootnote{A N H; \emph{om.}\ L}}\skt{ }\edtext{\skt{निर्दिशेत्}}{\lemma{\skt{निर्दिशेत्}}\Afootnote{A H; \skt{निर्दिषेत्}\ N; \emph{om.}\ L}}\skt{~। }\pend
\pstart\versepar \skt{}\edtext{\skt{स्वे}}{\lemma{\skt{स्वे}}\Afootnote{A N H; \emph{om.}\ L}}\skt{ }\edtext{\skt{स्वे}}{\lemma{\skt{स्वे}}\Afootnote{A N H; \emph{om.}\ L}}\skt{ }\edtext{\skt{द्रव्येषु}}{\lemma{\skt{द्रव्येषु}}\Afootnote{N H; \skt{द्रव्ये तु}\ A; \emph{om.}\ L}}\skt{ }\edtext{\skt{सर्व्वेषां}}{\lemma{\skt{सर्व्वेषां}}\Afootnote{A N H; \emph{om.}\ L}}\skt{ }\edtext{\edtext{\skt{व्यक्त\-ल\-क्ष\-ण\-मि\-ष्यते}}{\lemma{\skt{व्यक्तलक्षणम्}}\Afootnote{H; \skt{व्यक्तं लक्षणम्}\ A N; \emph{om.}\ L}}}{\lemma{\skt{इष्यते}}\Afootnote{A N H; \emph{om.}\ L}}\skt{~॥}\pend

\pstart\versepar \ledsidenote{\passref{3.1.22}}\skt{अष्टौ }\edtext{\skt{प्रकृतयः}}{\lemma{\skt{प्रकृतयः}}\Afootnote{A N H; \skt{प्रकृतिय}\ L}}\skt{ }\edtext{\skt{प्रोक्ता}}{\lemma{\skt{प्रोक्ता}}\Afootnote{A N H; \skt{प्रोक्तः}\ L}}\skt{ }\edtext{\skt{विकाराश्चैव}}{\lemma{\skt{चैव}}\Afootnote{N H; \skt{षोड}\ A L}}\skt{ }\edtext{\skt{षोडश}}{\lemma{\skt{षोडश}}\Afootnote{H; \skt{शैव तु}\ A L; \skt{षोडशः}\ N}}\skt{~। }\pend
\pstart\versepar \skt{}\edtext{\skt{क्षेत्रज्ञश्च}}{\lemma{\skt{क्षेत्रज्ञश्च}}\Afootnote{A N; \skt{क्षेत्रज्ञाश्च}\ H; \skt{क्षेत्रज्ञस्तु}\ L}}\skt{ समासेन }\edtext{\skt{स्वत\-न्त्र\-प\-र\-त\-न्त्रत}}{\lemma{\skt{स्वत\appbrk{}न्त्र\appbrk{}प\appbrk{}र\appbrk{}त\appbrk{}न्त्रत}}\Afootnote{H; \skt{स्वत\appbrk{}न्त्र\appbrk{}प\appbrk{}र\appbrk{}त\appbrk{}न्त्रयोः}\ A; \skt{सुतन्त्र परतन्त्रत}\ N; \skt{पर\appbrk{}त\appbrk{}न्त्र\appbrk{}स्व\appbrk{}त\appbrk{}न्त्रतः}\ L}}\skt{ }\edtext{\skt{इति}}{\lemma{\skt{इति}}\Afootnote{N H; \emph{om.}\ A L}}\skt{~॥}\pend

\pstart \ledsidenote{\passref{3.1.trailer}}\edtext{\skt{शारीरस्य}}{\lemma{\skt{शारीरस्य}}\Afootnote{H; \skt{इति सुश्रु\appbrk{}त\appbrk{}सं\appbrk{}हि\appbrk{}तायां शारीरस्थाने सर्व\appbrk{}भू\appbrk{}त\appbrk{}चि\appbrk{}न्ता\appbrk{}शा\appbrk{}रीरं नाम}\ A; \skt{शा}\ N}}\skt{ }\edtext{\edtext{\skt{प्रथमोऽध्यायः}}{\lemma{\skt{प्रथमो}}\Afootnote{A H; \emph{om.}\ N}}}{\lemma{\skt{ऽध्यायः}}\Afootnote{A; \emph{om.}\ N; \skt{ऽध्यायः~॥ ओ}\ H}}\skt{~॥}\pend

\endnumbering
\beginnumbering
\eledsection*{Adhyāya 2}
\pstart \ledsidenote{\passref{3.2.1}}\edtext{\skt{अथातः}}{\lemma{\skt{अथातः}}\Afootnote{A N H; \emph{om.}\ L}}\skt{ }\edtext{\skt{शुक्र\-शो\-णि\-त\-वि\-शुद्धिं}}{\lemma{\skt{शुक्र\appbrk{}शो\appbrk{}णि\appbrk{}त\appbrk{}वि\appbrk{}शुद्धिं}}\Afootnote{N H; \skt{शुक्र\appbrk{}शो\appbrk{}णि\appbrk{}त\appbrk{}शुद्धि}\ A; \skt{सु}\ L}}\skt{ }\edtext{\skt{शारीरं}}{\lemma{\skt{शारीरं}}\Afootnote{A H; \emph{om.}\ N; \skt{श्रुतः}\ L}}\skt{ }\edtext{\skt{व्याख्यास्यामः}}{\lemma{\skt{व्याख्यास्यामः}}\Afootnote{A N H; \emph{om.}\ L}}\skt{~॥}\footnoteA{Witness N omits \emph{śārīraṃ}, which offers a clearer syntax. Note that the 1931 and 1938 editions of the vulgate text offer two different solutions to the awkward phrase.}\pend

\pstart \ledsidenote{\passref{3.2.3}}\edtext{\skt{वात\-पि\-त्त\-श्ले\-ष्म\-शो\-णि\-त\-कु\-ण\-प\-ग्र\-न्थि\-पू\-ति\-पू\-य\-क्षी\-ण\-मू\-त्र\-पु\-री\-ष\-रे\-तांसि}}{\lemma{\skt{वात\appbrk{}पि\appbrk{}त्त\appbrk{}श्ले\appbrk{}ष्म\appbrk{}शो\appbrk{}णि\appbrk{}त\appbrk{}कु\appbrk{}ण\appbrk{}प\appbrk{}ग्र\appbrk{}न्थि\appbrk{}पू\appbrk{}ति\appbrk{}पू\appbrk{}य\appbrk{}क्षी\appbrk{}ण\appbrk{}मू\appbrk{}त्र\appbrk{}पु\appbrk{}री\appbrk{}ष\appbrk{}रे\appbrk{}तांसि}}\Afootnote{N H; \skt{वात\appbrk{}पि\appbrk{}त्त\appbrk{}श्ले\appbrk{}ष्म\appbrk{}कु\appbrk{}णप ग्रन्थि\appbrk{}पू\appbrk{}ति\appbrk{}पू\appbrk{}य\appbrk{}क्षी\appbrk{}ण\appbrk{}मू\appbrk{}त्र\appbrk{}पु\appbrk{}री\appbrk{}ष\appbrk{}रे\appbrk{}तसः}\ A; \skt{वातपित्तश्लेष्म एकु\appbrk{}ण\appbrk{}प\appbrk{}ग्र\appbrk{}न्थि\appbrk{}पू\appbrk{}य\appbrk{}क्षी\appbrk{}ण\appbrk{}पु\appbrk{}री\appbrk{}ष\appbrk{}रे\appbrk{}तांसि}\ L}}\footnoteA{The 1931 and 1938 editions of the SS by Ācārya have different readings at this point. \emph{śonita} is included in the 1931 edition but omitted in the 1938 edition. There also seems to be some textual confusion between \emph{granthi} and \emph{gandhi} in the next passage.}\skt{ }\edtext{\edtext{\skt{प्रजो\-त्पा\-द\-ने\-ष्व\-स\-म\-र्थानि}}{\lemma{\skt{प्रजोत्पादनेष्व्}}\Afootnote{N H L; \skt{प्रजोत्पादने}\ A}}}{\lemma{\skt{असमर्थानि}}\Afootnote{H; \skt{न समर्था}\ A; \skt{अशगर्था}\ N; \skt{असमन्यानि}\ L}}\skt{ }\edtext{\skt{भवन्ति}}{\lemma{\skt{भवन्ति}}\Afootnote{A H L; \emph{om.}\ N}}\skt{~।}\pend

\pstart \ledsidenote{\passref{3.2.4}}\edtext{\skt{तेषु}}{\lemma{\skt{तेषु}}\Afootnote{A H; \emph{om.}\ N L}}\skt{ }\edtext{\skt{वात\-व\-र्ण्ण\-वे\-दनं}}{\lemma{\skt{वात\appbrk{}व\appbrk{}र्ण्ण\appbrk{}वे\appbrk{}दनं}}\Afootnote{A H; \emph{om.}\ N; \skt{तेषु दोष}\ L}}\skt{ }\edtext{\skt{वातेन}}{\lemma{\skt{वातेन}}\Afootnote{A H; \skt{तेन}\ N; \emph{om.}\ L}}\skt{~। }\edtext{\skt{पित्त\-व\-र्ण्ण\-वे\-दनं}}{\lemma{\skt{पित्त\appbrk{}व\appbrk{}र्ण्ण\appbrk{}वे\appbrk{}दनं}}\Afootnote{A H; \skt{पित्त\appbrk{}व\appbrk{}र्ण्ण\appbrk{}वे\appbrk{}द\appbrk{}नाम्}\ N; \skt{नि}\ L}}\skt{ }\edtext{\skt{पित्तेन}}{\lemma{\skt{पित्तेन}}\Afootnote{A H; \skt{पत्तेन}\ N; \skt{मित्तानि}\ L}}\skt{~। }\edtext{\skt{श्लेष्म\-व\-र्ण्ण\-वे\-दनं}}{\lemma{\skt{श्लेष्म\appbrk{}व\appbrk{}र्ण्ण\appbrk{}वे\appbrk{}दनं}}\Afootnote{A H; \skt{श्लेष्म\appbrk{}व\appbrk{}र्ण्ण\appbrk{}वे\appbrk{}दन}\ N; \skt{स्वर्णवेद}\ L}}\skt{ }\edtext{\skt{श्लेष्मणा}}{\lemma{\skt{श्लेष्मणा}}\Afootnote{A N H; \skt{ना}\ L}}\skt{~। }\edtext{\edtext{\skt{शोणि\-त\-व\-र्ण्ण\-म्पि\-त्त\-वे\-दनं}}{\lemma{\skt{शोणितवर्ण्णम्}}\Afootnote{H; \emph{om.}\ A L; \skt{शोणितवर्ण्ण}\ N}}}{\lemma{\skt{पित्तवेदनं}}\Afootnote{N H; \emph{om.}\ A; \skt{भिर्वि}\ L}}\skt{ }\edtext{\skt{शोणितेन}}{\lemma{\skt{शोणितेन}}\Afootnote{N H; \emph{om.}\ A; \skt{ज्ञेयानि}\ L}}\skt{~। }\edtext{\skt{कुणपञ्च}}{\lemma{\skt{कुणपञ्च}}\Afootnote{N H; \skt{कुण\appbrk{}प\appbrk{}ग\appbrk{}न्ध्य\appbrk{}नल्पं च}\ A; \skt{कुण\appbrk{}प\appbrk{}ग\appbrk{}न्ध्य\appbrk{}न\appbrk{}ल्पम्\_}\ L}}\skt{ रक्तेन~। ग्रन्थिभूतं श्लेष्मवाताभ्यां~। }\edtext{\skt{पूतिपूयनिभं}}{\lemma{\skt{पूतिपूयनिभं}}\Afootnote{A N H; \skt{पूतिपूयाभम्\_}\ L}}\skt{ }\edtext{\skt{पित्तशोणिताभ्यां}}{\lemma{\skt{पित्तशोणिताभ्यां}}\Afootnote{N H L; \skt{पित्तश्लेष्मभ्यां}\ A}}\skt{~। }\edtext{\skt{क्षीणशुक्रं}}{\lemma{\skt{क्षीणशुक्रं}}\Afootnote{H; \skt{क्षीणं प्रागुक्तं}\ A L; \skt{क्षीणशुक्र}\ N}}\skt{ }\edtext{\skt{पित्तमारुताभ्यां}}{\lemma{\skt{पित्तमारुताभ्यां}}\Afootnote{A N H; \emph{om.}\ L}}\skt{~। }\edtext{\skt{मूत्र\-पु\-री\-ष\-गन्धि}}{\lemma{\skt{मूत्र\appbrk{}पु\appbrk{}री\appbrk{}ष\appbrk{}गन्धि}}\Afootnote{A H; \skt{मूत्रा\appbrk{}पु\appbrk{}री\appbrk{}ष\appbrk{}गन्धि}\ N; \skt{पुरिषगन्धि}\ L}}\skt{ }\edtext{\skt{सन्निपातेन}}{\lemma{\skt{सन्निपातेन}}\Afootnote{N H; \skt{सन्निपातेने}\ A L}}\skt{~। }\edtext{\skt{तेषु}}{\lemma{\skt{तेषु}}\Afootnote{N H; \skt{ति~। ते}\ A L}}\skt{ }\edtext{\skt{कुण\-प\-ग्र\-न्थि\-पू\-ति\-पू\-य\-रे\-तसः}}{\lemma{\skt{कुण\appbrk{}प\appbrk{}ग्र\appbrk{}न्थि\appbrk{}पू\appbrk{}ति\appbrk{}पू\appbrk{}य\appbrk{}रे\appbrk{}तसः}}\Afootnote{H; \skt{षु कुण\appbrk{}प\appbrk{}ग्र\appbrk{}न्थि\appbrk{}पू\appbrk{}ति\appbrk{}पू\appbrk{}य\appbrk{}क्षी\appbrk{}ण\appbrk{}रे\appbrk{}तसः}\ A; \skt{कुल\appbrk{}ण\appbrk{}प\appbrk{}ग्र\appbrk{}न्थि\appbrk{}पू\appbrk{}ति\appbrk{}पू\appbrk{}र\appbrk{}य\appbrk{}रे\appbrk{}तसः}\ N\textsuperscript{ac}; \skt{कुण\appbrk{}ण\appbrk{}प\appbrk{}ग्र\appbrk{}न्थि\appbrk{}पू\appbrk{}ति\appbrk{}पू\appbrk{}र\appbrk{}य\appbrk{}रे\appbrk{}तसः}\ N\textsuperscript{pc}; \skt{षु कुण\appbrk{}प\appbrk{}ग\appbrk{}न्धि\appbrk{}पू\appbrk{}य\appbrk{}रे\appbrk{}तांसि}\ L}}\skt{ कृच्छ्रसाध्याः~। }\edtext{\edtext{\skt{मूत्र\-पु\-री\-ष\-रे\-त\-स\-स्त्व\-साध्या}}{\lemma{\skt{मूत्र\appbrk{}पु\appbrk{}री\appbrk{}ष\appbrk{}रे\appbrk{}त\appbrk{}सस्}}\Afootnote{A H; \skt{मूत्र\appbrk{}पु\appbrk{}री\appbrk{}ष\appbrk{}रेत…}\ N; \skt{नि पुरिषरेतस्}\ L}}}{\lemma{\skt{त्वसाध्या}}\Afootnote{A H; \emph{om.}\ N; \skt{त्वसाद्ध्यम्}\ L}}\skt{ }\edtext{\skt{इति}}{\lemma{\skt{इति}}\Afootnote{A H L; \emph{om.}\ N}}\skt{~।}\pend

\pstart \ledsidenote{\passref{3.2.5}}\edtext{\edtext{\skt{आर्त्तवमपि}}{\lemma{\skt{आर्त्तवम्}}\Afootnote{A H L; \emph{om.}\ N}}}{\lemma{\skt{अपि}}\Afootnote{A H; \emph{om.}\ N; \skt{इति}\ L}}\skt{ }\edtext{\skt{त्रिभिर्द्दोषैः}}{\lemma{\skt{त्रिभिर्}}\Afootnote{A H L; \skt{भिर्}\ N}}\skt{ }\edtext{\edtext{\skt{शोणि\-त\-च\-तु\-र्थै\-ह्पृ\-थ\-क्द्व\-न्द्वै\-स्त्रिभिः}}{\lemma{\skt{शोणि\appbrk{}त\appbrk{}च\appbrk{}तु\appbrk{}र्थै\appbrk{}ह्पृ\appbrk{}थ\appbrk{}क्द्व\appbrk{}न्द्वैस्}}\Afootnote{H; \skt{शोणितचतुर्थैः पृथग्द्वन्द्वैः}\ A; \skt{शोणितचतुर्थैः पृथक्द्वन्द्वैस्}\ N; \skt{शोणि\appbrk{}त\appbrk{}च\appbrk{}तु\appbrk{}र्त्थै\appbrk{}र्द्द्व\appbrk{}न्द्वैस्}\ L}}}{\lemma{\skt{त्रिभिः}}\Afootnote{N H; \emph{om.}\ A L}}\skt{ }\edtext{\edtext{\skt{सम\-स्तै\-श्चो\-प\-सृ\-ष्ट\-म\-बी\-ज\-न्भ\-वति}}{\lemma{\skt{चोपसृष्टम्}}\Afootnote{A N H; \skt{चोपसृष्ट}\ L}}}{\lemma{\skt{अबीजन्}}\Afootnote{\skt{अबीजं}\ A N; \skt{अवीजन्न}\ H; \skt{बीजम्\_ बीजम्}\ L}}\skt{~। तदपि }\edtext{\skt{दोष\-व\-र्ण्ण\-वे\-द\-ना\-भि\-र्व्वि\-ज्ञेयं}}{\lemma{\skt{दोष\appbrk{}व\appbrk{}र्ण्ण\appbrk{}वे\appbrk{}द\appbrk{}ना\appbrk{}भिर्}}\Afootnote{H L; \skt{दोष\appbrk{}व\appbrk{}र्ण\appbrk{}वे\appbrk{}द\appbrk{}ना\appbrk{}दि\appbrk{}भिर्}\ A; \skt{दोष\appbrk{}व\appbrk{}र्ण्ण\appbrk{}वे\appbrk{}द\appbrk{}ना\appbrk{}ति\appbrk{}भिर्}\ N}}\skt{~। तेषु }\edtext{\edtext{\skt{कुण\-प\-ग्र\-न्थि\-पू\-ति\-पू\-य\-प्र\-का\-श\-न्त\-द\-साध्यं}}{\lemma{\skt{कुण\appbrk{}प\appbrk{}ग्र\appbrk{}न्थि\appbrk{}पू\appbrk{}ति\appbrk{}पू\appbrk{}य\appbrk{}प्र\appbrk{}का\appbrk{}शन्}}\Afootnote{N H; \skt{कुण\appbrk{}प\appbrk{}ग्र\appbrk{}न्थि\appbrk{}पू\appbrk{}ति\appbrk{}पू\appbrk{}य\appbrk{}क्षी\appbrk{}ण\appbrk{}मू\appbrk{}त्र\appbrk{}पु\appbrk{}री\appbrk{}ष\appbrk{}प्र\appbrk{}का\appbrk{}शम्}\ A; \skt{कुण\appbrk{}प\appbrk{}ग्र\appbrk{}न्धि\appbrk{}पू\appbrk{}य\appbrk{}प्र\appbrk{}काशं}\ L}}}{\lemma{\skt{तदसाध्यं}}\Afootnote{H; \skt{असाध्यं}\ A; \skt{त}\ N; \skt{कृच्छ्रसाद्ध्यं}\ L}}\skt{ }\edtext{\edtext{\edtext{\skt{साध्यमन्यच्चेति}}{\lemma{\skt{साध्यम्}}\Afootnote{A H; \skt{दसाध्यं साध्यम्}\ N; \skt{पुरि\appbrk{}षा\appbrk{}र्त्त\appbrk{}व\appbrk{}म\appbrk{}साद्ध्यं}\ L}}}{\lemma{\skt{अन्यच्}}\Afootnote{A N\textsuperscript{ac} H; \skt{अन्यति}\ N\textsuperscript{pc}; \emph{om.}\ L}}}{\lemma{\skt{चेति}}\Afootnote{A N\textsuperscript{ac} H; \emph{om.}\ N\textsuperscript{pc} L}}\skt{~॥}\pend

\pstart\versepar \ledsidenote{\passref{3.2.6}}\edtext{\skt{भवति}}{\lemma{\skt{भवति}}\Afootnote{H; \skt{भवन्ति}\ A N; \emph{om.}\ L}}\skt{ }\edtext{\skt{चात्र}}{\lemma{\skt{चात्र}}\Afootnote{A N H; \skt{भका~।}\ L}}\skt{~॥ }\pend
\pstart\versepar \skt{}\edtext{\edtext{\skt{तेष्वा\-द्या\-ञ्शु\-क्र\-दो\-षां\-स्त्री\-न्स्ने\-ह\-स्वे\-दा\-दि\-भि\-र्ज्ज\-येत्}}{\lemma{\skt{आद्या\appbrk{}ञ्शु\appbrk{}क्र\appbrk{}दो\appbrk{}षांस्}}\Afootnote{A; \skt{आद्यां शुक्रदोषांस्}\ N; \skt{आद्याः शुक्रदोषास्}\ H; \skt{आद्या\appbrk{}ञ्च्छु\appbrk{}क्र\appbrk{}दो\appbrk{}षांस्}\ L}}}{\lemma{\skt{त्रीन्स्ने\appbrk{}ह\appbrk{}स्वे\appbrk{}दा\appbrk{}दि\appbrk{}भिर्}}\Afootnote{A; \skt{त्री स्नेह\appbrk{}स्वे\appbrk{}दा\appbrk{}दि\appbrk{}भिर्}\ N H; \skt{त्रीं स्नेह\appbrk{}न\appbrk{}स्वे\appbrk{}दा\appbrk{}दि\appbrk{}भिर्}\ L}}\skt{~। }\pend
\pstart\versepar \skt{}\edtext{\skt{क्रिया\-वि\-शे\-षै\-र्म्म\-ति\-मा\-न्त\-थै\-वो\-त्त\-र\-ब\-स्तिभिः}}{\lemma{\skt{तथै\appbrk{}वो\appbrk{}त्त\appbrk{}र\appbrk{}ब\appbrk{}स्तिभिः}}\Afootnote{N H L; \skt{तथा चोत्तरबस्तिभिः}\ A}}\skt{~॥}\pend

\pstart \ledsidenote{\passref{3.2.6.1}}\skt{तत्र वातात्मके रेतसि बिल्व\-वि\-दा\-री\-क्षी\-र\-मा\-स्था\-प\-नम्~। }\edtext{\skt{मधू\-क\-रा\-स्ना\-दे\-व\-दा\-रु\-स\-र\-ल\-वि\-प\-क्व\-न्तै\-ल\-न्त\-दु\-त्त\-र\-ब\-स्तिषु}}{\lemma{\skt{तदु\appbrk{}त्त\appbrk{}र\appbrk{}ब\appbrk{}स्तिषु}}\Afootnote{H; \skt{उत्तरबस्तिषु}\ N}}\skt{ विदध्यात्~॥ }\edtext{\skt{दाडि\-म\-मा\-तु\-लु\-ङ्ग\-फ\-ल\-सै\-न्ध\-व\-क्षा\-र\-व\-सु\-क\-व\-सि\-र\-सिद्धं}}{\lemma{\skt{दाडि\appbrk{}म\appbrk{}मा\appbrk{}तु\appbrk{}लु\appbrk{}ङ्ग\appbrk{}फ\appbrk{}ल\appbrk{}सै\appbrk{}न्ध\appbrk{}व\appbrk{}क्षा\appbrk{}र\appbrk{}व\appbrk{}सु\appbrk{}क\appbrk{}व\appbrk{}सि\appbrk{}र\appbrk{}सिद्धं}}\Afootnote{N; \skt{दाडि\appbrk{}म\appbrk{}मा\appbrk{}तु\appbrk{}ल\appbrk{}ङ्ग\appbrk{}फ\appbrk{}ल\appbrk{}सै\appbrk{}न्ध\appbrk{}व\appbrk{}क्षा\appbrk{}र\appbrk{}व\appbrk{}सु\appbrk{}क\appbrk{}व-\appbrk{}सि\appbrk{}र\appbrk{}सिद्धं}\ H}}\skt{ सर्प्पिः पाययेत्~॥}\pend

\pstart \ledsidenote{\passref{3.2.6.2}}\edtext{\skt{पित्तात्मके}}{\lemma{\skt{पित्तात्मके}}\Afootnote{H; \skt{पिड्रात्मक}\ N}}\skt{ रेतसि }\edtext{\skt{पय\-स्या\-श्री\-प\-र्ण्णी\-म\-धु\-क\-वि\-पक्वं}}{\lemma{\skt{पय\appbrk{}स्या\appbrk{}श्री\appbrk{}प\appbrk{}र्ण्णी\appbrk{}म\appbrk{}धु\appbrk{}क\appbrk{}वि\appbrk{}पक्वं}}\Afootnote{H; \skt{पय\appbrk{}स्या।\appbrk{}श्री\appbrk{}प\appbrk{}र्ण्णी\appbrk{}म\appbrk{}धु\appbrk{}क\appbrk{}वि\appbrk{}पर्कं}\ N}}\skt{ क्षीरमास्थापनं~। }\edtext{\skt{सर्ज्जधवकल्कं}}{\lemma{\skt{सर्ज्जधवकल्कं}}\Afootnote{H; \emph{om.}\ N}}\skt{ }\edtext{\skt{योन्या}}{\lemma{\skt{योन्या}}\Afootnote{H; \emph{om.}\ N}}\skt{ }\edtext{\skt{धारयेत्}}{\lemma{\skt{धारयेत्}}\Afootnote{H; \emph{om.}\ N}}\skt{~। }\edtext{\skt{मधुकविपक्वं}}{\lemma{\skt{मधुकविपक्वं}}\Afootnote{H; \skt{मधुकविपर्क}\ N}}\skt{ तैल\-म\-नु\-वा\-सनं~। तदे\-वो\-त्त\-र\-ब\-स्तिञ्च }\edtext{\skt{विदध्यात्}}{\lemma{\skt{विदध्यात्}}\Afootnote{N\textsuperscript{pc} H; \skt{विदद्याध्यात्}\ N\textsuperscript{ac}}}\skt{~। }\edtext{\skt{काण्डे\-क्षु\-श्व\-दं\-ष्त्रा\-गु\-डू\-ची\-म\-धु\-प\-र्ण्णी\-भृ\-ङ्ग\-प\-ञ्च\-मू\-ल\-सिद्धं}}{\lemma{\skt{काण्डे\appbrk{}क्षु\appbrk{}श्व\appbrk{}दं\appbrk{}ष्त्रा\appbrk{}गु\appbrk{}डू\appbrk{}ची\appbrk{}म\appbrk{}धु\appbrk{}प\appbrk{}र्ण्णी\appbrk{}भृ\appbrk{}ङ्ग\appbrk{}प\appbrk{}ञ्च\appbrk{}मू\appbrk{}ल\appbrk{}सिद्धं}}\Afootnote{\skt{कण्ठे\appbrk{}क्षु\appbrk{}श्व\appbrk{}दं\appbrk{}ष्ट्रा\appbrk{}गु\appbrk{}डू\appbrk{}ची\appbrk{}म\appbrk{}धु\appbrk{}प\appbrk{}र्ण्णी\appbrk{}भृ\appbrk{}ण\appbrk{}प\appbrk{}ञ्च\appbrk{}मू\appbrk{}ल\appbrk{}ले\appbrk{}सिद्धं}\ N; \skt{काण्डे\appbrk{}क्षु\appbrk{}श्व\appbrk{}दं\appbrk{}ष्त्रा\appbrk{}गु\appbrk{}डू\appbrk{}ची\appbrk{}म\appbrk{}धु\appbrk{}प\appbrk{}र्ण्णी\appbrk{}भृ\appbrk{}ण\appbrk{}प\appbrk{}ञ्च\appbrk{}मू\appbrk{}ल\appbrk{}सिद्धं}\ H}}\skt{ सर्प्पिः }\edtext{\skt{पाययेत्}}{\lemma{\skt{पाययेत्}}\Afootnote{H; \skt{पापयेत्}\ N}}\skt{~॥}\pend

\pstart \ledsidenote{\passref{3.2.6.3}}\skt{श्लेष्मात्मके रेतसि राज\-वृ\-क्ष\-क\-षा\-य\-मा\-स्था\-पनं~। }\edtext{\edtext{\skt{पिप्प\-ली\-वि\-ड\-ङ्ग\-म\-धु\-वि\-प\-क्व\-न्तै\-ल\-म\-नु\-वा\-स\-नन्}}{\lemma{\skt{पिप्प\appbrk{}ली\appbrk{}वि\appbrk{}ड\appbrk{}ङ्ग\appbrk{}म\appbrk{}धु\appbrk{}वि\appbrk{}प\appbrk{}क्वन्}}\Afootnote{H; \skt{पिप्प\appbrk{}ली\appbrk{}वि\appbrk{}ड\appbrk{}ङ्ग\appbrk{}म\appbrk{}धु\appbrk{}वि\appbrk{}प\appbrk{}र्कन्}\ N}}}{\lemma{\skt{अनुवासनन्}}\Afootnote{H; \skt{अनुवासन}\ N}}\skt{~। तदे\-वो\-त्त\-र\-ब\-स्तिञ्च विदध्यात्~। }\edtext{\skt{पाष\-ण\-भे\-द\-का\-श्म\-र्या\-म\-ल\-क\-पि\-प्प\-ली\-व\-सु\-क\-व\-सि\-र\-सिद्धं}}{\lemma{\skt{पाष\appbrk{}ण\appbrk{}भे\appbrk{}द\appbrk{}का\appbrk{}श्म\appbrk{}र्या\appbrk{}म\appbrk{}ल\appbrk{}क\appbrk{}पि\appbrk{}प्प\appbrk{}ली\appbrk{}व\appbrk{}सु\appbrk{}क\appbrk{}व\appbrk{}सि\appbrk{}र\appbrk{}सिद्धं}}\Afootnote{\skt{पाषा\appbrk{}ण\appbrk{}भे\appbrk{}द\appbrk{}क\appbrk{}का\appbrk{}श्म\appbrk{}या\appbrk{}म\appbrk{}ल\appbrk{}क\appbrk{}पि\appbrk{}प्प\appbrk{}ल्ती\appbrk{}व\appbrk{}शु\appbrk{}क\appbrk{}व\appbrk{}सि\appbrk{}र\appbrk{}सिद्धं}\ N; \skt{पाषा\appbrk{}ण\appbrk{}भे\appbrk{}द\appbrk{}का\appbrk{}श्म\appbrk{}र्या\appbrk{}म\appbrk{}ल\appbrk{}क\appbrk{}पि\appbrk{}प्प\appbrk{}ली\appbrk{}व\appbrk{}सु\appbrk{}क\appbrk{}व\appbrk{}सि\appbrk{}र\appbrk{}सिद्धं}\ H}}\skt{ घृत\-म्प्र\-य\-च्छे\-त्पा\-नार्थं~॥}\pend

\pstart\versepar \ledsidenote{\passref{3.2.7}}\edtext{\skt{भवन्ति}}{\lemma{\skt{भवन्ति}}\Afootnote{H; \emph{om.}\ A N}}\skt{ }\edtext{\skt{चात्र}}{\lemma{\skt{चात्र}}\Afootnote{H; \emph{om.}\ A N}}\skt{ }\edtext{\skt{श्लोकाः}}{\lemma{\skt{श्लोकाः}}\Afootnote{H; \emph{om.}\ A; \skt{श्लोक}\ N}}\skt{~॥ }\pend
\pstart\versepar \skt{}\edtext{\edtext{\skt{पाययेन्मनुजं}}{\lemma{\skt{पाययेन्}}\Afootnote{H; \skt{पायये}\ A; \skt{पाययेत्}\ N}}}{\lemma{\skt{मनुजं}}\Afootnote{H; \skt{त नरं}\ A; \skt{मनुजन्}\ N}}\skt{ }\edtext{\edtext{\skt{सर्पि\-र्भि\-ष\-क्शो\-णि\-त\-रे\-तसि}}{\lemma{\skt{भिषक्}}\Afootnote{A; \skt{भिषज्}\ N; \skt{भ्भिषक्}\ H}}}{\lemma{\skt{शोणितरेतसि}}\Afootnote{H; \skt{कुणपरेतसि}\ A; \skt{छोणितरेतसि}\ N}}\skt{~। }\pend
\pstart\versepar \skt{}\edtext{\skt{धात\-की\-पु\-ष्प\-ख\-दि\-र\-दा\-डि\-मा\-र्ज्जु\-न\-सा\-धि\-तम्}}{\lemma{\skt{धात\appbrk{}की\appbrk{}पु\appbrk{}ष्प\appbrk{}ख\appbrk{}दि\appbrk{}र\appbrk{}दा\appbrk{}डि\appbrk{}मा\appbrk{}र्ज्जु\appbrk{}न\appbrk{}सा\appbrk{}धि\appbrk{}तम्}}\Afootnote{A H; \skt{धात\appbrk{}की\appbrk{}पु\appbrk{}ष्प\appbrk{}ख\appbrk{}दि\appbrk{}र\appbrk{}दा\appbrk{}डि\appbrk{}मा\appbrk{}र्ज्ज\appbrk{}न\appbrk{}सा\appbrk{}धि\appbrk{}तम्}\ N}}\skt{~॥}\pend

\pstart\versepar \ledsidenote{\passref{3.2.8}}\edtext{\skt{कुणपाख्ये}}{\lemma{\skt{कुणपाख्ये}}\Afootnote{H L; \skt{पायये}\ A; \skt{कुणपाख्यो}\ N}}\skt{ }\edtext{\skt{पिबेत्सर्प्पिः}}{\lemma{\skt{पिबेत्}}\Afootnote{H L; \skt{दथवा}\ A; \skt{पिबेत्}\ N}}\skt{ }\edtext{\skt{शाल\-सा\-रा\-दि\-सा\-धितं}}{\lemma{\skt{शाल\appbrk{}सा\appbrk{}रा\appbrk{}दि\appbrk{}सा\appbrk{}धितं}}\Afootnote{A N H; \skt{साल\appbrk{}सा\appbrk{}रा\appbrk{}दि\appbrk{}सा\appbrk{}धि\appbrk{}तम्}\ L}}\skt{~॥ }\pend
\pstart\versepar \skt{}\edtext{\skt{ग्रन्थिभूतेऽश्मभिः}}{\lemma{\skt{ऽश्मभिः}}\Afootnote{H; \skt{शटी}\ A; \skt{स्मभि}\ N; \skt{श्मआभिस्}\ L\textsuperscript{ac}; \skt{श्मभिस्}\ L\textsuperscript{pc}}}\skt{ सिद्धं }\edtext{\skt{पालाशे}}{\lemma{\skt{पालाशे}}\Afootnote{A N H; \skt{पालाशैर्}\ L}}\skt{ }\edtext{\skt{चापि}}{\lemma{\skt{चापि}}\Afootnote{N H; \skt{वापि}\ A L}}\skt{ }\edtext{\skt{भस्मनि}}{\lemma{\skt{भस्मनि}}\Afootnote{A N H; \skt{भस्मभिः}\ L}}\skt{~॥}\pend

\pstart\versepar \ledsidenote{\passref{3.2.9}}\edtext{\edtext{\skt{परू\-ष\-क\-व\-टा\-दि\-भ्या\-म्पू\-याभे}}{\lemma{\skt{परू\appbrk{}ष\appbrk{}क\appbrk{}व\appbrk{}टा\appbrk{}दि\appbrk{}भ्याम्}}\Afootnote{A H; \skt{परू\appbrk{}ष\appbrk{}क\appbrk{}टा\appbrk{}टा\appbrk{}दि\appbrk{}भ्याम्}\ N}}}{\lemma{\skt{पूयाभे}}\Afootnote{N H; \skt{पूयप्रख्ये}\ A}}\skt{ }\edtext{\skt{चापि}}{\lemma{\skt{चापि}}\Afootnote{N H; \skt{च}\ A}}\skt{ साधितम्~॥ }\pend
\pstart\versepar \skt{प्रागुक्तं वक्ष्यते यच्च }\edtext{\skt{तत्कार्यं}}{\lemma{\skt{कार्यं}}\Afootnote{A H; \skt{ञायं}\ N}}\skt{ क्षीणरेतसि~॥}\pend

\pstart \ledsidenote{\passref{3.2.10ab}}\skt{विट्प्रभे }\edtext{\skt{पाय\-ये\-त्स\-र्प्पि\-श्चि\-त्र\-को\-शी\-र\-हि\-ङ्गुभिः}}{\lemma{\skt{सर्प्पिश्}}\Afootnote{N H; \skt{सिद्धं}\ A}}\skt{~।}\footnoteA{The vulgate and Nepalese texts diverge again this point. The vulgate adds a half-śloka \emph{snigdhaṃ vāntaṃ ...} that is not present in the Nepalese version. The Nepalese version then has about eleven passages that are absent in the vulgate. The vulgate's verses 3.2.11cd then appears after 3.2.32.}\pend

\pstart \ledsidenote{\passref{3.2.10.add1}}\skt{स्नेहादीनां क्रमं }\edtext{\edtext{\edtext{\skt{कुर्यात्षट्स्वेतेषु}}{\lemma{\skt{कुर्यात्}}\Afootnote{H; \skt{कुयात्}\ N}}}{\lemma{\skt{षट्स्व्}}\Afootnote{H; \skt{षट्खप्य्}\ N}}}{\lemma{\skt{एतेषु}}\Afootnote{\skt{एडेषु}\ N; \skt{अतेषु}\ H}}\skt{ पण्डितः~॥}\footnoteA{Although not printed in the vulgate editions, the editor Y. T. Ācārya noted that words to this effect were present in some sources (1931 ed., p. 290, note 2; 1938 ed., p. 345, note 3).}\pend

\pstart\versepar \ledsidenote{\passref{3.2.10.add2}}\skt{असेवनाच्चिरं स्त्रीणां क्रिया\-यो\-गा\-त्त\-थैव च~। }\pend
\pstart\versepar \skt{कषा\-य\-क\-टु\-ति\-क्तानां }\edtext{\skt{द्रव्या\-णा\-म\-ति\-से\-व\-नात्}}{\lemma{\skt{द्रव्याणाम्}}\Afootnote{N; \skt{द्र्व्याणाम्}\ H}}\skt{~॥}\pend

\pstart\versepar \ledsidenote{\passref{3.2.10.add3}}\edtext{\skt{तथाम्लं लवणं}}{\lemma{\skt{तथाम्लं लवणं}}\Afootnote{\skt{तथाम्ब्ललवणं}\ N; \skt{तथास्लंलवणं}\ H}}\skt{ रूक्षं शुक्तं }\edtext{\skt{पर्युषितन्तथा}}{\lemma{\skt{पर्युषितन्}}\Afootnote{H; \skt{पयुषितन्}\ N}}\skt{~। }\pend
\pstart\versepar \skt{}\edtext{\skt{वेगाघाताच्च}}{\lemma{\skt{वेगाघाताच्च}}\Afootnote{N; \skt{वेगाद्याताच्च}\ H}}\skt{ }\edtext{\skt{योनीनां}}{\lemma{\skt{योनीनां}}\Afootnote{N; \skt{योगीनां}\ H}}\skt{ गमनाच्च प्रदुष्यति~॥}\pend

\pstart\versepar \ledsidenote{\passref{3.2.10.add4}}\edtext{\skt{तथार्त्तवदोषे च}}{\lemma{\skt{तथार्त्तवदोषे च}}\Afootnote{H; \skt{तथावदोषे च}\ N}}\skt{ }\edtext{\edtext{\skt{स्नेहा\-दि\-क\-र्म\-मु\-प\-दि\-शेत्}}{\lemma{\skt{स्नेहादिकर्मम्}}\Afootnote{H; \skt{स्नेहादिकर्म}\ N}}}{\lemma{\skt{उपदिशेत्}}\Afootnote{H; \skt{स पदसेत्}\ N}}\skt{~। }\pend
\pstart\versepar \skt{यथै\-वो\-क्त\-न्त\-थै\-वो\-त्त\-र\-ब\-स्तिञ्च विदध्यात्~।}\pend

\pstart\versepar \ledsidenote{\passref{3.2.10.add5}}\edtext{\skt{तत्रा\-श\-ण\-वा\-ता\-त्मके}}{\lemma{\skt{तत्रा\appbrk{}श\appbrk{}ण\appbrk{}वा\appbrk{}ता\appbrk{}त्मके}}\Afootnote{H\textsuperscript{pc}; \skt{तत्र वातात्मके}\ N; \skt{तत्रावातात्मके}\ H\textsuperscript{ac}}}\skt{ }\edtext{\skt{भङ्गा\-भ\-द्र\-दा\-रु\-का\-श्म\-र्य\-फ\-ल\-सिद्धं}}{\lemma{\skt{भङ्गा\appbrk{}भ\appbrk{}द्र\appbrk{}दा\appbrk{}रु\appbrk{}का\appbrk{}श्म\appbrk{}र्य\appbrk{}फ\appbrk{}ल\appbrk{}सिद्धं}}\Afootnote{H; \skt{भङ्गा\appbrk{}भ\appbrk{}द्र\appbrk{}दा\appbrk{}रु\appbrk{}का\appbrk{}श्य\appbrk{}प\appbrk{}फ\appbrk{}ल\appbrk{}सिन्ध}\ N}}\skt{ घृतम्पाययेत्~। }\pend
\pstart\versepar \skt{पयस्या }\edtext{\skt{काश्म\-र्य\-फ\-ल\-क्षी\-र\-वि\-दा\-री\-क्षी\-रो\-द\-क\-सिद्धं वा}}{\lemma{\skt{काश्म\appbrk{}र्य\appbrk{}फ\appbrk{}ल\appbrk{}क्षी\appbrk{}र\appbrk{}वि\appbrk{}दा\appbrk{}री\appbrk{}क्षी\appbrk{}रो\appbrk{}द\appbrk{}क\appbrk{}सिद्धं वा}}\Afootnote{H; \skt{हाथा\appbrk{}र्य\appbrk{}फ\appbrk{}ल\appbrk{}क्षी\appbrk{}र\appbrk{}वि\appbrk{}दा\appbrk{}री\appbrk{}क्षी\appbrk{}रोद}\ N}}\skt{ }\edtext{\skt{कृस\-र\-पा\-य\-स\-पिण्डं}}{\lemma{\skt{कृस\appbrk{}र\appbrk{}पा\appbrk{}य\appbrk{}स\appbrk{}पिण्डं}}\Afootnote{H; \skt{… पायसपिण्डं}\ N}}\skt{ }\edtext{\skt{योन्यान्धारयेत्}}{\lemma{\skt{धारयेत्}}\Afootnote{H; \skt{क्षारयेत्}\ N}}\skt{~।}\pend

\pstart \ledsidenote{\passref{3.2.10.add6}}\skt{मधु\-क\-मु\-ङ्ग\-प\-र्ण्णी\-क\-षा\-य\-मा\-च\-मनं पित्तात्मके चार्त्तवे~।}\pend

\pstart \ledsidenote{\passref{3.2.10.add7}}\edtext{\skt{काको\-ली\-क्षी\-र\-वि\-दा\-री\-क्वाथं}}{\lemma{\skt{काको\appbrk{}ली\appbrk{}क्षी\appbrk{}र\appbrk{}वि\appbrk{}दा\appbrk{}री\appbrk{}क्वाथं}}\Afootnote{H; \skt{काको\appbrk{}ली\appbrk{}क्षी\appbrk{}र\appbrk{}वि\appbrk{}दा\appbrk{}री\appbrk{}त्काथं}\ N}}\skt{ }\edtext{\skt{शर्क्क\-रा\-यु\-क्त\-म्पा\-य\-येत्}}{\lemma{\skt{पाययेत्}}\Afootnote{N\textsuperscript{pc} H; \skt{पाययेयोत्}\ N\textsuperscript{ac}}}\skt{~। }\edtext{\edtext{\skt{मधू\-क\-पु\-ष्प\-का\-श्म\-र्य\-फ\-ल\-क्वा\-थ\-म्फा\-नि\-त\-यु\-क्तम्वा}}{\lemma{\skt{मधू\appbrk{}क\appbrk{}पु\appbrk{}ष्प\appbrk{}का\appbrk{}श्म\appbrk{}र्य\appbrk{}फ\appbrk{}ल\appbrk{}क्वा\appbrk{}थम्}}\Afootnote{N; \skt{मधु\appbrk{}क\appbrk{}पु\appbrk{}ष्प-\appbrk{}का\appbrk{}श्म\appbrk{}र्य\appbrk{}फ\appbrk{}ल\appbrk{}क्वा\appbrk{}थम्}\ H}}}{\lemma{\skt{फानितयुक्तम्वा}}\Afootnote{H; \skt{फाणितयुक्तं}\ N}}\skt{ चन्द\-न\-प\-य\-स्या\-कल्कं }\edtext{\edtext{\skt{योन्यान्धारयेत्}}{\lemma{\skt{योन्यान्}}\Afootnote{H; \emph{om.}\ N}}}{\lemma{\skt{धारयेत्}}\Afootnote{H; \skt{पाययेत्}\ N}}\skt{~। }\edtext{\edtext{\skt{निम्ब\-गै\-रि\-क\-क\-षा\-य\-मा\-च\-मनं}}{\lemma{\skt{निम्ब\appbrk{}गै\appbrk{}रि\appbrk{}क\appbrk{}क\appbrk{}षा\appbrk{}यम्}}\Afootnote{H; \skt{निग्व\appbrk{}रौ\appbrk{}रि\appbrk{}क\appbrk{}क\appbrk{}षा\appbrk{}यम्}\ N}}}{\lemma{\skt{आचमनं}}\Afootnote{H; \skt{अचिमनं}\ N}}\skt{~॥ श्लेष्मात्मके }\edtext{\edtext{\skt{भद्र\-श्रि\-य\-क\-टु\-क\-रो\-हि\-णी\-क\-षा\-य\-म्पा\-य\-येत्}}{\lemma{\skt{भद्र\appbrk{}श्रि\appbrk{}य\appbrk{}क\appbrk{}टु\appbrk{}क\appbrk{}रो\appbrk{}हि\appbrk{}णी\appbrk{}क\appbrk{}षा\appbrk{}यम्}}\Afootnote{H; \skt{भद्र\appbrk{}श्रि\appbrk{}य\appbrk{}क\appbrk{}टु\appbrk{}क\appbrk{}रो\appbrk{}हि\appbrk{}णी\appbrk{}कषा…}\ N}}}{\lemma{\skt{पाययेत्}}\Afootnote{H; \skt{ययेत्}\ N}}\skt{~। क्षीर\-वृ\-क्ष\-प्र\-बा\-ल\-क\-ल्कम्वा~। }\edtext{\skt{तिन्दु\-क\-क\-पि\-थ\-श\-ला\-टु\-चू\-र्ण्णम्वा}}{\lemma{\skt{तिन्दु\appbrk{}क\appbrk{}क\appbrk{}पि\appbrk{}थ\appbrk{}श\appbrk{}ला\appbrk{}टु\appbrk{}चू\appbrk{}र्ण्णम्वा}}\Afootnote{H; \skt{तिन्दु\appbrk{}क\appbrk{}क\appbrk{}पि\appbrk{}त्थ\appbrk{}स\appbrk{}ला\appbrk{}दु\appbrk{}चू\appbrk{}र्ण्णम्वा}\ N}}\skt{ क्षौद्रेण लेहयेत्~। }\edtext{\skt{सर्ज्जधवकल्कं}}{\lemma{\skt{सर्ज्जधवकल्कं}}\Afootnote{H; \skt{सर्जथवकल्कं}\ N}}\skt{ }\edtext{\skt{योन्यान्धारयेत्}}{\lemma{\skt{धारयेत्}}\Afootnote{H; \skt{वारयेत्}\ N}}\skt{~। }\edtext{\skt{लोध्र\-ति\-न्दु\-क\-क\-षा\-य\-मा\-च\-मनं}}{\lemma{\skt{लोध्र\appbrk{}ति\appbrk{}न्दु\appbrk{}क\appbrk{}क\appbrk{}षा\appbrk{}यम्}}\Afootnote{\skt{लोध्र\appbrk{}ति\appbrk{}न्दु\appbrk{}क\appbrk{}षा\appbrk{}यम्}\ N; \skt{लाध्र\appbrk{}ति\appbrk{}न्दु\appbrk{}क\appbrk{}क\appbrk{}षा\appbrk{}यम्}\ H}}\skt{~॥}\pend

\pstart \ledsidenote{\passref{3.2.10.add8}}\skt{ग्रन्थिभूते }\edtext{\skt{शृङ्ग\-वे\-र\-पा\-ठा\-चू\-र्ण्ण\-युक्तं}}{\lemma{\skt{शृङ्ग\appbrk{}वे\appbrk{}र\appbrk{}पा\appbrk{}ठा\appbrk{}चू\appbrk{}र्ण्ण\appbrk{}युक्तं}}\Afootnote{H; \skt{शृङ्ग\appbrk{}वे\appbrk{}र\appbrk{}पा\appbrk{}ठा\appbrk{}चू\appbrk{}र्ण्ण\appbrk{}युक्तां}\ N}}\skt{ }\edtext{\skt{श्वेत\-सु\-र\-सा\-म्प्र\-य\-च्छेत्}}{\lemma{\skt{प्रयच्छेत्}}\Afootnote{H; \skt{प्रयच्छेतुआ}\ N}}\skt{~। कुष्ठ\-भ\-द्र\-दा\-रु\-कल्कं }\edtext{\skt{योन्यान्धारयेत्}}{\lemma{\skt{योन्यान्धारयेत्}}\Afootnote{H; \skt{योन्याम्वारयेत्}\ N}}\skt{~। तेषामेव कषा\-य\-मा\-च\-मनं~॥}\pend

\pstart \ledsidenote{\passref{3.2.10.add9}}\skt{कुणपाख्ये }\edtext{\skt{मञ्जि\-ष्ठा\-भ\-द्र\-श्रि\-य\-क\-षायं}}{\lemma{\skt{मञ्जि\appbrk{}ष्ठा\appbrk{}भ\appbrk{}द्र\appbrk{}श्रि\appbrk{}य\appbrk{}क\appbrk{}षायं}}\Afootnote{H; \skt{मंजि\appbrk{}ष्ठा\appbrk{}श्र\appbrk{}य\appbrk{}क\appbrk{}षा\appbrk{}यम्}\ N}}\skt{ पाययेत्~॥ }\edtext{\skt{कुट\-ज\-फ\-ल\-भ\-द्र\-श्रि\-य\-च\-न्द\-नम्वा}}{\lemma{\skt{कुट\appbrk{}ज\appbrk{}फ\appbrk{}ल\appbrk{}भ\appbrk{}द्र\appbrk{}श्रि\appbrk{}य\appbrk{}च\appbrk{}न्द\appbrk{}नम्वा}}\Afootnote{H; \skt{कुटु\appbrk{}ज\appbrk{}फ\appbrk{}ल\appbrk{}भ\appbrk{}द्र\appbrk{}प्रि\appbrk{}य\appbrk{}च\appbrk{}न्दनं वा}\ N}}\skt{ शर्क्करायुक्तं~॥ पूयाभे वा }\edtext{\skt{तदेव}}{\lemma{\skt{तदेव}}\Afootnote{H; \skt{तदेव}\ N}}\skt{ प्रयच्छेत्~।}\pend

\pstart \ledsidenote{\passref{3.2.10.add10}}\skt{खदि\-रा\-र्ज्जु\-न\-कल्कं योन्यां }\edtext{\skt{हारयेत्}}{\lemma{\skt{हारयेत्}}\Afootnote{H; \skt{वारयेत्}\ N}}\skt{~। तयोरेव कषा\-य\-मा\-च\-मनं~॥ क्रिमिजाते }\edtext{\skt{विद\-ञ्ग\-पा\-रि\-भ\-द्र\-क\-म\-ञ्जि\-ष्ठा\-क\-षा\-य\-म्पा\-य\-येत्}}{\lemma{\skt{विद\appbrk{}ञ्ग\appbrk{}पा\appbrk{}रि\appbrk{}भ\appbrk{}द्र\appbrk{}क\appbrk{}म\appbrk{}ञ्जि\appbrk{}ष्ठा\appbrk{}क\appbrk{}षा\appbrk{}यम्}}\Afootnote{H; \skt{विड\appbrk{}ङ्ग\appbrk{}पा\appbrk{}रि\appbrk{}भ\appbrk{}द्र\appbrk{}क\appbrk{}म\appbrk{}ञ्जि\appbrk{}ष्ठा\appbrk{}क\appbrk{}षायं}\ N}}\skt{~।}\pend

\pstart \ledsidenote{\passref{3.2.10.add11}}\skt{क्षौद्रयुक्तं }\edtext{\edtext{\skt{सुरा\-ष्टा\-ष्ट्रा\-रो\-च\-ना\-क\-ल्क\-म्भ\-द्र\-श्रि\-य\-कल्कं}}{\lemma{\skt{सुरा\appbrk{}ष्टा\appbrk{}ष्ट्रा\appbrk{}रो\appbrk{}च\appbrk{}ना\appbrk{}क\appbrk{}ल्कम्}}\Afootnote{H; \skt{सुरा\appbrk{}ष्ठा\appbrk{}रो\appbrk{}च\appbrk{}ना\appbrk{}क\appbrk{}ल्कम्}\ N}}}{\lemma{\skt{भद्रश्रियकल्कं}}\Afootnote{H; \skt{भद्र\appbrk{}श्रि\appbrk{}य\appbrk{}क\appbrk{}ल्क\appbrk{}युक्तं}\ N}}\skt{ }\edtext{\skt{योन्यान्धार्येत्}}{\lemma{\skt{धार्येत्}}\Afootnote{H; \skt{वारयेत्}\ N}}\skt{~। }\edtext{\edtext{\skt{तयोरेव}}{\lemma{\skt{तयोर्}}\Afootnote{H; \skt{त}\ N}}}{\lemma{\skt{एव}}\Afootnote{H; \skt{…व}\ N}}\skt{ }\edtext{\skt{कषा\-य\-मा\-च\-म\-न\-म्भ\-वति}}{\lemma{\skt{भवति}}\Afootnote{H; \emph{om.}\ N}}\skt{~॥}\pend

\pstart\versepar \ledsidenote{\passref{3.2.12cd}}\edtext{\skt{भवति}}{\lemma{\skt{भवति}}\Afootnote{N H; \emph{om.}\ A L}}\skt{ }\edtext{\skt{चात्र}}{\lemma{\skt{चात्र}}\Afootnote{N H; \emph{om.}\ A L}}\skt{~॥ }\pend
\pstart\versepar \skt{}\edtext{\skt{विधि\-मु\-त्त\-र\-ब\-स्त्यन्तं}}{\lemma{\skt{उत्तरबस्त्यन्तं}}\Afootnote{A N H; \skt{उत्तरवस्त्युक्तं}\ L}}\skt{ कुर्या\-दा\-र्त\-व\-शु\-द्धये~।}\pend

\pstart \ledsidenote{\passref{3.2.13cd}}\edtext{\edtext{\skt{कुर्या\-त्क\-ल्क\-म्पि\-चूं\-श्चापि}}{\lemma{\skt{कल्कम्}}\Afootnote{N H; \skt{कल्कान्}\ A; \skt{कल्पान्}\ L}}}{\lemma{\skt{चापि}}\Afootnote{A H L; \skt{चांपि}\ N}}\skt{ }\edtext{\edtext{\skt{पथ्या\-न्या\-च\-म\-नानि}}{\lemma{\skt{पथ्यान्य्}}\Afootnote{A H L; \skt{पथ्या}\ N}}}{\lemma{\skt{आचमनानि}}\Afootnote{A H L; \skt{मनानिनि}\ N\textsuperscript{ac}; \skt{मनानि}\ N\textsuperscript{pc}}}\skt{ च~॥}\pend

\pstart \ledsidenote{\passref{3.2.14cd}}\edtext{\skt{दुर्ग्ग\-न्ध\-पू\-ति\-सं\-काशे}}{\lemma{\skt{दुर्ग्ग\appbrk{}न्ध\appbrk{}पू\appbrk{}ति\appbrk{}सं\appbrk{}काशे}}\Afootnote{H; \skt{दुर्ग\appbrk{}न्धि\appbrk{}पू\appbrk{}य\appbrk{}स\appbrk{}ङ्काशे}\ A; \skt{दुग\appbrk{}र्त्व\appbrk{}पू\appbrk{}ति\appbrk{}स\appbrk{}ङ्काशे}\ N; \skt{दुर्गन्धे पूतिसंकाशे}\ L}}\skt{ }\edtext{\skt{मज्जाभे}}{\lemma{\skt{मज्जाभे}}\Afootnote{N H; \skt{मज्ज}\ A; \skt{मज्जोभे}\ L}}\skt{ }\edtext{\skt{वापि}}{\lemma{\skt{वापि}}\Afootnote{H; \emph{om.}\ A; \skt{चापि}\ N; \skt{वा}\ L}}\skt{ }\edtext{\skt{शोणिते}}{\lemma{\skt{शोणिते}}\Afootnote{N H L; \skt{तुल्ये तथार्तवे}\ A}}\skt{~।}\pend

\pstart \ledsidenote{\passref{3.2.15ab}}\edtext{\edtext{\skt{पिबे\-द्भ\-द्र\-श्रि\-य\-क्वा\-थ\-ञ्च\-न्द\-न\-क्वा\-थ\-मेव}}{\lemma{\skt{भद्र\appbrk{}श्रि\appbrk{}य\appbrk{}क्वा\appbrk{}थञ्}}\Afootnote{H; \skt{भद्रश्रियः क्वाथं}\ A; \skt{भद्र\appbrk{}श्रि\appbrk{}य\appbrk{}त्का\appbrk{}थञ्}\ N}}}{\lemma{\skt{एव}}\Afootnote{A H; \skt{एठ}\ N}}\skt{ }\edtext{\skt{वा}}{\lemma{\skt{वा}}\Afootnote{N H; \skt{च}\ A}}\skt{~॥}\pend

\pstart \ledsidenote{\passref{3.2.14ab}}\skt{ग्रन्थिभूते }\edtext{\skt{पिबेत्पाठां}}{\lemma{\skt{पाठां}}\Afootnote{A H; \skt{पाठा}\ N L}}\skt{ }\edtext{\skt{त्र्यूषणं}}{\lemma{\skt{त्र्यूषणं}}\Afootnote{A H; \skt{भ्यूषणा}\ N; \emph{om.}\ L}}\skt{ }\edtext{\skt{वृक्षकाणि च}}{\lemma{\skt{वृक्षकाणि च}}\Afootnote{A N H; \skt{वृक्षकं कटुभिस्त्रिभिः}\ L}}\skt{~।}\pend

\pstart \ledsidenote{\passref{3.2.14.add1}}\edtext{\skt{क्षार\-ना\-ग\-र\-हि\-ङ्गू\-ना\-न्निः\-क्वाथं}}{\lemma{\skt{निःक्वाथं}}\Afootnote{\skt{निःक्वाथ}\ N; \skt{निक्वाथं}\ H}}\skt{ सुरसाम्पिबेत्~॥}\pend

\pstart \ledsidenote{\passref{3.2.24}}\edtext{\skt{एव\-म\-दु\-ष्ट\-शुक्रः}}{\lemma{\skt{अदुष्टशुक्रः}}\Afootnote{A N H; \skt{दुष्टशुकश्}\ L}}\skt{ }\edtext{\skt{शुद्धार्तवा}}{\lemma{\skt{शुद्धार्तवा}}\Afootnote{\skt{शुद्धार्तवा च}\ A; \skt{शुद्धार्त्तवाम्}\ N H L}}\skt{~॥ }\footnoteA{We emend against all the manuscript witnesses. MS H adds a Newa comma after ṛtau in the next section, suggesting that the scribe thought that was the start of a new section of text.}\pend

\pstart\versepar \ledsidenote{\passref{3.2.25}}\skt{ऋतौ }\edtext{\skt{प्रथ\-म\-दि\-व\-सा\-त्प्र\-भृति}}{\lemma{\skt{प्रथ\appbrk{}म\appbrk{}दि\appbrk{}व\appbrk{}सात्}}\Afootnote{A H L; \skt{थमदिवसात्}\ N}}\skt{ }\edtext{\skt{ब्रह्मचारिणी}}{\lemma{\skt{ब्रह्मचारिणी}}\Afootnote{N; \skt{ब्रह्मचारि}\ A; \skt{ब्रह्मचारिणीं}\ H; \skt{ब्रह्मचा}\ L}}\skt{ }\edtext{\skt{स्नाना\-नु\-ले\-प\-ना\-ल\-ङ्का\-र\-वि\-ले\-ख\-न\-वि\-र\-हिता}}{\lemma{\skt{स्नाना\appbrk{}नु\appbrk{}ले\appbrk{}प\appbrk{}ना\appbrk{}ल\appbrk{}ङ्का\appbrk{}र\appbrk{}वि\appbrk{}ले\appbrk{}ख\appbrk{}न\appbrk{}वि\appbrk{}र\appbrk{}हिता}}\Afootnote{\skt{णी}\ A; \skt{स्नान\appbrk{}नु\appbrk{}ले\appbrk{}प\appbrk{}ना\appbrk{}ङ्ग\appbrk{}ल\appbrk{}ङ्का\appbrk{}र\appbrk{}वि\appbrk{}ले\appbrk{}ख\appbrk{}ख\appbrk{}न\appbrk{}वि\appbrk{}र\appbrk{}हिता}\ N\textsuperscript{ac}; \skt{स्नान\appbrk{}नु\appbrk{}ले\appbrk{}प\appbrk{}ना\appbrk{}ल\appbrk{}ङ्का\appbrk{}र\appbrk{}वि\appbrk{}ले\appbrk{}ख\appbrk{}न\appbrk{}वि\appbrk{}र\appbrk{}हिता}\ N\textsuperscript{pc}; \skt{स्नाना\appbrk{}नु\appbrk{}ले\appbrk{}प\appbrk{}ना\appbrk{}ल\appbrk{}ङ्का\appbrk{}र\appbrk{}वि\appbrk{}ले\appbrk{}ख\appbrk{}न\appbrk{}वि\appbrk{}र\appbrk{}हि\appbrk{}तान्}\ H; \emph{om.}\ L}}\skt{ }\edtext{\skt{दिवा\-स्व\-प्ना\-ञ्ज\-ना\-श्रु\-पा\-ता\-भ्य\-ङ्ग\-न\-ख\-च्छे\-द\-न\-प्र\-धा\-व\-न\-ह\-स\-न\-क\-थ\-ना\-ति\-श\-ब्द\-श्र\-व\-णा\-या\-स\-म्प\-रि\-ह\-रेत्}}{\lemma{\skt{दिवा\appbrk{}स्व\appbrk{}प्ना\appbrk{}ञ्ज\appbrk{}ना\appbrk{}श्रु\appbrk{}पा\appbrk{}ता\appbrk{}भ्य\appbrk{}ङ्ग\appbrk{}न\appbrk{}ख\appbrk{}च्छे\appbrk{}द\appbrk{}न\appbrk{}प्र\appbrk{}धा\appbrk{}व\appbrk{}न\appbrk{}ह\appbrk{}स\appbrk{}न\appbrk{}क\appbrk{}थ\appbrk{}ना\appbrk{}ति\appbrk{}श\appbrk{}ब्द\appbrk{}श्र\appbrk{}व\appbrk{}णा\appbrk{}या\appbrk{}सम्}}\Afootnote{\skt{दिवा\appbrk{}स्व\appbrk{}प्ना\appbrk{}ञ्ज\appbrk{}ना\appbrk{}श्रु\appbrk{}पा\appbrk{}त\appbrk{}स्ना\appbrk{}ना\appbrk{}नु\appbrk{}ले\appbrk{}प\appbrk{}ना\appbrk{}भ्य\appbrk{}ङ्ग\appbrk{}न\appbrk{}ख\appbrk{}च्छे\appbrk{}द\appbrk{}न\appbrk{}प्र\appbrk{}धा\appbrk{}व\appbrk{}न\appbrk{}ह\appbrk{}स\appbrk{}न\appbrk{}क\appbrk{}थ\appbrk{}ना\appbrk{}ति\appbrk{}श\appbrk{}ब्द\appbrk{}श्र\appbrk{}व\appbrk{}णा\appbrk{}व\appbrk{}ले\appbrk{}ख\appbrk{}ना\appbrk{}नि\appbrk{}ला\appbrk{}या\appbrk{}सान्}\ A; \skt{दिवा\appbrk{}स्व\appbrk{}प्नां\appbrk{}ज\appbrk{}ना\appbrk{}श्रू\appbrk{}पा\appbrk{}ता\appbrk{}भ्य\appbrk{}ङ्ग\appbrk{}न\appbrk{}ख\appbrk{}च्छे\appbrk{}द\appbrk{}न\appbrk{}प्र\appbrk{}धा\appbrk{}व\appbrk{}न\appbrk{}ह\appbrk{}स\appbrk{}न\appbrk{}क\appbrk{}थ\appbrk{}ना\appbrk{}ति\appbrk{}श\appbrk{}ब्द\appbrk{}श्र\appbrk{}व\appbrk{}णा\appbrk{}या\appbrk{}सम्}\ N; \skt{दिवा\appbrk{}स्व\appbrk{}प्ना\appbrk{}ञ्ज\appbrk{}ना\appbrk{}श्रु\appbrk{}पा\appbrk{}ता\appbrk{}भ्य\appbrk{}ङ्ग\appbrk{}न\appbrk{}ख\appbrk{}च्छे\appbrk{}द\appbrk{}न\appbrk{}प्र\appbrk{}धा\appbrk{}व\appbrk{}न\appbrk{}ह\appbrk{}स\appbrk{}न\appbrk{}क\appbrk{}थ\appbrk{}ना\appbrk{}ति\appbrk{}श\appbrk{}ब्द\appbrk{}श्र\appbrk{}व\appbrk{}णा\appbrk{}या\appbrk{}साम्}\ H; \emph{om.}\ L}}\skt{~। }\pend
\pstart\versepar \skt{किं }\edtext{\skt{कारणं}}{\lemma{\skt{कारणं}}\Afootnote{A N H; \emph{om.}\ L}}\skt{~। }\edtext{\skt{दिवा}}{\lemma{\skt{दिवा}}\Afootnote{A N H; \emph{om.}\ L}}\skt{ }\edtext{\skt{स्वपत्या}}{\lemma{\skt{स्वपत्या}}\Afootnote{H; \skt{स्वपन्त्याः स्वा}\ A; \skt{स्वपन्त्या}\ N; \skt{णीम्\_}\ L}}\skt{ }\edtext{\skt{बधिरः}}{\lemma{\skt{बधिरः}}\Afootnote{N; \skt{पशीलः}\ A; \emph{om.}\ H L}}\skt{~। }\edtext{\skt{अञ्जनादन्धः}}{\lemma{\skt{अञ्जनादन्धः}}\Afootnote{A; \skt{अजेरादन्धः}\ N; \emph{om.}\ H L}}\skt{~। }\edtext{\edtext{\skt{रोद\-ना\-द्वि\-कृ\-त\-दृष्टिः}}{\lemma{\skt{रोदनाद्}}\Afootnote{A N; \emph{om.}\ H L}}}{\lemma{\skt{विकृतदृष्टिः}}\Afootnote{A N H; \emph{om.}\ L}}\skt{~। }\edtext{\edtext{\skt{स्नाना\-नु\-ले\-प\-ना\-द्दुः\-शीलः}}{\lemma{\skt{स्नाना\appbrk{}नु\appbrk{}ले\appbrk{}प\appbrk{}नाद्}}\Afootnote{A N H; \skt{स्नानानुलेप}\ L}}}{\lemma{\skt{दुःशीलः}}\Afootnote{N H; \skt{दुः}\ A; \emph{om.}\ L}}\skt{~। }\edtext{\edtext{\skt{तैलाभ्यङ्गात्कुष्ठी}}{\lemma{\skt{तैलाभ्यङ्गात्}}\Afootnote{\skt{खशीलः , तैलाभ्यङ्गात्}\ A; \skt{तैलाभ्यङ्गा}\ N; \skt{वधिरः~। अञ्जनादन्धः~। रोद\appbrk{}ना\appbrk{}द्तै\appbrk{}ला\appbrk{}भ्य\appbrk{}ङ्गात्}\ H; \emph{om.}\ L}}}{\lemma{\skt{कुष्ठी}}\Afootnote{A N H; \emph{om.}\ L}}\skt{~। }\edtext{\edtext{\skt{नखा\-व\-क\-र्त्त\-ना\-त्कु\-नखः}}{\lemma{\skt{नखा\appbrk{}व\appbrk{}क\appbrk{}र्त्त\appbrk{}नात्}}\Afootnote{N H; \skt{नखा\appbrk{}प\appbrk{}क\appbrk{}र्त\appbrk{}नात्}\ A; \emph{om.}\ L}}}{\lemma{\skt{कुनखः}}\Afootnote{N H; \skt{कुनखी, प्रधा\appbrk{}व\appbrk{}ना\appbrk{}च्च\appbrk{}ञ्चलः, हस\appbrk{}ना\appbrk{}च्छ्या\appbrk{}व\appbrk{}द\appbrk{}न्तौ\appbrk{}ष्ठ\appbrk{}ता\appbrk{}लु\appbrk{}जिह्वः}\ A; \skt{ना}\ L}}\skt{~। }\edtext{\edtext{\skt{प्रले\-प\-ना\-त्ख\-लतिः}}{\lemma{\skt{प्रलेपनात्}}\Afootnote{H; \skt{प्रलापी चातिकथनात् , अति\appbrk{}श\appbrk{}ब्द\appbrk{}श्र\appbrk{}व\appbrk{}णा\appbrk{}द्ब\appbrk{}धिरः, अवलेखनात्}\ A; \skt{प्रलेखना…}\ N; \skt{लंका\appbrk{}रा\appbrk{}व\appbrk{}ले\appbrk{}ख\appbrk{}न\appbrk{}र\appbrk{}हितां दर्भ\appbrk{}सं\appbrk{}स्त\appbrk{}र\appbrk{}शा\appbrk{}यि\appbrk{}निम्}\ L}}}{\lemma{\skt{खलतिः}}\Afootnote{A H; \skt{…लति}\ N; \emph{om.}\ L}}\skt{~। }\edtext{\skt{मारु\-ता\-या\-स\-से\-व\-ना\-दु\-न्मत्तो}}{\lemma{\skt{मारु\appbrk{}ता\appbrk{}या\appbrk{}स\appbrk{}से\appbrk{}व\appbrk{}ना\appbrk{}दु\appbrk{}न्मत्तो}}\Afootnote{A; \skt{मारु\appbrk{}तो\appbrk{}या\appbrk{}ष\appbrk{}से\appbrk{}व\appbrk{}ना\appbrk{}दु\appbrk{}न्मत्तो}\ N; \skt{मारु\appbrk{}ता\appbrk{}या\appbrk{}स\appbrk{}से\appbrk{}व\appbrk{}ना\appbrk{}दु\appbrk{}त्मत्तो}\ H; \emph{om.}\ L}}\skt{ }\edtext{\skt{गर्भ्भो}}{\lemma{\skt{गर्भ्भो}}\Afootnote{H; \skt{गर्भो}\ A N; \emph{om.}\ L}}\skt{ }\edtext{\edtext{\edtext{\skt{भव\-ती\-त्ये\-व\-मे\-ता\-न्प\-रि\-ह\-रेत्}}{\lemma{\skt{भवतीत्य्}}\Afootnote{A N H; \emph{om.}\ L}}}{\lemma{\skt{एवम्}}\Afootnote{A N H; \emph{om.}\ L}}}{\lemma{\skt{एता\appbrk{}न्प\appbrk{}रि\appbrk{}ह\appbrk{}रेत्}}\Afootnote{A; \skt{एता\appbrk{}म्प\appbrk{}रि\appbrk{}ह\appbrk{}रेत्}\ N H; \emph{om.}\ L}}\footnoteA{We emend from \emph{etām} to the vulgate's \emph{etān} for sense.}\skt{~॥ }\pend
\pstart\versepar \skt{}\edtext{\skt{दर्भ\-सं\-स्त\-र\-ण\-शा\-यिनीं}}{\lemma{\skt{दर्भ\appbrk{}सं\appbrk{}स्त\appbrk{}र\appbrk{}ण\appbrk{}शा\appbrk{}यिनीं}}\Afootnote{\skt{दर्भ\appbrk{}सं\appbrk{}स्त\appbrk{}र\appbrk{}शा\appbrk{}यिनीं}\ A; \skt{दर्भ\appbrk{}सं\appbrk{}स्त\appbrk{}रं\appbrk{}शा\appbrk{}र्पिनी}\ N; \skt{दर्भ्भ\appbrk{}सं\appbrk{}स्त\appbrk{}र\appbrk{}ण\appbrk{}शा\appbrk{}यिनीं}\ H; \emph{om.}\ L}}\skt{ }\edtext{\edtext{\skt{कर\-त\-ल\-श\-रा\-व\-प\-र्ण्णा\-ना\-म\-न्य\-त\-म\-भो\-जिनीं}}{\lemma{\skt{कर\appbrk{}त\appbrk{}ल\appbrk{}श\appbrk{}रा\appbrk{}व\appbrk{}प\appbrk{}र्ण्णा\appbrk{}नाम्}}\Afootnote{H; \skt{कर\appbrk{}त\appbrk{}ल\appbrk{}श\appbrk{}रा\appbrk{}व\appbrk{}पर्णा}\ A; \skt{कर\appbrk{}त\appbrk{}ल\appbrk{}स\appbrk{}रो\appbrk{}व\appbrk{}प\appbrk{}र्ण्णा\appbrk{}नाम्}\ N; \skt{कर\appbrk{}त\appbrk{}ल\appbrk{}श\appbrk{}रा\appbrk{}व\appbrk{}प\appbrk{}र्णाव्}\ L}}}{\lemma{\skt{अन्य\appbrk{}त\appbrk{}म\appbrk{}भो\appbrk{}जिनीं}}\Afootnote{N H; \skt{न्यतमभोजिनीं}\ A; \skt{यन्न्य\appbrk{}त\appbrk{}म\appbrk{}भो\appbrk{}जि\appbrk{}निम्}\ L}}\skt{ }\edtext{\skt{हविष्यं}}{\lemma{\skt{हविष्यं}}\Afootnote{A N H; \skt{भविष्य}\ L}}\skt{ }\edtext{\skt{त्र्यहं}}{\lemma{\skt{त्र्यहं}}\Afootnote{N H L; \skt{त्र्यहं च}\ A}}\skt{ }\edtext{\skt{भर्त्ता}}{\lemma{\skt{भर्त्ता}}\Afootnote{N H L; \skt{भर्तुः}\ A}}\skt{ }\edtext{\skt{रक्षेत्}}{\lemma{\skt{रक्षेत्}}\Afootnote{N H; \skt{संरक्षेत्}\ A; \skt{रषिते}\ L}}\skt{~। ततः }\edtext{\skt{शुद्धस्नातान्तु}}{\lemma{\skt{शुद्धस्नातान्तु}}\Afootnote{N H; \skt{शुद्धस्नातां}\ A L}}\skt{ }\edtext{\edtext{\edtext{\skt{चतु\-र्थ्या\-म\-ना\-ह\-त\-व\-स्त्र\-वा\-सि\-नी\-म\-ल\-ङ्कृतां}}{\lemma{\skt{चतुर्थ्याम्}}\Afootnote{N H L; \skt{चतुर्थेऽ}\ A}}}{\lemma{\skt{अना\appbrk{}ह\appbrk{}त\appbrk{}व\appbrk{}स्त्र\appbrk{}वा\appbrk{}सि\appbrk{}नीम्}}\Afootnote{N H; \skt{हन्य\appbrk{}ह\appbrk{}त\appbrk{}वा\appbrk{}सःस}\ A; \skt{हतवासिनम्}\ L}}}{\lemma{\skt{अलङ्कृतां}}\Afootnote{N H L; \skt{मलङ्कृतां}\ A}}\skt{ }\edtext{\edtext{\skt{कृत\-म\-ङ्ग\-ल\-स्व\-स्ति\-वा\-च\-ना\-म्भ\-र्त्ता\-र\-न्द\-र्श\-येद्}}{\lemma{\skt{कृत\appbrk{}म\appbrk{}ङ्ग\appbrk{}ल\appbrk{}स्व\appbrk{}स्ति\appbrk{}वा\appbrk{}च\appbrk{}नाम्}}\Afootnote{A H; \skt{कृतमङ्गलं स्वस्तिवाचनाम्}\ N; \skt{कृत\appbrk{}मं\appbrk{}ग\appbrk{}ल\appbrk{}स्व\appbrk{}स्ति\appbrk{}वा\appbrk{}चकां}\ L}}}{\lemma{\skt{भर्त्तारन्}}\Afootnote{A H L; \skt{भर्तुर्}\ N}}\skt{~। }\pend
\pstart\versepar \skt{अपि तत्कस्य हेतोः~॥}\edtext{}{\lemma{\skt{परिहरेत्~। किं}}\Afootnote{A N H; \skt{रि}\ L}}\edtext{}{\lemma{\skt{दर्शयेद्~। अपि}}\Afootnote{N H L; \skt{दर्शयेत्}\ A}}\pend

\pstart\versepar \ledsidenote{\passref{3.2.26}}\skt{भवति चात्र~॥ }\pend
\pstart\versepar \skt{पूर्व्व\-म्प\-श्ये\-दृ\-तु\-स्नाता }\edtext{\edtext{\skt{यादृ\-श\-न्न\-र\-म\-ङ्गना}}{\lemma{\skt{यादृशन्}}\Afootnote{A H L; \skt{यासन्}\ N}}}{\lemma{\skt{अङ्गना}}\Afootnote{A H L; \skt{अ…ना}\ N}}\skt{~। }\pend
\pstart\versepar \skt{}\edtext{\edtext{\skt{तादृ\-श\-ञ्ज\-न\-ये\-त्पु\-त्र\-म्भ\-र्त्ता\-र\-न्द\-र्श\-ये\-त्ततः}}{\lemma{\skt{तादृशञ्}}\Afootnote{A N H; \skt{तदृशञ्}\ L}}}{\lemma{\skt{दर्शयेत्}}\Afootnote{N H; \skt{दर्शये}\ A L}}\skt{~॥}\edtext{}{\lemma{\skt{भवति चात्र~॥ पूर्व्वम्}}\Afootnote{N; \skt{पूर्वं}\ A L H}}\pend

\pstart\versepar \ledsidenote{\passref{3.2.27}}\skt{ततो }\edtext{\skt{विधेयं}}{\lemma{\skt{विधेयं}}\Afootnote{H; \skt{विधानं}\ A L; \skt{विधेय}\ N}}\skt{ }\edtext{\skt{पुत्री\-य\-मु\-पा\-ध्यायः}}{\lemma{\skt{पुत्रीयम्}}\Afootnote{A H; \skt{षुंत्रीयम्}\ N; \skt{पुत्रियम्}\ L}}\skt{ }\edtext{\skt{समाचरेत्}}{\lemma{\skt{समाचरेत्}}\Afootnote{A N H; \skt{समारभेत्}\ L}}\skt{~। }\pend
\pstart\versepar \skt{}\edtext{\skt{कर्मान्ते च}}{\lemma{\skt{कर्मान्ते च}}\Afootnote{A H; \skt{कर्मात्ते च}\ N; \skt{कर्मन्ते च}\ L}}\skt{ }\edtext{\skt{क्रमं }\edtext{\edtext{\skt{ह्येतमवेक्ष्येत}}{\lemma{\skt{एतम्}}\Afootnote{N H; \skt{एनम्}\ A; \skt{चेदम्}\ L}}}{\lemma{\skt{अवेक्ष्येत}}\Afootnote{N H; \skt{आरभेत}\ A; \skt{उपेक्षेत}\ L}}}{\lemma{\skt{क्रमं ह्य्}}\Afootnote{A H; \skt{क्रंम ह्य्}\ N; \skt{क्रमञ्}\ L}}\skt{ विचक्षणः~॥}\pend

\pstart \ledsidenote{\passref{3.2.28}}\edtext{\skt{ततोऽपराह्णे}}{\lemma{\skt{ऽपराह्णे}}\Afootnote{A N H; \skt{पराह्णे च}\ L}}\skt{ }\edtext{\edtext{\skt{पुमा\-न्स\-र्प्पिः\-क्षी\-राभ्यां}}{\lemma{\skt{पुमान्}}\Afootnote{H L; \skt{पुमान्मासं ब्रह्मचारी सर्पिःस्निग्धः}\ A; \skt{षुमां}\ N}}}{\lemma{\skt{सर्प्पिःक्षीराभ्यां}}\Afootnote{A N H; \skt{सर्प्पिक्षीराभ्याम्}\ L}}\skt{ }\edtext{\skt{शाल्यो\-द\-न\-म\-भि\-प्राश्य}}{\lemma{\skt{अभिप्राश्य}}\Afootnote{H L; \skt{भु}\ A; \skt{अभिप्राश्}\ N}}\skt{ }\edtext{\skt{मासं}}{\lemma{\skt{मासं}}\Afootnote{H L; \skt{क्त्वा मासं}\ A; \skt{यमासम्}\ N}}\skt{ }\edtext{\skt{ब्रह्मचारी}}{\lemma{\skt{ब्रह्मचारी}}\Afootnote{N H L; \skt{ब्रह्मचारिणीं}\ A}}\skt{ }\edtext{\edtext{\skt{तैल\-मा\-षो\-त्ता\-हा\-रा\-न्ना\-री\-मु\-प\-ग\-च्छे\-द्रात्रौ}}{\lemma{\skt{तैल\appbrk{}मा\appbrk{}षो\appbrk{}त्ता\appbrk{}हा\appbrk{}रान्}}\Afootnote{\skt{तैलस्निग्धां तैल\appbrk{}मा\appbrk{}षो\appbrk{}त्त\appbrk{}रा\appbrk{}हारां}\ A; \skt{तैल\appbrk{}मा\appbrk{}षो\appbrk{}त्त\appbrk{}रा\appbrk{}हा\appbrk{}रान्}\ N H; \skt{तैल\appbrk{}मा\appbrk{}षो\appbrk{}त्त\appbrk{}रान्}\ L}}}{\lemma{\skt{उपगच्छेद्}}\Afootnote{N H; \skt{उपेयाद्}\ A L}}\skt{ }\edtext{\edtext{\skt{साम\-भि\-र\-भि\-वि\-श्वास्य}}{\lemma{\skt{सामभिर्}}\Afootnote{N H; \skt{सामादिभिर्}\ A; \skt{साम}\ L}}}{\lemma{\skt{अभिविश्वास्य}}\Afootnote{A N H; \skt{भिर्विश्वास्या}\ L}}\skt{ }\edtext{\edtext{\skt{विकल्प्यमेवं}}{\lemma{\skt{विकल्प्यम्}}\Afootnote{N; \skt{विकल्प्यै}\ A; \skt{विकल्पम्}\ H; \skt{संकल्प्य}\ L}}}{\lemma{\skt{एवं}}\Afootnote{N H; \skt{वं}\ A; \emph{om.}\ L}}\skt{ }\edtext{\skt{चतुर्थ्यां}}{\lemma{\skt{चतुर्थ्यां}}\Afootnote{A N H; \skt{च एवम्}\ L}}\skt{ }\edtext{\edtext{\edtext{\skt{षाष्ठ्या\-म\-ष्ट\-म्या\-न्द\-शम्यां}}{\lemma{\skt{षाष्ठ्याम्}}\Afootnote{H; \skt{षष्ठ्याम्}\ A L N}}}{\lemma{\skt{अष्टम्यान्}}\Afootnote{A H L; \skt{अष्ट}\ N}}}{\lemma{\skt{दशम्यां}}\Afootnote{A H L; \skt{म्}\ N}}\skt{ }\edtext{\edtext{\skt{द्वाद\-श्या\-ञ्चो\-पे\-याद्}}{\lemma{\skt{द्वादश्याञ्}}\Afootnote{A H L; \emph{om.}\ N}}}{\lemma{\skt{चोपेयाद्}}\Afootnote{H; \skt{चोपेयादिति पुत्रकामः}\ A; \skt{यान् …। चोपेयाद्}\ N; \skt{वोपेयात्}\ L}}\skt{~।}\footnoteA{Witness N inserts a misplaced folio at this point, image A 1267\_11\_67.jpg, covering SS.3.2.22bis, 23, 33, 34, 35, 30.}\pend

\pstart \ledsidenote{\passref{3.2.31b}}\skt{अतः परं}\footnoteA{Witnesses H and N make a saut du même au même from \emph{ataḥ paraṃ} at the start of 30 to the end of 31.}\skt{ }\edtext{\skt{मासादुपेयाल्}}{\lemma{\skt{मासादुपेयाल्}}\Afootnote{N H; \skt{मासादुपेयात्}\ A}}\skt{~।}\pend

\pstart \ledsidenote{\passref{3.2.32}}\edtext{\skt{लब्धगर्भा}\footnoteA{Witness N is faint and written close to the abraded top of the folio so that superscript vowels are uncertain. Witness H has a scribal correction of the vowel \emph{i}. The reading adopted here, with a feminine sentence subject, is neither the Nepalese nor the vulgate and is almost certainly unoriginal. Without further evidence it is not possible to recover the oldest reading.}\skt{ }\edtext{\skt{तेष्वेवाहःसु}}{\lemma{\skt{एवाहःसु}}\Afootnote{N H; \skt{अहःसु}\ A L}}}{\lemma{\skt{लब्धगर्भा तेष्व्}}\Afootnote{\skt{लब्ध\appbrk{}ग\appbrk{}र्भा\appbrk{}या\appbrk{}श्चै\appbrk{}तेष्व्}\ A; \skt{लब्धगर्भन्नि व च तेष्व्}\ N H; \skt{लब्धगर्भान्न चैवेतेष्व्}\ L}}\skt{ }\edtext{\edtext{\edtext{\skt{लक्ष्म\-णा\-व\-ट\-शु\-ङ्गा\-स\-ह\-दे\-वा\-वि\-श्व\-दे\-वा\-ना\-म\-न्य\-त\-म\-म\-भि\-षू\-यात्}}{\lemma{\skt{लक्ष्म\appbrk{}णा\appbrk{}व\appbrk{}ट\appbrk{}शु\appbrk{}ङ्गा\appbrk{}स\appbrk{}ह\appbrk{}दे\appbrk{}वा\appbrk{}वि\appbrk{}श्व\appbrk{}दे\appbrk{}वा\appbrk{}नाम्}}\Afootnote{H L; \skt{लक्ष्म\appbrk{}णा\appbrk{}व\appbrk{}ट\appbrk{}शु\appbrk{}ङ्ग\appbrk{}स\appbrk{}ह\appbrk{}दे\appbrk{}वा\appbrk{}वि\appbrk{}श्व\appbrk{}दे\appbrk{}वा\appbrk{}नाम्}\ A; \skt{लक्ष्म\appbrk{}णा\appbrk{}व\appbrk{}ट\appbrk{}शु\appbrk{}ङ्गा\appbrk{}स\appbrk{}ह\appbrk{}सु\appbrk{}दे\appbrk{}वा\appbrk{}वि\appbrk{}श्व\appbrk{}दे\appbrk{}वा\appbrk{}नाम्}\ N}}}{\lemma{\skt{अन्यतमम्}}\Afootnote{N H L; \skt{अन्यतमां}\ A}}\skt{~। त्रींश्चतुरो}}{\lemma{\skt{अभिषूयात्~। त्रींश्}}\Afootnote{H; \skt{क्षीरेणाभिषुत्य त्रींश्}\ A; \skt{अभिषूयात्~। त्रिंश्}\ N; \skt{अभिषुय त्रींश्}\ L}}\skt{ वा बिन्दू\-न्द\-द्या\-द्द\-क्षिणे नासापुटे }\edtext{\skt{पुत्रकामायै}}{\lemma{\skt{पुत्रकामायै}}\Afootnote{N H L; \skt{पुत्र}\ A}}\skt{ }\edtext{\skt{वामे}}{\lemma{\skt{वामे}}\Afootnote{N H L; \emph{om.}\ A}}\skt{ }\edtext{\skt{स्त्रीकामायै न}}{\lemma{\skt{स्त्रीकामायै न}}\Afootnote{N H L; \skt{कामायै , न}\ A}}\skt{ }\edtext{\edtext{\skt{चैन\-न्नि\-ष्ठी\-वे\-दिति}}{\lemma{\skt{चैनन्}}\Afootnote{N H; \skt{च तान्}\ A; \skt{चैवेतन्}\ L}}}{\lemma{\skt{निष्ठीवेदिति}}\Afootnote{N H; \skt{निष्ठीवेत्}\ A; \skt{निष्ठीवेदिति भातका}\ L}}\skt{~॥}\pend

\pstart\versepar \ledsidenote{\passref{3.2.11cd}}\edtext{\skt{श्लोकाः}}{\lemma{\skt{श्लोकाः}}\Afootnote{H; \emph{om.}\ A L; \skt{श्लोक~॥ श्लोका}\ N}}\skt{~॥ }\pend
\pstart\versepar \skt{}\edtext{\skt{स्फटिकाभं}}{\lemma{\skt{स्फटिकाभं}}\Afootnote{A H L; \skt{स्फेटि}\ N}}\skt{ }\edtext{\skt{द्रवं}}{\lemma{\skt{द्रवं}}\Afootnote{A H L; \emph{om.}\ N}}\skt{ }\edtext{\skt{स्निग्धं}}{\lemma{\skt{स्निग्धं}}\Afootnote{A H; \emph{om.}\ N; \skt{स्निग्धम}\ L}}\skt{ }\edtext{\edtext{\skt{मधु\-र\-म्म\-धु\-गन्धि}}{\lemma{\skt{मधुरम्}}\Afootnote{A H L; \emph{om.}\ N}}\skt{ च}}{\lemma{\skt{मधुगन्धि च}}\Afootnote{A H L; \skt{वसि\appbrk{}सृ\appbrk{}स\appbrk{}धु\appbrk{}भ\appbrk{}म्य\appbrk{}युप …}\ N}}\skt{~।}\pend

\pstart \ledsidenote{\passref{3.2.12ab}}\skt{शुक्रमिच्छन्ति केचित्तु }\edtext{\skt{क्षौद्र\-तै\-ल\-स\-म\-प्र\-भम्}}{\lemma{\skt{क्षौद्र\appbrk{}तै\appbrk{}ल\appbrk{}स\appbrk{}म\appbrk{}प्र\appbrk{}भम्}}\Afootnote{H; \skt{तैलक्षौद्रनिभं तथा}\ A; \skt{क्षौद्र\appbrk{}तै\appbrk{}ल\appbrk{}नि\appbrk{}भ\appbrk{}न्तथा}\ L}}\skt{~॥}\pend

\pstart\versepar \ledsidenote{\passref{3.2.17}}\edtext{\skt{शशासृक्प्रतिमं}}{\lemma{\skt{शशासृक्प्रतिमं}}\Afootnote{A H; \skt{सासृक्षुतिमं}\ N; \skt{शशा\appbrk{}सृ\appbrk{}क्\_\appbrk{}स\appbrk{}दृशं}\ L}}\skt{ }\edtext{\skt{यत्तु}}{\lemma{\skt{यत्तु}}\Afootnote{A H L; \skt{यन्तु}\ N}}\skt{ यद्वा }\edtext{\skt{लाक्षारसोपमं}}{\lemma{\skt{लाक्षारसोपमं}}\Afootnote{A H; \skt{लाक्षारसोपसं}\ N; \skt{लाक्षा\appbrk{}र\appbrk{}सो\appbrk{}प\appbrk{}मम}\ L}}\skt{~। }\pend
\pstart\versepar \skt{}\edtext{\skt{तदा\-र्त\-व\-म्प्र\-श\-न्सन्ति}}{\lemma{\skt{प्रशन्सन्ति}}\Afootnote{A H L; \skt{प्रसुसन्ति}\ N}}\skt{ }\edtext{\skt{यच्च न}}{\lemma{\skt{च न}}\Afootnote{\skt{वासो न}\ A; \skt{चाध्भ्यु न}\ N; \skt{चा न}\ H; \skt{वाप्सु न}\ L}}\skt{ }\edtext{\skt{विरज्यते}}{\lemma{\skt{विरज्यते}}\Afootnote{N H; \skt{विरञ्जयेत्}\ A; \skt{विरज्यति}\ L}}\skt{~॥}\pend

\pstart\versepar \ledsidenote{\passref{3.2.18}}\edtext{\skt{तदे\-वा\-ति\-प्र\-स\-ङ्गेन}}{\lemma{\skt{तदे\appbrk{}वा\appbrk{}ति\appbrk{}प्र\appbrk{}स\appbrk{}ङ्गेन}}\Afootnote{A N H; \skt{तदे\appbrk{}व\appbrk{}ति\appbrk{}प्र\appbrk{}सं\appbrk{}गेन}\ L}}\skt{ }\edtext{\edtext{\skt{प्रवृ\-त्त\-म\-नृ\-ता\-वपि}}{\lemma{\skt{प्रवृत्तम्}}\Afootnote{A H L; \skt{प्रवृत्त}\ N}}}{\lemma{\skt{अनृताव्}}\Afootnote{A H; \skt{ननृताव्}\ N; \skt{अन्नृताव्}\ L}}\skt{~। }\pend
\pstart\versepar \skt{}\edtext{\edtext{\edtext{\skt{असृ\-ग्द\-र\-म्वि\-जा\-नी\-या\-द्र\-क्त\-ल\-क्ष\-ण\-ल\-क्षितं}}{\lemma{\skt{असृग्दरम्}}\Afootnote{A H L; \skt{असृघरं}\ N}}}{\lemma{\skt{विजानीयाद्}}\Afootnote{H L; \skt{विजा\appbrk{}नी\appbrk{}या\appbrk{}द\appbrk{}तोऽ\appbrk{}न्यद्}\ A; \skt{विजानीयाद्रक्}\ N}}}{\lemma{\skt{रक्त\appbrk{}ल\appbrk{}क्ष\appbrk{}ण\appbrk{}ल\appbrk{}क्षितं}}\Afootnote{H L; \skt{रक्तलक्षणात्}\ A; \skt{षक्त\appbrk{}ल\appbrk{}क्ष\appbrk{}थ\appbrk{}ल\appbrk{}क्षितं}\ N}}\skt{~॥}\pend

\pstart\versepar \ledsidenote{\passref{3.2.19}}\edtext{\skt{असृग्दरं}}{\lemma{\skt{असृग्दरं}}\Afootnote{N H; \skt{असृग्दरो}\ A}}\skt{ }\edtext{\edtext{\edtext{\skt{भवे\-त्स\-र्व्व\-न्सा\-ङ्ग\-म\-र्द्द\-न्स\-वे\-दनं}}{\lemma{\skt{सर्व्वन्}}\Afootnote{H; \skt{सर्वः}\ A; \skt{सर्वं}\ N}}}{\lemma{\skt{साङ्गमर्द्दन्}}\Afootnote{H; \skt{साङ्गमर्दः}\ A; \skt{साङ्गमर्द्द}\ N}}}{\lemma{\skt{सवेदनं}}\Afootnote{N H; \skt{सवेदनः}\ A}}\skt{~। }\pend
\pstart\versepar \skt{}\edtext{\skt{अतिप्रवृत्तौ}}{\lemma{\skt{अतिप्रवृत्तौ}}\Afootnote{N H; \skt{तस्यातिवृत्तौ}\ A}}\skt{ }\edtext{\skt{दौर्ब्ब\-ल्य\-म्भ्र\-म\-मूर्च्छा}}{\lemma{\skt{भ्रममूर्च्छा}}\Afootnote{N H; \skt{भ्रमो मूर्च्छा}\ A}}\skt{ तमस्तृषा~॥}\pend

\pstart\versepar \ledsidenote{\passref{3.2.20}}\skt{दाहः प्रलापः पाण्डुत्वन्तन्द्रा रोगाश्च वातजाः~। }\pend
\pstart\versepar \skt{तरुण्या }\edtext{\edtext{\edtext{\skt{हित\-से\-वि\-न्या\-स्त\-द\-ल्पो\-प\-द्र\-व\-म्भ\-वेत्}}{\lemma{\skt{हितसेविन्यास्}}\Afootnote{A H; \skt{हि…सेविन्यास्}\ N}}}{\lemma{\skt{तद\appbrk{}ल्पो\appbrk{}प\appbrk{}द्र\appbrk{}वम्}}\Afootnote{H; \skt{तमल्पोपद्रवं}\ A; \skt{तल्पोपद्रवम्}\ N\textsuperscript{ac}; \skt{तद\appbrk{}ल्पो\appbrk{}प\appbrk{}द्र\appbrk{}वम्}\ N\textsuperscript{pc}}}}{\lemma{\skt{भवेत्}}\Afootnote{N H; \skt{भिषक्}\ A}}\skt{~॥}\pend

\pstart \ledsidenote{\passref{3.2.21cd}}\edtext{\skt{दोषै\-रा\-वृ\-त\-मा\-र्ग्ग\-त्वा\-दा\-र्त्त\-व\-न्न\-श्यति}}{\lemma{\skt{आवृ\appbrk{}त\appbrk{}मा\appbrk{}र्ग्ग\appbrk{}त्वा\appbrk{}दा\appbrk{}र्त्त\appbrk{}वन्}}\Afootnote{A; \skt{आवृ\appbrk{}त\appbrk{}मा\appbrk{}र्ग\appbrk{}न्वा\appbrk{}दा\appbrk{}र्त्तष}\ N; \skt{आवृ\appbrk{}त\appbrk{}मा\appbrk{}र्ग्ग\appbrk{}त्वा\appbrk{}दा\appbrk{}र्त्त\appbrk{}वन्}\ H}}\skt{ }\edtext{\skt{स्त्रियाः}}{\lemma{\skt{स्त्रियाः}}\Afootnote{A H; \skt{स्त्रियः}\ N}}\skt{~।}\pend

\pstart\versepar \ledsidenote{\passref{3.2.22}}\skt{तत्र }\edtext{\skt{मांस\-कु\-ल\-त्था\-म्ल\-ति\-ल\-मा\-ष\-सुरा}}{\lemma{\skt{मांस\appbrk{}कु\appbrk{}ल\appbrk{}त्था\appbrk{}म्ल\appbrk{}ति\appbrk{}ल\appbrk{}मा\appbrk{}ष\appbrk{}सुरा}}\Afootnote{H; \skt{मत्स्य\appbrk{}कु\appbrk{}ल\appbrk{}त्था\appbrk{}म्ल\appbrk{}ति\appbrk{}ल\appbrk{}मा\appbrk{}ष\appbrk{}सुरा}\ A; \skt{मान्स\appbrk{}कु\appbrk{}ल\appbrk{}त्था\appbrk{}म्ब्ल\appbrk{}ति\appbrk{}ल\appbrk{}मा\appbrk{}ष\appbrk{}सुरा}\ N}}\skt{ }\edtext{\skt{हिताः}}{\lemma{\skt{हिताः}}\Afootnote{A H; \skt{हिता}\ N}}\skt{~॥ }\pend
\pstart\versepar \skt{}\edtext{\skt{पाने}}{\lemma{\skt{पाने}}\Afootnote{A H; \skt{पारें}\ N}}\skt{ }\edtext{\skt{मूत्रमुदश्विच्च}}{\lemma{\skt{उदश्विच्च}}\Afootnote{A H; \skt{उदत्स्विर्व}\ N}}\skt{ दधिशुक्तञ्च भोजने~।}\pend

\pstart\versepar \ledsidenote{\passref{3.2.23}}\skt{क्षीणम्प्रागीरितं रक्तं सल\-क्ष-\-ण\-चि\-कि\-त्सितं~॥ }\pend
\pstart\versepar \skt{}\edtext{\skt{तथाप्यस्य}}{\lemma{\skt{अस्य}}\Afootnote{N H; \skt{अत्र}\ A}}\skt{ विधा\-त\-व्य\-म्वि\-धानं }\edtext{\skt{रक्तपित्तवत्}}{\lemma{\skt{रक्तपित्तवत्}}\Afootnote{N H; \skt{नष्टरक्तवत्}\ A}}\skt{~।}\pend

\pstart\versepar \ledsidenote{\passref{3.2.33}}\edtext{\edtext{\skt{ध्रुवञ्चतुर्ण्णां}}{\lemma{\skt{ध्रुवञ्}}\Afootnote{A H L; \skt{ध्रूवम्}\ N}}}{\lemma{\skt{चतुर्ण्णां}}\Afootnote{A H; \skt{चतुर्ण्णा}\ N; \skt{त्रयाणां}\ L}}\skt{ }\edtext{\skt{सान्निध्ये गर्भ्भः}}{\lemma{\skt{सान्निध्ये गर्भ्भः}}\Afootnote{H; \skt{सान्निध्याद्गर्भः}\ A; \skt{सानिण्यगर्भं}\ N; \skt{सान्निद्ध्यात्गर्भ}\ L}}\skt{ }\edtext{\skt{स्याद्विधिपूर्व्वकः}}{\lemma{\skt{विधिपूर्व्वकः}}\Afootnote{N H; \skt{विधिपूर्वकम्}\ A L}}\skt{~॥ }\pend
\pstart\versepar \skt{}\edtext{\skt{क्षेत्त्र\-बी\-जो\-द\-क\-तृणां}}{\lemma{\skt{क्षेत्त्र\appbrk{}बी\appbrk{}जो\appbrk{}द\appbrk{}क\appbrk{}तृणां}}\Afootnote{H; \skt{ऋतु\appbrk{}क्षे\appbrk{}त्रा\appbrk{}म्बु\appbrk{}बी\appbrk{}जानां}\ A; \skt{क्षत्र\appbrk{}बी\appbrk{}जो\appbrk{}द\appbrk{}क\appbrk{}र्क्कु\appbrk{}नात्}\ N; \skt{क्षेत्र\appbrk{}म्बी\appbrk{}जो\appbrk{}द\appbrk{}का\appbrk{}दीरां}\ L}}\skt{ }\edtext{\skt{सामग्र्यादङ्कुरो}}{\lemma{\skt{सामग्र्यादङ्कुरो}}\Afootnote{A H; \skt{सामग्र्यादंज्जु}\ N; \skt{सान्निद्ध्यादङ्कुरो}\ L}}\skt{ }\edtext{\skt{यथा}}{\lemma{\skt{यथा}}\Afootnote{A H L; \emph{om.}\ N}}\skt{~॥}\pend

\pstart\versepar \ledsidenote{\passref{3.2.34}}\skt{एवं जाता }\edtext{\skt{रूपवन्तो}}{\lemma{\skt{रूपवन्तो}}\Afootnote{A H; \skt{रूपवत्ता}\ N; \skt{रूपावन्तो}\ L}}\skt{ }\edtext{\skt{महा\-स\-त्त्वा\-श्चि\-रा\-युषः}}{\lemma{\skt{महासत्त्वाश्}}\Afootnote{H L; \skt{सत्त्ववन्तश्}\ A; \skt{महासन्धाश्}\ N}}\skt{~॥ }\pend
\pstart\versepar \skt{}\edtext{\edtext{\skt{भव\-न्त्यृ\-ण\-प्र\-मो\-क्ता\-र\-स्स\-पुत्राः}}{\lemma{\skt{भव\appbrk{}न्त्यृ\appbrk{}ण\appbrk{}प्र\appbrk{}मो\appbrk{}क्ता\appbrk{}रस्}}\Afootnote{H; \skt{भवन्त्यृणस्य मोक्तारः}\ A; \skt{भव\appbrk{}न्त्या\appbrk{}ण\appbrk{}प्र\appbrk{}मो\appbrk{}क्तारः}\ N L}}}{\lemma{\skt{सपुत्राः}}\Afootnote{N H; \skt{सत्पुत्राः}\ A L}}\skt{ }\edtext{\skt{पुत्रिणां}}{\lemma{\skt{पुत्रिणां}}\Afootnote{N H; \skt{पुत्रिणे}\ A; \skt{पुत्रिणा}\ L}}\skt{ हिताः~॥}\pend

\pstart \ledsidenote{\passref{3.2.35}}\skt{तत्र }\edtext{\edtext{\edtext{\skt{तेजो\-धा\-तु\-र्व\-र्ण्णा\-ना\-म्प्र\-भवः}}{\lemma{\skt{तेजोधातुर्}}\Afootnote{A H; \skt{तेजोधातु}\ N L}}}{\lemma{\skt{वर्ण्णानाम्}}\Afootnote{A N H; \skt{वर्णानाम}\ L}}}{\lemma{\skt{प्रभवः}}\Afootnote{N H L; \skt{प्रभवः}\ A}}\skt{ }\edtext{\skt{इति}}{\lemma{\skt{इति}}\Afootnote{N H L; \emph{om.}\ A}}\skt{ }\edtext{\skt{कृत्वा}}{\lemma{\skt{कृत्वा}}\Afootnote{N H L; \skt{स}\ A}}\skt{ }\edtext{\skt{यदा}}{\lemma{\skt{यदा}}\Afootnote{A N H; \skt{दया}\ L}}\skt{ }\edtext{\edtext{\edtext{\skt{गर्भो\-त्प\-त्ता\-व\-पा\-न्धा\-तु\-प्रा\-य\-मा\-हा\-र\-म्भ\-वति}}{\lemma{\skt{गर्भोत्पत्ताव्}}\Afootnote{A N; \skt{गर्ब्भोत्पत्तावपा}\ H; \skt{गर्भोतोअत्ताव्}\ L}}}{\lemma{\skt{अपा\appbrk{}न्धा\appbrk{}तु\appbrk{}प्रा\appbrk{}यम्}}\Afootnote{N; \skt{अब्धातुप्रायो}\ A L; \skt{अब्धातुप्रायम्}\ H}}}{\lemma{\skt{आहारम्}}\Afootnote{N H; \emph{om.}\ A L}}\skt{ तदा }\edtext{\skt{गर्भ्भस्य}}{\lemma{\skt{गर्भ्भस्य}}\Afootnote{H; \skt{गर्भं}\ A; \skt{गर्भस्य}\ N L}}\skt{ }\edtext{\skt{गौरत्वं}}{\lemma{\skt{गौरत्वं}}\Afootnote{N H; \skt{गौरं क}\ A; \skt{गौरवं क}\ L}}\skt{~।}\footnoteA{Note the slightly unusual form \emph{apān-dhātu-prāyam} with the inflected genitive pl. of the first member of the compound (Whitney 1250e).}\skt{ }\edtext{\skt{पृथि\-वी\-धा\-तु\-प्रायः}}{\lemma{\skt{पृथि\appbrk{}वी\appbrk{}धा\appbrk{}तु\appbrk{}प्रायः}}\Afootnote{N H; \skt{रोति, पृथि\appbrk{}वी\appbrk{}धा\appbrk{}तु\appbrk{}प्रायः}\ A; \skt{रोति~। पृथिवीधातुः प्रायः}\ L}}\skt{ }\edtext{\skt{कृष्णत्वं}}{\lemma{\skt{कृष्णत्वं}}\Afootnote{N H; \skt{कृष्णं}\ A; \emph{om.}\ L}}\skt{~। }\edtext{\skt{पृथि\-व्या\-का\-श\-धा\-तु\-प्रायः}}{\lemma{\skt{पृथि\appbrk{}व्या\appbrk{}का\appbrk{}श\appbrk{}धा\appbrk{}तु\appbrk{}प्रायः}}\Afootnote{N H; \skt{पृथि\appbrk{}व्या\appbrk{}का\appbrk{}श\appbrk{}धा\appbrk{}तु\appbrk{}प्रायः कृष्णश्यामं, तोया\appbrk{}का\appbrk{}श\appbrk{}धा\appbrk{}तु\appbrk{}प्रायो}\ A; \skt{काळबं पृथि\appbrk{}व्या\appbrk{}का\appbrk{}श\appbrk{}तो\appbrk{}य\appbrk{}धातुः प्राय}\ L\textsuperscript{ac}; \skt{काळबं पृथि\appbrk{}व्या\appbrk{}का\appbrk{}श\appbrk{}तो\appbrk{}य\appbrk{}धातुः प्रा}\ L\textsuperscript{pc}}}\skt{ }\edtext{\skt{श्यामत्वं}}{\lemma{\skt{श्यामत्वं}}\Afootnote{N H; \skt{गौरश्यामम्}\ A; \skt{स्यश्यामत्वम्}\ L\textsuperscript{ac}; \skt{य श्यामत्वम्}\ L\textsuperscript{pc}}}\skt{~। }\edtext{\edtext{\edtext{\skt{यादृ\-द्व\-र्ण्ण\-मा\-हा\-र\-मु\-प\-से\-वते}}{\lemma{\skt{यादृद्वर्ण्णम्}}\Afootnote{\skt{यादृग्वर्णम्}\ A N; \skt{यादृक्वर्ण्णम्}\ H; \skt{वादृशवर्ण्णं}\ L}}}{\lemma{\skt{आहारम्}}\Afootnote{A N H; \skt{अन्नम्}\ L}}}{\lemma{\skt{उपसेवते}}\Afootnote{A H L; \skt{उपसेवेत}\ N}}\skt{ }\edtext{\skt{गर्भ्भिणी}}{\lemma{\skt{गर्भ्भिणी}}\Afootnote{H; \skt{गर्भिणी}\ A N L}}\skt{ }\edtext{\skt{तद्वर्ण्णः}}{\lemma{\skt{तद्वर्ण्णः}}\Afootnote{N H; \skt{तादृग्वर्ण}\ A; \skt{तादृकवर्ण्णं}\ L}}\skt{ }\edtext{\skt{प्रसवो}}{\lemma{\skt{प्रसवो}}\Afootnote{H; \skt{प्रसवा}\ A; \skt{प्रस}\ N; \skt{प्रसू}\ L}}\skt{ }\edtext{\skt{भवतीत्येके}}{\lemma{\skt{भवतीत्य्}}\Afootnote{H; \skt{भवति’ इत्य्}\ A; \skt{तीत्य्}\ N; \skt{त इत्य्}\ L}}\skt{ भाषन्ते~। }\edtext{\skt{यथा}}{\lemma{\skt{यथा}}\Afootnote{H; \emph{om.}\ A L; \skt{यधा}\ N}}\skt{ }\edtext{\skt{कृष्ण\-पी\-त\-श्वे\-तासु भूमिषु}}{\lemma{\skt{कृष्ण\appbrk{}पी\appbrk{}त\appbrk{}श्वे\appbrk{}तासु भूमिषु}}\Afootnote{N H; \emph{om.}\ A L}}\skt{ }\edtext{\skt{सर्प्प\-वृ\-श्चि\-क\-ग\-ल\-गो\-डि\-का\-दयः}}{\lemma{\skt{सर्प्प\appbrk{}वृ\appbrk{}श्चि\appbrk{}क\appbrk{}ग\appbrk{}ल\appbrk{}गो\appbrk{}डि\appbrk{}का\appbrk{}दयः}}\Afootnote{N; \emph{om.}\ A L; \skt{सर्प्प\appbrk{}वृ\appbrk{}श्चि\appbrk{}क\appbrk{}ग\appbrk{}ल\appbrk{}गो\appbrk{}ण्डि\appbrk{}का\appbrk{}दयः}\ H}}\footnoteA{Witness N breaks off at this point and is unavailable for the remainder of the Suśrutasaṃhitā.}\skt{ }\edtext{\skt{सत्वाः कृष्णपीतश्वेता इति~॥}\footnoteA{\emph{galagoṇḍikādayaḥ} in witness H is likely to be a variant pronunciation rather than a scribal error.}\skt{ तत्र}}{\lemma{\skt{सत्वाः कृष्णपीतश्वेता इति~॥ तत्र}}\Afootnote{H; \skt{तत्र}\ A L; \emph{om.}\ N}}\skt{ }\edtext{\edtext{\edtext{\edtext{\skt{दृष्टि\-भा\-ग\-म\-प्र\-ति\-प\-न्न\-न्ते\-जो\-धा\-तु\-र्ज्जा\-त्यन्धं}}{\lemma{\skt{दृष्टिभागम्}}\Afootnote{A H L; \emph{om.}\ N}}}{\lemma{\skt{अप्रतिपन्नन्}}\Afootnote{A H; \emph{om.}\ N; \skt{प्रतिपन्नस्}\ L}}}{\lemma{\skt{तेजोधातुर्}}\Afootnote{H L; \skt{तेजो}\ A; \emph{om.}\ N}}}{\lemma{\skt{ज्जात्यन्धं}}\Afootnote{A H L; \emph{om.}\ N}}\skt{ }\edtext{\skt{करोति}}{\lemma{\skt{करोति}}\Afootnote{A H L; \emph{om.}\ N}}\skt{~। }\edtext{\skt{तदेव}}{\lemma{\skt{तदेव}}\Afootnote{A H; \emph{om.}\ N; \skt{स एव}\ L}}\skt{ }\edtext{\skt{रक्तानुगतं}}{\lemma{\skt{रक्तानुगतं}}\Afootnote{A H; \emph{om.}\ N; \skt{रक्तानुगतो}\ L}}\skt{ }\edtext{\skt{रक्ताक्षं}}{\lemma{\skt{रक्ताक्षं}}\Afootnote{H L; \skt{रक्ताक्षं, पित्तानुगतं पिङ्गाक्षं}\ A; \emph{om.}\ N}}\skt{~। }\edtext{\skt{श्लेष्मानुगतं}}{\lemma{\skt{श्लेष्मानुगतं}}\Afootnote{A H; \emph{om.}\ N; \skt{श्लेष्मानुगतः}\ L}}\skt{ }\edtext{\skt{शुक्लाक्षं}}{\lemma{\skt{शुक्लाक्षं}}\Afootnote{H L; \skt{शुक्ला}\ A; \emph{om.}\ N}}\skt{~। }\edtext{\edtext{\skt{पित्ता\-नु\-ग\-त\-म्पि\-ङ्गाक्षं}}{\lemma{\skt{पित्तानुगतम्}}\Afootnote{H; \emph{om.}\ A N; \skt{पित्तानुगतः}\ L}}}{\lemma{\skt{पिङ्गाक्षं}}\Afootnote{H L; \skt{क्षं}\ A; \emph{om.}\ N}}\skt{~। }\edtext{\edtext{\skt{वाता\-नु\-ग\-त\-म्वि\-कृ\-ताक्षं}}{\lemma{\skt{वातानुगतम्}}\Afootnote{A H; \emph{om.}\ N; \skt{वातानुगतो}\ L}}}{\lemma{\skt{विकृताक्षं}}\Afootnote{H; \skt{विकृताक्षमिति}\ A; \emph{om.}\ N; \skt{विकृताक्षमिति भकः}\ L}}\skt{~॥}\pend

\pstart\versepar \ledsidenote{\passref{3.2.35.1}}\skt{भवन्ति चात्र~॥ }\pend
\pstart\versepar \skt{विशुद्धः पवनो यस्य नयने प्रतिपद्यते~। }\pend
\pstart\versepar \skt{}\edtext{\skt{जायेते}}{\lemma{\skt{जायेते}}\Afootnote{\skt{जायते}\ H}}\skt{ नयने तस्य निम्ने }\edtext{\skt{कृष्णेऽसिते}}{\lemma{\skt{ऽसिते}}\Afootnote{\skt{सिते}\ H}}\skt{ तथा}\footnoteA{Witness H appears to read \emph{jāyate} but we emend to the dual.}\skt{~॥}\pend

\pstart\versepar \ledsidenote{\passref{3.2.35.2}}\skt{पित्तं कफयुतं यस्य निर्म्मलं याति लोचने~। }\pend
\pstart\versepar \skt{हरिपिङ्ग्लसम्ज्ञा वै जायेते तस्य लोचने~॥}\pend

\pstart\versepar \ledsidenote{\passref{3.2.35.3}}\skt{विनिर्द्धूतमलः श्लेष्मा यदा व्रजति लोचने~। }\pend
\pstart\versepar \skt{तस्य नेत्रे प्रकाशेते शुक्ल\-म\-ण्ड\-ल\-म\-ण्डले~॥}\pend

\pstart\versepar \ledsidenote{\passref{3.2.35.4}}\skt{रक्तङ्कफयुतं यस्य यदा व्रजति लोचने~॥ }\pend
\pstart\versepar \skt{कपोतनीले नयने रक्ते वा }\edtext{\skt{जायेते}}{\lemma{\skt{जायेते}}\Afootnote{\skt{जायते}\ H}}\skt{ नृणां}\footnoteA{Witness H appears to read \emph{jāyate} but we emend to the dual.}\skt{~॥}\pend

\pstart\versepar \ledsidenote{\passref{3.2.36}}\edtext{\skt{घृतकुम्भो}}{\lemma{\skt{घृतकुम्भो}}\Afootnote{H; \skt{भवन्ति चात्र -- घृतपिण्डो}\ A}}\skt{ }\edtext{\skt{यथा }\edtext{\skt{ह्यग्निमाश्रितः}}{\lemma{\skt{अग्निम्}}\Afootnote{H; \skt{ग्निम्}\ A}}}{\lemma{\skt{यथा ह्य्}}\Afootnote{H; \skt{यथैवा}\ A}}\skt{ प्रविलीयते~। }\pend
\pstart\versepar \skt{}\edtext{\skt{प्रस\-र्प्पे\-दा\-र्त्त\-व\-न्नार्याः}}{\lemma{\skt{प्रस\appbrk{}र्प्पे\appbrk{}दा\appbrk{}र्त्त\appbrk{}वन्}}\Afootnote{H; \skt{विसर्पत्यार्तवं}\ A}}\skt{ }\edtext{\skt{पुरुषस्य}}{\lemma{\skt{पुरुषस्य}}\Afootnote{H; \skt{तथा}\ A}}\skt{ }\edtext{\skt{समागमात्}}{\lemma{\skt{समागमात्}}\Afootnote{H; \skt{पुंसां समागमे}\ A}}\skt{~॥}\pend

\pstart\versepar \ledsidenote{\passref{3.2.37}}\edtext{\skt{बीजे}}{\lemma{\skt{बीजे}}\Afootnote{A H; \skt{बीजे न्}\ L}}\skt{ }\edtext{\skt{तु वायुना}}{\lemma{\skt{तु वायुना}}\Afootnote{H; \skt{ऽन्तर्वायुना}\ A; \skt{तर्वायुना}\ L}}\skt{ भिन्ने द्वौ जीवौ कुक्षिमागतौ~। }\pend
\pstart\versepar \skt{}\edtext{\skt{यमा\-वि\-त्य\-भि\-धी\-येते}}{\lemma{\skt{अभिधीयेते}}\Afootnote{A H; \skt{अभिधीयेत}\ L}}\skt{ }\edtext{\skt{धर्म्मेतरपुरा कृतौ}}{\lemma{\skt{धर्म्मेतरपुरा कृतौ}}\Afootnote{H; \skt{धर्मे\appbrk{}त\appbrk{}र\appbrk{}पु\appbrk{}रः\appbrk{}सरौ}\ A; \skt{थर्म्मे\appbrk{}त\appbrk{}र\appbrk{}पु\appbrk{}र\appbrk{}स्कृतौ}\ L}}\skt{~॥}\pend

\pstart\versepar \ledsidenote{\passref{3.2.37.1}}\skt{मिश्रीभावे यदा पुंसो रेत\-स्सु\-ब\-हु\-नि\-र्म्मलं~। }\pend
\pstart\versepar \skt{तदा प्रसूयते नारी गर्भ्भिणी दारकद्वयं~॥}\pend

\pstart\versepar \ledsidenote{\passref{3.2.37.2}}\skt{मिश्रीभावे स्त्रियाश्चैव बहुशुक्रं यदा भवेत्~। }\pend
\pstart\versepar \skt{कन्याद्वयन्तदा नारी सूयते नात्र संशयः~॥}\pend

\pstart\versepar \ledsidenote{\passref{3.2.38}}\edtext{\skt{आसेक्यो}}{\lemma{\skt{आसेक्यो}}\Afootnote{H L; \emph{om.}\ A}}\skt{ }\edtext{\skt{नाम}}{\lemma{\skt{नाम}}\Afootnote{H L; \emph{om.}\ A}}\skt{ }\edtext{\skt{भवति}}{\lemma{\skt{भवति}}\Afootnote{H L; \emph{om.}\ A}}\skt{ }\edtext{\skt{स्त्रीपुं\-सो\-र\-ल्प\-बी\-जयोः}}{\lemma{\skt{स्त्रीपुंसोर्}}\Afootnote{H L; \skt{पित्रोर्}\ A}}\skt{~॥ }\pend
\pstart\versepar \skt{स }\edtext{\skt{शुक्लं}}{\lemma{\skt{शुक्लं}}\Afootnote{H; \skt{शुक्रं}\ A L}}\skt{ प्राश्य लभते }\edtext{\skt{ध्वजो\-च्छ्रा\-य\-म\-सं\-शयः}}{\lemma{\skt{असंशयः}}\Afootnote{H; \skt{असंशयम्}\ A L\textsuperscript{pc}; \skt{असंशयुयम्}\ L\textsuperscript{ac}}}\skt{~॥}\edtext{}{\lemma{\skt{अल्पबीजयोः~॥ स}}\Afootnote{H L; \skt{अत्य\appbrk{}ल्प\appbrk{}बी\appbrk{}ज\appbrk{}त्वा\appbrk{}दा\appbrk{}सेक्यः पुरुषो भवेत्~। स}\ A}}\pend

\pstart\versepar \ledsidenote{\passref{3.2.39}}\edtext{\skt{यः पूतियोनौ}}{\lemma{\skt{यः पूतियोनौ}}\Afootnote{A H; \skt{यः पूतियोन्याञ्}\ L}}\skt{ जायेत स सौग\-न्धि\-क\-सं\-ज्ञितः~॥ }\pend
\pstart\versepar \skt{स }\edtext{\skt{योनि\-शे\-फ\-सो\-र्ग्ग\-न्ध\-मा\-घ्राय}}{\lemma{\skt{योनिशेफसोर्}}\Afootnote{A H; \skt{योशेफसोर्}\ L}}\skt{ लभते बलम्~॥}\pend

\pstart\versepar \ledsidenote{\passref{3.2.40}}\skt{स्वे गुदे }\edtext{\skt{ब्रह्मचर्याद्यः}}{\lemma{\skt{ब्रह्मचर्याद्यः}}\Afootnote{\skt{ऽब्रह्मचर्याद्यः}\ A; \skt{ब्रह्मचर्याद्या}\ H}}\skt{ स्त्रीषु }\edtext{\skt{स्वासु}}{\lemma{\skt{स्वासु}}\Afootnote{H; \skt{पुंवत्}\ A}}\skt{ प्रवर्तते~। }\pend
\pstart\versepar \skt{कुम्भीकः स तु विज्ञेय ईर्ष्यकं शृणु चापरं~॥}\pend

\pstart\versepar \ledsidenote{\passref{3.2.41}}\skt{दृष्ट्वा व्यवायमन्येषां }\edtext{\skt{व्यवायाय}}{\lemma{\skt{व्यवायाय}}\Afootnote{H; \skt{व्यवाये यः}\ A}}\skt{ प्रवर्त्तते~। }\pend
\pstart\versepar \skt{}\edtext{\skt{ईर्ष्यकः स तु}}{\lemma{\skt{ईर्ष्यकः स तु}}\Afootnote{H; \skt{ईर्ष्यकः स च}\ A}}\skt{ विज्ञेयः }\edtext{\skt{षण्ढकं}}{\lemma{\skt{षण्ढकं}}\Afootnote{H; \skt{षण्डकं}\ A}}\skt{ शृणु पञ्चमं~॥}\pend

\pstart\versepar \ledsidenote{\passref{3.2.42}}\edtext{\skt{कौमार्यां }\edtext{\skt{यस्त्वृतौ}}{\lemma{\skt{त्वृतौ}}\Afootnote{H L; \skt{ऋतौ}\ A}}}{\lemma{\skt{कौमार्यां यस्}}\Afootnote{H; \skt{यो भार्यायाम्}\ A; \skt{भार्य्याय स}\ L}}\skt{ }\edtext{\skt{मोहादङ्गनेव}}{\lemma{\skt{मोहादङ्गनेव}}\Afootnote{A H; \skt{मोहा\appbrk{}दं\appbrk{}ग\appbrk{}ना\appbrk{}वत्\_}\ L}}\skt{ प्रवर्तते~। }\pend
\pstart\versepar \skt{}\edtext{\skt{तत्र}}{\lemma{\skt{तत्र}}\Afootnote{H L; \skt{ततः}\ A}}\skt{ }\edtext{\skt{स्त्रीचेष्टिताकारो}}{\lemma{\skt{स्त्रीचेष्टिताकारो}}\Afootnote{A H; \skt{स्त्रीचेष्टिताचारोः}\ L}}\skt{ जायते }\edtext{\skt{षण्ढसंज्ञकः}}{\lemma{\skt{षण्ढसंज्ञकः}}\Afootnote{H; \skt{षण्ढसंज्ञितः}\ A; \skt{षण्डसंज्ञितः}\ L}}\skt{~॥}\pend

\pstart\versepar \ledsidenote{\passref{3.2.43}}\skt{ऋतौ }\edtext{\skt{पुरुषवच्चापि}}{\lemma{\skt{चापि}}\Afootnote{H L; \skt{वापि}\ A}}\skt{ }\edtext{\skt{यदि}}{\lemma{\skt{यदि}}\Afootnote{H; \emph{om.}\ A L}}\skt{ }\edtext{\edtext{\skt{योषित्प्रपद्यते}}{\lemma{\skt{योषित्}}\Afootnote{H; \emph{om.}\ A L}}}{\lemma{\skt{प्रपद्यते}}\Afootnote{H; \skt{प्रवर्तेताङ्गना यदि}\ A L}}\skt{~। }\pend
\pstart\versepar \skt{तत्र कन्या यदि भवेत्सा भवे\-न्न\-र\-चे\-ष्टिता~॥}\pend

\pstart\versepar \ledsidenote{\passref{3.2.44}}\skt{आसेक्यश्च सुगन्धी च कुम्भी\-क\-श्चे\-र्ष्य\-क\-स्तथा~। }\pend
\pstart\versepar \skt{}\edtext{\skt{सशुक्रकास्त्वमी}}{\lemma{\skt{सशुक्रकास्}}\Afootnote{H; \skt{सरेतसस्}\ A}}\skt{ }\edtext{\skt{ज्ञेया ह्यशुक्रः}}{\lemma{\skt{ज्ञेया ह्य्}}\Afootnote{H; \skt{ज्ञेया}\ A}}\skt{ }\edtext{\skt{षण्ढसंज्ञकः}}{\lemma{\skt{षण्ढसंज्ञकः}}\Afootnote{H; \skt{षण्ढसंज्ञितः}\ A}}\skt{~॥}\pend

\pstart\versepar \ledsidenote{\passref{3.2.45}}\edtext{\skt{अन\-यो\-र्व्वि\-प्र\-कृत्या}}{\lemma{\skt{अनयोर्}}\Afootnote{H; \skt{अनया}\ A}}\skt{ तु तेषां }\edtext{\skt{शुक्रवहा}}{\lemma{\skt{शुक्रवहा}}\Afootnote{H; \skt{शुक्रवहाः}\ A}}\skt{ सिरा~॥ }\pend
\pstart\versepar \skt{}\edtext{\skt{हर्षा\-त्स्फु\-ट\-त्व\-मा\-याति}}{\lemma{\skt{आयाति}}\Afootnote{\skt{आयान्ति}\ A H}}\skt{ }\edtext{\skt{ध्वजोच्छ्रायन्ततो}}{\lemma{\skt{ध्वजोच्छ्रायन्}}\Afootnote{H; \skt{ध्वजोच्छ्रायस्}\ A}}\skt{ भवेत्~॥}\footnoteA{We emend H's reading \emph{āyānti} to the singular. There are two semen-carrying vessels in the Suśrutasaṃhitā's anatomical theory (śārīrasthāna 9.12), not many.}\pend

\pstart\versepar \ledsidenote{\passref{3.2.46}}\edtext{\skt{आका\-रा\-चा\-र\-चि\-न्ता\-भि\-र्या\-दृ\-शीभिः}}{\lemma{\skt{आका\appbrk{}रा\appbrk{}चा\appbrk{}र\appbrk{}चि\appbrk{}न्ता\appbrk{}भिर्}}\Afootnote{H; \skt{आहा\appbrk{}रा\appbrk{}चा\appbrk{}र\appbrk{}चे\appbrk{}ष्टा\appbrk{}भिर्}\ A}}\skt{ समन्वितौ~। }\pend
\pstart\versepar \skt{}\edtext{\skt{स्त्रीपुंसौ तु}}{\lemma{\skt{स्त्रीपुंसौ तु}}\Afootnote{H; \skt{स्त्रीपुंसौ स}\ A}}\skt{ }\edtext{\skt{भवेयातां}}{\lemma{\skt{भवेयातां}}\Afootnote{H; \skt{मुपेयातां}\ A}}\skt{ }\edtext{\skt{तयोर्ग्गर्ब्भोऽपि}}{\lemma{\skt{ग्गर्ब्भो}}\Afootnote{H; \skt{पुत्रो}\ A}}\skt{ तादृशः~॥}\pend

\pstart\versepar \ledsidenote{\passref{3.2.47}}\edtext{\skt{यदा तु}}{\lemma{\skt{यदा तु}}\Afootnote{H; \skt{यदा}\ A}}\skt{ }\edtext{\skt{नारी}}{\lemma{\skt{नारी}}\Afootnote{H; \emph{om.}\ A}}\skt{ }\edtext{\skt{नारी च}}{\lemma{\skt{नारी च}}\Afootnote{H; \skt{नार्याव्}\ A}}\skt{ }\edtext{\skt{मैथु\-ना\-यो\-प\-क\-ल्पतः}}{\lemma{\skt{मैथु\appbrk{}ना\appbrk{}यो\appbrk{}प\appbrk{}क\appbrk{}ल्पतः}}\Afootnote{H; \skt{उपेयातां वृषस्यन्त्यौ कथञ्चन}\ A}}\skt{~। }\pend
\pstart\versepar \skt{}\edtext{\skt{मुञ्चतः}}{\lemma{\skt{मुञ्चतः}}\Afootnote{A; \skt{मुञ्चन्तः}\ H}}\skt{ शुक्र\-म\-न्यो\-न्य\-म\-न\-स्थि\-स्तत्र जायते~॥}\pend

\pstart\versepar \ledsidenote{\passref{3.2.50}}\edtext{\skt{सर्प्प\-वृ\-श्चि\-क\-कु\-ष्मा\-ण्ड\-वि\-कृ\-ता\-कृ\-त\-यस्तु ते}}{\lemma{\skt{सर्प्प\appbrk{}वृ\appbrk{}श्चि\appbrk{}क\appbrk{}कु\appbrk{}ष्मा\appbrk{}ण्ड\appbrk{}वि\appbrk{}कृ\appbrk{}ता\appbrk{}कृ\appbrk{}त\appbrk{}यस्तु ते}}\Afootnote{H; \skt{सर्प\appbrk{}वृ\appbrk{}श्चि\appbrk{}क\appbrk{}कू\appbrk{}ष्मा\appbrk{}ण्ड\appbrk{}वि\appbrk{}कृ\appbrk{}ता\appbrk{}कृ\appbrk{}त\appbrk{}यश्च ये}\ A; \skt{सर्प\appbrk{}वृ\appbrk{}श्चि\appbrk{}क\appbrk{}कूं\appbrk{}भा\appbrk{}ण्ड\appbrk{}वि\appbrk{}कृ\appbrk{}ता\appbrk{}वि\appbrk{}कृ\appbrk{}त\appbrk{}आ\appbrk{}यस्तु ये}\ L\textsuperscript{ac}; \skt{सर्प\appbrk{}वृ\appbrk{}श्चि\appbrk{}क\appbrk{}कूं\appbrk{}भा\appbrk{}ण्ड\appbrk{}वि\appbrk{}कृ\appbrk{}ता\appbrk{}कृ\appbrk{}त\appbrk{}यस्तु ये}\ L\textsuperscript{pc}}}\skt{~। }\pend
\pstart\versepar \skt{}\edtext{\skt{गर्भ्भा}}{\lemma{\skt{गर्भ्भा}}\Afootnote{H; \skt{गर्भास्त्वेते}\ A; \skt{गर्भास्त्वेता}\ L}}\skt{ }\edtext{\edtext{\skt{एवम्विधाश्चान्ये}}{\lemma{\skt{एवम्विधाश्}}\Afootnote{H; \skt{स्त्रियाश्}\ A; \skt{स्त्रीयश्}\ L}}\skt{ ज्ञेयाः}}{\lemma{\skt{चान्ये ज्ञेयाः}}\Afootnote{H; \skt{चैव ज्ञेयाः}\ A L}}\skt{ पापकृतो भृशं~॥}\pend

\pstart\versepar \ledsidenote{\passref{3.2.51}}\edtext{\skt{गर्भ्भो}}{\lemma{\skt{गर्भ्भो}}\Afootnote{H; \skt{गर्भो}\ A L}}\skt{ }\edtext{\skt{वातप्रकोपेण}}{\lemma{\skt{वातप्रकोपेण}}\Afootnote{A H; \skt{वातप्रकोपेन}\ L}}\skt{ }\edtext{\skt{दौहृदेन}}{\lemma{\skt{दौहृदेन}}\Afootnote{H; \skt{दौहृदे वा}\ A; \skt{दौहृदे चा}\ L}}\skt{ }\edtext{\skt{विमानितः}}{\lemma{\skt{विमानितः}}\Afootnote{H; \skt{वमानिते}\ A L}}\skt{~। }\pend
\pstart\versepar \skt{}\edtext{\skt{भवेत्कुब्जोऽथवा}}{\lemma{\skt{ऽथवा}}\Afootnote{H; \emph{om.}\ A L}}\skt{ }\edtext{\skt{कूनिः}}{\lemma{\skt{कूनिः}}\Afootnote{\skt{कुणिः}\ A L; \skt{कूनि}\ H}}\skt{ }\edtext{\edtext{\edtext{\skt{पङ्गुर्म्मूकोऽथ}}{\lemma{\skt{पङ्गुर्}}\Afootnote{A H; \skt{पंकुर्}\ L}}}{\lemma{\skt{म्मूको}}\Afootnote{A H; \skt{म्मूखो}\ L}}}{\lemma{\skt{ऽथ}}\Afootnote{H; \emph{om.}\ A L}}\skt{ }\edtext{\skt{मिर्म्मिणः}}{\lemma{\skt{मिर्म्मिणः}}\Afootnote{H; \skt{मिन्मिण एव वा}\ A; \skt{दुर्म्मन एव वा}\ L}}\skt{~॥}\pend

\pstart\versepar \ledsidenote{\passref{3.2.52}}\edtext{\edtext{\skt{माता\-पि\-त्रो\-र\-ना\-चा\-रा\-द\-शु\-भैश्च}}{\lemma{\skt{मातापित्रोर्}}\Afootnote{A H; \skt{मात्रापित्रोश्}\ L}}}{\lemma{\skt{अना\appbrk{}चा\appbrk{}रा\appbrk{}द\appbrk{}शु\appbrk{}भैश्च}}\Afootnote{H; \skt{तु नास्ति\appbrk{}क्या\appbrk{}द\appbrk{}शु\appbrk{}भैश्च}\ A; \skt{च नास्तिक्यात्\_ अशुभैश्च}\ L}}\skt{ पुराकृतैः~। }\pend
\pstart\versepar \skt{}\edtext{\edtext{\skt{वाता\-दी\-ना\-म्प्र\-को\-पेण}}{\lemma{\skt{वातादीनाम्}}\Afootnote{A H; \skt{वातादिनाञ्}\ L}}}{\lemma{\skt{प्रकोपेण}}\Afootnote{A H; \skt{च कोपेन}\ L}}\skt{ }\edtext{\skt{गर्भ्भो}}{\lemma{\skt{गर्भ्भो}}\Afootnote{H; \skt{गर्भो}\ A; \skt{गर्भ}\ L}}\skt{ }\edtext{\skt{वैकृ\-त\-मा\-प्नु\-यात्}}{\lemma{\skt{आप्नुयात्}}\Afootnote{A H; \skt{उच्यते}\ L}}\skt{~॥}\pend

\pstart\versepar \ledsidenote{\passref{3.2.53}}\skt{मला\-ल्प\-त्वा\-द\-यो\-गाच्च }\edtext{\skt{वायोः}}{\lemma{\skt{वायोः}}\Afootnote{A H\textsuperscript{pc} L; \skt{वायो}\ H\textsuperscript{ac}}}\skt{ }\edtext{\skt{पक्वाशयस्य च}}{\lemma{\skt{पक्वाशयस्य च}}\Afootnote{A H; \skt{पक्वाशयस्य व}\ L}}\skt{~। }\pend
\pstart\versepar \skt{वात\-मू\-त्र\-पु\-री\-षाणि न }\edtext{\skt{गर्भ्भस्थः}}{\lemma{\skt{गर्भ्भस्थः}}\Afootnote{H; \skt{गर्भस्थः}\ A L}}\skt{ करोति }\edtext{\skt{ह}}{\lemma{\skt{ह}}\Afootnote{H L; \skt{हि}\ A}}\skt{~॥}\pend

\pstart\versepar \ledsidenote{\passref{3.2.54}}\edtext{\skt{जरायुना}}{\lemma{\skt{जरायुना}}\Afootnote{H; \skt{जरायुणा}\ A; \skt{जरायणा}\ L}}\skt{ मुखे च्छन्ने कण्ठे च कफवेष्टिते~। }\pend
\pstart\versepar \skt{}\edtext{\edtext{\skt{वायो\-र्ग्ग\-ति\-नि\-रो\-धाच्च}}{\lemma{\skt{वायोर्}}\Afootnote{A H; \skt{वायो}\ L}}\skt{ न}}{\lemma{\skt{ग्गतिनिरोधाच्च न}}\Afootnote{H; \skt{मार्गनिरोधाच्च न}\ A; \skt{मार्गप्ररोधाच्च न}\ L}}\skt{ }\edtext{\skt{गर्भ्भस्थः}}{\lemma{\skt{गर्भ्भस्थः}}\Afootnote{H; \skt{गर्भस्थः}\ A L}}\skt{ }\edtext{\skt{प्ररोदिति}}{\lemma{\skt{प्ररोदिति}}\Afootnote{A H; \skt{प्ररोदति}\ L}}\skt{~॥}\pend

\pstart\versepar \ledsidenote{\passref{3.2.55}}\edtext{\edtext{\edtext{\skt{निश्वा\-सो\-च्छ्वा\-स\-सं\-क्षो\-भ\-स्व\-प्ना\-ङ्ग\-र्भ्भोऽथ}}{\lemma{\skt{निश्वा\appbrk{}सो\appbrk{}च्छ्वा\appbrk{}स\appbrk{}सं\appbrk{}क्षो\appbrk{}भ\appbrk{}स्व\appbrk{}प्नाङ्}}\Afootnote{\skt{निःश्वा\appbrk{}सो\appbrk{}च्छ्वा\appbrk{}स\appbrk{}स\appbrk{}ङ्क्षो\appbrk{}भ\appbrk{}स्व\appbrk{}प्नान्}\ A; \skt{निश्वा\appbrk{}सो\appbrk{}च्छ्वा\appbrk{}स\appbrk{}सं\appbrk{}क्षोभः स्वप्नाङ्}\ H}}}{\lemma{\skt{गर्भ्भो}}\Afootnote{H; \skt{गर्भो}\ A}}}{\lemma{\skt{ऽथ}}\Afootnote{H; \skt{ऽधि}\ A}}\skt{ गच्छति~। }\pend
\pstart\versepar \skt{}\edtext{\skt{मातु\-र्न्निः\-श्व\-सि\-तो\-च्छ्वा\-स\-सं\-क्षो\-भ\-स्व\-प्न\-सं\-भ\-वान्}}{\lemma{\skt{न्निःश्व\appbrk{}सि\appbrk{}तो\appbrk{}च्छ्वा\appbrk{}स\appbrk{}सं\appbrk{}क्षो\appbrk{}भ\appbrk{}स्व\appbrk{}प्न\appbrk{}सं\appbrk{}भ\appbrk{}वान्}}\Afootnote{H; \skt{निश्व\appbrk{}सि\appbrk{}तो\appbrk{}च्छ्वा\appbrk{}स\appbrk{}स\appbrk{}ङ्क्षो\appbrk{}भ\appbrk{}स्व\appbrk{}प्न\appbrk{}सं\appbrk{}भ\appbrk{}वान्}\ A}}\skt{~॥}\pend

\pstart\versepar \ledsidenote{\passref{3.2.56}}\skt{सन्निवेशः शरीराणां दन्तानां }\edtext{\skt{पतनोद्भवौ}}{\lemma{\skt{पतनोद्भवौ}}\Afootnote{A H; \skt{पातनोक्त्\_भवौ}\ L\textsuperscript{ac}; \skt{पातनोत्\_भवौ}\ L\textsuperscript{pc}}}\skt{~। }\pend
\pstart\versepar \skt{तलेष्वसंभवो }\edtext{\skt{यच्च}}{\lemma{\skt{यच्च}}\Afootnote{H L; \skt{यश्च}\ A}}\skt{ }\edtext{\skt{रोम्ना\-मे\-त\-त्स्व\-भा\-वतः}}{\lemma{\skt{रोम्नाम्}}\Afootnote{H; \skt{रोम्णाम्}\ A L}}\skt{~॥}\pend

\pstart\versepar \ledsidenote{\passref{3.2.57}}\edtext{\skt{भाविताः}}{\lemma{\skt{भाविताः}}\Afootnote{A L; \skt{भाविता}\ H}}\skt{ पूर्व्वदेहेषु सततं शास्त्रबुद्धयः~। }\pend
\pstart\versepar \skt{भवन्ति }\edtext{\skt{सत्त्वबहुलाः}}{\lemma{\skt{सत्त्वबहुलाः}}\Afootnote{L; \skt{सत्त्वभूयिष्ठाः}\ A; \skt{स्त्तानबहुलाः}\ H}}\skt{ }\edtext{\skt{पूर्व्वजातिस्मरा}}{\lemma{\skt{पूर्व्वजातिस्मरा}}\Afootnote{A H; \skt{पुरोजातिस्मरा}\ L}}\skt{ नराः~॥ }\footnoteA{Witness H has a strange conjunct consonant after \emph{bhavanti}, followed by \emph{na}. The whole looks like \emph{sttāna}, which is not visually similar to \emph{sattva}. We have followed witness L, which at least make good sense.}\pend

\pstart \ledsidenote{\passref{3.2.trailer}}\edtext{\skt{इति}}{\lemma{\skt{इति}}\Afootnote{H L; \skt{इति सुश्रु\appbrk{}त\appbrk{}सं\appbrk{}हि\appbrk{}तायां शारीरस्थाने शुक्र\appbrk{}शो\appbrk{}णि\appbrk{}त\appbrk{}शुद्धि}\ A}}\skt{ }\edtext{\skt{शारीरं}}{\lemma{\skt{शारीरं}}\Afootnote{A H; \skt{शरीरे}\ L}}\skt{ }\edtext{\edtext{\skt{द्वितीयोऽध्यायः}}{\lemma{\skt{द्वितीयो}}\Afootnote{H; \skt{नाम द्वितीयो}\ A; \skt{द्वितियो}\ L}}}{\lemma{\skt{ऽध्यायः}}\Afootnote{A H; \skt{द्ध्यायः}\ L}}\skt{~॥~॥}\pend

\endnumbering
\beginnumbering
\eledsection*{Adhyāya 3}
\pstart \ledsidenote{\passref{3.3.1}}\skt{अथातो }\edtext{\skt{गर्ब्भावक्रान्तिं}}{\lemma{\skt{गर्ब्भावक्रान्तिं}}\Afootnote{H; \skt{गर्भावक्रान्तिं}\ A; \skt{गर्भावकान्}\ L}}\skt{ }\edtext{\skt{शारीरं}}{\lemma{\skt{शारीरं}}\Afootnote{A H; \skt{तीयम}\ L}}\skt{ व्याख्यास्यामः~॥}\pend

\pstart \ledsidenote{\passref{3.3.3}}\skt{सौम्यं }\edtext{\edtext{\edtext{\skt{शुक्र\-मा\-र्त्त\-व\-मा\-ग्ने\-य\-मि\-त\-रे\-षा\-म\-प्यत्र}}{\lemma{\skt{आर्त्तवम्}}\Afootnote{A H; \skt{आग्नेयम्}\ L}}}{\lemma{\skt{आग्नेयम्}}\Afootnote{A H; \skt{आर्त्तवम्}\ L}}}{\lemma{\skt{अप्य्}}\Afootnote{A H; \skt{अत्र प्य्}\ L}}\skt{ भूतानां }\edtext{\skt{सान्नि\-ध्य\-म\-स्त्य\-णुना}}{\lemma{\skt{अणुना}}\Afootnote{A; \skt{अनुना}\ H; \skt{अणत्रनाम्}\ L}}\skt{ विशेषेण }\edtext{\skt{पर\-स्प\-रो\-प\-का\-रा\-त्प\-र\-स्प\-रा\-नु\-प्र\-वे\-शाच्च}}{\lemma{\skt{पर\appbrk{}स्प\appbrk{}रो\appbrk{}प\appbrk{}का\appbrk{}रात्}}\Afootnote{H L; \skt{पर\appbrk{}स्प\appbrk{}रो\appbrk{}प\appbrk{}का\appbrk{}रात् (प\appbrk{}र\appbrk{}स्प\appbrk{}रा\appbrk{}नु\appbrk{}ग्र\appbrk{}हात्,)}\ A}}\skt{~॥ }\footnoteA{In both Nepalese witnesses, the dental \emph{na} in \emph{anunā} is unambiguous, but we have to emend to the retroflex for grammar and sense. In Suśrutasaṃhitā sūtrasthāna 14.16, a similar phrase, \emph{aṇunā viśeṣeṇa-} has the retroflex in the witnesses K and N, so the emendation is certain.}\pend

\pstart \ledsidenote{\passref{3.3.4}}\edtext{\skt{तत्र}}{\lemma{\skt{तत्र}}\Afootnote{A H; \skt{तत्तु}\ L}}\skt{ }\edtext{\skt{दम्पत्योः}}{\lemma{\skt{दम्पत्योः}}\Afootnote{H; \skt{स्त्रीपुंसयोः}\ A; \skt{शुक\appbrk{}म्स्त्री\appbrk{}पुं\appbrk{}स\appbrk{}योस्}\ L}}\skt{ }\edtext{\skt{सम्प्रयोगे}}{\lemma{\skt{सम्प्रयोगे}}\Afootnote{H L; \skt{संयोगे}\ A}}\skt{ }\edtext{\skt{तेजः}}{\lemma{\skt{तेजः}}\Afootnote{A H; \emph{om.}\ L}}\skt{ }\edtext{\edtext{\edtext{\skt{शरी\-रा\-द्वा\-यु\-स्स\-मी\-र\-यति}}{\lemma{\skt{शरीराद्}}\Afootnote{A H; \emph{om.}\ L}}}{\lemma{\skt{वायुस्}}\Afootnote{A H; \emph{om.}\ L}}}{\lemma{\skt{समीरयति}}\Afootnote{H; \skt{उदीरयति}\ A; \emph{om.}\ L}}\skt{~। }\edtext{\skt{तत्र}}{\lemma{\skt{तत्र}}\Afootnote{H; \skt{ततस्}\ A; \emph{om.}\ L}}\skt{ }\edtext{\edtext{\skt{तेजो\-नि\-ल\-स\-न्नि\-पा\-ता\-च्छुक्रं}}{\lemma{\skt{तेजो\appbrk{}नि\appbrk{}ल\appbrk{}स\appbrk{}न्नि\appbrk{}पा\appbrk{}ताच्}}\Afootnote{A H; \emph{om.}\ L}}}{\lemma{\skt{छुक्रं}}\Afootnote{A H; \emph{om.}\ L}}\skt{ }\edtext{\skt{च्युतं}}{\lemma{\skt{च्युतं}}\Afootnote{A H; \emph{om.}\ L}}\skt{ }\edtext{\edtext{\skt{योनि\-म\-भि\-प्र\-प\-द्यते}}{\lemma{\skt{योनिम्}}\Afootnote{A H; \emph{om.}\ L}}}{\lemma{\skt{अभिप्रपद्यते}}\Afootnote{H; \skt{अभि\appbrk{}प्र\appbrk{}ति\appbrk{}प\appbrk{}द्यते}\ A; \emph{om.}\ L}}\skt{~। }\edtext{\skt{संसृज्यते}}{\lemma{\skt{संसृज्यते}}\Afootnote{A H; \emph{om.}\ L}}\skt{ }\edtext{\skt{चार्त्तवेन}}{\lemma{\skt{चार्त्तवेन}}\Afootnote{A H; \emph{om.}\ L}}\skt{~। }\edtext{\edtext{\edtext{\skt{ततोऽ\-ग्नी\-षो\-म\-सं\-यो\-गा\-त्सं\-सृ\-ज्य\-मानो}}{\lemma{\skt{ततो}}\Afootnote{A H; \emph{om.}\ L}}}{\lemma{\skt{ऽग्नी\appbrk{}षो\appbrk{}म\appbrk{}सं\appbrk{}यो\appbrk{}गात्}}\Afootnote{A H; \emph{om.}\ L}}}{\lemma{\skt{संसृज्यमानो}}\Afootnote{A H; \emph{om.}\ L}}\skt{ }\edtext{\edtext{\skt{गर्ब्भा\-श\-य\-म\-नु\-प्र\-प\-द्यते}}{\lemma{\skt{गर्ब्भाशयम्}}\Afootnote{H; \skt{गर्भाशयम्}\ A; \emph{om.}\ L}}}{\lemma{\skt{अनुप्रपद्यते}}\Afootnote{H; \skt{अनु\appbrk{}प्र\appbrk{}ति\appbrk{}प\appbrk{}द्यते}\ A; \emph{om.}\ L}}\skt{~। }\edtext{\edtext{\skt{क्षेत्त्र\-ज्ञ\-श्चे\-त\-यिता}}{\lemma{\skt{क्षेत्त्रज्ञश्}}\Afootnote{A H; \emph{om.}\ L}}}{\lemma{\skt{चेतयिता}}\Afootnote{H; \skt{वेदयिता}\ A; \emph{om.}\ L}}\skt{ }\edtext{\skt{स्प्रष्टा}}{\lemma{\skt{स्प्रष्टा}}\Afootnote{A H; \emph{om.}\ L}}\skt{ }\edtext{\skt{घ्राता}}{\lemma{\skt{घ्राता}}\Afootnote{A H; \emph{om.}\ L}}\skt{ }\edtext{\skt{द्रष्टा}}{\lemma{\skt{द्रष्टा}}\Afootnote{A H; \emph{om.}\ L}}\skt{ }\edtext{\skt{श्रोता}}{\lemma{\skt{श्रोता}}\Afootnote{A H; \emph{om.}\ L}}\skt{ }\edtext{\skt{रसयिता}}{\lemma{\skt{रसयिता}}\Afootnote{A H; \emph{om.}\ L}}\skt{ }\edtext{\skt{पुरुषो}}{\lemma{\skt{पुरुषो}}\Afootnote{A H; \emph{om.}\ L}}\skt{ }\edtext{\skt{गन्ता}}{\lemma{\skt{गन्ता}}\Afootnote{H; \skt{स्रष्टा गन्ता}\ A; \skt{योनौ र}\ L}}\skt{ }\edtext{\skt{साक्षी}}{\lemma{\skt{साक्षी}}\Afootnote{A H; \emph{om.}\ L}}\skt{ }\edtext{\skt{धाता}}{\lemma{\skt{धाता}}\Afootnote{A H; \emph{om.}\ L}}\skt{ }\edtext{\skt{वक्ता यः }\edtext{\edtext{\skt{कोऽसावाद्य}}{\lemma{\skt{ऽसाव्}}\Afootnote{A H; \skt{सा}\ L}}}{\lemma{\skt{आद्य}}\Afootnote{H; \emph{om.}\ A; \skt{व्रिसु}\ L}}}{\lemma{\skt{वक्ता यः को}}\Afootnote{A H; \skt{ज}\ L}}\skt{ }\edtext{\edtext{\edtext{\edtext{\skt{आयु\-रि\-त्ये\-व\-मा\-दिभिः}}{\lemma{\skt{आयुर्}}\Afootnote{H; \emph{om.}\ A; \skt{ष्}\ L}}}{\lemma{\skt{इत्य्}}\Afootnote{A H; \emph{om.}\ L}}}{\lemma{\skt{एवम्}}\Afootnote{A H; \skt{टम्}\ L}}}{\lemma{\skt{आदिभिः}}\Afootnote{A H; \emph{om.}\ L}}\skt{ }\edtext{\edtext{\edtext{\skt{पर्या\-य\-वा\-च\-कै\-र्न्ना\-म\-भि\-र\-भि\-धी\-यते}}{\lemma{\skt{पर्या\appbrk{}य\appbrk{}वा\appbrk{}च\appbrk{}कैर्}}\Afootnote{A H; \emph{om.}\ L}}}{\lemma{\skt{न्नामभिर्}}\Afootnote{A H; \emph{om.}\ L}}}{\lemma{\skt{अभिधीयते}}\Afootnote{A H; \emph{om.}\ L}}\skt{~। }\edtext{\edtext{\edtext{\skt{दैव\-सं\-यो\-गा\-द\-क्ष\-योऽ\-व्य\-योऽ\-चिन्त्यो}}{\lemma{\skt{दैव\appbrk{}सं\appbrk{}यो\appbrk{}गा\appbrk{}द\appbrk{}क्षयो}}\Afootnote{A H; \skt{अन्वक्षम्}\ L}}}{\lemma{\skt{ऽव्ययो}}\Afootnote{A H; \skt{एव}\ L}}}{\lemma{\skt{ऽचिन्त्यो}}\Afootnote{A H; \emph{om.}\ L}}\skt{ }\edtext{\skt{भूतात्मा}}{\lemma{\skt{भूतात्मा}}\Afootnote{H; \skt{भूतात्मना}\ A; \skt{भूतात्मना च}\ L}}\skt{ }\edtext{\edtext{\skt{सत्त्व\-र\-ज\-स्त\-मो\-भि\-र्द्दे\-वा\-सु\-रै\-र\-प\-रैश्च}}{\lemma{\skt{सत्त्व\appbrk{}र\appbrk{}ज\appbrk{}स्त\appbrk{}मो\appbrk{}भिर्}}\Afootnote{H; \skt{सहान्वक्षं सत्त्व\appbrk{}र\appbrk{}ज\appbrk{}स्त\appbrk{}मो\appbrk{}भिर्}\ A; \skt{सह भूत्वा सत्व\appbrk{}र\appbrk{}ज\appbrk{}स्त\appbrk{}मो\appbrk{}भिर्}\ L}}}{\lemma{\skt{द्देवासुरैर्}}\Afootnote{H L; \skt{दैवासुरैर्}\ A}}\skt{ }\edtext{\skt{सह}}{\lemma{\skt{सह}}\Afootnote{H; \emph{om.}\ A L}}\skt{ भावैर्व्वायुना }\edtext{\skt{प्रेर्यमाणो}}{\lemma{\skt{प्रेर्यमाणो}}\Afootnote{H; \skt{भिप्रेर्यमाणो}\ A; \skt{भिप्रेर्य्यमाणं}\ L}}\skt{ }\edtext{\edtext{\skt{गर्ब्भा\-श\-य\-म\-नु\-प्र\-वि\-शति}}{\lemma{\skt{गर्ब्भाशयम्}}\Afootnote{H; \skt{गर्भाशयम्}\ A L}}}{\lemma{\skt{अनुप्रविशति}}\Afootnote{H; \skt{अनु\appbrk{}प्र\appbrk{}वि\appbrk{}श्या\appbrk{}व\appbrk{}ति\appbrk{}ष्ठते}\ A; \skt{आत्\_\appbrk{}म्या\appbrk{}व\appbrk{}ति\appbrk{}ष्ठते}\ L}}\skt{~॥}\pend

\pstart \ledsidenote{\passref{3.3.5}}\skt{तत्र शुक्र\-बा\-हु\-ल्या\-त्पु\-मा\-ना\-र्त्त\-व\-बा\-हु\-ल्यात्स्त्री }\edtext{\skt{साम्या\-दु\-भ\-यो\-र्न\-पुं\-सकं}}{\lemma{\skt{नपुंसकं}}\Afootnote{H L; \skt{नपुंसकमिति}\ A}}\skt{~॥}\pend

\pstart \ledsidenote{\passref{3.3.6ab}}\edtext{\skt{तत्र}}{\lemma{\skt{तत्र}}\Afootnote{H; \emph{om.}\ A L}}\skt{ }\edtext{\skt{ऋतु\-द्वा\-द\-श\-रात्रं}}{\lemma{\skt{ऋतु\appbrk{}द्वा\appbrk{}द\appbrk{}श\appbrk{}रात्रं}}\Afootnote{H; \skt{ऋतुस्तु द्वादशरात्रं}\ A L}}\skt{ }\edtext{\skt{भवति}}{\lemma{\skt{भवति}}\Afootnote{H; \skt{भवति दृष्टार्तवः}\ A; \skt{दृष्टार्तवो भवति}\ L}}\skt{~।}\pend

\pstart \ledsidenote{\passref{3.3.6.1}}\skt{तत्र प्रथमे दिवसे ऋतुमत्यां मैथु\-न\-ग\-म\-न\-म\-ना\-यु\-ष्य\-म्पुंसो भवति~।}\footnoteA{\emph{puṃso} (m.) \emph{bhavati} needs to be read as a separate sentence because of the gender difference with \emph{anāyuṣyam} (n.).}\skt{ यच्च तत्राधीयते गर्ब्भः सप्रसवतो }\edtext{\skt{विपद्यते~। †द्वितीये}}{\lemma{\skt{विपद्यते~। †द्वितीये}}\Afootnote{\skt{विपद्यते~। द्वितीये}\ H}}\skt{ }\edtext{\skt{गृहे}}{\lemma{\skt{गृहे}}\Afootnote{\skt{तिवागृहे}\ H}}\skt{ }\edtext{\skt{विपद्यते†}}{\lemma{\skt{विपद्यते†}}\Afootnote{\skt{विपद्यते}\ H}}\skt{~। तृती\-येऽ\-प्ये\-व\-म\-स\-म्पू\-र्ण्णोऽ\-ङ्गोऽ\-ल्पा\-युश्च भवति~। तस्मा\-त्त्रि\-रा\-त्र\-म्प\-रि\-ह\-रेत्~। न च प्रवर्त्तति रक्तबीजं प्रदिष्टं गुणकरं भवति~। यथा }\edtext{\skt{नद्याम् †प्र\-ति\-स्रो\-त\-सि\-चि†\-द्रव्यं}}{\lemma{\skt{नद्याम् †प्र\appbrk{}ति\appbrk{}स्रो\appbrk{}त\appbrk{}सि\appbrk{}चि†\appbrk{}द्रव्यं}}\Afootnote{\skt{नद्या\appbrk{}म्प्र\appbrk{}ति\appbrk{}स्रो\appbrk{}तसि प्ल चिद्रव्यं}\ H}}\skt{ प्रतिक्षिप्तं न प्रतिनिवर्त्तते~। तद्वद्बीजं द्रष्टव्यन्~। तस्मा\-त्त्रि\-रा\-त्र\-न्नि\-य\-म\-व\-न्ती\-म्प\-रि\-ह\-रेत्~। तत्र ऋतु\-द्वा\-द\-श\-रा\-त्र\-न्दृ\-ष्ट्वा\-ना\-र्त्तवा भवति~॥}\pend

\pstart \ledsidenote{\passref{3.3.6cd}}\edtext{\edtext{\skt{अदृष्टार्त्तवास्ता}}{\lemma{\skt{अदृष्टार्त्तवास्}}\Afootnote{H; \skt{अदृष्टार्त}\ A; \skt{अदृष्टार्त्तवो}\ L}}\skt{ एके}}{\lemma{\skt{ता एके}}\Afootnote{\skt{वाप्यस्तीत्येके}\ A; \skt{तानेके}\ H; \skt{प्यस्तीति तदर्वागेके}\ L}}\skt{ }\edtext{\skt{भाषन्ते}}{\lemma{\skt{भाषन्ते}}\Afootnote{A H; \skt{हाषन्ते}\ L}}\skt{~॥}\pend

\pstart\versepar \ledsidenote{\passref{3.3.9}}\edtext{\skt{भवति}}{\lemma{\skt{भवति}}\Afootnote{H; \emph{om.}\ A}}\skt{ }\edtext{\skt{चात्र}}{\lemma{\skt{चात्र}}\Afootnote{H; \emph{om.}\ A}}\skt{~॥ }\pend
\pstart\versepar \skt{}\edtext{\skt{दिने}}{\lemma{\skt{दिने}}\Afootnote{H; \emph{om.}\ A}}\skt{ }\edtext{\skt{व्यतीते}}{\lemma{\skt{व्यतीते}}\Afootnote{H; \emph{om.}\ A}}\skt{ }\edtext{\skt{नियमं}}{\lemma{\skt{नियमं}}\Afootnote{H\textsuperscript{pc}; \skt{नियतं}\ A H\textsuperscript{ac}}}\skt{ }\edtext{\skt{संकुचत्यम्बुजं}}{\lemma{\skt{संकुचत्य्}}\Afootnote{H; \skt{दिवसेऽतीते सङ्कुचत्य्}\ A}}\skt{ यथा~। }\pend
\pstart\versepar \skt{ऋतौ व्यतीते नार्यास्तु योनिः }\edtext{\skt{संकुचतिस्तथा}}{\lemma{\skt{संकुचतिस्}}\Afootnote{H; \skt{संव्रियते}\ A}}\skt{~॥}\pend

\pstart\versepar \ledsidenote{\passref{3.3.7}}\edtext{\skt{पीन\-प्र\-स\-न्न\-व\-द\-ना\-म्प्र\-क्लि\-न्ना\-त्म\-मु\-ख\-द्विजां}}{\lemma{\skt{पीन\appbrk{}प्र\appbrk{}स\appbrk{}न्न\appbrk{}व\appbrk{}द\appbrk{}नाम्}}\Afootnote{H; \skt{भवन्ति चात्र पीन\appbrk{}प्र\appbrk{}स\appbrk{}न्न\appbrk{}व\appbrk{}दनां}\ A}}\skt{~। }\pend
\pstart\versepar \skt{नर\-का\-मा\-म्प्रि\-य\-कथां स्रस्त\-कु\-क्ष्य\-क्षि\-मू\-र्द्ध\-जाम्~॥}\pend

\pstart\versepar \ledsidenote{\passref{3.3.8}}\skt{स्फुर\-द्भु\-ज\-कु\-च\-श्रो\-णि\-ना\-भ्यू\-रु\-ज\-घ\-न\-स्फिचां~॥ }\pend
\pstart\versepar \skt{}\edtext{\edtext{\skt{हर्षौ\-त्सु\-क्य\-प\-रां\-श्चैव}}{\lemma{\skt{हर्षौ\appbrk{}त्सु\appbrk{}क्य\appbrk{}प\appbrk{}रांश्}}\Afootnote{H; \skt{हर्षौत्सुक्यपरां}\ A}}}{\lemma{\skt{चैव}}\Afootnote{H; \skt{चापि}\ A}}\skt{ विद्यादृतुमतीं }\edtext{\skt{स्त्रियम्}}{\lemma{\skt{स्त्रियम्}}\Afootnote{H; \skt{इति}\ A}}\skt{~॥}\pend

\pstart\versepar \ledsidenote{\passref{3.3.10}}\skt{मासेनोपचितं काले धमनीभ्यां }\edtext{\skt{यदागतं}}{\lemma{\skt{यदागतं}}\Afootnote{H; \skt{तदार्तवम्}\ A}}\skt{~॥ }\pend
\pstart\versepar \skt{}\edtext{\skt{ईष\-त्कृ\-ष्ण\-म्वि\-ग\-न्धञ्च}}{\lemma{\skt{विगन्धञ्च}}\Afootnote{H; \skt{विवर्णं च}\ A}}\skt{ वायु\-र्यो\-नि\-मु\-ख\-न्न\-येत्~॥}\pend

\pstart\versepar \ledsidenote{\passref{3.3.11}}\skt{तद्व\-र्षा\-द्द्वा\-द\-शा\-त्काले वर्त्त\-मा\-न\-म\-सृ\-क्पुनः~। }\pend
\pstart\versepar \skt{}\edtext{\edtext{\skt{जरा\-प\-क्व\-श\-री\-रा\-णा\-म्प\-ञ्चा\-शा\-द्वि\-नि\-व\-र्त्तते}}{\lemma{\skt{पञ्चाशाद्}}\Afootnote{H; \skt{याति पञ्चाशतः}\ A}}}{\lemma{\skt{विनिवर्त्तते}}\Afootnote{H; \skt{क्षयम्}\ A}}\skt{~॥}\pend

\pstart\versepar \ledsidenote{\passref{3.3.12}}\skt{युग्मेषु तु पुमान्प्रोक्तो दिव\-से\-ष्व\-न्य\-था\-बला~। }\pend
\pstart\versepar \skt{पुष्पकाले शुचि\-स्त\-स्मा\-द\-प\-त्यार्थी स्त्रियं व्रजेत्~॥}\pend

\pstart \ledsidenote{\passref{3.3.13}}\skt{तत्र }\edtext{\skt{सद्यो\-गृ\-ही\-त\-ग\-र्ब्भायाः}}{\lemma{\skt{सद्यो\appbrk{}गृ\appbrk{}ही\appbrk{}त\appbrk{}ग\appbrk{}र्ब्भायाः}}\Afootnote{H; \skt{सद्यो\appbrk{}गृ\appbrk{}ही\appbrk{}त\appbrk{}ग\appbrk{}र्भाया लिङ्गानि}\ A; \skt{सद्यो\appbrk{}ग्र\appbrk{}ही\appbrk{}त\appbrk{}ग\appbrk{}र्भायाः}\ L}}\skt{ श्रमो }\edtext{\skt{ग्लानिः}}{\lemma{\skt{ग्लानिः}}\Afootnote{A H; \skt{र्गानिः}\ L}}\skt{ पिपासा }\edtext{\edtext{\skt{सक्थि\-स\-द\-न\-मा\-टोपः}}{\lemma{\skt{सक्थिसदनम्}}\Afootnote{A H; \skt{सक्थिसीदनम्}\ L}}}{\lemma{\skt{आटोपः}}\Afootnote{H; \emph{om.}\ A; \skt{आटोप}\ L}}\skt{ }\edtext{\edtext{\skt{शुक्र\-शो\-णि\-त\-यो\-र\-नु\-बन्धः}}{\lemma{\skt{शुक्र\appbrk{}शो\appbrk{}णि\appbrk{}त\appbrk{}योर्}}\Afootnote{A H; \skt{शुक\appbrk{}शो\appbrk{}णि\appbrk{}त\appbrk{}योर्}\ L}}}{\lemma{\skt{अनुबन्धः}}\Afootnote{H; \skt{अवबन्धः}\ A; \skt{अननुबन्धः}\ L}}\skt{ }\edtext{\skt{खरस्फुरणं}}{\lemma{\skt{खरस्फुरणं}}\Afootnote{H; \skt{स्फुरणं च}\ A; \skt{गुरुता स्फुरणञ्च}\ L}}\skt{ }\edtext{\skt{योन्याः}}{\lemma{\skt{योन्याः}}\Afootnote{H; \skt{योनेः}\ A; \skt{योन्या भवन्ति}\ L}}\skt{~॥}\pend

\pstart\versepar \ledsidenote{\passref{3.3.14ab}}\edtext{\skt{भवति}}{\lemma{\skt{भवति}}\Afootnote{H; \emph{om.}\ A L}}\skt{ }\edtext{\skt{चात्र}}{\lemma{\skt{चात्र}}\Afootnote{H; \emph{om.}\ A; \skt{भकः}\ L}}\skt{~॥ }\pend
\pstart\versepar \skt{स्तनयोः कृष्णमुखता }\edtext{\skt{रोम\-रा\-ज्यु\-द्ग\-म\-स्तथा}}{\lemma{\skt{रोम\appbrk{}रा\appbrk{}ज्यु\appbrk{}द्ग\appbrk{}मस्}}\Afootnote{A; \skt{रोम\appbrk{}रा\appbrk{}ज्यु\appbrk{}ङ्ग\appbrk{}मस्}\ H; \skt{रोम\appbrk{}रा\appbrk{}ज्यु\appbrk{}त्भ\appbrk{}वस्}\ L}}\skt{~।}\pend

\pstart \ledsidenote{\passref{3.3.15cd}}\skt{प्रसेकः }\edtext{\skt{सदनञ्चापि}}{\lemma{\skt{सदनञ्}}\Afootnote{A H; \skt{सीदनञ्}\ L}}\skt{ }\edtext{\skt{गर्ब्भिण्या}}{\lemma{\skt{गर्ब्भिण्या}}\Afootnote{H; \skt{गर्भिण्या}\ A L}}\skt{ }\edtext{\skt{लिङ्गमुच्यते}}{\lemma{\skt{उच्यते}}\Afootnote{A H; \skt{उद्ध्यते}\ L}}\skt{~॥}\pend

\pstart \ledsidenote{\passref{3.3.16}}\edtext{\skt{ततः}}{\lemma{\skt{ततः}}\Afootnote{H L; \skt{तदा}\ A}}\skt{ }\edtext{\skt{प्रभृत्येव च}}{\lemma{\skt{एव च}}\Afootnote{H L; \emph{om.}\ A}}\skt{ व्यवायं }\edtext{\skt{व्याया\-म\-म\-ति\-क\-र्शनं}}{\lemma{\skt{अतिकर्शनं}}\Afootnote{\skt{अति\appbrk{}त\appbrk{}र्प\appbrk{}ण\appbrk{}म\appbrk{}ति\appbrk{}क\appbrk{}र्शनं}\ A; \skt{अतिकर्षणं}\ H; \skt{अपकर्}\ L}}\skt{ }\edtext{\skt{दिवास्वप्नं}}{\lemma{\skt{दिवास्वप्नं}}\Afootnote{A H; \emph{om.}\ L}}\skt{ }\edtext{\skt{रात्रिजागरणं}}{\lemma{\skt{रात्रिजागरणं}}\Afootnote{A H; \emph{om.}\ L}}\skt{ }\edtext{\skt{शोकं}}{\lemma{\skt{शोकं}}\Afootnote{A H; \emph{om.}\ L}}\skt{ }\edtext{\edtext{\edtext{\edtext{\skt{याना\-रो\-ह\-ण\-म्भ\-य\-मु\-त्क\-ट\-का\-स\-न\-ञ्चै\-का\-न्ततः}}{\lemma{\skt{यानारोहणम्}}\Afootnote{A H; \emph{om.}\ L}}}{\lemma{\skt{भयम्}}\Afootnote{A H; \emph{om.}\ L}}}{\lemma{\skt{उत्क\appbrk{}ट\appbrk{}का\appbrk{}स\appbrk{}नञ्}}\Afootnote{H; \skt{उत्कुटुकासनं}\ A; \skt{शनम्}\ L}}}{\lemma{\skt{चैकान्ततः}}\Afootnote{A H; \emph{om.}\ L}}\skt{ }\edtext{\skt{स्नेहादिक्रमं}}{\lemma{\skt{स्नेहादिक्रमं}}\Afootnote{H L; \skt{स्नेहादिक्रियां}\ A}}\skt{ }\edtext{\edtext{\skt{शोणि\-त\-मो\-क्ष\-ण\-ञ्चा\-काले}}{\lemma{\skt{शोणि\appbrk{}त\appbrk{}मो\appbrk{}क्ष\appbrk{}णञ्}}\Afootnote{A H; \skt{शोणितमोक्ष}\ L}}}{\lemma{\skt{चाकाले}}\Afootnote{A H; \emph{om.}\ L}}\skt{ }\edtext{\skt{न}}{\lemma{\skt{न}}\Afootnote{\skt{वेगविधारणं च न}\ A; \skt{ना}\ H; \skt{णञ्च~। न}\ L}}\skt{ सेवेत~।}\footnoteA{We emend \emph{atikarṣaṇaṃ} in witness H for sense and because of the instability of \emph{ṣa} and \emph{ṇa} in H.}\pend

\pstart\versepar \ledsidenote{\passref{3.3.18}}\edtext{\skt{तत्र}}{\lemma{\skt{तत्र}}\Afootnote{A H; \skt{ततः}\ L}}\skt{ प्रथमे }\edtext{\skt{मासे}}{\lemma{\skt{मासे}}\Afootnote{H; \skt{मासि}\ A L}}\skt{ }\edtext{\skt{कललञ्जायते}}{\lemma{\skt{कललञ्}}\Afootnote{A H; \skt{क्लेदो}\ L}}\skt{~॥ }\pend
\pstart\versepar \skt{द्वितीये }\edtext{\skt{शोणितोष्णानिलैः}}{\lemma{\skt{शोणितोष्णानिलैः}}\Afootnote{H; \skt{शीतोष्मानिलैर्}\ A; \skt{मासि~। शीतोष्णानिलैर्}\ L}}\skt{ }\edtext{\skt{पच्यमानो}}{\lemma{\skt{पच्यमानो}}\Afootnote{H; \skt{अभि\appbrk{}प्र\appbrk{}प\appbrk{}च्य\appbrk{}मा\appbrk{}नानां}\ A; \skt{अभिपच्यमानो}\ L}}\skt{ महाभूतानां }\edtext{\skt{संजातो}}{\lemma{\skt{संजातो}}\Afootnote{H; \skt{संघातो}\ A; \skt{घ}\ L}}\skt{ }\edtext{\skt{घनञ्जायते}}{\lemma{\skt{घनञ्}}\Afootnote{H; \skt{घनः सञ्}\ A; \skt{नपिण्डस्सञ्}\ L}}\skt{~। }\edtext{\edtext{\skt{ग्रन्थिर्यदि}}{\lemma{\skt{ग्रन्थिर्}}\Afootnote{H; \emph{om.}\ A L}}}{\lemma{\skt{यदि}}\Afootnote{A H; \emph{om.}\ L}}\skt{ }\edtext{\skt{पुमान्}}{\lemma{\skt{पुमान्}}\Afootnote{H; \skt{पिण्डः पुमान्}\ A; \skt{यदि पुमान्\_}\ L}}\skt{~। }\edtext{\skt{स्त्री}}{\lemma{\skt{स्त्री}}\Afootnote{A H; \skt{स्त्रि}\ L}}\skt{ चेत्पेशी~। नपुंसकं }\edtext{\edtext{\skt{चेद\-र्ब्बु\-द\-म्भ\-वति}}{\lemma{\skt{चेदर्ब्बुदम्}}\Afootnote{A H; \skt{चेदबुदन्}\ L}}}{\lemma{\skt{भवति}}\Afootnote{H; \skt{इति}\ A; \emph{om.}\ L}}\skt{~॥ }\pend
\pstart\versepar \skt{तृतीये हस्त\-पा\-द\-शि\-र\-सा\-म्पञ्च }\edtext{\skt{पिटका}}{\lemma{\skt{पिटका}}\Afootnote{H; \skt{पिण्डका}\ A; \skt{पिण्डकानाम्}\ L}}\skt{ }\edtext{\skt{अभिनिर्वर्त्तन्ते}}{\lemma{\skt{अभिनिर्वर्त्तन्ते}}\Afootnote{H; \skt{निर्वर्तन्ते}\ A; \skt{अभिनिर्वृ}\ L}}\skt{~। }\edtext{\skt{सर्व्वा\-ङ्ग\-प्र\-त्य\-ङ्ग\-वि\-भा\-गश्च}}{\lemma{\skt{सर्व्वा\appbrk{}ङ्ग\appbrk{}प्र\appbrk{}त्य\appbrk{}ङ्ग\appbrk{}वि\appbrk{}भा\appbrk{}गश्च}}\Afootnote{H; \skt{ऽङ्ग\appbrk{}प्र\appbrk{}त्य\appbrk{}ङ्ग\appbrk{}वि\appbrk{}भा\appbrk{}गश्च}\ A; \skt{त्ति अंग\appbrk{}प्र\appbrk{}त्यं\appbrk{}ग\appbrk{}वि\appbrk{}भा\appbrk{}गश्च}\ L}}\skt{ सूक्ष्मो भवति~॥ }\pend
\pstart\versepar \skt{चतुर्थे }\edtext{\skt{सर्व्वा\-ङ्ग\-प्र\-त्य\-ङ्ग\-वि\-भा\-गश्च}}{\lemma{\skt{सर्व्वा\appbrk{}ङ्ग\appbrk{}प्र\appbrk{}त्य\appbrk{}ङ्ग\appbrk{}वि\appbrk{}भा\appbrk{}गश्च}}\Afootnote{H; \emph{om.}\ A L}}\skt{ }\edtext{\skt{प्रव्यक्तो}}{\lemma{\skt{प्रव्यक्तो}}\Afootnote{H; \emph{om.}\ A L}}\skt{ }\edtext{\skt{भवति}}{\lemma{\skt{भवति}}\Afootnote{H; \emph{om.}\ A L}}\skt{~॥ }\pend
\pstart\versepar \skt{}\edtext{\skt{पञ्चमे}}{\lemma{\skt{पञ्चमे}}\Afootnote{H; \emph{om.}\ A L}}\skt{ }\edtext{\skt{सर्व्वा\-ङ्ग\-प्र\-त्य\-ङ्ग\-वि\-भा\-गश्च}}{\lemma{\skt{सर्व्वा\appbrk{}ङ्ग\appbrk{}प्र\appbrk{}त्य\appbrk{}ङ्ग\appbrk{}वि\appbrk{}भा\appbrk{}गश्च}}\Afootnote{H L; \skt{सर्वा\appbrk{}ङ्ग\appbrk{}प्र\appbrk{}त्य\appbrk{}ङ्ग\appbrk{}वि\appbrk{}भागः}\ A}}\skt{ }\edtext{\skt{प्रव्यक्ततरो}}{\lemma{\skt{प्रव्यक्ततरो}}\Afootnote{H L; \skt{प्रव्यक्तो}\ A}}\skt{ }\edtext{\skt{भवति}}{\lemma{\skt{भवति}}\Afootnote{A H; \emph{om.}\ L}}\skt{~। }\edtext{\edtext{\skt{गर्ब्भ\-हृ\-द\-प्र\-व्य\-क्ति\-भा\-वा\-च्चे\-त\-ना\-धा\-तु\-र\-भि\-व्यक्तो}}{\lemma{\skt{गर्ब्भ\appbrk{}हृ\appbrk{}द\appbrk{}प्र\appbrk{}व्य\appbrk{}क्ति\appbrk{}भा\appbrk{}वाच्}}\Afootnote{H; \skt{गर्भ\appbrk{}हृ\appbrk{}द\appbrk{}य\appbrk{}प्र\appbrk{}व्य\appbrk{}क्ति\appbrk{}भा\appbrk{}वाच्}\ A; \skt{गर्भ\appbrk{}हृ\appbrk{}द\appbrk{}य\appbrk{}प्र\appbrk{}व्य\appbrk{}क्ति\appbrk{}हा\appbrk{}वाच्}\ L}}}{\lemma{\skt{अभिव्यक्तो}}\Afootnote{A; \skt{अभिव्यक्}\ H; \skt{अभिबुद्धतरो}\ L}}\skt{ }\footnoteA{We emend H's reading \emph{abhivyakte} to \emph{abhivyakto} for sense and following the vulgate. MS H, our only witness to the Nepalese version at this point, is badly damaged here and while it looks as if the akṣara is \emph{kte}, it is just possible that a vertical line after the character has been destroyed. The damaged portion may have supported the word \emph{bhavati}}\skt{~। }\edtext{\skt{कस्मात्}}{\lemma{\skt{कस्मात्}}\Afootnote{\skt{भवति, कस्मात्}\ A; \skt{ते … कस्मात्}\ H; \skt{भवति}\ L}}\skt{~। }\edtext{\skt{तत्स्थानत्वाद्}}{\lemma{\skt{तत्स्थानत्वाद्}}\Afootnote{H; \skt{तत्स्था\appbrk{}न\appbrk{}त्वा\appbrk{}त्त\appbrk{}स्माद्}\ A L}}\skt{~। }\edtext{\skt{गर्ब्भश्चतुर्थे}}{\lemma{\skt{गर्ब्भश्}}\Afootnote{H; \skt{गर्भश्}\ A L}}\skt{ }\edtext{\skt{मासे}}{\lemma{\skt{मासे}}\Afootnote{H; \skt{मास्य्}\ A L}}\skt{ }\edtext{\skt{अभि\-प्रा\-य\-मि\-न्द्रि\-या\-र्थेषु}}{\lemma{\skt{अभिप्रायम्}}\Afootnote{A H; \skt{आभिप्रायम्}\ L}}\skt{ करोति~।}\footnoteA{Note the absence of sandhi between \emph{māse} and \emph{abhiprāyam}.}\skt{ }\edtext{\skt{द्विहृदया च}}{\lemma{\skt{द्विहृदया च}}\Afootnote{H; \skt{द्विहृदयां च}\ A L}}\skt{ }\edtext{\skt{नारी}}{\lemma{\skt{नारी}}\Afootnote{H; \skt{नारीं}\ A L}}\skt{ }\edtext{\skt{भवति}}{\lemma{\skt{भवति}}\Afootnote{\skt{दौहृदिनीम्}\ A; \skt{भवति …}\ H; \skt{दौहृदि}\ L}}\skt{ }\edtext{\skt{निमि\-त्ता\-न्या\-च\-क्षते}}{\lemma{\skt{निमित्तान्य्}}\Afootnote{H; \emph{om.}\ A; \skt{नीत्य्}\ L}}\skt{~। }\edtext{\edtext{\skt{द्वौहृ\-द\-य\-वि\-मा\-न\-ना\-त्कु\-ब्ज\-कूनिं}}{\lemma{\skt{द्वौहृ\appbrk{}द\appbrk{}य\appbrk{}वि\appbrk{}मा\appbrk{}न\appbrk{}नात्}}\Afootnote{\skt{दौहृ\appbrk{}द\appbrk{}वि\appbrk{}मा\appbrk{}न\appbrk{}नात्}\ A; \skt{द्वौहृ\appbrk{}द\appbrk{}य\appbrk{}वि\appbrk{}मा\appbrk{}न\appbrk{}नान्}\ H; \emph{om.}\ L}}}{\lemma{\skt{कुब्जकूनिं}}\Afootnote{\skt{कुब्जं कुणिं}\ A; \skt{जव्जकू}\ H; \emph{om.}\ L}}\skt{ }\edtext{\edtext{\edtext{\skt{षण्ड\-म्वा\-म\-न\-वि\-कृ\-ता\-क्ष\-म\-न\-क्षम्वा}}{\lemma{\skt{षण्डम्}}\Afootnote{\skt{खञ्जं जडं}\ A; \skt{नीं षण्डम्}\ H; \emph{om.}\ L}}}{\lemma{\skt{वाम\appbrk{}न\appbrk{}वि\appbrk{}कृ\appbrk{}ता\appbrk{}क्षम्}}\Afootnote{H; \skt{वामनं विकृताक्ष}\ A; \emph{om.}\ L}}}{\lemma{\skt{अनक्षम्वा}}\Afootnote{H; \skt{मनक्षं वा}\ A; \emph{om.}\ L}}\skt{ }\edtext{\skt{नारी}}{\lemma{\skt{नारी}}\Afootnote{A H; \emph{om.}\ L}}\skt{ }\edtext{\edtext{\skt{सुतञ्जनयति}}{\lemma{\skt{सुतञ्}}\Afootnote{A H; \emph{om.}\ L}}}{\lemma{\skt{जनयति}}\Afootnote{A H; \emph{om.}\ L}}\skt{~। }\edtext{\skt{तस्मात्सा}}{\lemma{\skt{तस्मात्सा}}\Afootnote{A H; \emph{om.}\ L}}\skt{ }\edtext{\edtext{\edtext{\skt{यदिच्छेत्तत्तस्मै}}{\lemma{\skt{यदिच्छेत्}}\Afootnote{H; \skt{यद्यदिच्छेत्}\ A; \emph{om.}\ L}}}{\lemma{\skt{तत्}}\Afootnote{A H; \emph{om.}\ L}}}{\lemma{\skt{तस्मै}}\Afootnote{H; \skt{तस्यै}\ A; \emph{om.}\ L}}\skt{ }\edtext{\skt{दापयेत्}}{\lemma{\skt{दापयेत्}}\Afootnote{A H; \emph{om.}\ L}}\skt{~। }\edtext{\skt{लब्धदौहृदा हि}}{\lemma{\skt{लब्धदौहृदा हि}}\Afootnote{A H; \emph{om.}\ L}}\skt{ }\edtext{\skt{वीर्यवन्तं}}{\lemma{\skt{वीर्यवन्तं}}\Afootnote{A H; \emph{om.}\ L}}\skt{ }\edtext{\edtext{\skt{चिरा\-यु\-ष\-ञ्ज\-न\-यति}}{\lemma{\skt{चिरायुषञ्}}\Afootnote{H; \skt{चिरायुषं च}\ A; \emph{om.}\ L}}}{\lemma{\skt{जनयति}}\Afootnote{H; \skt{पुत्रं जनयति}\ A; \emph{om.}\ L}}\skt{~॥}\pend

\pstart\versepar \ledsidenote{\passref{3.3.19}}\edtext{\skt{भवति}}{\lemma{\skt{भवति}}\Afootnote{H; \skt{भवन्ति}\ A}}\skt{ चात्र~॥ }\pend
\pstart\versepar \skt{इन्द्रियार्थांस्तु }\edtext{\skt{यान्यां सा}}{\lemma{\skt{यान्यां सा}}\Afootnote{H; \skt{यान्यान्सा}\ A}}\skt{ भोक्तुमिच्छति }\edtext{\skt{गर्ब्भिणी}}{\lemma{\skt{गर्ब्भिणी}}\Afootnote{H; \skt{गर्भिणी}\ A}}\skt{~। }\pend
\pstart\versepar \skt{}\edtext{\skt{गर्ब्भा\-बा\-ध\-भ\-या\-त्तां\-स्ता\-म्भि\-ष\-गा\-हृत्य}}{\lemma{\skt{गर्ब्भा\appbrk{}बा\appbrk{}ध\appbrk{}भ\appbrk{}या\appbrk{}त्तां\appbrk{}स्ताम्}}\Afootnote{H; \skt{गर्भा\appbrk{}बा\appbrk{}ध\appbrk{}भ\appbrk{}या\appbrk{}त्तां\appbrk{}स्तान्}\ A}}\skt{ दापयेत्~॥}\pend

\pstart\versepar \ledsidenote{\passref{3.3.20}}\edtext{\skt{सम्प्राप्तदौहृदा}}{\lemma{\skt{सम्प्राप्तदौहृदा}}\Afootnote{\skt{सा प्राप्तदौहृदा}\ A; \skt{सम्प्रा\appbrk{}प्त\appbrk{}दौ\appbrk{}हृ\appbrk{}दम्}\ H}}\skt{ }\edtext{\skt{पुत्रम्प्रजायेत}}{\lemma{\skt{प्रजायेत}}\Afootnote{H; \skt{जनयेत}\ A}}\skt{ गुणान्वितं~। }\pend
\pstart\versepar \skt{अलब्धदौहृदा }\edtext{\skt{गर्ब्भे}}{\lemma{\skt{गर्ब्भे}}\Afootnote{H; \skt{गर्भे}\ A}}\skt{ }\edtext{\skt{लभते}}{\lemma{\skt{लभते}}\Afootnote{H; \skt{लभे}\ A}}\skt{ }\edtext{\skt{चात्मनि वा}}{\lemma{\skt{चात्मनि वा}}\Afootnote{H; \skt{तात्मनि वा}\ A}}\skt{ भयम्~॥}\footnoteA{We emend to the feminine \emph{samprāptadauhṛdā} for sense and because of symmetry with the second half of the verse.}\pend

\pstart\versepar \ledsidenote{\passref{3.3.21}}\skt{येषु येष्विन्द्रियार्थेषु }\edtext{\skt{दौहृदे या}}{\lemma{\skt{दौहृदे या}}\Afootnote{H; \skt{दौहृदे वै}\ A}}\skt{ }\edtext{\skt{विमानिता}}{\lemma{\skt{विमानिता}}\Afootnote{H; \skt{विमानना}\ A}}\skt{~। }\pend
\pstart\versepar \skt{प्रजायेत }\edtext{\skt{सुतं}}{\lemma{\skt{सुतं}}\Afootnote{H; \skt{सुत}\ A}}\skt{ }\edtext{\edtext{\skt{सार्त्ति\-न्त\-स्मि\-न्त\-स्मि\-न्त\-थे\-न्द्रिये}}{\lemma{\skt{सार्त्तिन्}}\Afootnote{H; \skt{स्यार्तिस्}\ A}}}{\lemma{\skt{तथेन्द्रिये}}\Afootnote{A; \skt{तथैन्द्रिये}\ H}}\skt{~॥}\pend

\pstart\versepar \ledsidenote{\passref{3.3.22}}\edtext{\skt{राजसन्दर्शने}}{\lemma{\skt{राजसन्दर्शने}}\Afootnote{A; \skt{राजा सन्दर्शने}\ H}}\skt{ यस्या }\edtext{\skt{दौहृदञ्जायते}}{\lemma{\skt{दौहृदञ्}}\Afootnote{H; \skt{दौर्हृदं}\ A}}\skt{ }\edtext{\skt{स्त्रियाः}}{\lemma{\skt{स्त्रियाः}}\Afootnote{A; \skt{स्त्रियः}\ H}}\skt{~। }\pend
\pstart\versepar \skt{अर्थ\-व\-न्त\-म्म\-हा\-भागं कुमारं सा प्रसूयते~॥}\footnoteA{We emend to \emph{rāja-} and \emph{striyāḥ} for sense, following the vulgate.}\pend

\pstart\versepar \ledsidenote{\passref{3.3.23}}\edtext{\skt{दुकू\-ल\-प\-ट्ट\-कौ\-षे\-य\-भू\-ष\-णा\-दिषु}}{\lemma{\skt{दुकू\appbrk{}ल\appbrk{}प\appbrk{}ट्ट\appbrk{}कौ\appbrk{}षे\appbrk{}य\appbrk{}भू\appbrk{}ष\appbrk{}णा\appbrk{}दिषु}}\Afootnote{H; \skt{दुकू\appbrk{}ल\appbrk{}प\appbrk{}ट्ट\appbrk{}कौ\appbrk{}शे\appbrk{}य\appbrk{}भू\appbrk{}ष\appbrk{}णा\appbrk{}दिषु}\ A}}\skt{ दौहृदा~। }\pend
\pstart\versepar \skt{}\edtext{\skt{अल\-ङ्का\-रै\-षि\-ण\-म्पुत्रं}}{\lemma{\skt{अल\appbrk{}ङ्का\appbrk{}रै\appbrk{}षि\appbrk{}णम्}}\Afootnote{\skt{त्~। अलङ्कारैषिणं}\ A; \skt{अल\appbrk{}ङ्का\appbrk{}रे\appbrk{}षि\appbrk{}णम्}\ H}}\skt{ ललितं सा प्रसूयते~॥}\pend

\pstart\versepar \ledsidenote{\passref{3.3.24}}\skt{आश्रमे संयतात्मानं धर्म्मशीलं }\edtext{\skt{प्रजायते}}{\lemma{\skt{प्रजायते}}\Afootnote{H; \skt{प्रसूयते}\ A}}\skt{~। }\pend
\pstart\versepar \skt{देव\-ता\-प्र\-ति\-मायां तु }\edtext{\skt{प्रसूते}}{\lemma{\skt{प्रसूते}}\Afootnote{A; \skt{प्रसूतो}\ H}}\skt{ }\edtext{\skt{हर्षदोपमम्}}{\lemma{\skt{हर्षदोपमम्}}\Afootnote{H; \skt{पार्षदोपमम्}\ A}}\skt{~॥ }\pend
\pstart\versepar \skt{दर्शने }\edtext{\skt{व्याडजातीनां}}{\lemma{\skt{व्याडजातीनां}}\Afootnote{H; \skt{व्यालजातीनां}\ A}}\skt{ हिंसाशीलं }\edtext{\skt{प्रजायते}}{\lemma{\skt{प्रजायते}}\Afootnote{H; \skt{प्रसूयते}\ A}}\skt{~॥}\pend

\pstart\versepar \ledsidenote{\passref{3.3.25}}\edtext{\skt{गोधामांसाशिनी}}{\lemma{\skt{गोधामांसाशिनी}}\Afootnote{H; \skt{गोधामांसाशने}\ A}}\skt{ पुत्रं }\edtext{\edtext{\skt{सुषु\-प्सु\-र्म्मा\-र\-णा\-त्मकं}}{\lemma{\skt{सुषुप्सुर्}}\Afootnote{H; \skt{सुषुप्सुं}\ A}}}{\lemma{\skt{म्मारणात्मकं}}\Afootnote{H; \skt{धारणात्मकम्}\ A}}\skt{~॥ }\pend
\pstart\versepar \skt{गवां }\edtext{\skt{मांसेन}}{\lemma{\skt{मांसेन}}\Afootnote{H; \skt{मांसे तु}\ A}}\skt{ बलिनं सर्व्वक्लेशसहं }\edtext{\skt{सुतं}}{\lemma{\skt{सुतं}}\Afootnote{H; \skt{तथा}\ A}}\skt{~।}\pend

\pstart \ledsidenote{\passref{3.3.26ab}}\skt{माहिषे }\edtext{\skt{दौर्हृदाच्छूरं}}{\lemma{\skt{दौर्हृदाच्}}\Afootnote{A; \skt{दौहृदाच्}\ H}}\skt{ रक्ताक्षं }\edtext{\skt{लोमशं}}{\lemma{\skt{लोमशं}}\Afootnote{H; \skt{लोमसं}\ A}}\skt{ }\edtext{\skt{सुतं}}{\lemma{\skt{सुतं}}\Afootnote{H; \skt{युतम्}\ A}}\skt{~॥}\pend

\pstart\versepar \ledsidenote{\passref{3.3.28}}\edtext{\skt{अतोऽनुक्तेषु या}}{\lemma{\skt{ऽनुक्तेषु या}}\Afootnote{A; \skt{नुक्तेषु या}\ H}}\skt{ }\edtext{\edtext{\skt{नारी\-म\-भि\-ध्या\-यति}}{\lemma{\skt{नारीम्}}\Afootnote{H; \skt{नारी}\ A}}}{\lemma{\skt{अभिध्यायति}}\Afootnote{H; \skt{समभिध्याति}\ A}}\skt{ }\edtext{\skt{दौहृदम्}}{\lemma{\skt{दौहृदम्}}\Afootnote{H; \skt{दौर्हृदम्}\ A}}\skt{~॥ }\pend
\pstart\versepar \skt{}\edtext{\skt{शरी\-रा\-हा\-र\-शीलैः सा}}{\lemma{\skt{शरी\appbrk{}रा\appbrk{}हा\appbrk{}र\appbrk{}शीलैः सा}}\Afootnote{H; \skt{शरी\appbrk{}रा\appbrk{}चा\appbrk{}र\appbrk{}शीलैः सा}\ A}}\skt{ समा\-न\-ञ्ज\-न\-यि\-ष्यति~॥}\pend

\pstart\versepar \ledsidenote{\passref{3.3.29}}\skt{कर्म्मणा }\edtext{\edtext{\skt{चोदि\-त\-ञ्ज\-न्तो\-र्ब्भ\-वि\-त\-व्य\-म्पु\-न\-र्ब्भ\-वेत्}}{\lemma{\skt{ब्भवितव्यम्}}\Afootnote{H; \skt{भवितव्यं}\ A}}}{\lemma{\skt{ब्भवेत्}}\Afootnote{H; \skt{भवेत्}\ A}}\skt{~। }\pend
\pstart\versepar \skt{यथा तथा }\edtext{\skt{दैव\-यो\-गा\-द्दौ\-हृ\-द\-ञ्ज\-न\-ये\-द्धृदि}}{\lemma{\skt{दौहृदञ्}}\Afootnote{H; \skt{दौर्हृदं}\ A}}\skt{~॥}\pend

\pstart \ledsidenote{\passref{3.3.30}}\skt{पञ्चमे मनः प्रति\-बु\-द्ध\-त\-र\-म्भ\-वति~॥ षष्ठे बुद्धिः~॥ सप्तमे }\edtext{\skt{सर्व्वप्रव्यक्तं शरीरं}}{\lemma{\skt{सर्व्वप्रव्यक्तं शरीरं}}\Afootnote{H; \skt{सर्वा\appbrk{}ङ्ग\appbrk{}प्र\appbrk{}त्य\appbrk{}ङ्ग\appbrk{}वि\appbrk{}भागः प्रव्यक्ततरः}\ A; \skt{सर्व\appbrk{}प्र\appbrk{}व्य\appbrk{}क्त\appbrk{}श\appbrk{}रीरः}\ L}}\skt{~॥ }\edtext{\edtext{\skt{अष्ट\-मेऽ\-स्थि\-री\-भ\-व\-त्यो\-ज\-स्तत्र}}{\lemma{\skt{ऽस्थिरीभवत्य्}}\Afootnote{A; \skt{स्थिरीभवत्य्}\ H L}}}{\lemma{\skt{ओजस्}}\Afootnote{A H; \emph{om.}\ L}}\skt{ }\edtext{\skt{जातो न}}{\lemma{\skt{जातो न}}\Afootnote{\skt{जातश्चेन्न}\ A; \skt{जातो}\ H; \skt{जातश्च न}\ L}}\skt{ }\edtext{\skt{जीवति~॥ नैरृ\-त\-भा\-ग\-धे\-य\-त्वा\-त्ततो}}{\lemma{\skt{जीवति~॥ नैरृ\appbrk{}त\appbrk{}भा\appbrk{}ग\appbrk{}धे\appbrk{}य\appbrk{}त्वा\appbrk{}त्ततो}}\Afootnote{\skt{जीवे\appbrk{}न्नि\appbrk{}रो\appbrk{}ज\appbrk{}स्त्वा\appbrk{}न्नै\appbrk{}रृ\appbrk{}त\appbrk{}भा\appbrk{}ग\appbrk{}त्वाच्च , ततो}\ A; \skt{न जीवति~॥ नैरृ\appbrk{}त\appbrk{}भा\appbrk{}ग\appbrk{}धे\appbrk{}य\appbrk{}त्वा\appbrk{}त्ततो}\ H; \skt{जीवे\appbrk{}न्नै\appbrk{}नृ\appbrk{}ति\appbrk{}भाग}\ L\textsuperscript{ac}; \skt{जीवेन्नैनृ}\ L\textsuperscript{pc}}}\skt{ }\edtext{\edtext{\skt{बलिम्मांसोदनं}}{\lemma{\skt{बलिम्}}\Afootnote{A H; \emph{om.}\ L}}}{\lemma{\skt{मांसोदनं}}\Afootnote{H; \skt{मांसौदनम्}\ A; \emph{om.}\ L}}\skt{ }\edtext{\skt{भागधेयन्तस्मै}}{\lemma{\skt{भागधेयन्तस्मै}}\Afootnote{H; \skt{अस्मै}\ A; \skt{त्वा}\ L\textsuperscript{ac}; \skt{तिभागत्}\ L\textsuperscript{pc}}}\skt{ }\edtext{\skt{दापयेत्}}{\lemma{\skt{दापयेत्}}\Afootnote{A H; \skt{त्}\ L\textsuperscript{ac}; \skt{वात्}\ L\textsuperscript{pc}}}\skt{~॥ }\edtext{\skt{नव\-म\-द\-श\-मै\-का\-द\-श\-द्वा\-द\-शा\-ना\-म\-न्य\-त\-म\-स्मि\-ञ्जा\-यते}}{\lemma{\skt{नव\appbrk{}म\appbrk{}द\appbrk{}श\appbrk{}मै\appbrk{}का\appbrk{}द\appbrk{}श\appbrk{}द्वा\appbrk{}द\appbrk{}शा\appbrk{}नाम्}}\Afootnote{A H; \skt{नव\appbrk{}म\appbrk{}न्द\appbrk{}श\appbrk{}मै\appbrk{}का\appbrk{}दश~। द्वादशानाम्}\ L}}\skt{~। }\edtext{\edtext{\skt{अतोऽन्यथा}}{\lemma{\skt{अतो}}\Afootnote{A; \skt{अथो}\ H L}}}{\lemma{\skt{ऽन्यथा}}\Afootnote{A; \skt{न्यथा}\ H L}}\skt{ विकारी भवति~॥}\pend

\pstart \ledsidenote{\passref{3.3.31}}\skt{मातुस्तु खलु }\edtext{\edtext{\skt{रस\-व\-हा\-या\-न्नाभ्यां}}{\lemma{\skt{रसवहायान्}}\Afootnote{A H; \skt{गर्भयरायान्}\ L}}}{\lemma{\skt{नाभ्यां}}\Afootnote{H; \skt{नाड्यां}\ A L}}\skt{ }\edtext{\skt{गर्ब्भनाडी}}{\lemma{\skt{गर्ब्भनाडी}}\Afootnote{H; \skt{गर्भनाभिनाडी}\ A; \skt{गर्भनाभिनाळी}\ L}}\skt{ प्रतिबद्धा }\edtext{\skt{भवति~। सास्य}}{\lemma{\skt{भवति~। सास्य}}\Afootnote{H; \skt{सास्य}\ A L}}\skt{ }\edtext{\skt{मातु\-रा\-हा\-र\-र\-स\-वी\-र्य\-म\-भ्य\-व\-ह\-रति}}{\lemma{\skt{अभ्यवहरति}}\Afootnote{H; \skt{अभिवहति}\ A; \skt{अभिवहन्ति}\ L}}\skt{~॥ }\edtext{\edtext{\skt{तेनो\-प\-स्ने\-हे\-ना\-भि\-वृ\-द्धि\-रस्य}}{\lemma{\skt{तेनो\appbrk{}प\appbrk{}स्ने\appbrk{}हे\appbrk{}ना\appbrk{}भि\appbrk{}वृ\appbrk{}द्धिर्}}\Afootnote{H L; \skt{तेनो\appbrk{}प\appbrk{}स्ने\appbrk{}हे\appbrk{}ना\appbrk{}स्या\appbrk{}भि\appbrk{}वृ\appbrk{}द्धिर्}\ A}}}{\lemma{\skt{अस्य}}\Afootnote{H; \emph{om.}\ A; \skt{अस्य च}\ L}}\skt{ }\edtext{\skt{भवति}}{\lemma{\skt{भवति}}\Afootnote{A H; \skt{भवन्ति}\ L}}\skt{~। }\edtext{\edtext{\skt{अस\-ञ्जा\-ता\-ङ्ग\-वि\-भा\-ग\-मपि}}{\lemma{\skt{अस\appbrk{}ञ्जा\appbrk{}ता\appbrk{}ङ्ग\appbrk{}वि\appbrk{}भा\appbrk{}गम्}}\Afootnote{H L\textsuperscript{pc}; \skt{अस\appbrk{}ञ्जा\appbrk{}ता\appbrk{}ङ्ग\appbrk{}प्र\appbrk{}त्य\appbrk{}ङ्ग\appbrk{}प्र\appbrk{}वि\appbrk{}भा\appbrk{}गम्}\ A; \skt{अस\appbrk{}ञ्जा\appbrk{}तां\appbrk{}ग\appbrk{}भि\appbrk{}वि\appbrk{}भा\appbrk{}गम्}\ L\textsuperscript{ac}}}\skt{ निःषेकात्प्रभृति}}{\lemma{\skt{अपि निःषेकात्}}\Afootnote{\skt{आ निषेकात्}\ A; \skt{अपि निःषेकान्}\ H; \skt{अभिनिषेक्त्}\ L\textsuperscript{ac}; \skt{अभिनिषेकात्}\ L\textsuperscript{pc}}}\skt{ }\edtext{\skt{सर्व्व\-श\-री\-रा\-नु\-सा\-रि\-णीनां}}{\lemma{\skt{सर्व्व\appbrk{}श\appbrk{}री\appbrk{}रा\appbrk{}नु\appbrk{}सा\appbrk{}रि\appbrk{}णीनां}}\Afootnote{H L; \skt{सर्व\appbrk{}श\appbrk{}री\appbrk{}रा\appbrk{}व\appbrk{}य\appbrk{}वा\appbrk{}नु\appbrk{}सा\appbrk{}रि\appbrk{}णीनां}\ A}}\skt{ }\edtext{\skt{रसवहानां}}{\lemma{\skt{रसवहानां}}\Afootnote{A H; \skt{रस\appbrk{}वा\appbrk{}हि\appbrk{}नी\appbrk{}नाम्\_}\ L}}\skt{ }\edtext{\edtext{\skt{तिर्य\-ग्ग\-ती\-ना\-न्ध\-म\-नी\-ना\-मु\-प\-स्ने\-हेन}}{\lemma{\skt{तिर्यग्गतीनान्}}\Afootnote{H; \skt{तिर्यग्गतानां}\ A L}}}{\lemma{\skt{उपस्नेहेन}}\Afootnote{H; \skt{उपस्नेहो}\ A L}}\skt{ जीवयति~॥}\footnoteA{We emend the H reading \emph{niḥṣekān} to the ablative, following the sense of the passage and the vulgate reading.}\pend

\pstart \ledsidenote{\passref{3.3.32}}\edtext{\skt{गर्ब्भस्य}}{\lemma{\skt{गर्ब्भस्य}}\Afootnote{H; \skt{गर्भस्य}\ A L}}\skt{ खलु }\edtext{\skt{सम्भवतः}}{\lemma{\skt{सम्भवतः}}\Afootnote{A H; \emph{om.}\ L}}\skt{~॥ }\edtext{\skt{पूर्व्वं}}{\lemma{\skt{पूर्व्वं}}\Afootnote{A H; \emph{om.}\ L}}\skt{ }\edtext{\skt{शिरः}}{\lemma{\skt{शिरः}}\Afootnote{A H; \emph{om.}\ L}}\skt{ }\edtext{\edtext{\skt{सम्भवतीत्याह}}{\lemma{\skt{सम्भवतीत्य्}}\Afootnote{A H; \skt{सम्भव}\ L}}}{\lemma{\skt{आह}}\Afootnote{A H; \emph{om.}\ L}}\skt{ }\edtext{\edtext{\skt{शौन\-क\-स्त\-न्मू\-ल\-त्वात्}}{\lemma{\skt{शौनकस्}}\Afootnote{A H; \emph{om.}\ L}}}{\lemma{\skt{तन्मूलत्वात्}}\Afootnote{H; \skt{शिरोमूलत्वात्}\ A; \emph{om.}\ L}}\skt{~॥ }\edtext{\skt{प्रधानेन्द्रियाणां}}{\lemma{\skt{प्रधानेन्द्रियाणां}}\Afootnote{A H; \emph{om.}\ L}}\skt{ }\edtext{\edtext{\skt{हृदयमिति}}{\lemma{\skt{हृदयम्}}\Afootnote{A H; \emph{om.}\ L}}}{\lemma{\skt{इति}}\Afootnote{A H; \emph{om.}\ L}}\skt{ }\edtext{\skt{कृतवीर्यः}}{\lemma{\skt{कृतवीर्यः}}\Afootnote{A H; \emph{om.}\ L}}\skt{ }\edtext{\edtext{\skt{बुद्धेर्मनसश्च}}{\lemma{\skt{बुद्धेर्}}\Afootnote{A H; \emph{om.}\ L}}}{\lemma{\skt{मनसश्च}}\Afootnote{A H; \emph{om.}\ L}}\skt{ }\edtext{\skt{स्थानं}}{\lemma{\skt{स्थानं}}\Afootnote{H; \skt{स्थान}\ A; \emph{om.}\ L}}\skt{~॥ }\edtext{\edtext{\skt{नाभिरिति}}{\lemma{\skt{नाभिर्}}\Afootnote{H; \skt{त्वात् ; नाभिर्}\ A; \emph{om.}\ L}}}{\lemma{\skt{इति}}\Afootnote{A H; \emph{om.}\ L}}\skt{ }\edtext{\edtext{\skt{पाराशर्यस्ततो}}{\lemma{\skt{पाराशर्यस्}}\Afootnote{A H; \emph{om.}\ L}}}{\lemma{\skt{ततो}}\Afootnote{H; \skt{ततो हि}\ A; \skt{तस्}\ L}}\skt{ }\edtext{\skt{वर्द्धते}}{\lemma{\skt{वर्द्धते}}\Afootnote{A H; \emph{om.}\ L}}\skt{ }\edtext{\skt{देहदेहिनः}}{\lemma{\skt{देहदेहिनः}}\Afootnote{H; \skt{देहो देहिनः}\ A; \emph{om.}\ L}}\skt{ }\edtext{\skt{प्राणः}}{\lemma{\skt{प्राणः}}\Afootnote{H; \emph{om.}\ A L}}\skt{~। }\edtext{\edtext{\skt{पाणिपादमिति}}{\lemma{\skt{पाणिपादम्}}\Afootnote{A H; \emph{om.}\ L}}}{\lemma{\skt{इति}}\Afootnote{A H; \emph{om.}\ L}}\skt{ }\edtext{\edtext{\edtext{\skt{मार्क्क\-ण्डे\-य\-स्त\-न्मू\-ल\-त्वा\-च्चे\-ष्टायाः}}{\lemma{\skt{मार्क्कण्डेयस्}}\Afootnote{A H; \emph{om.}\ L}}}{\lemma{\skt{तन्मूलत्वाच्}}\Afootnote{A H; \emph{om.}\ L}}}{\lemma{\skt{चेष्टायाः}}\Afootnote{A H; \emph{om.}\ L}}\skt{~। }\edtext{\skt{गर्ब्भस्य}}{\lemma{\skt{गर्ब्भस्य}}\Afootnote{H; \skt{गर्भस्य; मध्य}\ A; \emph{om.}\ L}}\skt{ }\edtext{\edtext{\skt{शरीरमिति}}{\lemma{\skt{शरीरम्}}\Afootnote{A H; \emph{om.}\ L}}}{\lemma{\skt{इति}}\Afootnote{A H; \emph{om.}\ L}}\skt{ }\edtext{\edtext{\skt{सुभू\-ति\-गौ\-त\-म\-स्तत्र}}{\lemma{\skt{सुभू\appbrk{}ति\appbrk{}गौ\appbrk{}त\appbrk{}मस्}}\Afootnote{H; \skt{सुभूतिर्गौतमः}\ A; \emph{om.}\ L}}}{\lemma{\skt{तत्र}}\Afootnote{H; \skt{तन्}\ A; \emph{om.}\ L}}\skt{ }\edtext{\edtext{\skt{निब\-न्ध\-त्वा\-त्स\-र्व्व\-गा\-त्र\-स\-म्भ\-वस्य}}{\lemma{\skt{निबन्धत्वात्}}\Afootnote{H; \skt{निबद्धत्वात्}\ A; \emph{om.}\ L}}}{\lemma{\skt{सर्व्व\appbrk{}गा\appbrk{}त्र\appbrk{}स\appbrk{}म्भ\appbrk{}वस्य}}\Afootnote{A H; \emph{om.}\ L}}\skt{~। }\edtext{\skt{तत्तु न}}{\lemma{\skt{तत्तु न}}\Afootnote{A; \skt{तन्तु न}\ H; \emph{om.}\ L}}\skt{ }\edtext{\skt{सम्यक्}}{\lemma{\skt{सम्यक्}}\Afootnote{A H; \emph{om.}\ L}}\skt{~। सर्व्वा\-ण्य\-ङ्ग\-प्र\-त्य\-ङ्गानि }\edtext{\edtext{\skt{युग\-प\-त्स\-म्भ\-वन्ति}}{\lemma{\skt{युगपत्}}\Afootnote{A H; \skt{युग}\ L}}}{\lemma{\skt{सम्भवन्ति}}\Afootnote{H; \skt{संभवन्तीत्य्}\ A; \emph{om.}\ L}}\skt{~। }\edtext{\skt{काले}}{\lemma{\skt{काले}}\Afootnote{H; \skt{आह धन्वन्तरिः}\ A; \emph{om.}\ L}}\skt{ }\edtext{\skt{गर्ब्भस्य}}{\lemma{\skt{गर्ब्भस्य}}\Afootnote{H; \skt{गर्भस्य}\ A; \emph{om.}\ L}}\skt{ }\edtext{\edtext{\skt{सूक्ष्म\-त्वा\-न्नो\-प\-ल\-भ्यन्ते}}{\lemma{\skt{सूक्ष्मत्वान्}}\Afootnote{A H; \emph{om.}\ L}}}{\lemma{\skt{नोपलभ्यन्ते}}\Afootnote{A H; \emph{om.}\ L}}\skt{~। }\edtext{\edtext{\skt{वंशा\-ङ्कु\-र\-व\-च्चू\-त\-फ\-ल\-वच्च}}{\lemma{\skt{वंशाङ्कुरवच्}}\Afootnote{A H; \emph{om.}\ L}}}{\lemma{\skt{चूतफलवच्च}}\Afootnote{A H; \emph{om.}\ L}}\skt{~। }\edtext{\edtext{\skt{तद्यथा}}{\lemma{\skt{तद्}}\Afootnote{A H; \emph{om.}\ L}}}{\lemma{\skt{यथा}}\Afootnote{A H; \emph{om.}\ L}}\skt{ }\edtext{\skt{चूतफले}}{\lemma{\skt{चूतफले}}\Afootnote{A H; \emph{om.}\ L}}\skt{ }\edtext{\skt{विपक्वे}}{\lemma{\skt{विपक्वे}}\Afootnote{H; \skt{परिपक्वे}\ A; \emph{om.}\ L}}\skt{ }\edtext{\skt{केश\-र\-मां\-सा\-स्थि\-म\-ज्जानः}}{\lemma{\skt{केश\appbrk{}र\appbrk{}मां\appbrk{}सा\appbrk{}स्थि\appbrk{}म\appbrk{}ज्जानः}}\Afootnote{A H\textsuperscript{pc}; \skt{केश\appbrk{}मा\appbrk{}न्सा\appbrk{}स्थि\appbrk{}म\appbrk{}ज्जानः}\ H\textsuperscript{ac}; \emph{om.}\ L}}\skt{ }\edtext{\edtext{\edtext{\skt{पृथ\-क्पृ\-थ\-ग्दृ\-श्यन्ते}}{\lemma{\skt{पृथक्}}\Afootnote{A H; \emph{om.}\ L}}}{\lemma{\skt{पृथग्}}\Afootnote{A H; \emph{om.}\ L}}}{\lemma{\skt{दृश्यन्ते}}\Afootnote{A H; \emph{om.}\ L}}\skt{ }\edtext{\edtext{\edtext{\skt{काल\-प्र\-क\-र्षा\-त्ता\-न्येव}}{\lemma{\skt{कालप्रकर्षात्}}\Afootnote{A H; \emph{om.}\ L}}}{\lemma{\skt{तान्य्}}\Afootnote{A H; \emph{om.}\ L}}}{\lemma{\skt{एव}}\Afootnote{A H; \emph{om.}\ L}}\skt{ }\edtext{\skt{तरुणे}}{\lemma{\skt{तरुणे}}\Afootnote{A H; \emph{om.}\ L}}\skt{ }\edtext{\skt{नोपलभ्यन्ते}}{\lemma{\skt{नोपलभ्यन्ते}}\Afootnote{A H; \skt{प}\ L}}\skt{ }\edtext{\skt{सूक्ष्मत्वात्}}{\lemma{\skt{सूक्ष्मत्वात्}}\Afootnote{A H; \emph{om.}\ L}}\skt{~। }\edtext{\skt{तेषां}}{\lemma{\skt{तेषां}}\Afootnote{A H; \emph{om.}\ L}}\skt{ }\edtext{\skt{सूक्ष्माणां}}{\lemma{\skt{सूक्ष्माणां}}\Afootnote{A H; \emph{om.}\ L}}\skt{ }\edtext{\edtext{\skt{केश\-रा\-दी\-ना\-म्प्र\-व्यक्तं}}{\lemma{\skt{केशरादीनाम्}}\Afootnote{A H; \emph{om.}\ L}}}{\lemma{\skt{प्रव्यक्तं}}\Afootnote{H; \skt{कालः प्रव्यक्त}\ A; \emph{om.}\ L}}\skt{ }\edtext{\skt{कालः}}{\lemma{\skt{कालः}}\Afootnote{H; \skt{तां}\ A; \emph{om.}\ L}}\skt{ }\edtext{\skt{करोति}}{\lemma{\skt{करोति}}\Afootnote{A H; \emph{om.}\ L}}\skt{~। }\edtext{\skt{एतेनैव}}{\lemma{\skt{एतेनैव}}\Afootnote{A H; \emph{om.}\ L}}\skt{ }\edtext{\skt{वंशाङ्कुरो}}{\lemma{\skt{वंशाङ्कुरो}}\Afootnote{A H; \emph{om.}\ L}}\skt{ }\edtext{\skt{व्याख्यातः}}{\lemma{\skt{व्याख्यातः}}\Afootnote{H; \skt{ऽपि व्याख्यातः}\ A; \emph{om.}\ L}}\skt{~। }\edtext{\skt{एवं}}{\lemma{\skt{एवं}}\Afootnote{A H; \emph{om.}\ L}}\skt{ }\edtext{\skt{गर्ब्भस्य}}{\lemma{\skt{गर्ब्भस्य}}\Afootnote{H; \skt{गर्भस्य}\ A; \emph{om.}\ L}}\skt{ }\edtext{\skt{तारुण्ये}}{\lemma{\skt{तारुण्ये}}\Afootnote{A H; \emph{om.}\ L}}\skt{ }\edtext{\edtext{\skt{सर्व्वे\-ष्व\-ङ्ग\-प्र\-त्य\-ङ्गेषु}}{\lemma{\skt{सर्व्वेष्व्}}\Afootnote{A H; \skt{दे}\ L}}}{\lemma{\skt{अङ्गप्रत्यङ्गेषु}}\Afootnote{A H; \skt{व सम्}\ L}}\skt{ }\edtext{\edtext{\skt{सत्स्वपि}}{\lemma{\skt{सत्स्व्}}\Afootnote{A H; \emph{om.}\ L}}}{\lemma{\skt{अपि}}\Afootnote{A H; \emph{om.}\ L}}\skt{ }\edtext{\edtext{\edtext{\skt{सौक्ष्म्या\-द\-नु\-प\-ल\-ब्धि\-स्ता\-न्येव}}{\lemma{\skt{सौक्ष्म्या\appbrk{}द\appbrk{}नु\appbrk{}प\appbrk{}ल\appbrk{}ब्धिस्}}\Afootnote{A H; \emph{om.}\ L}}}{\lemma{\skt{तान्य्}}\Afootnote{A H; \emph{om.}\ L}}}{\lemma{\skt{एव}}\Afootnote{A H; \emph{om.}\ L}}\skt{ }\edtext{\edtext{\skt{काल\-प्र\-क\-र्षा\-त्प्र\-व्य\-क्तानि}}{\lemma{\skt{कालप्रकर्षात्}}\Afootnote{A H; \emph{om.}\ L}}}{\lemma{\skt{प्रव्यक्तानि}}\Afootnote{A H; \emph{om.}\ L}}\skt{ भवन्ति~। }\edtext{\skt{पूर्व्वो\-त्त\-र\-का\-ल\-योगो}}{\lemma{\skt{पूर्व्वो\appbrk{}त्त\appbrk{}र\appbrk{}का\appbrk{}ल\appbrk{}योगो}}\Afootnote{H L; \emph{om.}\ A}}\skt{ }\edtext{\skt{नास्तीति}}{\lemma{\skt{नास्तीति}}\Afootnote{H; \emph{om.}\ A; \skt{नास्ति}\ L}}\skt{ }\edtext{\skt{केवलन्तु}}{\lemma{\skt{केवलन्तु}}\Afootnote{H L; \emph{om.}\ A}}\skt{ }\edtext{\edtext{\skt{सौक्ष्म्या\-न्ना\-भि\-व्य\-ज्यन्ते}}{\lemma{\skt{सौक्ष्म्यान्}}\Afootnote{H L; \emph{om.}\ A}}}{\lemma{\skt{नाभिव्यज्यन्ते}}\Afootnote{H; \emph{om.}\ A; \skt{नाभिव्यज्यते}\ L}}\skt{~। }\edtext{\skt{यथाकालं}}{\lemma{\skt{यथाकालं}}\Afootnote{H L; \emph{om.}\ A}}\skt{ }\edtext{\skt{प्रव्यक्तानि}}{\lemma{\skt{प्रव्यक्तानि}}\Afootnote{H L; \emph{om.}\ A}}\skt{ }\edtext{\skt{भवन्ति}}{\lemma{\skt{भवन्ति}}\Afootnote{H L; \emph{om.}\ A}}\skt{~।}\footnoteA{The last two sentences, which are not present in the vulgate, paraphrase the preceding sentence and may be a gloss.}\pend

\pstart \ledsidenote{\passref{3.3.33}}\edtext{\skt{तत्र}}{\lemma{\skt{तत्र}}\Afootnote{A; \emph{om.}\ H L}}\skt{ }\edtext{\skt{पितृ\-ज\-मा\-तृ\-ज\-र\-स\-जा\-त्म\-ज\-स\-त्त्व\-सा\-त्म्य\-जानि}}{\lemma{\skt{पितृ\appbrk{}ज\appbrk{}मा\appbrk{}तृ\appbrk{}ज\appbrk{}र\appbrk{}स\appbrk{}जा\appbrk{}त्म\appbrk{}ज\appbrk{}स\appbrk{}त्त्व\appbrk{}सा\appbrk{}त्म्य\appbrk{}जानि}}\Afootnote{\skt{गर्भस्य पितृ\appbrk{}ज\appbrk{}मा\appbrk{}तृ\appbrk{}ज\appbrk{}र\appbrk{}स\appbrk{}जा\appbrk{}त्म\appbrk{}ज\appbrk{}स\appbrk{}त्त्व\appbrk{}ज\appbrk{}सा\appbrk{}त्म्य\appbrk{}जानि}\ A; \skt{तत्र पितृ\appbrk{}ज\appbrk{}मा\appbrk{}तृ\appbrk{}ज\appbrk{}र\appbrk{}स\appbrk{}जा\appbrk{}त्म\appbrk{}ज\appbrk{}स\appbrk{}त्व\appbrk{}सा\appbrk{}त्म्य\appbrk{}जानि}\ H; \emph{om.}\ L}}\skt{ }\edtext{\skt{शरीरे}}{\lemma{\skt{शरीरे}}\Afootnote{H; \skt{शरीर}\ A; \emph{om.}\ L}}\skt{ }\edtext{\skt{लक्षणानि}}{\lemma{\skt{लक्षणानि}}\Afootnote{A H; \emph{om.}\ L}}\skt{ }\edtext{\skt{वक्ष्यामः}}{\lemma{\skt{वक्ष्यामः}}\Afootnote{H; \skt{व्याख्यास्यामः}\ A; \emph{om.}\ L}}\skt{~॥ }\edtext{\skt{तत्र}}{\lemma{\skt{तत्र}}\Afootnote{H; \skt{गर्भस्य}\ A; \skt{गर्भस्य तु}\ L}}\skt{ }\edtext{\skt{केश\-श्म\-श्रु\-द\-न्त\-न\-ख\-रो\-मा\-स्थि\-रे\-तः\-प्र\-भृ\-तीनि}}{\lemma{\skt{केश\appbrk{}श्म\appbrk{}श्रु\appbrk{}द\appbrk{}न्त\appbrk{}न\appbrk{}ख\appbrk{}रो\appbrk{}मा\appbrk{}स्थि\appbrk{}रे\appbrk{}तः\appbrk{}प्र\appbrk{}भृ\appbrk{}तीनि}}\Afootnote{H; \skt{केश\appbrk{}श्म\appbrk{}श्रु\appbrk{}लो\appbrk{}मा\appbrk{}स्थि\appbrk{}न\appbrk{}ख\appbrk{}द\appbrk{}न्त\appbrk{}सि\appbrk{}रा\appbrk{}स्ना\appbrk{}यु\appbrk{}ध\appbrk{}म\appbrk{}नी\appbrk{}रे\appbrk{}तः\appbrk{}प्र\appbrk{}भृ\appbrk{}तीनि}\ A; \skt{केश\appbrk{}श्म\appbrk{}श्रु\appbrk{}लो\appbrk{}चु\appbrk{}न\appbrk{}खा\appbrk{}स्थि\appbrk{}प्र\appbrk{}भृ\appbrk{}तीनि}\ L}}\skt{ }\edtext{\skt{पितृजानि}}{\lemma{\skt{पितृजानि}}\Afootnote{H; \skt{स्थिराणि पितृजानि}\ A L}}\skt{~॥ }\edtext{\skt{मांस\-शो\-णि\-त\-मे\-दो\-म\-ज्जा\-हृ\-न्ना\-भि\-य\-कृ\-त्प्ली\-हा\-न्त्र\-गु\-द\-प्र\-भृ\-तीनि}}{\lemma{\skt{मांस\appbrk{}शो\appbrk{}णि\appbrk{}त\appbrk{}मे\appbrk{}दो\appbrk{}म\appbrk{}ज्जा\appbrk{}हृ\appbrk{}न्ना\appbrk{}भि\appbrk{}य\appbrk{}कृ\appbrk{}त्प्ली\appbrk{}हा\appbrk{}न्त्र\appbrk{}गु\appbrk{}द\appbrk{}प्र\appbrk{}भृ\appbrk{}तीनि}}\Afootnote{H; \skt{मांस\appbrk{}शो\appbrk{}णि\appbrk{}त\appbrk{}मे\appbrk{}दो\appbrk{}म\appbrk{}ज्ज\appbrk{}हृ\appbrk{}न्ना\appbrk{}भि\appbrk{}य\appbrk{}कृ\appbrk{}त्प्ली\appbrk{}हा\appbrk{}न्त्र\appbrk{}गु\appbrk{}द\appbrk{}प्र\appbrk{}भृ\appbrk{}तीनि}\ A; \skt{मांस\appbrk{}शो\appbrk{}णि\appbrk{}त।\appbrk{}मे\appbrk{}दो\appbrk{}म\appbrk{}ज्ज\appbrk{}प्र\appbrk{}भृ\appbrk{}तिनी}\ L}}\skt{ मृदूनि मातृजानि~। शरीरोपचयो }\edtext{\edtext{\skt{बल\-वृ\-द्धि\-र्व्व\-र्ण्ण\-स्थि\-ति\-हा\-निश्च}}{\lemma{\skt{बलवृद्धिर्}}\Afootnote{H; \skt{बलं}\ A; \skt{बलं प्}\ L}}}{\lemma{\skt{व्वर्ण्ण\appbrk{}स्थि\appbrk{}ति\appbrk{}हा\appbrk{}निश्च}}\Afootnote{H; \skt{वर्णः स्थितिर्हानिश्च}\ A; \skt{रद्धि स्थितिर्हानिश्च}\ L}}\skt{ }\edtext{\skt{रसजानि}}{\lemma{\skt{रसजानि}}\Afootnote{A H; \skt{रजसानि}\ L}}\skt{~। }\edtext{\skt{इन्द्रियाणि}}{\lemma{\skt{इन्द्रियाणि}}\Afootnote{A H; \emph{om.}\ L}}\skt{ }\edtext{\edtext{\skt{ज्ञानमायुः}}{\lemma{\skt{ज्ञानम्}}\Afootnote{H; \skt{ज्ञानं विज्ञानम्}\ A; \skt{आयुरारो}\ L}}}{\lemma{\skt{आयुः}}\Afootnote{A H; \skt{ग्यं मेधा}\ L}}\skt{ }\edtext{\skt{सुखदुःखादिकं}}{\lemma{\skt{सुखदुःखादिकं}}\Afootnote{A H; \skt{सुखदुःखञ्}\ L}}\skt{ }\edtext{\skt{चात्मनः}}{\lemma{\skt{चात्मनः}}\Afootnote{H; \skt{चात्मजानि}\ A L}}\skt{~। }\edtext{\edtext{\skt{सत्त्व\-जा\-न्यु\-त्त\-रत्र}}{\lemma{\skt{सत्त्वजान्य्}}\Afootnote{A H; \emph{om.}\ L}}}{\lemma{\skt{उत्तरत्र}}\Afootnote{A H; \emph{om.}\ L}}\skt{ }\edtext{\skt{वक्ष्यामः}}{\lemma{\skt{वक्ष्यामः}}\Afootnote{A H; \emph{om.}\ L}}\skt{~। }\edtext{\edtext{\edtext{\skt{वीर्य\-मा\-रो\-ग्य\-म्बलं}}{\lemma{\skt{वीर्यम्}}\Afootnote{A H; \emph{om.}\ L}}}{\lemma{\skt{आरोग्यम्}}\Afootnote{A H; \emph{om.}\ L}}}{\lemma{\skt{बलं}}\Afootnote{H; \skt{बल}\ A; \emph{om.}\ L}}\skt{ }\edtext{\skt{वर्ण्णो}}{\lemma{\skt{वर्ण्णो}}\Afootnote{H; \skt{वर्णौ}\ A; \emph{om.}\ L}}\skt{ }\edtext{\skt{मेधा च}}{\lemma{\skt{मेधा च}}\Afootnote{A H; \emph{om.}\ L}}\skt{ }\edtext{\skt{सात्म्यजं}}{\lemma{\skt{सात्म्यजं}}\Afootnote{H; \skt{सात्म्यजानि}\ A; \emph{om.}\ L}}\skt{~॥}\pend

\pstart \ledsidenote{\passref{3.3.34}}\skt{तत्र यस्या दक्षिणे स्तने }\edtext{\skt{प्राक्पयसो}}{\lemma{\skt{पयसो}}\Afootnote{H L; \skt{पयो}\ A}}\skt{ दर्शनं }\edtext{\edtext{\skt{दक्षि\-ण\-कु\-क्षि\-म\-ह\-त्वञ्च}}{\lemma{\skt{दक्षि\appbrk{}ण\appbrk{}कु\appbrk{}क्षि\appbrk{}म\appbrk{}ह\appbrk{}त्वञ्}}\Afootnote{H; \skt{भवति दक्षिणाक्षिम}\ A; \skt{दक्षिणपार्श्वं सं}\ L}}}{\lemma{\skt{च}}\Afootnote{H; \skt{हत्त्वं च}\ A; \skt{श्रितश्च}\ L}}\skt{ }\edtext{\skt{पूर्व्वदक्षिणं}}{\lemma{\skt{पूर्व्वदक्षिणं}}\Afootnote{H; \skt{पूर्वं च दक्षिणं}\ A; \skt{गर्भः ग्रतौ दक्षिणां}\ L}}\skt{ }\edtext{\skt{सक्थि}}{\lemma{\skt{सक्थि}}\Afootnote{A H; \skt{सक्थ्}\ L}}\skt{ }\edtext{\skt{उत्कर्षयति}}{\lemma{\skt{उत्कर्षयति}}\Afootnote{H; \skt{उत्कर्षति}\ A; \skt{युत्कर्षति}\ L}}\skt{ }\edtext{\skt{बाहुल्याच्च}}{\lemma{\skt{बाहुल्याच्च}}\Afootnote{A H; \skt{बाहुश्च}\ L}}\skt{ }\edtext{\skt{पुन्नामधेयेषु}}{\lemma{\skt{पुन्नामधेयेषु}}\Afootnote{A H; \skt{पुन्नामधे}\ L}}\skt{ }\edtext{\skt{द्रव्येषु}}{\lemma{\skt{द्रव्येषु}}\Afootnote{A H; \skt{येषु}\ L}}\skt{ }\edtext{\edtext{\skt{दौहृ\-द\-म\-भि\-ध्या\-यति}}{\lemma{\skt{दौहृदम्}}\Afootnote{H; \skt{दौर्हृदम्}\ A; \skt{दौर्हृदेषु या}\ L}}}{\lemma{\skt{अभिध्यायति}}\Afootnote{A H; \skt{साभिलाषा भवति}\ L}}\skt{ }\edtext{\skt{स्वप्ने}}{\lemma{\skt{स्वप्ने}}\Afootnote{H L; \skt{स्वप्नेषु}\ A}}\skt{ चोपलभते }\edtext{\skt{पद्मो\-त्प\-ल\-कु\-मु\-दा\-म्रा\-म्रा\-त\-का\-दीनि}}{\lemma{\skt{पद्मो\appbrk{}त्प\appbrk{}ल\appbrk{}कु\appbrk{}मु\appbrk{}दा\appbrk{}म्रा\appbrk{}म्रा\appbrk{}त\appbrk{}का\appbrk{}दीनि}}\Afootnote{H; \skt{पद्मो\appbrk{}त्प\appbrk{}ल\appbrk{}कु\appbrk{}मु\appbrk{}दा\appbrk{}म्रा\appbrk{}त\appbrk{}का\appbrk{}दीनि}\ A; \emph{om.}\ L}}\skt{ }\edtext{\edtext{\skt{पुन्नामन्येव}}{\lemma{\skt{पुन्नामन्य्}}\Afootnote{H; \skt{पुन्नामान्य्}\ A; \skt{पुन्नामान्}\ L}}\skt{ वा}}{\lemma{\skt{एव वा}}\Afootnote{H; \skt{एव}\ A; \skt{येव}\ L}}\skt{ }\edtext{\skt{प्रस\-न्न\-मु\-ख\-वर्ण्णा च}}{\lemma{\skt{प्रस\appbrk{}न्न\appbrk{}मु\appbrk{}ख\appbrk{}वर्ण्णा च}}\Afootnote{A H; \skt{प्रसन्नो मुखवर्ण्णश्च}\ L}}\skt{ }\edtext{\skt{भवेतां}}{\lemma{\skt{भवेतां}}\Afootnote{H; \skt{भवति तां}\ A; \skt{ताम्\_}\ L}}\skt{ }\edtext{\skt{ब्रूया\-त्पु\-त्र\-मि\-य\-म्प्र\-ज\-न\-यि\-ष्य\-तीति}}{\lemma{\skt{प्रज\appbrk{}न\appbrk{}यि\appbrk{}ष्य\appbrk{}तीति}}\Afootnote{H; \skt{जनयिष्यतीति}\ A L}}\skt{~। तद्विपर्यये }\edtext{\skt{कन्यां}}{\lemma{\skt{कन्यां}}\Afootnote{A H; \skt{कन्न्याम्}\ L}}\skt{~। }\edtext{\skt{यस्याः}}{\lemma{\skt{यस्याः}}\Afootnote{A H; \emph{om.}\ L}}\skt{ }\edtext{\edtext{\edtext{\edtext{\edtext{\skt{पार्श्व\-द्व\-य\-म\-व\-न\-त\-म्पु\-र\-स्ता\-न्नि\-र्ग्ग\-त\-मु\-दरं}}{\lemma{\skt{पार्श्वद्वयम्}}\Afootnote{A H; \emph{om.}\ L}}}{\lemma{\skt{अवनतम्}}\Afootnote{H; \skt{उन्नतं}\ A; \emph{om.}\ L}}}{\lemma{\skt{पुरस्तान्}}\Afootnote{A H; \emph{om.}\ L}}}{\lemma{\skt{निर्ग्गतम्}}\Afootnote{A H; \emph{om.}\ L}}}{\lemma{\skt{उदरं}}\Afootnote{A H; \emph{om.}\ L}}\skt{ }\edtext{\skt{प्रागभिहितञ्च}}{\lemma{\skt{प्रागभिहितञ्च}}\Afootnote{A H; \emph{om.}\ L}}\skt{ }\edtext{\edtext{\skt{लक्षणन्तस्याः}}{\lemma{\skt{लक्षणन्}}\Afootnote{H; \skt{लक्षणं च}\ A; \emph{om.}\ L}}}{\lemma{\skt{तस्याः}}\Afootnote{A H; \emph{om.}\ L}}\skt{ }\edtext{\edtext{\skt{नपुंसकमिति}}{\lemma{\skt{नपुंसकम्}}\Afootnote{A H; \emph{om.}\ L}}}{\lemma{\skt{इति}}\Afootnote{A H; \emph{om.}\ L}}\skt{ }\edtext{\skt{विद्यात्}}{\lemma{\skt{विद्यात्}}\Afootnote{A H; \emph{om.}\ L}}\skt{~। }\edtext{\skt{यस्या}}{\lemma{\skt{यस्या}}\Afootnote{A H; \emph{om.}\ L}}\skt{ }\edtext{\skt{मध्ये}}{\lemma{\skt{मध्ये}}\Afootnote{A H; \emph{om.}\ L}}\skt{ }\edtext{\skt{निम्नो}}{\lemma{\skt{निम्नो}}\Afootnote{H; \skt{निम्नं}\ A; \emph{om.}\ L}}\skt{ }\edtext{\edtext{\skt{द्रोणी\-प्र\-भू\-त\-मु\-दरं}}{\lemma{\skt{द्रोणीप्रभूतम्}}\Afootnote{\skt{द्रोणीभूतम्}\ A; \skt{द्रोणीप्राभूतम्}\ H; \emph{om.}\ L}}\skt{ सा}}{\lemma{\skt{उदरं सा}}\Afootnote{A; \skt{उद…}\ H; \emph{om.}\ L}}\skt{ }\edtext{\skt{युग्मं}}{\lemma{\skt{युग्मं}}\Afootnote{A H; \emph{om.}\ L}}\skt{ }\edtext{\skt{प्रसूयत}}{\lemma{\skt{प्रसूयत}}\Afootnote{A H; \emph{om.}\ L}}\skt{ }\edtext{\skt{इति}}{\lemma{\skt{इति}}\Afootnote{A H; \emph{om.}\ L}}\skt{~॥}\pend

\pstart\versepar \ledsidenote{\passref{3.3.35}}\edtext{\skt{भवति}}{\lemma{\skt{भवति}}\Afootnote{H; \skt{भवन्ति}\ A L}}\skt{ चात्र~॥ }\pend
\pstart\versepar \skt{}\edtext{\skt{देव\-ता\-ब्रा\-ह्म\-ण\-पराः}}{\lemma{\skt{देव\appbrk{}ता\appbrk{}ब्रा\appbrk{}ह्म\appbrk{}ण\appbrk{}पराः}}\Afootnote{A H; \skt{श्लोकः~। दैव\appbrk{}त\appbrk{}ब्रा\appbrk{}ह्म\appbrk{}ण\appbrk{}परा}\ L}}\skt{ }\edtext{\skt{शुच्यो}}{\lemma{\skt{शुच्यो}}\Afootnote{H; \skt{शौचाचा}\ A; \skt{शुद्ध्या}\ L}}\skt{ }\edtext{\skt{हितमिताशनाः}}{\lemma{\skt{हितमिताशनाः}}\Afootnote{H; \skt{रहिते रताः}\ A; \skt{हितमिताशनि}\ L}}\skt{~।}\footnoteA{The form \emph{śucyaḥ} appears to be a contraction of \emph{śucayaḥ}, \emph{metri causa}. It could be seen as a transfer of an \emph{-i} stem to an \emph{-ī} stem (Oberlies 2003: 83-84).}\skt{ }\pend
\pstart\versepar \skt{}\edtext{\skt{महागुणाः}}{\lemma{\skt{महागुणाः}}\Afootnote{H; \skt{महागुणान्}\ A; \skt{महागुणम्\_}\ L}}\skt{ }\edtext{\skt{प्रसूयन्ते}}{\lemma{\skt{प्रसूयन्ते}}\Afootnote{A H; \skt{प्रसूयेत}\ L}}\skt{ }\edtext{\skt{विपरीतास्तु}}{\lemma{\skt{विपरीतास्तु}}\Afootnote{A H; \skt{विपरीता च}\ L}}\skt{ }\edtext{\skt{निर्ग्गुणाः}}{\lemma{\skt{निर्ग्गुणाः}}\Afootnote{H; \skt{निर्गुणान्}\ A; \skt{निर्गुणम्}\ L}}\skt{~॥}\pend

\pstart\versepar \ledsidenote{\passref{3.3.36}}\edtext{\skt{अङ्ग\-प्र\-त्य\-ङ्ग\-नि\-र्वृत्तिः}}{\lemma{\skt{अङ्ग\appbrk{}प्र\appbrk{}त्य\appbrk{}ङ्ग\appbrk{}नि\appbrk{}र्वृत्तिः}}\Afootnote{A H; \skt{अंग\appbrk{}प्र\appbrk{}त्यं\appbrk{}ग\appbrk{}नि\appbrk{}र्वृत्ति}\ L}}\skt{ स्वभावादेव जायते~। }\pend
\pstart\versepar \skt{}\edtext{\skt{अङ्ग\-प्र\-त्य\-ङ्ग\-नि\-र्वृत्ता ये}}{\lemma{\skt{अङ्ग\appbrk{}प्र\appbrk{}त्य\appbrk{}ङ्ग\appbrk{}नि\appbrk{}र्वृत्ता ये}}\Afootnote{H; \skt{अङ्ग\appbrk{}प्र\appbrk{}त्य\appbrk{}ङ्ग\appbrk{}नि\appbrk{}र्वृत्तौ ये}\ A; \emph{om.}\ L}}\skt{ }\edtext{\skt{भवन्ति}}{\lemma{\skt{भवन्ति}}\Afootnote{A H; \emph{om.}\ L}}\skt{ गुणागुणाः~॥ }\pend
\pstart\versepar \skt{ते वै }\edtext{\skt{गर्ब्भस्य}}{\lemma{\skt{गर्ब्भस्य}}\Afootnote{H; \skt{गर्भस्य}\ A; \skt{वस्य}\ L}}\skt{ }\edtext{\skt{विज्ञेया}}{\lemma{\skt{विज्ञेया}}\Afootnote{A H; \skt{ते ज्ञेया}\ L}}\skt{ धर्म्मा\-ध\-र्म्म\-नि\-मि\-त्तजाः~॥}\edtext{}{\lemma{\skt{गुणागुणाः~॥ ते वै}}\Afootnote{H; \skt{गुणागुणाः~। ते ते}\ A; \skt{गुण ये त्}\ L}}\pend

\pstart \ledsidenote{\passref{3.3.trailer}}\edtext{\skt{इति}}{\lemma{\skt{इति}}\Afootnote{H; \skt{इति सुश्रु\appbrk{}त\appbrk{}सं\appbrk{}हि\appbrk{}तायां शारीरस्थाने गर्भावक्रान्ति}\ A}}\skt{ }\edtext{\skt{शारी\-र\-तृ\-ती\-योऽ\-ध्यायः}}{\lemma{\skt{शारीरतृतीयो}}\Afootnote{H; \skt{शारीरं नाम तृतीयो}\ A}}\skt{~॥}\pend

\endnumbering
\beginnumbering
\eledsection*{Adhyāya 4}
\pstart \ledsidenote{\passref{3.4.1}}\skt{अथातो }\edtext{\skt{गर्भव्याकरणं}}{\lemma{\skt{गर्भव्याकरणं}}\Afootnote{A; \skt{गर्ब्भव्याकरणं}\ H}}\skt{ शारीरं व्याख्यास्यामः~॥}\pend

\pstart \ledsidenote{\passref{3.4.3}}\edtext{\skt{अग्नीषोमो}}{\lemma{\skt{अग्नीषोमो}}\Afootnote{H; \skt{अग्निः सोमो}\ A}}\skt{ वायुः सत्त्वं रजस्तमः पञ्चेन्द्रियाणि }\edtext{\skt{भूतात्मा}}{\lemma{\skt{भूतात्मा}}\Afootnote{H; \skt{भूतात्}\ A}}\skt{ }\edtext{\edtext{\skt{मनश्चेति}}{\lemma{\skt{मनश्}}\Afootnote{H; \emph{om.}\ A}}}{\lemma{\skt{चेति}}\Afootnote{H; \skt{मेति}\ A}}\skt{ प्राणाः~।}\footnoteA{Gayadāsa's commentary discusses the word \emph{manas}, which is not present in the vulgate. See MS Cambridge Add.2491, f.34r.}\pend

\pstart \ledsidenote{\passref{3.4.4}}\skt{तस्य }\edtext{\skt{खलु}}{\lemma{\skt{खलु}}\Afootnote{H; \skt{खल्}\ A}}\skt{ }\edtext{\skt{शुक्र\-शो\-णि\-त\-स्या\-भि\-प\-च्य\-मा\-नस्य}}{\lemma{\skt{शुक्र\appbrk{}शो\appbrk{}णि\appbrk{}त\appbrk{}स्या\appbrk{}भि\appbrk{}प\appbrk{}च्य\appbrk{}मा\appbrk{}नस्य}}\Afootnote{H; \skt{वेवंप्रवृत्तस्य शुक्र\appbrk{}शो\appbrk{}णि\appbrk{}त\appbrk{}स्या\appbrk{}भि\appbrk{}प\appbrk{}च्य\appbrk{}मा\appbrk{}नस्य}\ A}}\skt{ }\edtext{\skt{क्षीरस्येव हि}}{\lemma{\skt{क्षीरस्येव हि}}\Afootnote{H; \skt{क्षीरस्येव}\ A}}\skt{ सन्तानिकाः सप्त त्वचो भवन्ति~॥ तासां }\edtext{\skt{प्रथ\-मा\-व\-भा\-सिनी}}{\lemma{\skt{प्रथ\appbrk{}मा\appbrk{}व\appbrk{}भा\appbrk{}सिनी}}\Afootnote{A; \skt{प्रथमो वभासिनी}\ H}}\skt{ नाम या }\edtext{\skt{सर्व्व\-व\-र्ण्णा\-न\-व\-भा\-स\-यति}}{\lemma{\skt{सर्व्व\appbrk{}व\appbrk{}र्ण्णा\appbrk{}न\appbrk{}व\appbrk{}भा\appbrk{}स\appbrk{}यति}}\Afootnote{H; \skt{सर्वा\appbrk{}न्व\appbrk{}र्णा\appbrk{}न\appbrk{}व\appbrk{}भा\appbrk{}स\appbrk{}यति}\ A}}\skt{~। पञ्चविधाञ्च च्छाया\-म्प्र\-का\-श\-यति~। सा व्रीहे\-र\-ष्टा\-द\-श\-भा\-ग\-प्र\-माणा~। }\edtext{\skt{द्वितीया}}{\lemma{\skt{द्वितीया}}\Afootnote{H; \skt{सिध्म\appbrk{}प\appbrk{}द्म\appbrk{}क\appbrk{}ण्ट\appbrk{}का\appbrk{}धि\appbrk{}ष्ठाना; द्वितीया}\ A}}\skt{ लोहिता नाम षोड\-श\-भा\-ग\-प्र\-माणा~। }\edtext{\skt{तृतीया}}{\lemma{\skt{तृतीया}}\Afootnote{H; \skt{तिल\appbrk{}का\appbrk{}ल\appbrk{}क\appbrk{}न्य\appbrk{}च्छ\appbrk{}व्य\appbrk{}ङ्गा\appbrk{}धि\appbrk{}ष्ठाना; तृतीया}\ A}}\skt{ श्वेता }\edtext{\skt{द्वाद\-श\-भा\-ग\-प्र\-माणा}}{\lemma{\skt{द्वाद\appbrk{}श\appbrk{}भा\appbrk{}ग\appbrk{}प्र\appbrk{}माणा}}\Afootnote{H; \skt{नाम, द्वाद\appbrk{}श\appbrk{}भा\appbrk{}ग\appbrk{}प्र\appbrk{}माणा}\ A}}\skt{~। }\edtext{\skt{चरुर्थी}}{\lemma{\skt{चरुर्थी}}\Afootnote{H; \skt{चर्म\appbrk{}द\appbrk{}ला\appbrk{}ज\appbrk{}ग\appbrk{}ल्ली\appbrk{}म\appbrk{}ष\appbrk{}का\appbrk{}धि\appbrk{}ष्ठाना; चतुर्थी}\ A}}\skt{ }\edtext{\skt{ताम्रा\-ष्ट\-भा\-ग\-प्र\-माणा}}{\lemma{\skt{ताम्रा\appbrk{}ष्ट\appbrk{}भा\appbrk{}ग\appbrk{}प्र\appbrk{}माणा}}\Afootnote{H; \skt{ताम्रा नामा\appbrk{}ष्ट\appbrk{}भा\appbrk{}ग\appbrk{}प्र\appbrk{}माणा}\ A}}\skt{~। }\edtext{\skt{पञ्चमी}}{\lemma{\skt{पञ्चमी}}\Afootnote{H; \skt{विवि\appbrk{}ध\appbrk{}कि\appbrk{}ला\appbrk{}स\appbrk{}कु\appbrk{}ष्ठा\appbrk{}धि\appbrk{}ष्ठाना; पञ्चमी}\ A}}\skt{ }\edtext{\skt{वेदनी}}{\lemma{\skt{वेदनी}}\Afootnote{H; \skt{वेदिनी}\ A}}\skt{ नाम पञ्च\-भा\-ग\-प्र\-माणा~। }\edtext{\skt{षष्ठी}}{\lemma{\skt{षष्ठी}}\Afootnote{H; \skt{कुष्ठ\appbrk{}वि\appbrk{}स\appbrk{}र्पा\appbrk{}धि\appbrk{}ष्ठाना; षष्ठी}\ A}}\skt{ रोहिणी नाम व्रीहिप्रमाणा~। }\edtext{\skt{सप्तमी}}{\lemma{\skt{सप्तमी}}\Afootnote{H; \skt{ग्रन्थ्य\appbrk{}प\appbrk{}च्य\appbrk{}र्बु\appbrk{}द\appbrk{}श्ली\appbrk{}प\appbrk{}द\appbrk{}ग\appbrk{}ल\appbrk{}ग\appbrk{}ण्डा\appbrk{}धि\appbrk{}ष्ठाना; सप्तमी}\ A}}\skt{ }\edtext{\skt{शुक्रधरा}}{\lemma{\skt{शुक्रधरा}}\Afootnote{H; \skt{मांसधरा}\ A}}\skt{ नाम व्रीहि\-द्व\-य\-प्र\-माणा~॥ }\edtext{\skt{यतो}}{\lemma{\skt{यतो}}\Afootnote{H; \skt{भग\appbrk{}न्द\appbrk{}र\appbrk{}वि\appbrk{}द्र\appbrk{}ध्य\appbrk{}र्शोऽ\appbrk{}धि\appbrk{}ष्ठाना~। यदेतत्प्रमाणं निर्दिष्टं तन्मां\appbrk{}स\appbrk{}ले\appbrk{}ष्व\appbrk{}व\appbrk{}का\appbrk{}शेषु, न ललाटे सूक्ष्माङ्गुल्यादिषु च; यतो}\ A}}\skt{ }\edtext{\skt{वक्ष्यत्युदरे }\edtext{\edtext{\skt{"व्री\-हि\-मु\-खे\-ना\-ङ्गु\-ल\-म\-ङ्गु\-ष्ठो\-द\-र\-मा\-त्र\-म्वा\-व\-गा\-ह्येति"}}{\lemma{\skt{अङ्गु\appbrk{}ष्ठो\appbrk{}द\appbrk{}र\appbrk{}मा\appbrk{}त्रम्}}\Afootnote{H; \skt{गुष्ठो\appbrk{}द\appbrk{}र\appbrk{}प्र\appbrk{}मा\appbrk{}णम्}\ A}}}{\lemma{\skt{वावगाह्येति"}}\Afootnote{H; \skt{अवगाढं विध्येत्’ (चि.१४) इति}\ A}}}{\lemma{\skt{उदरे "व्री\appbrk{}हि\appbrk{}मु\appbrk{}खे\appbrk{}ना\appbrk{}ङ्गु\appbrk{}लम्}}\Afootnote{H; \skt{उदरेषु --- ‘व्री\appbrk{}हि\appbrk{}मु\appbrk{}खे\appbrk{}नाङ्}\ A}}\skt{~।}\footnoteA{This is a citation from Cikitsāsthāna 14.18.}\pend

\pstart \ledsidenote{\passref{3.4.5}}\skt{कलाः खल्वपि सप्त }\edtext{\skt{सम्भवन्ति}}{\lemma{\skt{सम्भवन्ति}}\Afootnote{H; \skt{भवन्ति}\ A}}\skt{ }\edtext{\skt{धात्वा\-न्त\-रा\-श\-य\-म\-र्यादाः}}{\lemma{\skt{धात्वा\appbrk{}न्त\appbrk{}रा\appbrk{}श\appbrk{}य\appbrk{}म\appbrk{}र्यादाः}}\Afootnote{\skt{धात्वा\appbrk{}श\appbrk{}या\appbrk{}न्त\appbrk{}र\appbrk{}म\appbrk{}र्यादाः}\ A; \skt{धात्वा\appbrk{}न्त\appbrk{}रा\appbrk{}श\appbrk{}य\appbrk{}म\appbrk{}र्यादाः~॥ श्लोकौ}\ H}}\skt{~॥}\pend

\pstart\versepar \ledsidenote{\passref{3.4.6}}\edtext{\skt{श्लोकौ}}{\lemma{\skt{श्लोकौ}}\Afootnote{\skt{भवतश्चात्र}\ A; \emph{om.}\ H}}\skt{~॥ }\pend
\pstart\versepar \skt{यथा हि सारः काष्ठेषु छिद्यमानेषु दृश्यते~। }\pend
\pstart\versepar \skt{तथा हि धातुर्म्मांसेषु च्छिद्यमानेषु दृश्यते~॥}\pend

\pstart\versepar \ledsidenote{\passref{3.4.7}}\skt{स्नायुभिश्च }\edtext{\skt{प्रति\-च्छ\-न्ना\-न्स\-न्त\-तांश्च}}{\lemma{\skt{प्रतिच्छन्नान्}}\Afootnote{A; \skt{प्रतिच्छन्नां}\ H}}\skt{ जरायुणा~। }\pend
\pstart\versepar \skt{श्लेष्मणा वेष्टितांश्चापि }\edtext{\edtext{\skt{कला\-भा\-गा\-न्वि\-दु\-र्ब्बुधाः}}{\lemma{\skt{कलाभागान्}}\Afootnote{\skt{कलाभागांस्तु तान्}\ A; \skt{कलाभागाम्}\ H}}}{\lemma{\skt{ब्बुधाः}}\Afootnote{H; \emph{om.}\ A}}\skt{~॥}\pend

\pstart \ledsidenote{\passref{3.4.8}}\edtext{\skt{तासां}}{\lemma{\skt{तासां}}\Afootnote{A; \skt{तासं}\ H}}\skt{ प्रथमा }\edtext{\skt{मांसधारा}}{\lemma{\skt{मांसधारा}}\Afootnote{H; \skt{मांसधरा}\ A}}\skt{ }\edtext{\skt{नाम}}{\lemma{\skt{नाम}}\Afootnote{H; \emph{om.}\ A}}\skt{~। }\edtext{\skt{तस्या}}{\lemma{\skt{तस्या}}\Afootnote{H; \skt{यस्यां}\ A}}\skt{ }\edtext{\skt{मांसे}}{\lemma{\skt{मांसे}}\Afootnote{A; \skt{मांस}\ H}}\skt{ सिरा\-स्ना\-यु\-ध\-म\-नी\-स्रो\-तसां }\edtext{\skt{नव\-न\-व\-प्र\-ता\-नानि}}{\lemma{\skt{नव\appbrk{}न\appbrk{}व\appbrk{}प्र\appbrk{}ता\appbrk{}नानि}}\Afootnote{H; \skt{प्रताना}\ A}}\skt{ भवन्ति~॥}\footnoteA{The word \emph{pratāna} is normally masculine, but here and in the Aṣṭāṅgasaṅgraha it is neuter (As. śārīrasthāna 5.20).}\footnoteA{We emend to \emph{māṃse} for sense, with the vulgate, and in the light of the of the \emph{māṃse} in 3.4.9.}\pend

\pstart\versepar \ledsidenote{\passref{3.4.9}}\edtext{\skt{भवति}}{\lemma{\skt{भवति}}\Afootnote{H; \emph{om.}\ A}}\skt{ चात्र~॥ }\pend
\pstart\versepar \skt{यथा }\edtext{\skt{बिसमृणालानि}}{\lemma{\skt{बिसमृणालानि}}\Afootnote{A; \skt{विषमृणालानि}\ H}}\skt{ विवर्द्धन्ते समन्ततः~। }\pend
\pstart\versepar \skt{भूमौ पङ्कोदकस्थानि तथा मांसे सिरादयः~॥}\edtext{}{\lemma{\skt{चात्र~॥ यथा}}\Afootnote{\skt{यथा}\ A H}}\pend

\pstart \ledsidenote{\passref{3.4.10}}\skt{द्वितीया रक्तधरा }\edtext{\skt{नाम}}{\lemma{\skt{नाम}}\Afootnote{H; \emph{om.}\ A}}\skt{ }\edtext{\skt{मांस\-स्या\-भ्य\-न्त\-र\-त\-स्तस्याः}}{\lemma{\skt{तस्याः}}\Afootnote{\skt{तस्यां}\ A; \skt{तस्या}\ H}}\skt{ }\edtext{\skt{शोणि\-त\-म्वि\-शे\-षतः}}{\lemma{\skt{विशेषतः}}\Afootnote{H; \skt{विशेषतश्च}\ A}}\skt{ }\edtext{\skt{सिरासु च}}{\lemma{\skt{सिरासु च}}\Afootnote{H; \skt{सिरासु}\ A}}\skt{ यकृत्प्लीह्नोश्च भवति~॥}\pend

\pstart\versepar \ledsidenote{\passref{3.4.11}}\skt{भवति चात्र~॥ }\pend
\pstart\versepar \skt{वृक्षाद्यथा हि प्रहतात्क्षीरिणः }\edtext{\edtext{\skt{क्षीरमास्रवेत्}}{\lemma{\skt{क्षीरम्}}\Afootnote{H; \skt{क्षीरिम्}\ A}}}{\lemma{\skt{आस्रवेत्}}\Afootnote{\skt{आवहेत्}\ A; \skt{आस्रयेत्}\ H}}\skt{~। }\pend
\pstart\versepar \skt{मांसादेवं क्षतात्क्षिप्रं शोणितं }\edtext{\skt{सम्प्रवर्तते}}{\lemma{\skt{सम्प्रवर्तते}}\Afootnote{H; \skt{संप्रसिच्यते}\ A}}\skt{~॥}\edtext{}{\lemma{\skt{भवति चात्र~॥ वृक्षाद्यथा हि}}\Afootnote{\skt{वृक्षाद्यथाभि}\ A; \skt{॥ भवति चात्र~॥ वृक्षाद्यथा हि}\ H}}\pend

\pstart \ledsidenote{\passref{3.4.12}}\skt{तृतीया मेदोधरा }\edtext{\skt{नाम}}{\lemma{\skt{नाम}}\Afootnote{H; \emph{om.}\ A}}\skt{ }\edtext{\skt{मेदस्तु}}{\lemma{\skt{मेदस्तु}}\Afootnote{H; \skt{मेदो}\ A}}\skt{ }\edtext{\skt{खलु}}{\lemma{\skt{खलु}}\Afootnote{H; \skt{हि}\ A}}\skt{ }\edtext{\edtext{\skt{सर्व्व\-भू\-ता\-ना\-मु\-द\-रे\-ष्व\-स्थिषु}}{\lemma{\skt{उदरेष्व्}}\Afootnote{H; \skt{उदरस्थमण्}\ A}}\skt{ च}}{\lemma{\skt{अस्थिषु च}}\Afootnote{H; \skt{वस्थिषु च}\ A}}\skt{ महत्सु }\edtext{\edtext{\skt{चान्तर्म्मज्जा}}{\lemma{\skt{चान्तर्}}\Afootnote{H; \skt{च}\ A}}\skt{ च}}{\lemma{\skt{म्मज्जा च}}\Afootnote{H; \skt{मज्जा}\ A}}\skt{ भवति~॥}\pend

\pstart\versepar \ledsidenote{\passref{3.4.13}}\skt{भवति चात्र~॥ }\pend
\pstart\versepar \skt{स्थूलास्थिषु विशेषेण मज्जा त्वभ्यन्तराश्रितः~। }\pend
\pstart\versepar \skt{अथेतरेषु सर्व्वेषु सरक्तं मेद }\edtext{\skt{उच्यते}}{\lemma{\skt{उच्यते}}\Afootnote{H; \skt{उच्यते~। शुद्धमांसस्य यः स्नेहः सा वसा परिकीर्तिता}\ A}}\skt{~॥}\edtext{}{\lemma{\skt{भवति चात्र~॥ स्थूलास्थिषु}}\Afootnote{\skt{स्थूलास्थिषु}\ A H}}\pend

\pstart \ledsidenote{\passref{3.4.14}}\skt{चतुर्थी श्लेष्मधरा }\edtext{\skt{नाम या}}{\lemma{\skt{नाम या}}\Afootnote{H; \emph{om.}\ A}}\skt{ सर्वसन्धिषु }\edtext{\skt{सर्व्व\-प्रा\-ण\-भृ\-ता\-म्भ\-वति}}{\lemma{\skt{सर्व्व\appbrk{}प्रा\appbrk{}ण\appbrk{}भृ\appbrk{}ताम्}}\Afootnote{\skt{प्राणभृतां}\ A; \skt{सर्व्व\appbrk{}प्रा\appbrk{}ण\appbrk{}प्र\appbrk{}ताम्}\ H}}\skt{~॥}\pend

\pstart\versepar \ledsidenote{\passref{3.4.15}}\skt{भवति चात्र~॥ }\pend
\pstart\versepar \skt{स्नेहाभ्यक्ते यथा ह्यक्षे चक्रं साधु प्रवर्त्तते~। }\pend
\pstart\versepar \skt{सन्धयः साधु वर्त्तन्ते }\edtext{\skt{संश्लिष्टाः}}{\lemma{\skt{संश्लिष्टाः}}\Afootnote{A; \skt{संश्लिष्टा}\ H}}\skt{ श्लेष्मणा तथा~॥}\edtext{}{\lemma{\skt{भवति चात्र~॥ स्नेहाभ्यक्ते}}\Afootnote{\skt{स्नेहाभ्यक्ते}\ A H}}\pend

\pstart \ledsidenote{\passref{3.4.16}}\skt{पञ्चमी पुरीषधरा }\edtext{\skt{नाम या}}{\lemma{\skt{नाम या}}\Afootnote{H; \skt{नाम}\ A}}\skt{ }\edtext{\edtext{\edtext{\edtext{\edtext{\edtext{\skt{चतु\-र्व्वि\-ध\-म\-प्य\-न्न\-मा\-मा\-श\-या\-त्प्र\-च्यु\-त\-म्प\-क्वा\-श\-यो\-प\-स्थितं}}{\lemma{\skt{चतुर्व्विधम्}}\Afootnote{H; \emph{om.}\ A}}}{\lemma{\skt{अप्य्}}\Afootnote{H; \emph{om.}\ A}}}{\lemma{\skt{अन्नम्}}\Afootnote{H; \emph{om.}\ A}}}{\lemma{\skt{आमाशयात्}}\Afootnote{H; \skt{यान्}\ A}}}{\lemma{\skt{प्रच्युतम्}}\Afootnote{H; \emph{om.}\ A}}}{\lemma{\skt{पक्वा\appbrk{}श\appbrk{}यो\appbrk{}प\appbrk{}स्थितं}}\Afootnote{H; \emph{om.}\ A}}\skt{ }\edtext{\skt{धारयति}}{\lemma{\skt{धारयति}}\Afootnote{H; \skt{तः}\ A}}\skt{~। }\edtext{\edtext{\skt{कोष्ठम्विभजति}}{\lemma{\skt{कोष्ठम्}}\Afootnote{H; \skt{कोष्ठे मलम्}\ A}}}{\lemma{\skt{विभजति}}\Afootnote{H; \skt{अभिविभजते}\ A}}\skt{ पक्वाशयस्था~॥}\pend

\pstart\versepar \ledsidenote{\passref{3.4.17}}\skt{भवति चात्र~॥ }\pend
\pstart\versepar \skt{यकृ\-त्स\-म\-न्ता\-त्को\-ष्ठञ्च तथान्त्राणि समाश्रिता~। }\pend
\pstart\versepar \skt{}\edtext{\skt{उण्डु\-क\-स्थ\-म्वि\-भ\-जते}}{\lemma{\skt{उण्डुकस्थम्}}\Afootnote{H; \skt{उण्डु(न्दु)कस्थं}\ A}}\skt{ }\edtext{\skt{मलम्मलधरा}}{\lemma{\skt{मलधरा}}\Afootnote{A; \skt{मलधराᳵ}\ H}}\skt{ }\edtext{\skt{कला}}{\lemma{\skt{कला}}\Afootnote{A; \skt{कलाः}\ H}}\skt{~॥}\footnoteA{We emend from plurals to feminine singular, with the vulgate, for \emph{kalā} and its qualifiers because of sense, because other passages describe single \emph{kalā}s, because \emph{vibhajate} is singular and because a plural would break the metre.}\edtext{}{\lemma{\skt{भवति चात्र~॥ यकृत्}}\Afootnote{\skt{यकृत्}\ A H}}\pend

\pstart \ledsidenote{\passref{3.4.18}}\skt{षष्ठी पित्तधरा }\edtext{\skt{नाम या}}{\lemma{\skt{नाम या}}\Afootnote{H; \skt{या}\ A}}\skt{ }\edtext{\edtext{\skt{चतु\-र्व्वि\-ध\-म\-शि\-त\-पी\-त\-खा\-दि\-त\-ली\-ढ\-मिति}}{\lemma{\skt{अशि\appbrk{}त\appbrk{}पी\appbrk{}त\appbrk{}खा\appbrk{}दि\appbrk{}त\appbrk{}ली\appbrk{}ढम्}}\Afootnote{\skt{अन्न\appbrk{}पा\appbrk{}न\appbrk{}मा\appbrk{}मा\appbrk{}श\appbrk{}या\appbrk{}त्प्र\appbrk{}च्युतं पक्वा\appbrk{}श\appbrk{}यो\appbrk{}प\appbrk{}स्थितं}\ A; \skt{अशि\appbrk{}त\appbrk{}पी\appbrk{}त\appbrk{}का\appbrk{}दित लीढम्}\ H}}}{\lemma{\skt{इति}}\Afootnote{H; \skt{धा}\ A}}\skt{ }\edtext{\skt{पचति}}{\lemma{\skt{पचति}}\Afootnote{H; \skt{रयति}\ A}}\skt{~॥}\pend

\pstart\versepar \ledsidenote{\passref{3.4.19}}\edtext{\skt{भवति}}{\lemma{\skt{भवति}}\Afootnote{H; \emph{om.}\ A}}\skt{ }\edtext{\skt{चात्र}}{\lemma{\skt{चात्र}}\Afootnote{H; \emph{om.}\ A}}\skt{~॥ }\pend
\pstart\versepar \skt{अशितं खादितं पीतं लीढं }\edtext{\skt{यत्कोष्ठगतं}}{\lemma{\skt{यत्}}\Afootnote{H; \emph{om.}\ A}}\skt{~॥ }\pend
\pstart\versepar \skt{}\edtext{\skt{तज्जीर्यति}}{\lemma{\skt{तज्}}\Afootnote{H; \skt{नृणाम्~। तज्}\ A}}\skt{ यथाकालं शोषि\-त\-म्पि\-त्त\-ते\-जसा~॥}\pend

\pstart \ledsidenote{\passref{3.4.20}}\skt{सप्तमी शुक्रधरा }\edtext{\skt{नाम या}}{\lemma{\skt{नाम या}}\Afootnote{H; \skt{या}\ A}}\skt{ सर्व्वप्राणिनां सर्व्व\-श\-री\-र\-व्या\-पिनी }\edtext{\skt{भवति}}{\lemma{\skt{भवति}}\Afootnote{H; \emph{om.}\ A}}\skt{~॥}\pend

\pstart\versepar \ledsidenote{\passref{3.4.21}}\edtext{\skt{भवति}}{\lemma{\skt{भवति}}\Afootnote{H; \emph{om.}\ A}}\skt{ }\edtext{\skt{चात्र}}{\lemma{\skt{चात्र}}\Afootnote{H; \emph{om.}\ A}}\skt{~॥ }\pend
\pstart\versepar \skt{यथा पयसि सर्प्पिस्तु }\edtext{\edtext{\skt{गुडश्चेक्षुरसे}}{\lemma{\skt{गुडश्}}\Afootnote{H; \skt{गूढश्}\ A}}}{\lemma{\skt{चेक्षुरसे}}\Afootnote{H; \skt{चेक्षौ रसो}\ A}}\skt{ यथा~। }\pend
\pstart\versepar \skt{शरीरेषु तथा शुक्रं नृणा\-म्वि\-द्या\-द्भि\-ष\-ग्वरः~॥}\pend

\pstart\versepar \ledsidenote{\passref{3.4.22}}\skt{द्व्यङ्गुले दक्षिणे पार्श्वे बस्तिद्वारस्य }\edtext{\skt{चाप्यवाङ्}}{\lemma{\skt{अवाङ्}}\Afootnote{H; \skt{अधः}\ A}}\skt{~। }\pend
\pstart\versepar \skt{}\edtext{\skt{मूत्रस्रोतस्ततः}}{\lemma{\skt{ततः}}\Afootnote{H; \skt{पथाच्}\ A}}\skt{ }\edtext{\skt{शुक्रं}}{\lemma{\skt{शुक्रं}}\Afootnote{H; \skt{छुक्रं}\ A}}\skt{ पुरुषस्य प्रवर्त्तते~॥}\pend

\pstart \ledsidenote{\passref{3.4.24}}\edtext{\skt{गृही\-त\-ग\-र्ब्भा\-णा\-मा\-र्त्त\-व\-व\-हानां}}{\lemma{\skt{गृही\appbrk{}त\appbrk{}ग\appbrk{}र्ब्भा\appbrk{}णाम्}}\Afootnote{H; \skt{गृही\appbrk{}त\appbrk{}ग\appbrk{}र्भा\appbrk{}णाम्}\ A}}\skt{ स्रोतसां }\edtext{\edtext{\skt{वर्त्मा\-ण्यु\-प\-रु\-द्धानि}}{\lemma{\skt{वर्त्माण्य्}}\Afootnote{H; \skt{वर्त्मान्य्}\ A}}}{\lemma{\skt{उपरुद्धानि}}\Afootnote{H; \skt{अवरुध्यन्ते}\ A}}\skt{ गर्भेण }\edtext{\skt{भवन्ति}}{\lemma{\skt{भवन्ति}}\Afootnote{H; \emph{om.}\ A}}\skt{~। }\edtext{\edtext{\skt{तस्मा\-द्गृ\-ही\-त\-ग\-र्ब्भा\-णा\-मा\-र्त्तवं}}{\lemma{\skt{गृही\appbrk{}त\appbrk{}ग\appbrk{}र्ब्भा\appbrk{}णाम्}}\Afootnote{H; \skt{गृही\appbrk{}त\appbrk{}ग\appbrk{}र्भा\appbrk{}णाम्}\ A}}\skt{ न तु}}{\lemma{\skt{आर्त्तवं न तु}}\Afootnote{H; \skt{आर्तवं न}\ A}}\skt{ दृश्यते~। }\edtext{\edtext{\edtext{\skt{तत\-स्त\-द\-धः\-प्र\-ति\-ह\-त\-ग\-त्यू\-र्द्ध्व\-भा\-ग\-मु\-परि}}{\lemma{\skt{तद\appbrk{}धः\appbrk{}प्र\appbrk{}ति\appbrk{}ह\appbrk{}त\appbrk{}गत्य्}}\Afootnote{H; \skt{तदधः प्रतिहतम्}\ A}}}{\lemma{\skt{ऊर्द्ध्वभागम्}}\Afootnote{H; \skt{ऊर्ध्वमागतम्}\ A}}}{\lemma{\skt{उपरि}}\Afootnote{H; \skt{अपरं}\ A}}\skt{ }\edtext{\skt{चोप\-ची\-य\-मा\-न\-मु\-दरे}}{\lemma{\skt{उदरे}}\Afootnote{H; \skt{अपरे}\ A}}\skt{ }\edtext{\skt{मुच्यते}}{\lemma{\skt{मुच्यते}}\Afootnote{H; \skt{त्यभिधीयते}\ A}}\skt{~। शेषं }\edtext{\edtext{\skt{चोर्द्ध्व\-त\-र\-भा\-ग\-मा\-ग\-त\-म्प\-यो\-ध\-रा\-व\-भि\-प्र\-प\-द्यते}}{\lemma{\skt{चोर्द्ध्व\appbrk{}त\appbrk{}र\appbrk{}भा\appbrk{}गम्}}\Afootnote{\skt{चोर्ध्वतरम्}\ A; \skt{चोर्द्ध\appbrk{}त\appbrk{}र\appbrk{}भा\appbrk{}गम्}\ H}}}{\lemma{\skt{अभिप्रपद्यते}}\Afootnote{\skt{अभि\appbrk{}प्र\appbrk{}ति\appbrk{}प\appbrk{}द्यते}\ A; \skt{अभिप्रपद्येते}\ H}}\skt{~। }\edtext{\skt{तस्माद्गर्ब्भिण्यः}}{\lemma{\skt{गर्ब्भिण्यः}}\Afootnote{H; \skt{गर्भिण्यः}\ A}}\skt{ पीनो\-न्न\-त\-प\-यो\-धरा भवन्ति~॥}\pend

\pstart \ledsidenote{\passref{3.4.25}}\edtext{\skt{यकृत्प्लीहानौ}}{\lemma{\skt{यकृत्प्लीहानौ}}\Afootnote{H; \skt{गर्भस्य यकृत्प्लीहानौ}\ A}}\skt{ शोणितजौ~। शोणि\-त\-फे\-न\-प्र\-भवः फुप्फुसः~। शोणि\-त\-कि\-ट्ट\-प्र\-भवः उण्डुकः~॥}\pend

\pstart\versepar \ledsidenote{\passref{3.4.26}}\edtext{\skt{भवति}}{\lemma{\skt{भवति}}\Afootnote{H; \emph{om.}\ A}}\skt{ }\edtext{\skt{चात्र}}{\lemma{\skt{चात्र}}\Afootnote{H; \emph{om.}\ A}}\skt{~॥ }\pend
\pstart\versepar \skt{असृजः श्लेष्मणश्चापि यः प्रसादः परो मतः~। }\pend
\pstart\versepar \skt{तं पच्यमानं पित्तेन }\edtext{\skt{वात\-श्चा\-प्य\-नु\-धा\-वति}}{\lemma{\skt{वातश्}}\Afootnote{H; \skt{वायुश्}\ A}}\skt{~॥}\pend

\pstart\versepar \ledsidenote{\passref{3.4.27}}\skt{ततोऽस्यान्त्राणि जायन्ते }\edtext{\skt{गुदं}}{\lemma{\skt{गुदं}}\Afootnote{A; \skt{गुदे}\ H}}\skt{ बस्तिश्च }\edtext{\skt{देहिनः}}{\lemma{\skt{देहिनः}}\Afootnote{H; \skt{देहि}\ A}}\skt{~। }\pend
\pstart\versepar \skt{}\edtext{\skt{तत्रास्य}}{\lemma{\skt{तत्रास्य}}\Afootnote{H; \emph{om.}\ A}}\skt{ }\edtext{\skt{मथ्यमानस्य}}{\lemma{\skt{मथ्यमानस्य}}\Afootnote{H; \skt{नः~। उदरे}\ A}}\skt{ }\edtext{\skt{धाम्यमानस्य रुक्मवत्}}{\lemma{\skt{धाम्यमानस्य रुक्मवत्}}\Afootnote{\skt{पच्य\appbrk{}मा\appbrk{}ना\appbrk{}ना\appbrk{}मा\appbrk{}ध्मा\appbrk{}ना\appbrk{}द्रु\appbrk{}क्म\appbrk{}सा\appbrk{}र\appbrk{}वत्}\ A; \skt{धाम्यमानस्य रुग्मवत्}\ H}}\skt{~॥}\footnoteA{We emend H's reading \emph{rugmat} because the sandhi \emph{-km-} is awkward and likely to have been corrected by a scribe. \emph{Pace} the unique occurrence of \emph{rugma} in the Mañjūśrīmūlakalpa (Edgerton, Dictionary, 455).}\pend

\pstart \ledsidenote{\passref{3.4.28ab}}\edtext{\skt{जन्तोः}}{\lemma{\skt{जन्तोः}}\Afootnote{H; \skt{कफ\appbrk{}शो\appbrk{}णि\appbrk{}त\appbrk{}मांसा}\ A}}\skt{ }\edtext{\skt{संजायते}}{\lemma{\skt{संजायते}}\Afootnote{H; \skt{नां साराज्}\ A}}\skt{ जिह्वा }\edtext{\skt{यया}}{\lemma{\skt{यया}}\Afootnote{H; \skt{प्र}\ A}}\skt{ }\edtext{\skt{वेदयते}}{\lemma{\skt{वेदयते}}\Afootnote{H; \skt{जायते}\ A}}\skt{ }\edtext{\skt{रसान्}}{\lemma{\skt{रसान्}}\Afootnote{H; \emph{om.}\ A}}\skt{~॥}\pend

\pstart \ledsidenote{\passref{3.4.31.2}}\skt{शोणि\-त\-क\-फ\-प्र\-सा\-दजं हृदयं~। यदाश्रया हि धमन्यः }\edtext{\skt{प्राण\-व\-हा\-स्त\-द्वि\-शे\-षेण}}{\lemma{\skt{तद्विशेषेण}}\Afootnote{H; \skt{तस्याधो वामतः प्लीहा फुप्फुसश्च, दक्षिणतो यकृत्क्लोम च; तद्विशेषेण}\ A}}\skt{ }\edtext{\skt{चेत\-ना\-स्था\-न\-मु\-क्तम्}}{\lemma{\skt{उक्तम्}}\Afootnote{H; \skt{अतस्}\ A}}\skt{~॥ तस्मिं\-स्त\-म\-सा\-वृते सर्व्वप्राणिनः स्वपन्ति~॥}\pend

\pstart \ledsidenote{\passref{3.4.33}}\skt{निद्रान्तु वैष्णवीं पाप्मा\-न\-मु\-प\-दि\-शन्ति~। सा }\edtext{\skt{स्वभावतः}}{\lemma{\skt{स्वभावतः}}\Afootnote{H; \skt{स्वभावत}\ A}}\skt{ }\edtext{\skt{सर्व्वप्राणिनो}}{\lemma{\skt{सर्व्वप्राणिनो}}\Afootnote{H; \skt{एव सर्वप्राणिनो}\ A}}\skt{ }\edtext{\skt{निस्पृशति}}{\lemma{\skt{निस्पृशति}}\Afootnote{H; \skt{ऽभिस्पृशति}\ A}}\skt{~। तत्र यदा संज्ञावहानि स्रोतान्सि तमोभूयिष्ठः श्लेष्मा प्रतिपद्यते तदा तामसी नाम निद्रा }\edtext{\skt{भवत्य्}}{\lemma{\skt{भवत्य्}}\Afootnote{H; \skt{संभवत्य्}\ A}}\skt{~। अनवबोधिनी सा प्रलयकाले~। }\edtext{\skt{तमो\-भू\-यि\-ष्ठा\-ना\-म\-हःसु च}}{\lemma{\skt{अहःसु च}}\Afootnote{H; \skt{अहःसु}\ A}}\skt{ }\edtext{\skt{निशासु}}{\lemma{\skt{निशासु}}\Afootnote{H; \skt{निशासु च}\ A}}\skt{ भवति~। }\edtext{\skt{रजो\-भू\-यि\-ष्ठा\-ना\-म\-नि\-मित्तं}}{\lemma{\skt{अनिमित्तं}}\Afootnote{A; \skt{अनिमित्त}\ H}}\skt{~। सत्त्व\-भू\-यि\-ष्ठा\-ना\-म\-र्ध\-रात्रे~। }\edtext{\skt{क्षीण\-श्ले\-ष्म\-णा\-म\-नि\-ल\-ब\-हु\-ला\-ना\-म्म\-नो\-भि\-घा\-ताच्च}}{\lemma{\skt{मनोभिघाताच्च}}\Afootnote{H; \skt{मनः\appbrk{}श\appbrk{}री\appbrk{}रा\appbrk{}भि\appbrk{}ता\appbrk{}प\appbrk{}वतां च नैव, सा}\ A}}\skt{ }\edtext{\skt{वैकारिकी}}{\lemma{\skt{वैकारिकी}}\Afootnote{H; \skt{वैकारिकी भवति}\ A}}\skt{~॥}\pend

\pstart\versepar \ledsidenote{\passref{3.4.34}}\edtext{\skt{भवन्ति}}{\lemma{\skt{भवन्ति}}\Afootnote{H; \skt{भवतश्}\ A}}\skt{ चात्र~॥ }\pend
\pstart\versepar \skt{हृद\-य\-ञ्चे\-त\-ना\-स्था\-न\-मुक्तं सुश्रुत देहिनाम्~। }\pend
\pstart\versepar \skt{}\edtext{\skt{तमोभिभूते}}{\lemma{\skt{तमोभिभूते}}\Afootnote{H; \skt{तमोऽभिभूते}\ A}}\skt{ तस्मिंस्तु निद्रा विशति }\edtext{\skt{देहिनः}}{\lemma{\skt{देहिनः}}\Afootnote{H; \skt{देहिनम्}\ A}}\skt{~॥}\pend

\pstart\versepar \ledsidenote{\passref{3.4.35}}\skt{निद्राहेतुस्तमः सत्त्वं बोधने हेतुरुच्यते~। }\pend
\pstart\versepar \skt{स्वभाव एव वा हेतु\-र्ग्ग\-री\-या\-न्प\-रि\-की\-र्त्यते~॥}\pend

\pstart\versepar \ledsidenote{\passref{3.4.36}}\skt{पूर्व्व\-दे\-हा\-नु\-भू\-तांस्तु भूतात्मा }\edtext{\skt{स्वपिति}}{\lemma{\skt{स्वपिति}}\Afootnote{H; \skt{स्वपतः}\ A}}\skt{ प्रभुः~॥ }\pend
\pstart\versepar \skt{रजोयुक्तेन मनसा }\edtext{\edtext{\edtext{\skt{गृह्णी\-तेऽ\-र्था\-ञ्शु\-भा\-शु\-भान्}}{\lemma{\skt{गृह्णीते}}\Afootnote{H; \skt{गृह्णात्य्}\ A}}}{\lemma{\skt{ऽर्थाञ्}}\Afootnote{A; \skt{र्थां}\ H}}}{\lemma{\skt{शुभाशुभान्}}\Afootnote{A; \skt{शुभाशुभां}\ H}}\skt{~॥}\pend

\pstart\versepar \ledsidenote{\passref{3.4.37}}\skt{करणानां तु }\edtext{\skt{वैगुण्ये}}{\lemma{\skt{वैगुण्ये}}\Afootnote{H; \skt{वैकल्ये}\ A}}\skt{ }\edtext{\skt{तम\-सा\-भि\-प्र\-व\-र्त्तिते}}{\lemma{\skt{तम\appbrk{}सा\appbrk{}भि\appbrk{}प्र\appbrk{}व\appbrk{}र्त्तिते}}\Afootnote{\skt{तम\appbrk{}सा\appbrk{}भि\appbrk{}प्र\appbrk{}व\appbrk{}र्धिते}\ A; \skt{तम\appbrk{}सा\appbrk{}भि\appbrk{}प्र\appbrk{}व\appbrk{}र्त्तते}\ H}}\skt{~। }\pend
\pstart\versepar \skt{अस्वपन्नपि भूतात्मा प्रसुप्त }\edtext{\skt{इति}}{\lemma{\skt{इति}}\Afootnote{H; \skt{इव}\ A}}\skt{ चोच्यते~॥}\pend

\pstart \ledsidenote{\passref{3.4.38}}\edtext{\skt{सर्व्वर्त्तुषु च}}{\lemma{\skt{सर्व्वर्त्तुषु च}}\Afootnote{H; \skt{सर्वर्तुषु}\ A}}\skt{ }\edtext{\skt{दिवास्वप्नः}}{\lemma{\skt{दिवास्वप्नः}}\Afootnote{\skt{दिवास्वापः}\ A; \skt{दिवास्वप्नम्}\ H}}\skt{ प्रतिषिद्धोऽन्यत्र ग्रीष्मात्~। प्रतिषिद्धेष्वपि }\edtext{\skt{प्रतिषिद्धं}}{\lemma{\skt{प्रतिषिद्धं}}\Afootnote{\skt{तु}\ A; \skt{प्रतिषिद्धो}\ H}}\skt{ }\edtext{\skt{बाल\-वृ\-द्ध\-क्ष\-त\-क्षी\-ण\-म\-द्य\-पी\-त\-स्त्री\-या\-न\-वा\-ह\-ना\-ध्व\-क\-र्म्म\-प\-रि\-श्रा\-न्ता\-ना\-म\-भु\-क्त\-वतां}}{\lemma{\skt{बाल\appbrk{}वृ\appbrk{}द्ध\appbrk{}क्ष\appbrk{}त\appbrk{}क्षी\appbrk{}ण\appbrk{}म\appbrk{}द्य\appbrk{}पी\appbrk{}त\appbrk{}स्त्री\appbrk{}या\appbrk{}न\appbrk{}वा\appbrk{}ह\appbrk{}ना\appbrk{}ध्व\appbrk{}क\appbrk{}र्म्म\appbrk{}प\appbrk{}रि\appbrk{}श्रा\appbrk{}न्ता\appbrk{}नाम्}}\Afootnote{H; \skt{बाल\appbrk{}वृ\appbrk{}द्ध\appbrk{}स्त्री\appbrk{}क\appbrk{}र्शि\appbrk{}त\appbrk{}क्ष\appbrk{}त\appbrk{}क्षी\appbrk{}ण\appbrk{}म\appbrk{}द्य\appbrk{}नि\appbrk{}त्य\appbrk{}या\appbrk{}न\appbrk{}वा\appbrk{}ह\appbrk{}ना\appbrk{}ध्व\appbrk{}क\appbrk{}र्म\appbrk{}प\appbrk{}रि\appbrk{}श्रा\appbrk{}न्ता\appbrk{}नाम्}\ A}}\skt{ }\edtext{\skt{मेदो\-मा\-रु\-त\-क\-फ\-र\-स\-क्षी\-णानां}}{\lemma{\skt{मेदो\appbrk{}मा\appbrk{}रु\appbrk{}त\appbrk{}क\appbrk{}फ\appbrk{}र\appbrk{}स\appbrk{}क्षी\appbrk{}णानां}}\Afootnote{H; \skt{मेदः\appbrk{}स्वे\appbrk{}द\appbrk{}क\appbrk{}फ\appbrk{}र\appbrk{}स\appbrk{}र\appbrk{}क्त\appbrk{}क्षी\appbrk{}णा\appbrk{}ना\appbrk{}म\appbrk{}जी\appbrk{}र्णिनां च}\ A}}\skt{ }\edtext{\skt{मुहू\-र्त्त\-म्प्र\-शं\-सन्ति}}{\lemma{\skt{प्रशंसन्ति}}\Afootnote{H; \skt{दिवा\appbrk{}स्व\appbrk{}प\appbrk{}न\appbrk{}म\appbrk{}प्र\appbrk{}ति\appbrk{}षि\appbrk{}द्धम्}\ A}}\skt{~। रात्रावपि }\edtext{\skt{खलु}}{\lemma{\skt{खलु}}\Afootnote{H; \emph{om.}\ A}}\skt{ जागरितवतां }\edtext{\edtext{\skt{जाग\-र\-ण\-का\-ला\-र्द्ध\-न्दिवा}}{\lemma{\skt{जाग\appbrk{}र\appbrk{}ण\appbrk{}का\appbrk{}ला\appbrk{}र्द्धन्}}\Afootnote{H; \skt{जाग\appbrk{}रि\appbrk{}त\appbrk{}का\appbrk{}ला\appbrk{}द\appbrk{}र्धम्}\ A}}}{\lemma{\skt{दिवा}}\Afootnote{H; \skt{इष्यते दिवा}\ A}}\skt{ }\edtext{\skt{स्वपेत्}}{\lemma{\skt{स्वपेत्}}\Afootnote{H; \skt{स्वपनम्}\ A}}\skt{~। }\edtext{\edtext{\skt{विकृतिस्तु}}{\lemma{\skt{विकृतिस्}}\Afootnote{H; \skt{विकृति}\ A}}}{\lemma{\skt{तु}}\Afootnote{H; \skt{र्हि}\ A}}\skt{ दिवास्वप्नो नाम~। तत्र स्वपतामधर्म्मः सर्व्व\-दो\-ष\-प्र\-को\-पश्च }\edtext{\edtext{\skt{तत्प्र\-को\-पा\-त्का\-स\-श्वा\-स\-प्र\-ति\-श्या\-य\-शि\-रो\-गौ\-र\-वा\-ङ्ग\-म\-र्द्दा\-रो\-च\-का\-ग्नि\-दौ\-र्व्व\-ल्यानि}}{\lemma{\skt{तत्प्रकोपात्}}\Afootnote{H; \skt{तत्प्रकोपाच्च}\ A}}}{\lemma{\skt{कास\appbrk{}श्वा\appbrk{}स\appbrk{}प्र\appbrk{}ति\appbrk{}श्या\appbrk{}य\appbrk{}शि\appbrk{}रो\appbrk{}गौ\appbrk{}र\appbrk{}वा\appbrk{}ङ्ग\appbrk{}म\appbrk{}र्द्दा\appbrk{}रो\appbrk{}च\appbrk{}का\appbrk{}ग्नि\appbrk{}दौ\appbrk{}र्व्व\appbrk{}ल्यानि}}\Afootnote{H; \skt{कास\appbrk{}श्वा\appbrk{}स\appbrk{}प्र\appbrk{}ति\appbrk{}श्या\appbrk{}य\appbrk{}शि\appbrk{}रो\appbrk{}गौ\appbrk{}र\appbrk{}वा\appbrk{}ङ्ग\appbrk{}म\appbrk{}र्दा\appbrk{}रो\appbrk{}च\appbrk{}क\appbrk{}ज्व\appbrk{}रा\appbrk{}ग्नि\appbrk{}दौ\appbrk{}र्ब\appbrk{}ल्यानि}\ A}}\skt{ भवन्ति~। रात्रावपि }\edtext{\skt{खलु}}{\lemma{\skt{खलु}}\Afootnote{H; \emph{om.}\ A}}\skt{ }\edtext{\edtext{\skt{जाग\-रि\-त\-व\-ता\-न्त\-न्नि\-मि\-त्तास्त}}{\lemma{\skt{तन्निमित्तास्}}\Afootnote{H; \skt{वात\appbrk{}पि\appbrk{}त्त\appbrk{}नि\appbrk{}मित्ता}\ A}}}{\lemma{\skt{त}}\Afootnote{H; \skt{स्त}\ A}}\skt{ }\edtext{\skt{एव}}{\lemma{\skt{एव}}\Afootnote{H; \skt{एवो}\ A}}\skt{ }\edtext{\skt{दोषाः}}{\lemma{\skt{दोषाः}}\Afootnote{H; \skt{प}\ A}}\skt{ }\edtext{\skt{प्रादुर्भवन्ति}}{\lemma{\skt{प्रादुर्भवन्ति}}\Afootnote{\skt{द्रवा भवन्ति}\ A; \skt{प्रादुर्ब्भवन्ति}\ H}}\skt{~॥}\pend

\pstart\versepar \ledsidenote{\passref{3.4.39}}\skt{भवन्ति चात्र }\edtext{\skt{श्लोकाः}}{\lemma{\skt{श्लोकाः}}\Afootnote{H; \emph{om.}\ A}}\skt{~॥ }\pend
\pstart\versepar \skt{तस्मान्न जागृयाद्रात्रौ दिवास्वप्नञ्च वर्ज्जयेत्~। }\pend
\pstart\versepar \skt{ज्ञात्वा }\edtext{\edtext{\skt{दोषकरावेतौ}}{\lemma{\skt{दोषकराव्}}\Afootnote{H; \skt{दोषकरा}\ A}}}{\lemma{\skt{एतौ}}\Afootnote{H; \skt{वेतौ}\ A}}\skt{ बुधः }\edtext{\skt{स्वप्नं}}{\lemma{\skt{स्वप्नं}}\Afootnote{A; \skt{स्वप्न}\ H}}\skt{ }\edtext{\skt{मितञ्चरेत्}}{\lemma{\skt{मितञ्}}\Afootnote{A; \skt{मितश्}\ H}}\skt{~॥}\pend

\pstart\versepar \ledsidenote{\passref{3.4.40}}\edtext{\skt{तथा ह्यरोगः}}{\lemma{\skt{तथा ह्य्}}\Afootnote{H; \emph{om.}\ A}}\skt{ सुमना }\edtext{\skt{बलवर्ण्णान्वितो}}{\lemma{\skt{बलवर्ण्णान्वितो}}\Afootnote{\skt{ह्येवं बलवर्णान्वितो}\ A; \skt{वलवर्ण्णान्विता}\ H}}\skt{ }\edtext{\skt{भृशं}}{\lemma{\skt{भृशं}}\Afootnote{H; \skt{वृषः}\ A}}\skt{~। }\pend
\pstart\versepar \skt{}\edtext{\skt{नातिस्थूलकृशः}}{\lemma{\skt{नातिस्थूलकृशः}}\Afootnote{A; \skt{नातिस्थूलकृषः}\ H}}\skt{ श्रीमान्नरो }\edtext{\skt{जीवेत्समाशतं}}{\lemma{\skt{समाशतं}}\Afootnote{H; \skt{समाः शतम्}\ A}}\skt{~॥}\pend

\pstart\versepar \ledsidenote{\passref{3.4.56}}\skt{मूर्च्छा पित्ततमःप्राया रजः\-पि\-त्ता\-नि\-ला\-द्भ्रमः~। }\pend
\pstart\versepar \skt{}\edtext{\skt{तमो\-वा\-त\-क\-फा\-त्तन्द्रा}}{\lemma{\skt{तमो\appbrk{}वा\appbrk{}त\appbrk{}क\appbrk{}फा\appbrk{}त्तन्द्रा}}\Afootnote{A; \skt{तमोवातकफा तन्द्रा}\ H}}\skt{ निद्रा }\edtext{\skt{श्लेष्मतमोभवा}}{\lemma{\skt{श्लेष्मतमोभवा}}\Afootnote{A; \skt{श्लेष्मतमोद्भवा}\ H}}\skt{~॥}\pend

\pstart\versepar \ledsidenote{\passref{3.4.42}}\skt{निद्रानाशो }\edtext{\edtext{\skt{निला\-त्पि\-त्ता\-न्म\-न\-स्ता\-पा\-त्क्ष\-या\-दपि}}{\lemma{\skt{निलात्}}\Afootnote{H; \skt{ऽनिलात्}\ A}}}{\lemma{\skt{पित्तान्}}\Afootnote{A; \skt{पित्तात्}\ H}}\skt{~। }\pend
\pstart\versepar \skt{}\edtext{\skt{अभिघाताच्च}}{\lemma{\skt{अभिघाताच्च}}\Afootnote{H; \skt{संभ\appbrk{}व\appbrk{}त्य\appbrk{}भि\appbrk{}घा\appbrk{}ताच्च}\ A}}\skt{ }\edtext{\skt{भवति}}{\lemma{\skt{भवति}}\Afootnote{H; \emph{om.}\ A}}\skt{ प्रत्यनीकैः }\edtext{\skt{स}}{\lemma{\skt{स}}\Afootnote{H; \skt{प्र}\ A}}\skt{ शाम्यति~॥}\pend

\pstart\versepar \ledsidenote{\passref{3.4.47}}\skt{निद्रातियोगे वमनं हितं संशोधनानि च~। }\pend
\pstart\versepar \skt{लङ्घनं रक्तमोक्षश्च मनो\-व्या\-कु\-ल\-नानि च~॥}\pend

\pstart\versepar \ledsidenote{\passref{3.4.48}}\skt{कफ\-मे\-दो\-वि\-षा\-र्त्तानां रात्रौ जागरणं हितम्~। }\pend
\pstart\versepar \skt{दिवास्वप्नश्च }\edtext{\skt{तृट्छू\-ल\-हि\-क्का\-जी\-र्ण्णा\-ति\-सा\-रिणां}}{\lemma{\skt{तृट्छू\appbrk{}ल\appbrk{}हि\appbrk{}क्का\appbrk{}जी\appbrk{}र्ण्णा\appbrk{}ति\appbrk{}सा\appbrk{}रिणां}}\Afootnote{H; \skt{तृट्शू\appbrk{}ल\appbrk{}हि\appbrk{}क्का\appbrk{}जी\appbrk{}र्णा\appbrk{}ति\appbrk{}सा\appbrk{}रि\appbrk{}णाम्}\ A}}\skt{~॥}\pend

\pstart\versepar \ledsidenote{\passref{3.4.49}}\edtext{\edtext{\skt{इन्द्रि\-या\-र्थे\-ष्व\-सं\-वि\-त्ति\-र्ज्जृ\-म्भनं}}{\lemma{\skt{असंवित्तिर्}}\Afootnote{H; \skt{असंप्राप्तिर्}\ A}}}{\lemma{\skt{ज्जृम्भनं}}\Afootnote{H; \skt{गौरवं}\ A}}\skt{ }\edtext{\skt{गौरवं}}{\lemma{\skt{गौरवं}}\Afootnote{H; \skt{जृम्भणं}\ A}}\skt{ क्लमः~। }\pend
\pstart\versepar \skt{}\edtext{\skt{निद्रार्त्तस्यैव}}{\lemma{\skt{निद्रार्त्तस्यैव}}\Afootnote{H; \skt{निद्रार्तस्येव}\ A}}\skt{ }\edtext{\skt{यस्येह}}{\lemma{\skt{यस्येह}}\Afootnote{H; \skt{यस्येहा}\ A}}\skt{ तस्य }\edtext{\skt{तन्द्रां}}{\lemma{\skt{तन्द्रां}}\Afootnote{A; \skt{तन्द्रा}\ H}}\skt{ विनिर्द्दिशेत्~॥}\footnoteA{The vulgate has an additional six verses at this point that do not appear in witness H. Or they may be a scribal omission by a saut du même au même by the scribe of witness H from the word \emph{vinirdiśet} at the end of 49 to the same word at the end of 55. The sense of these verses does not fit well in the context; in the eleventh--twelfth centuries the commentators Gayadāsa and Ḍalhaṇa recorded doubts about their inclusion. Further evidence from manuscripts that transmit early readings may clarify this issue.}\pend

\pstart \ledsidenote{\passref{3.4.57}}\edtext{\skt{गर्भस्य}}{\lemma{\skt{गर्भस्य}}\Afootnote{A; \skt{गर्ब्भस्य}\ H}}\skt{ खलु }\edtext{\edtext{\skt{रस\-गु\-ण\-नि\-मि\-त्ता\-न्मा\-तु\-रा\-हा\-र\-नि\-मि\-त्ताच्च}}{\lemma{\skt{रस\appbrk{}गु\appbrk{}ण\appbrk{}नि\appbrk{}मि\appbrk{}त्ता\appbrk{}न्मा\appbrk{}तुर्}}\Afootnote{\skt{रसनिमित्ता मारुता}\ A; \skt{रस\appbrk{}गु\appbrk{}ण\appbrk{}नि\appbrk{}मि\appbrk{}त्ता\appbrk{}त्मा\appbrk{}त्मा\appbrk{}तुर्}\ H}}}{\lemma{\skt{आहा\appbrk{}र\appbrk{}नि\appbrk{}मि\appbrk{}त्ताच्च}}\Afootnote{H; \skt{ध्माननिमित्ता च}\ A}}\skt{ }\edtext{\skt{परि\-वृ\-द्धि\-र्भ\-वति}}{\lemma{\skt{भवति}}\Afootnote{A; \skt{ब्भवति}\ H}}\skt{~॥ }\footnoteA{We emend the final consonant \emph{-t} of H in \emph{rasaguṇanimittān} for correct sandhi.}\pend

\pstart\versepar \ledsidenote{\passref{3.4.60}}\edtext{\skt{भवति}}{\lemma{\skt{भवति}}\Afootnote{H; \emph{om.}\ A}}\skt{ }\edtext{\skt{चात्र}}{\lemma{\skt{चात्र}}\Afootnote{H; \emph{om.}\ A}}\skt{~॥ }\pend
\pstart\versepar \skt{दृष्टिश्च }\edtext{\skt{रोमकूपाश्च}}{\lemma{\skt{रोमकूपाश्च}}\Afootnote{H; \skt{रोमकूपाश्च न}\ A}}\skt{ }\edtext{\skt{वर्द्धन्ते न}}{\lemma{\skt{वर्द्धन्ते न}}\Afootnote{H; \skt{वर्धन्ते}\ A}}\skt{ कदाचन~। }\pend
\pstart\versepar \skt{ध्रुवाण्येतानि मर्त्त्यानामिति धन्वन्तरेर्म्मतं~॥}\pend

\pstart\versepar \ledsidenote{\passref{3.4.61}}\skt{शरीरे }\edtext{\skt{क्षीयमाणेऽपि}}{\lemma{\skt{ऽपि}}\Afootnote{A; \skt{पि}\ H}}\skt{ }\edtext{\skt{वर्द्धेते}}{\lemma{\skt{वर्द्धेते}}\Afootnote{A; \skt{वर्द्धते}\ H}}\skt{ }\edtext{\edtext{\skt{द्वाविमौ}}{\lemma{\skt{द्वाव्}}\Afootnote{H; \skt{द्वा}\ A}}}{\lemma{\skt{इमौ}}\Afootnote{H; \skt{विमौ}\ A}}\skt{ सदा~। }\pend
\pstart\versepar \skt{स्वभा\-व\-म्प्र\-कृ\-ति\-ङ्कृत्वा नखकेशाविति स्थितिः~॥}\pend

\pstart \ledsidenote{\passref{3.4.62}}\edtext{\skt{तिस्रः}}{\lemma{\skt{तिस्रः}}\Afootnote{H; \skt{सप्त}\ A}}\skt{ प्रकृतयो भवन्ति }\edtext{\skt{वात\-पि\-त्त\-क\-फ\-नि\-मित्ताः}}{\lemma{\skt{वात\appbrk{}पि\appbrk{}त्त\appbrk{}क\appbrk{}फ\appbrk{}नि\appbrk{}मित्ताः}}\Afootnote{H; \skt{दोषैः पृथक्, द्विशः, समस्तैश्च}\ A}}\skt{~॥}\pend

\pstart\versepar \ledsidenote{\passref{3.4.63}}\edtext{\skt{शुक्ल\-शो\-णि\-त\-सं\-योगे यो}}{\lemma{\skt{शुक्ल\appbrk{}शो\appbrk{}णि\appbrk{}त\appbrk{}सं\appbrk{}योगे यो}}\Afootnote{\skt{शुक्र\appbrk{}शो\appbrk{}णि\appbrk{}त\appbrk{}सं\appbrk{}योगे यो}\ A; \skt{भवति चात्र शुक्ल\appbrk{}शो\appbrk{}णि\appbrk{}त\appbrk{}सं\appbrk{}योगे यो}\ H}}\skt{ }\edtext{\edtext{\skt{दोषस्तूत्कटो}}{\lemma{\skt{दोषस्}}\Afootnote{H; \skt{भवेद्दोष}\ A}}}{\lemma{\skt{तूत्कटो}}\Afootnote{H; \skt{उत्कटः}\ A}}\skt{ }\edtext{\skt{भवेत्}}{\lemma{\skt{भवेत्}}\Afootnote{H; \emph{om.}\ A}}\skt{~। }\pend
\pstart\versepar \skt{प्रकृतिर्ज्जायते तेन }\edtext{\skt{तासां मे}}{\lemma{\skt{तासां मे}}\Afootnote{H; \skt{तस्या मे}\ A}}\skt{ लक्षणं शृणु~॥}\pend

\pstart \ledsidenote{\passref{3.4.64}}\skt{तत्र }\edtext{\edtext{\skt{वात\-प्र\-कृ\-ति\-र्जा\-ग\-रूकः}}{\lemma{\skt{वातप्रकृतिर्}}\Afootnote{A; \skt{वात प्रकृति}\ H}}}{\lemma{\skt{जागरूकः}}\Afootnote{H; \skt{प्रजागरूकः}\ A}}\skt{ शीतद्वेषी }\edtext{\edtext{\skt{दुर्भ\-ग\-स्ते\-नोऽ\-नार्यो}}{\lemma{\skt{दुर्भगस्तेनो}}\Afootnote{\skt{दुर्भगः स्तेनो}\ A; \skt{दुर्ब्भगस्तेनो}\ H}}}{\lemma{\skt{ऽनार्यो}}\Afootnote{\emph{om.}\ A; \skt{नार्यो}\ H}}\skt{ }\edtext{\skt{मत्सरो}}{\lemma{\skt{मत्सरो}}\Afootnote{\skt{मत्सर्यनार्यो}\ A; \skt{मस्तरो}\ H}}\skt{ }\edtext{\skt{गान्धर्व्वः}}{\lemma{\skt{गान्धर्व्वः}}\Afootnote{H; \skt{गान्धर्वचित्तः}\ A}}\skt{ }\edtext{\skt{स्फुटितचरणो}}{\lemma{\skt{स्फुटितचरणो}}\Afootnote{\skt{स्फुटि\appbrk{}त\appbrk{}क\appbrk{}र\appbrk{}च\appbrk{}रणो}\ A; \skt{स्फुटितचर}\ H}}\skt{ }\edtext{\skt{रूक्षा\-ल्प\-के\-श\-स्मश्रुः}}{\lemma{\skt{रूक्षा\appbrk{}ल्प\appbrk{}के\appbrk{}श\appbrk{}स्मश्रुः}}\Afootnote{\skt{ऽल्परूक्षश्मश्रु}\ A; \skt{णौ रूक्षा\appbrk{}ल्प\appbrk{}के\appbrk{}श\appbrk{}स्मश्रुः}\ H}}\skt{ }\edtext{\skt{कठिनक्रोधी}}{\lemma{\skt{कठिनक्रोधी}}\Afootnote{H; \skt{नखकेशः क्राथी}\ A}}\skt{ दन्तखादी }\edtext{\edtext{\skt{अल्प\-ब\-लोऽ\-ल्पा\-युश्च}}{\lemma{\skt{अल्पबलो}}\Afootnote{H; \emph{om.}\ A}}}{\lemma{\skt{ऽल्पायुश्च}}\Afootnote{H; \skt{च}\ A}}\skt{ भवति~॥}\pend

\pstart\versepar \ledsidenote{\passref{3.4.65}}\edtext{\skt{भवति}}{\lemma{\skt{भवति}}\Afootnote{H; \emph{om.}\ A}}\skt{ }\edtext{\skt{चात्र}}{\lemma{\skt{चात्र}}\Afootnote{H; \emph{om.}\ A}}\skt{~॥ }\pend
\pstart\versepar \skt{}\edtext{\skt{अधृ\-ति\-र\-दृ\-ढ\-सौ\-हृदः}}{\lemma{\skt{अदृढसौहृदः}}\Afootnote{A; \skt{अदृ\appbrk{}ढा\appbrk{}सौ\appbrk{}हृ\appbrk{}दयः}\ H}}\skt{ कृतघ्नः}\footnoteA{We emend to the vulgate's \emph{-asauhṛdaḥ} for the \emph{Puṣpitāgrā} metre.}\skt{ }\pend
\pstart\versepar \skt{}\edtext{\skt{कृषपरुषो}}{\lemma{\skt{कृषपरुषो}}\Afootnote{H; \skt{कृशपरुषो}\ A}}\skt{ धमनीततः प्रलापी~। }\pend
\pstart\versepar \skt{}\edtext{\skt{द्रुत\-ग\-ति\-र\-ट\-नोऽ\-न\-व\-स्थि\-तात्मा}}{\lemma{\skt{ऽनवस्थितात्मा}}\Afootnote{A; \skt{नवस्थितात्मा}\ H}}\skt{ }\pend
\pstart\versepar \skt{}\edtext{\skt{नभसि च}}{\lemma{\skt{नभसि च}}\Afootnote{H; \skt{वियति च}\ A}}\skt{ गच्छति संभ्रमेण सुप्तः~॥}\pend

\pstart\versepar \ledsidenote{\passref{3.4.66}}\edtext{\skt{अव्य\-व\-स्थि\-त\-म\-ति\-श्च\-ल\-चित्तो}}{\lemma{\skt{चलचित्तो}}\Afootnote{H; \skt{चलदृष्टिर्}\ A}}\skt{ }\pend
\pstart\versepar \skt{}\edtext{\skt{मन्द\-द\-न्त\-ध\-न\-स\-ञ्च\-य\-मित्रः}}{\lemma{\skt{मन्द\appbrk{}द\appbrk{}न्त\appbrk{}ध\appbrk{}न\appbrk{}स\appbrk{}ञ्च\appbrk{}य\appbrk{}मित्रः}}\Afootnote{H; \skt{मन्द\appbrk{}र\appbrk{}त्न\appbrk{}ध\appbrk{}न\appbrk{}सं\appbrk{}च\appbrk{}य\appbrk{}मित्रः}\ A}}\skt{~॥ }\pend
\pstart\versepar \skt{किञ्चिदेव विल\-प\-त्य\-नि\-बद्धं }\pend
\pstart\versepar \skt{मारु\-त\-प्र\-कृ\-ति\-रेष मनुष्यः~॥}\pend

\pstart \ledsidenote{\passref{3.4.68}}\skt{पित्तप्रकृतिस्तु स्वेदनः }\edtext{\skt{शीतसहो}}{\lemma{\skt{शीतसहो}}\Afootnote{H; \emph{om.}\ A}}\skt{ }\edtext{\skt{दुर्ग्गन्धी}}{\lemma{\skt{दुर्ग्गन्धी}}\Afootnote{H; \skt{दुर्गन्धः}\ A}}\skt{ }\edtext{\skt{पीतः}}{\lemma{\skt{पीतः}}\Afootnote{H; \skt{पीत}\ A}}\skt{ }\edtext{\skt{कालो वा}}{\lemma{\skt{कालो वा}}\Afootnote{H; \emph{om.}\ A}}\skt{~। }\edtext{\edtext{\skt{शिथि\-ला\-ङ्ग\-स्ता\-म्र\-न\-ख\-न\-य\-न\-ता\-लु\-जि\-ह्वौ\-ष्ठ\-पा\-द\-पा\-णि\-तलो}}{\lemma{\skt{शिथिलाङ्गस्}}\Afootnote{H; \skt{शिथिलाङ्ग}\ A}}}{\lemma{\skt{ताम्र\appbrk{}न\appbrk{}ख\appbrk{}न\appbrk{}य\appbrk{}न\appbrk{}ता\appbrk{}लु\appbrk{}जि\appbrk{}ह्वौ\appbrk{}ष्ठ\appbrk{}पा\appbrk{}द\appbrk{}पा\appbrk{}णि\appbrk{}तलो}}\Afootnote{H; \skt{स्ताम्र\appbrk{}न\appbrk{}ख\appbrk{}न\appbrk{}य\appbrk{}न\appbrk{}ता\appbrk{}लु\appbrk{}जि\appbrk{}ह्वौ\appbrk{}ष्ठ\appbrk{}पा\appbrk{}णि\appbrk{}पा\appbrk{}द\appbrk{}तलो}\ A}}\skt{ }\edtext{\skt{दुर्भगो}}{\lemma{\skt{दुर्भगो}}\Afootnote{A; \skt{दुर्ब्भगो}\ H}}\skt{ }\edtext{\skt{वली\-प\-लि\-त\-सृष्टो}}{\lemma{\skt{वली\appbrk{}प\appbrk{}लि\appbrk{}त\appbrk{}सृष्टो}}\Afootnote{\skt{वली\appbrk{}प\appbrk{}लि\appbrk{}त\appbrk{}खा\appbrk{}लि\appbrk{}त्य\appbrk{}जुष्टो}\ A; \skt{वलीपलितो सृष्टो}\ H}}\skt{ बहु\-भु\-गु\-ष्ण\-द्वेषी }\edtext{\skt{क्षिप्र\-क्रो\-ध\-प्र\-सा\-द\-शीलो}}{\lemma{\skt{क्षिप्र\appbrk{}क्रो\appbrk{}ध\appbrk{}प्र\appbrk{}सा\appbrk{}द\appbrk{}शीलो}}\Afootnote{H; \skt{क्षिप्र\appbrk{}को\appbrk{}प\appbrk{}प्र\appbrk{}सादो}\ A}}\skt{ }\edtext{\skt{मध्यमायुश्च}}{\lemma{\skt{मध्यमायुश्च}}\Afootnote{H; \skt{मध्यबलो मध्यायुश्च}\ A}}\skt{ भवति~॥}\pend

\pstart\versepar \ledsidenote{\passref{3.4.69}}\edtext{\skt{भवति}}{\lemma{\skt{भवति}}\Afootnote{H; \emph{om.}\ A}}\skt{ }\edtext{\skt{चात्र}}{\lemma{\skt{चात्र}}\Afootnote{H; \emph{om.}\ A}}\skt{~॥ }\pend
\pstart\versepar \skt{मेधावी }\edtext{\skt{निपु\-ण\-म\-ति\-र्न्नि\-गृ\-ह्य\-वक्ता}}{\lemma{\skt{न्निगृह्यवक्ता}}\Afootnote{H; \skt{विगृह्य वक्ता}\ A}}\skt{ }\pend
\pstart\versepar \skt{तेजस्वी समितिषु दुर्न्निवारवीर्यः~। }\pend
\pstart\versepar \skt{सुप्तः }\edtext{\edtext{\skt{सन्क\-न\-क\-प\-ला\-श\-क\-र्ण्णि\-का\-रान्}}{\lemma{\skt{सन्}}\Afootnote{A; \skt{सं}\ H}}}{\lemma{\skt{कन\appbrk{}क\appbrk{}प\appbrk{}ला\appbrk{}श\appbrk{}क\appbrk{}र्ण्णि\appbrk{}का\appbrk{}रान्}}\Afootnote{A; \skt{कन\appbrk{}क\appbrk{}प\appbrk{}ला\appbrk{}श\appbrk{}क\appbrk{}र्ण्णि\appbrk{}कारां}\ H}}\skt{ }\pend
\pstart\versepar \skt{संपश्येदपि च }\edtext{\skt{हुता\-श\-वि\-द्यु\-दुल्काः}}{\lemma{\skt{हुता\appbrk{}श\appbrk{}वि\appbrk{}द्यु\appbrk{}दुल्काः}}\Afootnote{A; \skt{हुता\appbrk{}श\appbrk{}वि\appbrk{}द्यु\appbrk{}दुल्कां}\ H}}\skt{~॥}\pend

\pstart\versepar \ledsidenote{\passref{3.4.70}}\skt{न }\edtext{\edtext{\skt{भया\-त्प्र\-ण\-मे\-त्प्र\-ण\-तेषु}}{\lemma{\skt{प्रणमेत्}}\Afootnote{H; \skt{प्रणमे}\ A}}}{\lemma{\skt{प्रणतेषु}}\Afootnote{H; \skt{दनतेष्व्}\ A}}\skt{ }\edtext{\skt{मृदुः}}{\lemma{\skt{मृदुः}}\Afootnote{H; \skt{अमृदुः}\ A}}\skt{ }\pend
\pstart\versepar \skt{}\edtext{\skt{प्रणतेषु च}}{\lemma{\skt{प्रणतेषु च}}\Afootnote{H; \skt{प्रणतेष्वपि}\ A}}\skt{ सान्त्व\-न\-दा\-न\-रुचिः~। }\pend
\pstart\versepar \skt{भवतीह सदा व्यथितास्यगतिः }\pend
\pstart\versepar \skt{स भवेदिह पित्त\-कृ\-त\-प्र\-कृतिः~॥}\pend

\pstart \ledsidenote{\passref{3.4.72}}\skt{श्लेष्मप्रकृतिस्तु }\edtext{\skt{दूर्व्वे\-न्दी\-व\-र\-नि\-स्त्रिं\-शा\-र्द्द्रा\-रि\-ष्ट\-क\-श\-र\-का\-ण्डा\-ना\-म\-न्य\-त\-म\-वर्णः}}{\lemma{\skt{दूर्व्वे\appbrk{}न्दी\appbrk{}व\appbrk{}र\appbrk{}नि\appbrk{}स्त्रिं\appbrk{}शा\appbrk{}र्द्द्रा\appbrk{}रि\appbrk{}ष्ट\appbrk{}क\appbrk{}श\appbrk{}र\appbrk{}का\appbrk{}ण्डा\appbrk{}नाम्}}\Afootnote{A; \skt{दुर्व्वे\appbrk{}दी\appbrk{}स्व\appbrk{}र\appbrk{}नि\appbrk{}स्त्रिं\appbrk{}शा\appbrk{}र्द्द्रा\appbrk{}रि\appbrk{}ष्ट\appbrk{}क\appbrk{}श\appbrk{}र\appbrk{}का\appbrk{}ण्डा\appbrk{}नाम्}\ H}}\skt{ सुभगः प्रियदर्शनो मधुरप्रियः~। कृतज्ञो धृति\-मा\-न्स\-हि\-ष्णु\-र\-लो\-लुपो बल\-वां\-श्चि\-र\-ग्राही दृढवैरो }\edtext{\skt{दीर्घायुश्च}}{\lemma{\skt{दीर्घायुश्च}}\Afootnote{H; \skt{च}\ A}}\skt{ भवति~॥}\pend

\pstart\versepar \ledsidenote{\passref{3.4.73}}\edtext{\skt{स्निघ्दाङ्गः}}{\lemma{\skt{स्निघ्दाङ्गः}}\Afootnote{H; \skt{शुक्लाक्षः}\ A}}\skt{ }\edtext{\skt{स्थिर\-सु\-वि\-भ\-क्त\-चा\-रु\-गात्रो}}{\lemma{\skt{स्थिर\appbrk{}सु\appbrk{}वि\appbrk{}भ\appbrk{}क्त\appbrk{}चा\appbrk{}रु\appbrk{}गात्रो}}\Afootnote{H; \skt{स्थिर\appbrk{}कु\appbrk{}टि\appbrk{}ला\appbrk{}लि\appbrk{}नी\appbrk{}ल\appbrk{}केशो}\ A}}\skt{ }\pend
\pstart\versepar \skt{लक्ष्मी\-वा\-न्ज\-ल\-द\-मृ\-द\-ङ्ग\-सिं\-ह\-घोषः~। }\pend
\pstart\versepar \skt{सुप्तः सन्स\-क\-म\-ल\-हं\-स\-च\-क्र\-वा\-कान् }\pend
\pstart\versepar \skt{सम्पश्येदपि च }\edtext{\edtext{\skt{जला\-श्र\-या\-न्म\-नो\-ज्ञान्}}{\lemma{\skt{जलाश्रयान्}}\Afootnote{H; \skt{जलाशयान्}\ A}}}{\lemma{\skt{मनोज्ञान्}}\Afootnote{A; \skt{मनोज्ञां}\ H}}\skt{~॥}\footnoteA{Cf. Aṣṭāṅgasaṅgraha Śārīrasthāna 8.14cd: \emph{svapne sapadmān savihaṅgamālāṃs toyāśayān paśyati toyadāṃś ca ||} and Aṣṭāṅgahṛdaya Śārīrasthāna 3.102cd: \emph{svapne sapadmān savihaṅgamālāṃs toyāśayān paśyati toyadāṃś ca ||}.}\pend

\pstart\versepar \ledsidenote{\passref{3.4.75}}\skt{दृढशास्त्रमतिः }\edtext{\skt{स्थिरमित्रमनाः}}{\lemma{\skt{स्थिरमित्रमनाः}}\Afootnote{\skt{स्थिरमित्रधनः}\ A; \skt{स्थिरमित्रम मिः}\ H\textsuperscript{ac}; \skt{स्थिरमित्रम नः}\ H\textsuperscript{pc}}}\skt{ }\pend
\pstart\versepar \skt{परिगण्य }\edtext{\skt{चिराच्च}}{\lemma{\skt{च}}\Afootnote{H; \skt{प्र}\ A}}\skt{ ददाति }\edtext{\skt{बहु}}{\lemma{\skt{बहु}}\Afootnote{A; \skt{वहुः}\ H}}\skt{~। }\pend
\pstart\versepar \skt{}\edtext{\skt{परि\-नि\-ष्ठि\-त\-वा\-क्य\-पदः}}{\lemma{\skt{परि\appbrk{}नि\appbrk{}ष्ठि\appbrk{}त\appbrk{}वा\appbrk{}क्य\appbrk{}पदः}}\Afootnote{\skt{परि\appbrk{}नि\appbrk{}श्चि\appbrk{}त\appbrk{}वा\appbrk{}क्य\appbrk{}पदः}\ A; \skt{परि\appbrk{}नि\appbrk{}ष्ट\appbrk{}त\appbrk{}वा\appbrk{}क्य\appbrk{}पदः}\ H}}\skt{ सततं }\pend
\pstart\versepar \skt{गुरु\-मा\-न\-क\-रश्च }\edtext{\skt{कफप्रकृतिः}}{\lemma{\skt{कफप्रकृतिः}}\Afootnote{H; \skt{भवेत्स सदा}\ A}}\skt{~॥}\pend

\pstart\versepar \ledsidenote{\passref{3.4.77.add}}\skt{दृश्यते प्रकृतौ रूपं यदा दोषद्वयस्य तु~॥ }\pend
\pstart\versepar \skt{संसर्ग्गन्तं विजा\-नी\-या\-त्सं\-स\-र्ग्ग\-स्त्रि\-वि\-धश्च यः~॥}\pend

\pstart\versepar \ledsidenote{\passref{3.4.78}}\skt{प्रकोपो }\edtext{\skt{वान्यतोभावः}}{\lemma{\skt{वान्यतोभावः}}\Afootnote{\skt{वान्यथाभावो}\ A; \skt{वान्यतोभावो}\ H}}\skt{ क्षयो वा नोपजायते~।}\footnoteA{All witnesses support the reading \emph{bhāvo kṣayo}. This should be read \emph{bhāvo 'kṣayo}, which makes no sense.}\skt{ }\pend
\pstart\versepar \skt{प्रकृतीनां स्वभावेन }\edtext{\skt{ज्ञायन्ते तु}}{\lemma{\skt{ज्ञायन्ते तु}}\Afootnote{H; \skt{जायते तु}\ A}}\skt{ गतायुषः~॥}\pend

\pstart\versepar \ledsidenote{\passref{3.4.79}}\skt{विषजातो यथा }\edtext{\skt{कीटो}}{\lemma{\skt{कीटो}}\Afootnote{H; \skt{कीटो न}\ A}}\skt{ }\edtext{\skt{विषेण न}}{\lemma{\skt{विषेण न}}\Afootnote{H; \skt{विषेण}\ A}}\skt{ विपद्यते~। }\pend
\pstart\versepar \skt{}\edtext{\edtext{\edtext{\skt{तद्व\-त्प्र\-कृ\-ति\-भि\-र्द्दे\-ह\-स्त\-ज्जा\-त\-त्वान्न}}{\lemma{\skt{प्रकृतिभिर्}}\Afootnote{H; \skt{प्रकृतयो}\ A}}}{\lemma{\skt{द्देहस्}}\Afootnote{H; \skt{मर्त्यं}\ A}}}{\lemma{\skt{तज्जातत्वान्न}}\Afootnote{H; \skt{शक्नुवन्ति न}\ A}}\skt{ }\edtext{\skt{बाध्यते}}{\lemma{\skt{बाध्यते}}\Afootnote{\skt{बाधितुम्}\ A; \skt{वाधते}\ H}}\skt{~॥ }\footnoteA{We have emended the reading of H, \emph{bādhate}, to \emph{bādhyate} as the intransative use of the former is not clearly attested.}\pend

\pstart\versepar \ledsidenote{\passref{3.4.80}}\edtext{\edtext{\skt{प्रकृतिरिह}}{\lemma{\skt{प्रकृतिर्}}\Afootnote{H; \skt{प्रकृति}\ A}}}{\lemma{\skt{इह}}\Afootnote{H; \skt{मिह}\ A}}\skt{ नराणां }\edtext{\skt{भौतिकी}}{\lemma{\skt{भौतिकी}}\Afootnote{H; \skt{भौतिकीं}\ A}}\skt{ केचिदाहुः }\pend
\pstart\versepar \skt{पव\-न\-द\-ह\-न\-तोयैः }\edtext{\skt{कीर्त्तितास्तास्तु}}{\lemma{\skt{कीर्त्तितास्तास्तु}}\Afootnote{A; \skt{कीर्त्तितास्सस्तु}\ H}}\skt{ तिस्रः~। }\pend
\pstart\versepar \skt{स्थिर\-वि\-पु\-ल\-श\-रीरः }\edtext{\skt{पार्थि\-व\-श्चे\-त्क्ष\-मा\-वान्}}{\lemma{\skt{चेत्}}\Afootnote{H; \skt{च}\ A}}\skt{ }\pend
\pstart\versepar \skt{शुचिरथ चिरजीवी नाभसः खैर्म्महद्भिः~॥}\pend

\pstart\versepar \ledsidenote{\passref{3.4.81}}\edtext{\skt{शौच\-मा\-स्ति\-क्य\-म\-भ्यासो}}{\lemma{\skt{अभ्यासो}}\Afootnote{A; \skt{आभ्यासो}\ H}}\skt{ वेदेषु गुरुपूजनम्~। }\pend
\pstart\versepar \skt{प्रिया\-ति\-थि\-त्व\-मिज्या च ब्रह्मकायस्य लक्षणम्~॥}\pend

\pstart\versepar \ledsidenote{\passref{3.4.82}}\edtext{\skt{शौर्यमाज्ञा}}{\lemma{\skt{शौर्यम्}}\Afootnote{H; \skt{माहात्म्यं शौर्यम्}\ A}}\skt{ }\edtext{\skt{महाभाग्यं}}{\lemma{\skt{महाभाग्यं}}\Afootnote{H; \skt{च}\ A}}\skt{ सततं शास्त्रबुद्धिता~। }\pend
\pstart\versepar \skt{भृत्यानां }\edtext{\skt{भरणञ्चैव}}{\lemma{\skt{चैव}}\Afootnote{H; \skt{चापि}\ A}}\skt{ माहेन्द्रं कायलक्षणम्~॥}\pend

\pstart\versepar \ledsidenote{\passref{3.4.83}}\edtext{\skt{सहिष्णुतता}}{\lemma{\skt{सहिष्णुतता}}\Afootnote{H; \emph{om.}\ A}}\skt{ शीतसेवा }\edtext{\skt{पैङ्गल्यं}}{\lemma{\skt{पैङ्गल्यं}}\Afootnote{\skt{सहिष्णुत्वं पैङ्गल्यं}\ A; \skt{पैङ्गाख्यं}\ H}}\skt{ हरिकेशता~। }\pend
\pstart\versepar \skt{}\edtext{\edtext{\skt{प्रियोदकत्वञ्च}}{\lemma{\skt{प्रियोदकत्वञ्}}\Afootnote{H; \skt{प्रियवादित्वम्}\ A}}}{\lemma{\skt{च}}\Afootnote{H; \skt{इत्ये}\ A}}\skt{ }\edtext{\skt{तथा}}{\lemma{\skt{तथा}}\Afootnote{H; \skt{तद्}\ A}}\skt{ वारु\-ण\-ङ्का\-य\-ल\-क्ष\-णम्~॥}\pend

\pstart\versepar \ledsidenote{\passref{3.4.84}}\skt{मध्यस्थता सहि\-ष्णु\-त्व\-म\-र्थ\-स्या\-ग\-म\-स\-ञ्चयौ~। }\pend
\pstart\versepar \skt{}\edtext{\edtext{\skt{महा\-प्र\-स\-व\-श\-क्तिश्च}}{\lemma{\skt{महा\appbrk{}प्र\appbrk{}स\appbrk{}व\appbrk{}श\appbrk{}क्तिश्}}\Afootnote{H; \skt{महा\appbrk{}प्र\appbrk{}स\appbrk{}व\appbrk{}शक्ति}\ A}}}{\lemma{\skt{च}}\Afootnote{H; \skt{त्वं}\ A}}\skt{ कौबेरं कायलक्षणम्~॥}\pend

\pstart\versepar \ledsidenote{\passref{3.4.85}}\skt{गन्ध\-मा\-ल्य\-प्रि\-य\-त्वञ्च }\edtext{\skt{नृत्य\-वा\-दि\-त्र\-का\-मता}}{\lemma{\skt{नृत्य\appbrk{}वा\appbrk{}दि\appbrk{}त्र\appbrk{}का\appbrk{}मता}}\Afootnote{H; \skt{नृत्य\appbrk{}वा\appbrk{}दि\appbrk{}त्र\appbrk{}का\appbrk{}मिता}\ A}}\skt{~। }\pend
\pstart\versepar \skt{विहारशीलता चैव गान्ध\-र्व्व\-ङ्का\-य\-ल\-क्ष\-णम्~॥}\pend

\pstart\versepar \ledsidenote{\passref{3.4.86}}\skt{प्राप्तकारी दृढोत्थानो }\edtext{\skt{धृति\-मा\-न्स्मृ\-ति\-मा\-ञ्छुचिः}}{\lemma{\skt{धृतिमान्}}\Afootnote{H; \skt{निर्भयः}\ A}}\skt{~। }\pend
\pstart\versepar \skt{}\edtext{\edtext{\skt{राग\-द्वे\-ष\-भ\-या\-ज्ञा\-नै\-र्व्व\-र्ज्जितो}}{\lemma{\skt{राग\appbrk{}द्वे\appbrk{}ष\appbrk{}भ\appbrk{}या\appbrk{}ज्ञा\appbrk{}नैर्}}\Afootnote{H; \skt{राग\appbrk{}मो\appbrk{}ह\appbrk{}म\appbrk{}द\appbrk{}द्वे\appbrk{}षैर्}\ A}}}{\lemma{\skt{व्वर्ज्जितो}}\Afootnote{A; \skt{व्विर्ज्जितो}\ H}}\skt{ याम्यसत्त्ववान्~॥}\pend

\pstart\versepar \ledsidenote{\passref{3.4.87}}\edtext{\skt{ब्रह्म\-च\-र्य\-व्र\-त\-ज\-प\-हो\-मा\-ध्य\-य\-न\-से\-वि\-नम्}}{\lemma{\skt{ब्रह्म\appbrk{}च\appbrk{}र्य\appbrk{}व्र\appbrk{}त\appbrk{}ज\appbrk{}प\appbrk{}हो\appbrk{}मा\appbrk{}ध्य\appbrk{}य\appbrk{}न\appbrk{}से\appbrk{}वि\appbrk{}नम्}}\Afootnote{\skt{जप\appbrk{}व्र\appbrk{}त\appbrk{}ब्र\appbrk{}ह्म\appbrk{}च\appbrk{}र्य\appbrk{}हो\appbrk{}मा\appbrk{}ध्य\appbrk{}य\appbrk{}न\appbrk{}से\appbrk{}वि\appbrk{}नम्}\ A; \skt{व्रह्म\appbrk{}च\appbrk{}र्य\appbrk{}व्र\appbrk{}त\appbrk{}जपं होमा\appbrk{}ध्य\appbrk{}य\appbrk{}न\appbrk{}से\appbrk{}विनं}\ H}}\skt{~। }\pend
\pstart\versepar \skt{ज्ञान\-वि\-ज्ञा\-न\-स\-म्प\-न्न\-मृ\-षि\-सत्त्वं }\edtext{\edtext{\skt{विदुर्नरम्}}{\lemma{\skt{विदुर्}}\Afootnote{H; \emph{om.}\ A}}}{\lemma{\skt{नरम्}}\Afootnote{H; \skt{नरं विदुः}\ A}}\skt{~॥}\pend

\pstart\versepar \ledsidenote{\passref{3.4.88}}\skt{सप्तैते सात्त्विकाः काया }\edtext{\skt{राजसांस्तु}}{\lemma{\skt{राजसांस्तु}}\Afootnote{A; \skt{रजसास्तु}\ H}}\skt{ निबोध मे~। }\pend
\pstart\versepar \skt{ऐश्वर्यवन्तं रौद्रञ्च शूरं चण्ड\-म\-सू\-य\-कम्~।}\footnoteA{The following typology is paralleled in the Carakasaṃhitā, Śārīrasthāna 5, v. 38 (Ācārya 1941: 323--324).}\pend

\pstart \ledsidenote{\passref{3.4.89ab}}\skt{एकाशिनं }\edtext{\skt{त्वौप\-धि\-क\-मा\-सुरं}}{\lemma{\skt{त्वौपधिकम्}}\Afootnote{H; \skt{चौदरिकम्}\ A}}\skt{ सत्त्वमीदृशं~॥}\pend

\pstart \ledsidenote{\passref{3.4.91cd}}\skt{एकान्तग्राहिता }\edtext{\edtext{\skt{रौद्र्य\-म\-सू\-या\-ध\-र्म्म\-वा\-क्यता}}{\lemma{\skt{रौद्र्यम्}}\Afootnote{H; \skt{रौद्रम्}\ A}}}{\lemma{\skt{असू\appbrk{}या\appbrk{}ध\appbrk{}र्म्म\appbrk{}वा\appbrk{}क्यता}}\Afootnote{H; \skt{असूया धर्मबाह्यता}\ A}}\skt{~।}\pend

\pstart\versepar \ledsidenote{\passref{3.4.92}}\edtext{\edtext{\skt{भृश\-मा\-त्म\-स्त\-व\-श्चापि}}{\lemma{\skt{भृशम्}}\Afootnote{H; \skt{भृश}\ A}}}{\lemma{\skt{आत्मस्तवश्}}\Afootnote{H; \skt{मात्रं तमश्}\ A}}\skt{ राक्षसं कायलक्षणं~॥ }\pend
\pstart\versepar \skt{}\edtext{\skt{उत्सृष्टाचारता}}{\lemma{\skt{उत्सृष्टाचारता}}\Afootnote{\skt{उच्छिष्टाहारता}\ A; \skt{उत्सृष्ठाचारता}\ H}}\skt{ }\edtext{\skt{रौक्ष्ण्यं}}{\lemma{\skt{रौक्ष्ण्यं}}\Afootnote{H; \skt{तैक्ष्ण्यं}\ A}}\skt{ साहसप्रियता तथा~॥}\pend

\pstart \ledsidenote{\passref{3.4.93ab}}\edtext{\skt{स्त्रीलोलुपत्वं}}{\lemma{\skt{स्त्रीलोलुपत्वं}}\Afootnote{A; \skt{स्त्रीलौलुपत्वं}\ H}}\skt{ नैर्ल्लज्ज्यं }\edtext{\skt{पैशाचो}}{\lemma{\skt{पैशाचो}}\Afootnote{H; \skt{पैशाचं}\ A}}\skt{ }\edtext{\skt{गुणसङ्ग्रहः}}{\lemma{\skt{गुणसङ्ग्रहः}}\Afootnote{H; \skt{कायलक्षणम्}\ A}}\skt{~॥}\pend

\pstart \ledsidenote{\passref{3.4.89cd}}\edtext{\skt{तीक्ष्ण\-मा\-या\-स\-ब\-हुलं}}{\lemma{\skt{आयासबहुलं}}\Afootnote{\skt{आया}\ A; \skt{आयासवकुलं}\ H}}\skt{ }\edtext{\skt{क्रोधनं}}{\lemma{\skt{क्रोधनं}}\Afootnote{H; \emph{om.}\ A}}\skt{ }\edtext{\skt{सत्त्वभीरुकम्}}{\lemma{\skt{सत्त्वभीरुकम्}}\Afootnote{H; \skt{सिनं भीरुं चण्डं मायान्वितं तथा}\ A}}\skt{~।}\footnoteA{Witness H reads \emph{āyāsabakulam} which we take as a mistranscription of \emph{āyāsabahulam}. Given the context, one wonders whether there is a historical connection or confusion with Hindī \emph{ऐयाश} "debauched, lewd person."}\pend

\pstart \ledsidenote{\passref{3.4.90ab}}\edtext{\skt{विहा\-रा\-हा\-र\-च\-पलं}}{\lemma{\skt{विहा\appbrk{}रा\appbrk{}हा\appbrk{}र\appbrk{}च\appbrk{}पलं}}\Afootnote{H; \skt{विहा\appbrk{}रा\appbrk{}चा\appbrk{}र\appbrk{}च\appbrk{}पलं}\ A}}\skt{ सर्प्प\-स\-त्त्व\-म्वि\-दु\-र्न्नरं~॥}\pend

\pstart \ledsidenote{\passref{3.4.93cd}}\edtext{\edtext{\skt{अस\-म्वि\-भा\-ग\-म\-ल\-स\-न्दुः\-ख\-शी\-ल\-म\-सू\-य\-क\-म्दुः\-ख\-शी\-ल\-म\-सू\-यकं}}{\lemma{\skt{दुःखशीलम्}}\Afootnote{A NE-RORI\_6838; \skt{दुःशीलम्}\ H}}}{\lemma{\skt{असूयकम्}}\Afootnote{A NE-RORI\_6838; \skt{अनुसूचकम्}\ H}}\skt{~।}\footnoteA{We read against witness H and with the vulgate and the mūla as recorded in MS Jaipur RORI 6838 of Gayadāsa's Nyāyacandrikā. This makes better sense and avoids the problematic lexeme \emph{anusūcakam} in spite of lectio dificilior.}\pend

\pstart \ledsidenote{\passref{3.4.94ab}}\skt{लोलु\-प\-ञ्चा\-प्य\-दा\-तारं प्रेत\-स\-त्त्व\-म्वि\-दु\-र्न्न\-रम्~॥}\pend

\pstart \ledsidenote{\passref{3.4.90cd}}\edtext{\skt{प्रबु\-द्ध\-का\-म\-सेवी}}{\lemma{\skt{प्रबु\appbrk{}द्ध\appbrk{}का\appbrk{}म\appbrk{}सेवी}}\Afootnote{H; \skt{प्रवृ\appbrk{}द्ध\appbrk{}का\appbrk{}म\appbrk{}सेवी}\ A}}\skt{ }\edtext{\edtext{\skt{चाजस्रमाहार}}{\lemma{\skt{चाजस्रम्}}\Afootnote{\skt{चाप्यजस्रा}\ A; \skt{वाजस्राम्}\ H\textsuperscript{ac}; \skt{वाजस्रम्}\ H\textsuperscript{pc}}}}{\lemma{\skt{आहार}}\Afootnote{H; \skt{हार}\ A}}\skt{ एव च~।}\pend

\pstart \ledsidenote{\passref{3.4.91ab}}\edtext{\skt{अम\-र्ष\-णोऽ\-न\-व\-स्थानः}}{\lemma{\skt{ऽनवस्थानः}}\Afootnote{\skt{ऽनवस्थायी}\ A; \skt{नवस्थानः}\ H}}\skt{ शाकुनं }\edtext{\edtext{\skt{सत्त्वमुच्यते}}{\lemma{\skt{सत्त्वम्}}\Afootnote{H; \skt{का}\ A}}}{\lemma{\skt{उच्यते}}\Afootnote{H; \skt{यलक्षणम्}\ A}}\skt{~॥}\pend

\pstart \ledsidenote{\passref{3.4.94cd}}\skt{षडेते राजसाः }\edtext{\skt{कायास्तामसास्तु}}{\lemma{\skt{तामसास्तु}}\Afootnote{H; \skt{तामसांस्तु}\ A}}\skt{ निबोध मे~॥}\pend

\pstart\versepar \ledsidenote{\passref{3.4.95}}\skt{दुर्मेधस्त्वं मन्दता च }\edtext{\skt{स्वप्न\-मै\-थु\-न\-नि\-त्यता}}{\lemma{\skt{स्वप्न\appbrk{}मै\appbrk{}थु\appbrk{}न\appbrk{}नि\appbrk{}त्यता}}\Afootnote{H; \skt{स्वप्ने मैथुननित्यता}\ A}}\skt{~। }\pend
\pstart\versepar \skt{}\edtext{\skt{निरा\-क\-रि\-ष्णु\-ता\-शूरा}}{\lemma{\skt{निरा\appbrk{}क\appbrk{}रि\appbrk{}ष्णु\appbrk{}ता\appbrk{}शूरा}}\Afootnote{\skt{निराकरिष्णुता चैव}\ A; \skt{निरा\appbrk{}क\appbrk{}ष्णु\appbrk{}ता\appbrk{}शूरा}\ H}}\skt{ विज्ञेयाः पाशवा गुणाः~॥}\pend

\pstart\versepar \ledsidenote{\passref{3.4.96}}\skt{अनवस्थितता मौर्ख्यं भीरुत्वं }\edtext{\skt{तोयकामता}}{\lemma{\skt{तोयकामता}}\Afootnote{H; \skt{सलिलार्थिता}\ A}}\skt{~॥}\footnoteA{The reading \emph{toyakāmatā} is close to the \emph{toyakāmam} of Carakasaṃhitā 4.5.39(2) (p.324) while the vulgate's \emph{salilārthitā} is departs from both.}\skt{ }\pend
\pstart\versepar \skt{}\edtext{\skt{पर\-स्प\-रा\-भि\-म\-र्षश्च}}{\lemma{\skt{पर\appbrk{}स्प\appbrk{}रा\appbrk{}भि\appbrk{}म\appbrk{}र्षश्च}}\Afootnote{H; \skt{पर\appbrk{}स्प\appbrk{}रा\appbrk{}भि\appbrk{}म\appbrk{}र्दश्च}\ A}}\skt{ }\edtext{\skt{मत्स्यसत्त्वात्मका}}{\lemma{\skt{मत्स्यसत्त्वात्मका}}\Afootnote{\skt{मत्स्यसत्त्व}\ A; \skt{मत्स्यसर्व्वात्मका}\ H}}\skt{ }\edtext{\skt{गुणाः}}{\lemma{\skt{गुणाः}}\Afootnote{H; \skt{स्य लक्षणम्}\ A}}\skt{~॥}\footnoteA{We read with the vulgate against witness H's \emph{sarvvātmaka} for sense and because of the parallel sentence structure in the surrounding context. Also, \emph{𑐬𑑂𑐰} and \emph{ 𑐟𑑂𑐰 } are similar in early Nepalese orthography and may have been confused by the scribe of H or his exemplar.}\pend

\pstart\versepar \ledsidenote{\passref{3.4.97}}\edtext{\skt{सर्व\-बु\-द्ध्य\-ङ्ग\-हीनः स्यादाहारे}}{\lemma{\skt{सर्व\appbrk{}बु\appbrk{}द्ध्य\appbrk{}ङ्ग\appbrk{}हीनः स्यादाहारे}}\Afootnote{\skt{एक\appbrk{}स्था\appbrk{}न\appbrk{}र\appbrk{}ति\appbrk{}र्नि\appbrk{}त्य\appbrk{}मा\appbrk{}हारे}\ A; \skt{सत्व\appbrk{}वु\appbrk{}द्धा\appbrk{}ङ्ग\appbrk{}हीनः स्यादाहारे}\ H}}\skt{ केवले }\edtext{\skt{रतः}}{\lemma{\skt{रतः}}\Afootnote{H; \emph{om.}\ A}}\skt{~।}\footnoteA{The reading \emph{sarvabuddhyaṅga} of Carakasaṃhitā 4.5.39(2) (p.324) is preferred to the reading of H, \emph{sattvabudhāṅga}.}\skt{ }\pend
\pstart\versepar \skt{}\edtext{\skt{प्रिया\-र्थ\-का\-म\-र\-हितः}}{\lemma{\skt{प्रिया\appbrk{}र्थ\appbrk{}का\appbrk{}म\appbrk{}र\appbrk{}हितः}}\Afootnote{H; \skt{रतः}\ A}}\skt{ }\edtext{\skt{सत्त्वे}}{\lemma{\skt{सत्त्वे}}\Afootnote{H; \emph{om.}\ A}}\skt{ }\edtext{\skt{वानस्पते}}{\lemma{\skt{वानस्पते}}\Afootnote{H; \skt{वानस्पत्}\ A}}\skt{ }\edtext{\skt{स्थितः}}{\lemma{\skt{स्थितः}}\Afootnote{H; \skt{यो नरः सत्त्व\appbrk{}ध\appbrk{}र्म\appbrk{}का\appbrk{}मा\appbrk{}र्थ\appbrk{}व\appbrk{}र्जितः}\ A}}\skt{~॥}\pend

\pstart\versepar \ledsidenote{\passref{3.4.98}}\skt{इत्येते त्रिविधाः कायाः प्रोक्ता वै }\edtext{\skt{राजसादयः}}{\lemma{\skt{राजसादयः}}\Afootnote{H; \skt{तामसास्त}\ A}}\skt{~। }\pend
\pstart\versepar \skt{}\edtext{\skt{विज्ञाय}}{\lemma{\skt{विज्ञाय}}\Afootnote{H; \skt{था का}\ A}}\skt{ }\edtext{\edtext{\skt{काय\-प्र\-कृ\-ति\-म\-नु\-रूपां}}{\lemma{\skt{कायप्रकृतिम्}}\Afootnote{H; \skt{यानां प्रकृतीर्}\ A}}}{\lemma{\skt{अनुरूपां}}\Afootnote{H; \skt{ज्ञात्वा त्वनुरूपां}\ A}}\skt{ क्रियां }\edtext{\skt{चरेद्~॥ इति}}{\lemma{\skt{चरेद्~॥ इति}}\Afootnote{H; \skt{चरेत्}\ A}}\skt{~॥}\footnoteA{The vulgate verse \emph{mahāprakṛtayas}... that is not in the Nepalese version at this point was known to Gayadāsa and may even have been originated by him and then been inserted into the vulgate by Ḍalhaṇa.}\pend

\pstart \ledsidenote{\passref{3.4.trailer}}\edtext{\skt{इति शारीरं}}{\lemma{\skt{इति शारीरं}}\Afootnote{H; \skt{इति सुश्रु\appbrk{}त\appbrk{}सं\appbrk{}हि\appbrk{}तायां शारीरस्थाने गर्भव्याकरणं शारीरं}\ A}}\skt{ }\edtext{\skt{चतुर्थो}}{\lemma{\skt{चतुर्थो}}\Afootnote{H; \skt{नाम चतुर्थो}\ A}}\skt{ }\edtext{\skt{ध्यायः}}{\lemma{\skt{ध्यायः}}\Afootnote{H; \skt{ऽध्यायः}\ A}}\skt{~॥}\pend

\endnumbering
\beginnumbering
\eledsection*{Adhyāya 5}
\pstart \ledsidenote{\passref{3.5.1}}\skt{अथातः }\edtext{\skt{शरी\-र\-व्या\-क\-रणं}}{\lemma{\skt{शरी\appbrk{}र\appbrk{}व्या\appbrk{}क\appbrk{}रणं}}\Afootnote{H; \skt{शरी\appbrk{}र\appbrk{}सं\appbrk{}ख्या\appbrk{}व्या\appbrk{}क\appbrk{}रणं}\ A}}\skt{ }\edtext{\skt{शारी\-र\-म्व्या\-ख्या\-स्यामः}}{\lemma{\skt{शारीरम्}}\Afootnote{A; \skt{शारीतम्}\ H}}\skt{~॥}\pend

\pstart \ledsidenote{\passref{3.5.3}}\edtext{\edtext{\skt{शुक्र\-शो\-णि\-त\-ङ्ग\-र्भा\-श\-य\-स्थ\-मा\-त्म\-प्र\-कृ\-ति\-वि\-का\-र\-स\-म्मू\-र्च्छितं}}{\lemma{\skt{गर्भाशयस्थम्}}\Afootnote{A H\textsuperscript{pc}; \skt{गर्भासस्थम्}\ H\textsuperscript{ac}}}}{\lemma{\skt{आत्म\appbrk{}प्र\appbrk{}कृ\appbrk{}ति\appbrk{}वि\appbrk{}का\appbrk{}र\appbrk{}स\appbrk{}म्मू\appbrk{}र्च्छितं}}\Afootnote{A; \skt{आत्म\appbrk{}प्र\appbrk{}कृ\appbrk{}ति\appbrk{}वि\appbrk{}का\appbrk{}र\appbrk{}स\appbrk{}म्मू\appbrk{}र्च्छितो}\ H}}\skt{ गर्भ }\edtext{\skt{इत्युच्यते~। }\edtext{\edtext{\skt{तम\-चे\-त\-ना\-व\-स्थि\-त\-म्वा\-यु\-र्व्वि\-भ\-जते}}{\lemma{\skt{अचे\appbrk{}त\appbrk{}ना\appbrk{}व\appbrk{}स्थि\appbrk{}तम्}}\Afootnote{\skt{चेतनावस्थितं}\ A; \skt{दचे\appbrk{}त\appbrk{}ना\appbrk{}व\appbrk{}स्थि\appbrk{}तम्}\ H}}}{\lemma{\skt{व्विभजते}}\Afootnote{H; \skt{विभजति}\ A}}}{\lemma{\skt{उच्यते~। तम्}}\Afootnote{A; \skt{उच्यते~॥ त}\ H}}\skt{~। तेज एनम्पचति~। आपः }\edtext{\skt{क्लेदयन्ति}}{\lemma{\skt{क्लेदयन्ति}}\Afootnote{A; \skt{क्लेदयति}\ H}}\skt{~। पृथिवी संहन्ति~। आका\-श\-म्वि\-व\-र्द्ध\-यति~। एवम्विवर्द्धितः स यदा }\edtext{\edtext{\skt{हस्त\-पा\-द\-जि\-ह्वा\-घ्रा\-ण\-व\-द\-न\-क\-र्ण्णा\-दि\-भि\-रु\-पे\-त\-स्तदा}}{\lemma{\skt{हस्त\appbrk{}पा\appbrk{}द\appbrk{}जि\appbrk{}ह्वा\appbrk{}घ्रा\appbrk{}ण\appbrk{}व\appbrk{}द\appbrk{}न\appbrk{}क\appbrk{}र्ण्णा\appbrk{}दि\appbrk{}भिर्}}\Afootnote{\skt{हस्त\appbrk{}पा\appbrk{}द\appbrk{}जि\appbrk{}ह्वा\appbrk{}घ्रा\appbrk{}ण\appbrk{}क\appbrk{}र्ण\appbrk{}नि\appbrk{}त\appbrk{}म्बा\appbrk{}दि\appbrk{}भिर्}\ A; \skt{हस्त\appbrk{}पा\appbrk{}द\appbrk{}जि\appbrk{}ह्वा\appbrk{}घ्रा\appbrk{}ण\appbrk{}व\appbrk{}द\appbrk{}न\appbrk{}क\appbrk{}र्ण्णादि}\ H}}}{\lemma{\skt{उपेतस्}}\Afootnote{\skt{अङ्गैरुपेतस्}\ A; \skt{भिरुपेतस्}\ H}}\skt{ शरीरमिति संज्ञां लभते~। तच्च षडङ्गं }\edtext{\skt{शाखाश्चत्वारः}}{\lemma{\skt{चत्वारः}}\Afootnote{H; \skt{चतस्रो}\ A}}\skt{~। मध्यं }\edtext{\skt{पञ्चमं}}{\lemma{\skt{पञ्चमं}}\Afootnote{A; \emph{om.}\ H}}\skt{ षष्ठं शिर इति~॥}\pend

\pstart\versepar \ledsidenote{\passref{3.5.4}}\edtext{\skt{अतः}}{\lemma{\skt{अतः}}\Afootnote{A H\textsuperscript{pc}; \skt{अत्रः}\ H\textsuperscript{ac}}}\skt{ परम्प्रत्यङ्गानि वक्ष्यन्ते~। }\edtext{\skt{मस्त\-को\-द\-र\-पृ\-ष्ठ\-ना\-भि\-ल\-ला\-ट\-ना\-सा\-चि\-वु\-क\-ग्रीवा}}{\lemma{\skt{मस्त\appbrk{}को\appbrk{}द\appbrk{}र\appbrk{}पृ\appbrk{}ष्ठ\appbrk{}ना\appbrk{}भि\appbrk{}ल\appbrk{}ला\appbrk{}ट\appbrk{}ना\appbrk{}सा\appbrk{}चि\appbrk{}वु\appbrk{}क\appbrk{}ग्रीवा}}\Afootnote{H; \skt{मस्त\appbrk{}को\appbrk{}द\appbrk{}र\appbrk{}पृ\appbrk{}ष्ठ\appbrk{}ना\appbrk{}भि\appbrk{}ल\appbrk{}ला\appbrk{}ट\appbrk{}ना\appbrk{}सा\appbrk{}चि\appbrk{}बु\appbrk{}क\appbrk{}ब\appbrk{}स्ति\appbrk{}ग्रीवा}\ A}}\skt{ }\edtext{\skt{इत्येते}}{\lemma{\skt{एते}}\Afootnote{H; \skt{एता}\ A}}\skt{ एकैकाः~॥ }\edtext{\skt{कर्ण्ण\-ने\-त्र\-भ्रू\-श\-ङ्खां\-स\-ग\-ण्ड\-क\-क्ष\-स्त\-न\-वं\-क्ष\-ण\-वृ\-ष\-ण\-पा\-र्श्व\-स्फि\-क्जा\-नु\-कू\-र्प्प\-र\-वा\-हु\-प्र\-भृ\-तयो}}{\lemma{\skt{कर्ण्ण\appbrk{}ने\appbrk{}त्र\appbrk{}भ्रू\appbrk{}श\appbrk{}ङ्खां\appbrk{}स\appbrk{}ग\appbrk{}ण्ड\appbrk{}क\appbrk{}क्ष\appbrk{}स्त\appbrk{}न\appbrk{}वं\appbrk{}क्ष\appbrk{}ण\appbrk{}वृ\appbrk{}ष\appbrk{}ण\appbrk{}पा\appbrk{}र्श्व\appbrk{}स्फि\appbrk{}क्जा\appbrk{}नु\appbrk{}कू\appbrk{}र्प्प\appbrk{}र\appbrk{}वा\appbrk{}हु\appbrk{}प्र\appbrk{}भृ\appbrk{}तयो}}\Afootnote{\skt{कर्ण\appbrk{}ने\appbrk{}त्र\appbrk{}भ्रू\appbrk{}श\appbrk{}ङ्खां\appbrk{}स\appbrk{}ग\appbrk{}ण्ड\appbrk{}क\appbrk{}क्ष\appbrk{}स्त\appbrk{}न\appbrk{}व\appbrk{}ङ्क्ष\appbrk{}ण\appbrk{}वृ\appbrk{}ष\appbrk{}ण\appbrk{}पा\appbrk{}र्श्व\appbrk{}स्फि\appbrk{}ग्जा\appbrk{}नु\appbrk{}कू\appbrk{}र्प\appbrk{}र\appbrk{}बा\appbrk{}हू\appbrk{}रु\appbrk{}प्र\appbrk{}भृ\appbrk{}तयो}\ A; \skt{कर्ण्ण\appbrk{}ने\appbrk{}त्र\appbrk{}भ्रू\appbrk{}वां\appbrk{}च\appbrk{}स\appbrk{}ग\appbrk{}ण्ड\appbrk{}क\appbrk{}क्ष\appbrk{}स्त\appbrk{}न\appbrk{}वं\appbrk{}क्ष\appbrk{}ण\appbrk{}वृ\appbrk{}ष\appbrk{}ण\appbrk{}पा\appbrk{}र्श्व\appbrk{}स्फि\appbrk{}क्जा\appbrk{}नु\appbrk{}कु\appbrk{}र्प्प\appbrk{}र\appbrk{}वा\appbrk{}हु\appbrk{}प्र\appbrk{}भृ\appbrk{}तयो}\ H}}\skt{ द्वे द्वे~। }\edtext{\skt{विंशतिरङ्गुल्यः}}{\lemma{\skt{अङ्गुल्यः}}\Afootnote{H; \skt{अङ्गुलयः}\ A}}\skt{ स्रोतान्सि }\edtext{\skt{वक्ष्यन्ते}}{\lemma{\skt{वक्ष्यन्ते}}\Afootnote{H; \skt{वक्ष्यमा}\ A}}\skt{~॥}\footnoteA{We emend witness H's \emph{vāṃcasa} to the vulgate's reading, in the absence of a good idea about what H's reading could be.}\skt{ }\pend
\pstart\versepar \skt{}\edtext{\edtext{\edtext{\edtext{\skt{इत्ये\-त\-त्प्र\-त्य\-ङ्ग\-वि\-भा\-ग\-मुक्तं}}{\lemma{\skt{इत्य्}}\Afootnote{H; \skt{णानि}\ A}}}{\lemma{\skt{एतत्}}\Afootnote{H; \skt{एष}\ A}}}{\lemma{\skt{प्रत्य\appbrk{}ङ्ग\appbrk{}वि\appbrk{}भा\appbrk{}गम्}}\Afootnote{H; \skt{प्रत्यङ्गविभाग}\ A}}}{\lemma{\skt{उक्तं}}\Afootnote{H; \skt{उक्तः}\ A}}\skt{~।}\footnoteA{One would expect masculine \emph{vibhāgaḥ} rather than the neuter.}\pend

\pstart \ledsidenote{\passref{3.5.7}}\edtext{\edtext{\skt{त्वचोऽभिहिताः}}{\lemma{\skt{त्वचो}}\Afootnote{H; \skt{विस्तारोऽत ऊर्ध्वं --- त्वचो}\ A}}}{\lemma{\skt{ऽभिहिताः}}\Afootnote{A; \skt{भिहिताः}\ H}}\skt{ कला धातवो }\edtext{\skt{मलधराश्च}}{\lemma{\skt{मलधराश्च}}\Afootnote{H; \skt{मला दोषा यकृत्प्लीहानौ फुप्फुस उण्डुको हृदयं वृक्कौ च}\ A}}\skt{~॥}\pend

\pstart \ledsidenote{\passref{3.5.8}}\skt{आशयास्तु }\edtext{\skt{वात\-पि\-त्त\-श्ले\-ष्मा\-म\-प\-क्वा\-ग्नि\-मू\-त्राणां}}{\lemma{\skt{वात\appbrk{}पि\appbrk{}त्त\appbrk{}श्ले\appbrk{}ष्मा\appbrk{}म\appbrk{}प\appbrk{}क्वा\appbrk{}ग्नि\appbrk{}मू\appbrk{}त्राणां}}\Afootnote{H; \skt{वाताशयः, पित्ताशयः, श्लेष्माशयो, रक्ताशय, आमाशयः, पक्वाशयो, मूत्रा}\ A}}\skt{ }\edtext{\skt{सप्त}}{\lemma{\skt{सप्त}}\Afootnote{H; \skt{शयः}\ A}}\skt{~। स्त्रीणां गर्भाशयोऽष्टम इति~।}\pend

\pstart \ledsidenote{\passref{3.5.9}}\edtext{\skt{सार्द्ध\-त्रि\-व्या\-या\-मा\-न्य\-न्त्राणि}}{\lemma{\skt{सार्द्ध\appbrk{}त्रि\appbrk{}व्या\appbrk{}या\appbrk{}मा\appbrk{}न्य\appbrk{}न्त्राणि}}\Afootnote{\skt{सार्ध\appbrk{}त्रि\appbrk{}व्या\appbrk{}मा\appbrk{}न्य\appbrk{}न्त्राणि}\ A; \skt{सार्द्ध\appbrk{}त्रि\appbrk{}व्या\appbrk{}या\appbrk{}मा\appbrk{}न्य\appbrk{}न्त्राणि}\ H}}\skt{ }\edtext{\skt{पुंषां}}{\lemma{\skt{पुंषां}}\Afootnote{H; \skt{पुंसां}\ A}}\skt{ }\edtext{\edtext{\skt{स्त्रीणा\-म\-र्द्ध\-व्या\-या\-म\-ही\-नानि}}{\lemma{\skt{स्त्रीणाम्}}\Afootnote{H; \skt{स्त्रीणा}\ A}}}{\lemma{\skt{अर्द्ध\appbrk{}व्या\appbrk{}या\appbrk{}म\appbrk{}ही\appbrk{}नानि}}\Afootnote{H; \skt{मर्ध\appbrk{}व्या\appbrk{}म\appbrk{}ही\appbrk{}नानि}\ A}}\skt{~।}\footnoteA{Witness H's \emph{vyāyāma} is accepted, since dictionaries commonly include the sense {\symfont `}{\symfont `}fathom.'' The vulgate's \emph{vyāma} is historically and etymologically the more correct term.}\pend

\pstart \ledsidenote{\passref{3.5.10}}\skt{श्रव\-ण\-न\-य\-न\-व\-द\-न\-घ्रा\-ण\-गु\-द\-मे\-ढ्राणि नव स्रोतान्सि नरा\-णा\-म्ब\-हि\-र्म्मु\-खानि~। }\edtext{\skt{तान्येव}}{\lemma{\skt{तान्य्}}\Afootnote{H; \skt{एतान्य्}\ A}}\skt{ स्त्रीणामपराणि च त्रीणि~। द्वे }\edtext{\edtext{\skt{स्तन\-यो\-र\-प\-र\-म\-ध\-स्ता\-द्र\-क्त\-वाहि}}{\lemma{\skt{अपरम्}}\Afootnote{H; \emph{om.}\ A}}}{\lemma{\skt{रक्तवाहि}}\Afootnote{H; \skt{रक्तवहं}\ A}}\skt{ च~।}\pend

\pstart \ledsidenote{\passref{3.5.11}}\skt{षोडश कण्ड\-रा\-स्ता\-सा\-ञ्च\-तस्रः }\edtext{\skt{पादयोस्तावन्त}}{\lemma{\skt{तावन्त}}\Afootnote{H; \skt{तावत्}\ A}}\skt{ }\edtext{\skt{एव}}{\lemma{\skt{एव}}\Afootnote{H; \skt{यो}\ A}}\skt{ हस्त\-ग्री\-वा\-पृ\-ष्ठेषु~। तत्र हस्त\-पा\-द\-ग\-ता\-ना\-ङ्क\-ण्ड\-राणां }\edtext{\skt{नखा}}{\lemma{\skt{नखा}}\Afootnote{A; \skt{नव}\ H}}\skt{ }\edtext{\skt{अग्रप्ररोहा}}{\lemma{\skt{अग्रप्ररोहा}}\Afootnote{A; \skt{नव प्ररोहा}\ H}}\footnoteA{We emend witness H's \emph{nava nava} to the vulgate's \emph{nakhā agra-} for sense.}\skt{ }\edtext{\edtext{\skt{ग्रीवा\-हृ\-द\-य\-नि\-ब\-न्ध\-नी\-ना\-म\-धो\-मु\-खानां}}{\lemma{\skt{निबन्धनीनाम्}}\Afootnote{H; \skt{निबन्धिनीनाम्}\ A}}}{\lemma{\skt{अधोमुखानां}}\Afootnote{H; \skt{अधो\appbrk{}भा\appbrk{}ग\appbrk{}ग\appbrk{}तानां}\ A}}\skt{ }\edtext{\skt{मेढ्रं}}{\lemma{\skt{मेढ्रं}}\Afootnote{A; \skt{मेढ्र}\ H}}\footnoteA{We emend witness H's \emph{meḍhra} to the vulgate's \emph{meḍhraṃ} for sense.}\skt{ }\edtext{\edtext{\edtext{\edtext{\edtext{\edtext{\edtext{\skt{श्रोणी\-पृ\-ष्ठ\-नि\-ब\-न्ध\-नी\-ना\-म\-धो\-ग\-ता\-ना\-म्नि\-त\-म्ब\-मू\-रु\-व\-क्षोऽ\-क्ष\-पि\-ण्ड्या\-दी\-नाञ्च}}{\lemma{\skt{श्रोणीपृष्ठ}}\Afootnote{H; \skt{श्रोणिपृष्ठ}\ A}}}{\lemma{\skt{निबन्धनीनाम्}}\Afootnote{H; \skt{निबन्धिनीनाम्}\ A}}}{\lemma{\skt{अधोगतानाम्}}\Afootnote{H; \skt{अधो\appbrk{}भा\appbrk{}ग\appbrk{}ग\appbrk{}तानां}\ A}}}{\lemma{\skt{नितम्बम्}}\Afootnote{H; \skt{बिम्बं}\ A}}}{\lemma{\skt{ऊरु}}\Afootnote{H; \skt{मूर्धोरु}\ A}}}{\lemma{\skt{वक्षोऽक्ष}}\Afootnote{\skt{वक्षोऽंस}\ A; \skt{वन्धोत्क}\ H}}}{\lemma{\skt{पिण्ड्यादीनाञ्}}\Afootnote{H; \skt{पिण्डादीनां}\ A}}\footnoteA{Y. T. Ācārya wrote several footnotes pointing to variant readings with \emph{-akṣapiṇḍa-} in the manuscripts he used for the text and Ḍalhaṇa's commentary (1938: 364--365); in each case he removed that reading. We emend back to the reading for \emph{-vakṣākṣa-} following Ācārya's manuscript readings as well as Gayadāsa's reading \emph{ūruvakṣokṣapiṃḍa-} {\symfont `}{\symfont `}thighs, back and eyeballs" (MS Jaipur RORI 6838, f. 322v).}\skt{~॥}\pend

\pstart \ledsidenote{\passref{3.5.12}}\edtext{\skt{मांस\-सि\-रा\-स्ना\-य्व\-स्थि\-ध\-म\-न्नी\-जा\-लानि}}{\lemma{\skt{धमन्नी}}\Afootnote{H; \emph{om.}\ A}}\skt{ प्रत्येकं }\edtext{\skt{चत्वारि}}{\lemma{\skt{चत्वारि}}\Afootnote{A; \skt{चत्वारि चत्वारि}\ H}}\skt{ तानि }\edtext{\skt{मणि\-ब\-न्ध\-न\-गु\-ल्फ\-सं\-श्रि\-तानि}}{\lemma{\skt{मणि\appbrk{}ब\appbrk{}न्ध\appbrk{}न\appbrk{}गु\appbrk{}ल्फ\appbrk{}सं\appbrk{}श्रि\appbrk{}तानि}}\Afootnote{H; \skt{मणि\appbrk{}ब\appbrk{}न्ध\appbrk{}गु\appbrk{}ल्फ\appbrk{}सं\appbrk{}श्रि\appbrk{}तानि पर\appbrk{}स्प\appbrk{}र\appbrk{}नि\appbrk{}ब\appbrk{}द्धानि पर\appbrk{}स्प\appbrk{}र\appbrk{}ग\appbrk{}वा\appbrk{}क्षि\appbrk{}तानि चेति}\ A}}\skt{~। यैर्ग\-वा\-क्षि\-त\-मिदं शरीरम्~।}\pend

\pstart \ledsidenote{\passref{3.5.13}}\skt{षट्कूर्च्चास्ते }\edtext{\skt{हस्त\-पा\-द\-ग्री\-वा\-मे\-ढ्रेषु}}{\lemma{\skt{हस्त\appbrk{}पा\appbrk{}द\appbrk{}ग्री\appbrk{}वा\appbrk{}मे\appbrk{}ढ्रेषु}}\Afootnote{H; \skt{हस्त\appbrk{}पा\appbrk{}द\appbrk{}ग्री\appbrk{}वा\appbrk{}मे\appbrk{}ढ्रेषु; हस्तयोर्द्वौ, पादयोर्द्वौ, ग्रीवा\appbrk{}मे\appbrk{}ढ्र\appbrk{}यो\appbrk{}रे\appbrk{}कैकः}\ A}}\skt{~।}\pend

\pstart \ledsidenote{\passref{3.5.14}}\skt{महत्यो }\edtext{\skt{मांस\-र\-ज्ज\-व\-श्च\-तस्रः}}{\lemma{\skt{मांसरज्जवश्}}\Afootnote{A; \skt{मान्सरज्जवल्स्}\ H}}\skt{ }\edtext{\edtext{\edtext{\skt{पेशी\-नि\-ब\-न्ध\-ना\-र्थ\-म्पृ\-ष्ठ\-वं\-श\-मु\-भ\-यतो}}{\lemma{\skt{पेशी\appbrk{}नि\appbrk{}ब\appbrk{}न्ध\appbrk{}ना\appbrk{}र्थम्}}\Afootnote{H; \skt{पृष्ठ\appbrk{}वं\appbrk{}श\appbrk{}मु\appbrk{}भ\appbrk{}यतः पेशी\appbrk{}नि\appbrk{}ब\appbrk{}न्ध\appbrk{}नार्थं}\ A}}}{\lemma{\skt{पृष्ठवंशम्}}\Afootnote{H; \emph{om.}\ A}}}{\lemma{\skt{उभयतो}}\Afootnote{H; \emph{om.}\ A}}\skt{ द्वे बाह्ये }\edtext{\skt{अभ्यन्तरे}}{\lemma{\skt{अभ्यन्तरे}}\Afootnote{H; \skt{आभ्यन्तरे}\ A}}\skt{ च द्वे~।}\pend

\pstart \ledsidenote{\passref{3.5.15}}\skt{सप्त }\edtext{\edtext{\skt{सीमन्यस्तत्र}}{\lemma{\skt{सीमन्यस्}}\Afootnote{\skt{सेवन्यः}\ A; \skt{सीमन्य}\ H}}}{\lemma{\skt{तत्र}}\Afootnote{\skt{सिरसि}\ A; \skt{स्तत्र}\ H}}\skt{ विभक्ताः }\edtext{\skt{शिरसि}}{\lemma{\skt{शिरसि}}\Afootnote{H; \emph{om.}\ A}}\skt{ पञ्च }\edtext{\skt{जिह्वा\-शे\-प\-सो\-रे\-कै\-कास्ताः}}{\lemma{\skt{जिह्वा\appbrk{}शे\appbrk{}प\appbrk{}सो\appbrk{}रे\appbrk{}कै\appbrk{}कास्ताः}}\Afootnote{H; \skt{जिह्वा\appbrk{}शे\appbrk{}फ\appbrk{}सो\appbrk{}रे\appbrk{}कैका; ताः}\ A}}\skt{ परिहर्तव्याः शस्त्रेण~।}\pend

\pstart \ledsidenote{\passref{3.5.16}}\edtext{\edtext{\skt{चतु\-र्द्द\-शा\-स्सं\-ह\-ता\-स्तेषां}}{\lemma{\skt{चतुर्द्दशास्}}\Afootnote{H; \skt{चतुर्दशा}\ A}}}{\lemma{\skt{संहतास्}}\Afootnote{H; \skt{स्थ्नां संघाताः}\ A}}\skt{ त्रयो गुल्फ\-जा\-नु\-वं\-क्ष\-णेषु~। एते\-ने\-त\-र\-सक्थि }\edtext{\skt{वाहु च}}{\lemma{\skt{वाहु च}}\Afootnote{H; \skt{बाहू च}\ A}}\skt{ व्याख्यातौ~। }\edtext{\skt{त्रिक\-शि\-र\-सो\-रे\-कैकः}}{\lemma{\skt{त्रिक\appbrk{}शि\appbrk{}र\appbrk{}सो\appbrk{}रे\appbrk{}कैकः}}\Afootnote{H; \skt{त्रिक\appbrk{}शि\appbrk{}र\appbrk{}सो\appbrk{}रे\appbrk{}कैकः}\ A}}\skt{~।}\pend

\pstart \ledsidenote{\passref{3.5.17}}\skt{चतुर्द्दशैव }\edtext{\skt{सीमन्ताः~। }\edtext{\skt{यैर्म्मुक्तः}}{\lemma{\skt{म्मुक्तः}}\Afootnote{H; \skt{ह्युक्ताः}\ A}}}{\lemma{\skt{सीमन्ताः~। यैर्}}\Afootnote{H; \skt{सीमन्ताः; ते चास्थि\appbrk{}स\appbrk{}ङ्घा\appbrk{}त\appbrk{}व\appbrk{}द्ग\appbrk{}ण\appbrk{}नीयाः, यतस्तैर्युक्ता अस्थिसंघाताः; ये}\ A}}\skt{ }\edtext{\skt{संहातास्ते}}{\lemma{\skt{संहातास्ते}}\Afootnote{H; \skt{संघातास्ते}\ A}}\skt{ }\edtext{\skt{खल्व\-क\-द\-शै\-केषां}}{\lemma{\skt{खल्व\appbrk{}क\appbrk{}द\appbrk{}शै\appbrk{}केषां}}\Afootnote{H; \skt{खल्व\appbrk{}ष्टा\appbrk{}द\appbrk{}शै\appbrk{}के\appbrk{}षाम्}\ A}}\skt{~।}\pend

\pstart \ledsidenote{\passref{3.5.18}}\edtext{\skt{त्रीणि}}{\lemma{\skt{त्रीणि}}\Afootnote{A; \skt{त्रिणि}\ H}}\skt{ }\edtext{\skt{षष्ट्या}}{\lemma{\skt{षष्ट्या}}\Afootnote{H; \skt{सषष्}\ A}}\skt{ }\edtext{\edtext{\edtext{\skt{शतान्यस्थामिति}}{\lemma{\skt{शतान्य्}}\Afootnote{H; \skt{टान्य्}\ A}}}{\lemma{\skt{अस्थाम्}}\Afootnote{H; \skt{अस्थि}\ A}}}{\lemma{\skt{इति}}\Afootnote{H; \skt{शतानि}\ A}}\skt{ }\edtext{\skt{वेद\-वा\-दि\-नोऽ\-प\-भा\-षन्ते}}{\lemma{\skt{ऽपभाषन्ते}}\Afootnote{H; \skt{भाषन्ते}\ A}}\skt{~। }\edtext{\skt{शल्यतन्त्रेषु}}{\lemma{\skt{शल्यतन्त्रेषु}}\Afootnote{H; \skt{शल्यतन्त्रेषु तु}\ A}}\skt{ }\edtext{\skt{त्रीण्येव}}{\lemma{\skt{एव}}\Afootnote{A; \skt{अव}\ H}}\skt{ शतानि~। तेषां }\edtext{\skt{सविं\-श\-त्य\-स्थि\-शतं}}{\lemma{\skt{सविं\appbrk{}श\appbrk{}त्य\appbrk{}स्थि\appbrk{}शतं}}\Afootnote{\skt{सविं\appbrk{}श\appbrk{}म\appbrk{}स्थि\appbrk{}शतं}\ A; \skt{सविंशत्य्स्थिशतं}\ H}}\skt{ शाखासु~। सप्तदशोत्तरं शतं }\edtext{\skt{श्रोणी\-पा\-र्श्व\-पृ\-ष्ठो\-रस्सु}}{\lemma{\skt{श्रोणी\appbrk{}पा\appbrk{}र्श्व\appbrk{}पृ\appbrk{}ष्ठो\appbrk{}रस्सु}}\Afootnote{H; \skt{श्रोणि\appbrk{}पा\appbrk{}र्श्व\appbrk{}पृ\appbrk{}ष्ठो\appbrk{}रःसु}\ A}}\skt{ }\edtext{\skt{ग्रीवायां}}{\lemma{\skt{ग्रीवायां}}\Afootnote{H; \skt{ग्रीवां}\ A}}\skt{ }\edtext{\edtext{\skt{प्रत्यू\-र्द्ध्व\-न्त्रि\-ष\-ष्ठिर्}}{\lemma{\skt{ऊर्द्ध्वन्}}\Afootnote{A; \skt{ऊर्द्द्वन्}\ H}}}{\lemma{\skt{त्रिषष्ठिर्}}\Afootnote{H; \skt{त्रिषष्टिः, एवमस्थ्नां त्रीणि शतानि पूर्यन्ते}\ A}}\skt{~।}\footnoteA{See Hoernle's edition and discussion of this passage (Hoernle 1907: 218--219).}\pend

\pstart \ledsidenote{\passref{3.5.19}}\edtext{\skt{एकैकस्यां}}{\lemma{\skt{एकैकस्यां}}\Afootnote{H; \skt{एकैकस्यां तु}\ A}}\skt{ पादाङ्गुल्यां त्रीणि }\edtext{\skt{त्रिणि}}{\lemma{\skt{त्रिणि}}\Afootnote{H; \skt{त्रीणि}\ A}}\skt{ तानि पञ्चदश तल\-कू\-र्च्च\-गु\-ल्फ\-सं\-श्रि\-तानि~। दश पार्ष्ण्यामेकं जङ्घायां द्वे }\edtext{\skt{जानु\-न्ये\-क\-मे\-क\-मू\-रूभ्यां}}{\lemma{\skt{ऊरूभ्यां}}\Afootnote{H; \skt{ऊराविति}\ A}}\skt{~। }\edtext{\edtext{\skt{त्रिंश\-दे\-व\-मे\-क\-स्मि\-न्स\-क्थिनि}}{\lemma{\skt{त्रिंशदेवम्}}\Afootnote{H; \skt{त्रिंशदेवम्}\ A}}}{\lemma{\skt{सक्थिनि}}\Afootnote{H; \skt{सक्थ्नि}\ A}}\skt{ भवन्ति~। }\edtext{\skt{एतेन}}{\lemma{\skt{एतेन}}\Afootnote{H; \skt{एतेने}\ A}}\skt{ }\edtext{\skt{इतरसक्थि}}{\lemma{\skt{इतरसक्थि}}\Afootnote{H; \skt{तरसक्थि}\ A}}\skt{ }\edtext{\skt{वाहु}}{\lemma{\skt{वाहु}}\Afootnote{H; \skt{बाहू च}\ A}}\skt{ व्याख्यातौ~॥ श्रोण्याम्पञ्च~। तेषां गुद\-भ\-ग\-नि\-त\-म्वेषु चत्वारि त्रिक\-सं\-श्रि\-त\-मेकं पार्श्वे }\edtext{\skt{षट्त्रिं\-श\-दे\-क\-स्मि\-न्ता\-व\-न्त्येव}}{\lemma{\skt{तावन्त्येव}}\Afootnote{H; \emph{om.}\ A}}\skt{ }\edtext{\skt{द्वितीये}}{\lemma{\skt{द्वितीये}}\Afootnote{H; \skt{द्वितीयेऽप्येवं}\ A}}\skt{~। पृष्ठे }\edtext{\skt{द्वात्रिं\-श\-दे\-वा\-ष्ठा\-वु\-रसि}}{\lemma{\skt{द्वात्रिं\appbrk{}श\appbrk{}दे\appbrk{}वा\appbrk{}ष्ठा\appbrk{}वु\appbrk{}रसि}}\Afootnote{H; \skt{त्रिंशत्, अष्टावुरसि}\ A}}\skt{~। द्वे }\edtext{\skt{अक्षकसंज्ञे}}{\lemma{\skt{अक्षकसंज्ञे}}\Afootnote{H; \skt{अंसफलके}\ A}}\skt{ ग्रीवायां नव }\edtext{\skt{कण्टनाड्यां}}{\lemma{\skt{कण्टनाड्यां}}\Afootnote{H; \skt{कण्ठनाड्यां}\ A}}\skt{ }\edtext{\skt{चत्वाति}}{\lemma{\skt{चत्वाति}}\Afootnote{H; \skt{चत्वारि}\ A}}\skt{~। द्वे }\edtext{\skt{हनुसंज्ञे}}{\lemma{\skt{हनुसंज्ञे}}\Afootnote{H; \skt{हन्वोः}\ A}}\skt{ दन्ता }\edtext{\edtext{\skt{त्रिंश\-न्ना\-सा\-या\-न्त्री\-ण्ये\-क\-ता\-लुनि}}{\lemma{\skt{त्रिंशन्}}\Afootnote{H; \skt{द्वात्रिंशत्}\ A}}}{\lemma{\skt{एकतालुनि}}\Afootnote{H; \skt{एकं तालुनि}\ A}}\skt{ गण्ड\-क\-र्ण्ण\-शं\-खे\-ष्वे\-कैकं }\edtext{\skt{षट्छिरसीति}}{\lemma{\skt{छिरसीति}}\Afootnote{H; \skt{शिरसीति}\ A}}\skt{~।}\pend

\pstart \ledsidenote{\passref{3.5.20}}\edtext{\skt{तानेतानस्थीनि}}{\lemma{\skt{तानेतानस्थीनि}}\Afootnote{H; \skt{एतानि}\ A}}\skt{ पंचविधानि भवन्ति~॥ तद्यथा~॥ }\edtext{\skt{कपा\-ल\-रू\-च\-क\-त\-रु\-ण\-व\-ल\-य\-न\-ल\-क\-सं\-ज्ञानि}}{\lemma{\skt{कपा\appbrk{}ल\appbrk{}रू\appbrk{}च\appbrk{}क\appbrk{}त\appbrk{}रु\appbrk{}ण\appbrk{}व\appbrk{}ल\appbrk{}य\appbrk{}न\appbrk{}ल\appbrk{}क\appbrk{}सं\appbrk{}ज्ञानि}}\Afootnote{H; \skt{कपा\appbrk{}ल\appbrk{}रु\appbrk{}च\appbrk{}क\appbrk{}त\appbrk{}रु\appbrk{}ण\appbrk{}व\appbrk{}ल\appbrk{}य\appbrk{}न\appbrk{}ल\appbrk{}क\appbrk{}सं\appbrk{}ज्ञानि}\ A}}\skt{ }\edtext{\skt{भवन्ति}}{\lemma{\skt{भवन्ति}}\Afootnote{H; \emph{om.}\ A}}\skt{~। तेषां }\edtext{\skt{जानु\-र्प्प\-र\-नि\-तम्वाः}}{\lemma{\skt{जानु\appbrk{}र्प्प\appbrk{}र\appbrk{}नि\appbrk{}तम्वाः}}\Afootnote{H; \skt{जानुनितम्बां}\ A}}\skt{ }\edtext{\skt{सग\-ण्डु\-ता\-लु\-म\-ध्य\-शिरः सु}}{\lemma{\skt{सग\appbrk{}ण्डु\appbrk{}ता\appbrk{}लु\appbrk{}म\appbrk{}ध्य\appbrk{}शिरः सु}}\Afootnote{H; \skt{सग\appbrk{}ण्ड\appbrk{}ता\appbrk{}लु\appbrk{}श\appbrk{}ङ्ख\appbrk{}शि\appbrk{}रःसु}\ A}}\skt{ कपालानि~। दशनास्तु }\edtext{\skt{रुचकाः}}{\lemma{\skt{रुचकाः}}\Afootnote{H; \skt{रुचकानि}\ A}}\skt{~। }\edtext{\skt{घ्राण\-क\-ण्ठ\-क\-र्ण्णा\-क्षि\-कू\-टेषु}}{\lemma{\skt{घ्राण\appbrk{}क\appbrk{}ण्ठ\appbrk{}क\appbrk{}र्ण्णा\appbrk{}क्षि\appbrk{}कू\appbrk{}टेषु}}\Afootnote{H; \skt{घ्राण\appbrk{}क\appbrk{}र्ण\appbrk{}ग्री\appbrk{}वा\appbrk{}क्षि\appbrk{}को\appbrk{}षेषु}\ A}}\skt{ तरुणानि~। }\edtext{\skt{पाद\-ह\-स्त\-ग्री\-वा\-पृ\-ष्ठेषु}}{\lemma{\skt{पाद\appbrk{}ह\appbrk{}स्त\appbrk{}ग्री\appbrk{}वा\appbrk{}पृ\appbrk{}ष्ठेषु}}\Afootnote{H; \skt{पार्श्वपृष्ठोरःसु}\ A}}\skt{ }\edtext{\skt{वलयानि}}{\lemma{\skt{वलयानि}}\Afootnote{\skt{वलयानि , शेषाणि नलकसंज्ञानि}\ A; \skt{वलयानि~। भवति चात्र}\ H}}\skt{~।}\pend

\pstart\versepar \ledsidenote{\passref{3.5.21}}\edtext{\skt{भवति}}{\lemma{\skt{भवति}}\Afootnote{\skt{भवन्ति}\ A; \emph{om.}\ H}}\skt{ }\edtext{\skt{चात्र}}{\lemma{\skt{चात्र}}\Afootnote{A; \emph{om.}\ H}}\skt{~॥ }\pend
\pstart\versepar \skt{}\edtext{\skt{यथा हि}}{\lemma{\skt{यथा हि}}\Afootnote{H; \emph{om.}\ A}}\skt{ }\edtext{\skt{सारो}}{\lemma{\skt{सारो}}\Afootnote{H; \emph{om.}\ A}}\skt{ }\edtext{\skt{वृक्षस्य}}{\lemma{\skt{वृक्षस्य}}\Afootnote{H; \emph{om.}\ A}}\skt{ }\edtext{\edtext{\skt{तिष्ठ\-त्य\-भ्य\-न्त\-रा\-श्रितः}}{\lemma{\skt{तिष्ठत्य्}}\Afootnote{H; \emph{om.}\ A}}}{\lemma{\skt{अभ्यन्तराश्रितः}}\Afootnote{H; \skt{अभ्यन्तरगतैः सारैर्यथा तिष्ठन्ति}\ A}}\skt{~। }\pend
\pstart\versepar \skt{}\edtext{\edtext{\skt{एवमस्थीनि}}{\lemma{\skt{एवम्}}\Afootnote{H; \skt{भू}\ A}}}{\lemma{\skt{अस्थीनि}}\Afootnote{H; \skt{रुहाः~। अस्थिसा}\ A}}\skt{ }\edtext{\skt{जन्तूनां}}{\lemma{\skt{जन्तूनां}}\Afootnote{H; \skt{रैस्तथा}\ A}}\skt{ }\edtext{\edtext{\skt{शरीरन्ध्रियते}}{\lemma{\skt{शरीरन्}}\Afootnote{H; \skt{देहा}\ A}}\skt{ तु तैः}}{\lemma{\skt{ध्रियते तु तैः}}\Afootnote{H; \skt{ध्रियन्ते देहिनां ध्रुवम्}\ A}}\skt{~॥}\pend

\pstart\versepar \ledsidenote{\passref{3.5.22}}\edtext{\skt{तस्मा\-च्चि\-र\-वि\-न\-ष्टेऽपि}}{\lemma{\skt{ऽपि}}\Afootnote{\skt{षु}\ A; \skt{पि}\ H}}\skt{ }\edtext{\skt{त्वग्मांसे तु}}{\lemma{\skt{त्वग्मांसे तु}}\Afootnote{H; \skt{त्वङ्मांसेषु}\ A}}\skt{ शरीरिणां~। }\pend
\pstart\versepar \skt{अस्थीनि न }\edtext{\skt{विश्यन्ति}}{\lemma{\skt{विश्यन्ति}}\Afootnote{H; \skt{विनश्यन्ति}\ A}}\skt{ }\edtext{\skt{सारोष}}{\lemma{\skt{सारोष}}\Afootnote{H; \skt{सारा}\ A}}\skt{ }\edtext{\skt{विधीयते}}{\lemma{\skt{विधीयते}}\Afootnote{H; \skt{ण्येतानि देहिनाम्}\ A}}\skt{~॥}\pend

\pstart\versepar \ledsidenote{\passref{3.5.23}}\skt{मांसान्यत्र निबद्धानि सिराभिः स्नायुभिस्तथा~। }\pend
\pstart\versepar \skt{अस्थीन्यालम्बनं कृत्वा न }\edtext{\skt{शीर्यन्ति}}{\lemma{\skt{शीर्यन्ति}}\Afootnote{H; \skt{शीर्यन्ते}\ A}}\skt{ पतन्ति वा~॥}\pend

\pstart \ledsidenote{\passref{3.5.24}}\edtext{\skt{सन्धयो}}{\lemma{\skt{सन्धयो}}\Afootnote{H; \skt{सन्धयस्तु}\ A}}\skt{ द्विवि\-धा\-श्चे\-ष्टा\-वन्तः स्थिराश्च~।}\pend

\pstart \ledsidenote{\passref{3.5.26}}\edtext{\skt{संख्यातास्तु}}{\lemma{\skt{संख्यातास्तु}}\Afootnote{H; \skt{सङ्ख्यातस्तु}\ A}}\skt{ }\edtext{\skt{द्वे}}{\lemma{\skt{द्वे}}\Afootnote{H; \emph{om.}\ A}}\skt{ }\edtext{\skt{शते}}{\lemma{\skt{शते}}\Afootnote{H; \emph{om.}\ A}}\skt{ दशोत्तरे }\edtext{\skt{तेषां}}{\lemma{\skt{तेषां}}\Afootnote{H; \skt{द्वे शते~। तेषां}\ A}}\skt{ शाखा }\edtext{\skt{अष्ट\-ष\-ष्टि\-रे\-को\-न\-षष्टिः}}{\lemma{\skt{अष्टषष्टिर्}}\Afootnote{H; \skt{स्वष्टषष्टिः}\ A}}\skt{ कोष्ठे }\edtext{\skt{ग्रीवा}}{\lemma{\skt{ग्रीवा}}\Afootnote{H; \skt{ग्रीवां}\ A}}\skt{ }\edtext{\edtext{\edtext{\skt{यांत्युर्द्धन्त्र्याशीतिः}}{\lemma{\skt{यांत्य्}}\Afootnote{H; \skt{प्र}\ A}}}{\lemma{\skt{उर्द्धन्}}\Afootnote{H; \skt{त्यूर्ध्वं}\ A}}}{\lemma{\skt{त्र्याशीतिः}}\Afootnote{H; \skt{त्र्यशीतिः}\ A}}\skt{~। }\edtext{\skt{एकैकस्यान्तु}}{\lemma{\skt{एकैकस्यान्तु}}\Afootnote{H; \skt{एकैकस्यां}\ A}}\skt{ पादाङ्गुल्यां त्रयस्त्रयः~। द्वावङ्गुष्ठे ते }\edtext{\skt{चतु\-र्द्द\-श\-गु\-ल्फ\-जा\-नु\-वं\-क्ष्य\-णे\-ष्ठे\-कैकः}}{\lemma{\skt{चतु\appbrk{}र्द्द\appbrk{}श\appbrk{}गु\appbrk{}ल्फ\appbrk{}जा\appbrk{}नु\appbrk{}वं\appbrk{}क्ष्य\appbrk{}णे\appbrk{}ष्ठे\appbrk{}कैकः}}\Afootnote{H; \skt{चतुर्दश; जानु\appbrk{}गु\appbrk{}ल्फ\appbrk{}व\appbrk{}ङ्क्ष\appbrk{}णे\appbrk{}ष्वे\appbrk{}कैकः}\ A}}\skt{ }\edtext{\edtext{\skt{सप्त\-द\-शै\-क\-स्मि\-न्स\-क्थिनि}}{\lemma{\skt{सप्त\appbrk{}द\appbrk{}शै\appbrk{}क\appbrk{}स्मिन्}}\Afootnote{H; \skt{एवं सप्त\appbrk{}द\appbrk{}शै\appbrk{}क\appbrk{}स्मिन्}\ A}}}{\lemma{\skt{सक्थिनि}}\Afootnote{H; \skt{सक्थ्नि}\ A}}\skt{ भवन्ति~। }\edtext{\skt{एतेन}}{\lemma{\skt{एतेन}}\Afootnote{H; \skt{एतेने}\ A}}\skt{ }\edtext{\skt{इतरसक्थि}}{\lemma{\skt{इतरसक्थि}}\Afootnote{H; \skt{तरसक्थि}\ A}}\skt{ }\edtext{\skt{वाहु च}}{\lemma{\skt{वाहु च}}\Afootnote{H; \skt{बाहू च}\ A}}\skt{ व्याख्यातौ~॥ त्रयः कटीकपालेषु चतुर्व्विंशतिः पृष्ठवंशे~। }\edtext{\skt{तावन्त्य}}{\lemma{\skt{तावन्त्य}}\Afootnote{H; \skt{तावन्त}\ A}}\skt{ एव पार्श्वयोः~। }\edtext{\skt{अष्टावूरसि}}{\lemma{\skt{अष्टावूरसि}}\Afootnote{H; \skt{उरस्यष्टौ}\ A}}\skt{~। }\edtext{\skt{तावन्त्य}}{\lemma{\skt{तावन्त्य}}\Afootnote{H; \skt{तावन्त}\ A}}\skt{ एव }\edtext{\skt{ग्रीवायान्तु}}{\lemma{\skt{ग्रीवायान्तु}}\Afootnote{H; \skt{ग्रीवायां, त्}\ A}}\skt{ }\edtext{\skt{यः}}{\lemma{\skt{यः}}\Afootnote{H; \skt{रयः}\ A}}\skt{ }\edtext{\skt{कण्ठनाण्ड्यां}}{\lemma{\skt{कण्ठनाण्ड्यां}}\Afootnote{H; \skt{कण्ठे, नाडीषु}\ A}}\skt{ }\edtext{\skt{हृदयक्लोमास}}{\lemma{\skt{हृदयक्लोमास}}\Afootnote{H; \skt{हृदयक्लोमनि}\ A}}\skt{ }\edtext{\edtext{\skt{नाडी\-ष्व\-ष्टा\-द\-श\-ता\-स्वे\-व\-नि\-बद्धा}}{\lemma{\skt{नाडीष्व्}}\Afootnote{H; \skt{बद्धास्}\ A}}}{\lemma{\skt{अष्टा\appbrk{}द\appbrk{}श\appbrk{}ता\appbrk{}स्वे\appbrk{}व\appbrk{}नि\appbrk{}बद्धा}}\Afootnote{H; \skt{वष्टादश}\ A}}\skt{~। दन्तपरिमाणा }\edtext{\skt{दन्त\-मू\-ले\-ष्वे\-कैकः}}{\lemma{\skt{एकैकः}}\Afootnote{H; \skt{एकः}\ A}}\skt{~। काकलके }\edtext{\skt{नासायाञ्चैक्केकः}}{\lemma{\skt{चैक्केकः}}\Afootnote{H; \skt{च}\ A}}\skt{~। }\edtext{\skt{द्वौ}}{\lemma{\skt{द्वौ}}\Afootnote{H; \emph{om.}\ A}}\skt{ }\edtext{\skt{द्वौवर्त्म्ममण्डल }\edtext{\skt{योर्न्नेत्राश्रितौ}}{\lemma{\skt{न्नेत्राश्रितौ}}\Afootnote{H; \skt{नेत्राश्रयौ}\ A}}}{\lemma{\skt{द्वौवर्त्म्ममण्डल योर्}}\Afootnote{H; \skt{द्वौ वर्त्म\appbrk{}म\appbrk{}ण्ड\appbrk{}ल\appbrk{}योर्}\ A}}\skt{~। गण्ड\-क\-र्ण्ण\-शं\-खे\-ष्वे\-कैकः~। द्वौ हनुसंधी }\edtext{\edtext{\skt{द्वावु\-प\-रि\-ष्टा\-द्भुवौ}}{\lemma{\skt{द्वावुपरिष्टाद्}}\Afootnote{H; \skt{द्वावुपरिष्टाद्}\ A}}}{\lemma{\skt{भुवौ}}\Afootnote{H; \skt{भ्रुवोः}\ A}}\skt{ शंखयोश्च पञ्च शिरः कपालेष्वेको }\edtext{\skt{मूर्द्ध्रि}}{\lemma{\skt{मूर्द्ध्रि}}\Afootnote{H; \skt{मूर्ध्नि}\ A}}\skt{~।}\pend

\pstart \ledsidenote{\passref{3.5.27}}\skt{त एते }\edtext{\skt{संधयोष्टविधा}}{\lemma{\skt{संधयोष्टविधा}}\Afootnote{H; \skt{सन्धयोऽष्टविधाः}\ A}}\skt{ }\edtext{\skt{भवन्ति}}{\lemma{\skt{भवन्ति}}\Afootnote{H; \emph{om.}\ A}}\skt{~। }\edtext{\edtext{\skt{तद्यथा}}{\lemma{\skt{तद्}}\Afootnote{H; \emph{om.}\ A}}}{\lemma{\skt{यथा}}\Afootnote{H; \emph{om.}\ A}}\skt{~। }\edtext{\skt{कोरो\-दू\-ख\-ल\-सा\-मु\-द्ग\-प्र\-त\-र\-तू\-न\-सी\-वनी}}{\lemma{\skt{कोरो\appbrk{}दू\appbrk{}ख\appbrk{}ल\appbrk{}सा\appbrk{}मु\appbrk{}द्ग\appbrk{}प्र\appbrk{}त\appbrk{}र\appbrk{}तू\appbrk{}न\appbrk{}सी\appbrk{}वनी}}\Afootnote{H; \skt{कोरो\appbrk{}लू\appbrk{}ख\appbrk{}ल\appbrk{}सा\appbrk{}मु\appbrk{}द्ग\appbrk{}प्र\appbrk{}त\appbrk{}र\appbrk{}तु\appbrk{}न्न\appbrk{}से\appbrk{}वनी}\ A}}\skt{ वाय\-स\-तु\-ण्ड\-म\-ण्ड\-ल\-शं\-खा\-वर्ताः~। तेषा\-म\-ङ्गु\-लि\-म\-णि\-व\-न्ध\-गु\-ल्फ\-जा\-नु\-कू\-र्प्प\-रेषु कोराः संधयः~। }\edtext{\skt{कक्ष\-वं\-क्ष\-ण\-द\-श\-नेषू}}{\lemma{\skt{कक्ष\appbrk{}वं\appbrk{}क्ष\appbrk{}ण\appbrk{}द\appbrk{}श\appbrk{}नेषू}}\Afootnote{H; \skt{कक्षा\appbrk{}व\appbrk{}ङ्क्ष\appbrk{}ण\appbrk{}द\appbrk{}श\appbrk{}नेषू}\ A}}\skt{ }\edtext{\skt{दूखलसंज्ञा}}{\lemma{\skt{दूखलसंज्ञा}}\Afootnote{H; \skt{लूखलाः}\ A}}\skt{~। }\edtext{\skt{अंस\-म्पी\-ठ\-गु\-द\-भ\-ग\-नि\-त\-म्वेषु}}{\lemma{\skt{अंसम्}}\Afootnote{H; \skt{अंस}\ A}}\skt{ }\edtext{\skt{सामुद्राः}}{\lemma{\skt{सामुद्राः}}\Afootnote{H; \skt{सामुद्गाः}\ A}}\skt{~। }\edtext{\skt{ग्रीवा\-पृ\-ष्ठ\-वं\-शेषु}}{\lemma{\skt{ग्रीवा\appbrk{}पृ\appbrk{}ष्ठ\appbrk{}वं\appbrk{}शेषु}}\Afootnote{H; \skt{ग्रीवा\appbrk{}पृ\appbrk{}ष्ठ\appbrk{}वं\appbrk{}शयोः}\ A}}\skt{ प्रतराः~। शिरः कटीकपालेषु }\edtext{\skt{तून}}{\lemma{\skt{तून}}\Afootnote{H; \skt{तुन्न}\ A}}\skt{ }\edtext{\skt{सीवनी}}{\lemma{\skt{सीवनी}}\Afootnote{H; \skt{सेवन्}\ A}}\skt{ }\edtext{\skt{संज्ञाः}}{\lemma{\skt{संज्ञाः}}\Afootnote{H; \skt{यः}\ A}}\skt{~। }\edtext{\edtext{\skt{हन्वो\-रू\-भ\-य\-तस्तु}}{\lemma{\skt{हन्वो\appbrk{}रू\appbrk{}भ\appbrk{}य\appbrk{}तस्}}\Afootnote{\skt{हन्वो\appbrk{}रु\appbrk{}भ\appbrk{}य\appbrk{}तस्}\ A; \skt{हन्वोरूभय}\ H\textsuperscript{ac}; \skt{हन्वोरूभय त}\ H\textsuperscript{pc}}}}{\lemma{\skt{तु}}\Afootnote{A; \skt{स्तु}\ H}}\skt{ वायस तुण्डाः~। }\edtext{\skt{कण्ठ\-हृ\-द\-य\-ने\-त्र\-व\-र्त्म्म\-क्लो\-म\-ना\-डीषु}}{\lemma{\skt{कण्ठ\appbrk{}हृ\appbrk{}द\appbrk{}य\appbrk{}ने\appbrk{}त्र\appbrk{}व\appbrk{}र्त्म्म\appbrk{}क्लो\appbrk{}म\appbrk{}ना\appbrk{}डीषु}}\Afootnote{H; \skt{कण्ठ\appbrk{}हृ\appbrk{}द\appbrk{}य\appbrk{}ने\appbrk{}त्र\appbrk{}क्लो\appbrk{}म\appbrk{}ना\appbrk{}डीषु}\ A}}\skt{ }\edtext{\skt{मण्डलसंज्ञाः}}{\lemma{\skt{मण्डलसंज्ञाः}}\Afootnote{H; \skt{मण्डलाः}\ A}}\skt{~। }\edtext{\skt{स्रोत\-शृ\-ङ्गा\-ट\-केषु}}{\lemma{\skt{स्रोत\appbrk{}शृ\appbrk{}ङ्गा\appbrk{}ट\appbrk{}केषु}}\Afootnote{H; \skt{श्रोत्र\appbrk{}शृ\appbrk{}ङ्गा\appbrk{}ट\appbrk{}केषु}\ A}}\skt{ शंखावर्ताः~॥ तेषां नाम\-भि\-रे\-वा\-कृ\-तयः प्रायेण व्याख्याताः~॥}\pend

\pstart\versepar \ledsidenote{\passref{3.5.28}}\edtext{\skt{भवति}}{\lemma{\skt{भवति}}\Afootnote{H; \emph{om.}\ A}}\skt{ }\edtext{\skt{चात्र}}{\lemma{\skt{चात्र}}\Afootnote{H; \emph{om.}\ A}}\skt{~॥ }\pend
\pstart\versepar \skt{}\edtext{\skt{अस्न्थां तु}}{\lemma{\skt{अस्न्थां तु}}\Afootnote{H; \skt{अस्थ्नां तु}\ A}}\skt{ संधयो }\edtext{\skt{ह्येते}}{\lemma{\skt{ह्येते}}\Afootnote{\skt{ह्येते}\ A H}}\skt{ }\edtext{\skt{केवलं}}{\lemma{\skt{केवलं}}\Afootnote{H; \skt{केवलाः}\ A}}\skt{ परिकीर्तिताः~। }\pend
\pstart\versepar \skt{पेशी\-स्ना\-यु\-सि\-रा\-णान्तु }\edtext{\skt{सन्धिसंख्यान्न}}{\lemma{\skt{सन्धिसंख्यान्न}}\Afootnote{H; \skt{सन्धिसङ्ख्या न}\ A}}\skt{ विद्यते~॥}\pend

\pstart \ledsidenote{\passref{3.5.29}}\skt{नव स्नायु शतानि }\edtext{\skt{भवन्ति}}{\lemma{\skt{भवन्ति}}\Afootnote{H; \emph{om.}\ A}}\skt{~। }\edtext{\skt{तेषां}}{\lemma{\skt{तेषां}}\Afootnote{H; \skt{तासां}\ A}}\skt{ शाखासु }\edtext{\skt{षट्छतानि}}{\lemma{\skt{छतानि}}\Afootnote{H; \skt{शतानि}\ A}}\skt{~। द्वे शते त्रिंशच्च कोष्ठे }\edtext{\skt{ग्रीवायां}}{\lemma{\skt{ग्रीवायां}}\Afootnote{H; \skt{ग्रीवां}\ A}}\skt{ }\edtext{\skt{प्रत्युर्ध्वं}}{\lemma{\skt{प्रत्युर्ध्वं}}\Afootnote{H; \skt{प्रत्यूर्ध्वं}\ A}}\skt{ सप्ततिः~। }\edtext{\skt{एकैकस्यां}}{\lemma{\skt{एकैकस्यां}}\Afootnote{H; \skt{एकैकस्यां तु}\ A}}\skt{ पादाङ्गुल्यां }\edtext{\skt{षट्निचितानि}}{\lemma{\skt{षट्निचितानि}}\Afootnote{H; \skt{षण्निचितास्तास्}\ A}}\skt{ }\edtext{\skt{त्रिंशच्च}}{\lemma{\skt{त्रिंशच्च}}\Afootnote{H; \skt{त्रिंशत्}\ A}}\skt{~। तावत्य एव तल\-कू\-र्च्च\-गु\-ल्फेषु तावत्य एव जङ्घायां~। दश जानुनि~। }\edtext{\skt{चत्वारिंशदूर्व्वाः}}{\lemma{\skt{चत्वारिंशदूर्व्वाः}}\Afootnote{H; \skt{चत्वारिंशदूरौ}\ A}}\skt{~। दशवंक्षणे~। }\edtext{\edtext{\edtext{\skt{पञ्चा\-श\-त्छ\-त\-मे\-व\-मे\-क\-स्मि\-न्स\-क्छिनि}}{\lemma{\skt{पञ्चाशत्}}\Afootnote{H; \skt{शतमध्}\ A}}}{\lemma{\skt{छतम्}}\Afootnote{H; \skt{यर्धम्}\ A}}}{\lemma{\skt{सक्छिनि}}\Afootnote{H; \skt{सक्थ्नि}\ A}}\skt{ भवन्ति~। }\edtext{\skt{एतेन}}{\lemma{\skt{एतेन}}\Afootnote{H; \skt{एतेने}\ A}}\skt{ }\edtext{\skt{इतरसक्थि}}{\lemma{\skt{इतरसक्थि}}\Afootnote{H; \skt{तरसक्थि}\ A}}\skt{ }\edtext{\skt{वहू}}{\lemma{\skt{वहू}}\Afootnote{H; \skt{बहू च}\ A}}\skt{ व्याख्यातौ~॥ षष्टिः }\edtext{\skt{कट्याः}}{\lemma{\skt{कट्याः}}\Afootnote{H; \skt{कट्यां}\ A}}\skt{ }\edtext{\skt{पृष्ठेत्वशीतिः}}{\lemma{\skt{पृष्ठेत्वशीतिः}}\Afootnote{H; \skt{पृष्ठेऽशीतिः}\ A}}\skt{~। पार्श्वयोः~। षष्टिः~। उरसि त्रिंशत्~। }\edtext{\skt{षट्त्रिंशद्ग्रीवायां}}{\lemma{\skt{षट्त्रिंशद्}}\Afootnote{H; \skt{षट्त्रिंशद्}\ A}}\skt{~। }\edtext{\skt{मूर्द्ध्रि}}{\lemma{\skt{मूर्द्ध्रि}}\Afootnote{H; \skt{मूर्ध्नि}\ A}}\skt{ }\edtext{\skt{चतुस्त्रिंशत्~। एव}}{\lemma{\skt{चतुस्त्रिंशत्~। एव}}\Afootnote{H; \skt{चतुस्त्रिंशत्; एवं}\ A}}\skt{ नव\-स्ना\-यु\-श\-तानि व्याख्यातानि भवन्ति~। }\edtext{\skt{महास्नायुषु}}{\lemma{\skt{महास्नायुषु}}\Afootnote{H; \emph{om.}\ A}}\skt{ }\edtext{\skt{कण्डरेति}}{\lemma{\skt{कण्डरेति}}\Afootnote{H; \emph{om.}\ A}}\skt{ }\edtext{\skt{संज्ञा}}{\lemma{\skt{संज्ञा}}\Afootnote{H; \emph{om.}\ A}}\skt{~॥}\pend

\pstart\versepar \ledsidenote{\passref{3.5.30}}\skt{भवन्ति चात्र }\edtext{\skt{श्लोकाः}}{\lemma{\skt{श्लोकाः}}\Afootnote{H; \skt{स्नायूश्}\ A}}\skt{~॥ }\edtext{\skt{चतुर्व्विधानि}}{\lemma{\skt{चतुर्व्विधानि}}\Afootnote{H; \skt{चतुर्विधा वि}\ A}}\skt{ }\edtext{\skt{स्नायूनि}}{\lemma{\skt{स्नायूनि}}\Afootnote{H; \skt{द्यात्तास्}\ A}}\skt{ }\edtext{\skt{नास\-ते\-षा\-न्नि\-बोध}}{\lemma{\skt{नासतेषान्}}\Afootnote{H; \skt{तु सर्वा}\ A}}\skt{ मे~। }\pend
\pstart\versepar \skt{}\edtext{\skt{प्रतानवन्ति}}{\lemma{\skt{प्रतानवन्ति}}\Afootnote{H; \skt{प्रतानवत्यो}\ A}}\skt{ }\edtext{\skt{वृतानि}}{\lemma{\skt{वृतानि}}\Afootnote{H; \skt{वृत्ताश्}\ A}}\skt{ }\edtext{\skt{सुसिराणि}}{\lemma{\skt{सुसिराणि}}\Afootnote{H; \emph{om.}\ A}}\skt{ }\edtext{\skt{पृथूनि च}}{\lemma{\skt{पृथूनि च}}\Afootnote{H; \skt{च पृथ्व्यश्च शुषिरास्तथा}\ A}}\skt{~॥}\pend

\pstart\versepar \ledsidenote{\passref{3.5.31}}\edtext{\skt{प्रतानवन्ति}}{\lemma{\skt{प्रतानवन्ति}}\Afootnote{H; \skt{प्रतानवत्यः}\ A}}\skt{ शाखासु }\edtext{\skt{सर्व्वसंधिषु}}{\lemma{\skt{सर्व्वसंधिषु}}\Afootnote{A H\textsuperscript{pc}; \skt{सर्व्वसंधि}\ H\textsuperscript{ac}}}\skt{ चाप्यथ~। }\pend
\pstart\versepar \skt{}\edtext{\skt{वृतानि}}{\lemma{\skt{वृतानि}}\Afootnote{H; \skt{वृत्तास्तु}\ A}}\skt{ कण्डराः सर्व्वा }\edtext{\skt{विज्ञेया}}{\lemma{\skt{विज्ञेया}}\Afootnote{H; \skt{विज्ञेयाः}\ A}}\skt{ }\edtext{\edtext{\skt{कुशलैरिह}}{\lemma{\skt{कुशलैर्}}\Afootnote{H; \skt{कुशलै}\ A}}}{\lemma{\skt{इह}}\Afootnote{H; \skt{रिह}\ A}}\skt{~॥}\pend

\pstart\versepar \ledsidenote{\passref{3.5.32}}\edtext{\skt{आम\-प\-क्वा\-स\-या\-न्तेषु}}{\lemma{\skt{आम\appbrk{}प\appbrk{}क्वा\appbrk{}स\appbrk{}या\appbrk{}न्तेषु}}\Afootnote{H; \skt{आम\appbrk{}प\appbrk{}क्वा\appbrk{}श\appbrk{}या\appbrk{}न्तेषु}\ A}}\skt{ }\edtext{\skt{सुसिराणि}}{\lemma{\skt{सुसिराणि}}\Afootnote{H; \skt{बस्तौ च शुषिराः}\ A}}\skt{ }\edtext{\skt{पृथूनि च}}{\lemma{\skt{पृथूनि च}}\Afootnote{H; \skt{खलु}\ A}}\skt{~। }\pend
\pstart\versepar \skt{पार्श्वोरसि तथा पृष्ठे }\edtext{\skt{पृथूलानि}}{\lemma{\skt{पृथूलानि}}\Afootnote{H; \skt{पृथुलाश्च}\ A}}\skt{ }\edtext{\skt{शिरश्च}}{\lemma{\skt{शिरश्च}}\Afootnote{H; \skt{शिरस्यथ}\ A}}\skt{~॥}\pend

\pstart\versepar \ledsidenote{\passref{3.5.33}}\edtext{\skt{नौर्यथा}}{\lemma{\skt{नौर्}}\Afootnote{H; \skt{नौर्}\ A}}\skt{ }\edtext{\edtext{\skt{फलकास्तीर्ण्णा}}{\lemma{\skt{फलकास्}}\Afootnote{H; \skt{फलका}\ A}}}{\lemma{\skt{तीर्ण्णा}}\Afootnote{H; \skt{स्तीर्णा}\ A}}\skt{ }\edtext{\edtext{\skt{वन्ध\-नै\-र्व\-हु\-भि\-र्युताः}}{\lemma{\skt{वन्धनैर्}}\Afootnote{H; \skt{बन्धनै}\ A}}}{\lemma{\skt{वहुभिर्}}\Afootnote{H; \skt{र्बहुभिर्}\ A}}\skt{~। }\pend
\pstart\versepar \skt{}\edtext{\edtext{\skt{भवेद्भारसहा}}{\lemma{\skt{भवेद्}}\Afootnote{H; \skt{भारक्षमा भवेद}\ A}}}{\lemma{\skt{भारसहा}}\Afootnote{H; \skt{प्सु नृयुक्ता सुसमा}\ A}}\skt{ }\edtext{\skt{दृढ\-स\-न्धि\-स\-मा\-श्रिताः}}{\lemma{\skt{दृढ\appbrk{}स\appbrk{}न्धि\appbrk{}स\appbrk{}मा\appbrk{}श्रिताः}}\Afootnote{H; \skt{हिता}\ A}}\skt{~॥}\pend

\pstart\versepar \ledsidenote{\passref{3.5.34}}\edtext{\skt{एव,}}{\lemma{\skt{एव}}\Afootnote{H; \skt{एवम्}\ A}}\skt{ एव }\edtext{\skt{शरी\-रे\-स्मि\-न्या\-वन्तः}}{\lemma{\skt{शरीरेस्मिन्}}\Afootnote{H; \skt{शरीरेऽस्मिन्}\ A}}\skt{ सन्धयः स्मृताः~। }\pend
\pstart\versepar \skt{}\edtext{\skt{स्नायु\-भि\-र्व्व\-हु\-भि\-र्व्व\-द्दा\-स्तेन}}{\lemma{\skt{व्वद्दास्तेन}}\Afootnote{H; \skt{बद्धास्तेन}\ A}}\skt{ }\edtext{\skt{भारसहाः}}{\lemma{\skt{भारसहाः}}\Afootnote{H; \skt{भारसहा नराः}\ A}}\skt{ }\edtext{\skt{स्मृता}}{\lemma{\skt{स्मृता}}\Afootnote{H; \emph{om.}\ A}}\skt{~॥}\pend

\pstart\versepar \ledsidenote{\passref{3.5.35}}\skt{न }\edtext{\skt{ह्यस्थीनि न}}{\lemma{\skt{अस्थीनि न}}\Afootnote{H; \skt{अस्थीनि न वा}\ A}}\skt{ }\edtext{\skt{पेश्यो वा न}}{\lemma{\skt{पेश्यो वा न}}\Afootnote{H; \skt{पेश्यो न}\ A}}\skt{ सिरा न च सन्धयः~। }\pend
\pstart\versepar \skt{व्यापादितास्तथा }\edtext{\edtext{\skt{हन्युर्यथा}}{\lemma{\skt{हन्युर्}}\Afootnote{H; \skt{हन्यु}\ A}}}{\lemma{\skt{यथा}}\Afootnote{H; \skt{र्यथा}\ A}}\skt{ }\edtext{\skt{स्नायूनि}}{\lemma{\skt{स्नायूनि}}\Afootnote{H; \skt{स्नायुः श}\ A}}\skt{ }\edtext{\skt{देहिनां}}{\lemma{\skt{देहिनां}}\Afootnote{H; \skt{रीरिणम्}\ A}}\skt{~॥}\pend

\pstart\versepar \ledsidenote{\passref{3.5.36}}\edtext{\skt{यः स्नायूनि}}{\lemma{\skt{यः स्नायूनि}}\Afootnote{H; \skt{यः स्नायूः प्र}\ A}}\skt{ }\edtext{\skt{विजानानि स}}{\lemma{\skt{विजानानि स}}\Afootnote{H; \skt{विजानाति}\ A}}\skt{ वाह्यास्यन्तराणि च~। }\pend
\pstart\versepar \skt{स }\edtext{\skt{गुढां}}{\lemma{\skt{गुढां}}\Afootnote{H; \skt{गूढं}\ A}}\skt{ }\edtext{\edtext{\skt{शल्य\-मु\-द्व\-र्न्तु\-न्दे\-हा\-च्छ\-क्रोति}}{\lemma{\skt{शल्यमुद्वर्न्तुन्}}\Afootnote{H; \skt{शल्यमाहर्तुं}\ A}}}{\lemma{\skt{छक्रोति}}\Afootnote{H; \skt{अक्नोति}\ A}}\skt{ }\edtext{\skt{मानवः}}{\lemma{\skt{मानवः}}\Afootnote{H; \skt{देहिनाम्}\ A}}\skt{~॥}\edtext{}{\lemma{\skt{वाह्यास्यन्तराणि च~। स}}\Afootnote{H; \skt{बाह्या\appbrk{}श्चा\appbrk{}भ्य\appbrk{}न्त\appbrk{}रा\appbrk{}स्तथा~। स}\ A}}\pend

\pstart \ledsidenote{\passref{3.5.37}}\skt{पञ्च पेशीशतानि }\edtext{\skt{तासाञ्च}}{\lemma{\skt{तासाञ्च}}\Afootnote{H; \skt{भवन्ति~। तासां}\ A}}\skt{ }\edtext{\skt{चत्वातिं}}{\lemma{\skt{चत्वातिं}}\Afootnote{H; \skt{चत्वारि}\ A}}\skt{ }\edtext{\skt{शक्तानि}}{\lemma{\skt{शक्तानि}}\Afootnote{H; \skt{शतानि}\ A}}\skt{ शाखासु~। कोष्ठे }\edtext{\skt{षष्ठिः}}{\lemma{\skt{षष्ठिः}}\Afootnote{H; \skt{षट्षष्टिः}\ A}}\skt{~। }\edtext{\skt{ग्रीवायां}}{\lemma{\skt{ग्रीवायां}}\Afootnote{H; \skt{ग्रीवां}\ A}}\skt{ }\edtext{\edtext{\skt{प्रत्यु\-र्द्ध\-ञ्च\-त्वा\-रिं\-शत्}}{\lemma{\skt{प्रत्युर्द्धञ्}}\Afootnote{H; \skt{प्रत्यूर्ध्वं}\ A}}\skt{~॥ एकैकस्यां}}{\lemma{\skt{चत्वारिंशत्~॥ एकैकस्यां}}\Afootnote{H; \skt{चतुस्त्रिंशत्~। एकैकस्यां तु}\ A}}\skt{ पादाङ्गुल्यां }\edtext{\skt{तिस्रस्तिस्रस्ताः}}{\lemma{\skt{तिस्रस्तिस्रस्ताः}}\Afootnote{H; \skt{तिस्रस्तिस्रस्ताः}\ A}}\skt{ पञ्चदश~। }\edtext{\skt{पञ्च\-द\-श\-प्र\-पेदे}}{\lemma{\skt{पञ्च\appbrk{}द\appbrk{}श\appbrk{}प्र\appbrk{}पेदे}}\Afootnote{H; \skt{दश प्रपदे}\ A}}\skt{ पादोपरि~। }\edtext{\edtext{\skt{कुर्च्च\-स\-न्नि\-वि\-ष्टा\-स्ता\-वन्त्य}}{\lemma{\skt{कुर्च्च\appbrk{}स\appbrk{}न्नि\appbrk{}वि\appbrk{}ष्टास्}}\Afootnote{H; \skt{कूर्चसन्निविष्टा}\ A}}}{\lemma{\skt{तावन्त्य}}\Afootnote{H; \skt{स्तावत्य}\ A}}\skt{ एव~। दश\-गु\-ल्फ\-त\-लयोः~। }\edtext{\skt{गुल्फजान्वोन्तरे}}{\lemma{\skt{गुल्फजान्वोन्तरे}}\Afootnote{H; \skt{गुल्फजान्वन्तरे}\ A}}\skt{ विंशतिः~। पञ्च }\edtext{\skt{जानिनि}}{\lemma{\skt{जानिनि}}\Afootnote{H; \skt{जानुनि}\ A}}\skt{~। }\edtext{\skt{ऊर्वोद्दश}}{\lemma{\skt{ऊर्वोद्दश}}\Afootnote{H; \skt{विंशतिरूरौ , दश}\ A}}\skt{~। वंक्षणे }\edtext{\skt{दश}}{\lemma{\skt{दश}}\Afootnote{H; \skt{शतम्}\ A}}\skt{~। }\edtext{\edtext{\skt{शत\-मे\-क\-स्मि\-न्स\-क्थिनि}}{\lemma{\skt{शतम्}}\Afootnote{H; \skt{एवम्}\ A}}}{\lemma{\skt{सक्थिनि}}\Afootnote{H; \skt{सक्थ्नि}\ A}}\skt{ }\edtext{\skt{भवन्त्ये न}}{\lemma{\skt{भवन्त्ये न}}\Afootnote{H; \skt{भवति; एते}\ A}}\skt{ }\edtext{\skt{इतर}}{\lemma{\skt{इतर}}\Afootnote{H; \skt{नेतर}\ A}}\skt{ सक्थि वाहू च व्याख्यातौ~॥ }\edtext{\skt{एका}}{\lemma{\skt{एका}}\Afootnote{H; \skt{तिस्रः पायौ, एका}\ A}}\skt{ }\edtext{\skt{मेढ्र}}{\lemma{\skt{मेढ्र}}\Afootnote{H; \skt{मेढ्रे}\ A}}\skt{ सेवन्यां~। }\edtext{\skt{अपरे}}{\lemma{\skt{अपरे}}\Afootnote{H; \skt{चापरा}\ A}}\skt{ द्वे }\edtext{\skt{वृषशयोः}}{\lemma{\skt{वृषशयोः}}\Afootnote{H; \skt{वृषणयोः}\ A}}\skt{~। }\edtext{\skt{गुदे}}{\lemma{\skt{गुदे}}\Afootnote{H; \emph{om.}\ A}}\skt{ }\edtext{\skt{तिस्रः}}{\lemma{\skt{तिस्रः}}\Afootnote{H; \emph{om.}\ A}}\skt{~। }\edtext{\skt{स्फिचोश्च}}{\lemma{\skt{स्फिचोश्च}}\Afootnote{H; \skt{स्फिचोः}\ A}}\skt{ पञ्च पञ्च~। द्वे बस्ति शिरसि~। }\edtext{\skt{सप्तोदरे}}{\lemma{\skt{सप्तोदरे}}\Afootnote{H; \skt{पञ्चोदरे}\ A}}\skt{ नाभ्यामेका~। }\edtext{\skt{पृष्ठो\-र्द्ध\-सं\-नि\-विष्टाः}}{\lemma{\skt{पृष्ठो\appbrk{}र्द्ध\appbrk{}सं\appbrk{}नि\appbrk{}विष्टाः}}\Afootnote{H; \skt{पृष्ठो\appbrk{}र्ध्व\appbrk{}स\appbrk{}न्नि\appbrk{}विष्टाः}\ A}}\skt{ पञ्च पञ्च दीर्घाः~। }\edtext{\skt{पार्श्वयोर्दश}}{\lemma{\skt{पार्श्वयोर्}}\Afootnote{H; \skt{षट्पार्श्वयोः}\ A}}\skt{~। }\edtext{\skt{द्वे}}{\lemma{\skt{द्वे}}\Afootnote{H; \emph{om.}\ A}}\skt{ वक्षसि~। अक्षकां }\edtext{\skt{संप्र\-ति\-स\-मं\-ता\-त्सप्त}}{\lemma{\skt{संप्र\appbrk{}ति\appbrk{}स\appbrk{}मं\appbrk{}तात्}}\Afootnote{H; \skt{सौ प्रति समन्तात्}\ A}}\skt{~। द्वे }\edtext{\skt{हृद्यामासये}}{\lemma{\skt{हृद्यामासये}}\Afootnote{H; \skt{हृदयामाशय}\ A}}\skt{~। }\edtext{\skt{एकाभ्यन्तरतः}}{\lemma{\skt{एकाभ्यन्तरतः}}\Afootnote{H; \skt{योः, षट्य}\ A}}\skt{ }\edtext{\skt{कोष्ठे}}{\lemma{\skt{कोष्ठे}}\Afootnote{H; \skt{कृत्प्ली\appbrk{}हो\appbrk{}ण्डु(\appbrk{}न्दु)\appbrk{}केषु}\ A}}\skt{~। ग्रीवायाञ्च }\edtext{\skt{तस्र}}{\lemma{\skt{तस्र}}\Afootnote{H; \skt{तस्रः}\ A}}\skt{~। }\edtext{\skt{एकैका}}{\lemma{\skt{एकैका}}\Afootnote{H; \skt{अष्टौ हन्वोः , एकैका का}\ A}}\skt{ }\edtext{\skt{गलकालयोः}}{\lemma{\skt{गलकालयोः}}\Afootnote{H; \skt{कलकगलयोः}\ A}}\skt{~। द्वे तालुनि~। }\edtext{\skt{जिह्वायां}}{\lemma{\skt{जिह्वायां}}\Afootnote{H; \skt{एका जिह्वायां}\ A}}\skt{ }\edtext{\skt{षट्}}{\lemma{\skt{षट्}}\Afootnote{H; \emph{om.}\ A}}\skt{~। द्वे }\edtext{\skt{जोष्ठयोः}}{\lemma{\skt{जोष्ठयोः}}\Afootnote{H; \skt{ओष्ठयोः}\ A}}\skt{~। }\edtext{\skt{नासायां}}{\lemma{\skt{नासायां}}\Afootnote{H; \skt{द्वे नासायां}\ A}}\skt{ द्वे~। }\edtext{\edtext{\skt{गण्ड\-यो\-श्च\-त\-स्र\-श्च\-तस्रः}}{\lemma{\skt{गण्डयोश्}}\Afootnote{H; \skt{नेत्रयोः, गण्डयोश्}\ A}}}{\lemma{\skt{चतस्रश्}}\Afootnote{H; \emph{om.}\ A}}\skt{~। }\edtext{\skt{नेत्रयोर्द्ध}}{\lemma{\skt{नेत्रयोर्द्ध}}\Afootnote{H; \emph{om.}\ A}}\skt{~। कर्ण्णयोर्द्वे~। चतस्रो ललाटे~। }\edtext{\edtext{\skt{षट्छिरसीति}}{\lemma{\skt{षट्}}\Afootnote{H; \skt{एका}\ A}}}{\lemma{\skt{छिरसीति}}\Afootnote{H; \skt{शिरसीति}\ A}}\skt{~। }\edtext{\skt{तान्येतानि}}{\lemma{\skt{तान्य्}}\Afootnote{H; \skt{एवम्}\ A}}\skt{ पञ्च पेशीशतानि }\edtext{\edtext{\skt{पुंसाम्भवन्ति}}{\lemma{\skt{पुंसाम्}}\Afootnote{H; \emph{om.}\ A}}}{\lemma{\skt{भवन्ति}}\Afootnote{H; \emph{om.}\ A}}\skt{~।}\pend

\pstart \ledsidenote{\passref{3.5.39}}\edtext{\skt{स्त्रीणा\-ञ्च\-तु\-र्द्द\-शा\-स्य\-धिका}}{\lemma{\skt{चतु\appbrk{}र्द्द\appbrk{}शा\appbrk{}स्य\appbrk{}धिका}}\Afootnote{H; \skt{तु विंशतिरधिका}\ A}}\skt{~। दश तासां स्तन\-यो\-रे\-कै\-क\-स्मि\-न्पञ्च }\edtext{\skt{पञ्च}}{\lemma{\skt{पञ्च}}\Afootnote{H; \skt{पञ्चेति}\ A}}\skt{~। यौवने तासां }\edtext{\skt{परि\-वृ\-द्धि\-र्भ\-वति}}{\lemma{\skt{भवति}}\Afootnote{H; \emph{om.}\ A}}\skt{~। अपत्यपथे चतस्रस्तासां }\edtext{\skt{द्वे}}{\lemma{\skt{द्वे}}\Afootnote{H; \emph{om.}\ A}}\skt{ }\edtext{\skt{प्रसृ\-तेऽ\-भ्य\-न्त\-रतः}}{\lemma{\skt{ऽभ्यन्तरतः}}\Afootnote{A; \skt{भ्यन्तरतः}\ H}}\skt{~। द्वे }\edtext{\skt{मुखास्रिते}}{\lemma{\skt{मुखास्रिते}}\Afootnote{H; \skt{मुखाश्रिते}\ A}}\skt{ }\edtext{\skt{च}}{\lemma{\skt{च}}\Afootnote{H; \skt{बाह्ये च}\ A}}\skt{~। वृत्ते द्वे }\edtext{\skt{गर्भमाश्रिताः}}{\lemma{\skt{गर्भमाश्रिताः}}\Afootnote{\skt{गर्भ\appbrk{}च्छि\appbrk{}द्र\appbrk{}सं\appbrk{}श्रिता}\ A; \skt{गर्ब्भमाश्रिताः}\ H}}\skt{~। }\edtext{\skt{तिस्रः}}{\lemma{\skt{तिस्रः}}\Afootnote{H; \skt{स्तिस्रः}\ A}}\skt{ }\edtext{\edtext{\skt{शुक्रार्भवः}}{\lemma{\skt{शुक्रार्}}\Afootnote{H; \skt{शुक्रा}\ A}}}{\lemma{\skt{भवः}}\Afootnote{H; \skt{र्तव}\ A}}\skt{ प्रवेशिन्यः~। तिस्र एव }\edtext{\skt{गर्भशय्यायां}}{\lemma{\skt{गर्भशय्यायां}}\Afootnote{\skt{पित्त\appbrk{}प\appbrk{}क्वा\appbrk{}श\appbrk{}य\appbrk{}यो\appbrk{}र्मध्ये गर्भशय्या}\ A; \skt{गर्ब्भशय्यायां}\ H}}\skt{ यत्र }\edtext{\skt{गर्भस्तिष्ठति}}{\lemma{\skt{गर्भस्}}\Afootnote{A; \skt{गर्ब्भस्}\ H}}\skt{~।}\pend

\pstart \ledsidenote{\passref{3.5.42}}\edtext{\edtext{\skt{मर्म्म\-सि\-रा\-ध\-म\-नी\-स्रो\-त\-सा\-म\-न्यत्र}}{\lemma{\skt{मर्म्म\appbrk{}सि\appbrk{}रा\appbrk{}ध\appbrk{}म\appbrk{}नी\appbrk{}स्रो\appbrk{}त\appbrk{}साम्}}\Afootnote{H; \skt{मर्म\appbrk{}सि\appbrk{}रा\appbrk{}ध\appbrk{}म\appbrk{}नी\appbrk{}स्रो\appbrk{}तसा}\ A}}}{\lemma{\skt{अन्यत्र}}\Afootnote{H; \skt{मन्यत्र प्र}\ A}}\skt{ विभागः~।}\pend

\pstart\versepar \ledsidenote{\passref{3.5.43}}\edtext{\skt{भवति}}{\lemma{\skt{भवति}}\Afootnote{H; \emph{om.}\ A}}\skt{ }\edtext{\skt{चात्र}}{\lemma{\skt{चात्र}}\Afootnote{H; \emph{om.}\ A}}\skt{~॥ }\pend
\pstart\versepar \skt{}\edtext{\skt{शंखनाभ्यो}}{\lemma{\skt{शंखनाभ्यो}}\Afootnote{H; \skt{शङ्खनाभ्या}\ A}}\skt{ }\edtext{\edtext{\edtext{\edtext{\skt{कृति\-र्यो\-नि\-स्त्रि\-रा\-वर्ता}}{\lemma{\skt{कृतिर्}}\Afootnote{H; \skt{कृति}\ A}}}{\lemma{\skt{योनिस्}}\Afootnote{H; \skt{र्योनि}\ A}}}{\lemma{\skt{त्रिर्}}\Afootnote{H; \skt{स्त्र्}\ A}}}{\lemma{\skt{आवर्ता}}\Afootnote{H; \skt{यावर्ता सा}\ A}}\skt{ प्रकीर्तिता~। }\pend
\pstart\versepar \skt{}\edtext{\edtext{\skt{तस्यास्तृतीये}}{\lemma{\skt{तस्यास्}}\Afootnote{H; \skt{तस्या}\ A}}}{\lemma{\skt{तृतीये}}\Afootnote{H; \skt{स्तृतीये}\ A}}\skt{ त्वावर्ते }\edtext{\skt{गर्भशय्या}}{\lemma{\skt{गर्भशय्या}}\Afootnote{A; \skt{गर्ब्भशय्या}\ H}}\skt{ प्रतिष्ठिता~॥}\pend

\pstart\versepar \ledsidenote{\passref{3.5.44}}\skt{यथा रोहितमत्स्यस्य मुखम्भवति रूपतः~। }\pend
\pstart\versepar \skt{}\edtext{\skt{तत्स्थानञ्च}}{\lemma{\skt{स्थानञ्च}}\Afootnote{H; \skt{संस्थानां}\ A}}\skt{ }\edtext{\skt{यथारूपं}}{\lemma{\skt{यथारूपं}}\Afootnote{H; \skt{तथारूपां}\ A}}\skt{ }\edtext{\edtext{\skt{गर्भ\-श\-य्या\-म्वि\-दु\-र्बुधा}}{\lemma{\skt{गर्भशय्याम्}}\Afootnote{A; \skt{गर्ब्भशय्याम्}\ H}}}{\lemma{\skt{विदुर्बुधा}}\Afootnote{\skt{विदुर्बुधाः}\ A; \skt{विदुर्व्वुधा}\ H}}\skt{~॥}\pend

\pstart\versepar \ledsidenote{\passref{3.5.45}}\skt{आभुग्नोऽभिमुखः शेते }\edtext{\skt{गर्ब्भो}}{\lemma{\skt{गर्ब्भो}}\Afootnote{H; \skt{गर्भो}\ A}}\skt{ }\edtext{\skt{गर्ब्भाशये}}{\lemma{\skt{गर्ब्भाशये}}\Afootnote{H; \skt{गर्भाशये}\ A}}\skt{ स्त्रियाः~। }\pend
\pstart\versepar \skt{स योनिं }\edtext{\skt{याति}}{\lemma{\skt{याति}}\Afootnote{H; \emph{om.}\ A}}\skt{ शिरसा }\edtext{\skt{स्वभा\-वा\-त्प्र\-स\-व\-म्प्रति}}{\lemma{\skt{स्वभावात्}}\Afootnote{H; \skt{याति स्वभावात्}\ A}}\skt{~॥}\pend

\pstart\versepar \ledsidenote{\passref{3.5.46}}\skt{त्वक्पर्यन्तस्य देहस्य }\edtext{\skt{योऽय\-म\-ङ्ग\-वि\-नि\-श्चयः}}{\lemma{\skt{ऽयम्}}\Afootnote{A; \skt{यम्}\ H}}\skt{~। }\pend
\pstart\versepar \skt{शल्यज्ञानादृते }\edtext{\skt{नायं}}{\lemma{\skt{नायं}}\Afootnote{H; \skt{नैष}\ A}}\skt{ }\edtext{\edtext{\skt{वर्ण्यतेऽङ्गेषु}}{\lemma{\skt{वर्ण्यते}}\Afootnote{A; \skt{वर्ण्यते ङ्}\ H}}}{\lemma{\skt{ऽङ्गेषु}}\Afootnote{A; \skt{गेषु}\ H}}\skt{ केषुचित्~॥}\pend

\pstart\versepar \ledsidenote{\passref{3.5.47}}\edtext{\skt{तस्मान्निःसंशयं}}{\lemma{\skt{निःसंशयं}}\Afootnote{A; \skt{निसंशयं}\ H}}\skt{ ज्ञानं }\edtext{\skt{शल्य}}{\lemma{\skt{शल्य}}\Afootnote{H; \skt{हर्त्रा शल्य}\ A}}\skt{ }\edtext{\skt{हर्त्ताभिवाञ्च्छता}}{\lemma{\skt{हर्त्ताभिवाञ्च्छता}}\Afootnote{H; \skt{स्य वाञ्छता}\ A}}\skt{~। }\pend
\pstart\versepar \skt{}\edtext{\skt{धावयित्वा}}{\lemma{\skt{धावयित्वा}}\Afootnote{H; \skt{शोधयित्वा}\ A}}\skt{ मृतं सम्य\-ग्द्र\-ष्ट\-व्योऽ\-ङ्ग\-वि\-नि\-श्चयः~॥}\pend

\pstart\versepar \ledsidenote{\passref{3.5.48}}\skt{प्रत्यक्षतो हि यद्दृष्टं शास्त्रदृष्टञ्च यद्भवेत्~। }\pend
\pstart\versepar \skt{}\edtext{\skt{समा\-ग\-त\-न्त\-दु\-भ\-य\-म्भूयो}}{\lemma{\skt{समागतन्}}\Afootnote{H; \skt{समासतस्}\ A}}\skt{ }\edtext{\edtext{\skt{ज्ञान\-म्वि\-व\-र्द्ध\-येत्}}{\lemma{\skt{ज्ञानम्}}\Afootnote{H; \skt{ज्ञान}\ A}}}{\lemma{\skt{विवर्द्धयेत्}}\Afootnote{H; \skt{विवर्धनम्}\ A}}\skt{~॥}\pend

\pstart \ledsidenote{\passref{3.5.49}}\edtext{\skt{तस्मा\-त्स\-म\-स्त\-गा\-त्र\-म\-वि\-ह\-त\-म\-दी\-र्घ\-व्या\-धि\-पी\-डितं}}{\lemma{\skt{अविहतम्}}\Afootnote{H; \skt{अवि\appbrk{}षो\appbrk{}प\appbrk{}ह\appbrk{}तम्}\ A}}\skt{ }\edtext{\edtext{\skt{पुरु\-ष\-म\-प\-हृ\-त्या\-प\-गायां}}{\lemma{\skt{पुरुषम्}}\Afootnote{H; \skt{अवर्षशतिकं निःसृ\appbrk{}ष्टा\appbrk{}न्त्र\appbrk{}पु\appbrk{}रीषं पुरुषम्}\ A}}}{\lemma{\skt{अप\appbrk{}हृ\appbrk{}त्या\appbrk{}प\appbrk{}गायां}}\Afootnote{H; \skt{आव\appbrk{}ह\appbrk{}न्त्या\appbrk{}मा\appbrk{}प\appbrk{}गायां}\ A}}\skt{~। }\edtext{\skt{विबद्धं}}{\lemma{\skt{विबद्धं}}\Afootnote{H; \skt{निबद्धं}\ A}}\skt{ }\edtext{\skt{मुञ्जबल्वजैः}}{\lemma{\skt{मुञ्जबल्वजैः}}\Afootnote{H; \emph{om.}\ A}}\skt{ }\edtext{\skt{पञ्ज\-र\-स्थ\-म\-प्र\-काशे}}{\lemma{\skt{अप्रकाशे}}\Afootnote{H; \skt{मुञ्ज\appbrk{}व\appbrk{}ल्क\appbrk{}ल\appbrk{}कु\appbrk{}श\appbrk{}श\appbrk{}णा\appbrk{}दी\appbrk{}ना\appbrk{}म\appbrk{}न्य\appbrk{}त\appbrk{}मे\appbrk{}ना\appbrk{}वे\appbrk{}ष्टि\appbrk{}ता\appbrk{}ङ्ग\appbrk{}म\appbrk{}प्र\appbrk{}काशे}\ A}}\skt{ देशे कोथयेत्~। }\edtext{\skt{सम्य\-क्प्र\-कु\-थि\-त\-ञ्चो\-द्धृ\-त्या\-य\-त\-देहं}}{\lemma{\skt{चोद्धृ\appbrk{}त्या\appbrk{}य\appbrk{}त\appbrk{}देहं}}\Afootnote{H; \skt{चोद्धृत्य, ततो देहं}\ A}}\skt{ }\edtext{\skt{कृत्वो\-शी\-र\-तृ\-ण\-वे\-णु\-कु\-र्च्चा\-दी\-ना\-म\-न्य\-त\-मेन}}{\lemma{\skt{कृत्वो\appbrk{}शी\appbrk{}र\appbrk{}तृ\appbrk{}ण\appbrk{}वे\appbrk{}णु\appbrk{}कु\appbrk{}र्च्चा\appbrk{}दी\appbrk{}नाम्}}\Afootnote{H; \skt{सप्त\appbrk{}रा\appbrk{}त्रा\appbrk{}दु\appbrk{}शी\appbrk{}र\appbrk{}बा\appbrk{}ल\appbrk{}वे\appbrk{}णु\appbrk{}ब\appbrk{}ल्व\appbrk{}ज\appbrk{}कू\appbrk{}र्चा\appbrk{}नाम्}\ A}}\skt{ शनैः }\edtext{\skt{शनैरेव}}{\lemma{\skt{एव}}\Afootnote{H; \skt{अव}\ A}}\skt{ }\edtext{\skt{घृष्य}}{\lemma{\skt{घृष्य}}\Afootnote{H; \skt{घर्षयंस्}\ A}}\skt{ त्वगा\-दी\-न्स\-र्व्वा\-नेव }\edtext{\skt{बाह्या\-भ्य\-न्त\-रा\-न\-ङ्ग\-प्र\-त्य\-ङ्ग\-वि\-शे\-षा\-न्य\-थो\-क्ता\-नु\-प\-ल\-क्ष\-ये\-च्च\-क्षुषा}}{\lemma{\skt{यथो\appbrk{}क्ता\appbrk{}नु\appbrk{}प\appbrk{}ल\appbrk{}क्ष\appbrk{}येच्}}\Afootnote{H; \skt{यथो\appbrk{}क्ता\appbrk{}न्ल\appbrk{}क्ष\appbrk{}येच्}\ A}}\skt{~॥}\footnoteA{On \emph{ghṛṣya} for \emph{ghṛṣṭvā}, see Oberlies (2003): 9.7.2 (p. 283).}\pend

\pstart\versepar \ledsidenote{\passref{3.5.50}}\skt{श्लोकौ~॥ }\pend
\pstart\versepar \skt{न }\edtext{\skt{शक्यञ्चक्षुषा}}{\lemma{\skt{शक्यञ्}}\Afootnote{H; \skt{शक्यश्}\ A}}\skt{ द्रष्टुं }\edtext{\skt{सौक्ष्म्यासौक्ष्म्यतरं}}{\lemma{\skt{सौक्ष्म्यासौक्ष्म्यतरं}}\Afootnote{H; \skt{देहे सूक्ष्मतमो}\ A}}\skt{ }\edtext{\skt{विदुः}}{\lemma{\skt{विदुः}}\Afootnote{H; \skt{विभुः}\ A}}\skt{~। }\pend
\pstart\versepar \skt{दृश्यते }\edtext{\edtext{\skt{ज्ञान\-च\-क्षु\-र्भि\-स्त\-प\-श्च\-क्षु\-र्भि\-रेव}}{\lemma{\skt{तपश्चक्षुर्भिर्}}\Afootnote{\skt{तपश्चक्षुर्भिर्}\ A; \skt{तपश्चक्षुर्ब्भिर्}\ H}}}{\lemma{\skt{एव}}\Afootnote{A H\textsuperscript{pc}; \skt{ए वा}\ H\textsuperscript{ac}}}\skt{ }\edtext{\skt{वा}}{\lemma{\skt{वा}}\Afootnote{H\textsuperscript{pc}; \skt{च}\ A; \emph{om.}\ H\textsuperscript{ac}}}\skt{~॥}\edtext{}{\lemma{\skt{श्लोकौ~॥ न}}\Afootnote{H; \skt{न}\ A}}\pend

\pstart\versepar \ledsidenote{\passref{3.5.51}}\skt{शरीरे }\edtext{\skt{चापि}}{\lemma{\skt{चापि}}\Afootnote{H; \skt{चैव}\ A}}\skt{ शास्त्रे च }\edtext{\skt{दृष्टार्थस्तु}}{\lemma{\skt{तु}}\Afootnote{H; \skt{स्याद्}\ A}}\skt{ विशारदः~। }\pend
\pstart\versepar \skt{दृष्टश्रुताभ्यां सन्देहं }\edtext{\skt{व्यपोह्यारभते}}{\lemma{\skt{व्यपोह्यारभते}}\Afootnote{H; \skt{अवा\appbrk{}पो\appbrk{}ह्या\appbrk{}च\appbrk{}रेत्}\ A}}\skt{ }\edtext{\edtext{\skt{क्रियामिति}}{\lemma{\skt{क्रियाम्}}\Afootnote{H; \skt{क्रियाः}\ A}}}{\lemma{\skt{इति}}\Afootnote{H; \emph{om.}\ A}}\skt{~॥~॥}\pend

\pstart \ledsidenote{\passref{3.5.trailer}}\edtext{\skt{इति शारीरे}}{\lemma{\skt{इति शारीरे}}\Afootnote{H; \skt{इति सुश्रु\appbrk{}त\appbrk{}सं\appbrk{}हि\appbrk{}तायां शारीरस्थाने शरी\appbrk{}र\appbrk{}सं\appbrk{}ख्या\appbrk{}व्या\appbrk{}क\appbrk{}र\appbrk{}ण\appbrk{}शा\appbrk{}रीरं}\ A}}\skt{ }\edtext{\skt{पञ्चमो}}{\lemma{\skt{पञ्चमो}}\Afootnote{H; \skt{नाम पञ्चमो}\ A}}\skt{ }\edtext{\skt{ध्यायः}}{\lemma{\skt{ध्यायः}}\Afootnote{H; \skt{ऽध्यायः}\ A}}\skt{~॥~॥}\pend

\endnumbering
\beginnumbering
\eledsection*{Adhyāya 6}
\pstart \ledsidenote{\passref{3.6.1}}\skt{अथातः }\edtext{\skt{प्रत्ये\-क\-म\-र्म्म\-नि\-र्द्देशं}}{\lemma{\skt{प्रत्ये\appbrk{}क\appbrk{}म\appbrk{}र्म्म\appbrk{}नि\appbrk{}र्द्देशं}}\Afootnote{H; \skt{प्रत्ये\appbrk{}क\appbrk{}म\appbrk{}र्म\appbrk{}नि\appbrk{}र्देशं}\ A}}\skt{ शारी\-र\-म्व्या\-ख्या\-स्यामः~॥}\pend

\pstart \ledsidenote{\passref{3.6.3}}\skt{सप्तो\-त्त\-र\-म्म\-र्म्म\-शतं~। }\edtext{\skt{तानि}}{\lemma{\skt{तानि}}\Afootnote{H; \skt{ता}\ A}}\skt{ }\edtext{\skt{पञ्चविधानि}}{\lemma{\skt{पञ्चविधानि}}\Afootnote{H; \skt{नि}\ A}}\skt{ मर्म्माणि }\edtext{\skt{भवन्ति}}{\lemma{\skt{भवन्ति}}\Afootnote{H; \skt{पञ्चात्मकानि भवन्ति}\ A}}\skt{~। तद्यथा~। मांसमर्म्माणि~। }\edtext{\skt{शिरामर्म्माणि}}{\lemma{\skt{शिरामर्म्माणि}}\Afootnote{H; \skt{सिरामर्माणि}\ A}}\skt{~। स्नायुमर्म्माणि, अस्थिमर्म्माणि~। सन्धिमर्म्माणि~। }\edtext{\skt{न}}{\lemma{\skt{न}}\Afootnote{H; \skt{चेति~। न}\ A}}\skt{ खलु }\edtext{\skt{मांस\-शि\-रा\-स्ना\-य्व\-स्थि\-स\-न्धि\-व्य\-ति\-रे\-के\-णा\-न्यानि}}{\lemma{\skt{मांस\appbrk{}शि\appbrk{}रा\appbrk{}स्ना\appbrk{}य्व\appbrk{}स्थि\appbrk{}स\appbrk{}न्धि\appbrk{}व्य\appbrk{}ति\appbrk{}रे\appbrk{}के\appbrk{}णा\appbrk{}न्यानि}}\Afootnote{H; \skt{मांस\appbrk{}सि\appbrk{}रा\appbrk{}स्ना\appbrk{}य्व\appbrk{}स्थि\appbrk{}स\appbrk{}न्धि\appbrk{}व्य\appbrk{}ति\appbrk{}रे\appbrk{}के\appbrk{}णा\appbrk{}न्यानि}\ A}}\skt{ मर्म्माणि भवन्ति~॥ यस्मा\-न्नो\-प\-ल\-भ्यन्ते~॥}\pend

\pstart \ledsidenote{\passref{3.6.4}}\skt{तत्रैकादश मांसमर्म्माणि, }\edtext{\edtext{\skt{एक\-च\-त्वा\-रिं\-श\-च्छि\-रा\-म\-र्म्माणि}}{\lemma{\skt{एक\appbrk{}च\appbrk{}त्वा\appbrk{}रिं\appbrk{}शच्}}\Afootnote{H; \skt{एक\appbrk{}च\appbrk{}त्वा\appbrk{}रिं\appbrk{}शत्}\ A}}}{\lemma{\skt{छिरामर्म्माणि}}\Afootnote{H; \skt{सिरामर्माणि}\ A}}\skt{~। }\edtext{\skt{सप्ताविंशतिः}}{\lemma{\skt{सप्ताविंशतिः}}\Afootnote{H; \skt{सप्तविंशतिः}\ A}}\skt{ स्नायुमर्म्माणि~। अष्टा\-व\-स्थि\-म\-र्म्माणि~। विंशतिः }\edtext{\skt{सन्धिमर्म्माणीति}}{\lemma{\skt{सन्धिमर्म्माणीति}}\Afootnote{H; \skt{सन्धिमर्माणि चेति}\ A}}\skt{~। तदे\-त\-त्स\-प्तो\-त्त\-र\-म्म\-र्म्म\-शतं~॥}\pend

\pstart \ledsidenote{\passref{3.6.5}}\edtext{\edtext{\skt{एषा\-मे\-का\-द\-शै\-क\-स्मि\-न्स\-क्थिनि}}{\lemma{\skt{एषाम्}}\Afootnote{H; \skt{तेषाम्}\ A}}}{\lemma{\skt{सक्थिनि}}\Afootnote{H; \skt{सक्थ्नि}\ A}}\skt{ भवन्ति~। एते\-ने\-त\-र\-सक्थि बाहू च }\edtext{\skt{व्याख्यातौ च}}{\lemma{\skt{व्याख्यातौ च}}\Afootnote{H\textsuperscript{ac}; \skt{व्याख्यातौ}\ A H\textsuperscript{pc}}}\skt{ }\edtext{\skt{भवतः}}{\lemma{\skt{भवतः}}\Afootnote{H; \emph{om.}\ A}}\skt{~। उद\-रो\-र\-सो\-र्द्द्वा\-दश~। चतुर्द्दश पृष्ठे~। }\edtext{\skt{ग्रीवायाम्प्रत्यूर्द्ध्वं}}{\lemma{\skt{ग्रीवायाम्}}\Afootnote{H; \skt{ग्रीवां}\ A}}\skt{ सप्तत्रिंशत्~।}\pend

\pstart \ledsidenote{\passref{3.6.6}}\skt{तत्र सक्थिमर्म्माणि~। }\edtext{\skt{क्षिप्र\-त\-ल\-हृ\-द\-य\-कु\-र्च्च\-कू\-र्च्च\-शि\-रो\-गु\-ल्फे\-न्द्र\-ब\-स्ति\-जा\-नु\-श्चा\-ण्यू\-र्वी\-लो\-हि\-ता\-क्षाणि}}{\lemma{\skt{क्षिप्र\appbrk{}त\appbrk{}ल\appbrk{}हृ\appbrk{}द\appbrk{}य\appbrk{}कु\appbrk{}र्च्च\appbrk{}कू\appbrk{}र्च्च\appbrk{}शि\appbrk{}रो\appbrk{}गु\appbrk{}ल्फे\appbrk{}न्द्र\appbrk{}ब\appbrk{}स्ति\appbrk{}जा\appbrk{}नु\appbrk{}श्चा\appbrk{}ण्यू\appbrk{}र्वी\appbrk{}लो\appbrk{}हि\appbrk{}ता\appbrk{}क्षाणि}}\Afootnote{H; \skt{क्षिप्र\appbrk{}त\appbrk{}ल\appbrk{}हृ\appbrk{}द\appbrk{}य\appbrk{}कू\appbrk{}र्च\appbrk{}कू\appbrk{}र्च\appbrk{}शि\appbrk{}रो\appbrk{}गु\appbrk{}ल्फे\appbrk{}न्द्र\appbrk{}ब\appbrk{}स्ति\appbrk{}जा\appbrk{}न्वा\appbrk{}ण्यु\appbrk{}र्वि\appbrk{}लो\appbrk{}हि\appbrk{}ता\appbrk{}क्षाणि}\ A}}\skt{ }\edtext{\skt{चिटिपञ्चेति}}{\lemma{\skt{चिटिपञ्}}\Afootnote{H; \skt{विटपं}\ A}}\skt{~। }\edtext{\skt{एते\-ने\-त\-र\-त्स\-क्थि\-बाहू}}{\lemma{\skt{एते\appbrk{}ने\appbrk{}त\appbrk{}र\appbrk{}त्स\appbrk{}क्थि\appbrk{}बाहू}}\Afootnote{H; \skt{एते\appbrk{}ने\appbrk{}त\appbrk{}र\appbrk{}त्सक्थि}\ A}}\skt{ }\edtext{\skt{व्याख्यातौ}}{\lemma{\skt{व्याख्यातौ}}\Afootnote{H; \skt{व्याख्यातम्}\ A}}\skt{~॥ उदरोरसोस्तु }\edtext{\edtext{\skt{मर्म्मा\-ण्य\-नु\-व्या\-ख्या\-स्यामः}}{\lemma{\skt{मर्म्माण्य्}}\Afootnote{H; \emph{om.}\ A}}}{\lemma{\skt{अनुव्याख्यास्यामः}}\Afootnote{H; \emph{om.}\ A}}\skt{~॥ }\edtext{\skt{गुदो}}{\lemma{\skt{गुदो}}\Afootnote{H; \skt{गुद}\ A}}\skt{ }\edtext{\edtext{\edtext{\skt{बस्ति\-र्न्ना\-भि\-र्हृ\-दयं}}{\lemma{\skt{बस्तिर्}}\Afootnote{H; \skt{बस्ति}\ A}}}{\lemma{\skt{न्नाभिर्}}\Afootnote{H; \skt{नाभि}\ A}}}{\lemma{\skt{हृदयं}}\Afootnote{H; \skt{हृदय}\ A}}\skt{ }\edtext{\skt{स्तनमूले}}{\lemma{\skt{स्तनमूले}}\Afootnote{H; \skt{स्तनमूल}\ A}}\skt{ }\edtext{\skt{स्तनलोहिते}}{\lemma{\skt{स्तनलोहिते}}\Afootnote{H; \skt{स्तनरोहिता}\ A}}\skt{ }\edtext{\skt{अपालाये}}{\lemma{\skt{अपालाये}}\Afootnote{H; \skt{पलापान्}\ A}}\skt{ }\edtext{\skt{अपस्तम्भे}}{\lemma{\skt{अपस्तम्भे}}\Afootnote{H; \skt{यपस्तम्भौ}\ A}}\skt{ चेति~॥ पृष्ठमर्म्माणि तु~। कटीकतरुण~। }\edtext{\skt{कुकु\-न्द\-र\-नि\-तम्बे}}{\lemma{\skt{कुकु\appbrk{}न्द\appbrk{}र\appbrk{}नि\appbrk{}तम्बे}}\Afootnote{H; \skt{कुकु\appbrk{}न्द\appbrk{}र\appbrk{}नि\appbrk{}तम्ब}\ A}}\skt{ }\edtext{\skt{पार्श्व\-स\-न्धि\-बृ\-ह\-त्यं\-स\-फ\-ल\-क्र\-क\-न्यासौ}}{\lemma{\skt{पार्श्व\appbrk{}स\appbrk{}न्धि\appbrk{}बृ\appbrk{}ह\appbrk{}त्यं\appbrk{}स\appbrk{}फ\appbrk{}ल\appbrk{}क्र\appbrk{}क\appbrk{}न्यासौ}}\Afootnote{H; \skt{पार्श्व\appbrk{}स\appbrk{}न्धि\appbrk{}बृ\appbrk{}ह\appbrk{}त्यं\appbrk{}स\appbrk{}फ\appbrk{}ल\appbrk{}का\appbrk{}न्यंसौ}\ A}}\skt{ }\edtext{\skt{चेति}}{\lemma{\skt{चेति}}\Afootnote{H; \skt{चेति~। बाहुमर्माणि तु शिप्र\appbrk{}त\appbrk{}ल\appbrk{}हृ\appbrk{}द\appbrk{}य\appbrk{}कू\appbrk{}र्च\appbrk{}कू\appbrk{}र्च\appbrk{}शि\appbrk{}रो\appbrk{}म\appbrk{}णि\appbrk{}ब\appbrk{}न्धे\appbrk{}न्द्र\appbrk{}ब\appbrk{}स्ति\appbrk{}कू\appbrk{}र्प\appbrk{}रा\appbrk{}ण्यु\appbrk{}र्वी\appbrk{}लो\appbrk{}हि\appbrk{}ता\appbrk{}क्षाणि कक्षधरं चेति}\ A}}\skt{~॥ }\edtext{\skt{जत्रुण}}{\lemma{\skt{जत्रुण}}\Afootnote{H; \skt{एतेनेतरो बाहुर्व्याख्यातः~। जत्रुण}\ A}}\skt{ ऊर्ध्वं चतस्रो धमन्यो ष्टौ मातृकाः~। द्वे कृकाटिके~। द्वे विधुरे~। द्वे फणे~। द्वावपाङ्गौ~। }\edtext{\skt{द्वौ}}{\lemma{\skt{द्वौ}}\Afootnote{H; \emph{om.}\ A}}\skt{ }\edtext{\skt{शंखौ}}{\lemma{\skt{शंखौ}}\Afootnote{H; \emph{om.}\ A}}\skt{~। द्वावावर्तौ~। द्वावुत्क्षेपौ~। }\edtext{\skt{एका}}{\lemma{\skt{एका}}\Afootnote{H; \skt{द्वौ शङ्खावेका}\ A}}\skt{ }\edtext{\skt{स्थपणी}}{\lemma{\skt{स्थपणी}}\Afootnote{H; \skt{स्थपनी}\ A}}\skt{~। पञ्च सीमन्ताः~। चत्वारि }\edtext{\skt{शृङ्गा\-ट\-का\-न्ये\-कोऽ\-धि\-प\-ति\-रिति}}{\lemma{\skt{शृङ्गाटकान्येको}}\Afootnote{H; \skt{शृङ्गाटकान्येको}\ A}}\skt{~॥}\pend

\pstart \ledsidenote{\passref{3.6.7}}\skt{तत्र तल\-हृ\-द\-ये\-न्द्र\-ब\-स्ति\-गु\-द\-स्त\-न\-रो\-हि\-तानि~। मांसमर्म्माणि~। नीलधमनी~। मातृका }\edtext{\skt{सिराः}}{\lemma{\skt{सिराः}}\Afootnote{H; \emph{om.}\ A}}\skt{ }\edtext{\edtext{\edtext{\skt{शृङ्गा\-ट\-का\-पा\-ङ्ग\-स्थ\-प\-णी\-फ\-ण\-स्त\-न\-मू\-ला\-प\-ला\-पा\-प\-स्त\-म्भ\-हृ\-द\-य\-ना\-भि\-पा\-र्श्व\-स\-न्धि\-बृ\-ह\-ती\-लो\-हि\-ता\-क्षा\-ण्यू\-र्व्व्या\-स्सि\-रा\-म\-र्म्माणि}}{\lemma{\skt{शृङ्गा\appbrk{}ट\appbrk{}का\appbrk{}पा\appbrk{}ङ्ग\appbrk{}स्थ\appbrk{}प\appbrk{}णी\appbrk{}फ\appbrk{}ण\appbrk{}स्त\appbrk{}न\appbrk{}मू\appbrk{}ला\appbrk{}प\appbrk{}ला\appbrk{}पा\appbrk{}प\appbrk{}स्त\appbrk{}म्भ\appbrk{}हृ\appbrk{}द\appbrk{}य\appbrk{}ना\appbrk{}भि\appbrk{}पा\appbrk{}र्श्व\appbrk{}स\appbrk{}न्धि\appbrk{}बृ\appbrk{}ह\appbrk{}ती\appbrk{}लो\appbrk{}हि\appbrk{}ता\appbrk{}क्षाण्य्}}\Afootnote{H; \skt{शृङ्गा\appbrk{}ट\appbrk{}का\appbrk{}पा\appbrk{}ङ्ग\appbrk{}स्थ\appbrk{}प\appbrk{}नी\appbrk{}फ\appbrk{}ण\appbrk{}स्त\appbrk{}न\appbrk{}मू\appbrk{}ला\appbrk{}प\appbrk{}ला\appbrk{}पा\appbrk{}प\appbrk{}स्त\appbrk{}म्भ\appbrk{}हृ\appbrk{}द\appbrk{}य\appbrk{}ना\appbrk{}भि\appbrk{}पा\appbrk{}र्श्व\appbrk{}स\appbrk{}न्धि\appbrk{}बृ\appbrk{}ह\appbrk{}ती\appbrk{}लो\appbrk{}हि\appbrk{}ताक्षो}\ A}}}{\lemma{\skt{ऊर्व्व्यास्}}\Afootnote{H; \skt{र्व्यः}\ A}}}{\lemma{\skt{सिरामर्म्माणि}}\Afootnote{H; \skt{सिरामर्माणि, आणी(णि)}\ A}}\skt{~। }\edtext{\skt{चिटि\-प\-क\-क्ष\-ध\-र\-कुर्च्चः}}{\lemma{\skt{चिटि\appbrk{}प\appbrk{}क\appbrk{}क्ष\appbrk{}ध\appbrk{}र\appbrk{}कुर्च्चः}}\Afootnote{H; \skt{विट\appbrk{}प\appbrk{}क\appbrk{}क्ष\appbrk{}ध\appbrk{}र\appbrk{}कूर्च}\ A}}\skt{~। }\edtext{\edtext{\skt{कूर्च्च\-शि\-रो\-ब\-स्ति\-कु\-क्षा\-ण्य\-न्स\-वि\-धु\-रो\-त्क्षेपाः}}{\lemma{\skt{कूर्च्च\appbrk{}शि\appbrk{}रो\appbrk{}ब\appbrk{}स्ति\appbrk{}कु\appbrk{}क्षाण्य्}}\Afootnote{H; \skt{कूर्च\appbrk{}शि\appbrk{}रो\appbrk{}ब\appbrk{}स्ति\appbrk{}क्षिप्}\ A}}}{\lemma{\skt{अन्स\appbrk{}वि\appbrk{}धु\appbrk{}रो\appbrk{}त्क्षेपाः}}\Afootnote{H; \skt{रांस\appbrk{}वि\appbrk{}धु\appbrk{}रो\appbrk{}त्क्षेपाः}\ A}}\skt{ स्नायुमर्म्माणि~। }\edtext{\edtext{\skt{कटी\-क\-त\-रु\-ण\-नि\-त\-म्बा\-न्स\-फ\-ल\-क\-शं\-खा\-न्य\-स्थि\-म\-र्म्माणि}}{\lemma{\skt{कटी\appbrk{}क\appbrk{}त\appbrk{}रु\appbrk{}ण\appbrk{}नि\appbrk{}त\appbrk{}म्बा\appbrk{}न्स\appbrk{}फ\appbrk{}ल\appbrk{}क\appbrk{}शं\appbrk{}खान्य्}}\Afootnote{H; \skt{कटी\appbrk{}क\appbrk{}त\appbrk{}रु\appbrk{}ण\appbrk{}नि\appbrk{}त\appbrk{}म्बां\appbrk{}स\appbrk{}फ\appbrk{}ल\appbrk{}क\appbrk{}श\appbrk{}ङ्खास्त्}\ A}}}{\lemma{\skt{अस्थिमर्म्माणि}}\Afootnote{H; \skt{वस्थिमर्माणि}\ A}}\skt{~। }\edtext{\skt{जानु\-कु\-र्प्प\-र\-सी\-म\-न्ताऽ\-धि\-प\-ति\-गु\-ल्फ\-म\-णि\-ब\-न्ध\-कु\-कु\-न्द\-रा\-व\-र्त्त\-कृ\-का\-टि\-का\-श्चेति}}{\lemma{\skt{जानु\appbrk{}कु\appbrk{}र्प्प\appbrk{}र\appbrk{}सी\appbrk{}म\appbrk{}न्ताऽ\appbrk{}धि\appbrk{}प\appbrk{}ति\appbrk{}गु\appbrk{}ल्फ\appbrk{}म\appbrk{}णि\appbrk{}ब\appbrk{}न्ध\appbrk{}कु\appbrk{}कु\appbrk{}न्द\appbrk{}रा\appbrk{}व\appbrk{}र्त्त\appbrk{}कृ\appbrk{}का\appbrk{}टि\appbrk{}काश्}}\Afootnote{H; \skt{जानु\appbrk{}कू\appbrk{}र्प\appbrk{}र\appbrk{}सी\appbrk{}म\appbrk{}न्ता\appbrk{}धि\appbrk{}प\appbrk{}ति\appbrk{}गु\appbrk{}ल्फ\appbrk{}म\appbrk{}णि\appbrk{}ब\appbrk{}न्ध\appbrk{}कु\appbrk{}कु\appbrk{}न्द\appbrk{}रा\appbrk{}व\appbrk{}र्त\appbrk{}कृ\appbrk{}का\appbrk{}टि\appbrk{}काश्}\ A}}\skt{~। सन्धिमर्म्माणि~॥}\pend

\pstart \ledsidenote{\passref{3.6.8}}\skt{तान्येतानि }\edtext{\skt{मर्म्माणि}}{\lemma{\skt{मर्म्माणि}}\Afootnote{H; \emph{om.}\ A}}\skt{ }\edtext{\skt{पञ्चविधानि}}{\lemma{\skt{पञ्चविधानि}}\Afootnote{H; \skt{पञ्चविकल्पानि}\ A}}\skt{ भवन्ति~॥ तद्यथा~। सद्यः\-प्रा\-ण\-ह\-राणि~। काला\-न्त\-र\-प्रा\-ण\-ह\-राणि~। विशल्यघ्नानि~। वैकल्यकराणि~। रुजाकराणि चेति~॥ तत्र सद्यः\-प्रा\-ण\-ह\-रा\-ण्ये\-को\-न\-विं\-शतिः~। काला\-न्त\-र\-प्रा\-ण\-ह\-राणि }\edtext{\skt{त्रयस्त्रिंशत्}}{\lemma{\skt{त्रयस्त्रिंशत्}}\Afootnote{H; \skt{त्रयस्त्रिंशत्, त्}\ A}}\skt{~। }\edtext{\skt{विशल्यघ्नानि}}{\lemma{\skt{विशल्यघ्नानि}}\Afootnote{H; \skt{रीणि विशल्यघ्नानि}\ A}}\skt{ }\edtext{\skt{त्रीणि}}{\lemma{\skt{त्रीणि}}\Afootnote{H; \emph{om.}\ A}}\skt{~। चतु\-श्च\-त्वा\-रिं\-श\-द्वै\-क\-ल्य\-क\-राणि~। अष्टौ रुजाकराणीति~॥}\pend

\pstart\versepar \ledsidenote{\passref{3.6.9}}\skt{भवन्ति चात्र~॥ }\pend
\pstart\versepar \skt{शृङ्गा\-ट\-का\-न्य\-धि\-पतिः शंखौ कण्ठसिरा }\edtext{\skt{गुदः}}{\lemma{\skt{गुदः}}\Afootnote{H; \skt{गुदम्}\ A}}\skt{~। }\pend
\pstart\versepar \skt{}\edtext{\skt{हृद्बस्ति}}{\lemma{\skt{हृद्बस्ति}}\Afootnote{H; \skt{हृदयं बस्ति}\ A}}\skt{ }\edtext{\skt{सह वै}}{\lemma{\skt{सह वै}}\Afootnote{H; \emph{om.}\ A}}\skt{ }\edtext{\skt{नाभ्यां}}{\lemma{\skt{नाभ्यां}}\Afootnote{H; \skt{नाभ्यौ(भी) च}\ A}}\skt{ घ्नन्ति सद्यो हतानि तु~॥}\pend

\pstart\versepar \ledsidenote{\passref{3.6.10}}\skt{वक्षोमर्म्माणि }\edtext{\edtext{\skt{सीमन्तास्तलः}}{\lemma{\skt{सीमन्तास्}}\Afootnote{H; \skt{सीमन्त}\ A}}}{\lemma{\skt{तलः}}\Afootnote{H; \skt{तल}\ A}}\skt{ क्षिप्रेन्द्रबस्तयः~। }\pend
\pstart\versepar \skt{कटीकतरुणे सन्धी }\edtext{\skt{पार्श्वजो}}{\lemma{\skt{पार्श्वजो}}\Afootnote{H; \skt{पार्श्वजौ}\ A}}\skt{ बृहती च }\edtext{\skt{ये}}{\lemma{\skt{ये}}\Afootnote{H; \skt{या}\ A}}\skt{~॥}\pend

\pstart\versepar \ledsidenote{\passref{3.6.11}}\skt{नितम्बाविति चैतानि काला\-न्त\-र\-ह\-राणि तु~। }\pend
\pstart\versepar \skt{उत्क्षेपौ }\edtext{\skt{स्थपणी}}{\lemma{\skt{स्थपणी}}\Afootnote{H; \skt{स्थपनी}\ A}}\skt{ चैव विशल्यघ्नानि निर्द्दिशेत्~॥}\pend

\pstart\versepar \ledsidenote{\passref{3.6.12}}\edtext{\skt{लोहिताक्षश्च}}{\lemma{\skt{लोहिताक्षश्च}}\Afootnote{H; \skt{लोहिताक्षाणि}\ A}}\skt{ }\edtext{\skt{जानू\-र्व्वी\-कु\-र्च्चा\-चि\-टि\-प\-कु\-र्प्परा}}{\lemma{\skt{जानू\appbrk{}र्व्वी\appbrk{}कु\appbrk{}र्च्चा\appbrk{}चि\appbrk{}टि\appbrk{}प\appbrk{}कु\appbrk{}र्प्परा}}\Afootnote{H; \skt{जानू\appbrk{}र्वी\appbrk{}कू\appbrk{}र्च\appbrk{}वि\appbrk{}ट\appbrk{}प\appbrk{}कू\appbrk{}र्पराः}\ A}}\skt{~। }\pend
\pstart\versepar \skt{कुकुन्दरे कक्षधरे विधुरे सकृकाटिके~॥}\pend

\pstart\versepar \ledsidenote{\passref{3.6.13}}\edtext{\skt{अन्सा\-न्स\-फ\-ल\-का\-पाङ्गे}}{\lemma{\skt{अन्सा\appbrk{}न्स\appbrk{}फ\appbrk{}ल\appbrk{}का\appbrk{}पाङ्गे}}\Afootnote{H; \skt{अंसां\appbrk{}स\appbrk{}फ\appbrk{}ल\appbrk{}का\appbrk{}पाङ्गा}\ A}}\skt{ नीले मन्ये }\edtext{\skt{फणे}}{\lemma{\skt{फणे}}\Afootnote{H; \skt{फणौ}\ A}}\skt{ तथा~। }\pend
\pstart\versepar \skt{वैक\-ल्य\-क\-र\-णा\-न्या\-हु\-रा\-वर्त्तौ द्वौ तथैव च~॥}\pend

\pstart\versepar \ledsidenote{\passref{3.6.14}}\skt{गुल्फौ द्वौ मणिबन्धौ द्वौ द्वे द्वे }\edtext{\skt{कुर्च्चशिरांसि च}}{\lemma{\skt{कुर्च्चशिरांसि च}}\Afootnote{H; \skt{कूर्चशिरांसि च}\ A}}\skt{~। }\pend
\pstart\versepar \skt{रुजाकराणि जानी\-या\-द\-ष्टा\-वे\-तानि बुद्धिमान्~॥ }\pend
\pstart\versepar \skt{}\edtext{\skt{पञ्चोक्तानि}}{\lemma{\skt{पञ्चोक्तानि}}\Afootnote{H; \emph{om.}\ A}}\skt{ }\edtext{\skt{विकल्पानि}}{\lemma{\skt{विकल्पानि}}\Afootnote{H; \emph{om.}\ A}}\skt{ }\edtext{\skt{मर्म्मणां}}{\lemma{\skt{मर्म्मणां}}\Afootnote{H; \emph{om.}\ A}}\skt{ }\edtext{\edtext{\skt{कर्म्मणाम्पृथक्}}{\lemma{\skt{कर्म्मणाम्}}\Afootnote{H; \emph{om.}\ A}}}{\lemma{\skt{पृथक्}}\Afootnote{H; \skt{क्षिप्राणि विद्धमात्राणि घ्नन्ति कालान्तरेण च}\ A}}\skt{~॥}\pend

\pstart \ledsidenote{\passref{3.6.15}}\skt{मर्म्माणि }\edtext{\skt{नाम}}{\lemma{\skt{नाम}}\Afootnote{H; \emph{om.}\ A}}\skt{ मांस\-सि\-रा\-स्ना\-य्व\-स्थि\-स\-न्धि\-स\-न्नि\-पाताः~। तेषु स्वभावत एव }\edtext{\skt{प्राणास्तिष्ठन्ति}}{\lemma{\skt{प्राणास्}}\Afootnote{H; \skt{विशेषेण प्राणास्}\ A}}\skt{~। }\edtext{\skt{मर्म्म\-स्व\-भि\-ह\-ता\-स्तां\-स्ता\-न्भा\-वा\-ना\-प\-द्यन्ते}}{\lemma{\skt{मर्म्मस्व्}}\Afootnote{H; \skt{तस्मान्मर्मस्व्}\ A}}\skt{~॥}\pend

\pstart \ledsidenote{\passref{3.6.16}}\skt{तत्र }\edtext{\skt{सद्यः\-प्रा\-ण\-ह\-राणि}}{\lemma{\skt{सद्यः\appbrk{}प्रा\appbrk{}ण\appbrk{}ह\appbrk{}राणि}}\Afootnote{H; \skt{सद्यः\appbrk{}प्रा\appbrk{}ण\appbrk{}ह\appbrk{}रा\appbrk{}ण्या\appbrk{}ग्ने\appbrk{}यानि, अग्निगुणेष्वाशु क्षीणेषु क्षपयन्ति; काला\appbrk{}न्त\appbrk{}र\appbrk{}प्रा\appbrk{}ण\appbrk{}ह\appbrk{}राणि}\ A}}\skt{~। }\edtext{\skt{आग्नेयानि}}{\lemma{\skt{आग्नेयानि}}\Afootnote{H; \skt{सौम्याग्नेयानि}\ A}}\skt{~। }\edtext{\skt{अग्निगुणेषु}}{\lemma{\skt{अग्निगुणेषु}}\Afootnote{H; \skt{अग्निगुणेष्वाशु क्षीणेषु}\ A}}\skt{ }\edtext{\edtext{\skt{क्षीणेष्वायुः}}{\lemma{\skt{क्षीणेष्व्}}\Afootnote{H; \skt{क्रमेण च सोमगुणेषु}\ A}}}{\lemma{\skt{आयुः}}\Afootnote{H; \skt{कालान्}\ A}}\skt{ }\edtext{\skt{क्षपयन्ति}}{\lemma{\skt{क्षपयन्ति}}\Afootnote{H; \skt{तरेण क्षपयन्ति}\ A}}\skt{~। }\edtext{\skt{विशल्यघ्नानि}}{\lemma{\skt{विशल्यघ्नानि}}\Afootnote{H; \skt{विश\appbrk{}ल्य\appbrk{}प्रा\appbrk{}ण\appbrk{}ह\appbrk{}राणि}\ A}}\skt{ }\edtext{\skt{वायवीयानि}}{\lemma{\skt{वायवीयानि}}\Afootnote{H; \skt{वायव्यानि}\ A}}\skt{ }\edtext{\skt{शस्त्र\-मु\-खा\-व\-रुद्धो}}{\lemma{\skt{शस्त्र\appbrk{}मु\appbrk{}खा\appbrk{}व\appbrk{}रुद्धो}}\Afootnote{H; \skt{शल्य\appbrk{}मु\appbrk{}खा\appbrk{}व\appbrk{}रुद्धो}\ A}}\skt{ }\edtext{\skt{यावदत्र}}{\lemma{\skt{यावदत्र}}\Afootnote{H; \skt{यावदन्तर्}\ A}}\skt{ वायुस्तिष्ठति }\edtext{\skt{ताव\-ज्जी\-व\-त्यु\-द्धृ\-त\-मात्रे}}{\lemma{\skt{तावज्}}\Afootnote{H; \skt{ताव}\ A}}\skt{ तु शल्ये~। }\edtext{\skt{मर्म्मस्थानाश्रयो}}{\lemma{\skt{मर्म्मस्थानाश्रयो}}\Afootnote{H; \skt{मर्मस्थानाश्रितो}\ A}}\skt{ }\edtext{\edtext{\skt{वायुर्न्निष्क्रामति}}{\lemma{\skt{वायुर्}}\Afootnote{H; \skt{वायु}\ A}}}{\lemma{\skt{न्निष्क्रामति}}\Afootnote{H; \skt{र्निष्क्रामति}\ A}}\skt{~। तस्मात्सशल्यो }\edtext{\edtext{\skt{जीव\-त्यु\-द्धृ\-त\-शल्यो}}{\lemma{\skt{जीवत्य्}}\Afootnote{H; \skt{जीवत्}\ A}}}{\lemma{\skt{उद्धृतशल्यो}}\Afootnote{H; \skt{युद्धृतशल्यो}\ A}}\skt{ म्रियते~। }\edtext{\skt{वैकल्यकराणि}}{\lemma{\skt{वैकल्यकराणि}}\Afootnote{H; \skt{(पा\appbrk{}का\appbrk{}त्प\appbrk{}ति\appbrk{}त\appbrk{}शल्यो वा जीवति); वैकल्यकराणि}\ A}}\skt{~। सौम्यानि~। सोमो हि }\edtext{\edtext{\skt{शैत्यात्स्थिरत्वाच्च}}{\lemma{\skt{शैत्यात्}}\Afootnote{H; \emph{om.}\ A}}}{\lemma{\skt{स्थिरत्वाच्च}}\Afootnote{H; \skt{स्थिरत्वाच्छैत्याच्च}\ A}}\skt{ प्राणावलम्बनं करोति~। }\edtext{\skt{रुजाकराणि}}{\lemma{\skt{रुजाकराणि}}\Afootnote{H; \skt{रुजाकराण्}\ A}}\skt{~। }\edtext{\skt{अग्नि\-मा\-रु\-त\-भू\-यि\-ष्ठानि}}{\lemma{\skt{अग्नि\appbrk{}मा\appbrk{}रु\appbrk{}त\appbrk{}भू\appbrk{}यि\appbrk{}ष्ठानि}}\Afootnote{H; \skt{यग्नि\appbrk{}वा\appbrk{}यु\appbrk{}गु\appbrk{}ण\appbrk{}भू\appbrk{}यि\appbrk{}ष्ठानि}\ A}}\skt{ }\edtext{\skt{विशेषेण तौ}}{\lemma{\skt{विशेषेण तौ}}\Afootnote{H; \skt{विशेषतश्च तौ}\ A}}\skt{ रुजाकरौ पाञ्चभौतिकीञ्च }\edtext{\skt{रुजामपरे}}{\lemma{\skt{अपरे}}\Afootnote{H; \emph{om.}\ A}}\skt{ }\edtext{\skt{त्वाहुः}}{\lemma{\skt{त्वाहुः}}\Afootnote{H; \skt{आहुरेके}\ A}}\skt{~।}\pend

\pstart \ledsidenote{\passref{3.6.17}}\edtext{\skt{केचि\-दा\-हु\-र्म्मां\-सा\-दीनां}}{\lemma{\skt{म्मांसादीनां}}\Afootnote{H; \skt{मांसादीनां पञ्चानाम्}\ A}}\skt{ }\edtext{\skt{समस्तानां}}{\lemma{\skt{समस्तानां}}\Afootnote{H; \skt{अपि समस्तानां}\ A}}\skt{ }\edtext{\skt{समृद्धानाञ्च}}{\lemma{\skt{समृद्धानाञ्च}}\Afootnote{H; \skt{विवृद्धानां}\ A}}\skt{ समवायात्~। सद्यः\-प्रा\-ण\-ह\-राणि~। }\edtext{\skt{एक\-ही\-ना\-ना\-म\-ल्पा\-नाञ्च}}{\lemma{\skt{अल्पानाञ्च}}\Afootnote{H; \skt{अल्पानां}\ A}}\skt{ }\edtext{\skt{सम\-वा\-य\-का\-ला\-न्त\-र\-प्रा\-ण\-ह\-राणि}}{\lemma{\skt{सम\appbrk{}वा\appbrk{}य\appbrk{}का\appbrk{}ला\appbrk{}न्त\appbrk{}र\appbrk{}प्रा\appbrk{}ण\appbrk{}ह\appbrk{}राणि}}\Afootnote{H; \skt{वा काला\appbrk{}न्त\appbrk{}र\appbrk{}प्रा\appbrk{}ण\appbrk{}ह\appbrk{}राणि}\ A}}\skt{~। }\edtext{\skt{द्विहीनानां तु}}{\lemma{\skt{द्विहीनानां तु}}\Afootnote{H; \skt{द्विहीनानां}\ A}}\skt{ विश\-ल्य\-प्रा\-ण\-ह\-राणि~। }\edtext{\skt{त्रिहीनानान्तु}}{\lemma{\skt{त्रिहीनानान्तु}}\Afootnote{H; \skt{त्रिहीनानां}\ A}}\skt{ }\edtext{\skt{वैगुण्यकराणि}}{\lemma{\skt{वैगुण्यकराणि}}\Afootnote{H; \skt{वैकल्यकराणि}\ A}}\skt{~। एकस्मिन्नेव रुजाकराणीति~॥ }\edtext{\edtext{\skt{यतश्चैवमतो}}{\lemma{\skt{यतश्}}\Afootnote{H; \emph{om.}\ A}}}{\lemma{\skt{चैवमतो}}\Afootnote{H; \skt{नैवं, यतो}\ A}}\skt{ }\edtext{\edtext{\skt{स्थिमर्म्मस्वपि}}{\lemma{\skt{स्थिमर्म्मस्व्}}\Afootnote{H; \skt{ऽस्थिमर्मस्}\ A}}}{\lemma{\skt{अपि}}\Afootnote{H; \skt{वप्}\ A}}\skt{ }\edtext{\skt{क्षतेषु}}{\lemma{\skt{क्षतेषु}}\Afootnote{H; \skt{यभिहतेषु}\ A}}\skt{ शोणि\-ता\-ग\-म\-न\-म्भ\-वति~॥}\pend

\pstart\versepar \ledsidenote{\passref{3.6.18}}\edtext{\skt{भवन्ति}}{\lemma{\skt{भवन्ति}}\Afootnote{H; \emph{om.}\ A}}\skt{ }\edtext{\skt{चात्र}}{\lemma{\skt{चात्र}}\Afootnote{H; \emph{om.}\ A}}\skt{~॥ }\pend
\pstart\versepar \skt{चतुर्व्विधा यास्तु सिराः शरीरे प्रायेण ता मर्म्मसु सन्निविष्टाः~। }\pend
\pstart\versepar \skt{स्नाय्वस्थिमांसानि तथैव }\edtext{\skt{सन्धी}}{\lemma{\skt{सन्धी}}\Afootnote{H; \skt{सन्धीन्}\ A}}\skt{ }\edtext{\skt{सन्त\-र्प्प\-य\-न्त्यो\-प\-च\-यन्ति}}{\lemma{\skt{सन्त\appbrk{}र्प्प\appbrk{}य\appbrk{}न्त्यो\appbrk{}प\appbrk{}च\appbrk{}यन्ति}}\Afootnote{H; \skt{सन्तर्प्य}\ A}}\skt{ }\edtext{\skt{देहं}}{\lemma{\skt{देहं}}\Afootnote{H; \skt{देहं प्रति यापयन्ति}\ A}}\skt{~॥}\pend

\pstart \ledsidenote{\passref{3.6.19}}\skt{ततः क्षते मर्म्मणि ताः प्रवृद्धः समन्ततो वायु\-र\-भि\-स्तृ\-णोति~। विवर्द्धमानस्तु स मातरिश्वा }\edtext{\skt{रुजाः}}{\lemma{\skt{रुजाः}}\Afootnote{H; \skt{रुजः}\ A}}\skt{ सुतीव्राः प्रतनोति काये~॥}\pend

\pstart \ledsidenote{\passref{3.6.20}}\edtext{\skt{रुजाभिभूतस्तु}}{\lemma{\skt{रुजाभिभूतस्तु}}\Afootnote{H; \skt{रुजाभिभूतं तु}\ A}}\skt{ ततः शरीरं प्रलीयते नश्यति चास्य }\edtext{\skt{संज्ञा}}{\lemma{\skt{संज्ञा}}\Afootnote{H; \skt{संज्ञा~। अतो हि शल्यं विनि\appbrk{}ह\appbrk{}र्तु\appbrk{}मि\appbrk{}च्छ\appbrk{}न्म\appbrk{}र्माणि यत्नेन परीक्ष्य कर्षेत्}\ A}}\skt{~॥}\pend

\pstart \ledsidenote{\passref{3.6.21}}\skt{एतेन शेषं व्याख्यातं~॥}\pend

\pstart \ledsidenote{\passref{3.6.22}}\skt{तत्र सद्यः\-प्रा\-ण\-ह\-र\-मन्ते विद्धं कालान्तरेण मारयति~। काला\-न्त\-र\-प्रा\-ण\-ह\-र\-मन्ते विद्धं }\edtext{\edtext{\skt{वैक\-ल्य\-क\-र\-म्भ\-वति}}{\lemma{\skt{वैकल्यकरम्}}\Afootnote{H; \skt{वैकल्यमापा}\ A}}}{\lemma{\skt{भवति}}\Afootnote{H; \skt{दयति}\ A}}\skt{~। }\edtext{\skt{विशल्यघ्नञ्च}}{\lemma{\skt{विशल्यघ्नञ्च}}\Afootnote{H; \skt{विशल्यघ्नं}\ A}}\skt{ }\edtext{\skt{वैक\-ल्य\-क\-र\-म्भ\-वति}}{\lemma{\skt{भवति}}\Afootnote{H; \emph{om.}\ A}}\skt{~। }\edtext{\skt{नानातीव्रां}}{\lemma{\skt{नानातीव्रां}}\Afootnote{H; \emph{om.}\ A}}\skt{ }\edtext{\edtext{\skt{वेद\-ना\-ञ्ज\-न\-यंति}}{\lemma{\skt{वेदनाञ्}}\Afootnote{H; \skt{च}\ A}}}{\lemma{\skt{जनयंति}}\Afootnote{H; \skt{भवति}\ A}}\skt{~। वैकल्यकरं कालान्तरेण क्लेशयति~। }\edtext{\skt{रुजाम्वा}}{\lemma{\skt{रुजाम्वा}}\Afootnote{H; \skt{रुजां च}\ A}}\skt{ करोति रुजाकरं }\edtext{\skt{चाल्पवेदनां}}{\lemma{\skt{चाल्पवेदनां}}\Afootnote{H; \skt{अतीव्रवेदनं}\ A}}\skt{ }\edtext{\skt{करोति}}{\lemma{\skt{करोति}}\Afootnote{H; \skt{भवति}\ A}}\skt{~।}\pend

\pstart \ledsidenote{\passref{3.6.23}}\skt{तत्र सद्यः\-प्रा\-ण\-ह\-राणि }\edtext{\skt{सप्त\-रा\-त्रा\-भ्य\-न्त\-रा\-त्मा\-र\-यन्ति}}{\lemma{\skt{सप्त\appbrk{}रा\appbrk{}त्रा\appbrk{}भ्य\appbrk{}न्त\appbrk{}रात्}}\Afootnote{H; \skt{सप्त\appbrk{}रा\appbrk{}त्रा\appbrk{}भ्य\appbrk{}न्त\appbrk{}रान्}\ A}}\skt{~। काला\-न्त\-र\-प्रा\-ण\-ह\-राणि }\edtext{\skt{पक्षात्मासाद्वा}}{\lemma{\skt{पक्षात्}}\Afootnote{H; \skt{पक्षान्}\ A}}\skt{~। तेष्वपि तु क्षिप्राणि कदाचिदाशु मारयन्ति~। विश\-ल्य\-प्रा\-ण\-ह\-राणि वैकल्यकराणि च कदा\-चि\-द\-त्य\-भि\-ह\-तानि मारयन्ति~॥}\pend

\pstart \ledsidenote{\passref{3.6.24}}\edtext{\skt{उत}}{\lemma{\skt{उत}}\Afootnote{H; \skt{अत}\ A}}\skt{ ऊर्ध्वं }\edtext{\skt{सक्थि\-म\-र्म्मा\-ण्य\-नु\-व्या\-ख्या\-स्यामः}}{\lemma{\skt{अनुव्याख्यास्यामः}}\Afootnote{H; \skt{व्याख्यास्यामः}\ A}}\skt{~॥ तत्र पाद\-स्या\-ङ्गु\-ष्ठा\-ङ्गु\-ल्यो\-र्म्मध्ये क्षिप्रं नाम }\edtext{\skt{मर्म}}{\lemma{\skt{मर्म}}\Afootnote{H; \skt{मर्म, तत्र विद्ध\appbrk{}स्या\appbrk{}क्षे\appbrk{}प\appbrk{}केण म}\ A}}\skt{ … }\edtext{\edtext{\skt{मध्य\-मा\-ङ्गु\-लि\-म\-नु\-पू\-र्व्वेण}}{\lemma{\skt{{\symfont ‡}}}\Afootnote{H; \skt{रणं}\ A}}}{\lemma{\skt{मध्यमाङ्गुलिम्}}\Afootnote{H; \skt{मध्यमाङ्गुलीम्}\ A}}\skt{ मध्ये पादतलस्य }\edtext{\skt{तल\-हृ\-द\-य\-न्नाम मर्म्म}}{\lemma{\skt{नाम}}\Afootnote{H; \skt{ना}\ A}}\skt{ तत्र }\edtext{\skt{रुजायाम्मरणं}}{\lemma{\skt{रुजायाम्}}\Afootnote{H; \skt{रुजाभिर्}\ A}}\skt{~। }\edtext{\edtext{\skt{क्षिप्र\-त्यो\-प\-रि\-ष्टा\-द\-ङ्गु\-ल\-मु\-भ\-यतः}}{\lemma{\skt{क्षिप्र\appbrk{}त्यो\appbrk{}प\appbrk{}रि\appbrk{}ष्टा\appbrk{}द\appbrk{}ङ्गु\appbrk{}लम्}}\Afootnote{H; \skt{क्षिप्रस्योपरिष्टा}\ A}}\skt{ कूर्च्चो}}{\lemma{\skt{उभयतः}}\Afootnote{H; \skt{दुभयतः}\ A}}\skt{ नाम~। }\edtext{\skt{…णवेपने भवतः}}{\lemma{\skt{{\symfont ‡}णवेपने}}\Afootnote{H; \skt{तत्र पादस्य भ्रमणवेपने}\ A}}\skt{~। }\edtext{\skt{गुल्फसन्धेरधः कूर्च्चशिरो}}{\lemma{\skt{अधः}}\Afootnote{H; \skt{अध उभयतः}\ A}}\skt{ }\edtext{\skt{नाम तत्र}}{\lemma{\skt{नाम}}\Afootnote{H; \emph{om.}\ A}}\skt{ रुजाशोफौ }\edtext{\skt{भवतः~॥}}{\lemma{\skt{भवतः}}\Afootnote{H; \emph{om.}\ A}}\skt{ }\edtext{\skt{पादजङ्घयोस्तु सन्धाने}}{\lemma{\skt{पादजङ्घयोस्तु}}\Afootnote{H; \skt{पादजङ्घयोः}\ A}}\skt{ गुल्फो }\edtext{\skt{नाम तत्र}}{\lemma{\skt{नाम}}\Afootnote{H; \emph{om.}\ A}}\skt{ }\edtext{\skt{रुक\-स्त\-ब्ध\-स\-क्थिता च कुणता}}{\lemma{\skt{रुक\appbrk{}स्त\appbrk{}ब्ध\appbrk{}स\appbrk{}क्थिता च}}\Afootnote{H; \skt{रुजः स्तब्धपाद}\ A}}\skt{ वा… }\edtext{\skt{पार्ष्णिम्प्रति}}{\lemma{\skt{वा{\symfont ‡}}}\Afootnote{H; \skt{वा}\ A}}\skt{ }\edtext{\edtext{\skt{जङ्घा\-म\-ध्ये\-न्द्र\-ब\-स्ति\-र्न्नाम}}{\lemma{\skt{जङ्घा\appbrk{}म\appbrk{}ध्ये\appbrk{}न्द्र\appbrk{}ब\appbrk{}स्तिर्}}\Afootnote{H; \skt{जङ्घामध्ये इन्द्रबस्तिः}\ A}}\skt{ मर्म्म}}{\lemma{\skt{न्नाम}}\Afootnote{H; \emph{om.}\ A}}\skt{ तत्र }\edtext{\skt{शोणितक्षये मरणं}}{\lemma{\skt{शोणितक्षये}}\Afootnote{H; \skt{शोणितक्षयेण}\ A}}\skt{~। जङ्घोर्व्वोः सन्धाने }\edtext{\edtext{\skt{जानुर्न्नाम}}{\lemma{\skt{जानुर्}}\Afootnote{H; \skt{जा}\ A}}\skt{ तत्र}}{\lemma{\skt{न्नाम}}\Afootnote{H; \skt{नु}\ A}}\skt{ खञ्जता }\edtext{\skt{भवति~। }\edtext{\edtext{\skt{जानो\-स्त्र्य\-ङ्गु\-ला\-दू\-र्द्ध्व…\-णि\-र्न्नाम}}{\lemma{\skt{जानोस्}}\Afootnote{H; \skt{तस्}\ A}}}{\lemma{\skt{त्र्यङ्गुलादूर्द्ध्व{\symfont ‡}णिर्}}\Afootnote{H; \skt{त्र्यङ्गु}\ A}}}{\lemma{\skt{भवति}}\Afootnote{H; \skt{जानुन ऊर्ध्वमुभय}\ A}}\skt{ }\edtext{\skt{मर्म्म तत्र}}{\lemma{\skt{मर्म्म}}\Afootnote{H; \skt{माणी}\ A}}\skt{ शोफाभिवृद्धिः स्तब्धसक्थिता च~। }\edtext{\skt{ऊरूमध्ये ऊर्व्वी}}{\lemma{\skt{ऊरूमध्ये}}\Afootnote{H; \skt{ऊरुमध्ये}\ A}}\skt{ नाम }\edtext{\edtext{\skt{मर्म्म}}{\lemma{\skt{नाम}}\Afootnote{H; \emph{om.}\ A}}\skt{ तत्र}}{\lemma{\skt{मर्म्म}}\Afootnote{H; \emph{om.}\ A}}\skt{ शोणि\-त\-क्ष\-या\-त्स\-क्थि\-शोषः~। }\edtext{\edtext{\skt{ऊर्व्व्यास्तूर्ध्वमधो}}{\lemma{\skt{ऊर्व्व्यास्}}\Afootnote{H; \skt{उर्व्या}\ A}}}{\lemma{\skt{तूर्ध्वम्}}\Afootnote{H; \skt{ऊर्ध्वम्}\ A}}\skt{ }\edtext{\skt{वंक्षणस…ले लोहिताक्षन्नाम}}{\lemma{\skt{वंक्षणस{\symfont ‡}ले}}\Afootnote{H; \skt{वङ्क्ष\appbrk{}ण\appbrk{}स\appbrk{}न्धे\appbrk{}रू\appbrk{}रु\appbrk{}मूले}\ A}}\skt{ }\edtext{\edtext{\skt{मर्म्म}}{\lemma{\skt{नाम}}\Afootnote{H; \emph{om.}\ A}}\skt{ तत्र}}{\lemma{\skt{मर्म्म}}\Afootnote{H; \emph{om.}\ A}}\skt{ }\edtext{\skt{शोणितक्षये मर\-ण\-म्प\-क्षा\-घातो}}{\lemma{\skt{शोणितक्षये}}\Afootnote{H; \skt{लोहितक्षयेण}\ A}}\skt{ वा~॥ वंक्ष\-ण\-वृ\-ष\-ण\-यो\-र\-न्तरे }\edtext{\skt{चिटिपन्नाम }\edtext{\skt{मर्म्म}}{\lemma{\skt{नाम}}\Afootnote{H; \emph{om.}\ A}}}{\lemma{\skt{चिटिपन्}}\Afootnote{H; \skt{विटपं}\ A}}\skt{ तत्र }\edtext{\edtext{\skt{सांढ्य\-म\-ल्प\-शु\-क्रता}}{\lemma{\skt{सांढ्यम्}}\Afootnote{H; \skt{षाण्ढ्यम्}\ A}}\skt{ च भवति}}{\lemma{\skt{अल्पशुक्रता च}}\Afootnote{H; \skt{अल्पशुक्रता वा}\ A}}\skt{~। एव\-मे\-ता\-न्ये\-का\-दश सक्थिमर्म्माणि व्याख्यातानि~। एते\-ने\-त\-र\-सक्थि बाहू च }\edtext{\skt{व्याख्याताः~॥ विशेषतस्तु}}{\lemma{\skt{व्याख्याताः}}\Afootnote{H; \skt{व्याख्यातौ}\ A}}\skt{ यानि }\edtext{\skt{सक्थिनि गुल्फ\-जा\-नु\-वं\-क्ष\-ण\-चि\-टि\-पानि}}{\lemma{\skt{सक्थिनि}}\Afootnote{H; \skt{सक्थ्नि}\ A}}\skt{ तानि बाहौ }\edtext{\skt{मणि\-ब\-न्ध\-कु\-र्प्प\-र\-क\-क्ष\-ध\-राणि भवन्ति}}{\lemma{\skt{मणि\appbrk{}ब\appbrk{}न्ध\appbrk{}कु\appbrk{}र्प्प\appbrk{}र\appbrk{}क\appbrk{}क्ष\appbrk{}ध\appbrk{}राणि}}\Afootnote{H; \skt{मणि\appbrk{}ब\appbrk{}न्ध\appbrk{}कू\appbrk{}र्प\appbrk{}र\appbrk{}क\appbrk{}क्ष\appbrk{}धरा}\ A}}\skt{~। यथा }\edtext{\skt{च पादजङ्घयोस्तु}}{\lemma{\skt{यथा च}}\Afootnote{H; \emph{om.}\ A}}\skt{ }\edtext{\skt{सन्धौ गुल्फौ}}{\lemma{\skt{सन्धौ}}\Afootnote{H; \emph{om.}\ A}}\skt{ नाम }\edtext{\edtext{\skt{मर्म्म}}{\lemma{\skt{नाम}}\Afootnote{H; \emph{om.}\ A}}\skt{~। एवं}}{\lemma{\skt{मर्म्म}}\Afootnote{H; \emph{om.}\ A}}\skt{ }\edtext{\edtext{\skt{हस्त\-त\-ल\-प्र\-को\-ष्ठयोः}}{\lemma{\skt{एवं}}\Afootnote{H; \emph{om.}\ A}}\skt{ सन्धौ}}{\lemma{\skt{हस्त\appbrk{}त\appbrk{}ल\appbrk{}प्र\appbrk{}को\appbrk{}ष्ठयोः}}\Afootnote{H; \emph{om.}\ A}}\skt{ }\edtext{\skt{मणिबन्धो नाम }\edtext{\skt{मर्म्म}}{\lemma{\skt{नाम}}\Afootnote{H; \emph{om.}\ A}}}{\lemma{\skt{मणिबन्धो}}\Afootnote{H; \skt{णि}\ A}}\skt{~। यथा }\edtext{\skt{च जङ्घोर्व्वोः}}{\lemma{\skt{यथा च}}\Afootnote{H; \emph{om.}\ A}}\skt{ }\edtext{\skt{सन्धाने जानु नाम}}{\lemma{\skt{सन्धाने}}\Afootnote{H; \emph{om.}\ A}}\skt{ }\edtext{\edtext{\skt{मर्म्म}}{\lemma{\skt{नाम}}\Afootnote{H; \emph{om.}\ A}}\skt{~। एवं}}{\lemma{\skt{मर्म्म}}\Afootnote{H; \emph{om.}\ A}}\skt{ }\edtext{\edtext{\skt{प्रहस्तबाह्वोः}}{\lemma{\skt{एवं}}\Afootnote{H; \emph{om.}\ A}}\skt{ सन्धाने}}{\lemma{\skt{प्रहस्तबाह्वोः}}\Afootnote{H; \emph{om.}\ A}}\skt{ }\edtext{\skt{कुर्प्परो नाम }\edtext{\skt{मर्म्म}}{\lemma{\skt{नाम}}\Afootnote{H; \emph{om.}\ A}}}{\lemma{\skt{कुर्प्परो}}\Afootnote{H; \emph{om.}\ A}}\skt{~। यथा }\edtext{\skt{च वंक्ष\-ण\-वृ\-ष\-ण\-यो\-र्म्मध्ये}}{\lemma{\skt{यथा च}}\Afootnote{H; \skt{यथा}\ A}}\skt{ }\edtext{\skt{चिटिपन्नाम }\edtext{\skt{मर्म्म}}{\lemma{\skt{नाम}}\Afootnote{H; \emph{om.}\ A}}}{\lemma{\skt{चिटिपन्}}\Afootnote{H; \skt{विटपम्}\ A}}\skt{~। }\edtext{\skt{एव\-म\-क्ष\-क\-क्ष\-यो\-र्म्मध्ये}}{\lemma{\skt{अक्षकक्षयोर्}}\Afootnote{H; \skt{वक्षःकक्षयोर्}\ A}}\skt{ }\edtext{\skt{कक्षधरो नाम }\edtext{\skt{मर्म्म}}{\lemma{\skt{नाम}}\Afootnote{H; \emph{om.}\ A}}}{\lemma{\skt{कक्षधरो}}\Afootnote{H; \skt{कक्षधरं}\ A}}\skt{~। तत्र }\edtext{\edtext{\skt{विद्धेषु}}{\lemma{\skt{तत्र}}\Afootnote{H; \skt{तस्मिन्}\ A}}\skt{ त एवोपद्रवा}}{\lemma{\skt{विद्धेषु त}}\Afootnote{H; \skt{विद्धे त}\ A}}\skt{ }\edtext{\skt{भवन्ति~। विशेषतस्तु}}{\lemma{\skt{भवन्ति}}\Afootnote{H; \emph{om.}\ A}}\skt{ मणिबन्धे कुण्ठता~। }\edtext{\skt{कुर्प्परे पक्षाघातः}}{\lemma{\skt{कुर्प्परे}}\Afootnote{H; \skt{कूर्पराख्ये कुणिः, कक्षधरे}\ A}}\skt{~। }\edtext{\skt{कुणित्वं कक्षधरे}}{\lemma{\skt{कुणित्वं}}\Afootnote{H; \emph{om.}\ A}}\skt{ चेति~। }\edtext{\skt{एवमेतानि}}{\lemma{\skt{चेति}}\Afootnote{H; \skt{तः}\ A}}\skt{ }\edtext{\skt{चतु\-श्च\-त्वा\-रिं\-श\-च्छा\-खा\-म\-र्म्माणि व्याख्यातानि}}{\lemma{\skt{छाखामर्म्माणि}}\Afootnote{H; \skt{छाखासु मर्माणि}\ A}}\skt{ }\edtext{\skt{भवन्ति~॥}}{\lemma{\skt{भवन्ति}}\Afootnote{H; \emph{om.}\ A}}\edtext{}{\lemma{\skt{मर्म्म}}\Afootnote{H; \skt{म}\ A}}\edtext{}{\lemma{\skt{कुणता}}\Afootnote{H; \skt{ता खञ्जता}\ A}}\edtext{}{\lemma{\skt{मर्म्म}}\Afootnote{H; \emph{om.}\ A}}\edtext{}{\lemma{\skt{न्नाम}}\Afootnote{H; \skt{ल}\ A}}\edtext{}{\lemma{\skt{ऊर्व्वी}}\Afootnote{H; \skt{उर्वी}\ A}}\edtext{}{\lemma{\skt{मर्म्म}}\Afootnote{H; \emph{om.}\ A}}\edtext{}{\lemma{\skt{गुल्फ\appbrk{}जा\appbrk{}नु\appbrk{}वं\appbrk{}क्ष\appbrk{}ण\appbrk{}चि\appbrk{}टि\appbrk{}पानि}}\Afootnote{H; \skt{गुल्फ\appbrk{}जा\appbrk{}नु\appbrk{}वि\appbrk{}ट\appbrk{}पानि}\ A}}\edtext{}{\lemma{\skt{भवन्ति}}\Afootnote{H; \emph{om.}\ A}}\edtext{}{\lemma{\skt{पादजङ्घयोस्तु}}\Afootnote{H; \emph{om.}\ A}}\edtext{}{\lemma{\skt{गुल्फौ}}\Afootnote{H; \emph{om.}\ A}}\edtext{}{\lemma{\skt{सन्धौ}}\Afootnote{H; \emph{om.}\ A}}\edtext{}{\lemma{\skt{मर्म्म}}\Afootnote{H; \emph{om.}\ A}}\edtext{}{\lemma{\skt{जङ्घोर्व्वोः}}\Afootnote{H; \emph{om.}\ A}}\edtext{}{\lemma{\skt{जानु}}\Afootnote{H; \emph{om.}\ A}}\edtext{}{\lemma{\skt{सन्धाने}}\Afootnote{H; \emph{om.}\ A}}\edtext{}{\lemma{\skt{मर्म्म}}\Afootnote{H; \emph{om.}\ A}}\edtext{}{\lemma{\skt{म्मध्ये}}\Afootnote{H; \skt{अन्तरे}\ A}}\edtext{}{\lemma{\skt{मर्म्म}}\Afootnote{H; \emph{om.}\ A}}\edtext{}{\lemma{\skt{मर्म्म}}\Afootnote{H; \emph{om.}\ A}}\edtext{}{\lemma{\skt{पक्षाघातः}}\Afootnote{H; \skt{पक्षाघा}\ A}}\edtext{}{\lemma{\skt{कक्षधरे}}\Afootnote{H; \emph{om.}\ A}}\pend

\pstart \ledsidenote{\passref{3.6.25}}\skt{अत }\edtext{\skt{ऊर्ध्व\-मु\-द\-रो\-र\-सो\-र्द्वा\-दश}}{\lemma{\skt{द्वादश}}\Afootnote{H; \emph{om.}\ A}}\skt{ मर्म्मा\-ण्य\-नु\-व्या\-ख्या\-स्यामः~॥ तत्र }\edtext{\skt{वात\-व\-र्च्चो\-निः\-स\-रणी}}{\lemma{\skt{वात\appbrk{}व\appbrk{}र्च्चो\appbrk{}निः\appbrk{}स\appbrk{}रणी}}\Afootnote{H; \skt{वात\appbrk{}व\appbrk{}र्चो\appbrk{}नि\appbrk{}र\appbrk{}सनं}\ A}}\skt{ }\edtext{\skt{स्थूला\-न्त्र\-प्र\-ति\-ब\-द्ध\-म्म\-हा\-गुदो}}{\lemma{\skt{महागुदो}}\Afootnote{H; \skt{गुदं}\ A}}\skt{ }\edtext{\skt{नाम}}{\lemma{\skt{नाम}}\Afootnote{H; \skt{नाम मर्म}\ A}}\skt{~। तत्र }\edtext{\skt{सद्यो}}{\lemma{\skt{सद्यो}}\Afootnote{H; \skt{सद्यो}\ A}}\skt{ मरणं~। }\edtext{\skt{अल्प\-मां\-स\-शो\-णि\-त\-म\-भ्य\-न्त\-रतः}}{\lemma{\skt{अल्प\appbrk{}मां\appbrk{}स\appbrk{}शो\appbrk{}णि\appbrk{}तम्}}\Afootnote{H; \skt{अल्प\appbrk{}मां\appbrk{}स\appbrk{}शो\appbrk{}णितो}\ A}}\skt{~। }\edtext{\edtext{\skt{कट्या\-मु\-त्रा\-स\-य\-यो\-र्ब\-स्ति\-र्न्नाम}}{\lemma{\skt{कट्या\appbrk{}मु\appbrk{}त्रा\appbrk{}स\appbrk{}य\appbrk{}योर्}}\Afootnote{H; \skt{कट्यां मूत्राशयो}\ A}}}{\lemma{\skt{न्नाम}}\Afootnote{H; \emph{om.}\ A}}\skt{ }\edtext{\skt{मर्म्म}}{\lemma{\skt{मर्म्म}}\Afootnote{H; \emph{om.}\ A}}\skt{~। तत्रापि }\edtext{\skt{सद्यो}}{\lemma{\skt{सद्यो}}\Afootnote{H; \skt{सद्यो}\ A}}\skt{ मरणं~। अश्म\-री\-व्र\-णा\-दृते तत्राप्युभयतो }\edtext{\skt{विद्धे}}{\lemma{\skt{विद्धे}}\Afootnote{H; \skt{भिन्ने}\ A}}\skt{ }\edtext{\skt{मरणमेकतो}}{\lemma{\skt{मरणम्}}\Afootnote{H; \skt{न जीवति}\ A}}\skt{ भिन्ने }\edtext{\skt{मूत्रास्रावी}}{\lemma{\skt{मूत्रास्रावी}}\Afootnote{H; \skt{मूत्रस्रावी}\ A}}\skt{ व्रणो भवति~। स तु यत्नेनोपक्रान्तो रोहति~॥ पक्वा\-मा\-श\-य\-यो\-र्म्मध्ये }\edtext{\skt{शिराप्रभावो}}{\lemma{\skt{शिराप्रभावो}}\Afootnote{H; \skt{सिराप्रभवा}\ A}}\skt{ }\edtext{\skt{नाभिर्न्नाम}}{\lemma{\skt{न्नाम}}\Afootnote{H; \emph{om.}\ A}}\skt{ }\edtext{\skt{मर्म्म}}{\lemma{\skt{मर्म्म}}\Afootnote{H; \emph{om.}\ A}}\skt{ तत्रापि सद्यो मरणं~। }\edtext{\edtext{\skt{स्तन\-यो\-र्म्म\-ध्य\-म\-धि\-ष्ठा\-यो\-रसि}}{\lemma{\skt{म्मध्यम्}}\Afootnote{H; \skt{मध्य}\ A}}}{\lemma{\skt{अधिष्ठायोरसि}}\Afootnote{H; \skt{मधिष्ठायोरस्}\ A}}\skt{ }\edtext{\skt{आमाशयद्वारं}}{\lemma{\skt{आमाशयद्वारं}}\Afootnote{H; \skt{यामाशयद्वारं}\ A}}\skt{ सत्त्व\-र\-ज\-स्त\-म\-सा\-म\-धि\-ष्ठानं }\edtext{\skt{हृदयन्नाम}}{\lemma{\skt{नाम}}\Afootnote{H; \emph{om.}\ A}}\skt{ }\edtext{\skt{मर्म्म}}{\lemma{\skt{मर्म्म}}\Afootnote{H; \emph{om.}\ A}}\skt{ तत्रापि }\edtext{\skt{सद्यो}}{\lemma{\skt{सद्यो}}\Afootnote{H; \skt{सद्य}\ A}}\skt{ }\edtext{\skt{मरणं}}{\lemma{\skt{मरणं}}\Afootnote{H; \skt{एव मरणं}\ A}}\skt{~। स्तन\-यो\-र\-ध\-स्ता\-द्द्व्य\-ङ्गु\-ल\-मु\-भ\-यतः स्तनमूले }\edtext{\skt{नाम}}{\lemma{\skt{नाम}}\Afootnote{H; \emph{om.}\ A}}\skt{ }\edtext{\skt{मर्मणी}}{\lemma{\skt{मर्मणी}}\Afootnote{H; \emph{om.}\ A}}\skt{~। तत्र कफ\-पू\-र्ण्ण\-को\-ष्ठ\-तया }\edtext{\skt{मरणं}}{\lemma{\skt{मरणं}}\Afootnote{H; \skt{(कासश्वासाभ्यां ) म्रियते}\ A}}\skt{~। }\edtext{\skt{स्तन\-चू\-चु\-क\-यो\-रूर्द्ध्वे}}{\lemma{\skt{ऊर्द्ध्वे}}\Afootnote{H; \skt{ऊर्ध्वं द्व्}\ A}}\skt{ }\edtext{\skt{स्तनरोहिते}}{\lemma{\skt{स्तनरोहिते}}\Afootnote{H; \skt{यङ्गु\appbrk{}ल\appbrk{}मु\appbrk{}भ\appbrk{}यतः स्तनरोहितौ}\ A}}\skt{ }\edtext{\skt{नाम}}{\lemma{\skt{नाम}}\Afootnote{H; \emph{om.}\ A}}\skt{ }\edtext{\skt{मर्म्मणी}}{\lemma{\skt{मर्म्मणी}}\Afootnote{H; \emph{om.}\ A}}\skt{~। तत्र }\edtext{\skt{रक्त\-पू\-र्ण्ण\-को\-ष्ठ\-तया}}{\lemma{\skt{रक्त\appbrk{}पू\appbrk{}र्ण्ण\appbrk{}को\appbrk{}ष्ठ\appbrk{}तया}}\Afootnote{H; \skt{लोहि\appbrk{}त\appbrk{}पू\appbrk{}र्ण\appbrk{}को\appbrk{}ष्ठ\appbrk{}तया}\ A}}\skt{ }\edtext{\skt{कास\-श्वा\-सा\-भ्या\-म्म\-रणं}}{\lemma{\skt{मरणं}}\Afootnote{H; \skt{च म्रियते}\ A}}\skt{~। }\edtext{\edtext{\edtext{\skt{अंस\-कू\-ट\-यो\-र\-ध\-स्ता\-त्पा\-र्श्व\-यो\-रु\-प\-रि\-भा\-ग\-यो\-र\-पा\-लापौ}}{\lemma{\skt{पार्श्वयोर्}}\Afootnote{H; \skt{पार्श्वो}\ A}}}{\lemma{\skt{उप\appbrk{}रि\appbrk{}भा\appbrk{}ग\appbrk{}योर्}}\Afootnote{H; \skt{परिभागयोर्}\ A}}}{\lemma{\skt{अपालापौ}}\Afootnote{H; \skt{अपलापौ}\ A}}\skt{ }\edtext{\skt{नाम}}{\lemma{\skt{नाम}}\Afootnote{H; \skt{ना}\ A}}\skt{ }\edtext{\skt{मर्म्मणी}}{\lemma{\skt{मर्म्मणी}}\Afootnote{H; \skt{म}\ A}}\skt{~। }\edtext{\skt{तत्र}}{\lemma{\skt{तत्र}}\Afootnote{H; \skt{तत्र र}\ A}}\skt{ }\edtext{\edtext{\skt{पूय\-भा\-व\-मा\-पन्ने}}{\lemma{\skt{पूयभावम्}}\Afootnote{H; \skt{क्तेन पूयभावं}\ A}}}{\lemma{\skt{आपन्ने}}\Afootnote{H; \emph{om.}\ A}}\skt{ }\edtext{\skt{शोणिते}}{\lemma{\skt{शोणिते}}\Afootnote{H; \skt{गतेन}\ A}}\skt{ मरणं~। उभयत्रोरसो }\edtext{\skt{नाड्यो}}{\lemma{\skt{नाड्यो}}\Afootnote{H; \skt{नाड्यौ}\ A}}\skt{ वातवहे~। अपस्तम्भौ }\edtext{\skt{नाम}}{\lemma{\skt{नाम}}\Afootnote{H; \skt{ना}\ A}}\skt{ }\edtext{\skt{मर्म्मणी}}{\lemma{\skt{मर्म्मणी}}\Afootnote{H; \skt{म}\ A}}\skt{~। तत्र वात\-पू\-र्ण्ण\-को\-ष्ठ\-तया }\edtext{\skt{कास\-श्वा\-सा\-भ्या\-म्म\-रणं}}{\lemma{\skt{कासश्वासाभ्याम्}}\Afootnote{H; \skt{कासश्वासाभ्यां च}\ A}}\skt{~। एव\-मे\-ता\-न्यु\-द\-रो\-र\-सो\-र्द्वा\-दश मर्म्माणि }\edtext{\skt{व्याख्यातानि}}{\lemma{\skt{व्याख्यातानि}}\Afootnote{A H\textsuperscript{pc}; \skt{व्याख्यानि}\ H\textsuperscript{ac}}}\skt{~॥}\pend

\pstart \ledsidenote{\passref{3.6.26}}\skt{अत }\edtext{\skt{ऊर्ध्व\-म्पृ\-ष्ठ\-म\-र्म्मा\-ण्य\-नु\-व्या\-ख्या\-स्यामः}}{\lemma{\skt{अनुव्याख्यास्यामः}}\Afootnote{H; \skt{व्याख्यास्यामः}\ A}}\skt{~॥ तत्र पृष्ठ\-वं\-श\-मु\-भ\-यतः }\edtext{\skt{प्रति\-श्रो\-णी\-क\-र्ण्णा\-व\-स्थिनी}}{\lemma{\skt{प्रति\appbrk{}श्रो\appbrk{}णी\appbrk{}क\appbrk{}र्ण्णा\appbrk{}व\appbrk{}स्थिनी}}\Afootnote{H; \skt{प्रति\appbrk{}श्रो\appbrk{}णि\appbrk{}का\appbrk{}ण्ड\appbrk{}म\appbrk{}स्थिनी}\ A}}\skt{ कटीकतरुणे }\edtext{\skt{नाम}}{\lemma{\skt{नाम}}\Afootnote{H; \emph{om.}\ A}}\skt{ }\edtext{\skt{मर्म्मणी}}{\lemma{\skt{मर्म्मणी}}\Afootnote{H; \emph{om.}\ A}}\skt{~। तत्र }\edtext{\skt{शोणि\-त\-क्ष\-या\-त्पा\-ण्डु\-वि\-वर्ण्णो}}{\lemma{\skt{पाण्डुविवर्ण्णो}}\Afootnote{H; \skt{पाण्डुर्विवर्णो}\ A}}\skt{ }\edtext{\skt{हीनरूपो वा}}{\lemma{\skt{हीनरूपो वा}}\Afootnote{H; \skt{हीनरूपश्च}\ A}}\skt{ }\edtext{\skt{मरणं}}{\lemma{\skt{मरणं}}\Afootnote{H; \skt{म्रियते}\ A}}\skt{~। }\edtext{\skt{पार्श्व\-यो\-र्ज्ज\-घ\-न\-ब\-हि\-र्ब्भागे}}{\lemma{\skt{ज्जघ\appbrk{}न\appbrk{}ब\appbrk{}हि\appbrk{}र्ब्भागे}}\Afootnote{H; \skt{जघ\appbrk{}न\appbrk{}ब\appbrk{}हि\appbrk{}र्भागे}\ A}}\skt{ }\edtext{\edtext{\skt{पृष्ठ\-वं\-श\-मु\-भ\-यतो}}{\lemma{\skt{पृष्ठवंशम्}}\Afootnote{H; \skt{पृष्ठवंश}\ A}}}{\lemma{\skt{उभयतो}}\Afootnote{H; \skt{मुभयतो}\ A}}\skt{ कुकुन्दरे }\edtext{\skt{नाम}}{\lemma{\skt{नाम}}\Afootnote{H; \emph{om.}\ A}}\skt{ }\edtext{\skt{मर्म्मणी}}{\lemma{\skt{मर्म्मणी}}\Afootnote{H; \emph{om.}\ A}}\skt{~। तत्र }\edtext{\edtext{\skt{स्पर्शा\-ज्ञा\-न\-म\-धः\-काये}}{\lemma{\skt{स्पर्शाज्ञानम्}}\Afootnote{H; \skt{स्पर्शाज्ञान}\ A}}}{\lemma{\skt{अधःकाये}}\Afootnote{H; \skt{मधःकाये}\ A}}\skt{ }\edtext{\skt{दौर्ब्बल्यं}}{\lemma{\skt{दौर्ब्बल्यं}}\Afootnote{H; \emph{om.}\ A}}\skt{ }\edtext{\skt{चेष्टानाशश्च}}{\lemma{\skt{चेष्टानाशश्च}}\Afootnote{H; \skt{चेष्टोपघातश्}\ A}}\skt{ }\edtext{\skt{भवति}}{\lemma{\skt{भवति}}\Afootnote{H; \skt{च}\ A}}\skt{~। }\edtext{\edtext{\edtext{\skt{श्रोणी\-क\-र्ण्ण\-यो\-रु\-प\-र्या\-मा\-श\-या\-च्छा\-दने}}{\lemma{\skt{श्रोणीकर्ण्णयोर्}}\Afootnote{H; \skt{श्रोणीकाण्डयोर्}\ A}}}{\lemma{\skt{आमाशयाच्}}\Afootnote{H; \skt{आशयाच्}\ A}}}{\lemma{\skt{छादने}}\Afootnote{H; \skt{छादनौ}\ A}}\skt{ }\edtext{\skt{पार्श्वा\-न्त\-र\-प्र\-ति\-वहे}}{\lemma{\skt{पार्श्वा\appbrk{}न्त\appbrk{}र\appbrk{}प्र\appbrk{}ति\appbrk{}वहे}}\Afootnote{H; \skt{पार्श्वा\appbrk{}न्त\appbrk{}र\appbrk{}प्र\appbrk{}ति\appbrk{}बद्धौ}\ A}}\skt{ }\edtext{\skt{नितम्बे}}{\lemma{\skt{नितम्बे}}\Afootnote{H; \skt{नितम्बौ}\ A}}\skt{ }\edtext{\skt{नाम}}{\lemma{\skt{नाम}}\Afootnote{H; \emph{om.}\ A}}\skt{ }\edtext{\skt{मर्म्मणी}}{\lemma{\skt{मर्म्मणी}}\Afootnote{H; \emph{om.}\ A}}\skt{~। तत्रा\-धः\-का\-य\-शोषो }\edtext{\edtext{\skt{दौर्ब्ब\-ल्य\-म्म\-र\-णञ्च}}{\lemma{\skt{दौर्ब्बल्यम्}}\Afootnote{H; \skt{दौर्बल्याच्च}\ A}}}{\lemma{\skt{मरणञ्च}}\Afootnote{H; \skt{मरणं}\ A}}\skt{~। }\edtext{\skt{अधः\-पा\-र्श्वा\-न्त\-र\-प्र\-ति\-बद्धे}}{\lemma{\skt{अधः\appbrk{}पा\appbrk{}र्श्वा\appbrk{}न्त\appbrk{}र\appbrk{}प्र\appbrk{}ति\appbrk{}बद्धे}}\Afootnote{H; \skt{अधः\appbrk{}पा\appbrk{}र्श्वा\appbrk{}न्त\appbrk{}र\appbrk{}प्र\appbrk{}ति\appbrk{}बद्धौ}\ A}}\skt{ }\edtext{\skt{जघ\-न\-पा\-र्श्व\-मध्ये}}{\lemma{\skt{जघ\appbrk{}न\appbrk{}पा\appbrk{}र्श्व\appbrk{}मध्ये}}\Afootnote{H; \skt{जघ\appbrk{}न\appbrk{}पा\appbrk{}र्श्व\appbrk{}म\appbrk{}ध्य\appbrk{}योस्}\ A}}\skt{ }\edtext{\skt{तिर्य\-गू\-र्द्ध्व\-ज\-घ\-ना\-त्पा\-र्श्व\-सन्धी}}{\lemma{\skt{तिर्य\appbrk{}गू\appbrk{}र्द्ध्व\appbrk{}ज\appbrk{}घ\appbrk{}नात्}}\Afootnote{H; \skt{तिर्यगूर्ध्वं च जघनात्}\ A}}\skt{ }\edtext{\skt{नाम}}{\lemma{\skt{नाम}}\Afootnote{H; \emph{om.}\ A}}\skt{ }\edtext{\skt{मर्म्मणी}}{\lemma{\skt{मर्म्मणी}}\Afootnote{H; \emph{om.}\ A}}\skt{~। तत्र }\edtext{\skt{शोणि\-त\-पू\-र्ण्ण\-को\-ष्ठ\-तया}}{\lemma{\skt{शोणि\appbrk{}त\appbrk{}पू\appbrk{}र्ण्ण\appbrk{}को\appbrk{}ष्ठ\appbrk{}तया}}\Afootnote{H; \skt{लोहि\appbrk{}त\appbrk{}पू\appbrk{}र्ण\appbrk{}को\appbrk{}ष्ठ\appbrk{}तया}\ A}}\skt{ म्रियते~। स्तन\-मू\-ला\-दृ\-जू\-भ\-यतः पृष्ठवंशस्य बृहती }\edtext{\skt{नाम}}{\lemma{\skt{नाम}}\Afootnote{H; \emph{om.}\ A}}\skt{ }\edtext{\skt{मर्म्मणी}}{\lemma{\skt{मर्म्मणी}}\Afootnote{H; \emph{om.}\ A}}\skt{~। तत्र शोणि\-ता\-ति\-प्र\-वृ\-त्ति\-नि\-मि\-त्तै\-रु\-प\-द्र\-वै\-र्म्म्रि\-यते~। पृष्ठोपरि पृष्ठ\-वं\-श\-मु\-भ\-यतः~। }\edtext{\edtext{\skt{त्रिक\-सं\-ब\-न्धेऽ\-न्स\-फले}}{\lemma{\skt{त्रिकसंबन्धे}}\Afootnote{H; \skt{त्रिकसंबद्धे}\ A}}}{\lemma{\skt{ऽन्सफले}}\Afootnote{H; \skt{अंसफ}\ A}}\skt{ }\edtext{\skt{नाम}}{\lemma{\skt{नाम}}\Afootnote{H; \emph{om.}\ A}}\skt{ }\edtext{\skt{मर्म्मणी}}{\lemma{\skt{मर्म्मणी}}\Afootnote{H; \skt{लके}\ A}}\skt{~। तत्र }\edtext{\edtext{\skt{बाहुश्चापशोषः}}{\lemma{\skt{बाहुश्}}\Afootnote{H; \skt{बाह्वोः}\ A}}}{\lemma{\skt{चापशोषः}}\Afootnote{H; \skt{स्वापशोषौ}\ A}}\skt{~। }\edtext{\edtext{\skt{बाहु\-शी\-र्ष\-ग्री\-वा\-म\-ध्येऽं\-श\-पी\-ठ\-स्क\-न्ध\-सं\-धाने}}{\lemma{\skt{बाहु\appbrk{}शी\appbrk{}र्ष\appbrk{}ग्री\appbrk{}वा\appbrk{}मध्ये}}\Afootnote{H; \skt{बाहु\appbrk{}मू\appbrk{}र्ध\appbrk{}ग्री\appbrk{}वा\appbrk{}मध्यें}\ A}}}{\lemma{\skt{ऽंश\appbrk{}पी\appbrk{}ठ\appbrk{}स्क\appbrk{}न्ध\appbrk{}सं\appbrk{}धाने}}\Afootnote{H; \skt{ऽस\appbrk{}पी\appbrk{}ठ\appbrk{}स्क\appbrk{}न्ध\appbrk{}बन्}\ A}}\skt{ }\edtext{\skt{क्रकन्यासे}}{\lemma{\skt{क्रकन्यासे}}\Afootnote{H; \skt{ध}\ A}}\skt{ }\edtext{\skt{नाम}}{\lemma{\skt{नाम}}\Afootnote{H; \skt{ना}\ A}}\skt{ }\edtext{\skt{मर्म्मणी}}{\lemma{\skt{मर्म्मणी}}\Afootnote{H; \skt{वंसौ}\ A}}\skt{~। तत्र }\edtext{\edtext{\skt{तब्ध\-बा\-हु\-र्ब्भ\-ग्न\-पृ\-ष्ठता}}{\lemma{\skt{तब्धबाहुर्}}\Afootnote{H; \skt{स्तब्धबाहु}\ A}}\skt{ वा}}{\lemma{\skt{ब्भग्नपृष्ठता वा}}\Afootnote{H; \emph{om.}\ A}}\skt{ }\edtext{\skt{भवति}}{\lemma{\skt{भवति}}\Afootnote{H; \skt{ता}\ A}}\skt{~। एवमेतानि चतुर्द्दश पृष्ठमर्म्माणि व्याख्यातानि~॥}\pend

\pstart \ledsidenote{\passref{3.6.27}}\skt{अत }\edtext{\edtext{\skt{ऊर्द्ध्व\-मू\-र्द्ध\-ज\-त्रु\-ग\-तानि}}{\lemma{\skt{ऊर्द्ध्वम्}}\Afootnote{H; \skt{ऊर्ध्व}\ A}}}{\lemma{\skt{ऊर्द्ध\appbrk{}ज\appbrk{}त्रु\appbrk{}ग\appbrk{}तानि}}\Afootnote{H; \skt{मूर्ध्व\appbrk{}ज\appbrk{}त्रु\appbrk{}ग\appbrk{}तानि}\ A}}\skt{ }\edtext{\edtext{\skt{मर्म्मा\-ण्य\-नु\-व्या\-ख्या\-स्यामः}}{\lemma{\skt{मर्म्माण्य्}}\Afootnote{H; \emph{om.}\ A}}}{\lemma{\skt{अनुव्याख्यास्यामः}}\Afootnote{H; \skt{व्याख्यास्यामः}\ A}}\skt{~॥ तत्र कण्ठ\-ना\-डी\-मु\-भ\-य\-त\-श्च\-तस्रो धमन्यो }\edtext{\skt{व्यत्यासेन}}{\lemma{\skt{व्यत्यासेन}}\Afootnote{H; \emph{om.}\ A}}\skt{~। द्वे नीले द्वे च मन्ये }\edtext{\skt{नाम}}{\lemma{\skt{नाम}}\Afootnote{H; \skt{व्य}\ A}}\skt{ }\edtext{\skt{मर्म्मणी}}{\lemma{\skt{मर्म्मणी}}\Afootnote{H; \skt{त्यासेन}\ A}}\skt{~। तत्र मूकता }\edtext{\skt{स्वरवैकृत्य}}{\lemma{\skt{स्वरवैकृत्य}}\Afootnote{H; \skt{स्वरवैकृ}\ A}}\skt{ }\edtext{\edtext{\skt{मर\-ण\-म\-र\-स\-ग्रा\-हिता}}{\lemma{\skt{मरणम्}}\Afootnote{H; \skt{तम}\ A}}\skt{ च}}{\lemma{\skt{अरसग्राहिता च}}\Afootnote{H; \skt{रसग्राहिता च}\ A}}\skt{ }\edtext{\skt{भवति}}{\lemma{\skt{भवति}}\Afootnote{H; \emph{om.}\ A}}\skt{~। ग्रीवा\-या\-मु\-भ\-य\-त\-श्च\-तस्रः }\edtext{\skt{सिरामातृका}}{\lemma{\skt{सिरामातृका}}\Afootnote{H; \skt{सिरा मातृकाः}\ A}}\skt{ }\edtext{\skt{नाम}}{\lemma{\skt{नाम}}\Afootnote{H; \emph{om.}\ A}}\skt{ }\edtext{\skt{मर्म्मणी}}{\lemma{\skt{मर्म्मणी}}\Afootnote{H; \emph{om.}\ A}}\skt{~। तत्र }\edtext{\skt{सद्यो}}{\lemma{\skt{सद्यो}}\Afootnote{H; \skt{सद्यो}\ A}}\skt{ मरणं~। शिरोग्रीवयोः }\edtext{\skt{कृकाटिका}}{\lemma{\skt{कृकाटिका}}\Afootnote{H; \skt{सन्धाने कृकाटिके}\ A}}\skt{ }\edtext{\skt{नाम}}{\lemma{\skt{नाम}}\Afootnote{H; \emph{om.}\ A}}\skt{ }\edtext{\skt{मर्म्मणी}}{\lemma{\skt{मर्म्मणी}}\Afootnote{H; \emph{om.}\ A}}\skt{~। तत्र }\edtext{\skt{चलशीर्षता}}{\lemma{\skt{चलशीर्षता}}\Afootnote{H; \skt{चलमूर्धता}\ A}}\skt{ }\edtext{\skt{मरणम्वा}}{\lemma{\skt{मरणम्वा}}\Afootnote{H; \emph{om.}\ A}}\skt{~। }\edtext{\skt{कर्ण्ण\-पृ\-ष्ठ\-तोऽ\-धः\-सं\-स्रिते}}{\lemma{\skt{ऽधःसंस्रिते}}\Afootnote{H; \skt{ऽधःसंश्रिते}\ A}}\skt{ विधुरे }\edtext{\skt{नाम}}{\lemma{\skt{नाम}}\Afootnote{H; \emph{om.}\ A}}\skt{ }\edtext{\skt{मर्म्मणी}}{\lemma{\skt{मर्म्मणी}}\Afootnote{H; \emph{om.}\ A}}\skt{~। तत्र बाधिर्यं~। घ्राण\-मा\-र्ग्ग\-मु\-भ\-यतः स्रोतो\-मा\-र्ग्ग\-प्र\-ति\-ब\-द्धेऽ\-भ्य\-न्त\-रतः फणे }\edtext{\skt{नाम}}{\lemma{\skt{नाम}}\Afootnote{H; \emph{om.}\ A}}\skt{ }\edtext{\skt{मर्म्मणी}}{\lemma{\skt{मर्म्मणी}}\Afootnote{H; \emph{om.}\ A}}\skt{~। तत्र }\edtext{\skt{गन्धा\-ज्ञा\-न\-मा\-पा\-द\-यति}}{\lemma{\skt{आपादयति}}\Afootnote{H; \emph{om.}\ A}}\skt{~। }\edtext{\skt{भ्रुवोः}}{\lemma{\skt{भ्रुवोः}}\Afootnote{H; \skt{भ्रू}\ A}}\skt{ पुच्छान्तयोरधः~। }\edtext{\edtext{\skt{नेत्रयोर्ब्बाह्यतः}}{\lemma{\skt{नेत्रयोर्}}\Afootnote{H; \skt{ऽक्ष्णो}\ A}}}{\lemma{\skt{ब्बाह्यतः}}\Afootnote{H; \skt{र्बाह्यतो}\ A}}\skt{~। अपाङ्गौ }\edtext{\skt{नाम}}{\lemma{\skt{नाम}}\Afootnote{H; \emph{om.}\ A}}\skt{ }\edtext{\skt{मर्म्मणी}}{\lemma{\skt{मर्म्मणी}}\Afootnote{H; \emph{om.}\ A}}\skt{~। }\edtext{\skt{तत्रा\-न्ध्य\-दृ\-ष्ट्यु\-प\-घातो वा}}{\lemma{\skt{तत्रा\appbrk{}न्ध्य\appbrk{}दृ\appbrk{}ष्ट्यु\appbrk{}प\appbrk{}घातो वा}}\Afootnote{H; \skt{तत्रान्ध्यं दृष्ट्युपघातो वा}\ A}}\skt{~। }\edtext{\skt{भ्रुवोः}}{\lemma{\skt{भ्रुवोः}}\Afootnote{H; \emph{om.}\ A}}\skt{ }\edtext{\edtext{\edtext{\skt{पुच्छ\-यो\-रु\-प\-र्य\-नु\-क\-र्ण्ण\-ल\-लाटो}}{\lemma{\skt{पुच्छयोर्}}\Afootnote{H; \emph{om.}\ A}}}{\lemma{\skt{उपर्य्}}\Afootnote{H; \emph{om.}\ A}}}{\lemma{\skt{अनु\appbrk{}क\appbrk{}र्ण्ण\appbrk{}ल\appbrk{}लाटो}}\Afootnote{H; \emph{om.}\ A}}\skt{ }\edtext{\skt{शंखौ}}{\lemma{\skt{शंखौ}}\Afootnote{H; \emph{om.}\ A}}\skt{ }\edtext{\skt{नाम}}{\lemma{\skt{नाम}}\Afootnote{H; \emph{om.}\ A}}\skt{ }\edtext{\skt{मर्म्मणी}}{\lemma{\skt{मर्म्मणी}}\Afootnote{H; \emph{om.}\ A}}\skt{~। }\edtext{\skt{तत्र}}{\lemma{\skt{तत्र}}\Afootnote{H; \emph{om.}\ A}}\skt{ }\edtext{\skt{सद्यो}}{\lemma{\skt{सद्यो}}\Afootnote{H; \emph{om.}\ A}}\skt{ }\edtext{\skt{मरणं}}{\lemma{\skt{मरणं}}\Afootnote{H; \emph{om.}\ A}}\skt{~। }\edtext{\skt{भ्रुवोरन्तरे}}{\lemma{\skt{अन्तरे}}\Afootnote{H; \emph{om.}\ A}}\skt{ उपरि निम्नयोरावर्त्तौ }\edtext{\skt{नाम}}{\lemma{\skt{नाम}}\Afootnote{H; \skt{ना}\ A}}\skt{ }\edtext{\skt{मर्म्मणी}}{\lemma{\skt{मर्म्मणी}}\Afootnote{H; \skt{म}\ A}}\skt{~। }\edtext{\edtext{\skt{तत्राप्यान्ध्यं}}{\lemma{\skt{तत्राप्य्}}\Afootnote{H; \skt{तत्राप्}\ A}}}{\lemma{\skt{आन्ध्यं}}\Afootnote{H; \skt{यान्ध्यं}\ A}}\skt{ दृष्ट्युपघातो वा~। }\edtext{\skt{आव\-र्त्त\-शं\-ख\-यो\-रु\-परि}}{\lemma{\skt{आव\appbrk{}र्त्त\appbrk{}शं\appbrk{}ख\appbrk{}योर्}}\Afootnote{H; \skt{भ्रुवो\appbrk{}र\appbrk{}न्त\appbrk{}यो\appbrk{}रु\appbrk{}परि कर्ण\appbrk{}ल\appbrk{}ला\appbrk{}ट\appbrk{}यो\appbrk{}र्मध्ये शङ्खौ, तत्र सद्योमरणं; शङ्खयोर्}\ A}}\skt{ }\edtext{\skt{केशान्ते}}{\lemma{\skt{केशान्ते}}\Afootnote{H; \skt{केशान्त}\ A}}\skt{ उत्क्षेपौ }\edtext{\skt{नाम}}{\lemma{\skt{नाम}}\Afootnote{H; \emph{om.}\ A}}\skt{ }\edtext{\skt{मर्म्मणी}}{\lemma{\skt{मर्म्मणी}}\Afootnote{H; \emph{om.}\ A}}\skt{~। तत्र सशल्यः }\edtext{\skt{पाका\-त्प\-ति\-त\-शल्यो}}{\lemma{\skt{पाकात्}}\Afootnote{H; \skt{जीवेत्पाकात्}\ A}}\skt{ वा }\edtext{\skt{जीवे}}{\lemma{\skt{जीवे}}\Afootnote{H; \emph{om.}\ A}}\skt{ नोद्धृतशल्यः~। भ्रुवोर्म्मध्ये }\edtext{\edtext{\skt{स्थपणिर्न्नाम}}{\lemma{\skt{स्थपणिर्}}\Afootnote{H; \skt{स्थप}\ A}}}{\lemma{\skt{न्नाम}}\Afootnote{H; \emph{om.}\ A}}\skt{ }\edtext{\skt{मर्म्म}}{\lemma{\skt{मर्म्म}}\Afootnote{H; \skt{नी}\ A}}\skt{~। तत्रो\-त्क्षे\-प\-व\-त्पञ्च सन्धयः शिरसि विभक्ताः सीमन्ता }\edtext{\skt{नाम}}{\lemma{\skt{नाम}}\Afootnote{H; \skt{ना}\ A}}\skt{ }\edtext{\skt{मर्म्मणी}}{\lemma{\skt{मर्म्मणी}}\Afootnote{H; \skt{म}\ A}}\skt{~। }\edtext{\skt{तत्रो\-त्मा\-द\-चि\-त्त\-ना\-साभ्यां}}{\lemma{\skt{तत्रो\appbrk{}त्मा\appbrk{}द\appbrk{}चि\appbrk{}त्त\appbrk{}ना\appbrk{}साभ्यां}}\Afootnote{H; \skt{तत्रो\appbrk{}न्मा\appbrk{}द\appbrk{}भ\appbrk{}य\appbrk{}चि\appbrk{}त्त\appbrk{}ना\appbrk{}शैर्}\ A}}\skt{ मरणं~। }\edtext{\skt{घ्राण\-स्रो\-त्रा\-क्षि\-स\-न्धि\-जि\-ह्वा\-त\-र्प्प\-णीनां}}{\lemma{\skt{घ्राण\appbrk{}स्रो\appbrk{}त्रा\appbrk{}क्षि\appbrk{}स\appbrk{}न्धि\appbrk{}जि\appbrk{}ह्वा\appbrk{}त\appbrk{}र्प्प\appbrk{}णीनां}}\Afootnote{H; \skt{घ्राण\appbrk{}श्रो\appbrk{}त्रा\appbrk{}क्षि\appbrk{}जि\appbrk{}ह्वा\appbrk{}स\appbrk{}न्त\appbrk{}र्प\appbrk{}णीनां}\ A}}\skt{ सिराणां मध्ये }\edtext{\skt{शिरः}}{\lemma{\skt{शिरः}}\Afootnote{H; \skt{सि}\ A}}\skt{ }\edtext{\skt{स्रोतः}}{\lemma{\skt{स्रोतः}}\Afootnote{H; \skt{रा}\ A}}\skt{ }\edtext{\skt{सन्निपातस्तानि}}{\lemma{\skt{तानि}}\Afootnote{H; \skt{शृङ्गाटकानि , तानि}\ A}}\skt{ चत्वारि }\edtext{\skt{शृङ्गाटकानि}}{\lemma{\skt{शृङ्गाटकानि}}\Afootnote{H; \emph{om.}\ A}}\skt{ मर्म्माणि~। तत्रापि }\edtext{\skt{सद्यो}}{\lemma{\skt{सद्यो}}\Afootnote{H; \skt{सद्यो}\ A}}\skt{ मरणं~। }\edtext{\edtext{\skt{मस्त\-का\-भ्य\-न्त\-र\-मु\-प\-रि\-ष्टा\-त्सि\-रा\-स\-न्धि\-स\-न्नि\-पातो}}{\lemma{\skt{मस्त\appbrk{}का\appbrk{}भ्य\appbrk{}न्त\appbrk{}रम्}}\Afootnote{H; \skt{मस्तकाभ्यन्तरो}\ A}}}{\lemma{\skt{उपरिष्टात्}}\Afootnote{H; \skt{परिष्ठात्}\ A}}\skt{ }\edtext{\skt{रोमावर्त्तः }\edtext{\skt{सोऽधिपतिर्न्नाम}}{\lemma{\skt{न्नाम}}\Afootnote{H; \emph{om.}\ A}}}{\lemma{\skt{रोमावर्त्तः सो}}\Afootnote{H; \skt{रोमावर्तो}\ A}}\skt{ }\edtext{\skt{मर्म्म}}{\lemma{\skt{मर्म्म}}\Afootnote{H; \emph{om.}\ A}}\skt{~। तत्रापि }\edtext{\skt{सद्यो}}{\lemma{\skt{सद्यो}}\Afootnote{H; \skt{सद्य}\ A}}\skt{ }\edtext{\skt{मरणं}}{\lemma{\skt{मरणं}}\Afootnote{H; \skt{एव}\ A}}\skt{~। }\edtext{\skt{एव\-मे\-ता\-न्यू\-र्द्ध्व\-ज\-त्रु\-ग\-तानि}}{\lemma{\skt{ऊर्द्ध्व\appbrk{}ज\appbrk{}त्रु\appbrk{}ग\appbrk{}तानि}}\Afootnote{H; \skt{सप्त\appbrk{}त्रिं\appbrk{}श\appbrk{}दू\appbrk{}र्ध्व\appbrk{}ज\appbrk{}त्रु\appbrk{}ग\appbrk{}तानि}\ A}}\skt{ }\edtext{\skt{सप्त\-त्रिं\-श\-द्म\-र्म्माणि}}{\lemma{\skt{सप्तत्रिंशद्}}\Afootnote{H; \emph{om.}\ A}}\skt{ व्याख्यातानि }\edtext{\skt{भवन्ति}}{\lemma{\skt{भवन्ति}}\Afootnote{H; \emph{om.}\ A}}\skt{~॥}\pend

\pstart\versepar \ledsidenote{\passref{3.6.28}}\skt{भवन्ति चात्र श्लोकाः~॥ }\pend
\pstart\versepar \skt{}\edtext{\skt{ऊर्व्व्यः}}{\lemma{\skt{ऊर्व्व्यः}}\Afootnote{H; \skt{उर्व्यः}\ A}}\skt{ शिरान्सि }\edtext{\skt{चिटिपे}}{\lemma{\skt{चिटिपे}}\Afootnote{H; \skt{विटपे}\ A}}\skt{ }\edtext{\skt{सहकक्षपार्श्वे}}{\lemma{\skt{सहकक्षपार्श्वे}}\Afootnote{H; \skt{च सकक्षपार्श्वे}\ A}}\skt{ }\edtext{\skt{एकै\-क\-म\-ङ्गु\-ल\-मिता}}{\lemma{\skt{एकै\appbrk{}क\appbrk{}म\appbrk{}ङ्गु\appbrk{}ल\appbrk{}मिता}}\Afootnote{H; \skt{एकै\appbrk{}क\appbrk{}म\appbrk{}ङ्गु\appbrk{}ल\appbrk{}मितं}\ A}}\skt{ }\edtext{\skt{स्तनमूलगुल्फाः}}{\lemma{\skt{स्तनमूलगुल्फाः}}\Afootnote{H; \skt{स्तन\appbrk{}पू\appbrk{}र्व\appbrk{}मू\appbrk{}लम्}\ A}}\skt{~। }\pend
\pstart\versepar \skt{}\edtext{\skt{ज्ञेया}}{\lemma{\skt{ज्ञेया}}\Afootnote{H; \skt{वि}\ A}}\skt{ }\edtext{\edtext{\skt{द्विरङ्गुलमिता}}{\lemma{\skt{द्विर्}}\Afootnote{H; \skt{द्ध्}\ A}}}{\lemma{\skt{अङ्गुलमिता}}\Afootnote{H; \skt{यङ्गु\appbrk{}ल\appbrk{}द्व\appbrk{}य\appbrk{}मितं}\ A}}\skt{ }\edtext{\skt{मणिबन्धषष्ठा}}{\lemma{\skt{मणिबन्धषष्ठा}}\Afootnote{H; \skt{मणिबन्धगुल्फं}\ A}}\skt{ त्रीण्येव }\edtext{\skt{जानुनि}}{\lemma{\skt{जानुनि}}\Afootnote{H; \skt{जानु स}\ A}}\skt{ परं सह }\edtext{\skt{कुर्प्पराभ्यां}}{\lemma{\skt{कुर्प्पराभ्यां}}\Afootnote{H; \skt{कूर्पराभ्याम्}\ A}}\skt{~॥}\pend

\pstart\versepar \ledsidenote{\passref{3.6.29}}\edtext{\skt{हृद्ब\-स्ति\-कु\-र्च्च\-गु\-द\-नाभि}}{\lemma{\skt{हृद्ब\appbrk{}स्ति\appbrk{}कु\appbrk{}र्च्च\appbrk{}गु\appbrk{}द\appbrk{}नाभि}}\Afootnote{H; \skt{हृद्ब\appbrk{}स्ति\appbrk{}कू\appbrk{}र्च\appbrk{}गु\appbrk{}द\appbrk{}नाभि}\ A}}\skt{ वदन्ति मूर्द्ध्नि चत्वारि पञ्च च गले दश यानि च द्वे~॥ }\pend
\pstart\versepar \skt{तानि }\edtext{\skt{श्चपा\-णि\-त\-ल\-कु\-ञ्चि\-त\-स\-न्मि\-तानि,}}{\lemma{\skt{श्चपा\appbrk{}णि\appbrk{}त\appbrk{}ल\appbrk{}कु\appbrk{}ञ्चि\appbrk{}त\appbrk{}स\appbrk{}न्मि\appbrk{}तानि}}\Afootnote{H; \skt{स्वपा\appbrk{}णि\appbrk{}त\appbrk{}ल\appbrk{}कु\appbrk{}ञ्चि\appbrk{}त\appbrk{}सं\appbrk{}मि\appbrk{}तानि}\ A}}\skt{ }\edtext{\edtext{\skt{शेषाण्यवैहि}}{\lemma{\skt{शेषाण्य्}}\Afootnote{H; \skt{शेषाण्}\ A}}}{\lemma{\skt{अवैहि}}\Afootnote{H; \skt{यवेहि}\ A}}\skt{ }\edtext{\skt{परिवर्त्ति}}{\lemma{\skt{परिवर्त्ति}}\Afootnote{H; \skt{परिविस्त}\ A}}\skt{ }\edtext{\skt{ततो}}{\lemma{\skt{ततो}}\Afootnote{H; \skt{रतो}\ A}}\skt{ }\edtext{\skt{गुलार्द्धं}}{\lemma{\skt{गुलार्द्धं}}\Afootnote{H; \skt{ऽङ्गुलार्धम्}\ A}}\skt{~॥}\pend

\pstart\versepar \ledsidenote{\passref{3.6.31}}\skt{च्छिन्नेषु पाणिचरणेषु सिरा }\edtext{\edtext{\edtext{\skt{मुखा\-ना\-मा\-कु\-ञ्च\-ना\-त्खलु}}{\lemma{\skt{मुखानाम्}}\Afootnote{H; \skt{नराणां}\ A}}}{\lemma{\skt{आकुञ्चनात्}}\Afootnote{H; \skt{सङ्कोच}\ A}}}{\lemma{\skt{खलु}}\Afootnote{H; \skt{मीयु}\ A}}\skt{ }\edtext{\edtext{\edtext{\skt{नृणा\-म\-सृ\-ग\-ल्प\-मेति}}{\lemma{\skt{नृणाम्}}\Afootnote{H; \skt{र}\ A}}}{\lemma{\skt{असृगल्पम्}}\Afootnote{H; \skt{सृगल्}\ A}}}{\lemma{\skt{एति}}\Afootnote{H; \emph{om.}\ A}}\skt{~। }\pend
\pstart\versepar \skt{}\edtext{\skt{छिन्नेषु}}{\lemma{\skt{छिन्नेषु}}\Afootnote{H; \emph{om.}\ A}}\skt{ }\edtext{\skt{पाणिचरणेषु}}{\lemma{\skt{पाणिचरणेषु}}\Afootnote{H; \emph{om.}\ A}}\skt{ }\edtext{\edtext{\skt{भवन्त्यतस्तु}}{\lemma{\skt{भवन्त्य्}}\Afootnote{H; \emph{om.}\ A}}}{\lemma{\skt{अतस्तु}}\Afootnote{H; \emph{om.}\ A}}\skt{ }\edtext{\skt{नालेषु}}{\lemma{\skt{नालेषु}}\Afootnote{H; \emph{om.}\ A}}\skt{ }\edtext{\skt{वृत्तमथितेषु}}{\lemma{\skt{वृत्तमथितेषु}}\Afootnote{H; \emph{om.}\ A}}\skt{ }\edtext{\skt{यथोत्पलानि}}{\lemma{\skt{यथोत्पलानि}}\Afootnote{H; \skt{पमतो निरेति~। प्राप्या\appbrk{}मि\appbrk{}त\appbrk{}व्य\appbrk{}स\appbrk{}न\appbrk{}मु\appbrk{}ग्र\appbrk{}मतो मनुष्याः संच्छि\appbrk{}न्न\appbrk{}शा\appbrk{}ख\appbrk{}त\appbrk{}रु\appbrk{}व\appbrk{}न्नि\appbrk{}धनं न यान्ति}\ A}}\skt{~॥}\pend

\pstart\versepar \ledsidenote{\passref{3.6.32}}\skt{क्षिप्रेषु }\edtext{\skt{वत्सतलजेषु च}}{\lemma{\skt{वत्सतलजेषु च}}\Afootnote{H; \skt{तत्र सतलेषु ह}\ A}}\skt{ }\edtext{\skt{भेद\-दो\-षा\-ङ्ग\-च्छ\-त्य\-सृ\-ग्ब\-हु\-रु\-जञ्च}}{\lemma{\skt{भेद\appbrk{}दो\appbrk{}षा\appbrk{}ङ्ग\appbrk{}च्छ\appbrk{}त्य\appbrk{}सृ\appbrk{}ग्ब\appbrk{}हु\appbrk{}रु\appbrk{}जञ्च}}\Afootnote{H; \skt{तेषु रक्तं गच्छत्यतीव पवनश्च}\ A}}\skt{ }\edtext{\skt{करोति}}{\lemma{\skt{करोति}}\Afootnote{H; \skt{रुजं करोति}\ A}}\skt{ }\edtext{\skt{वायुः}}{\lemma{\skt{वायुः}}\Afootnote{H; \emph{om.}\ A}}\skt{~। }\pend
\pstart\versepar \skt{एवं विना\-श\-मु\-प\-यान्ति हि तत्र विद्धा }\edtext{\skt{किञ्ज\-ल्क\-प\-त्र\-म\-थ\-ना\-दिव}}{\lemma{\skt{किञ्ज\appbrk{}ल्क\appbrk{}प\appbrk{}त्र\appbrk{}म\appbrk{}थ\appbrk{}ना\appbrk{}दिव}}\Afootnote{H; \emph{om.}\ A}}\skt{ }\edtext{\skt{पङ्कजानि}}{\lemma{\skt{पङ्कजानि}}\Afootnote{H; \skt{वृक्षा इवा\appbrk{}यु\appbrk{}ध\appbrk{}नि\appbrk{}पा\appbrk{}त\appbrk{}नि\appbrk{}कृ\appbrk{}त्त\appbrk{}मूलाः}\ A}}\skt{~॥}\pend

\pstart \ledsidenote{\passref{3.6.33}}\edtext{\skt{मर्म्माणि}}{\lemma{\skt{मर्म्माणि}}\Afootnote{H; \skt{तस्मा\appbrk{}त्त\appbrk{}यो\appbrk{}र\appbrk{}भि\appbrk{}ह\appbrk{}तस्य तु पाणिपादं छेत्तव्यमाशु मणि\appbrk{}ब\appbrk{}न्ध\appbrk{}न\appbrk{}गु\appbrk{}ल्फ\appbrk{}देशे मर्माणि}\ A}}\skt{ }\edtext{\skt{शल्य\-वि\-ष\-या\-द्र्ध\-मु\-दा\-ह\-रन्ति}}{\lemma{\skt{शल्य\appbrk{}वि\appbrk{}ष\appbrk{}या\appbrk{}द्र्धम्}}\Afootnote{H; \skt{शल्य\appbrk{}वि\appbrk{}ष\appbrk{}या\appbrk{}र्धम्}\ A}}\skt{ }\edtext{\skt{यस्माद्धि}}{\lemma{\skt{यस्माद्धि}}\Afootnote{H; \skt{यस्माच्च}\ A}}\skt{ मर्म्मसु हता न भवन्ति }\edtext{\skt{मर्त्त्याः}}{\lemma{\skt{मर्त्त्याः}}\Afootnote{H; \skt{सद्यः}\ A}}\skt{~।}\pend

\pstart\versepar \ledsidenote{\passref{3.6.34}}\skt{जीवन्ति तत्र यदि वैद्यगुणेन के चित्ते प्राप्नुवन्ति विक\-ल\-त्व\-म\-सं\-शयं हि~॥ }\pend
\pstart\versepar \skt{सम्भि\-न्न\-ज\-र्ज्ज\-रि\-त\-को\-ष्ठ\-शि\-रः\-क\-पाला जीवन्ति }\edtext{\skt{शस्त्रविहतैश्च}}{\lemma{\skt{शस्त्रविहतैश्च}}\Afootnote{H; \skt{शस्त्रनिहतैश्च}\ A}}\skt{ शरीरदेशैः~।}\pend

\pstart\versepar \ledsidenote{\passref{3.6.35}}\skt{छिन्नैश्च }\edtext{\skt{सक्थि\-भु\-ज\-पा\-द\-क\-रै\-र\-शेषं}}{\lemma{\skt{अशेषं}}\Afootnote{H; \skt{अशेषैर्}\ A}}\skt{ येषान्न }\edtext{\skt{मर्म्मपतिता}}{\lemma{\skt{मर्म्मपतिता}}\Afootnote{H; \skt{मर्मसु कृता}\ A}}\skt{ विविधाः }\edtext{\skt{प्रहाराः}}{\lemma{\skt{प्रहाराः}}\Afootnote{H; \skt{प्रहाराः}\ A}}\skt{~॥ }\pend
\pstart\versepar \skt{}\edtext{\skt{अग्नीषोमानिलाः}}{\lemma{\skt{अग्नीषोमानिलाः}}\Afootnote{H; \skt{सोम\appbrk{}मा\appbrk{}रु\appbrk{}त\appbrk{}ते\appbrk{}जांसि रजः}\ A}}\skt{ }\edtext{\skt{सत्त्वं}}{\lemma{\skt{सत्त्वं}}\Afootnote{H; \skt{सत्त्व}\ A}}\skt{ }\edtext{\skt{रजश्च}}{\lemma{\skt{रजश्च}}\Afootnote{H; \skt{तमांसि}\ A}}\skt{ }\edtext{\skt{तम}}{\lemma{\skt{तम}}\Afootnote{H; \skt{च~। म}\ A}}\skt{ }\edtext{\skt{एव च}}{\lemma{\skt{एव च}}\Afootnote{H; \skt{र्मसु}\ A}}\skt{~। }\pend
\pstart\versepar \skt{}\edtext{\skt{प्रायेण}}{\lemma{\skt{प्रायेण}}\Afootnote{H; \skt{प्रा}\ A}}\skt{ }\edtext{\skt{मर्म्मसु}}{\lemma{\skt{मर्म्मसु}}\Afootnote{H; \skt{यशः}\ A}}\skt{ }\edtext{\skt{नृणां}}{\lemma{\skt{नृणां}}\Afootnote{H; \emph{om.}\ A}}\skt{ }\edtext{\skt{सन्ति}}{\lemma{\skt{सन्ति}}\Afootnote{H; \skt{पुंसां}\ A}}\skt{ }\edtext{\skt{भूतात्मना}}{\lemma{\skt{भूतात्मना}}\Afootnote{H; \skt{भूतात्मा चा}\ A}}\skt{ }\edtext{\skt{सह}}{\lemma{\skt{सह}}\Afootnote{H; \skt{वतिष्ठते}\ A}}\skt{~॥}\pend

\pstart\versepar \ledsidenote{\passref{3.6.36}}\edtext{\edtext{\skt{तस्मा\-त्म\-र्म्म\-स्व\-भि\-हता}}{\lemma{\skt{तस्मात्}}\Afootnote{H; \emph{om.}\ A}}}{\lemma{\skt{अभिहता}}\Afootnote{H; \skt{अभिहतास्त}\ A}}\skt{ }\edtext{\skt{न}}{\lemma{\skt{न}}\Afootnote{H; \skt{स्मान्न}\ A}}\skt{ जीवन्ति शरीरिणः~। }\pend
\pstart\versepar \skt{}\edtext{\skt{इन्द्रि\-या\-र्थे\-ष्व\-स\-म्वि\-त्ति\-र्म्म\-नो\-बु\-द्धि\-वि\-प\-र्ययः}}{\lemma{\skt{असम्वित्तिर्}}\Afootnote{H; \skt{असंप्राप्तिर्}\ A}}\skt{~॥}\pend

\pstart\versepar \ledsidenote{\passref{3.6.37}}\edtext{\skt{रुजः}}{\lemma{\skt{रुजः}}\Afootnote{H; \skt{रुजश्च}\ A}}\skt{ }\edtext{\skt{सुतीव्रा}}{\lemma{\skt{सुतीव्रा}}\Afootnote{H; \skt{विविधास्तीव्रा}\ A}}\skt{ }\edtext{\skt{विविधा}}{\lemma{\skt{विविधा}}\Afootnote{H; \emph{om.}\ A}}\skt{ }\edtext{\skt{भवन्त्यसुहते}}{\lemma{\skt{असुहते}}\Afootnote{H; \skt{आशुहरे}\ A}}\skt{ }\edtext{\skt{नरे}}{\lemma{\skt{नरे}}\Afootnote{H; \skt{हते}\ A}}\skt{~॥ }\pend
\pstart\versepar \skt{हते }\edtext{\skt{कालान्तरे}}{\lemma{\skt{कालान्तरे}}\Afootnote{H; \skt{कालान्तर}\ A}}\skt{ }\edtext{\skt{चापि}}{\lemma{\skt{चापि}}\Afootnote{H; \skt{घ्ने तु}\ A}}\skt{ }\edtext{\skt{ध्रुवो}}{\lemma{\skt{ध्रुवो}}\Afootnote{H; \skt{ध्रुवं}\ A}}\skt{ धातुक्षयो नृणां~॥}\pend

\pstart\versepar \ledsidenote{\passref{3.6.38}}\skt{ततो धातु\-क्ष\-या\-ज्ज\-न्तु\-र्व्वे\-द\-ना\-भिश्च नश्यति~। }\pend
\pstart\versepar \skt{हते }\edtext{\skt{वैकल्यजनके}}{\lemma{\skt{वैकल्यजनके}}\Afootnote{H; \skt{वैकल्यजनने}\ A}}\skt{ केवलं वैद्यनैपुणात्~॥}\pend

\pstart\versepar \ledsidenote{\passref{3.6.39}}\edtext{\edtext{\skt{शरी\-र\-क्ष\-य\-मा\-साद्य}}{\lemma{\skt{शरीरक्षयम्}}\Afootnote{H; \skt{शरीरं क्रिय}\ A}}}{\lemma{\skt{आसाद्य}}\Afootnote{H; \skt{या युक्तं}\ A}}\skt{ विक\-ल\-त्व\-म\-वा\-प्नु\-यात्~। }\pend
\pstart\versepar \skt{}\edtext{\skt{विशल्यघ्नेषु}}{\lemma{\skt{विशल्यघ्नेषु}}\Afootnote{H; \skt{विशल्यघ्ने तु}\ A}}\skt{ }\edtext{\skt{विज्ञेयो}}{\lemma{\skt{विज्ञेयो}}\Afootnote{H; \skt{विज्ञेयं}\ A}}\skt{ }\edtext{\skt{यो}}{\lemma{\skt{यो}}\Afootnote{H; \skt{पूर्वोक्तं यच्च कारणम्}\ A}}\skt{ }\edtext{\skt{हेतुः}}{\lemma{\skt{हेतुः}}\Afootnote{H; \emph{om.}\ A}}\skt{ }\edtext{\skt{प्रागुदीरितः}}{\lemma{\skt{प्रागुदीरितः}}\Afootnote{H; \emph{om.}\ A}}\skt{~॥}\pend

\pstart\versepar \ledsidenote{\passref{3.6.40}}\skt{रुजाकराणि मर्म्माणि क्षतानि विविधा रुजः~। }\pend
\pstart\versepar \skt{कुर्व्वन्त्यन्ते च वैकल्यं कुवैद्यवशगो यदि~॥}\pend

\pstart\versepar \ledsidenote{\passref{3.6.41}}\skt{छेद\-भे\-दा\-भि\-घा\-तेभ्यो }\edtext{\edtext{\skt{दह\-ना\-द्दा\-र\-णा\-त्तथा}}{\lemma{\skt{दारणात्}}\Afootnote{H; \skt{दारणा}\ A}}}{\lemma{\skt{तथा}}\Afootnote{H; \skt{दपि}\ A}}\skt{~। }\pend
\pstart\versepar \skt{}\edtext{\skt{उप\-घा\-त\-म्वि\-जा\-नी\-या\-त्म\-र्म्मणां}}{\lemma{\skt{विजानीयात्}}\Afootnote{H; \skt{विजानीयान्}\ A}}\skt{ तुल्यलक्षणं~॥}\pend

\pstart\versepar \ledsidenote{\passref{3.6.42}}\edtext{\skt{मर्म्माभिघातस्तु च कश्चिदस्ति}}{\lemma{\skt{मर्म्माभिघातस्तु च कश्}}\Afootnote{H; \skt{मर्माभिघातस्तु न क}\ A}}\skt{ योऽल्पात्ययो वापि निरत्ययो वा~। }\pend
\pstart\versepar \skt{प्रायेण मर्म्म\-स्व\-भि\-ता\-डि\-तास्तु }\edtext{\skt{वैक\-ल्य\-म\-र्च्छ\-न्त्यथ}}{\lemma{\skt{अर्च्छन्त्य्}}\Afootnote{H; \skt{ऋच्छन्त्य्}\ A}}\skt{ वा म्रियन्ते~॥}\edtext{}{\lemma{\skt{चिदस्ति यो}}\Afootnote{H; \skt{श्चिदस्ति यो}\ A}}\pend

\pstart\versepar \ledsidenote{\passref{3.6.43}}\skt{मर्म्माण्यधिष्ठाय हि ये विकारा मूर्च्छन्ति काये विविधा नराणाम्~। }\pend
\pstart\versepar \skt{प्रायेण ते कृच्छ्रतमा भवन्ति नरस्य यत्नैरपि साध्यमानाः~॥~॥}\pend

\pstart \ledsidenote{\passref{3.6.trailer}}\edtext{\skt{इति शारीरे}}{\lemma{\skt{इति शारीरे}}\Afootnote{H; \skt{इति सुश्रु\appbrk{}त\appbrk{}सं\appbrk{}हि\appbrk{}तायां शारीरस्थाने प्रत्ये\appbrk{}क\appbrk{}म\appbrk{}र्म\appbrk{}नि\appbrk{}र्दे\appbrk{}श\appbrk{}शा\appbrk{}रीरं}\ A}}\skt{ }\edtext{\skt{षष्ठोऽध्यायः}}{\lemma{\skt{षष्ठो}}\Afootnote{H; \skt{नाम षष्ठो}\ A}}\skt{~॥ ०~॥}\pend

\endnumbering
\beginnumbering
\eledsection*{Adhyāya 7}
\pstart \ledsidenote{\passref{3.7.1}}\skt{अथातः }\edtext{\skt{सिरा\-व\-र्ण्ण\-न\-वि\-भक्तिं}}{\lemma{\skt{सिरा\appbrk{}व\appbrk{}र्ण्ण\appbrk{}न\appbrk{}वि\appbrk{}भक्तिं}}\Afootnote{H; \skt{सिरा\appbrk{}व\appbrk{}र्ण\appbrk{}वि\appbrk{}भक्ति}\ A}}\skt{ शारी\-र\-म्व्या\-ख्या\-स्यामः~॥}\pend

\pstart \ledsidenote{\passref{3.7.3}}\skt{सप्त }\edtext{\skt{शिराशतानि}}{\lemma{\skt{शिराशतानि}}\Afootnote{H; \skt{सिराशतानि}\ A}}\skt{ भवन्ति~। याभिरिदं शरीरमाराम इव }\edtext{\skt{जलहरणीभिः}}{\lemma{\skt{जलहरणीभिः}}\Afootnote{H; \skt{जलहारिणीभिः}\ A}}\skt{~। केदार इव च }\edtext{\skt{कुल्या\-भि\-रु\-प\-स्नि\-ह्यते}}{\lemma{\skt{उपस्निह्यते}}\Afootnote{H; \skt{उप\appbrk{}स्नि\appbrk{}ह्य\appbrk{}तेऽ\appbrk{}नु\appbrk{}गृ\appbrk{}ह्यते}\ A}}\skt{ }\edtext{\skt{चाकु\-ञ्च\-न\-प्र\-सा\-र\-णा\-भि\-र्व्वि\-शे\-षै\-र्द्द्रु\-म\-प\-त्र\-से\-व\-नी\-ना\-मिव}}{\lemma{\skt{चाकु\appbrk{}ञ्च\appbrk{}न\appbrk{}प्र\appbrk{}सा\appbrk{}र\appbrk{}णा\appbrk{}भिर्}}\Afootnote{H; \skt{चाकु\appbrk{}ञ्च\appbrk{}न\appbrk{}प्र\appbrk{}सा\appbrk{}र\appbrk{}णा\appbrk{}दि\appbrk{}भिर्}\ A}}\skt{ }\edtext{\edtext{\skt{तासा\-म्प्र\-ता\-ना\-म्तासां}}{\lemma{\skt{तासाम्}}\Afootnote{A; \skt{चासाम्}\ H}}}{\lemma{\skt{प्रतानाम्}}\Afootnote{H; \skt{प्रतानाः}\ A}}\skt{ }\edtext{\skt{खलु}}{\lemma{\skt{खलु}}\Afootnote{H; \emph{om.}\ A}}\skt{ }\edtext{\edtext{\skt{नाभिमूलन्ततः}}{\lemma{\skt{नाभिमूलन्}}\Afootnote{H; \skt{नाभिर्मूलं}\ A}}}{\lemma{\skt{ततः}}\Afootnote{H; \skt{ततश्च}\ A}}\skt{ }\edtext{\skt{प्रस\-र\-न्त्यू\-र्ध्व\-म\-ध\-स्ति\-र्यक्च}}{\lemma{\skt{ऊर्ध्वम्}}\Afootnote{A; \skt{उर्ध्वम्}\ H}}\skt{~॥}\pend

\pstart\versepar \ledsidenote{\passref{3.7.4}}\edtext{\skt{भवति}}{\lemma{\skt{भवति}}\Afootnote{H; \skt{भवतश्}\ A}}\skt{ चात्र~॥ }\pend
\pstart\versepar \skt{}\edtext{\skt{यावन्त्यस्तु}}{\lemma{\skt{यावन्त्यस्तु}}\Afootnote{H; \skt{श्लोकौ -- यावत्यस्तु}\ A}}\skt{ सिराः काये सम्भवन्ति }\edtext{\skt{शरीरिणः}}{\lemma{\skt{शरीरिणः}}\Afootnote{H; \skt{शरीरिणाम्}\ A}}\skt{~। }\pend
\pstart\versepar \skt{नाभ्यां }\edtext{\skt{सर्व्वाणि}}{\lemma{\skt{सर्व्वाणि}}\Afootnote{H; \skt{सर्वा नि}\ A}}\skt{ }\edtext{\edtext{\skt{बद्धास्ताः}}{\lemma{\skt{बद्धास्}}\Afootnote{H; \skt{बद्धा}\ A}}}{\lemma{\skt{ताः}}\Afootnote{H; \skt{स्ताः}\ A}}\skt{ }\edtext{\skt{प्रतन्वन्ति}}{\lemma{\skt{प्रतन्वन्ति}}\Afootnote{A; \skt{प्रतन्वति}\ H}}\skt{ समन्ततः~॥}\pend

\pstart \ledsidenote{\passref{3.7.6}}\skt{तासां }\edtext{\skt{खलु}}{\lemma{\skt{खलु}}\Afootnote{H; \emph{om.}\ A}}\skt{ }\edtext{\edtext{\skt{मूल\-सि\-रा\-श्च\-त्वा\-रि\-श\-त्तासु}}{\lemma{\skt{चत्वारिशत्}}\Afootnote{H; \skt{चत्वारिंशत्}\ A}}}{\lemma{\skt{तासु}}\Afootnote{H; \skt{तासां}\ A}}\skt{ }\edtext{\skt{दशवातवहाः}}{\lemma{\skt{दशवातवहाः}}\Afootnote{H; \skt{वातवाहिन्यो दश, पित्तवाहि}\ A}}\skt{~। }\edtext{\skt{दश}}{\lemma{\skt{दश}}\Afootnote{H; \skt{न्यो दश}\ A}}\skt{ }\edtext{\skt{पित्तवहाः}}{\lemma{\skt{पित्तवहाः}}\Afootnote{H; \skt{कफवाहि}\ A}}\skt{~। }\edtext{\skt{दश}}{\lemma{\skt{दश}}\Afootnote{H; \skt{न्यो दश}\ A}}\skt{ }\edtext{\skt{श्लेष्मवहाः}}{\lemma{\skt{श्लेष्मवहाः}}\Afootnote{H; \emph{om.}\ A}}\skt{~। दश }\edtext{\skt{रक्तवहा}}{\lemma{\skt{रक्तवहा}}\Afootnote{H; \skt{रक्तवाहि}\ A}}\skt{ }\edtext{\skt{इति}}{\lemma{\skt{इति}}\Afootnote{H; \skt{न्यः}\ A}}\skt{~। तासां }\edtext{\skt{खलु}}{\lemma{\skt{खलु}}\Afootnote{H; \skt{तु}\ A}}\skt{ }\edtext{\skt{वातवहानां}}{\lemma{\skt{वातवहानां}}\Afootnote{H; \skt{वातवाहिनी}\ A}}\skt{ }\edtext{\skt{सिराणां}}{\lemma{\skt{सिराणां}}\Afootnote{H; \skt{नां}\ A}}\skt{ वातस्थानग ताना\-म्प\-ञ्च\-स\-प्त\-ति\-शतं भवति~। }\edtext{\skt{तावन्त्येव}}{\lemma{\skt{तावन्त्य्}}\Afootnote{H; \skt{तावत्य}\ A}}\skt{ }\edtext{\skt{पित्तवहाः}}{\lemma{\skt{पित्तवहाः}}\Afootnote{\skt{पित्तवाहिन्यः पित्तस्थाने, कफवाहिन्यश्}\ A; \skt{पित्तवहाः पित्त}\ H}}\skt{ }\edtext{\skt{पित्तस्थाने}}{\lemma{\skt{पित्तस्थाने}}\Afootnote{\skt{च कफस्थाने}\ A; \skt{स्थाने}\ H}}\skt{~। }\edtext{\skt{रक्तवहाश्च}}{\lemma{\skt{रक्तवहाश्च}}\Afootnote{H; \skt{रक्तवाहिन्यश्च}\ A}}\skt{ }\edtext{\skt{यकृत्प्लीहाः}}{\lemma{\skt{यकृत्प्लीहाः}}\Afootnote{H; \skt{यकृत्प्ली}\ A}}\skt{~। }\edtext{\skt{कफवहाश्च}}{\lemma{\skt{कफवहाश्च}}\Afootnote{H; \skt{ह्}\ A}}\skt{ }\edtext{\skt{कफस्थाने}}{\lemma{\skt{कफस्थाने}}\Afootnote{H; \skt{नोः}\ A}}\skt{~। एवमेतानि सप्त }\edtext{\skt{शिराशतानि}}{\lemma{\skt{शिराशतानि}}\Afootnote{H; \skt{सिराशतानि}\ A}}\skt{ }\edtext{\skt{भवन्ति}}{\lemma{\skt{भवन्ति}}\Afootnote{H; \emph{om.}\ A}}\skt{~।}\pend

\pstart \ledsidenote{\passref{3.7.7}}\skt{तत्र }\edtext{\edtext{\edtext{\skt{वात\-व\-हा\-ना\-मे\-क\-स्मि\-न्स\-क्थिनि}}{\lemma{\skt{वातवहानाम्}}\Afootnote{H; \skt{वातवाहिन्यः}\ A}}}{\lemma{\skt{एकस्मिन्}}\Afootnote{H; \skt{सिरा एकस्मिन्}\ A}}}{\lemma{\skt{सक्थिनि}}\Afootnote{H; \skt{सक्थ्नि}\ A}}\skt{ पञ्चविंशतिः }\edtext{\skt{सिरा}}{\lemma{\skt{सिरा}}\Afootnote{H; \emph{om.}\ A}}\skt{ }\edtext{\skt{भवन्ति}}{\lemma{\skt{भवन्ति}}\Afootnote{H; \emph{om.}\ A}}\skt{~। }\edtext{\skt{एते\-ने\-त\-र\-सक्ति}}{\lemma{\skt{एते\appbrk{}ने\appbrk{}त\appbrk{}र\appbrk{}सक्ति}}\Afootnote{H; \skt{एते\appbrk{}ने\appbrk{}त\appbrk{}र\appbrk{}सक्थि}\ A}}\skt{ वाहू च व्याख्यातौ~। }\edtext{\skt{कोष्ठे}}{\lemma{\skt{कोष्ठे}}\Afootnote{H; \skt{विशेषतस्तु कोष्ठे}\ A}}\skt{ चतुस्त्रिंशत्तासां }\edtext{\skt{गुद\-मे\-ढ्र\-सं\-श्रिताः}}{\lemma{\skt{गुद\appbrk{}मे\appbrk{}ढ्र\appbrk{}सं\appbrk{}श्रिताः}}\Afootnote{H; \skt{गुदमेढ्राश्रिताः}\ A}}\skt{ }\edtext{\edtext{\skt{श्रोण्यामष्टौ}}{\lemma{\skt{श्रोण्याम्}}\Afootnote{H; \skt{श्रोण्या}\ A}}}{\lemma{\skt{अष्टौ}}\Afootnote{H; \skt{मष्टौ}\ A}}\skt{ द्वे द्वे पार्श्वयोः षट्पृष्ठे }\edtext{\skt{तावन्त्य}}{\lemma{\skt{तावन्त्य}}\Afootnote{H; \skt{तावत्य}\ A}}\skt{ एवोदरे दश वक्षसि~। }\edtext{\skt{जत्रुण}}{\lemma{\skt{जत्रुण}}\Afootnote{H; \skt{एक\appbrk{}च\appbrk{}त्वा\appbrk{}रिं\appbrk{}श\appbrk{}ज्ज\appbrk{}त्रुण}\ A}}\skt{ }\edtext{\edtext{\edtext{\skt{ऊर्द्ध्व\-मे\-क\-च\-त्वा\-रिं\-श\-त्ता\-सा\-मष्टौ}}{\lemma{\skt{एक\appbrk{}च\appbrk{}त्वा\appbrk{}रिं\appbrk{}शत्}}\Afootnote{H; \emph{om.}\ A}}}{\lemma{\skt{तासाम्}}\Afootnote{H; \emph{om.}\ A}}}{\lemma{\skt{अष्टौ}}\Afootnote{H; \emph{om.}\ A}}\skt{ }\edtext{\skt{जिह्वायां}}{\lemma{\skt{जिह्वायां}}\Afootnote{H; \skt{तासां}\ A}}\skt{~। }\edtext{\skt{नव}}{\lemma{\skt{नव}}\Afootnote{H; \skt{च}\ A}}\skt{ }\edtext{\skt{नासायां}}{\lemma{\skt{नासायां}}\Afootnote{H; \skt{तुर्}\ A}}\skt{~। }\edtext{\skt{द्वे}}{\lemma{\skt{द्वे}}\Afootnote{H; \skt{दश}\ A}}\skt{ ग्रीवायां~। }\edtext{\edtext{\skt{हन्वोश्चतस्रः}}{\lemma{\skt{हन्वोश्}}\Afootnote{H; \emph{om.}\ A}}}{\lemma{\skt{चतस्रः}}\Afootnote{H; \emph{om.}\ A}}\skt{~। }\edtext{\skt{नेत्रयोः}}{\lemma{\skt{नेत्रयोः}}\Afootnote{H; \emph{om.}\ A}}\skt{ }\edtext{\edtext{\skt{षट्ललाटे}}{\lemma{\skt{षट्}}\Afootnote{H; \emph{om.}\ A}}}{\lemma{\skt{ललाटे}}\Afootnote{H; \emph{om.}\ A}}\skt{ }\edtext{\skt{तिस्रः}}{\lemma{\skt{तिस्रः}}\Afootnote{H; \emph{om.}\ A}}\skt{ कर्ण्णयोश्चतस्रः~। }\edtext{\skt{द्वे}}{\lemma{\skt{द्वे}}\Afootnote{H; \skt{नव जिह्वायां, षण्नासिकायां, अष्टौ ने}\ A}}\skt{ }\edtext{\skt{शंखयोः}}{\lemma{\skt{शंखयोः}}\Afootnote{H; \skt{त्रयोः}\ A}}\skt{~। }\edtext{\skt{तिस्रः}}{\lemma{\skt{तिस्रः}}\Afootnote{H; \emph{om.}\ A}}\skt{ }\edtext{\skt{शिर\-सी\-त्ये\-व\-मे\-त\-त्प\-ञ्च\-स\-प्त\-ति\-शतं}}{\lemma{\skt{शिरसीत्येवम्}}\Afootnote{H; \skt{एवम्}\ A}}\skt{~। }\edtext{\skt{वातवहानां}}{\lemma{\skt{वातवहानां}}\Afootnote{H; \skt{वातवाहिनीनां}\ A}}\skt{ }\edtext{\skt{शिराणां}}{\lemma{\skt{शिराणां}}\Afootnote{H; \skt{सिराणां}\ A}}\skt{ व्याख्यातं~। }\edtext{\skt{एष}}{\lemma{\skt{एष}}\Afootnote{H; \skt{भवति~। एष}\ A}}\skt{ एव विभागः }\edtext{\skt{शेषा\-णा\-मे\-व\-मे\-तानि}}{\lemma{\skt{एवम्}}\Afootnote{H; \skt{अपि~। विशेषतस्तु पित्तवाहिन्यो नेत्रयोर्दश, कर्णयोर्द्वे; एवं रक्तवहाः कफवहाश्च~। एवम्}\ A}}\skt{ सप्त }\edtext{\skt{शिराशतानि}}{\lemma{\skt{शिराशतानि}}\Afootnote{H; \skt{सिराशतानि सविभागानि}\ A}}\skt{ व्याख्यातानि }\edtext{\skt{भवन्ति}}{\lemma{\skt{भवन्ति}}\Afootnote{H; \emph{om.}\ A}}\skt{~।}\pend

\pstart \ledsidenote{\passref{3.7.8}}\skt{भवन्ति }\edtext{\skt{चात्र}}{\lemma{\skt{चात्र}}\Afootnote{H; \skt{चात्र क्रिया\appbrk{}णा\appbrk{}म\appbrk{}प्र\appbrk{}ती\appbrk{}घा\appbrk{}त\appbrk{}म\appbrk{}मोहं बुद्धिकर्मणाम्~। करो\appbrk{}त्य\appbrk{}न्या\appbrk{}न्गु\appbrk{}णां\appbrk{}श्चापि स्वाः सिराः पवनश्चरन्}\ A}}\skt{~॥}\pend

\pstart\versepar \ledsidenote{\passref{3.7.8.0}}\skt{यदा वायुरदुष्टस्तु सेवते स्ववहाः सिराः~। }\pend
\pstart\versepar \skt{तदा\-स्तु\-व\-ल\-व\-र्ण्णौजः प्रसीदेच्च मनस्तथा~॥}\pend

\pstart\versepar \ledsidenote{\passref{3.7.8.1}}\skt{क्रिया\-ना\-म\-प्र\-ति\-घा\-त\-म\-मोहं बुद्धिकर्म्मणां~। }\pend
\pstart\versepar \skt{करो\-त्य\-न्या\-न्गु\-णां\-श्चापि सिराः पव\-न\-स्व\-श्च\-रन्~॥}\pend

\pstart\versepar \ledsidenote{\passref{3.7.9}}\skt{यदा तु कुपितो वायुः }\edtext{\skt{स्ताः}}{\lemma{\skt{स्ताः}}\Afootnote{H; \skt{स्वाः}\ A}}\skt{ }\edtext{\skt{शिरा}}{\lemma{\skt{शिरा}}\Afootnote{H; \skt{सिराः}\ A}}\skt{ प्रतिपद्यते~। }\pend
\pstart\versepar \skt{तदास्य }\edtext{\skt{रोगा}}{\lemma{\skt{रोगा}}\Afootnote{H; \skt{विविधा रोगा}\ A}}\skt{ जायन्ते }\edtext{\skt{विविधा}}{\lemma{\skt{विविधा}}\Afootnote{H; \emph{om.}\ A}}\skt{ वातसम्भवाः~॥}\pend

\pstart\versepar \ledsidenote{\passref{3.7.9.add}}\skt{यदा त्वकुपितम्पित्तं सेवते स्ववहाः सिराः~। }\pend
\pstart\versepar \skt{अव्या\-प\-न्न\-स्त\-दा\-ग्निस्तु सम्यक्चरति भोजनं~॥}\pend

\pstart\versepar \ledsidenote{\passref{3.7.10}}\edtext{\skt{भ्राजि\-ष्णु\-ता\-म\-न्न\-रु\-चि\-म\-ग्ने\-र्द्दी\-प्ति\-म\-रो\-गतां}}{\lemma{\skt{अग्नेर्}}\Afootnote{H; \skt{अग्नि}\ A}}\skt{~। }\pend
\pstart\versepar \skt{}\edtext{\edtext{\edtext{\edtext{\skt{करो\-त्य\-न्या\-न्गु\-णां\-श्चापि}}{\lemma{\skt{करोत्य्}}\Afootnote{H; \skt{संसर्पत्स्वाः सिराः पित्तं कुर्याच्}\ A}}}{\lemma{\skt{अन्यान्}}\Afootnote{H; \skt{चान्यान्}\ A}}}{\lemma{\skt{गुणांश्}}\Afootnote{H; \skt{गुणा}\ A}}}{\lemma{\skt{चापि}}\Afootnote{H; \skt{नपि}\ A}}\skt{ }\edtext{\edtext{\edtext{\skt{पित्त\-म\-त्म\-सि\-रा\-श्च\-रन्}}{\lemma{\skt{पित्तम्}}\Afootnote{H; \emph{om.}\ A}}}{\lemma{\skt{अत्मसिराश्}}\Afootnote{H; \emph{om.}\ A}}}{\lemma{\skt{चरन्}}\Afootnote{H; \emph{om.}\ A}}\skt{~॥}\pend

\pstart\versepar \ledsidenote{\passref{3.7.11}}\skt{यदा }\edtext{\skt{तु}}{\lemma{\skt{तु}}\Afootnote{H; \skt{प्र}\ A}}\skt{ कुपितम्पित्तं }\edtext{\skt{ताः}}{\lemma{\skt{ताः}}\Afootnote{H; \skt{सेवते स्ववहाः}\ A}}\skt{ सिराः }\edtext{\skt{प्रतिपद्यते}}{\lemma{\skt{प्रतिपद्यते}}\Afootnote{H; \emph{om.}\ A}}\skt{~। }\pend
\pstart\versepar \skt{तदास्य }\edtext{\skt{रोगा}}{\lemma{\skt{रोगा}}\Afootnote{H; \skt{विविधा रोगा}\ A}}\skt{ जायन्ते }\edtext{\skt{विविधा}}{\lemma{\skt{विविधा}}\Afootnote{H; \emph{om.}\ A}}\skt{ पित्तसम्भवाः~॥}\pend

\pstart\versepar \ledsidenote{\passref{3.7.11.1}}\skt{यदा त्वकुपितः श्लेष्मा स्वाः सिराः संप्रपद्यते~। }\pend
\pstart\versepar \skt{आशयाः सन्धयश्चैव वर्त्तन्तेस्य निरामयाः~॥}\pend

\pstart\versepar \ledsidenote{\passref{3.7.12}}\edtext{\edtext{\skt{स्नेहमङ्गेषु}}{\lemma{\skt{स्नेहम्}}\Afootnote{H; \skt{स्नेह}\ A}}}{\lemma{\skt{अङ्गेषु}}\Afootnote{H; \skt{मङ्गेषु}\ A}}\skt{ सन्धीनां }\edtext{\edtext{\skt{स्थैर्य\-म्व\-ल\-म\-दी\-नतां}}{\lemma{\skt{वलम्}}\Afootnote{H; \skt{बल}\ A}}}{\lemma{\skt{अदीनतां}}\Afootnote{H; \skt{मुदीर्णताम्}\ A}}\skt{~। }\pend
\pstart\versepar \skt{}\edtext{\edtext{\skt{करो\-त्य\-न्या\-न्गु\-णां\-श्चापि}}{\lemma{\skt{करोत्य्}}\Afootnote{H; \skt{करोत्}\ A}}}{\lemma{\skt{अन्यान्}}\Afootnote{H; \skt{यन्यान्}\ A}}\skt{ }\edtext{\skt{वलास}}{\lemma{\skt{वलास}}\Afootnote{H; \skt{बलासः}\ A}}\skt{ }\edtext{\skt{स्वसिराश्चरन्}}{\lemma{\skt{स्वसिराश्}}\Afootnote{H; \skt{स्वाः सिराश्}\ A}}\skt{~॥}\pend

\pstart\versepar \ledsidenote{\passref{3.7.13}}\skt{यदा तु कुपितः श्लेष्मा }\edtext{\skt{ताः}}{\lemma{\skt{ताः}}\Afootnote{H; \skt{स्वाः}\ A}}\skt{ सिराः प्रतिपद्यते~। }\pend
\pstart\versepar \skt{तदास्य }\edtext{\skt{रोगा}}{\lemma{\skt{रोगा}}\Afootnote{H; \skt{विविधा रोगा}\ A}}\skt{ जायन्ते }\edtext{\skt{विविधाः}}{\lemma{\skt{विविधाः}}\Afootnote{H; \emph{om.}\ A}}\skt{ श्लेष्मसम्भवाः~॥}\pend

\pstart\versepar \ledsidenote{\passref{3.7.13.1}}\skt{यदा त्वकुपितं रक्तं सेवते स्ववहाः सिराः~। }\pend
\pstart\versepar \skt{तदास्य सम्यग्जानाति स्पर्शानां तु शुभाशुभं~॥}\pend

\pstart\versepar \ledsidenote{\passref{3.7.14}}\edtext{\skt{वर्ण्णप्रसादनं}}{\lemma{\skt{वर्ण्णप्रसादनं}}\Afootnote{H; \emph{om.}\ A}}\skt{ }\edtext{\skt{स्थैर्यं}}{\lemma{\skt{स्थैर्यं}}\Afootnote{H; \emph{om.}\ A}}\skt{ धातूनां }\edtext{\edtext{\skt{पुष्टिमेव}}{\lemma{\skt{पुष्टिम्}}\Afootnote{H; \skt{पूरणं वर्णं स्पर्श\appbrk{}ज्ञा\appbrk{}न\appbrk{}म\appbrk{}सं\appbrk{}श\appbrk{}यम्}\ A}}}{\lemma{\skt{एव}}\Afootnote{\skt{स्वाः सिराः संच}\ A; \skt{एव च}\ H}}\skt{ च~। }\pend
\pstart\versepar \skt{}\edtext{\edtext{\edtext{\skt{करो\-त्य\-न्या\-न्गु\-णां\-श्चापि}}{\lemma{\skt{अन्यान्}}\Afootnote{\skt{चान्यान्}\ A; \skt{न्यान्}\ H}}}{\lemma{\skt{गुणांश्}}\Afootnote{H; \skt{गुणा}\ A}}}{\lemma{\skt{चापि}}\Afootnote{H; \skt{नपि}\ A}}\skt{ }\edtext{\edtext{\edtext{\skt{रक्त\-मा\-त्म\-सि\-रा\-श्च\-रन्}}{\lemma{\skt{रक्तम्}}\Afootnote{H; \emph{om.}\ A}}}{\lemma{\skt{आत्मसिराश्}}\Afootnote{H; \emph{om.}\ A}}}{\lemma{\skt{चरन्}}\Afootnote{H; \emph{om.}\ A}}\skt{~॥}\edtext{}{\lemma{\skt{च~। करोत्य्}}\Afootnote{\skt{रद्रक्तं कुर्याच्}\ A; \skt{करोत्य}\ H}}\pend

\pstart\versepar \ledsidenote{\passref{3.7.15}}\skt{यदा तु कुपितं रक्तं }\edtext{\skt{ताः}}{\lemma{\skt{ताः}}\Afootnote{H; \skt{सेवते स्ववहाः}\ A}}\skt{ सिराः }\edtext{\skt{प्रतिपद्यते}}{\lemma{\skt{प्रतिपद्यते}}\Afootnote{H; \emph{om.}\ A}}\skt{~। }\pend
\pstart\versepar \skt{तदास्य }\edtext{\skt{रोगा}}{\lemma{\skt{रोगा}}\Afootnote{H; \skt{विविधा रोगा}\ A}}\skt{ जायन्ते }\edtext{\skt{विविधा}}{\lemma{\skt{विविधा}}\Afootnote{H; \emph{om.}\ A}}\skt{ रक्तसम्भवाः~॥}\pend

\pstart\versepar \ledsidenote{\passref{3.7.15.1}}\skt{अप्रदुष्टाः शरीराणि धारयन्ति रसैर्हितैः~। }\pend
\pstart\versepar \skt{कल्पन्ते च विकाराय प्रदुष्टा नात्र संशयः~॥}\pend

\pstart\versepar \ledsidenote{\passref{3.7.16}}\skt{न हि }\edtext{\skt{वातसिराः}}{\lemma{\skt{वातसिराः}}\Afootnote{H; \skt{वातं सिराः}\ A}}\skt{ }\edtext{\skt{कश्चित्न}}{\lemma{\skt{कश्चित्न}}\Afootnote{H; \skt{काश्चिन्न}\ A}}\skt{ पित्तं केवलं }\edtext{\skt{सिराः}}{\lemma{\skt{सिराः}}\Afootnote{H; \skt{तथा}\ A}}\skt{~। }\pend
\pstart\versepar \skt{}\edtext{\skt{श्लेष्माणञ्च}}{\lemma{\skt{श्लेष्माणञ्च}}\Afootnote{H; \skt{श्लेष्माणं वा}\ A}}\skt{ }\edtext{\edtext{\skt{वहन्त्येताः}}{\lemma{\skt{वहन्त्य्}}\Afootnote{H; \skt{वहन्त्}\ A}}}{\lemma{\skt{एताः}}\Afootnote{H; \skt{येता}\ A}}\skt{ }\edtext{\skt{सर्व्वाः}}{\lemma{\skt{सर्व्वाः}}\Afootnote{H; \skt{अतः}\ A}}\skt{ सर्व्ववहाः }\edtext{\skt{सिराः}}{\lemma{\skt{सिराः}}\Afootnote{H; \skt{स्मृताः}\ A}}\skt{~॥}\pend

\pstart\versepar \ledsidenote{\passref{3.7.17}}\edtext{\skt{प्रदुष्टानान्तु}}{\lemma{\skt{प्रदुष्टानान्तु}}\Afootnote{H; \skt{प्रदुष्टानां हि}\ A}}\skt{ दोषाणां }\edtext{\skt{मूर्च्छि\-ता\-ना\-म्प्र\-धा\-न\-तान्}}{\lemma{\skt{प्रधानतान्}}\Afootnote{H; \skt{प्रधावताम्}\ A}}\skt{~। }\pend
\pstart\versepar \skt{ध्रुव\-मु\-न्मा\-र्ग्ग\-ग\-मनं }\edtext{\skt{सर्व्वाः}}{\lemma{\skt{सर्व्वाः}}\Afootnote{H; \skt{अतः}\ A}}\skt{ }\edtext{\skt{सर्व्ववहा}}{\lemma{\skt{सर्व्ववहा}}\Afootnote{H; \skt{सर्ववहाः}\ A}}\skt{ }\edtext{\skt{ह्यतः}}{\lemma{\skt{ह्यतः}}\Afootnote{H; \skt{स्मृताः}\ A}}\skt{~॥}\pend

\pstart\versepar \ledsidenote{\passref{3.7.18}}\skt{तत्रारुणा }\edtext{\skt{वातवहा}}{\lemma{\skt{वातवहा}}\Afootnote{H; \skt{वातवहाः पूर्यन्ते वायुना सिराः~। पित्ता}\ A}}\skt{ }\edtext{\skt{नीलाः}}{\lemma{\skt{नीलाः}}\Afootnote{H; \skt{दुष्णाश्च नीलाश्}\ A}}\skt{ }\edtext{\skt{पित्तवहाः सिराः}}{\lemma{\skt{पित्तवहाः सिराः}}\Afootnote{H; \skt{च, शीता गौर्यः स्थिराः}\ A}}\skt{~। }\pend
\pstart\versepar \skt{… }\pend
\pstart\versepar \skt{}\edtext{\edtext{\skt{असृग्वहाश्च}}{\lemma{\skt{{\symfont ‡}}}\Afootnote{H; \skt{कफा}\ A}}\skt{ रोहिण्यो}}{\lemma{\skt{असृग्वहाश्च}}\Afootnote{H; \skt{तसृग्वहास्तु}\ A}}\skt{ }\edtext{\skt{गौर्यः श्लेष्मवहाः}}{\lemma{\skt{गौर्यः}}\Afootnote{H; \skt{यः}\ A}}\skt{ }\edtext{\skt{सिराः~॥}}{\lemma{\skt{सिराः}}\Afootnote{H; \skt{सिरा नात्युष्णशीतलाः}\ A}}\edtext{}{\lemma{\skt{रोहिण्यो}}\Afootnote{H; \skt{रोहिण्}\ A}}\edtext{}{\lemma{\skt{श्लेष्मवहाः}}\Afootnote{H; \emph{om.}\ A}}\pend

\pstart \ledsidenote{\passref{3.7.19}}\skt{अत ऊर्द्ध्वं }\edtext{\skt{प्रवक्ष्यामि }\edtext{\skt{यास्त्ववध्याः}}{\lemma{\skt{त्ववध्याः}}\Afootnote{H; \skt{न विध्येद्याः}\ A}}}{\lemma{\skt{प्रवक्ष्यामि यास्}}\Afootnote{H; \skt{प्रवक्ष्यामि}\ A}}\skt{ सिरा }\edtext{\skt{नृणां}}{\lemma{\skt{नृणां}}\Afootnote{H; \skt{भिषक्~। वैकल्यं मरणं चापि व्यधात्तासां ध्रुवं भवेत्}\ A}}\skt{~।}\pend

\pstart\versepar \ledsidenote{\passref{3.7.21}}\skt{शाखासु }\edtext{\skt{षोडशा}}{\lemma{\skt{षोडशा}}\Afootnote{H; \skt{षोड}\ A}}\skt{ }\edtext{\skt{व्यध्याः}}{\lemma{\skt{व्यध्याः}}\Afootnote{H; \skt{श सिराः}\ A}}\skt{ कोष्ठे द्वात्रिंशदेव तु~। }\pend
\pstart\versepar \skt{}\edtext{\skt{जत्रोरूर्द्ध्वन्तु}}{\lemma{\skt{जत्रोरूर्द्ध्वन्तु}}\Afootnote{H; \skt{पञ्चा\appbrk{}श\appbrk{}ज्ज\appbrk{}त्रु\appbrk{}ण\appbrk{}श्चो\appbrk{}र्ध्वम्}\ A}}\skt{ }\edtext{\skt{पञ्चाशदव्यध्याः}}{\lemma{\skt{पञ्चाशदव्यध्याः}}\Afootnote{H; \skt{अव्यध्याः}\ A}}\skt{ परिकीर्त्तिताः~।}\pend

\pstart \ledsidenote{\passref{3.7.19cd}}\skt{वैकल्यं मरणम्वाशु व्यध्यात्तासां ध्रुवम्भवेत्~॥}\pend

\pstart \ledsidenote{\passref{3.7.22}}\skt{तत्र }\edtext{\skt{सिराः}}{\lemma{\skt{सिराः}}\Afootnote{H; \skt{सिरा}\ A}}\skt{ }\edtext{\skt{शत\-मे\-क\-स्मि\-न्स\-क्थिनि}}{\lemma{\skt{सक्थिनि}}\Afootnote{H; \skt{सक्थ्नि}\ A}}\skt{ }\edtext{\skt{भवन्ति}}{\lemma{\skt{भवन्ति}}\Afootnote{H; \skt{भवति}\ A}}\skt{~। तासां }\edtext{\skt{जलधारा}}{\lemma{\skt{जलधारा}}\Afootnote{H; \skt{जालधरा}\ A}}\skt{ त्वेका }\edtext{\skt{तिस्र\-श्चा\-भ्य\-न्त\-रतः}}{\lemma{\skt{चाभ्यन्तरतः}}\Afootnote{H; \skt{चाभ्यन्तराः}\ A}}\skt{~। तत्रोर्व्वीसंज्ञे द्वे लोहिताक्ष }\edtext{\skt{संज्ञका तास्त्वव्यध्याः}}{\lemma{\skt{संज्ञका तास्}}\Afootnote{H; \skt{संज्ञा चैका, तास्}\ A}}\skt{~। एते\-ने\-त\-र\-सक्थि वाहू च व्याख्यातौ~। }\edtext{\edtext{\skt{एव\-म\-श\-स्त्र\-कृत्याः}}{\lemma{\skt{एवम्}}\Afootnote{H; \skt{एव}\ A}}}{\lemma{\skt{अशस्त्रकृत्याः}}\Afootnote{H; \skt{मशस्त्रकृत्याः}\ A}}\skt{ षोडशशाखासु~। }\edtext{\skt{श्रोण्यां}}{\lemma{\skt{श्रोण्यां}}\Afootnote{H; \emph{om.}\ A}}\skt{ }\edtext{\skt{द्वात्रिंशत्}}{\lemma{\skt{द्वात्रिंशत्}}\Afootnote{H; \skt{द्वात्रिंशच्छ्रोण्यां}\ A}}\skt{~। तासा\-म\-ष्टा\-व\-श\-स्त्र\-कृत्याः~। द्वे द्वे }\edtext{\edtext{\skt{वंक्ष\-ण\-चि\-टि\-प\-यो\-रेवं}}{\lemma{\skt{वंक्ष\appbrk{}ण\appbrk{}चि\appbrk{}टि\appbrk{}प\appbrk{}योर्}}\Afootnote{H; \skt{विटपयोः}\ A}}}{\lemma{\skt{एवं}}\Afootnote{H; \emph{om.}\ A}}\skt{ कटी\-क\-त\-रु\-ण\-योश्च~। अष्टा\-व\-ष्टा\-वे\-कै\-क\-स्मि\-न्पार्श्वे }\edtext{\skt{तासामेकैक}}{\lemma{\skt{एकैक}}\Afootnote{H; \skt{एकैकाम्}\ A}}\skt{ }\edtext{\skt{ऊर्द्ध्वगे}}{\lemma{\skt{ऊर्द्ध्वगे}}\Afootnote{H; \skt{ऊर्ध्व}\ A}}\skt{ }\edtext{\skt{उभयतो}}{\lemma{\skt{उभयतो}}\Afootnote{H; \emph{om.}\ A}}\skt{ }\edtext{\skt{द्वे}}{\lemma{\skt{द्वे}}\Afootnote{H; \emph{om.}\ A}}\skt{ }\edtext{\skt{पार्श्वे}}{\lemma{\skt{पार्श्वे}}\Afootnote{H; \emph{om.}\ A}}\skt{ }\edtext{\skt{सन्धिसंश्रिते च}}{\lemma{\skt{सन्धिसंश्रिते च}}\Afootnote{H; \skt{गां}\ A}}\skt{ परिहरेत्~। }\edtext{\skt{चतुर्व्विंशतिः}}{\lemma{\skt{चतुर्व्विंशतिः}}\Afootnote{H; \skt{पार्श्वसन्धिगते च द्वे; चतस्रो विंशतिश्च}\ A}}\skt{ पृष्ठे }\edtext{\edtext{\skt{पृष्ठ\-वं\-श\-मु\-भ\-यतः}}{\lemma{\skt{पृष्ठवंशम्}}\Afootnote{H; \skt{पृष्ठवंश}\ A}}}{\lemma{\skt{उभयतः}}\Afootnote{H; \skt{मुभयतः}\ A}}\skt{~। }\edtext{\skt{तासां}}{\lemma{\skt{तासां}}\Afootnote{H; \skt{तासा}\ A}}\skt{ }\edtext{\skt{वृहती}}{\lemma{\skt{वृहती}}\Afootnote{H; \emph{om.}\ A}}\skt{ }\edtext{\skt{संज्ञे}}{\lemma{\skt{संज्ञे}}\Afootnote{H; \emph{om.}\ A}}\skt{ }\edtext{\skt{ऊर्द्ध्वगामिन्ये}}{\lemma{\skt{ऊर्द्ध्वगामिन्ये}}\Afootnote{H; \skt{मूर्ध्वगामिन्यौ द्वे}\ A}}\skt{ द्वे परिहरेत्~। }\edtext{\skt{तावन्त्य}}{\lemma{\skt{तावन्त्य}}\Afootnote{H; \skt{बृहतीसिरे; तावत्य}\ A}}\skt{ एवोदरे तासां मेढ्रोपरि रोम\-रा\-जी\-मु\-भ\-यतो द्वे द्वे परिहरेत्~। }\edtext{\skt{चत्वा\-रि\-शं\-द्व\-क्षसि}}{\lemma{\skt{चत्वारिशंद्}}\Afootnote{H; \skt{चत्वारिंशद्}\ A}}\skt{ तासां चतु\-र्द्द\-शा\-श\-स्त्र\-कृत्या }\edtext{\skt{हृदि}}{\lemma{\skt{हृदि}}\Afootnote{H; \skt{हृदये}\ A}}\skt{ }\edtext{\skt{द्वे}}{\lemma{\skt{द्वे}}\Afootnote{H; \skt{द्वे, द्वे द्वे}\ A}}\skt{ स्तनमूले }\edtext{\skt{संश्रिते तु}}{\lemma{\skt{संश्रिते तु}}\Afootnote{H; \emph{om.}\ A}}\skt{ }\edtext{\skt{द्वे}}{\lemma{\skt{द्वे}}\Afootnote{H; \emph{om.}\ A}}\skt{ }\edtext{\skt{द्वे}}{\lemma{\skt{द्वे}}\Afootnote{H; \emph{om.}\ A}}\skt{ }\edtext{\skt{परिहरेत्}}{\lemma{\skt{परिहरेत्}}\Afootnote{H; \emph{om.}\ A}}\skt{~। }\edtext{\edtext{\edtext{\skt{स्तन\-मु\-भ\-य\-त\-श्च\-तस्रः}}{\lemma{\skt{स्तनम्}}\Afootnote{H; \emph{om.}\ A}}}{\lemma{\skt{उभयतश्}}\Afootnote{H; \emph{om.}\ A}}}{\lemma{\skt{चतस्रः}}\Afootnote{H; \skt{स्त}\ A}}\skt{ }\edtext{\skt{परिहरेत्}}{\lemma{\skt{परिहरेत्}}\Afootnote{H; \skt{नरोहि}\ A}}\skt{~। }\edtext{\skt{द्वे}}{\lemma{\skt{द्वे}}\Afootnote{H; \emph{om.}\ A}}\skt{ }\edtext{\skt{द्वे}}{\lemma{\skt{द्वे}}\Afootnote{H; \emph{om.}\ A}}\skt{ }\edtext{\skt{चापालापे}}{\lemma{\skt{चापालापे}}\Afootnote{H; \skt{तापलाप}\ A}}\skt{ }\edtext{\edtext{\skt{स्तम्भेष्वित्ये}}{\lemma{\skt{स्तम्भेष्व्}}\Afootnote{H; \skt{स्तम्भेषू}\ A}}\skt{ व}}{\lemma{\skt{इत्ये व}}\Afootnote{H; \skt{भयतोऽष्टौ; एवं}\ A}}\skt{ }\edtext{\skt{द्वात्रिं\-श\-द\-श\-स्त्र\-कृत्याः}}{\lemma{\skt{द्वात्रिं\appbrk{}श\appbrk{}द\appbrk{}श\appbrk{}स्त्र\appbrk{}कृत्याः}}\Afootnote{H; \skt{द्वात्रिं\appbrk{}श\appbrk{}द\appbrk{}श\appbrk{}स्त्र\appbrk{}कृत्याः}\ A}}\skt{~। }\edtext{\skt{श्रोणी}}{\lemma{\skt{श्रोणी}}\Afootnote{H; \emph{om.}\ A}}\skt{ पृष्ठोदरोरःसु भवन्ति }\edtext{\skt{चतिःषष्टिशतं}}{\lemma{\skt{चतिःषष्टिशतं}}\Afootnote{\skt{चतुःषष्टं सिराशतं जत्रुण ऊर्ध्वं}\ A; \skt{चतिःषष्टिशतं मूर्}\ H}}\skt{ }\edtext{\skt{मूर्द्धनि तासां}}{\lemma{\skt{मूर्द्धनि तासां}}\Afootnote{\skt{भवति; तत्र षट्पञ्}\ A; \skt{द्धनि तासां}\ H}}\skt{ }\edtext{\skt{विंशतिः}}{\lemma{\skt{विंशतिः}}\Afootnote{H; \skt{चाशच्छ्}\ A}}\skt{~। }\edtext{\skt{शिरोधरायां}}{\lemma{\skt{शिरोधरायां}}\Afootnote{H; \skt{इरोधरायां}\ A}}\skt{ }\edtext{\edtext{\skt{तासामष्टौ}}{\lemma{\skt{तासाम्}}\Afootnote{H; \skt{तासा}\ A}}}{\lemma{\skt{अष्टौ}}\Afootnote{H; \skt{मष्टौ}\ A}}\skt{ चतस्रश्च }\edtext{\skt{मर्म्मसंज्ञाः}}{\lemma{\skt{मर्म्मसंज्ञाः}}\Afootnote{H; \skt{मर्मसंज्ञाः}\ A}}\skt{~। }\edtext{\skt{द्वे}}{\lemma{\skt{द्वे}}\Afootnote{H; \skt{परिहरेत्, द्वे}\ A}}\skt{ }\edtext{\skt{कृका\-टि\-क\-यो\-स्ता\-वन्त्य}}{\lemma{\skt{तावन्त्य}}\Afootnote{H; \emph{om.}\ A}}\skt{ }\edtext{\skt{एव}}{\lemma{\skt{एव}}\Afootnote{H; \skt{द्वे}\ A}}\skt{ विधुरयोः }\edtext{\edtext{\skt{हन्वोरुभयतो}}{\lemma{\skt{हन्वोर्}}\Afootnote{H; \skt{एवं ग्रीवायां षोडशाव्यध्याः; हन्वोर्}\ A}}}{\lemma{\skt{उभयतो}}\Afootnote{H; \skt{उभयतो}\ A}}\skt{ }\edtext{\skt{अष्टावष्टौ}}{\lemma{\skt{अष्टावष्टौ}}\Afootnote{H; \skt{ऽष्टावष्टौ}\ A}}\skt{ तासां तु }\edtext{\skt{सन्धिनिवंधन्यौ}}{\lemma{\skt{सन्धिनिवंधन्यौ}}\Afootnote{H; \skt{सन्धिधमन्यौ}\ A}}\skt{ द्वे द्वे परिहरेत्~। }\edtext{\skt{षट्त्रिंशज्जिह्वायां}}{\lemma{\skt{षट्त्रिंशज्}}\Afootnote{H; \skt{षट्त्रिंशज्}\ A}}\skt{ }\edtext{\edtext{\skt{तासामधः}}{\lemma{\skt{तासाम्}}\Afootnote{H; \skt{तासा}\ A}}}{\lemma{\skt{अधः}}\Afootnote{H; \skt{मधः}\ A}}\skt{ षोड\-शा\-श\-स्त्र\-कृत्या }\edtext{\skt{रसवाहिन्यः}}{\lemma{\skt{रसवाहिन्यः}}\Afootnote{H; \skt{रसवहे द्वे, वाग्वहे च द्वे; द्विर्द्वा}\ A}}\skt{~। }\edtext{\skt{चतु\-व्विं\-श\-ति\-र्न्ना\-सायां}}{\lemma{\skt{चतुव्विंशतिर्}}\Afootnote{H; \skt{दश}\ A}}\skt{ }\edtext{\skt{तासा\-मौ\-प\-ना\-स्या\-श्च\-तस्रः}}{\lemma{\skt{औपनास्याश्}}\Afootnote{H; \skt{औपनासिक्यश्}\ A}}\skt{ }\edtext{\skt{परिहरे}}{\lemma{\skt{परिहरे}}\Afootnote{H; \skt{परिहरेत्}\ A}}\skt{ तासामेव च }\edtext{\edtext{\skt{तालुन्येका}}{\lemma{\skt{तालुन्य्}}\Afootnote{H; \skt{तालुन्}\ A}}}{\lemma{\skt{एका}}\Afootnote{H; \skt{येकां मृदा}\ A}}\skt{ }\edtext{\skt{षट्त्रिं\-श\-दु\-भ\-यो\-र्न्ने\-त्रयोः}}{\lemma{\skt{षट्त्रिं\appbrk{}श\appbrk{}दु\appbrk{}भ\appbrk{}योर्}}\Afootnote{H; \skt{वुद्देशे; अष्ट\appbrk{}त्रिं\appbrk{}श\appbrk{}दु\appbrk{}भ\appbrk{}योर्}\ A}}\skt{~। }\edtext{\skt{तासा\-मे\-कै\-का\-म\-पा\-ङ्गयोः}}{\lemma{\skt{अपाङ्गयोः}}\Afootnote{H; \skt{अपाङ्गयोः परिहरेत्; कर्णयोर्}\ A}}\skt{ }\edtext{\skt{परिहरेत्}}{\lemma{\skt{परिहरेत्}}\Afootnote{H; \skt{दश, तासां शब्द\appbrk{}वा\appbrk{}हि\appbrk{}नी\appbrk{}मे\appbrk{}कैकां परिहरेत्; नासानेत्}\ A}}\skt{~। }\edtext{\skt{ललाटे}}{\lemma{\skt{ललाटे}}\Afootnote{H; \skt{रगतास्तु ललाटे}\ A}}\skt{ }\edtext{\skt{षष्टिस्तासां}}{\lemma{\skt{षष्टिस्तासां}}\Afootnote{H; \skt{षष्टिः, तासां}\ A}}\skt{ }\edtext{\skt{केशा\-न्ता\-नु\-ग\-ता\-श्च\-तस्रः}}{\lemma{\skt{चतस्रः}}\Afootnote{H; \skt{चतस्}\ A}}\skt{ }\edtext{\skt{परिहरेत्~। }\edtext{\skt{आव\-र्त्त\-यो\-रे\-कै\-का\-स्थ\-प\-ण्या\-श्चैका}}{\lemma{\skt{एकै\appbrk{}का\appbrk{}स्थ\appbrk{}प\appbrk{}ण्याश्}}\Afootnote{H; \skt{एकैका, स्थपन्यां}\ A}}}{\lemma{\skt{परिहरेत्~। आवर्त्तयोर्}}\Afootnote{H; \skt{र, आवर्तयोर्}\ A}}\skt{ }\edtext{\skt{परिहर्त्तव्याः}}{\lemma{\skt{परिहर्त्तव्याः}}\Afootnote{H; \skt{परिहर्तव्}\ A}}\skt{ }\edtext{\skt{कर्ण्णयोः पञ्च}}{\lemma{\skt{कर्ण्णयोः पञ्च}}\Afootnote{H; \emph{om.}\ A}}\skt{ }\edtext{\skt{पञ्च}}{\lemma{\skt{पञ्च}}\Afootnote{H; \emph{om.}\ A}}\skt{~। }\edtext{\skt{तासां}}{\lemma{\skt{तासां}}\Afootnote{H; \skt{या}\ A}}\skt{ }\edtext{\edtext{\skt{शब्द\-व\-हा\-मे\-कैकां}}{\lemma{\skt{शब्दवहाम्}}\Afootnote{H; \emph{om.}\ A}}}{\lemma{\skt{एकैकां}}\Afootnote{H; \emph{om.}\ A}}\skt{ }\edtext{\skt{परिहरेत्}}{\lemma{\skt{परिहरेत्}}\Afootnote{H; \emph{om.}\ A}}\skt{~। }\edtext{\skt{शंखयोश्च}}{\lemma{\skt{शंखयोश्च}}\Afootnote{H; \skt{शङ्खयोर्}\ A}}\skt{ }\edtext{\skt{तावन्त्य}}{\lemma{\skt{तावन्त्य}}\Afootnote{H; \skt{द}\ A}}\skt{ }\edtext{\skt{एव}}{\lemma{\skt{एव}}\Afootnote{H; \skt{श}\ A}}\skt{ तासां शंख\-स\-न्धि\-ग\-ता\-मे\-कै\-का\-म्प\-रि\-ह\-रेत्~। }\edtext{\skt{द्वादशमूर्द्धनि}}{\lemma{\skt{द्वादशमूर्द्धनि}}\Afootnote{H; \skt{द्वादश मूर्ध्नि}\ A}}\skt{ तासामुत्क्षेपयोः~। द्वे }\edtext{\skt{द्वे}}{\lemma{\skt{द्वे}}\Afootnote{H; \skt{परि}\ A}}\skt{ }\edtext{\skt{सीम}}{\lemma{\skt{सीम}}\Afootnote{H; \skt{हरेत्}\ A}}\skt{ }\edtext{\edtext{\skt{मन्ते\-ष्वे\-कै\-का\-म\-धि\-पतौ}}{\lemma{\skt{मन्तेष्व्}}\Afootnote{H; \skt{सीमन्तेष्व्}\ A}}}{\lemma{\skt{एकैकाम्}}\Afootnote{H; \skt{एकैकाम्, एकाम्}\ A}}\skt{ }\edtext{\skt{परिहरेदिति}}{\lemma{\skt{परिहरेदिति}}\Afootnote{H; \skt{इति; एव\appbrk{}म\appbrk{}श\appbrk{}स्त्र\appbrk{}कृत्याः पञ्चाशज्जत्रुण ऊर्ध्वमिति}\ A}}\skt{~॥}\pend

\pstart\versepar \ledsidenote{\passref{3.7.22.add1}}\skt{भवन्ति चात्र श्लोकाः~॥ }\pend
\pstart\versepar \skt{सिराशतानां सप्तानां शरीरेषु शरीरिणां~। }\pend
\pstart\versepar \skt{अशस्त्रकृत्या नवतिस्तथाष्टौ च विनिर्द्दिशेत्~।~॥}\pend

\pstart\versepar \ledsidenote{\passref{3.7.22.add2}}\skt{स्पन्दिन्यौ जलधारिण्यौ याश्च मर्म्म समाश्रिताः~। }\pend
\pstart\versepar \skt{अव्यध्यास्ता }\edtext{\edtext{\skt{विजा\-नी\-या\-त्स्ना\-यु\-स\-न्धि\-ग\-ताश्च}}{\lemma{\skt{स्नायु\appbrk{}स\appbrk{}न्धि\appbrk{}ग\appbrk{}ताश्}}\Afootnote{\skt{स्नायुसन्धिगता}\ H}}}{\lemma{\skt{च}}\Afootnote{\skt{श्च}\ H}}\skt{ याः~॥}\pend

\pstart \ledsidenote{\passref{3.7.22.add3}}\skt{वैक\-ल्प\-म\-र\-ण\-ञ्चापि व्यधादासां ध्रुवं भवेत्~।}\pend

\pstart\versepar \ledsidenote{\passref{3.7.23}}\edtext{\edtext{\skt{व्याप्नुयाद्वाखिलं}}{\lemma{\skt{व्याप्नुयाद्}}\Afootnote{H; \skt{भवति चात्र -- व्याप्नुवन्त्य्}\ A}}}{\lemma{\skt{वाखिलं}}\Afootnote{H; \skt{अभितो}\ A}}\skt{ देहं नाभितः प्रसृताः सिराः~॥ }\pend
\pstart\versepar \skt{प्रतानाः पद्मिनी कन्दा }\edtext{\skt{द्विसादीनां}}{\lemma{\skt{द्विसादीनां}}\Afootnote{H; \skt{द्बिसादीनां}\ A}}\skt{ यथाजलं~॥}\pend

\pstart \ledsidenote{\passref{3.7.trailer}}\edtext{\skt{इति}}{\lemma{\skt{इति}}\Afootnote{H; \skt{इति सुश्रु\appbrk{}त\appbrk{}सं\appbrk{}हि\appbrk{}तायां शारीरस्थाने सिरा\appbrk{}व\appbrk{}र्ण\appbrk{}वि\appbrk{}भक्ति}\ A}}\skt{ }\edtext{\skt{शारीरे}}{\lemma{\skt{शारीरे}}\Afootnote{H; \skt{शारीरं}\ A}}\skt{ }\edtext{\skt{सप्तमो}}{\lemma{\skt{सप्तमो}}\Afootnote{H; \skt{नाम सप्तमो}\ A}}\skt{ }\edtext{\skt{ध्यायः}}{\lemma{\skt{ध्यायः}}\Afootnote{H; \skt{ऽध्यायः}\ A}}\skt{~॥}\pend

\endnumbering
\beginnumbering
\eledsection*{Adhyāya 8}
\pstart \ledsidenote{\passref{3.8.1}}\skt{अथातो }\edtext{\skt{व्यधविधिं}}{\lemma{\skt{व्यधविधिं}}\Afootnote{H; \skt{सिराव्यधविधिं}\ A}}\skt{ शारीरं व्याख्यास्यामः~॥}\pend

\pstart \ledsidenote{\passref{3.8.3}}\edtext{\skt{तत्र}}{\lemma{\skt{तत्र}}\Afootnote{H; \emph{om.}\ A}}\skt{ }\edtext{\skt{बाल\-स्थ\-वि\-र\-भी\-रु\-रू\-क्ष\-क्ष\-त\-क्षी\-ण\-म\-द्या\-ध्व\-स्त्री\-क\-र्षि\-ता\-ना\-सि\-त\-वा\-न्त\-वि\-र\-क्त\-द\-न्त\-ब\-स्ति\-क्ली\-ब\-ग\-र्ब्भि\-णी\-का\-स\-श्वा\-स\-शो\-ष\-प्र\-वृ\-द्ध\-ज्व\-रा\-ती\-सा\-रा\-क्षे\-प\-क\-प\-क्षा\-घा\-तो\-प\-वा\-स\-पी\-डि\-तानां}}{\lemma{\skt{बाल\appbrk{}स्थ\appbrk{}वि\appbrk{}र\appbrk{}भी\appbrk{}रु\appbrk{}रू\appbrk{}क्ष\appbrk{}क्ष\appbrk{}त\appbrk{}क्षी\appbrk{}ण\appbrk{}म\appbrk{}द्या\appbrk{}ध्व\appbrk{}स्त्री\appbrk{}क\appbrk{}र्षि\appbrk{}ता\appbrk{}ना\appbrk{}सि\appbrk{}त\appbrk{}वा\appbrk{}न्त\appbrk{}वि\appbrk{}र\appbrk{}क्त\appbrk{}द\appbrk{}न्त\appbrk{}ब\appbrk{}स्ति\appbrk{}क्ली\appbrk{}ब\appbrk{}ग\appbrk{}र्ब्भि\appbrk{}णी\appbrk{}का\appbrk{}स\appbrk{}श्वा\appbrk{}स\appbrk{}शो\appbrk{}ष\appbrk{}प्र\appbrk{}वृ\appbrk{}द्ध\appbrk{}ज्व\appbrk{}रा\appbrk{}ती\appbrk{}सा\appbrk{}रा\appbrk{}क्षे\appbrk{}प\appbrk{}क\appbrk{}प\appbrk{}क्षा\appbrk{}घा\appbrk{}तो\appbrk{}प\appbrk{}वा\appbrk{}स\appbrk{}पी\appbrk{}डि\appbrk{}तानां}}\Afootnote{H; \skt{बाल\appbrk{}स्थ\appbrk{}वि\appbrk{}र\appbrk{}रू\appbrk{}क्ष\appbrk{}क्ष\appbrk{}त\appbrk{}क्षी\appbrk{}ण\appbrk{}भी\appbrk{}रु\appbrk{}प\appbrk{}रि\appbrk{}श्रा\appbrk{}न्त\appbrk{}म\appbrk{}द्य\appbrk{}पा\appbrk{}ध्व\appbrk{}स्त्री\appbrk{}क\appbrk{}र्शि\appbrk{}त\appbrk{}व\appbrk{}मि\appbrk{}त\appbrk{}वि\appbrk{}रि\appbrk{}क्ता\appbrk{}स्था\appbrk{}पि\appbrk{}ता\appbrk{}नु\appbrk{}वा\appbrk{}सि\appbrk{}त\appbrk{}जा\appbrk{}ग\appbrk{}रि\appbrk{}त\appbrk{}क्ली\appbrk{}ब\appbrk{}कृ\appbrk{}श\appbrk{}ग\appbrk{}र्भि\appbrk{}णीनां कास\appbrk{}श्वा\appbrk{}स\appbrk{}शो\appbrk{}ष\appbrk{}प्र\appbrk{}वृ\appbrk{}द्ध\appbrk{}ज्व\appbrk{}रा\appbrk{}क्षे\appbrk{}प\appbrk{}क\appbrk{}प\appbrk{}क्षा\appbrk{}घा\appbrk{}तो\appbrk{}प\appbrk{}वा\appbrk{}स\appbrk{}पि\appbrk{}पा\appbrk{}सा\appbrk{}मू\appbrk{}र्च्छा\appbrk{}प्र\appbrk{}पी\appbrk{}डि\appbrk{}तानां च}\ A}}\skt{ }\edtext{\skt{सिरानव्यध्याः~॥ किं}}{\lemma{\skt{सिरानव्यध्याः~॥ किं}}\Afootnote{H; \skt{सिरां न विध्येत् , याश्चाव्यध्}\ A}}\skt{ }\edtext{\skt{कारणमेषां तु}}{\lemma{\skt{कारणमेषां तु}}\Afootnote{H; \emph{om.}\ A}}\skt{ }\edtext{\skt{खलु}}{\lemma{\skt{खलु}}\Afootnote{H; \emph{om.}\ A}}\skt{ }\edtext{\skt{वेधा\-दि\-न्द्रि\-य\-सं\-मोहः}}{\lemma{\skt{वेधा\appbrk{}दि\appbrk{}न्द्रि\appbrk{}य\appbrk{}सं\appbrk{}मोहः}}\Afootnote{H; \emph{om.}\ A}}\skt{ }\edtext{\skt{शोषो वा}}{\lemma{\skt{शोषो वा}}\Afootnote{H; \emph{om.}\ A}}\skt{ }\edtext{\skt{भवति}}{\lemma{\skt{भवति}}\Afootnote{H; \emph{om.}\ A}}\skt{~। }\edtext{\skt{रक्तपित्तिनो}}{\lemma{\skt{रक्तपित्तिनो}}\Afootnote{H; \emph{om.}\ A}}\skt{ }\edtext{\edtext{\skt{रक्ता\-त्य\-या\-द्वा\-त\-कोपो}}{\lemma{\skt{रक्तात्ययाद्}}\Afootnote{H; \emph{om.}\ A}}}{\lemma{\skt{वातकोपो}}\Afootnote{H; \emph{om.}\ A}}\skt{ }\edtext{\skt{मृत्युर्व्वा}}{\lemma{\skt{मृत्युर्व्वा}}\Afootnote{H; \emph{om.}\ A}}\skt{~। }\edtext{\skt{गर्ब्भिण्यस्तु}}{\lemma{\skt{गर्ब्भिण्यस्तु}}\Afootnote{H; \emph{om.}\ A}}\skt{ }\edtext{\skt{गर्ब्भ\-शो\-ष\-प\-त\-न\-ही\-ना\-ङ्गा\-ना\-यु\-ष्कता वा}}{\lemma{\skt{गर्ब्भ\appbrk{}शो\appbrk{}ष\appbrk{}प\appbrk{}त\appbrk{}न\appbrk{}ही\appbrk{}ना\appbrk{}ङ्गा\appbrk{}ना\appbrk{}यु\appbrk{}ष्कता वा}}\Afootnote{H; \skt{याः}\ A}}\skt{ }\edtext{\skt{गर्ब्भस्य}}{\lemma{\skt{गर्ब्भस्य}}\Afootnote{\skt{व्यध्याश्चादृष्टाः, दृष्टाश्चायन्त्रिताः, यन्त्रि\appbrk{}ता\appbrk{}श्चा\appbrk{}नु\appbrk{}त्थिता इति}\ A; \skt{गर्ब्भस्य~॥ भवति चात्र}\ H}}\skt{~॥}\pend

\pstart\versepar \ledsidenote{\passref{3.8.3a}}\edtext{\skt{भवति}}{\lemma{\skt{भवति}}\Afootnote{\emph{om.}\ H}}\skt{ }\edtext{\skt{चात्र}}{\lemma{\skt{चात्र}}\Afootnote{\emph{om.}\ H}}\skt{~॥ }\pend
\pstart\versepar \skt{प्रति\-षि\-द्धा\-ना\-व\-शेकां तु शेषाणामपि देहिनां~। }\pend
\pstart\versepar \skt{}\edtext{\edtext{\edtext{\edtext{\skt{व्याधे\-र्व्वि\-वृ\-द्धि\-र्मृ\-त्युर्व्वा}}{\lemma{\skt{व्याधेर्}}\Afootnote{\skt{व्याधे}\ H}}}{\lemma{\skt{व्विवृद्धिर्}}\Afootnote{\skt{र्व्विवृद्धिर्}\ H}}}{\lemma{\skt{मृत्युर्}}\Afootnote{\skt{मृत्यु}\ H}}}{\lemma{\skt{व्वा}}\Afootnote{\skt{र्व्वा}\ H}}\skt{ सिरावेधेन सम्भवेत्~॥}\pend

\pstart \ledsidenote{\passref{3.8.3b}}\skt{याश्चाव्यध्या दृष्टा\-दृ\-ष्टा\-श्चा\-य\-न्त्रि\-ता\-य\-न्त्रिता नोत्थिता इति~॥}\pend

\pstart \ledsidenote{\passref{3.8.4}}\edtext{\skt{शोणि\-त\-मो\-क्ष\-सा\-ध्याश्च}}{\lemma{\skt{शोणि\appbrk{}त\appbrk{}मो\appbrk{}क्ष\appbrk{}सा\appbrk{}ध्याश्च}}\Afootnote{H; \skt{शोणि\appbrk{}ता\appbrk{}व\appbrk{}से\appbrk{}क\appbrk{}सा\appbrk{}ध्याश्च ये}\ A}}\skt{ विकाराः प्रागभिहिताः~। }\edtext{\skt{तेष्वपक्षेषु}}{\lemma{\skt{अपक्षेषु}}\Afootnote{H; \skt{चापक्वेष्व्}\ A}}\skt{ }\edtext{\skt{चान्येषु}}{\lemma{\skt{चान्येषु}}\Afootnote{H; \skt{अन्येषु}\ A}}\skt{ }\edtext{\skt{चानुभुक्तेषु}}{\lemma{\skt{चानुभुक्तेषु}}\Afootnote{H; \skt{चानुक्तेषु}\ A}}\skt{ }\edtext{\skt{यथा\-भ्या\-स\-म\-न\-त्य\-यञ्च}}{\lemma{\skt{अनत्ययञ्च}}\Afootnote{H; \skt{यथान्यायं च}\ A}}\skt{ सिरां }\edtext{\skt{व्यधयेत्}}{\lemma{\skt{व्यधयेत्}}\Afootnote{H; \skt{विध्येत्}\ A}}\skt{~॥}\pend

\pstart \ledsidenote{\passref{3.8.5}}\skt{प्रति\-षि\-द्धा\-ना\-मपि च विषोपसर्ग्गे }\edtext{\skt{सिरा\-व्य\-ध\-न\-म\-प्र\-ति\-षिद्धं}}{\lemma{\skt{सिराव्यधनम्}}\Afootnote{H; \skt{आत्ययिके च सिराव्यधनम्}\ A}}\pend

\pstart \ledsidenote{\passref{3.8.6}}\skt{तत्र स्निग्ध\-स्वि\-न्न\-मा\-तुरं यथा दोषप्रत्यनीकं द्रवप्रायमन्नं भुक्तवन्तं }\edtext{\skt{यवागूम्पीतम्वा}}{\lemma{\skt{पीतम्वा}}\Afootnote{H; \skt{पीतवन्तं वा}\ A}}\skt{ यथा\-का\-ल\-मु\-प\-स्था\-प्या\-सीनं स्थितम्वा प्राणा\-न\-बा\-ध\-मानः~। }\edtext{\skt{प्लोत\-प\-ट्ट\-श्च\-र्म्मा\-न्त\-र्व्व\-ल्क\-ल\-ल\-ता\-ना\-म\-न्य\-त\-मेन}}{\lemma{\skt{प्लोतपट्टश्}}\Afootnote{H; \skt{वस्त्रपट्ट}\ A}}\skt{ }\edtext{\skt{यन्त्रेण}}{\lemma{\skt{यन्त्रेण}}\Afootnote{H; \emph{om.}\ A}}\skt{ यन्त्रयित्वा नाति\-गा\-ढ\-न्ना\-ति\-शि\-थिलं शरी\-र\-प्र\-दे\-श\-मा\-साद्य }\edtext{\skt{विधि\-व\-च्छ\-स्त्र\-मा\-दाय}}{\lemma{\skt{विधि\appbrk{}व\appbrk{}च्छ\appbrk{}स्त्र\appbrk{}मा\appbrk{}दाय}}\Afootnote{H; \skt{प्राप्तं शस्त्रमादा}\ A}}\skt{ }\edtext{\skt{यथोक्तां}}{\lemma{\skt{यथोक्तां}}\Afootnote{H; \skt{य}\ A}}\skt{ सिरां }\edtext{\skt{व्यधयेत्}}{\lemma{\skt{व्यधयेत्}}\Afootnote{H; \skt{विध्येत्}\ A}}\skt{~।}\pend

\pstart \ledsidenote{\passref{3.8.7}}\edtext{\skt{तत्र}}{\lemma{\skt{तत्र}}\Afootnote{H; \emph{om.}\ A}}\skt{ }\edtext{\skt{व्यध्यसिरां}}{\lemma{\skt{व्यध्यसिरां}}\Afootnote{H; \emph{om.}\ A}}\skt{~। नैवातिशीते नात्युष्णे न }\edtext{\skt{प्रवाते न च}}{\lemma{\skt{प्रवाते न च}}\Afootnote{H; \skt{प्रवाते न चा}\ A}}\skt{ }\edtext{\skt{दुर्द्दिने}}{\lemma{\skt{दुर्द्दिने}}\Afootnote{H; \skt{भ्रिते}\ A}}\skt{~। सिराणां व्यधनं कार्यमरोगे वा कदाचन~॥}\pend

\pstart \ledsidenote{\passref{3.8.8}}\skt{तत्र }\edtext{\skt{व्यध्यसिरां}}{\lemma{\skt{व्यध्यसिरां}}\Afootnote{H; \skt{व्यध्यसिरं}\ A}}\skt{ पुरुषं }\edtext{\edtext{\skt{प्रत्या\-दि\-त्य\-मु\-ख\-म\-र\-न्नि\-मा\-त्रो\-च्छ्रिते}}{\lemma{\skt{प्रत्यादित्य्}}\Afootnote{H; \skt{प्रत्यादित्}\ A}}}{\lemma{\skt{अमु\appbrk{}ख\appbrk{}म\appbrk{}र\appbrk{}न्नि\appbrk{}मा\appbrk{}त्रो\appbrk{}च्छ्रिते}}\Afootnote{H; \skt{यमु\appbrk{}ख\appbrk{}म\appbrk{}र\appbrk{}त्नि\appbrk{}मा\appbrk{}त्रो\appbrk{}च्छ्रिते}\ A}}\skt{ }\edtext{\skt{निवेश्यासने}}{\lemma{\skt{निवेश्यासने}}\Afootnote{H; \skt{उपवेश्यासने}\ A}}\skt{ सक्थ्नो\-रा\-कु\-ञ्चि\-त\-यो\-र्न्नि\-वेश्य }\edtext{\skt{कुर्प्प\-र\-स\-न्धि\-द्वयं}}{\lemma{\skt{कुर्प्प\appbrk{}र\appbrk{}स\appbrk{}न्धि\appbrk{}द्वयं}}\Afootnote{H; \skt{कूर्परे सन्धिद्वय}\ A}}\skt{ }\edtext{\edtext{\skt{हस्ता\-व\-न्त\-गू\-ढा\-ङ्गु\-ष्ठ\-कृ\-त\-मु\-ष्टी\-म\-न्ययोः}}{\lemma{\skt{हस्ता\appbrk{}व\appbrk{}न्त\appbrk{}गू\appbrk{}ढा\appbrk{}ङ्गु\appbrk{}ष्ठ\appbrk{}कृ\appbrk{}त\appbrk{}मु\appbrk{}ष्टीम्}}\Afootnote{H; \skt{स्योपरि हस्ता\appbrk{}व\appbrk{}न्त\appbrk{}र्गू\appbrk{}ढा\appbrk{}ङ्गु\appbrk{}ष्ठ\appbrk{}कृ\appbrk{}त\appbrk{}मुष्टी}\ A}}}{\lemma{\skt{अन्ययोः}}\Afootnote{H; \skt{मन्ययोः}\ A}}\skt{ स्थापयित्वा }\edtext{\skt{यन्त्रशाटकेन}}{\lemma{\skt{यन्त्रशाटकेन}}\Afootnote{H; \skt{यन्त्रणशाटकं ग्रीवा\appbrk{}मु\appbrk{}ष्ट्यो\appbrk{}रु\appbrk{}परि}\ A}}\skt{ }\edtext{\skt{परिक्षिप्य}}{\lemma{\skt{परिक्षिप्य}}\Afootnote{H; \skt{परिक्षिप्}\ A}}\skt{ }\edtext{\edtext{\skt{ग्रीवायामन्येन}}{\lemma{\skt{ग्रीवायाम्}}\Afootnote{H; \skt{या}\ A}}}{\lemma{\skt{अन्येन}}\Afootnote{H; \skt{न्येन}\ A}}\skt{ }\edtext{\skt{पश्चात्स्थितेन}}{\lemma{\skt{पश्चात्}}\Afootnote{H; \skt{पुरुषेण पश्चात्}\ A}}\skt{ }\edtext{\skt{पुरुषेण}}{\lemma{\skt{पुरुषेण}}\Afootnote{H; \emph{om.}\ A}}\skt{ }\edtext{\skt{वामहस्तेन}}{\lemma{\skt{वामहस्तेन}}\Afootnote{H; \skt{वाम\appbrk{}ह\appbrk{}स्ते\appbrk{}नो\appbrk{}त्ता\appbrk{}नेन}\ A}}\skt{ शाटकान्तद्वयं ग्राहयित्वा }\edtext{\edtext{\skt{सिरो\-त्था\-प\-ना\-र्थ\-म\-सृ\-क्स्रा\-व\-णा\-र्थञ्च}}{\lemma{\skt{सिरो\appbrk{}त्था\appbrk{}प\appbrk{}ना\appbrk{}र्थम्}}\Afootnote{H; \skt{ततो ब्रूयात् --- दक्षिणहस्तेन सिरोत्थापनार्थं}\ A}}}{\lemma{\skt{असृ\appbrk{}क्स्रा\appbrk{}व\appbrk{}णा\appbrk{}र्थञ्च}}\Afootnote{H; \skt{नात्या\appbrk{}य\appbrk{}त\appbrk{}शि\appbrk{}थिलं यन्त्र\appbrk{}मा\appbrk{}वे\appbrk{}ष्ट\appbrk{}येति, असृक्स्रावणार्थं च}\ A}}\skt{ }\edtext{\skt{यन्त्रम्प्रमाणेन}}{\lemma{\skt{प्रमाणेन}}\Afootnote{H; \skt{पृष्ठमध्ये}\ A}}\skt{ }\edtext{\skt{पीडयेत्}}{\lemma{\skt{पीडयेत्}}\Afootnote{H; \skt{पीडयेति}\ A}}\skt{~। }\edtext{\skt{कर्म्मपुरुषश्च}}{\lemma{\skt{कर्म्मपुरुषश्च}}\Afootnote{H; \skt{कर्मपुरुषं च}\ A}}\skt{ }\edtext{\skt{मुखं}}{\lemma{\skt{मुखं}}\Afootnote{H; \emph{om.}\ A}}\skt{ }\edtext{\skt{वायुना}}{\lemma{\skt{वायुना}}\Afootnote{H; \skt{वायुपूर्णमुखं स्था}\ A}}\skt{ }\edtext{\skt{पूरयेत्~। }\edtext{\skt{उत्त\-मा\-ङ्ग\-ग\-ता\-ना\-मेष}}{\lemma{\skt{एष}}\Afootnote{H; \skt{अन्तर्मुखव}\ A}}}{\lemma{\skt{पूरयेत्~। उत्त\appbrk{}मा\appbrk{}ङ्ग\appbrk{}ग\appbrk{}ता\appbrk{}नाम्}}\Afootnote{H; \skt{पयेत् ; एष उत्त\appbrk{}मा\appbrk{}ङ्ग\appbrk{}ग\appbrk{}ता\appbrk{}नाम्}\ A}}\skt{ }\edtext{\skt{सिरो\-त्था\-प\-म\-य\-न्त्र\-विधिः~। }\edtext{\edtext{\edtext{\skt{स्थित\-स\-क्थि\-मी\-षे\-त्कु\-ञ्चि\-त\-पा\-द\-मुच्चैः}}{\lemma{\skt{ईषेत्}}\Afootnote{H; \emph{om.}\ A}}}{\lemma{\skt{कुञ्चितपादम्}}\Afootnote{H; \skt{य पादं}\ A}}}{\lemma{\skt{उच्चैः}}\Afootnote{H; \skt{समे स्थाने सुस्थितं}\ A}}}{\lemma{\skt{सिरो\appbrk{}त्था\appbrk{}प\appbrk{}म\appbrk{}य\appbrk{}न्त्र\appbrk{}विधिः~। स्थितसक्थिम्}}\Afootnote{H; \skt{र्जानां सिराणां व्यधने यन्त्रणविधिः~। पाद\appbrk{}व्य\appbrk{}ध्य\appbrk{}सि\appbrk{}रस्}\ A}}\skt{ }\edtext{\skt{स्थापयित्वा}}{\lemma{\skt{स्थापयित्वा}}\Afootnote{H; \skt{स्थापयित्वान्यं पाद\appbrk{}मी\appbrk{}ष\appbrk{}त्स\appbrk{}ङ्कु\appbrk{}चि\appbrk{}त\appbrk{}मुच्चैः कृत्वा}\ A}}\skt{ }\edtext{\skt{व्यध\-प्र\-दे\-श\-स्यो\-परि}}{\lemma{\skt{व्यध\appbrk{}प्र\appbrk{}दे\appbrk{}श\appbrk{}स्यो\appbrk{}परि}}\Afootnote{H; \skt{व्यध्यसिरं पादं जानुसन्धेरधः शाटकेनावेष्ट्य हस्ताभ्यां प्रपीड्य गुल्फं व्यध्य\appbrk{}प्र\appbrk{}दे\appbrk{}श\appbrk{}स्यो\appbrk{}परि}\ A}}\skt{ }\edtext{\edtext{\edtext{\skt{सम्यग्यन्त्रम्बध्वा}}{\lemma{\skt{सम्यग्}}\Afootnote{H; \skt{च}\ A}}}{\lemma{\skt{यन्त्रम्}}\Afootnote{H; \skt{तुरङ्}\ A}}}{\lemma{\skt{बध्वा}}\Afootnote{H; \skt{गुले प्लोता\appbrk{}दी\appbrk{}ना\appbrk{}म\appbrk{}न्य\appbrk{}त\appbrk{}मेन बद्ध्वा वा}\ A}}\skt{ }\edtext{\skt{पाद\-सि\-रा\-म्वि\-ध्येत्~। उपविष्टस्य}}{\lemma{\skt{विध्येत्~। उपविष्टस्य}}\Afootnote{H; \skt{विध्येत्~। अथोपरिष्ठा}\ A}}\skt{ }\edtext{\skt{हस्तं}}{\lemma{\skt{हस्तं}}\Afootnote{H; \skt{द्धस्तौ}\ A}}\skt{ }\edtext{\skt{गूढां\-ष्ठ\-गु\-ष्ठ\-नि\-ब\-द्ध\-मुष्टिं}}{\lemma{\skt{गूढां\appbrk{}ष्ठ\appbrk{}गु\appbrk{}ष्ठ\appbrk{}नि\appbrk{}ब\appbrk{}द्ध\appbrk{}मुष्टिं}}\Afootnote{H; \skt{गूढा\appbrk{}ङ्गु\appbrk{}ष्ठ\appbrk{}कृ\appbrk{}त\appbrk{}मुष्टी}\ A}}\skt{ सम्यगासने }\edtext{\skt{स्थापयित्वा}}{\lemma{\skt{स्थापयित्वा}}\Afootnote{H; \skt{स्थापयित्वा सुखोपविष्टस्य पूर्ववद्यन्त्रं बद्ध्वा}\ A}}\skt{ हस्तसिरां विध्येत्~। }\edtext{\skt{गृध्र\-सी\-वि\-श्व\-च्यो\-रा\-कु\-ञ्चि\-त\-जा\-नु\-कु\-र्प्प\-रस्य}}{\lemma{\skt{गृध्र\appbrk{}सी\appbrk{}वि\appbrk{}श्व\appbrk{}च्यो\appbrk{}रा\appbrk{}कु\appbrk{}ञ्चि\appbrk{}त\appbrk{}जा\appbrk{}नु\appbrk{}कु\appbrk{}र्प्प\appbrk{}रस्य}}\Afootnote{H; \skt{गृध्रसीविश्वाच्योः सङ्कु\appbrk{}चि\appbrk{}त\appbrk{}जा\appbrk{}नु\appbrk{}कू\appbrk{}र्प\appbrk{}रस्य}\ A}}\skt{ }\edtext{\skt{श्रोणी\-पृ\-ष्ठ\-स्क\-न्धे\-षू\-न्न\-मि\-त\-पृ\-ष्ठस्य}}{\lemma{\skt{श्रोणी\appbrk{}पृ\appbrk{}ष्ठ\appbrk{}स्क\appbrk{}न्धे\appbrk{}षू\appbrk{}न्न\appbrk{}मि\appbrk{}त\appbrk{}पृ\appbrk{}ष्ठस्य}}\Afootnote{H; \skt{श्रोणी\appbrk{}पृ\appbrk{}ष्ठ\appbrk{}स्क\appbrk{}न्धे\appbrk{}षू\appbrk{}न्ना\appbrk{}मि\appbrk{}त\appbrk{}पृ\appbrk{}ष्ठ\appbrk{}स्या\appbrk{}वा\appbrk{}क्शि\appbrk{}र\appbrk{}स्क\appbrk{}स्यो\appbrk{}प\appbrk{}वि\appbrk{}ष्टस्य~। विस्फू\appbrk{}र्जि\appbrk{}त\appbrk{}पृ\appbrk{}ष्ठस्य विध्ये}\ A}}\skt{ }\edtext{\skt{उदरोरसौ}}{\lemma{\skt{उदरोरसौ}}\Afootnote{H; \skt{त्~। उदरोरसोः}\ A}}\skt{ }\edtext{\skt{प्रसारितो}}{\lemma{\skt{प्रसारितो}}\Afootnote{H; \skt{प्रसा\appbrk{}रि\appbrk{}तो\appbrk{}र\appbrk{}स्क\appbrk{}स्यो\appbrk{}न्ना\appbrk{}मित}\ A}}\skt{ }\edtext{\skt{रसः}}{\lemma{\skt{रसः}}\Afootnote{H; \skt{शिरस्कस्य}\ A}}\skt{ }\edtext{\skt{स्फुर्ज्जितदेहस्य}}{\lemma{\skt{स्फुर्ज्जितदेहस्य}}\Afootnote{H; \skt{विस्फू\appbrk{}र्जि\appbrk{}त\appbrk{}दे\appbrk{}हस्य}\ A}}\skt{ }\edtext{\edtext{\skt{बाहू\-भ्या\-म\-व\-ल\-म्ब\-मा\-नस्य}}{\lemma{\skt{बाहूभ्याम्}}\Afootnote{H; \skt{बाहुभ्याम्}\ A}}}{\lemma{\skt{अव\appbrk{}ल\appbrk{}म्ब\appbrk{}मा\appbrk{}नस्य}}\Afootnote{H; \skt{अव\appbrk{}ल\appbrk{}म्ब\appbrk{}मा\appbrk{}न\appbrk{}दे\appbrk{}हस्य}\ A}}\skt{ पार्श्वयोः~। }\edtext{\skt{अव\-न\-मि\-त\-मे\-ढ्रस्य}}{\lemma{\skt{अव\appbrk{}न\appbrk{}मि\appbrk{}त\appbrk{}मे\appbrk{}ढ्रस्य}}\Afootnote{H; \skt{अना\appbrk{}मि\appbrk{}त\appbrk{}मे\appbrk{}ढ्रस्य}\ A}}\skt{ }\edtext{\skt{मेढ्रं}}{\lemma{\skt{मेढ्रं}}\Afootnote{H; \skt{मेढ्रे}\ A}}\skt{~। }\edtext{\skt{उन्ना\-मि\-त\-वि\-द\-ष्ट\-जि\-ह्वा\-स्या\-धो\-जि\-ह्वा\-याम्}}{\lemma{\skt{उन्ना\appbrk{}मि\appbrk{}त\appbrk{}वि\appbrk{}द\appbrk{}ष्ट\appbrk{}जि\appbrk{}ह्वा\appbrk{}स्या\appbrk{}धो\appbrk{}जि\appbrk{}ह्वा\appbrk{}याम्}}\Afootnote{H; \skt{उन्न\appbrk{}मि\appbrk{}त\appbrk{}वि\appbrk{}द\appbrk{}ष्ट\appbrk{}जि\appbrk{}ह्वा\appbrk{}ग्र\appbrk{}स्या\appbrk{}धो\appbrk{}जि\appbrk{}ह्वा\appbrk{}याम्}\ A}}\skt{~। }\edtext{\skt{अति\-व्या\-ता\-ना\-न\-नस्य}}{\lemma{\skt{अति\appbrk{}व्या\appbrk{}ता\appbrk{}ना\appbrk{}न\appbrk{}नस्य}}\Afootnote{H; \skt{अति\appbrk{}व्या\appbrk{}त्ता\appbrk{}न\appbrk{}नस्य}\ A}}\skt{ तालुनि }\edtext{\skt{दन्तमूले च}}{\lemma{\skt{दन्तमूले च}}\Afootnote{H; \skt{दन्तमूलेषु च}\ A}}\skt{~। एवं यन्त्रो\-पा\-या\-न\-न्यांश्च सिरो\-त्था\-प\-न\-हे\-तू\-न्बु\-द्ध्या\-वेक्ष्य शरीरवशेन व्याधिवशेन च }\edtext{\skt{विदध्यादिति}}{\lemma{\skt{विदध्यादिति}}\Afootnote{H; \skt{विदध्यात्}\ A}}\skt{~।}\pend

\pstart \ledsidenote{\passref{3.8.9}}\edtext{\skt{मांस\-ले\-ष्व\-व\-का\-सेषु}}{\lemma{\skt{अवकासेषु}}\Afootnote{H; \skt{अवकाशेषु}\ A}}\skt{ }\edtext{\skt{व्रीहिमुखेन}}{\lemma{\skt{व्रीहिमुखेन}}\Afootnote{H; \emph{om.}\ A}}\skt{ }\edtext{\skt{यवमात्रं}}{\lemma{\skt{यवमात्रं}}\Afootnote{H; \skt{यवमात्रं शस्त्रं}\ A}}\skt{ }\edtext{\edtext{\skt{निद\-ध्या\-द्व्री\-हि\-मा\-त्रम्वा}}{\lemma{\skt{निदध्याद्}}\Afootnote{H; \skt{निदध्या}\ A}}}{\lemma{\skt{व्रीहिमात्रम्वा}}\Afootnote{H; \emph{om.}\ A}}\skt{~। }\edtext{\skt{अतो}}{\lemma{\skt{अतो}}\Afootnote{H; \skt{त्, अतो}\ A}}\skt{ }\edtext{\edtext{\skt{न्यथा\-र्द्ध\-य\-व\-मा\-त्र\-म\-स्थ्ना\-मु\-परि}}{\lemma{\skt{न्यथा\appbrk{}र्द्ध\appbrk{}य\appbrk{}व\appbrk{}मा\appbrk{}त्रम्}}\Afootnote{H; \skt{ऽन्य\appbrk{}था\appbrk{}र्ध\appbrk{}य\appbrk{}व\appbrk{}मात्रं व्रीहिमात्रं वा}\ A}}}{\lemma{\skt{अस्थ्नाम्}}\Afootnote{H; \skt{व्रीहिमुखेन, अस्थ्नाम्}\ A}}\skt{ }\edtext{\skt{कुठारिकया}}{\lemma{\skt{कुठारिकया}}\Afootnote{H; \skt{कुठारिकया विध्ये\appbrk{}द\appbrk{}र्ध\appbrk{}य\appbrk{}व\appbrk{}मा\appbrk{}त्रम्}\ A}}\skt{~॥}\pend

\pstart \ledsidenote{\passref{3.8.10}}\edtext{\skt{भवंति}}{\lemma{\skt{भवंति}}\Afootnote{H; \emph{om.}\ A}}\skt{ }\edtext{\skt{चात्र}}{\lemma{\skt{चात्र}}\Afootnote{H; \emph{om.}\ A}}\skt{~॥ व्यभ्रे वर्षासु }\edtext{\skt{विध्येत}}{\lemma{\skt{विध्येत}}\Afootnote{H; \skt{विध्येत्तु}\ A}}\skt{ }\edtext{\skt{ग्रीष्मकालेषु}}{\lemma{\skt{ग्रीष्मकालेषु}}\Afootnote{H; \skt{ग्रीष्मकाले तु}\ A}}\skt{ शीतले~। हेमन्तकाले मध्याह्ने शस्त्रकालास्त्रयः स्मृताः~॥}\pend

\pstart \ledsidenote{\passref{3.8.11}}\edtext{\skt{द्विप्र\-का\-र\-प्र\-म\-थि\-त\-धा\-रया}}{\lemma{\skt{द्विप्र\appbrk{}का\appbrk{}र\appbrk{}प्र\appbrk{}म\appbrk{}थि\appbrk{}त\appbrk{}धा\appbrk{}रया}}\Afootnote{H; \skt{सम्य\appbrk{}क्श\appbrk{}स्त्र\appbrk{}नि\appbrk{}पा\appbrk{}तेन धारया}\ A}}\skt{ यास्रवेदसृक्~। मुहूर्त्तं }\edtext{\skt{रुद्ध्वा}}{\lemma{\skt{रुद्ध्वा}}\Afootnote{H; \skt{रुद्धा}\ A}}\skt{ तिष्ठेच्च सुविद्धान्तां विनिर्द्दिशेत्~॥}\pend

\pstart \ledsidenote{\passref{3.8.12}}\skt{यथा कुसुम्भपुष्पेभ्यः पूर्व्वं स्रवति पीतिका~। तथा सिरासु विद्धासु दुष्टमग्रे प्रवर्त्तते~॥}\pend

\pstart \ledsidenote{\passref{3.8.14}}\skt{क्षीणस्य बहुदोषस्य }\edtext{\skt{मूर्च्छ\-यो\-प\-द्रु\-तस्य वा}}{\lemma{\skt{मूर्च्छ\appbrk{}यो\appbrk{}प\appbrk{}द्रु\appbrk{}तस्य वा}}\Afootnote{H; \skt{मूर्च्छ\appbrk{}या\appbrk{}भि\appbrk{}ह\appbrk{}तस्य च}\ A}}\skt{~। भूयोऽपराह्णे }\edtext{\edtext{\edtext{\skt{विस्रा\-व्या\-त्तू\-त्त\-रेऽ\-हनि}}{\lemma{\skt{विस्राव्यात्}}\Afootnote{H; \skt{विस्राव्या सापरेद्युस्त्}\ A}}}{\lemma{\skt{तूत्तरे}}\Afootnote{H; \skt{र्यहे}\ A}}}{\lemma{\skt{ऽहनि}}\Afootnote{H; \skt{ऽपि}\ A}}\skt{ }\edtext{\skt{वायुनः}}{\lemma{\skt{वायुनः}}\Afootnote{H; \skt{वा}\ A}}\skt{~॥}\pend

\pstart \ledsidenote{\passref{3.8.15}}\edtext{\edtext{\skt{सशे\-ष\-दो\-ष\-रु\-धि\-र\-मपि}}{\lemma{\skt{सशे\appbrk{}ष\appbrk{}दो\appbrk{}ष\appbrk{}रु\appbrk{}धि\appbrk{}रम्}}\Afootnote{H; \skt{रक्तं सशेषदोषं तु}\ A}}}{\lemma{\skt{अपि}}\Afootnote{H; \emph{om.}\ A}}\skt{ }\edtext{\skt{कुर्याद्विचक्षणः}}{\lemma{\skt{कुर्याद्}}\Afootnote{H; \skt{कुर्यादपि}\ A}}\skt{~। न }\edtext{\skt{चातिनिसृतं}}{\lemma{\skt{चातिनिसृतं}}\Afootnote{H; \skt{चातिनिःस्रुतं}\ A}}\skt{ कुर्याच्छेषं संश\-म\-नै\-र्ज्ज\-येत्~॥}\pend

\pstart \ledsidenote{\passref{3.8.17}}\skt{तत्र }\edtext{\skt{पाद\-दा\-ह\-पा\-द\-ह\-र्ष\-चि\-प्य\-वा\-त\-शो\-णि\-त\-पा\-द\-क\-ण्ट\-क\-वि\-पा\-दि\-का\-पा\-द\-दा\-री\-प्र\-भृ\-तिषु}}{\lemma{\skt{पाद\appbrk{}दा\appbrk{}ह\appbrk{}पा\appbrk{}द\appbrk{}ह\appbrk{}र्ष\appbrk{}चि\appbrk{}प्य\appbrk{}वा\appbrk{}त\appbrk{}शो\appbrk{}णि\appbrk{}त\appbrk{}पा\appbrk{}द\appbrk{}क\appbrk{}ण्ट\appbrk{}क\appbrk{}वि\appbrk{}पा\appbrk{}दि\appbrk{}का\appbrk{}पा\appbrk{}द\appbrk{}दा\appbrk{}री\appbrk{}प्र\appbrk{}भृ\appbrk{}तिषु}}\Afootnote{H; \skt{पाद\appbrk{}दा\appbrk{}ह\appbrk{}पा\appbrk{}द\appbrk{}ह\appbrk{}र्ष\appbrk{}चि\appbrk{}प्प\appbrk{}वि\appbrk{}स\appbrk{}र्प\appbrk{}वा\appbrk{}त\appbrk{}शो\appbrk{}णि\appbrk{}त\appbrk{}वा\appbrk{}त\appbrk{}क\appbrk{}ण्ट\appbrk{}क\appbrk{}वि\appbrk{}च\appbrk{}र्चि\appbrk{}का\appbrk{}पा\appbrk{}द\appbrk{}दा\appbrk{}री\appbrk{}प्र\appbrk{}भृति}\ A}}\skt{ }\edtext{\skt{रोगेषु}}{\lemma{\skt{रोगेषु}}\Afootnote{H; \skt{षु}\ A}}\skt{ }\edtext{\skt{क्षिप्र\-म\-र्म्मो\-प\-रि\-ष्टा\-द्द्व्य\-ङ्गुले}}{\lemma{\skt{क्षिप्र\appbrk{}म\appbrk{}र्म्मो\appbrk{}प\appbrk{}रि\appbrk{}ष्टाद्}}\Afootnote{H; \skt{क्षिप्रमर्मण उपरिष्ठाद्}\ A}}\skt{ व्रीहिमुखेन सिरां }\edtext{\skt{व्यधयेत्}}{\lemma{\skt{व्यधयेत्}}\Afootnote{H; \skt{विध्येत्}\ A}}\skt{~। }\edtext{\skt{श्लीपदेषु}}{\lemma{\skt{श्लीपदेषु}}\Afootnote{H; \skt{श्लीपदे तच्चिकित्सिते}\ A}}\skt{ }\edtext{\skt{यथोक्तं}}{\lemma{\skt{यथोक्तं}}\Afootnote{H; \skt{यथा वक्ष्यते}\ A}}\skt{ }\edtext{\skt{क्रोष्टु\-क\-शि\-र्ष\-वा\-त\-वे\-द\-नासु}}{\lemma{\skt{क्रोष्टु\appbrk{}क\appbrk{}शि\appbrk{}र्ष\appbrk{}वा\appbrk{}त\appbrk{}वे\appbrk{}द\appbrk{}नासु}}\Afootnote{H; \skt{क्रोष्टु\appbrk{}क\appbrk{}शि\appbrk{}रः\appbrk{}ख\appbrk{}ञ्ज\appbrk{}प\appbrk{}ङ्गु\appbrk{}ल\appbrk{}वा\appbrk{}त\appbrk{}वे\appbrk{}द\appbrk{}नासु}\ A}}\skt{ }\edtext{\skt{सक्थिसंस्रितां}}{\lemma{\skt{सक्थिसंस्रितां}}\Afootnote{H; \skt{जङ्घायां}\ A}}\skt{~। }\edtext{\edtext{\skt{गुल्फ\-स\-न्धे\-रु\-परि}}{\lemma{\skt{गुल्फसन्धेर्}}\Afootnote{H; \skt{गुल्फस्यो}\ A}}}{\lemma{\skt{उपरि}}\Afootnote{H; \skt{परि}\ A}}\skt{ चतुरङ्गुले }\edtext{\edtext{\skt{अप\-च्या\-मि\-न्द्र\-ब\-स्ते\-र\-ध\-स्ता\-द्द्व्य\-ङ्गुले}}{\lemma{\skt{अप\appbrk{}च्या\appbrk{}मि\appbrk{}न्द्र\appbrk{}ब\appbrk{}स्तेर्}}\Afootnote{H; \skt{अप\appbrk{}च्या\appbrk{}मि\appbrk{}न्द्र\appbrk{}बस्ते}\ A}}}{\lemma{\skt{अधस्ताद्}}\Afootnote{H; \skt{रधस्ताद्}\ A}}\skt{~। }\edtext{\edtext{\skt{जानु\-स\-न्धे\-रु\-परि}}{\lemma{\skt{जानुसन्धेर्}}\Afootnote{H; \skt{जानुसन्धे}\ A}}}{\lemma{\skt{उपरि}}\Afootnote{H; \skt{रुपर्}\ A}}\skt{ }\edtext{\skt{चतुरङ्गुले}}{\lemma{\skt{चतुरङ्गुले}}\Afootnote{H; \skt{यधो वा चतुरङ्गुले}\ A}}\skt{ }\edtext{\skt{गृध्र\-स्या\-मू\-रु\-सं\-धि\-स\-न्नि\-विष्टां}}{\lemma{\skt{गृध्र\appbrk{}स्या\appbrk{}मू\appbrk{}रु\appbrk{}सं\appbrk{}धि\appbrk{}स\appbrk{}न्नि\appbrk{}विष्टां}}\Afootnote{H; \skt{गृध्रस्यां, ऊरु\appbrk{}मू\appbrk{}ल\appbrk{}सं\appbrk{}श्रितां}\ A}}\skt{~। गलगण्डे }\edtext{\edtext{\skt{श्रोणि\-प्र\-ति\-स\-म\-न्ता\-त्प्र\-वा\-हि\-कां\-यां\-शू\-लिन्यां}}{\lemma{\skt{श्रोणि\appbrk{}प्र\appbrk{}ति\appbrk{}स\appbrk{}म\appbrk{}न्तात्}}\Afootnote{H; \emph{om.}\ A}}}{\lemma{\skt{प्रवा\appbrk{}हि\appbrk{}कां\appbrk{}यां\appbrk{}शू\appbrk{}लिन्यां}}\Afootnote{H; \emph{om.}\ A}}\skt{ }\edtext{\skt{परि\-व\-र्त्ति\-को\-प\-दं\-श\-शू\-क\-व्या\-पत्सु}}{\lemma{\skt{परि\appbrk{}व\appbrk{}र्त्ति\appbrk{}को\appbrk{}प\appbrk{}दं\appbrk{}श\appbrk{}शू\appbrk{}क\appbrk{}व्या\appbrk{}पत्सु}}\Afootnote{H; \emph{om.}\ A}}\skt{ }\edtext{\skt{मेढ्रमध्ये}}{\lemma{\skt{मेढ्रमध्ये}}\Afootnote{H; \emph{om.}\ A}}\skt{~। }\edtext{\skt{एते\-ने\-त\-र\-स\-क्थि\-बा\-हू\-व्या\-ख्यातौ}}{\lemma{\skt{एते\appbrk{}ने\appbrk{}त\appbrk{}र\appbrk{}स\appbrk{}क्थि\appbrk{}बा\appbrk{}हू\appbrk{}व्या\appbrk{}ख्यातौ}}\Afootnote{H; \skt{एते\appbrk{}ने\appbrk{}त\appbrk{}र\appbrk{}सक्थि बाहू च व्याख्यातौ}\ A}}\skt{~॥ विशेषतस्तु }\edtext{\skt{वाम\-बा\-हु\-कु\-र्प्प\-र\-स\-न्धा\-व\-भ्य\-न्त\-रतो}}{\lemma{\skt{वाम\appbrk{}बा\appbrk{}हु\appbrk{}कु\appbrk{}र्प्प\appbrk{}र\appbrk{}स\appbrk{}न्धा\appbrk{}व\appbrk{}भ्य\appbrk{}न्त\appbrk{}रतो}}\Afootnote{H; \skt{वामबाहौ कूर्प\appbrk{}र\appbrk{}स\appbrk{}न्धे\appbrk{}र\appbrk{}भ्य\appbrk{}न्त\appbrk{}रतो}\ A}}\skt{ बाहुमध्ये }\edtext{\skt{प्लीह्नि}}{\lemma{\skt{प्लीह्नि}}\Afootnote{H; \skt{प्लीह्नि कनि}\ A}}\skt{ }\edtext{\edtext{\skt{एतामेव}}{\lemma{\skt{एताम्}}\Afootnote{H; \skt{ष्ठिका\appbrk{}ना\appbrk{}मि\appbrk{}क\appbrk{}यो\appbrk{}र्मध्ये वा}\ A}}}{\lemma{\skt{एव}}\Afootnote{H; \skt{एवं}\ A}}\skt{ }\edtext{\skt{दक्षि\-ण\-बा\-हु\-य\-कृ\-द्दाल्ये}}{\lemma{\skt{दक्षि\appbrk{}ण\appbrk{}बा\appbrk{}हु\appbrk{}य\appbrk{}कृ\appbrk{}द्दाल्ये}}\Afootnote{H; \skt{दक्षिणबाहौ यकृद्दाल्ये}\ A}}\skt{~। }\edtext{\skt{एतामेव}}{\lemma{\skt{एताम्}}\Afootnote{H; \skt{(कफोदरे च,) एताम्}\ A}}\skt{ च कास\-श्वा\-स\-यो\-र\-प्या\-दि\-शन्ति~। }\edtext{\skt{गृध्रस्यामिव}}{\lemma{\skt{गृध्रस्यामिव}}\Afootnote{H; \skt{गृध्रस्यामिव}\ A}}\skt{ }\edtext{\edtext{\skt{विपं\-च्या\-म्वा\-म\-पा\-र्श्व\-क\-क्ष\-स्त\-ना\-न्त\-र\-यो\-र\-न्त\-र\-स्था\-म्वि\-द्रधौ}}{\lemma{\skt{विपं\appbrk{}च्या\appbrk{}म्वा\appbrk{}म\appbrk{}पा\appbrk{}र्श्व\appbrk{}क\appbrk{}क्ष\appbrk{}स्त\appbrk{}ना\appbrk{}न्त\appbrk{}र\appbrk{}योर्}}\Afootnote{H; \skt{विश्वाच्यां, श्रोणिं प्रति समन्ताद्द्व्यङ्गुले प्रवाहिकायां शूलिन्यां, परि\appbrk{}व\appbrk{}र्ति\appbrk{}को\appbrk{}प\appbrk{}दं\appbrk{}श\appbrk{}शू\appbrk{}क\appbrk{}दो\appbrk{}ष\appbrk{}शु\appbrk{}क्र\appbrk{}व्या\appbrk{}पत्सु मेढ्रमध्ये, (वृषणयोः पार्श्वे मूत्रवृद्ध्यां, नाभे\appbrk{}र\appbrk{}ध\appbrk{}श्च\appbrk{}तु\appbrk{}र\appbrk{}ङ्गुले सेवन्या वामपार्श्वे दकोदरे ,) वामपार्श्वे कक्षास्तनयोर्}\ A}}}{\lemma{\skt{अन्तरस्थाम्}}\Afootnote{H; \skt{अन्तरेऽन्तर्}\ A}}\skt{ }\edtext{\skt{पार्श्वशूले च}}{\lemma{\skt{पार्श्वशूले च}}\Afootnote{H; \skt{पार्श्वशूले}\ A}}\skt{~। }\edtext{\edtext{\skt{एतामेव}}{\lemma{\skt{एताम्}}\Afootnote{H; \emph{om.}\ A}}\skt{ च}}{\lemma{\skt{एव च}}\Afootnote{H; \skt{च}\ A}}\skt{ }\edtext{\skt{कफोदरे}}{\lemma{\skt{कफोदरे}}\Afootnote{H; \emph{om.}\ A}}\skt{~। }\edtext{\skt{दक्षिणपार्श्वे}}{\lemma{\skt{दक्षिणपार्श्वे}}\Afootnote{H; \emph{om.}\ A}}\skt{ }\edtext{\edtext{\skt{बाहु\-शो\-षा\-प\-बा\-हु\-क\-यो\-र्ब्बा\-हु\-मध्ये}}{\lemma{\skt{बाहु\appbrk{}शो\appbrk{}षा\appbrk{}प\appbrk{}बा\appbrk{}हु\appbrk{}क\appbrk{}योर्}}\Afootnote{H; \skt{बाहु\appbrk{}शो\appbrk{}षा\appbrk{}व\appbrk{}बा\appbrk{}हु\appbrk{}क\appbrk{}योर्}\ A}}}{\lemma{\skt{ब्बाहुमध्ये}}\Afootnote{H; \skt{अप्ये}\ A}}\skt{ }\edtext{\skt{एके}}{\lemma{\skt{एके}}\Afootnote{H; \skt{के}\ A}}\skt{ }\edtext{\edtext{\skt{वदन्त्यंसमध्ये}}{\lemma{\skt{वदन्त्य्}}\Afootnote{H; \skt{वदन्त्}\ A}}}{\lemma{\skt{अंसमध्ये}}\Afootnote{H; \skt{यंसयोरन्तरे}\ A}}\skt{ त्रिक\-स\-न्धि\-म\-ध्य\-ग\-ता\-न्तृ\-ती\-यके~। }\edtext{\edtext{\skt{स्कन्ध\-स\-न्धि\-ग\-ता\-ना\-म\-न्य\-त\-र\-पा\-र्श्व\-सं\-स्थि\-ता\-ना\-ञ्च\-तु\-र्थके}}{\lemma{\skt{स्कन्ध\appbrk{}स\appbrk{}न्धि\appbrk{}ग\appbrk{}ता\appbrk{}नाम्}}\Afootnote{H; \skt{अधः\appbrk{}स्क\appbrk{}न्ध\appbrk{}स\appbrk{}न्धि\appbrk{}गता}\ A}}}{\lemma{\skt{अन्य\appbrk{}त\appbrk{}र\appbrk{}पा\appbrk{}र्श्व\appbrk{}सं\appbrk{}स्थि\appbrk{}ता\appbrk{}नाञ्}}\Afootnote{H; \skt{मन्य\appbrk{}त\appbrk{}र\appbrk{}पा\appbrk{}र्श्व\appbrk{}सं\appbrk{}स्थितां}\ A}}\skt{~। }\edtext{\edtext{\skt{हनु\-स\-न्धि\-म\-ध्य\-ग\-ता\-म\-प\-स्मारे}}{\lemma{\skt{हनु\appbrk{}स\appbrk{}न्धि\appbrk{}म\appbrk{}ध्य\appbrk{}ग\appbrk{}ताम्}}\Afootnote{H; \skt{हनु\appbrk{}स\appbrk{}न्धि\appbrk{}म\appbrk{}ध्य\appbrk{}गता}\ A}}}{\lemma{\skt{अपस्मारे}}\Afootnote{H; \skt{मपस्मारे}\ A}}\skt{~। }\edtext{\edtext{\skt{शंख\-स\-न्धि\-ग\-ता\-ना\-मु\-न्मादे}}{\lemma{\skt{शंख\appbrk{}स\appbrk{}न्धि\appbrk{}ग\appbrk{}ता\appbrk{}नाम्}}\Afootnote{H; \skt{शङ्ख\appbrk{}के\appbrk{}शा\appbrk{}न्त\appbrk{}स\appbrk{}न्धि\appbrk{}ग\appbrk{}ता\appbrk{}मु\appbrk{}रोऽ\appbrk{}पा\appbrk{}ङ्ग\appbrk{}ल\appbrk{}ला\appbrk{}टेषु}\ A}}}{\lemma{\skt{उन्मादे}}\Afootnote{H; \skt{चोन्मादे}\ A}}\skt{~। जिह्वा\-रो\-गे\-ष्व\-धो\-जि\-ह्वायां~। दन्तव्याधिषु }\edtext{\skt{दन्तमूले}}{\lemma{\skt{दन्तमूले}}\Afootnote{H; \skt{च}\ A}}\skt{~। तालुनि }\edtext{\skt{तालव्ये च}}{\lemma{\skt{तालव्ये च}}\Afootnote{H; \skt{तालव्येषु}\ A}}\skt{~। कर्ण्णयोरुपरि सम\-न्ता\-त्क\-र्ण्ण\-शूले }\edtext{\skt{तद्रोगे च}}{\lemma{\skt{तद्रोगे च}}\Afootnote{H; \skt{तद्रोगेषु च}\ A}}\skt{~। }\edtext{\skt{गन्धाघ्राणे}}{\lemma{\skt{गन्धाघ्राणे}}\Afootnote{H; \skt{गन्धाग्रहणे}\ A}}\skt{ }\edtext{\skt{नासारोगेषु}}{\lemma{\skt{नासारोगेषु}}\Afootnote{H; \skt{नासारोगेषु च}\ A}}\skt{ नासाग्रे~। }\edtext{\edtext{\skt{तिमि\-रा\-क्षि\-पा\-क\-प्र\-भृ\-ति\-ष्व\-क्ष्चा\-म\-ये\-षू\-प\-नास्यां}}{\lemma{\skt{तिमि\appbrk{}रा\appbrk{}क्षि\appbrk{}पा\appbrk{}क\appbrk{}प्र\appbrk{}भृ\appbrk{}तिष्व्}}\Afootnote{H; \skt{तिमि\appbrk{}रा\appbrk{}क्षि\appbrk{}पा\appbrk{}क\appbrk{}प्र\appbrk{}भृ\appbrk{}तिष्}\ A}}}{\lemma{\skt{अक्ष्चा\appbrk{}म\appbrk{}ये\appbrk{}षू\appbrk{}प\appbrk{}नास्यां}}\Afootnote{H; \skt{वक्ष्या\appbrk{}म\appbrk{}ये\appbrk{}षू\appbrk{}प\appbrk{}ना\appbrk{}सिके}\ A}}\skt{ }\edtext{\edtext{\skt{लालाटीमापाङ्गी}}{\lemma{\skt{लालाटीम्}}\Afootnote{H; \skt{लालाट्}\ A}}\skt{ वा}}{\lemma{\skt{आपाङ्गी वा}}\Afootnote{H; \skt{यामपाङ्ग्यां वा}\ A}}\skt{~। }\edtext{\skt{एतामेव}}{\lemma{\skt{एताम्}}\Afootnote{H; \skt{एता}\ A}}\skt{ }\edtext{\skt{शिरो\-रो\-गा\-धि\-म\-न्थ\-प्र\-भृ\-तिषु च}}{\lemma{\skt{शिरो\appbrk{}रो\appbrk{}गा\appbrk{}धि\appbrk{}म\appbrk{}न्थ\appbrk{}प्र\appbrk{}भृ\appbrk{}तिषु च}}\Afootnote{H; \skt{शिरो\appbrk{}रो\appbrk{}गा\appbrk{}धि\appbrk{}म\appbrk{}न्थ\appbrk{}प्र\appbrk{}भृ\appbrk{}तिषु}\ A}}\skt{ }\edtext{\skt{रोगेषु}}{\lemma{\skt{रोगेषु}}\Afootnote{H; \skt{रोगेष्विति}\ A}}\skt{~।}\pend

\pstart \ledsidenote{\passref{3.8.18}}\edtext{\skt{तत्र}}{\lemma{\skt{तत्र}}\Afootnote{H; \skt{दुष्टव्यधा विंशतिः}\ A}}\skt{ दुर्व्वि\-द्धा\-ति\-वि\-द्धा\-कु\-ञ्चिता पिच्चिता कुट्टिता }\edtext{\skt{अप्र\-स्रु\-ता\-नु\-तीर्ण्णा}}{\lemma{\skt{अप्र\appbrk{}स्रु\appbrk{}ता\appbrk{}नु\appbrk{}तीर्ण्णा}}\Afootnote{H; \skt{प्रस्रु\appbrk{}ता\appbrk{}त्यु\appbrk{}दी\appbrk{}र्णान्}\ A}}\skt{ }\edtext{\skt{अन्ते\-वि\-द्धा\-प\-रि\-शु\-ष्का\-कू\-णिता}}{\lemma{\skt{अन्ते\appbrk{}वि\appbrk{}द्धा\appbrk{}प\appbrk{}रि\appbrk{}शु\appbrk{}ष्का\appbrk{}कू\appbrk{}णिता}}\Afootnote{H; \skt{ते विद्धा परिशुष्का कूणिता}\ A}}\skt{ वेपिता }\edtext{\skt{अनुत्थितविद्धा}}{\lemma{\skt{अनुत्थितविद्धा}}\Afootnote{H; \skt{नुत्थित}\ A}}\skt{~। }\edtext{\skt{अव्यध्यविद्धा}}{\lemma{\skt{अव्यध्यविद्धा}}\Afootnote{H; \skt{विद्धा}\ A}}\skt{ शस्त्रहता तिर्यग्विद्धा~। }\edtext{\skt{अविद्धा}}{\lemma{\skt{अविद्धा}}\Afootnote{H; \skt{विद्धा}\ A}}\skt{ }\edtext{\skt{विद्रुता}}{\lemma{\skt{विद्रुता}}\Afootnote{H; \skt{व्यध्या विद्रुता}\ A}}\skt{ धेनुका पुनः पुनर्व्विद्धा }\edtext{\skt{मर्म्म\-सि\-रा\-स्ना\-य्व\-स्थि\-स\-न्धिषु}}{\lemma{\skt{मर्म्म\appbrk{}सि\appbrk{}रा\appbrk{}स्ना\appbrk{}य्व\appbrk{}स्थि\appbrk{}स\appbrk{}न्धिषु}}\Afootnote{H; \skt{सिरा\appbrk{}स्ना\appbrk{}य्व\appbrk{}स्थि\appbrk{}स\appbrk{}न्धिम}\ A}}\skt{ }\edtext{\skt{विद्धा}}{\lemma{\skt{विद्धा}}\Afootnote{H; \skt{र्मसु}\ A}}\skt{ चेति~। }\edtext{\skt{विंशतिं}}{\lemma{\skt{विंशतिं}}\Afootnote{H; \emph{om.}\ A}}\skt{ }\edtext{\skt{दुष्टव्यध्या}}{\lemma{\skt{दुष्टव्यध्या}}\Afootnote{H; \emph{om.}\ A}}\skt{ }\edtext{\skt{भवन्ति}}{\lemma{\skt{भवन्ति}}\Afootnote{H; \emph{om.}\ A}}\skt{~॥}\pend

\pstart \ledsidenote{\passref{3.8.19}}\edtext{\skt{तत्र}}{\lemma{\skt{तत्र}}\Afootnote{H; \skt{तत्र या}\ A}}\skt{ }\edtext{\edtext{\skt{सूक्ष्म\-श\-स्त्र\-प्र\-णि\-धा\-ना\-न्ना\-सृ\-ग्व्यक्तं}}{\lemma{\skt{सूक्ष्म\appbrk{}श\appbrk{}स्त्र\appbrk{}प्र\appbrk{}णि\appbrk{}धा\appbrk{}नान्}}\Afootnote{H; \skt{सूक्ष्मशस्त्रविद्धा}\ A}}}{\lemma{\skt{नासृग्व्यक्तं}}\Afootnote{H; \skt{व्यक्त}\ A}}\skt{ }\edtext{\skt{स्रवति}}{\lemma{\skt{स्रवति}}\Afootnote{H; \skt{मसृक्स्रवति}\ A}}\skt{~। }\edtext{\skt{तत्र}}{\lemma{\skt{तत्र}}\Afootnote{H; \emph{om.}\ A}}\skt{ रुजा }\edtext{\skt{शोफौ}}{\lemma{\skt{शोफौ}}\Afootnote{H; \skt{शोफ}\ A}}\skt{ }\edtext{\skt{भवतः~। सा}}{\lemma{\skt{भवतः~। सा}}\Afootnote{H; \skt{वती च सा}\ A}}\skt{ दुर्व्विद्धा }\edtext{\skt{नाम}}{\lemma{\skt{नाम}}\Afootnote{H; \emph{om.}\ A}}\skt{~। प्रमा\-णा\-ति\-रि\-क्त\-वि\-द्धा\-या\-मन्तः }\edtext{\edtext{\skt{शोणि\-त\-म्प्र\-वि\-श\-त्य\-ति\-प्र\-वृ\-त्तिर्व्वा}}{\lemma{\skt{शोणितम्}}\Afootnote{H; \emph{om.}\ A}}}{\lemma{\skt{प्रवि\appbrk{}श\appbrk{}त्य\appbrk{}ति\appbrk{}प्र\appbrk{}वृ\appbrk{}त्तिर्व्वा}}\Afootnote{H; \skt{प्रविशति शोणितं शोणि\appbrk{}ता\appbrk{}ति\appbrk{}प्र\appbrk{}वृ\appbrk{}त्तिर्वा}\ A}}\skt{ }\edtext{\skt{भवति}}{\lemma{\skt{भवति}}\Afootnote{H; \emph{om.}\ A}}\skt{ सातिविद्धा }\edtext{\skt{नाम}}{\lemma{\skt{नाम}}\Afootnote{H; \emph{om.}\ A}}\skt{ }\edtext{\edtext{\skt{कुञ्चि\-ता\-या\-म\-प्येवं}}{\lemma{\skt{कुञ्चितायाम्}}\Afootnote{H; \skt{कुञ्चिताया}\ A}}}{\lemma{\skt{अप्येवं}}\Afootnote{H; \skt{मप्येवं}\ A}}\skt{~। }\edtext{\edtext{\skt{कुण्ठ\-श\-स्त्र\-प्र\-म\-थि\-ता\-पृ\-थु\-ल\-भा\-व\-मा\-पन्ना}}{\lemma{\skt{कुण्ठ\appbrk{}श\appbrk{}स्त्र\appbrk{}प्र\appbrk{}म\appbrk{}थि\appbrk{}ता\appbrk{}पृ\appbrk{}थु\appbrk{}ल\appbrk{}भा\appbrk{}वम्}}\Afootnote{H; \skt{कुण्ठ\appbrk{}श\appbrk{}स्त्र\appbrk{}प्र\appbrk{}म\appbrk{}थिता पृथुलीभाव}\ A}}}{\lemma{\skt{आपन्ना}}\Afootnote{H; \skt{मापन्ना}\ A}}\skt{ पिच्चिता }\edtext{\skt{नाम}}{\lemma{\skt{नाम}}\Afootnote{H; \emph{om.}\ A}}\skt{ }\edtext{\skt{अना\-सा\-दि\-त\-म\-न्त\-योश्च}}{\lemma{\skt{अना\appbrk{}सा\appbrk{}दि\appbrk{}त\appbrk{}म\appbrk{}न्त\appbrk{}योश्च}}\Afootnote{H; \skt{अनासादिता पुनः पुनरन्तयोश्च}\ A}}\skt{ बहुशः~। }\edtext{\skt{कुट्टिताकुट्टिता}}{\lemma{\skt{कुट्टिताकुट्टिता}}\Afootnote{H; \skt{शस्त्राभिहता कुट्टिता}\ A}}\skt{ }\edtext{\skt{नाम}}{\lemma{\skt{नाम}}\Afootnote{H; \emph{om.}\ A}}\skt{ }\edtext{\skt{शीत\-भ\-या\-मू\-र्च्छा\-भि\-र\-प्र\-वृ\-त्त\-शो\-णिता}}{\lemma{\skt{शीत\appbrk{}भ\appbrk{}या\appbrk{}मू\appbrk{}र्च्छा\appbrk{}भिर्}}\Afootnote{H; \skt{शीत\appbrk{}भ\appbrk{}य\appbrk{}मू\appbrk{}र्च्छा\appbrk{}भिर्}\ A}}\skt{ }\edtext{\skt{अप्रस्रूता}}{\lemma{\skt{अप्रस्रूता}}\Afootnote{H; \skt{अप्रस्रुता}\ A}}\skt{ }\edtext{\skt{नाम}}{\lemma{\skt{नाम}}\Afootnote{H; \emph{om.}\ A}}\skt{ तीक्ष्ण\-म\-हा\-मु\-ख\-श\-स्त्र\-विद्धा }\edtext{\skt{अनुदीर्ण्णा}}{\lemma{\skt{अनुदीर्ण्णा}}\Afootnote{H; \skt{त्युदीर्णा}\ A}}\skt{ }\edtext{\skt{नाम}}{\lemma{\skt{नाम}}\Afootnote{H; \emph{om.}\ A}}\skt{ }\edtext{\skt{अल्प\-र\-क्त\-स्रा\-वि\-न्य\-न्ते\-विद्धा}}{\lemma{\skt{अल्प\appbrk{}र\appbrk{}क्त\appbrk{}स्रा\appbrk{}वि\appbrk{}न्य\appbrk{}न्ते\appbrk{}विद्धा}}\Afootnote{H; \skt{अल्प\appbrk{}र\appbrk{}क्त\appbrk{}स्रा\appbrk{}वि\appbrk{}ण्य\appbrk{}न्ते\appbrk{}विद्धा}\ A}}\skt{ }\edtext{\skt{नाम।}}{\lemma{\skt{नाम}}\Afootnote{H; \emph{om.}\ A}}\skt{ }\edtext{\skt{क्षीण\-शो\-णि\-ता\-नि\-ल\-पू\-र्ण्णा\-प\-रि\-शुष्का}}{\lemma{\skt{क्षीण\appbrk{}शो\appbrk{}णि\appbrk{}ता\appbrk{}नि\appbrk{}ल\appbrk{}पू\appbrk{}र्ण्णा\appbrk{}प\appbrk{}रि\appbrk{}शुष्का}}\Afootnote{H; \skt{क्षीण\appbrk{}शो\appbrk{}णि\appbrk{}त\appbrk{}स्या\appbrk{}नि\appbrk{}ल\appbrk{}पूर्णा परिशुष्का}\ A}}\skt{ }\edtext{\skt{नाम}}{\lemma{\skt{नाम}}\Afootnote{H; \emph{om.}\ A}}\skt{~। }\edtext{\skt{चतु\-र्थ\-भा\-ग\-मा\-सा\-दिता}}{\lemma{\skt{चतु\appbrk{}र्थ\appbrk{}भा\appbrk{}ग\appbrk{}मा\appbrk{}सा\appbrk{}दिता}}\Afootnote{H; \skt{चतु\appbrk{}र्भा\appbrk{}गा\appbrk{}सा\appbrk{}दिता}\ A}}\skt{ किञ्चि\-त्प्र\-वृ\-त्त\-शो\-णि\-ता\-कू\-णिता }\edtext{\skt{नाम}}{\lemma{\skt{नाम}}\Afootnote{H; \emph{om.}\ A}}\skt{~। }\edtext{\skt{दुत्था\-न\-ब\-न्धो\-द्वे\-षि\-ण्या\-शो\-णि\-त\-सं\-मोहो}}{\lemma{\skt{दुत्था\appbrk{}न\appbrk{}ब\appbrk{}न्धो\appbrk{}द्वे\appbrk{}षि\appbrk{}ण्या\appbrk{}शो\appbrk{}णि\appbrk{}त\appbrk{}सं\appbrk{}मोहो}}\Afootnote{\skt{दुःस्था\appbrk{}न\appbrk{}ब\appbrk{}न्ध\appbrk{}ना\appbrk{}द्वे\appbrk{}प\appbrk{}मा\appbrk{}नायाः शोणितसंमोहो}\ A; \skt{दुत्था\appbrk{}न\appbrk{}ब\appbrk{}न्धो\appbrk{}द्वे\appbrk{}पि\appbrk{}ण्या\appbrk{}शो\appbrk{}णि\appbrk{}त\appbrk{}सं\appbrk{}मोहो}\ H}}\skt{ भवति~। सा वेपिता }\edtext{\skt{नाम}}{\lemma{\skt{नाम}}\Afootnote{H; \emph{om.}\ A}}\skt{~। }\edtext{\edtext{\skt{अनु\-स्थि\-ता\-या\-म\-प्ये\-व\-म\-श\-स्त्र\-कृत्या}}{\lemma{\skt{अनुस्थितायाम्}}\Afootnote{H; \skt{अनु\appbrk{}त्थि\appbrk{}त\appbrk{}वि\appbrk{}द्धाया}\ A}}}{\lemma{\skt{अप्ये\appbrk{}व\appbrk{}म\appbrk{}श\appbrk{}स्त्र\appbrk{}कृत्या}}\Afootnote{H; \skt{मप्येवं}\ A}}\skt{ }\edtext{\skt{अव्यध्या}}{\lemma{\skt{अव्यध्या}}\Afootnote{H; \emph{om.}\ A}}\skt{ }\edtext{\skt{नाम}}{\lemma{\skt{नाम}}\Afootnote{H; \emph{om.}\ A}}\skt{~। }\edtext{\skt{छित्वा\-ति\-प्र\-वृ\-त्त\-शो\-णिता}}{\lemma{\skt{छित्वा\appbrk{}ति\appbrk{}प्र\appbrk{}वृ\appbrk{}त्त\appbrk{}शो\appbrk{}णिता}}\Afootnote{H; \skt{छिन्ना\appbrk{}ति\appbrk{}प्र\appbrk{}वृ\appbrk{}त्त\appbrk{}शो\appbrk{}णिता}\ A}}\skt{ }\edtext{\skt{क्रिया\-स\-ङ्ग्र\-ह\-करी}}{\lemma{\skt{क्रिया\appbrk{}स\appbrk{}ङ्ग्र\appbrk{}ह\appbrk{}करी}}\Afootnote{H; \skt{क्रियासङ्गकरी}\ A}}\skt{ शस्त्रहता }\edtext{\skt{नाम}}{\lemma{\skt{नाम}}\Afootnote{H; \emph{om.}\ A}}\skt{~। तिर्य\-क्प्र\-णि\-हि\-त\-श\-स्त्रा\-कि\-ञ्चि\-च्छे\-षा\-ति\-र्य\-ग्विद्धा }\edtext{\skt{नाम}}{\lemma{\skt{नाम}}\Afootnote{H; \emph{om.}\ A}}\skt{~। बहुशः क्षता हीन\-श\-स्त्र\-प्र\-णि\-धा\-ने\-ना\-विद्धा }\edtext{\skt{नाम}}{\lemma{\skt{नाम}}\Afootnote{H; \skt{अशस्त्रकृत्या अव्य}\ A}}\skt{~। }\edtext{\skt{अन\-व\-स्थि\-त\-विद्धा}}{\lemma{\skt{अन\appbrk{}व\appbrk{}स्थि\appbrk{}त\appbrk{}विद्धा}}\Afootnote{H; \skt{ध्या, अन\appbrk{}व\appbrk{}स्थि\appbrk{}त\appbrk{}विद्धा}\ A}}\skt{ विद्रुता }\edtext{\skt{नाम}}{\lemma{\skt{नाम}}\Afootnote{H; \emph{om.}\ A}}\skt{~। }\edtext{\skt{व्यधप्रदेशस्य}}{\lemma{\skt{व्यधप्रदेशस्य}}\Afootnote{H; \skt{प्रदेशस्य}\ A}}\skt{ }\edtext{\skt{बहु\-शो\-व\-घ\-ट्ट\-ना\-द\-सं\-आ\-रो\-ह\-मा\-ण\-व्यधा}}{\lemma{\skt{बहु\appbrk{}शो\appbrk{}व\appbrk{}घ\appbrk{}ट्ट\appbrk{}ना\appbrk{}द\appbrk{}सं\appbrk{}आ\appbrk{}रो\appbrk{}ह\appbrk{}मा\appbrk{}ण\appbrk{}व्यधा}}\Afootnote{H; \skt{बहु\appbrk{}शोऽ\appbrk{}व\appbrk{}घ\appbrk{}ट्ट\appbrk{}ना\appbrk{}दा\appbrk{}रो\appbrk{}ह\appbrk{}द्व्यधा}\ A}}\skt{ }\edtext{\skt{रक्त\-स्रा\-वि\-न्य\-त्यर्थं}}{\lemma{\skt{रक्त\appbrk{}स्रा\appbrk{}वि\appbrk{}न्य\appbrk{}त्यर्थं}}\Afootnote{H; \skt{मुहुर्मुहुः शोणितस्रावा}\ A}}\skt{ धेनुका }\edtext{\skt{नाम}}{\lemma{\skt{नाम}}\Afootnote{H; \skt{सूक्ष्म\appbrk{}श\appbrk{}स्त्र\appbrk{}व्य\appbrk{}ध\appbrk{}ना\appbrk{}द्ब\appbrk{}हुशो भिन्ना}\ A}}\skt{~। पुनः }\edtext{\skt{पुनः}}{\lemma{\skt{पुनः}}\Afootnote{H; \skt{पुन}\ A}}\skt{ }\edtext{\skt{प्रणिधाना}}{\lemma{\skt{प्रणिधाना}}\Afootnote{H; \skt{र्}\ A}}\skt{ }\edtext{\skt{सिरा\-म\-र्म्म\-स्ना\-य्व\-स्थि\-सं\-धिषु}}{\lemma{\skt{सिरा\appbrk{}म\appbrk{}र्म्म\appbrk{}स्ना\appbrk{}य्व\appbrk{}स्थि\appbrk{}सं\appbrk{}धिषु}}\Afootnote{H; \skt{विद्धा, स्नाय्व\appbrk{}स्थि\appbrk{}सि\appbrk{}रा\appbrk{}स\appbrk{}न्धि\appbrk{}म\appbrk{}र्मसु}\ A}}\skt{ विद्धा }\edtext{\skt{रुजाशोफौ}}{\lemma{\skt{रुजाशोफौ}}\Afootnote{H; \skt{रुजां शोफं}\ A}}\skt{ }\edtext{\skt{वैक\-ल्य\-म्म\-र\-णम्वा}}{\lemma{\skt{वा}}\Afootnote{H; \skt{चापा}\ A}}\skt{ }\edtext{\skt{भवति}}{\lemma{\skt{भवति}}\Afootnote{H; \skt{दयति}\ A}}\skt{~॥}\pend

\pstart \ledsidenote{\passref{3.8.20}}\skt{भवन्ति चात्र~॥ सिरासुशिक्षितो नास्ति चलाह्येताः स्वभावतः~। मत्स्य\-व\-त्प\-रि\-व\-र्त्तन्ते तस्माद्यत्नेन ताडयेत्~॥}\pend

\pstart \ledsidenote{\passref{3.8.21}}\skt{अजानता गृहीते तु शस्त्रे कायनिपातिते~। भवन्ति }\edtext{\skt{व्यापदस्त्वेता}}{\lemma{\skt{व्यापदस्त्वेता}}\Afootnote{H; \skt{व्यापदश्चैता}\ A}}\skt{ बह\-व\-श्चा\-प्यु\-प\-द्रवाः~॥}\pend

\pstart \ledsidenote{\passref{3.8.22}}\skt{स्नेहादिभिः }\edtext{\edtext{\skt{क्रियायोगैर्न्न}}{\lemma{\skt{क्रियायोगैर्}}\Afootnote{H; \skt{क्रियायोगै}\ A}}}{\lemma{\skt{न्न}}\Afootnote{H; \skt{र्न}\ A}}\skt{ तथा लेपनैरपि~। यांत्याशु व्याधयः शान्तिं यथा सम्य\-क्सि\-रा\-व्य\-धात्~॥}\pend

\pstart \ledsidenote{\passref{3.8.23}}\edtext{\skt{सिराव्यधं}}{\lemma{\skt{सिराव्यधं}}\Afootnote{H; \skt{सिराव्यध}\ A}}\skt{ }\edtext{\skt{चिकित्सार्थं}}{\lemma{\skt{चिकित्सार्थं}}\Afootnote{H; \skt{श्चिकित्सार्धं}\ A}}\skt{ शल्यतन्त्रे }\edtext{\skt{विधीयते}}{\lemma{\skt{विधीयते}}\Afootnote{H; \skt{प्रकीर्तितः}\ A}}\skt{~। यथा }\edtext{\edtext{\skt{सम्यक्प्रणिहिता}}{\lemma{\skt{सम्यक्}}\Afootnote{H; \emph{om.}\ A}}}{\lemma{\skt{प्रणिहिता}}\Afootnote{H; \skt{प्रणिहितः}\ A}}\skt{ }\edtext{\skt{बस्तिः}}{\lemma{\skt{बस्तिः}}\Afootnote{H; \skt{सम्यग्बस्तिः}\ A}}\skt{ कायचिकित्सिते~॥}\pend

\pstart \ledsidenote{\passref{3.8.24}}\skt{तत्र }\edtext{\edtext{\skt{स्निग्ध\-स्वि\-न्न\-वा\-न्त\-वि\-रि\-क्ता\-नु\-वा\-सि\-ता\-स्था\-पि\-त\-वि\-द्ध\-सि\-रै\-र्म्म\-नुष्यैः}}{\lemma{\skt{स्निग्ध\appbrk{}स्वि\appbrk{}न्न\appbrk{}वा\appbrk{}न्त\appbrk{}वि\appbrk{}रि\appbrk{}क्ता\appbrk{}नु\appbrk{}वा\appbrk{}सि\appbrk{}ता\appbrk{}स्था\appbrk{}पि\appbrk{}त\appbrk{}वि\appbrk{}द्ध\appbrk{}सि\appbrk{}रैर्}}\Afootnote{H; \skt{स्निग्ध\appbrk{}स्वि\appbrk{}न्न\appbrk{}वा\appbrk{}न्त\appbrk{}वि\appbrk{}रि\appbrk{}क्ता\appbrk{}स्था\appbrk{}पि\appbrk{}ता\appbrk{}नु\appbrk{}व\appbrk{}सित}\ A}}}{\lemma{\skt{म्मनुष्यैः}}\Afootnote{H; \skt{सिराविद्धैः}\ A}}\skt{ परिहर्त्तव्यानि~॥ }\edtext{\skt{क्रोधा\-या\-स\-मै\-थु\-न\-दि\-वा\-स्व\-प्नो\-च्चै\-र्ब्भा\-ष्य\-व्या\-या\-म\-या\-ना\-स\-ना\-नि\-स्था\-न\-चं\-क्र\-म\-ण\-शी\-त\-वा\-ता\-त\-प\-चि\-रु\-द्धा\-सा\-त्म्या\-जी\-र्ण्णा\-न्या\-च\-ल\-ला\-भा\-त्मा\-स\-मेके}}{\lemma{\skt{क्रोधा\appbrk{}या\appbrk{}स\appbrk{}मै\appbrk{}थु\appbrk{}न\appbrk{}दि\appbrk{}वा\appbrk{}स्व\appbrk{}प्नो\appbrk{}च्चै\appbrk{}र्ब्भा\appbrk{}ष्य\appbrk{}व्या\appbrk{}या\appbrk{}म\appbrk{}या\appbrk{}ना\appbrk{}स\appbrk{}ना\appbrk{}नि\appbrk{}स्था\appbrk{}न\appbrk{}चं\appbrk{}क्र\appbrk{}म\appbrk{}ण\appbrk{}शी\appbrk{}त\appbrk{}वा\appbrk{}ता\appbrk{}त\appbrk{}प\appbrk{}चि\appbrk{}रु\appbrk{}द्धा\appbrk{}सा\appbrk{}त्म्या\appbrk{}जी\appbrk{}र्ण्णा\appbrk{}न्या\appbrk{}च\appbrk{}ल\appbrk{}ला\appbrk{}भा\appbrk{}त्मा\appbrk{}स\appbrk{}मेके}}\Afootnote{\skt{क्रोधा\appbrk{}या\appbrk{}स\appbrk{}मै\appbrk{}थु\appbrk{}न\appbrk{}दि\appbrk{}वा\appbrk{}स्व\appbrk{}प्न\appbrk{}वा\appbrk{}ग्व्या\appbrk{}या\appbrk{}म\appbrk{}या\appbrk{}ना\appbrk{}ध्य\appbrk{}य\appbrk{}न\appbrk{}स्था\appbrk{}ना\appbrk{}स\appbrk{}न\appbrk{}च\appbrk{}ङ्क्र\appbrk{}म\appbrk{}ण\appbrk{}शी\appbrk{}त\appbrk{}वा\appbrk{}ता\appbrk{}त\appbrk{}प\appbrk{}वि\appbrk{}रु\appbrk{}द्धा\appbrk{}सा\appbrk{}त्म्या\appbrk{}जी\appbrk{}र्णा\appbrk{}न्या\appbrk{}ब\appbrk{}ल\appbrk{}ला\appbrk{}भात्, मासमेके}\ A; \skt{क्रोधा\appbrk{}या\appbrk{}स\appbrk{}मै\appbrk{}थु\appbrk{}न\appbrk{}दि\appbrk{}वा\appbrk{}स्व\appbrk{}प्नो\appbrk{}च्चै\appbrk{}र्ब्भा\appbrk{}ष्य\appbrk{}व्या\appbrk{}या\appbrk{}म\appbrk{}या\appbrk{}ना\appbrk{}स\appbrk{}नानि स्थान\appbrk{}चं\appbrk{}क्र\appbrk{}म\appbrk{}ण\appbrk{}शी\appbrk{}त\appbrk{}वा\appbrk{}ता\appbrk{}त\appbrk{}प\appbrk{}चि\appbrk{}रु\appbrk{}द्धा\appbrk{}सा\appbrk{}त्म्या\appbrk{}जी\appbrk{}र्ण्णा\appbrk{}न्या\appbrk{}ब\appbrk{}ल\appbrk{}ला\appbrk{}भा\appbrk{}त्मा\appbrk{}स\appbrk{}मेके}\ H}}\skt{ }\edtext{\skt{मन्यन्ते}}{\lemma{\skt{मन्यन्ते}}\Afootnote{H; \skt{मन्यन्ते~। एतेषां}\ A}}\skt{~। }\edtext{\edtext{\skt{विस्त\-र\-मे\-षा\-मु\-प\-रि\-ष्टा\-द्व\-क्ष्यामः}}{\lemma{\skt{एषाम्}}\Afootnote{H; \emph{om.}\ A}}}{\lemma{\skt{उपरिष्टाद्}}\Afootnote{H; \skt{उपरिष्ठाद्}\ A}}\skt{~॥}\pend

\pstart\versepar \ledsidenote{\passref{3.8.25}}\edtext{\skt{भवति}}{\lemma{\skt{भवति}}\Afootnote{H; \emph{om.}\ A}}\skt{ }\edtext{\skt{चात्र}}{\lemma{\skt{चात्र}}\Afootnote{H; \emph{om.}\ A}}\skt{~॥ }\pend
\pstart\versepar \skt{}\edtext{\skt{सिरा\-वि\-षा\-ण\-तु\-म्बी\-भि\-र्ज्ज\-लौ\-काभिः}}{\lemma{\skt{सिरा\appbrk{}वि\appbrk{}षा\appbrk{}ण\appbrk{}तु\appbrk{}म्बी\appbrk{}भिर्}}\Afootnote{H; \skt{सिरा\appbrk{}वि\appbrk{}षा\appbrk{}ण\appbrk{}तु\appbrk{}म्बैस्तु}\ A}}\skt{ पदैस्तथा~। अवगाढं यथापूर्व्वं }\edtext{\skt{क्रियते}}{\lemma{\skt{क्रियते}}\Afootnote{H; \skt{निर्हरे}\ A}}\skt{ }\edtext{\skt{दुष्टशोणितं}}{\lemma{\skt{दुष्टशोणितं}}\Afootnote{H; \skt{द्दुष्टशोणितम्}\ A}}\skt{~॥}\pend

\pstart \ledsidenote{\passref{3.8.trailer}}\edtext{\skt{इति }\edtext{\skt{शारी\-रेऽ\-ष्ट\-मोऽ\-ध्यायः}}{\lemma{\skt{ऽष्टमो}}\Afootnote{H; \skt{नामाष्टमो}\ A}}}{\lemma{\skt{इति शारीरे}}\Afootnote{H; \skt{इति सुश्रु\appbrk{}त\appbrk{}सं\appbrk{}हि\appbrk{}तायां शारीरस्थाने सिरा\appbrk{}व्य\appbrk{}ध\appbrk{}वि\appbrk{}धि\appbrk{}शा\appbrk{}रीरं}\ A}}\skt{~॥}\pend

\endnumbering
\beginnumbering
\eledsection*{Adhyāya 9}
\pstart \ledsidenote{\passref{3.9.1}}\skt{अथातो धम\-नी\-व्या\-क\-रणं शारीरं व्याख्यास्यामः~॥}\pend

\pstart \ledsidenote{\passref{3.9.3}}\skt{चतु\-र्व्विं\-श\-ति\-र्द्ध\-मन्यो }\edtext{\skt{भवन्ति~। तास्तु}}{\lemma{\skt{भवन्ति~। तास्तु}}\Afootnote{H; \emph{om.}\ A}}\skt{ }\edtext{\skt{नाभिप्रभवाः}}{\lemma{\skt{नाभिप्रभवाः}}\Afootnote{H; \skt{नाभिप्रभवा अभिहिताः}\ A}}\skt{~। तत्र }\edtext{\skt{केचि\-दा\-हु\-रा\-चार्याः}}{\lemma{\skt{आचार्याः}}\Afootnote{H; \emph{om.}\ A}}\skt{ }\edtext{\edtext{\skt{सिरा\-ध\-म\-नी\-स्रो\-त\-सा\-म\-वि\-भागः}}{\lemma{\skt{सिरा\appbrk{}ध\appbrk{}म\appbrk{}नी\appbrk{}स्रो\appbrk{}त\appbrk{}साम्}}\Afootnote{H; \skt{सिरा\appbrk{}ध\appbrk{}म\appbrk{}नी\appbrk{}स्रो\appbrk{}तसा}\ A}}}{\lemma{\skt{अविभागः}}\Afootnote{H; \skt{मविभागः}\ A}}\skt{ सिराविकारा एव हि धमन्यः स्रोतान्सि चेति~। तत्तु न }\edtext{\skt{सम्यगत्रोच्यते}}{\lemma{\skt{सम्यगत्रोच्यते}}\Afootnote{H; \skt{सम्यक्, अ}\ A}}\skt{~। }\edtext{\skt{अन्या}}{\lemma{\skt{अन्या}}\Afootnote{H; \skt{न्या}\ A}}\skt{ एव हि धमन्यः }\edtext{\skt{स्रोतान्सि}}{\lemma{\skt{स्रोतान्सि}}\Afootnote{H; \skt{स्रोतांसि च}\ A}}\skt{ }\edtext{\skt{सिराश्चेति}}{\lemma{\skt{सिराश्चेति}}\Afootnote{H; \skt{सिराभ्यः}\ A}}\skt{~। कस्माद्~। }\edtext{\skt{व्यञ्ज\-ना\-न्य\-त्वा\-त्मू\-ल\-स\-न्नि\-य\-मा\-त्क\-र्म्म\-वै\-शे\-ष्या\-दा\-ग\-माच्च}}{\lemma{\skt{व्यञ्जनान्यत्वात्}}\Afootnote{H; \skt{व्यञ्जनान्यत्वान्}\ A}}\skt{~। केवलं तु }\edtext{\edtext{\skt{पर\-स्प\-र\-प्र\-ता\-न\-स\-न्नि\-क\-र्षा\-त्स\-दृ\-श\-क\-र्म्म\-त्वा\-त्सौ\-क्ष्म्याच्च}}{\lemma{\skt{पर\appbrk{}स्प\appbrk{}र\appbrk{}प्र\appbrk{}ता\appbrk{}न\appbrk{}स\appbrk{}न्नि\appbrk{}क\appbrk{}र्षात्}}\Afootnote{H; \skt{पर\appbrk{}स्प\appbrk{}र\appbrk{}स\appbrk{}न्नि\appbrk{}क\appbrk{}र्षात्}\ A}}}{\lemma{\skt{सदृ\appbrk{}श\appbrk{}क\appbrk{}र्म्म\appbrk{}त्वात्}}\Afootnote{H; \skt{सदृ\appbrk{}शा\appbrk{}ग\appbrk{}म\appbrk{}क\appbrk{}र्म\appbrk{}त्वात्}\ A}}\skt{ विभ\-क्त\-क\-र्म्म\-णा\-म\-प्य\-वि\-भाग }\edtext{\skt{इव}}{\lemma{\skt{इव}}\Afootnote{A; \skt{एव}\ H}}\skt{ कर्म्मसु भवति~।}\pend

\pstart \ledsidenote{\passref{3.9.4}}\edtext{\skt{तासां}}{\lemma{\skt{तासां}}\Afootnote{H; \skt{तासां तु}\ A}}\skt{ खलु }\edtext{\edtext{\skt{नाभि\-प्र\-भ\-वा\-णा\-मू\-र्द्ध्वगा}}{\lemma{\skt{नाभि\appbrk{}प्र\appbrk{}भ\appbrk{}वा\appbrk{}णाम्}}\Afootnote{A; \skt{नाभिप्रभवा}\ H}}}{\lemma{\skt{ऊर्द्ध्वगा}}\Afootnote{\skt{धम\appbrk{}नी\appbrk{}ना\appbrk{}मू\appbrk{}र्ध्वगा}\ A; \skt{नामूर्द्ध्वगा}\ H}}\skt{ दश दश चाधो\-गा\-मि\-न्य\-श्च\-त\-स्र\-स्ति\-र्यग्गाः~।}\pend

\pstart \ledsidenote{\passref{3.9.5}}\skt{ऊर्द्ध्वगाः }\edtext{\skt{खलु}}{\lemma{\skt{खलु}}\Afootnote{H; \emph{om.}\ A}}\skt{ }\edtext{\skt{शब्द\-स्प\-र्श\-रू\-प\-र\-स\-ग\-न्ध\-प्र\-श्वा\-सो\-च्छ्वा\-स\-प्र\-क्षि\-त\-जृ\-म्भि\-त\-ह\-सि\-त\-क\-थि\-त\-रु\-दि\-ता\-दी\-न्वि\-शे\-षा\-न\-भि\-व\-हन्त्यः}}{\lemma{\skt{शब्द\appbrk{}स्प\appbrk{}र्श\appbrk{}रू\appbrk{}प\appbrk{}र\appbrk{}स\appbrk{}ग\appbrk{}न्ध\appbrk{}प्र\appbrk{}श्वा\appbrk{}सो\appbrk{}च्छ्वा\appbrk{}स\appbrk{}प्र\appbrk{}क्षि\appbrk{}त\appbrk{}जृ\appbrk{}म्भि\appbrk{}त\appbrk{}ह\appbrk{}सि\appbrk{}त\appbrk{}क\appbrk{}थि\appbrk{}त\appbrk{}रु\appbrk{}दि\appbrk{}ता\appbrk{}दीन्}}\Afootnote{\skt{शब्द\appbrk{}स्प\appbrk{}र्श\appbrk{}रू\appbrk{}प\appbrk{}र\appbrk{}स\appbrk{}ग\appbrk{}न्ध\appbrk{}प्र\appbrk{}श्वा\appbrk{}सो\appbrk{}च्छ्वा\appbrk{}स\appbrk{}जृ\appbrk{}म्भि\appbrk{}त\appbrk{}क्षु\appbrk{}द्ध\appbrk{}सि\appbrk{}त\appbrk{}क\appbrk{}थि\appbrk{}त\appbrk{}रु\appbrk{}दि\appbrk{}ता\appbrk{}दीन्}\ A; \skt{शब्द\appbrk{}स्प\appbrk{}र्श\appbrk{}रू\appbrk{}प\appbrk{}र\appbrk{}स\appbrk{}ग\appbrk{}न्ध\appbrk{}प्र\appbrk{}श्वा\appbrk{}सो\appbrk{}च्छ्वा\appbrk{}स\appbrk{}प्रे\appbrk{}क्षि\appbrk{}त\appbrk{}क्षु\appbrk{}त्जृ\appbrk{}म्भि\appbrk{}त\appbrk{}ह\appbrk{}सि\appbrk{}त\appbrk{}क\appbrk{}थि\appbrk{}त\appbrk{}रु\appbrk{}दि\appbrk{}ता\appbrk{}दीन्}\ H}}\skt{ शरी\-र\-न्धा\-र\-यन्ति~। तास्तु }\edtext{\edtext{\edtext{\skt{हृद\-य\-म\-भि\-प्र\-प\-न्ना\-स्त्रिधा}}{\lemma{\skt{हृदयम्}}\Afootnote{H; \skt{हृदय}\ A}}}{\lemma{\skt{अभिप्रपन्नास्}}\Afootnote{H; \skt{मभिप्रपन्ना}\ A}}}{\lemma{\skt{त्रिधा}}\Afootnote{H; \skt{स्}\ A}}\skt{ त्रिधा }\edtext{\skt{जायन्ते~। }\edtext{\skt{तास्त्रिंशत्तासां}}{\lemma{\skt{त्रिंशत्}}\Afootnote{H; \skt{स्त्रिंशत्}\ A}}}{\lemma{\skt{जायन्ते~। तास्}}\Afootnote{H; \skt{जायन्ते, ता}\ A}}\skt{ तु }\edtext{\edtext{\skt{वात\-पि\-त्त\-श्ले\-ष्म\-शो\-णि\-त\-र\-सान्द्वे}}{\lemma{\skt{वात\appbrk{}पि\appbrk{}त्त\appbrk{}श्ले\appbrk{}ष्म\appbrk{}शो\appbrk{}णि\appbrk{}त\appbrk{}र\appbrk{}सान्}}\Afootnote{H\textsuperscript{pc}; \skt{वात\appbrk{}पि\appbrk{}त्त\appbrk{}क\appbrk{}फ\appbrk{}शो\appbrk{}णि\appbrk{}त\appbrk{}र\appbrk{}सान्}\ A; \skt{वात\appbrk{}पि\appbrk{}त्त\appbrk{}श्ले\appbrk{}ष्म\appbrk{}शो\appbrk{}णि\appbrk{}त\appbrk{}र\appbrk{}सान्}\ H\textsuperscript{ac}}}}{\lemma{\skt{द्वे}}\Afootnote{A H\textsuperscript{pc}; \skt{वहत द्वे}\ H\textsuperscript{ac}}}\skt{ द्वे वहतस्ता दश~। }\edtext{\edtext{\skt{शब्द\-स्प\-र्श\-रू\-प\-र\-स\-ग\-न्धा\-न\-ष्टा\-भि\-र्गृ\-ह्णीते}}{\lemma{\skt{शब्द\appbrk{}स्प\appbrk{}र्श\appbrk{}रू\appbrk{}प\appbrk{}र\appbrk{}स\appbrk{}ग\appbrk{}न्धा\appbrk{}न\appbrk{}ष्टा\appbrk{}भिर्}}\Afootnote{H; \skt{शब्द\appbrk{}रू\appbrk{}प\appbrk{}र\appbrk{}स\appbrk{}ग\appbrk{}न्धा\appbrk{}न\appbrk{}ष्टा\appbrk{}भिर्}\ A}}}{\lemma{\skt{गृह्णीते}}\Afootnote{A; \skt{गृहीते}\ H}}\skt{~। द्वाभ्यां भाषते द्वाभ्यां घोषं करोति~। द्वाभ्यां स्वपिति द्वाभ्यां प्रतिबुध्यते~। द्वे }\edtext{\skt{अस्रवाहिन्यो}}{\lemma{\skt{अस्रवाहिन्यो}}\Afootnote{H; \skt{चाश्रुवाहिन्यौ}\ A}}\skt{ द्वे }\edtext{\skt{स्तन्यं}}{\lemma{\skt{स्तन्यं}}\Afootnote{A; \skt{स्तन्य}\ H}}\skt{ स्त्रिया वहतः }\edtext{\skt{स्तनसंश्रिता}}{\lemma{\skt{स्तनसंश्रिता}}\Afootnote{H; \skt{स्तनसंश्रिते, ते}\ A}}\skt{ एव~। शुक्रं नरस्य }\edtext{\edtext{\edtext{\edtext{\skt{स्तन\-यो\-र\-भि\-व\-ह\-त\-स्ता\-स्त्वे\-ता\-स्त्रिं\-श\-त्स\-वि\-भागा}}{\lemma{\skt{स्तनयोर्}}\Afootnote{H; \skt{स्तनाभ्या}\ A}}}{\lemma{\skt{अभि\appbrk{}व\appbrk{}ह\appbrk{}त\appbrk{}स्तास्}}\Afootnote{H; \skt{मभिवहतः, ता}\ A}}}{\lemma{\skt{त्वेतास्}}\Afootnote{H; \skt{स्त्वेता}\ A}}}{\lemma{\skt{त्रिंशत्}}\Afootnote{H; \skt{स्त्रिंशत्}\ A}}\skt{ व्याख्याताः~। एता\-भि\-रू\-र्द्ध्व\-न्ना\-भे\-रु\-द\-र\-पा\-र्श्व\-पृ\-ष्ठोरः }\edtext{\skt{स्कन्ध\-ग्री\-वा\-शी\-र्ष\-बा\-हवो}}{\lemma{\skt{स्कन्ध\appbrk{}ग्री\appbrk{}वा\appbrk{}शी\appbrk{}र्ष\appbrk{}बा\appbrk{}हवो}}\Afootnote{H; \skt{स्कन्ध\appbrk{}ग्री\appbrk{}वा\appbrk{}बा\appbrk{}हवो}\ A}}\skt{ धार्यन्ते }\edtext{\skt{याप्यन्ते}}{\lemma{\skt{याप्यन्ते}}\Afootnote{H; \skt{याप्यन्ते च}\ A}}\skt{~॥}\pend

\pstart\versepar \ledsidenote{\passref{3.9.11}}\edtext{\skt{भवति}}{\lemma{\skt{भवति}}\Afootnote{H; \emph{om.}\ A}}\skt{ }\edtext{\skt{चात्र}}{\lemma{\skt{चात्र}}\Afootnote{H; \emph{om.}\ A}}\skt{~॥ }\pend
\pstart\versepar \skt{}\edtext{\skt{पञ्चा\-धि\-भू\-ता\-स्त्वथ}}{\lemma{\skt{पञ्चाधिभूतास्}}\Afootnote{H; \skt{पञ्चाभिभूतास्}\ A}}\skt{ पञ्चकृत्वः }\pend
\pstart\versepar \skt{}\edtext{\skt{पञ्चे\-न्द्रि\-य\-म्प\-ञ्चाधा}}{\lemma{\skt{पञ्चाधा}}\Afootnote{\skt{पञ्चसु}\ A; \skt{पञ्चाध}\ H}}\skt{ भावयन्ति~। }\pend
\pstart\versepar \skt{पञ्चेन्द्रियं }\edtext{\skt{पञ्चाधा}}{\lemma{\skt{पञ्चाधा}}\Afootnote{\skt{पञ्चसु}\ A; \skt{पञ्चाध}\ H}}\skt{ भावयित्वा }\pend
\pstart\versepar \skt{पंचत्वमायान्ति विनाशकाले~॥}\pend

\pstart \ledsidenote{\passref{3.9.7}}\edtext{\skt{अधोगमास्तु}}{\lemma{\skt{अधोगमास्तु}}\Afootnote{A; \skt{अ व्या गमास्तु}\ H\textsuperscript{ac}; \skt{अ धो गमातास्तु}\ H\textsuperscript{pc}}}\skt{ }\edtext{\skt{वात\-मू\-त्र\-पु\-री\-ष\-शु\-क्रा\-र्त्त\-वा\-न\-भि\-व\-हन्ति~। तास्तु}}{\lemma{\skt{वात\appbrk{}मू\appbrk{}त्र\appbrk{}पु\appbrk{}री\appbrk{}ष\appbrk{}शु\appbrk{}क्रा\appbrk{}र्त्त\appbrk{}वा\appbrk{}न\appbrk{}भि\appbrk{}व\appbrk{}हन्ति~। तास्तु}}\Afootnote{H; \skt{वात\appbrk{}मू\appbrk{}त्र\appbrk{}पु\appbrk{}री\appbrk{}ष\appbrk{}शु\appbrk{}क्रा\appbrk{}र्त\appbrk{}वा\appbrk{}दी\appbrk{}न्यधो वहन्ति~। तास्तु}\ A}}\skt{ }\edtext{\edtext{\edtext{\edtext{\edtext{\skt{पित्ता\-श\-य\-म\-भि\-प्र\-प\-न्ना\-स्त\-त्र\-स्थ\-म\-न्न\-पा\-न\-र\-स\-म\-भि\-प\-क्व\-मौ\-ष्णा\-द्वि\-वे\-च\-य\-न्त्योऽ\-भि\-व\-ह\-न्त्य\-स्त\-र्प\-यन्ति}}{\lemma{\skt{तत्रस्थम्}}\Afootnote{A; \skt{तत्रस्थाम्}\ H}}}{\lemma{\skt{अन्न\appbrk{}पा\appbrk{}न\appbrk{}र\appbrk{}सम्}}\Afootnote{H; \skt{एवा\appbrk{}न्न\appbrk{}पा\appbrk{}न\appbrk{}रसं}\ A}}}{\lemma{\skt{अभिपक्वम्}}\Afootnote{H; \skt{विपक्वम्}\ A}}}{\lemma{\skt{औष्णाद्}}\Afootnote{H; \skt{औष्ण्याद्}\ A}}}{\lemma{\skt{तर्पयन्ति}}\Afootnote{H; \skt{शरीरं तर्पयन्ति, अर्पयन्ति}\ A}}\skt{ }\edtext{\edtext{\skt{ऊर्द्ध्व\-गा\-ना\-न्ति\-र्य\-ग्गा\-नाञ्च}}{\lemma{\skt{ऊर्द्ध्वगानान्}}\Afootnote{H; \skt{चोर्ध्वगानां}\ A}}}{\lemma{\skt{तिर्यग्गानाञ्च}}\Afootnote{H; \skt{तिर्यग्गाणां च}\ A}}\skt{ }\edtext{\skt{रस\-स्था\-ना\-ञ्चा\-भि\-पू\-र\-यन्ति}}{\lemma{\skt{रसस्थानाञ्}}\Afootnote{H; \skt{रसस्थानं}\ A}}\skt{~। }\edtext{\skt{तत्र}}{\lemma{\skt{तत्र}}\Afootnote{H; \emph{om.}\ A}}\skt{ मूत्र\-पु\-री\-ष\-स्वे\-दांश्च विवे\-च\-यं\-त्या\-म\-प\-क्वा\-श\-या\-न्तरे }\edtext{\skt{त्रिधा}}{\lemma{\skt{त्रिधा}}\Afootnote{H; \skt{च}\ A}}\skt{ त्रिधा }\edtext{\skt{प्रजायन्ते }\edtext{\skt{तांस्त्रिशत्}}{\lemma{\skt{त्रिशत्}}\Afootnote{H; \skt{स्त्रिंशत्}\ A}}}{\lemma{\skt{प्रजायन्ते तांस्}}\Afootnote{H; \skt{जायन्ते, ता}\ A}}\skt{~। तासान्तु वात\-पि\-त्त\-क\-फ\-शो\-णि\-त\-र\-सान्द्वे द्वे }\edtext{\edtext{\skt{वहतस्ता}}{\lemma{\skt{वहतस्}}\Afootnote{H; \skt{वहत}\ A}}}{\lemma{\skt{ता}}\Afootnote{H; \skt{स्ता}\ A}}\skt{ दश~। द्वे अन्नवाहिन्यौ~। अन्त्राश्रिते तोयवहे }\edtext{\skt{द्वे}}{\lemma{\skt{द्वे}}\Afootnote{H; \emph{om.}\ A}}\skt{ द्वे मूत्र\-ब\-स्ति\-म\-भि\-प्र\-पन्ने~। }\edtext{\skt{शुक्र\-प्रा\-दु\-र्ब्भा\-वाय च}}{\lemma{\skt{शुक्र\appbrk{}प्रा\appbrk{}दु\appbrk{}र्ब्भा\appbrk{}वाय च}}\Afootnote{H; \skt{मूत्रवहे द्वे, शुक्रवहे द्वे शुक्र\appbrk{}प्रा\appbrk{}दु\appbrk{}र्भा\appbrk{}वाय}\ A}}\skt{ द्वे~। }\edtext{\skt{एवं}}{\lemma{\skt{एवं}}\Afootnote{H; \skt{विसर्गाय, ते एव}\ A}}\skt{ }\edtext{\skt{रक्त\-म\-भि\-व\-हते}}{\lemma{\skt{अभिवहते}}\Afootnote{H; \skt{अभिवहतो विसृजतश्च}\ A}}\skt{~। }\edtext{\skt{नारी\-णा\-मा\-र्त्त\-व\-संज्ञे}}{\lemma{\skt{आर्त्तवसंज्ञे}}\Afootnote{H; \skt{आर्तवसंज्ञं}\ A}}\skt{ द्वे~। }\edtext{\skt{वर्च्चोनिःसरतः}}{\lemma{\skt{वर्च्चोनिःसरतः}}\Afootnote{H; \skt{वर्चो निरसन्यौ}\ A}}\skt{ स्थूला\-न्त्र\-प्र\-ति\-बद्धे }\edtext{\skt{द्वे}}{\lemma{\skt{द्वे}}\Afootnote{H; \emph{om.}\ A}}\skt{~। }\edtext{\edtext{\edtext{\skt{अष्टा\-व\-न्य\-स्ति\-र्य\-ग्गास्ताः}}{\lemma{\skt{अन्यस्}}\Afootnote{H; \skt{अन्यास्}\ A}}}{\lemma{\skt{तिर्यग्गास्}}\Afootnote{H; \skt{तिर्यग्गा}\ A}}}{\lemma{\skt{ताः}}\Afootnote{H; \skt{मिनीनां धमनीनां}\ A}}\skt{ स्वेदमर्प्पयन्ति~। }\edtext{\skt{तास्त्वे\-ता\-स्त्रिं\-श\-त्सं\-वि\-भागा}}{\lemma{\skt{संविभागा}}\Afootnote{H; \skt{सविभागा}\ A}}\skt{ }\edtext{\skt{व्याख्याताः}}{\lemma{\skt{व्याख्याताः}}\Afootnote{H; \skt{व्याख्याताः}\ A}}\skt{~। }\edtext{\edtext{\skt{ताभिरधो}}{\lemma{\skt{ताभिर्}}\Afootnote{H; \skt{एताभिर्}\ A}}}{\lemma{\skt{अधो}}\Afootnote{H; \skt{अधो}\ A}}\skt{ नाभेः }\edtext{\skt{पक्वा\-श\-यै\-क\-टी\-मू\-त्र\-ब\-स्ति\-गु\-द\-मे\-ढ्र\-स\-क्थीनि}}{\lemma{\skt{पक्वा\appbrk{}श\appbrk{}यै\appbrk{}क\appbrk{}टी\appbrk{}मू\appbrk{}त्र\appbrk{}ब\appbrk{}स्ति\appbrk{}गु\appbrk{}द\appbrk{}मे\appbrk{}ढ्र\appbrk{}स\appbrk{}क्थीनि}}\Afootnote{\skt{पक्वा\appbrk{}श\appbrk{}य\appbrk{}क\appbrk{}टी\appbrk{}मू\appbrk{}त्र\appbrk{}पु\appbrk{}री\appbrk{}ष\appbrk{}गु\appbrk{}द\appbrk{}ब\appbrk{}स्ति\appbrk{}मे\appbrk{}ढ्र\appbrk{}स\appbrk{}क्थीनि}\ A; \skt{पक्वा\appbrk{}स\appbrk{}य\appbrk{}क\appbrk{}टी\appbrk{}मू\appbrk{}त्र\appbrk{}व\appbrk{}स्ति\appbrk{}गु\appbrk{}द\appbrk{}मे\appbrk{}ढ्र\appbrk{}स\appbrk{}क्थीनि}\ H}}\skt{ }\edtext{\skt{धार्यन्ते}}{\lemma{\skt{धार्यन्ते}}\Afootnote{H; \skt{धार्यन्ते याप्यन्ते च}\ A}}\skt{~।}\pend

\pstart\versepar \ledsidenote{\passref{3.9.8}}\skt{भवति चात्र~॥ }\pend
\pstart\versepar \skt{}\edtext{\skt{अधोगमास्तु}}{\lemma{\skt{अधोगमास्तु}}\Afootnote{A; \skt{अधोगतास्तु}\ H\textsuperscript{ac}; \skt{अधोगतामास्तु}\ H\textsuperscript{pc}}}\skt{ कुर्व्वन्ति कर्म्माण्येतानि सर्व्वशः~। }\pend
\pstart\versepar \skt{}\edtext{\skt{तिर्यग्गास्त्वेव}}{\lemma{\skt{त्वेव}}\Afootnote{H; \skt{संप्र}\ A}}\skt{ वक्ष्यामि कर्म्म चासां }\edtext{\skt{यथातथं}}{\lemma{\skt{यथातथं}}\Afootnote{H; \skt{यथायथम्}\ A}}\skt{~॥}\pend

\pstart \ledsidenote{\passref{3.9.9}}\edtext{\skt{तिर्यग्गानां तु}}{\lemma{\skt{तिर्यग्गानां तु}}\Afootnote{H; \skt{तिर्यग्गाणां तु}\ A}}\skt{ }\edtext{\skt{चत\-सृ\-णा\-मे\-कैका}}{\lemma{\skt{एकैका}}\Afootnote{H; \skt{धम\appbrk{}नी\appbrk{}ना\appbrk{}मे\appbrk{}कैका}\ A}}\skt{ शतधा सहस्रधा }\edtext{\skt{च}}{\lemma{\skt{च}}\Afootnote{H; \skt{चोत्तरोत्तरं}\ A}}\skt{ }\edtext{\skt{भिद्यन्ते~। }\edtext{\skt{तास्त्वसंख्येया}}{\lemma{\skt{त्वसंख्येया}}\Afootnote{H; \skt{स्त्वसङ्ख्येयाः}\ A}}}{\lemma{\skt{भिद्यन्ते~। तास्}}\Afootnote{H; \skt{विभज्यन्ते, ता}\ A}}\skt{~। ताभिरिदं शरीरं गवाक्षितं विब\-द्ध\-मा\-त\-तञ्च~। }\edtext{\skt{तासाञ्च}}{\lemma{\skt{तासाञ्च}}\Afootnote{H; \skt{तासां}\ A}}\skt{ मुखानि रोम\-कू\-प\-प्र\-ति\-ब\-द्धानि यैः }\edtext{\skt{स्वेद}}{\lemma{\skt{स्वेद}}\Afootnote{H; \skt{स्वेदम्}\ A}}\skt{ }\edtext{\skt{प्रस्रवति}}{\lemma{\skt{प्रस्रवति}}\Afootnote{H; \skt{अभिवहन्ति}\ A}}\skt{ }\edtext{\edtext{\skt{रस\-श्चा\-भि\-स\-न्त\-र्प्प\-यति}}{\lemma{\skt{रसश्}}\Afootnote{H; \skt{रसं}\ A}}}{\lemma{\skt{चाभि\appbrk{}स\appbrk{}न्त\appbrk{}र्प्प\appbrk{}यति}}\Afootnote{H; \skt{चाभितर्पयन्त्य्}\ A}}\skt{ अन्तर्ब्बहिश्च~। }\edtext{\skt{ताभिरेव}}{\lemma{\skt{ताभिर्}}\Afootnote{H; \skt{तैर्}\ A}}\skt{ }\edtext{\edtext{\skt{चाभ्य\-ङ्ग\-प\-रि\-षे\-का\-व\-गा\-हा\-व\-ले\-प\-न\-वी\-र्या\-ण्य\-न्तः\-श\-री\-र\-म\-भि\-प्र\-प\-द्यन्ते}}{\lemma{\skt{चाभ्य\appbrk{}ङ्ग\appbrk{}प\appbrk{}रि\appbrk{}षे\appbrk{}का\appbrk{}व\appbrk{}गा\appbrk{}हा\appbrk{}व\appbrk{}ले\appbrk{}प\appbrk{}न\appbrk{}वी\appbrk{}र्याण्य्}}\Afootnote{\skt{चाभ्य\appbrk{}ङ्ग\appbrk{}प\appbrk{}रि\appbrk{}षे\appbrk{}का\appbrk{}व\appbrk{}गा\appbrk{}हा\appbrk{}ले\appbrk{}प\appbrk{}न\appbrk{}वी\appbrk{}र्याण्य्}\ A; \skt{चाभ्यं\appbrk{}ग\appbrk{}प\appbrk{}रि\appbrk{}शे\appbrk{}का\appbrk{}व\appbrk{}गा\appbrk{}हा\appbrk{}व\appbrk{}ले\appbrk{}प\appbrk{}न\appbrk{}वी\appbrk{}र्याण्य्}\ H}}}{\lemma{\skt{अभिप्रपद्यन्ते}}\Afootnote{H; \skt{अभि\appbrk{}प्र\appbrk{}ति\appbrk{}प\appbrk{}द्यन्ते}\ A}}\skt{ त्वचि विपक्वानि~। }\edtext{\skt{ताभिरेव}}{\lemma{\skt{ताभिर्}}\Afootnote{H; \skt{तैर्}\ A}}\skt{ च स्पर्शं }\edtext{\edtext{\skt{शुभमशुभम्वा}}{\lemma{\skt{शुभम्}}\Afootnote{H; \skt{सुखम्}\ A}}}{\lemma{\skt{अशुभम्वा}}\Afootnote{H; \skt{असुखं वा}\ A}}\skt{ }\edtext{\skt{गृह्णीते}}{\lemma{\skt{गृह्णीते}}\Afootnote{\skt{गृह्णीते, तास्}\ A; \skt{गृ}\ H}}\skt{~। }\edtext{\skt{ताश्चतस्रो}}{\lemma{\skt{ताश्}}\Afootnote{\skt{त्वेताश्}\ A; \skt{हीते ताश्}\ H}}\skt{ धमन्यः सर्व्वाङ्गगताः सविभागा व्याख्याताः~॥}\pend

\pstart\versepar \ledsidenote{\passref{3.9.9.add}}\skt{भवतश्चात्र श्लोकौ~॥ }\pend
\pstart\versepar \skt{धमन्यः सर्व्व }\edtext{\skt{एवैताः}}{\lemma{\skt{एवैताः}}\Afootnote{\skt{एवास्ताः}\ H}}\skt{ प्रतिबद्धा यथेन्द्रियं~। }\pend
\pstart\versepar \skt{जाग्रतां स्वपताम्वापि }\edtext{\skt{वह\-न्ती\-न्द्रि\-य\-गो\-चरे}}{\lemma{\skt{वह\appbrk{}न्ती\appbrk{}न्द्रि\appbrk{}य\appbrk{}गो\appbrk{}चरे}}\Afootnote{\skt{वह\appbrk{}ती\appbrk{}न्द्रि\appbrk{}य\appbrk{}गो\appbrk{}चरे~॥}\ H}}\skt{~॥}\footnoteA{This verse is unique to the Nepalese version and is supported in witness H. It is hard to parse, partly because of the plurals in the first pāda and the singular in the second. Some transmission problem is signalled by the need for the emendation of H's \emph{evās tāḥ} and \emph{vahatī}.}\pend

\pstart\versepar \ledsidenote{\passref{3.9.10}}\skt{यथा स्वभावतः खानि मृणालेषु बिसेषु च~। }\pend
\pstart\versepar \skt{}\edtext{\skt{तथैव}}{\lemma{\skt{तथैव}}\Afootnote{H; \skt{धम}\ A}}\skt{ }\edtext{\skt{तासां}}{\lemma{\skt{तासां}}\Afootnote{H; \skt{नीनां}\ A}}\skt{ }\edtext{\skt{द्वाराणि}}{\lemma{\skt{द्वाराणि}}\Afootnote{H; \skt{तथा खानि}\ A}}\skt{ रसो यैरुपचीयते~॥}\pend

\pstart \ledsidenote{\passref{3.9.12}}\skt{अत ऊर्द्ध्वं स्रोतसां }\edtext{\skt{मूलं}}{\lemma{\skt{मूलं}}\Afootnote{H; \skt{मूल}\ A}}\skt{ विद्धलक्षणं }\edtext{\skt{व्याख्यास्यामः}}{\lemma{\skt{व्याख्यास्यामः}}\Afootnote{H; \skt{उपदेक्ष्यामः}\ A}}\skt{~। तानि तु प्राणा\-न्नो\-द\-क\-र\-स\-र\-क्त\-मां\-स\-मे\-दो\-मू\-त्र\-पु\-री\-ष\-शु\-क्रा\-र्त्त\-व\-व\-हानि }\edtext{\skt{येष्वधीकारः}}{\lemma{\skt{अधीकारः}}\Afootnote{H; \skt{अधिकारः}\ A}}\skt{~। }\edtext{\skt{एषां}}{\lemma{\skt{एषां}}\Afootnote{H; \skt{एकेषां बहूनि; एतेषां}\ A}}\skt{ विशेषा बहवः~। }\edtext{\skt{तेषां}}{\lemma{\skt{तेषां}}\Afootnote{H; \skt{तत्र}\ A}}\skt{ }\edtext{\skt{प्राणवहे}}{\lemma{\skt{प्राणवहे}}\Afootnote{A; \emph{om.}\ H}}\skt{ }\edtext{\skt{द्वे,}}{\lemma{\skt{द्वे}}\Afootnote{A; \emph{om.}\ H}}\skt{ }\edtext{\skt{तयोर्मूलं}}{\lemma{\skt{तयोर्}}\Afootnote{A; \emph{om.}\ H}}\footnoteA{I insert \emph{ prāṇavahe dve, tayor} from the vulgate here, to make sense of the sentence and the corrollary in the next sentence.}\skt{ हृदयं }\edtext{\skt{रसवहिन्यश्च}}{\lemma{\skt{रसवहिन्यश्च}}\Afootnote{H; \skt{रसवाहिन्यश्च}\ A}}\skt{ धमन्यस्~। तत्र }\edtext{\skt{विद्धस्य}}{\lemma{\skt{विद्धस्य}}\Afootnote{H; \skt{विद्धस्या}\ A}}\skt{ }\edtext{\skt{क्रोश\-न\-वि\-न\-म\-न\-भ्र\-म\-णानि संज्ञानाशो}}{\lemma{\skt{क्रोश\appbrk{}न\appbrk{}वि\appbrk{}न\appbrk{}म\appbrk{}न\appbrk{}भ्र\appbrk{}म\appbrk{}णानि संज्ञानाशो}}\Afootnote{\skt{क्रोश\appbrk{}न\appbrk{}वि\appbrk{}न\appbrk{}म\appbrk{}न\appbrk{}मो\appbrk{}ह\appbrk{}न\appbrk{}भ्र\appbrk{}म\appbrk{}ण\appbrk{}वे\appbrk{}प\appbrk{}नानि मरणं वा भवति; अन्नवहे द्वे, तयो\appbrk{}र्मू\appbrk{}ल\appbrk{}मा\appbrk{}मा\appbrk{}श\appbrk{}योऽ\appbrk{}न्न\appbrk{}वा\appbrk{}हि\appbrk{}न्यश्च धमन्यः, तत्र विद्धस्याध्मानं शूलोऽ\appbrk{}न्न\appbrk{}द्वे\appbrk{}ष\appbrk{}श्छर्दिः पिपासान्ध्यं}\ A; \skt{क्रोश\appbrk{}न\appbrk{}वि\appbrk{}न\appbrk{}म\appbrk{}न\appbrk{}भ्र\appbrk{}म\appbrk{}णानि संज्ञानाशो मरणञ्च~। अन्नवहे द्वे तयो\appbrk{}र्म्मू\appbrk{}ल\appbrk{}मा\appbrk{}मा\appbrk{}सयो न्न}\ H}}\skt{ }\edtext{\skt{मरणञ्च}}{\lemma{\skt{मरणञ्च}}\Afootnote{A; \emph{om.}\ H}}\skt{~। }\edtext{\skt{उदकवहे}}{\lemma{\skt{उदकवहे}}\Afootnote{A; \emph{om.}\ H}}\skt{ }\edtext{\skt{द्वे}}{\lemma{\skt{द्वे}}\Afootnote{A; \emph{om.}\ H}}\skt{ }\edtext{\edtext{\edtext{\skt{तयोर्म्मूलम्द्वयं}}{\lemma{\skt{तयोर्}}\Afootnote{A; \emph{om.}\ H}}}{\lemma{\skt{म्मूलम्}}\Afootnote{A; \emph{om.}\ H}}\skt{ }\edtext{\skt{आमा\-श\-योऽ\-न्न\-वा\-हि\-न्यश्च}}{\lemma{\skt{ऽन्नवाहिन्यश्च}}\Afootnote{\skt{रसवाहिन्यश्च}\ A; \skt{वाहिन्यश्च}\ H}}}{\lemma{\skt{द्वयं आमाशयो}}\Afootnote{\skt{तालु क्लोम च, तत्र विद्धस्य पिपासा सद्योमरणं च; रसवहे द्वे, तयोर्मूलं हृदयं}\ A; \emph{om.}\ H}}\skt{ धमन्यः~। तत्र विद्धस्य }\edtext{\skt{प्राण\-व\-ह\-वि\-द्ध\-व\-त्म\-रणं}}{\lemma{\skt{प्राण\appbrk{}व\appbrk{}ह\appbrk{}वि\appbrk{}द्ध\appbrk{}वत्}}\Afootnote{H; \skt{शोषः\appbrk{}प्रा\appbrk{}ण\appbrk{}व\appbrk{}ह\appbrk{}वि\appbrk{}द्ध\appbrk{}वच्च}\ A}}\skt{ }\edtext{\skt{तल्लिङ्गश्च}}{\lemma{\skt{तल्लिङ्गश्च}}\Afootnote{H; \skt{तल्लिङ्गानि च}\ A}}\skt{~। रक्तवहे द्वे तयोर्मूलं यकृत्प्लीहानौ~। }\edtext{\skt{तत्र}}{\lemma{\skt{तत्र}}\Afootnote{H; \skt{रक्तवाहिन्यश्च धमन्यः , तत्र}\ A}}\skt{ विद्धस्य श्यावाङ्गता ज्वरः }\edtext{\skt{पाण्डुता}}{\lemma{\skt{पाण्डुता}}\Afootnote{H; \skt{दाहः पाण्डुता}\ A}}\skt{ }\edtext{\skt{दाहः}}{\lemma{\skt{दाहः}}\Afootnote{H; \emph{om.}\ A}}\skt{ }\edtext{\edtext{\skt{शोणि\-ता\-भि\-ग\-म\-नञ्च}}{\lemma{\skt{शोणि\appbrk{}ता\appbrk{}भि\appbrk{}ग\appbrk{}म\appbrk{}नञ्}}\Afootnote{H; \skt{शोणितागमनं}\ A}}}{\lemma{\skt{च}}\Afootnote{H; \skt{रक्तनेत्रता च}\ A}}\skt{~। मांसवहे द्वे तयोर्म्मूलं }\edtext{\skt{स्नायु\-स्त्व\-ग्र\-क्त\-व\-हाश्च}}{\lemma{\skt{स्नायु\appbrk{}स्त्व\appbrk{}ग्र\appbrk{}क्त\appbrk{}व\appbrk{}हाश्च}}\Afootnote{H; \skt{स्नायुत्वचं रक्तवहाश्च}\ A}}\skt{ }\edtext{\skt{नाड्यस्}}{\lemma{\skt{नाड्यस्}}\Afootnote{H; \skt{धमन्यः}\ A}}\skt{~। तत्र विद्धस्य श्वय\-थु\-र्म्मां\-स\-शोषः }\edtext{\skt{सिराग्रन्थयो}}{\lemma{\skt{सिराग्रन्थयो}}\Afootnote{A; \skt{शिराग्रन्थयोर्}\ H}}\skt{ }\edtext{\skt{म्मरणञ्च}}{\lemma{\skt{म्मरणञ्च}}\Afootnote{\skt{मरणं च}\ A H}}\skt{~। मेदोवहे द्वे तयोर्म्मूलं कटीवृक्कौ च~। तत्र विद्धस्य स्वेदागमनं }\edtext{\edtext{\edtext{\skt{स्निग्धा\-ङ्ग\-न्ता\-लु\-शो\-षश्च}}{\lemma{\skt{स्निग्धाङ्गन्}}\Afootnote{H; \skt{स्निग्धाङ्गता}\ A}}}{\lemma{\skt{तालुशोषश्}}\Afootnote{A; \skt{ता}\ H}}}{\lemma{\skt{च}}\Afootnote{\skt{स्थूलशोफता पिपासा च}\ A; \skt{लुशोषश्च}\ H}}\skt{~। मूत्रवहे द्वे तयोर्म्मूलं }\edtext{\skt{बस्तिमेढ्रञ्च}}{\lemma{\skt{बस्तिमेढ्रञ्च}}\Afootnote{H; \skt{बस्तिर्मेढ्रं च}\ A}}\skt{~। तत्र विद्ध\-स्या\-न\-द्ध\-ब\-स्तिता }\edtext{\skt{मूत्र\-नि\-रो\-ध\-स्त\-ब्ध\-मे\-ढ्रता च}}{\lemma{\skt{मूत्र\appbrk{}नि\appbrk{}रो\appbrk{}ध\appbrk{}स्त\appbrk{}ब्ध\appbrk{}मे\appbrk{}ढ्रता च}}\Afootnote{H; \skt{मूत्रनिरोधः स्तब्धमेढ्रता च}\ A}}\skt{~। पुरीषवहे द्वे तयो\-र्म्मू\-ल\-म्प\-क्वा\-शयो }\edtext{\skt{गुदश्च}}{\lemma{\skt{गुदश्च}}\Afootnote{H; \skt{गुदं च}\ A}}\skt{~। तत्र विद्धस्यानाहो दुर्ग्गन्धता ग्रथितान्त्रता च~। शुक्रवहे द्वे तयोर्म्मूलं स्तनौ वृषणौ च~। तत्र विद्धस्य }\edtext{\skt{क्लैब्यं}}{\lemma{\skt{क्लैब्यं}}\Afootnote{H; \skt{क्लीबता}\ A}}\skt{ }\edtext{\skt{चिराद्वा}}{\lemma{\skt{चिराद्वा}}\Afootnote{H; \skt{चिरात्}\ A}}\skt{ }\edtext{\skt{प्रसेकः}}{\lemma{\skt{प्रसेकः}}\Afootnote{H; \emph{om.}\ A}}\skt{ }\edtext{\skt{प्रसेके}}{\lemma{\skt{प्रसेके}}\Afootnote{H; \skt{प्रसेको}\ A}}\skt{ }\edtext{\skt{चाल्प\-शु\-क्र\-द\-र्शनं}}{\lemma{\skt{चाल्प\appbrk{}शु\appbrk{}क्र\appbrk{}द\appbrk{}र्शनं}}\Afootnote{H; \skt{रक्तशुक्रता च ; आर्तवहे द्वे, तयोर्मूलं गर्भाशय आर्त\appbrk{}व\appbrk{}वा\appbrk{}हि\appbrk{}न्यश्च धमन्यः, तत्र विद्धाया वन्ध्यात्वं मैथु\appbrk{}ना\appbrk{}स\appbrk{}हि\appbrk{}ष्णु\appbrk{}त्व\appbrk{}मा\appbrk{}र्त\appbrk{}व\appbrk{}ना\appbrk{}शश्च}\ A}}\skt{ }\edtext{\skt{सेव\-नी\-च्छे\-दा\-द्रु\-जा\-प्रा\-दु\-र्ब्भावो}}{\lemma{\skt{सेव\appbrk{}नी\appbrk{}च्छे\appbrk{}दा\appbrk{}द्रु\appbrk{}जा\appbrk{}प्रा\appbrk{}दु\appbrk{}र्ब्भावो}}\Afootnote{\skt{सेव\appbrk{}नी\appbrk{}च्छे\appbrk{}दा\appbrk{}द्रु\appbrk{}जा\appbrk{}प्रा\appbrk{}दु\appbrk{}र्भावः}\ A; \skt{सेवनीच्छेदा प्रादुर्ब्भावो}\ H}}\skt{ }\edtext{\edtext{\skt{मरणञ्च}}{\lemma{\skt{मरणञ्}}\Afootnote{\skt{बस्ति\appbrk{}गु\appbrk{}द\appbrk{}वि\appbrk{}द्ध\appbrk{}ल\appbrk{}क्षणं}\ A; \skt{म}\ H}}}{\lemma{\skt{च}}\Afootnote{\skt{प्रागुक्तमिति}\ A; \skt{रणञ्च}\ H}}\skt{~।}\footnoteA{We emend witness H's \emph{sevanīcheddā prādurbhāvo} to the vulgate's reading. Perhaps there was an eye-skip from \emph{[chedā]-ā} to \emph{[ruj]-ā}.}\skt{ स्रोतोविद्धन्तु }\edtext{\skt{प्रत्या\-ख्या\-यो\-प\-क्र\-मेत}}{\lemma{\skt{प्रत्या\appbrk{}ख्या\appbrk{}यो\appbrk{}प\appbrk{}क्र\appbrk{}मेत}}\Afootnote{H; \skt{प्रत्या\appbrk{}ख्या\appbrk{}यो\appbrk{}प\appbrk{}च\appbrk{}रेत्, उ}\ A}}\skt{~। }\edtext{\skt{उद्धृतशल्यं तु}}{\lemma{\skt{उद्धृतशल्यं तु}}\Afootnote{H; \skt{द्धृतशल्यं तु}\ A}}\skt{ }\edtext{\skt{क्षत\-वि\-धा\-ने\-नो\-प\-क्र\-मे\-दिति}}{\lemma{\skt{क्षत\appbrk{}वि\appbrk{}धा\appbrk{}ने\appbrk{}नो\appbrk{}प\appbrk{}क्र\appbrk{}मे\appbrk{}दिति}}\Afootnote{H; \skt{क्षत\appbrk{}वि\appbrk{}धा\appbrk{}ने\appbrk{}नो\appbrk{}प\appbrk{}च\appbrk{}रेत्}\ A}}\skt{~।}\pend

\pstart\versepar \ledsidenote{\passref{3.9.13.add1}}\skt{श्लोकौ~॥ }\pend
\pstart\versepar \skt{य एवं न प्रजानाति स्रोत\-सा\-म्मू\-ल\-नि\-श्चयं~। }\pend
\pstart\versepar \skt{मह\-द्भ\-य\-म\-वा\-प्नोति न स कर्म्मसु सिद्ध्यति~॥}\pend

\pstart\versepar \ledsidenote{\passref{3.9.13.add2}}\skt{यः सम्यगेतं जानीयात्स भवे\-द्रा\-ज\-स\-म्मतः~। }\pend
\pstart\versepar \skt{पूजार्हो भिषजां ह्येष इति धन्वन्तरेर्मतं~॥}\pend

\pstart \ledsidenote{\passref{3.9.trailer}}\skt{इति }\edtext{\skt{धम\-नी\-व्या\-क\-रणं}}{\lemma{\skt{धम\appbrk{}नी\appbrk{}व्या\appbrk{}क\appbrk{}रणं}}\Afootnote{\skt{सुश्रु\appbrk{}त\appbrk{}सं\appbrk{}हि\appbrk{}तायां शारीरस्थाने धम\appbrk{}नी\appbrk{}व्या\appbrk{}क\appbrk{}रण}\ A}}\skt{ }\edtext{\skt{शारी\-र\-न्न\-व\-मो\-ध्यायः}}{\lemma{\skt{शारी\appbrk{}र\appbrk{}न्न\appbrk{}व\appbrk{}मो\appbrk{}ध्यायः}}\Afootnote{\skt{शारीरं नाम नवमोऽध्यायः}\ A}}\skt{~॥}\pend

\endnumbering
\beginnumbering
\eledsection*{Adhyāya 10}
\pstart \ledsidenote{\passref{3.10.1}}\skt{॥ अथातो }\edtext{\skt{गर्ब्भि\-णी\-व्या\-क\-रणं}}{\lemma{\skt{गर्ब्भि\appbrk{}णी\appbrk{}व्या\appbrk{}क\appbrk{}रणं}}\Afootnote{H; \skt{गर्भि\appbrk{}णी\appbrk{}व्या\appbrk{}क\appbrk{}रणं}\ A}}\skt{ शारी\-र\-म्व्या\-ख्या\-स्यामः~॥}\pend

\pstart \ledsidenote{\passref{3.10.3}}\edtext{\skt{तत्र}}{\lemma{\skt{तत्र}}\Afootnote{H; \emph{om.}\ A}}\skt{ }\edtext{\edtext{\skt{गर्भि\-णी\-म्प्र\-थ\-म\-दि\-व\-सा\-त्प्र\-भृ\-त्यादि}}{\lemma{\skt{गर्भिणीम्}}\Afootnote{\skt{गर्भिणी}\ A; \skt{गर्ब्भिणीम्}\ H}}}{\lemma{\skt{प्रभृत्यादि}}\Afootnote{H; \skt{प्रभृति}\ A}}\skt{ }\edtext{\skt{सेन्नित्यं}}{\lemma{\skt{सेन्नित्यं}}\Afootnote{H; \skt{नित्यं}\ A}}\skt{ }\edtext{\skt{हृष्ट\-शु\-च्य\-लं\-कृ\-त\-वा\-ससा}}{\lemma{\skt{हृष्ट\appbrk{}शु\appbrk{}च्य\appbrk{}लं\appbrk{}कृ\appbrk{}त\appbrk{}वा\appbrk{}ससा}}\Afootnote{H; \skt{प्रहृष्टा शुच्यलङ्कृता शुक्लवसना}\ A}}\skt{ शान्ति\-म\-ङ्ग\-ल\-दे\-व\-ता\-ब्रा\-ह्म\-ण\-गु\-रु\-परा च }\edtext{\skt{भवितव्यं}}{\lemma{\skt{भवितव्यं}}\Afootnote{H; \skt{भवेत्}\ A}}\skt{~। }\edtext{\skt{मलि\-न\-वि\-कृ\-त\-ही\-न\-गा\-त्राणि च न}}{\lemma{\skt{मलि\appbrk{}न\appbrk{}वि\appbrk{}कृ\appbrk{}त\appbrk{}ही\appbrk{}न\appbrk{}गा\appbrk{}त्राणि च न}}\Afootnote{H; \skt{मलि\appbrk{}न\appbrk{}वि\appbrk{}कृ\appbrk{}त\appbrk{}ही\appbrk{}न\appbrk{}गा\appbrk{}त्राणि न}\ A}}\skt{ स्पृशेत्~। }\edtext{\skt{दुर्ग्ग\-न्ध\-दु\-र्द्द\-श\-नानि च}}{\lemma{\skt{दुर्ग्ग\appbrk{}न्ध\appbrk{}दु\appbrk{}र्द्द\appbrk{}श\appbrk{}नानि च}}\Afootnote{H; \skt{दुर्ग\appbrk{}न्ध\appbrk{}दु\appbrk{}र्द\appbrk{}र्श\appbrk{}नानि}\ A}}\skt{ }\edtext{\skt{दूरत}}{\lemma{\skt{दूरत}}\Afootnote{H; \emph{om.}\ A}}\skt{ }\edtext{\skt{एव}}{\lemma{\skt{एव}}\Afootnote{H; \emph{om.}\ A}}\skt{ }\edtext{\skt{परिहरेत्}}{\lemma{\skt{परिहरेत्}}\Afootnote{H; \emph{om.}\ A}}\skt{~। }\edtext{\edtext{\edtext{\skt{वहि\-न्न्निः\-कृ\-टा\-शू\-न्या\-गा\-र\-चै\-त्य\-श्म\-शा\-न\-चै\-त्य\-वृ\-क्ष\-से\-वाश्च}}{\lemma{\skt{वहि\appbrk{}न्न्निः\appbrk{}कृ\appbrk{}टा\appbrk{}शू\appbrk{}न्या\appbrk{}गा\appbrk{}र\appbrk{}चैत्य्}}\Afootnote{H; \emph{om.}\ A}}}{\lemma{\skt{अश्मशानचैत्य्}}\Afootnote{H; \emph{om.}\ A}}}{\lemma{\skt{अवृक्षसेवाश्च}}\Afootnote{H; \emph{om.}\ A}}\skt{ }\edtext{\skt{क्रोधा\-या\-स\-क\-रांश्च}}{\lemma{\skt{क्रोधा\appbrk{}या\appbrk{}स\appbrk{}क\appbrk{}रांश्च}}\Afootnote{H; \emph{om.}\ A}}\skt{ }\edtext{\edtext{\skt{भावा\-नु\-च्चै\-र्भा\-ष्या\-दिकं}}{\lemma{\skt{भावानुच्चैर्}}\Afootnote{H; \emph{om.}\ A}}\skt{ च}}{\lemma{\skt{भाष्यादिकं च}}\Afootnote{\emph{om.}\ A; \skt{ब्भाष्यादिकं च}\ H}}\skt{ परिहरेत्~। उद्वेजनीयाश्च कथाः शुष्कम्पर्युषितं कुथितं क्लिन्नं }\edtext{\skt{चान्न\-नो\-प\-यु\-ञ्जीत}}{\lemma{\skt{चान्न\appbrk{}नो\appbrk{}प\appbrk{}यु\appbrk{}ञ्जीत}}\Afootnote{H; \skt{चान्नं नोपभुञ्जीत}\ A}}\skt{ }\edtext{\skt{यानि च}}{\lemma{\skt{यानि च}}\Afootnote{H; \skt{बहिर्निष्क्रमणं शून्या\appbrk{}गा\appbrk{}र\appbrk{}चै\appbrk{}त्य\appbrk{}श्म\appbrk{}शा\appbrk{}न\appbrk{}वृ\appbrk{}क्षा\appbrk{}श्र\appbrk{}या\appbrk{}न्क्रो\appbrk{}ध\appbrk{}म\appbrk{}य\appbrk{}श\appbrk{}स्क\appbrk{}रांश्च भावा\appbrk{}नु\appbrk{}च्चै\appbrk{}र्भा\appbrk{}ष्या\appbrk{}दिकं च परिहरेद्यानि च}\ A}}\skt{ }\edtext{\edtext{\skt{गर्भ\-म्व्या\-पा\-द\-यन्ति}}{\lemma{\skt{गर्भम्}}\Afootnote{A; \skt{गर्ब्भम्}\ H}}\skt{~। न च}}{\lemma{\skt{व्यापादयन्ति~। न च}}\Afootnote{H; \skt{व्यापादयन्ति, न चा}\ A}}\skt{ }\edtext{\skt{तैला\-भ्य\-ङ्गो\-त्सा\-द\-ना\-दीनि}}{\lemma{\skt{तैला\appbrk{}भ्य\appbrk{}ङ्गो\appbrk{}त्सा\appbrk{}द\appbrk{}ना\appbrk{}दीनि}}\Afootnote{H; \skt{भीक्ष्णं तैला\appbrk{}भ्य\appbrk{}ङ्गो\appbrk{}त्सा\appbrk{}द\appbrk{}नादी}\ A}}\skt{ निसेवेत~। न चाया\-स\-ये\-च्छ\-रीरं पूर्व्वोक्तानि च परिहरेत्~। }\edtext{\skt{शयनासनञ्च}}{\lemma{\skt{शयनासनञ्च}}\Afootnote{H; \skt{शयनासनं मृ}\ A}}\skt{ }\edtext{\skt{समृद्धास्तरणं}}{\lemma{\skt{समृद्धास्}}\Afootnote{H; \skt{द्वास्}\ A}}\skt{ }\edtext{\edtext{\edtext{\skt{नात्यु\-च्छ्रि\-त\-म\-पा\-श्र\-यो\-पे\-त\-म\-सं\-वा\-ध\-म्वि\-द\-ध्यात्}}{\lemma{\skt{नात्युच्छ्रितम्}}\Afootnote{H; \skt{नात्युच्चम्}\ A}}}{\lemma{\skt{अपा\appbrk{}श्र\appbrk{}यो\appbrk{}पे\appbrk{}तम्}}\Afootnote{H; \skt{अपाश्रयोपेत}\ A}}}{\lemma{\skt{असंवाधम्}}\Afootnote{H; \skt{मसंबाधं च}\ A}}\skt{~। हृद्यं }\edtext{\edtext{\skt{द्रव\-म\-धु\-र\-प्रा\-य\-स्नि\-ग्ध\-न्दी\-प\-नीयं}}{\lemma{\skt{द्रव\appbrk{}म\appbrk{}धु\appbrk{}र\appbrk{}प्रा\appbrk{}य\appbrk{}स्नि\appbrk{}ग्धन्}}\Afootnote{H; \skt{द्रव\appbrk{}म\appbrk{}धु\appbrk{}र\appbrk{}प्रायं स्निग्धं}\ A}}}{\lemma{\skt{दीपनीयं}}\Afootnote{H; \skt{दीपनीय}\ A}}\skt{ }\edtext{\skt{संस्कृ\-त\-भो\-ज\-नञ्च}}{\lemma{\skt{संस्कृ\appbrk{}त\appbrk{}भो\appbrk{}ज\appbrk{}नञ्च}}\Afootnote{H; \skt{संस्कृतं च भोजनं}\ A}}\skt{ भोज\-ये\-त्सा\-मा\-न्य\-मे\-तदा }\edtext{\skt{प्रसवकालात्}}{\lemma{\skt{प्रसवकालात्}}\Afootnote{H; \skt{प्रसवात्}\ A}}\skt{~।}\pend

\pstart \ledsidenote{\passref{3.10.4}}\skt{विशेषतस्तु }\edtext{\skt{गर्ब्भिणी}}{\lemma{\skt{गर्ब्भिणी}}\Afootnote{H; \skt{गर्भिणी}\ A}}\skt{ }\edtext{\skt{प्रथ\-म\-द्वि\-ती\-य\-तृ\-ती\-येषु}}{\lemma{\skt{प्रथ\appbrk{}म\appbrk{}द्वि\appbrk{}ती\appbrk{}य\appbrk{}तृ\appbrk{}ती\appbrk{}येषु}}\Afootnote{H; \skt{प्रथ\appbrk{}म\appbrk{}द्वि\appbrk{}ती\appbrk{}य\appbrk{}तृ\appbrk{}तीय}\ A}}\skt{ मासेषु }\edtext{\skt{मधु\-र\-शी\-त\-द्र\-व\-प्रा\-य\-मा\-हा\-रमा}}{\lemma{\skt{मधु\appbrk{}र\appbrk{}शी\appbrk{}त\appbrk{}द्र\appbrk{}व\appbrk{}प्रा\appbrk{}य\appbrk{}मा\appbrk{}हा\appbrk{}रमा}}\Afootnote{H; \skt{मधु\appbrk{}र\appbrk{}शी\appbrk{}त\appbrk{}द्र\appbrk{}व\appbrk{}प्रा\appbrk{}य\appbrk{}मा\appbrk{}हा\appbrk{}रमु}\ A}}\skt{ }\edtext{\skt{आसेवेत}}{\lemma{\skt{आसेवेत}}\Afootnote{H; \skt{पसेवेत}\ A}}\skt{~। }\edtext{\skt{चतुर्थे}}{\lemma{\skt{चतुर्थे}}\Afootnote{H; \skt{विशेषतस्तु तृतीये षष्टिकौदनं पयसा भोजयेत् , चतुर्थे}\ A}}\skt{ }\edtext{\skt{मासे}}{\lemma{\skt{मासे}}\Afootnote{H; \skt{दध्ना, पञ्चमे पयसा, षष्ठे सर्पिषेत्येके; चतुर्थे}\ A}}\skt{ }\edtext{\skt{पयो}}{\lemma{\skt{पयो}}\Afootnote{H; \skt{पयो}\ A}}\skt{ }\edtext{\edtext{\edtext{\skt{नव\-नी\-त\-मा\-हा\-र\-मा\-ह\-रेत्}}{\lemma{\skt{नवनीतम्}}\Afootnote{H; \skt{नव\appbrk{}नी\appbrk{}त\appbrk{}सं\appbrk{}सृ\appbrk{}ष्ट\appbrk{}मा\appbrk{}हा\appbrk{}र\appbrk{}ये\appbrk{}ज्जा\appbrk{}ङ्ग\appbrk{}ल\appbrk{}मां\appbrk{}स\appbrk{}स\appbrk{}हितं}\ A}}}{\lemma{\skt{आहारम्}}\Afootnote{H; \skt{हृद्यमन्नं}\ A}}}{\lemma{\skt{आहरेत्}}\Afootnote{H; \skt{च भोजयेत्}\ A}}\skt{~। पञ्चमे क्षीरसर्प्पिः संसृष्टं~। षष्ठे }\edtext{\skt{स्वदंष्ट्रासिद्धस्य}}{\lemma{\skt{स्वदंष्ट्रासिद्धस्य}}\Afootnote{H; \skt{श्वदंष्ट्रासिद्धस्य}\ A}}\skt{ }\edtext{\skt{सर्प्पिषा}}{\lemma{\skt{सर्प्पिषा}}\Afootnote{H; \skt{सर्पि}\ A}}\skt{ }\edtext{\skt{पक्षमात्रं}}{\lemma{\skt{पक्षमात्रं}}\Afootnote{H; \skt{षो मात्रां}\ A}}\skt{ }\edtext{\skt{युक्त्या}}{\lemma{\skt{युक्त्या}}\Afootnote{H; \emph{om.}\ A}}\skt{ }\edtext{\skt{यवा\-गू\-म्पा\-य\-येत्}}{\lemma{\skt{यवागूम्}}\Afootnote{H; \emph{om.}\ A}}\skt{~। }\edtext{\skt{सप्तमे}}{\lemma{\skt{सप्तमे}}\Afootnote{H; \skt{यवागूं वा, सप्तमे}\ A}}\skt{ }\edtext{\edtext{\edtext{\skt{पृथ\-क्प\-र्ण्या\-दि\-सि\-द्ध\-मे\-व\-मा\-प्या\-य्यते}}{\lemma{\skt{पृथक्}}\Afootnote{H; \skt{सर्पिः पृथक्}\ A}}}{\lemma{\skt{एवम्}}\Afootnote{H; \skt{एव}\ A}}}{\lemma{\skt{आप्याय्यते}}\Afootnote{H; \skt{माप्यायते}\ A}}\skt{ }\edtext{\skt{गर्भः}}{\lemma{\skt{गर्भः}}\Afootnote{A; \skt{गर्ब्भः}\ H}}\skt{~। अष्टमे }\edtext{\skt{खल्वेनां}}{\lemma{\skt{खल्वेनां}}\Afootnote{H; \emph{om.}\ A}}\skt{ वदरोदकेन }\edtext{\edtext{\skt{वला\-ति\-व\-ला\-प\-यो\-द\-धि\-मस्तु}}{\lemma{\skt{वला\appbrk{}ति\appbrk{}व\appbrk{}ला\appbrk{}प\appbrk{}यो\appbrk{}द\appbrk{}धि\appbrk{}मस्}}\Afootnote{H; \skt{बला\appbrk{}ति\appbrk{}ब\appbrk{}ला\appbrk{}श\appbrk{}त\appbrk{}पु\appbrk{}ष्पा\appbrk{}प\appbrk{}ल\appbrk{}ल\appbrk{}प\appbrk{}यो\appbrk{}द\appbrk{}धिम}\ A}}}{\lemma{\skt{तु}}\Afootnote{H; \skt{स्तु}\ A}}\skt{ }\edtext{\edtext{\edtext{\skt{तैल\-म\-धु\-घृ\-त\-श\-त\-पु\-ष्पा\-मि\-श्रे\-णा\-स्था\-प\-ये\-त्पु\-रा\-ण\-पु\-री\-ष\-शु\-द्ध्य\-र्थ\-म\-नु\-लो\-मा\-र्थञ्च}}{\lemma{\skt{तैल\appbrk{}म\appbrk{}धु\appbrk{}घृ\appbrk{}त\appbrk{}श\appbrk{}त\appbrk{}पु\appbrk{}ष्पा\appbrk{}मि\appbrk{}श्रे\appbrk{}णा\appbrk{}स्था\appbrk{}प\appbrk{}येत्}}\Afootnote{H; \skt{तैल\appbrk{}ल\appbrk{}व\appbrk{}ण\appbrk{}म\appbrk{}द\appbrk{}न\appbrk{}फ\appbrk{}ल\appbrk{}म\appbrk{}धु\appbrk{}घृ\appbrk{}त\appbrk{}मि\appbrk{}श्रे\appbrk{}णा\appbrk{}स्था\appbrk{}प\appbrk{}येत्}\ A}}}{\lemma{\skt{पुरा\appbrk{}ण\appbrk{}पु\appbrk{}री\appbrk{}ष\appbrk{}शु\appbrk{}द्ध्य\appbrk{}र्थम्}}\Afootnote{H; \skt{पुरा\appbrk{}ण\appbrk{}पु\appbrk{}री\appbrk{}ष\appbrk{}शु\appbrk{}द्ध्यर्थ}\ A}}}{\lemma{\skt{अनुलोमार्थञ्च}}\Afootnote{H; \skt{मनुलोमनार्थं च}\ A}}\skt{ वायोस्ततः }\edtext{\skt{पयो}}{\lemma{\skt{पयो}}\Afootnote{H; \skt{पयो}\ A}}\skt{ }\edtext{\skt{मधुरसिद्धेन}}{\lemma{\skt{मधुरसिद्धेन}}\Afootnote{H; \skt{मधु\appbrk{}र\appbrk{}क\appbrk{}षा\appbrk{}य\appbrk{}सि\appbrk{}द्धेन}\ A}}\skt{ तैले\-ना\-नु\-वा\-स\-येत्~। अनुलोमे हि वायौ सुखम्प्रसूयते~। }\edtext{\skt{निरुपद्रवाश्च}}{\lemma{\skt{निरुपद्रवाश्च}}\Afootnote{H; \skt{निरुपद्रवा च}\ A}}\skt{ }\edtext{\skt{भवन्ति}}{\lemma{\skt{भवन्ति}}\Afootnote{H; \skt{भवति}\ A}}\skt{~। अत ऊर्द्ध्वं }\edtext{\edtext{\skt{स्निग्धा\-भि\-र्य\-वा\-गू\-भि\-र्ज्जा\-ङ्ग\-ल\-र\-सै\-श्चो\-प\-ह\-रेदा}}{\lemma{\skt{ज्जाङ्गलरसैश्}}\Afootnote{H; \skt{जाङ्गलरसै}\ A}}}{\lemma{\skt{चोपहरेदा}}\Afootnote{H; \skt{श्चोपक्रमेदा}\ A}}\skt{ प्रसवकालात्~। एव\-मु\-प\-क्रा\-न्ता\-स्नि\-ग्धा\-व\-ल\-व\-ती\-सु\-ख\-म\-नु\-प\-द्रवा प्रसूयते~।}\pend

\pstart \ledsidenote{\passref{3.10.5}}\skt{नवमे मासि }\edtext{\edtext{\skt{सूति\-का\-गा\-र\-मेनां}}{\lemma{\skt{सूतिकागारम्}}\Afootnote{H; \skt{सूतिकागार}\ A}}}{\lemma{\skt{एनां}}\Afootnote{H; \skt{मेनां}\ A}}\skt{ प्रवेशयेत्~। प्रशस्ते तिथ्यादौ~। }\edtext{\skt{तत्रेष्टागारं}}{\lemma{\skt{तत्रेष्टागारं}}\Afootnote{H; \skt{तत्रारिष्टं}\ A}}\skt{ ब्राह्म\-ण\-क्ष\-त्त्रि\-य\-वै\-श्य\-शू\-द्राणां }\edtext{\skt{श्वेत\-र\-क्त\-कृ\-ष्णेषु}}{\lemma{\skt{श्वेत\appbrk{}र\appbrk{}क्त\appbrk{}कृ\appbrk{}ष्णेषु}}\Afootnote{H; \skt{श्वेत\appbrk{}र\appbrk{}क्त\appbrk{}पी\appbrk{}त\appbrk{}कृ\appbrk{}ष्णेषु}\ A}}\skt{ भूमिप्रदेशेषु~। }\edtext{\skt{विल्व\-न्य\-ग्रो\-ध\-ति\-न्दु\-क\-नि\-र्म्मितं}}{\lemma{\skt{विल्व\appbrk{}न्य\appbrk{}ग्रो\appbrk{}ध\appbrk{}ति\appbrk{}न्दु\appbrk{}क\appbrk{}नि\appbrk{}र्म्मितं}}\Afootnote{H; \skt{बिल्व\appbrk{}न्य\appbrk{}ग्रो\appbrk{}ध\appbrk{}ति\appbrk{}न्दु\appbrk{}क\appbrk{}भ\appbrk{}ल्ला\appbrk{}त\appbrk{}क\appbrk{}नि\appbrk{}र्मित}\ A}}\skt{ }\edtext{\edtext{\skt{सर्व्वा\-गा\-र\-न्त\-त्म\-यञ्च}}{\lemma{\skt{तत्}}\Afootnote{H; \skt{यथासङ्ख्यं तन्}\ A}}}{\lemma{\skt{मयञ्च}}\Afootnote{H; \skt{मय}\ A}}\skt{ }\edtext{\skt{पर्यवनद्धं}}{\lemma{\skt{पर्यवनद्धं}}\Afootnote{H; \skt{पर्यङ्कं}\ A}}\skt{ }\edtext{\skt{समु\-प\-लि\-प्त\-भि\-त्ति\-प\-रि\-च्छदं}}{\lemma{\skt{समु\appbrk{}प\appbrk{}लि\appbrk{}प्त\appbrk{}भि\appbrk{}त्ति\appbrk{}प\appbrk{}रि\appbrk{}च्छदं}}\Afootnote{H; \skt{समु\appbrk{}प\appbrk{}लि\appbrk{}प्त\appbrk{}भित्तिं सुवि\appbrk{}भ\appbrk{}क्त\appbrk{}प\appbrk{}रि\appbrk{}च्छदं}\ A}}\skt{ }\edtext{\edtext{\skt{प्राग्द्वा\-र\-मु\-द\-ग्द्वा\-र\-म्वा\-ष्ट\-ह\-स्ता\-यत}}{\lemma{\skt{उदग्द्वारम्}}\Afootnote{H; \skt{दक्षिणद्वारं}\ A}}}{\lemma{\skt{वाष्टहस्तायत}}\Afootnote{H; \skt{वाष्टहस्तायतं}\ A}}\skt{ चतु\-र्ह\-स्त\-वि\-स्तृतं रक्षा\-म\-ङ्ग\-ल\-स\-म्प\-न्न\-म्वि\-धेयं~।}\pend

\pstart \ledsidenote{\passref{3.10.7}}\skt{तत्रो\-प\-स्थि\-त\-प्र\-स\-वायाः कटीपृष्ठं प्रति\-स\-म\-न्ता\-द्वे\-दना भवति~। }\edtext{\edtext{\skt{मूत्र\-म\-भी\-क्ष्ण\-म्प्र\-सि\-च्यते}}{\lemma{\skt{मूत्रमभीक्ष्णम्}}\Afootnote{H; \skt{अभीक्ष्णं}\ A}}}{\lemma{\skt{प्रसिच्यते}}\Afootnote{H; \skt{पुरी\appbrk{}ष\appbrk{}प्र\appbrk{}वृ\appbrk{}त्ति\appbrk{}र्मूत्रं प्रसिच्यते}\ A}}\skt{~। }\edtext{\skt{योनिमुखाच्च}}{\lemma{\skt{योनिमुखाच्च}}\Afootnote{H; \skt{योनिमुखाच्}\ A}}\skt{ }\edtext{\skt{श्लेष्मा}}{\lemma{\skt{श्लेष्मा}}\Afootnote{H; \skt{छ्लेष्मा च}\ A}}\pend

\pstart \ledsidenote{\passref{3.10.8}}\edtext{\skt{प्रजायिष्यमाणां}}{\lemma{\skt{प्रजायिष्यमाणां}}\Afootnote{H; \skt{प्रज\appbrk{}न\appbrk{}यि\appbrk{}ष्य\appbrk{}माणां}\ A}}\skt{ }\edtext{\edtext{\edtext{\skt{तुस्व\-भ्य\-क्ता\-मु\-ष्णो\-द\-क\-प\-रि\-षि\-क्त\-गा\-त्री\-ङ्कृ\-त\-म\-ङ्ग\-ला\-म\-थेनां}}{\lemma{\skt{तुस्वभ्यक्ताम्}}\Afootnote{H; \skt{कृत\appbrk{}म\appbrk{}ङ्ग\appbrk{}ल\appbrk{}स्व\appbrk{}स्ति\appbrk{}वा\appbrk{}चनां कुमा\appbrk{}र\appbrk{}प\appbrk{}रि\appbrk{}वृतां पुन्ना\appbrk{}म\appbrk{}फ\appbrk{}ल\appbrk{}हस्तां स्वभ्यक्ता}\ A}}}{\lemma{\skt{उष्णो\appbrk{}द\appbrk{}क\appbrk{}प\appbrk{}रि\appbrk{}षि\appbrk{}क्त\appbrk{}गा\appbrk{}त्रीङ्}}\Afootnote{H; \skt{मुष्णो\appbrk{}द\appbrk{}क\appbrk{}प\appbrk{}रि\appbrk{}षिक्}\ A}}\skt{ स}}{\lemma{\skt{कृत\appbrk{}म\appbrk{}ङ्ग\appbrk{}ला\appbrk{}म\appbrk{}थेनां स}}\Afootnote{H; \skt{तामथैनां सम्}\ A}}\skt{ }\edtext{\skt{घृतां}}{\lemma{\skt{घृतां}}\Afootnote{H; \skt{भृतां}\ A}}\skt{ }\edtext{\skt{यवागूं}}{\lemma{\skt{यवागूं}}\Afootnote{H; \skt{यवागू}\ A}}\skt{ }\edtext{\edtext{\skt{कण्ठात्पीतवती}}{\lemma{\skt{कण्ठात्}}\Afootnote{H; \skt{माकण्ठात्}\ A}}}{\lemma{\skt{पीतवती}}\Afootnote{H; \skt{पाययेत्}\ A}}\skt{ }\edtext{\skt{मृदुनि}}{\lemma{\skt{मृदुनि}}\Afootnote{H; \skt{ततः}\ A}}\skt{ }\edtext{\skt{कृतो\-प\-धा\-न\-वि\-स्तीर्ण्णे}}{\lemma{\skt{कृतो\appbrk{}प\appbrk{}धा\appbrk{}न\appbrk{}वि\appbrk{}स्तीर्ण्णे}}\Afootnote{H; \skt{कृतोपधाने मृदुनि विस्तीर्णे}\ A}}\skt{ शयने }\edtext{\edtext{\edtext{\skt{स्थिता\-मा\-भु\-ग्न\-स\-क्थि\-नी\-मु\-त्ता\-ना\-म\-श\-ङ्क\-नी\-या\-स्ति\-स्र\-श्च\-तस्रः}}{\lemma{\skt{स्थिता\appbrk{}मा\appbrk{}भु\appbrk{}ग्न\appbrk{}स\appbrk{}क्थि\appbrk{}नी\appbrk{}मु\appbrk{}त्ता\appbrk{}नाम्}}\Afootnote{\skt{स्थिता\appbrk{}मा\appbrk{}भु\appbrk{}ग्न\appbrk{}स\appbrk{}क्थी\appbrk{}मु\appbrk{}त्ताना}\ A; \skt{स्थिता\appbrk{}मा\appbrk{}भु\appbrk{}ग्न\appbrk{}स\appbrk{}क्थि\appbrk{}नी\appbrk{}मु\appbrk{}त्ता\appbrk{}नाम्}\ H}}}{\lemma{\skt{अशङ्कनीयास्}}\Afootnote{H; \skt{मशङ्कनीयाश्}\ A}}}{\lemma{\skt{तिस्रश्}}\Afootnote{H; \emph{om.}\ A}}\skt{ }\edtext{\skt{परि\-ण\-त\-व\-यसः}}{\lemma{\skt{परि\appbrk{}ण\appbrk{}त\appbrk{}व\appbrk{}यसः}}\Afootnote{H; \skt{स्त्रियः परि\appbrk{}ण\appbrk{}त\appbrk{}व\appbrk{}यसः}\ A}}\skt{~। }\edtext{\skt{व्रज\-न\-न\-कु\-शलाः}}{\lemma{\skt{व्रज\appbrk{}न\appbrk{}न\appbrk{}कु\appbrk{}शलाः}}\Afootnote{H; \skt{प्रज\appbrk{}न\appbrk{}न\appbrk{}कु\appbrk{}शलाः}\ A}}\skt{ }\edtext{\skt{कल्पितनख्यः}}{\lemma{\skt{कल्पितनख्यः}}\Afootnote{H; \skt{कर्तितनखाः}\ A}}\skt{ }\edtext{\skt{परि\-च\-रे\-यु\-रिति}}{\lemma{\skt{परि\appbrk{}च\appbrk{}रे\appbrk{}यु\appbrk{}रिति}}\Afootnote{H; \skt{परि\appbrk{}च\appbrk{}रे\appbrk{}यु\appbrk{}रिति}\ A}}\skt{~॥}\pend

\pstart \ledsidenote{\passref{3.10.9}}\skt{अथास्या }\edtext{\edtext{\edtext{\skt{विशि\-खा\-न्त\-र\-म\-नु\-लो\-मा\-या\-म\-नु\-सु\-ख\-म\-भ्य\-ञ्ज्यात्}}{\lemma{\skt{अनुलोमायाम्}}\Afootnote{H; \skt{अनुलोमम्}\ A}}}{\lemma{\skt{अनुसुखम्}}\Afootnote{H; \skt{अनुसुख}\ A}}}{\lemma{\skt{अभ्यञ्ज्यात्}}\Afootnote{H; \skt{मभ्यज्यानु}\ A}}\skt{~। }\edtext{\edtext{\skt{ब्रूयाच्चैनामेताः}}{\lemma{\skt{चैनाम्}}\Afootnote{H; \skt{चैताम्}\ A}}}{\lemma{\skt{एताः}}\Afootnote{H; \skt{एका}\ A}}\skt{ सुभगे }\edtext{\edtext{\skt{प्रवाहयश्चेति}}{\lemma{\skt{प्रवाहयश्}}\Afootnote{H; \skt{प्रवाहस्व्}\ A}}\skt{~। मा}}{\lemma{\skt{चेति~। मा}}\Afootnote{H; \skt{एति, न}\ A}}\skt{ }\edtext{\skt{चाप्राप्तावीः}}{\lemma{\skt{चाप्राप्तावीः}}\Afootnote{H; \skt{चाप्रा}\ A}}\skt{ }\edtext{\skt{प्रवाहिष्ठाः}}{\lemma{\skt{प्रवाहिष्ठाः}}\Afootnote{H; \emph{om.}\ A}}\skt{~। }\edtext{\skt{अनागता स्वाविषु}}{\lemma{\skt{अनागता स्वाविषु}}\Afootnote{H; \emph{om.}\ A}}\skt{ }\edtext{\skt{प्रवाहमाणाः}}{\lemma{\skt{प्रवाहमाणाः}}\Afootnote{H; \skt{प्}\ A}}\skt{~। }\edtext{\skt{बधिरं}}{\lemma{\skt{बधिरं}}\Afootnote{H; \emph{om.}\ A}}\skt{ }\edtext{\skt{मूकं}}{\lemma{\skt{मूकं}}\Afootnote{H; \emph{om.}\ A}}\skt{ }\edtext{\skt{स्रस्तहनुं}}{\lemma{\skt{स्रस्तहनुं}}\Afootnote{H; \emph{om.}\ A}}\skt{ }\edtext{\skt{मूर्द्धाभिघातिनं}}{\lemma{\skt{मूर्द्धाभिघातिनं}}\Afootnote{H; \emph{om.}\ A}}\skt{ }\edtext{\skt{कास\-श्वा\-स\-शो\-षो\-प\-द्रुतं}}{\lemma{\skt{कास\appbrk{}श्वा\appbrk{}स\appbrk{}शो\appbrk{}षो\appbrk{}प\appbrk{}द्रुतं}}\Afootnote{H; \emph{om.}\ A}}\skt{ }\edtext{\skt{कुव्जं}}{\lemma{\skt{कुव्जं}}\Afootnote{H; \emph{om.}\ A}}\skt{ }\edtext{\skt{विकटम्वा}}{\lemma{\skt{विकटम्वा}}\Afootnote{H; \emph{om.}\ A}}\skt{ }\edtext{\skt{जनयति~।त\-स्मा\-द\-ना\-ग\-तावी न}}{\lemma{\skt{जनयति~।त\appbrk{}स्मा\appbrk{}द\appbrk{}ना\appbrk{}ग\appbrk{}तावी न}}\Afootnote{H; \skt{तावी}\ A}}\skt{ }\edtext{\skt{प्रवाहेत}}{\lemma{\skt{प्रवाहेत}}\Afootnote{H; \skt{प्रवा}\ A}}\skt{~। }\edtext{\skt{स्रस्ते तु}}{\lemma{\skt{स्रस्ते तु}}\Afootnote{H; \skt{हस्व, त}\ A}}\skt{ }\edtext{\skt{कुक्षौ}}{\lemma{\skt{कुक्षौ}}\Afootnote{H; \skt{तो}\ A}}\skt{ विमुक्ते }\edtext{\skt{गर्भनाडी}}{\lemma{\skt{गर्भनाडी}}\Afootnote{A; \skt{गर्ब्भनाडी}\ H}}\skt{ प्रवन्धे }\edtext{\skt{समूले}}{\lemma{\skt{समूले}}\Afootnote{H; \skt{सशूलेषु}\ A}}\skt{ }\edtext{\skt{श्रोणी\-वं\-क्ष\-ण\-ब\-स्ति\-शिरः}}{\lemma{\skt{श्रोणी\appbrk{}वं\appbrk{}क्ष\appbrk{}ण\appbrk{}ब\appbrk{}स्ति\appbrk{}शिरः}}\Afootnote{H; \skt{श्रोणि\appbrk{}व\appbrk{}ङ्क्ष\appbrk{}ण\appbrk{}ब\appbrk{}स्ति\appbrk{}शिरः}\ A}}\skt{ }\edtext{\skt{सुप्रवाहेत}}{\lemma{\skt{सुप्रवाहेत}}\Afootnote{H; \skt{सु च प्रवाहेथाः}\ A}}\skt{ शनैः शनैस्ततो }\edtext{\skt{गर्भयोनिमुखं}}{\lemma{\skt{गर्भयोनिमुखं}}\Afootnote{\skt{गर्भनिर्गमे प्रगाढं, ततो गर्भे योनिमुखं}\ A; \skt{गर्ब्भयोनिमुखं}\ H}}\skt{ }\edtext{\edtext{\skt{प्रप\-न्न\-गा\-ढ\-त\-र\-मा\-वि\-श\-ल्य\-भा\-वात्}}{\lemma{\skt{प्रप\appbrk{}न्न\appbrk{}गा\appbrk{}ढ\appbrk{}त\appbrk{}रम्}}\Afootnote{H; \skt{प्रपन्ने गाढतर}\ A}}}{\lemma{\skt{आवि\appbrk{}श\appbrk{}ल्य\appbrk{}भा\appbrk{}वात्}}\Afootnote{H; \skt{मावि\appbrk{}श\appbrk{}ल्य\appbrk{}भा\appbrk{}वात्; अका\appbrk{}ल\appbrk{}प्र\appbrk{}वा\appbrk{}ह\appbrk{}णा\appbrk{}द्ब\appbrk{}धिरं मूकं कुब्जं व्यस्त\appbrk{}ह\appbrk{}नु\appbrk{}मू\appbrk{}र्ध्वा\appbrk{}भि\appbrk{}घा\appbrk{}तिनं कास\appbrk{}श्वा\appbrk{}स\appbrk{}शो\appbrk{}षो\appbrk{}प\appbrk{}द्रुतं विकटं वा जनयति}\ A}}\skt{~।}\pend

\pstart \ledsidenote{\passref{3.10.10}}\skt{तत्र }\edtext{\skt{प्रति\-लो\-म\-म\-नु\-पानं}}{\lemma{\skt{प्रति\appbrk{}लो\appbrk{}म\appbrk{}म\appbrk{}नु\appbrk{}पानं}}\Afootnote{\skt{प्रतिलोममनु}\ A; \skt{प्रति\appbrk{}लो\appbrk{}म\appbrk{}नु\appbrk{}पानं}\ H}}\skt{ }\edtext{\edtext{\skt{पाय\-ये\-त्शो\-स\-येत्}}{\lemma{\skt{पाययेत्}}\Afootnote{H; \emph{om.}\ A}}}{\lemma{\skt{शोसयेत्}}\Afootnote{H; \skt{लोमयेत्, प्राञ्ज\appbrk{}ल\appbrk{}मा\appbrk{}क\appbrk{}र्षेत्}\ A}}\skt{~।}\pend

\pstart \ledsidenote{\passref{3.10.12}}\edtext{\skt{जातस्योल्बं}}{\lemma{\skt{जातस्योल्बं}}\Afootnote{H; \skt{अथ जातस्योल्बं}\ A}}\skt{ }\edtext{\skt{विशोध्यास्यं}}{\lemma{\skt{विशोध्यास्यं}}\Afootnote{H; \skt{मुखं च}\ A}}\skt{ }\edtext{\skt{सैन्ध\-व\-स\-र्प्पि\-षा\-ष्टा\-ङ्गु\-ला\-व\-शेषां}}{\lemma{\skt{सैन्ध\appbrk{}व\appbrk{}स\appbrk{}र्प्पि\appbrk{}षा\appbrk{}ष्टा\appbrk{}ङ्गु\appbrk{}ला\appbrk{}व\appbrk{}शेषां}}\Afootnote{H; \skt{सैन्धवसर्पिषा विशोध्य, घृताक्तं मूर्ध्नि पिचुं दद्यात्; ततो नाभि\appbrk{}ना\appbrk{}डी\appbrk{}म\appbrk{}ष्टाङ्गु}\ A}}\skt{ }\edtext{\edtext{\skt{नाडीम्बद्धयित्वा}}{\lemma{\skt{नाडीम्}}\Afootnote{H; \emph{om.}\ A}}}{\lemma{\skt{बद्धयित्वा}}\Afootnote{H; \emph{om.}\ A}}\skt{ }\edtext{\skt{कुष्ठ\-तै\-ले\-ना\-भ्यज्य}}{\lemma{\skt{कुष्ठ\appbrk{}तै\appbrk{}ले\appbrk{}ना\appbrk{}भ्यज्य}}\Afootnote{H; \skt{लमायम्य}\ A}}\skt{ }\edtext{\skt{सूत्रेणाबद्ध्य}}{\lemma{\skt{सूत्रेणाबद्ध्य}}\Afootnote{H; \skt{सूत्रेण बद्ध्वा छेदयेत्, तत्सूत्रैकदेशं च}\ A}}\skt{ कुमारस्य }\edtext{\skt{ग्रीवा\-या\-म्व\-ध्नी\-यात्}}{\lemma{\skt{वध्नीयात्}}\Afootnote{H; \skt{सम्यग्बध्नीयात्}\ A}}\skt{~।}\pend

\pstart \ledsidenote{\passref{3.10.13}}\skt{अथ कुमारं शीता\-भि\-र\-द्भि\-रा\-श्वास्य जातकर्म्मणि }\edtext{\skt{मधु\-स\-र्प्पि\-र\-न\-न्ता\-मिश्रं}}{\lemma{\skt{मधु\appbrk{}स\appbrk{}र्प्पि\appbrk{}र\appbrk{}न\appbrk{}न्ता\appbrk{}मिश्रं}}\Afootnote{H; \skt{कृते मधु\appbrk{}स\appbrk{}र्पि\appbrk{}र\appbrk{}नन्}\ A}}\skt{ }\edtext{\skt{व्रह्मीरसेन}}{\lemma{\skt{व्रह्मीरसेन}}\Afootnote{H; \skt{त}\ A}}\skt{ }\edtext{\skt{सुवर्ण्णमङ्गुल्या}}{\lemma{\skt{सुवर्ण्णमङ्गुल्या}}\Afootnote{H; \skt{चूर्णमङ्गुल्या}\ A}}\skt{ }\edtext{\skt{नासिकया}}{\lemma{\skt{नासिकया}}\Afootnote{H; \skt{नामिकया}\ A}}\skt{ लेहयेत्~। }\edtext{\skt{सूतिकाञ्च}}{\lemma{\skt{सूतिकाञ्च}}\Afootnote{H; \skt{ततो ब}\ A}}\skt{ }\edtext{\edtext{\skt{घृत\-तै\-ल\-यो\-र\-न्य\-त\-रस्य}}{\lemma{\skt{घृततैलयोर्}}\Afootnote{H; \skt{लातैलेनाभ्यज्य, क्षीर\appbrk{}वृ\appbrk{}क्ष\appbrk{}क\appbrk{}षा\appbrk{}येण सर्व\appbrk{}ग\appbrk{}न्धो\appbrk{}द\appbrk{}केन वा रूप्यहे}\ A}}\skt{ मात्रां}}{\lemma{\skt{अन्यतरस्य मात्रां}}\Afootnote{H; \skt{मप्रतप्तेन वा वारिणा स्नापयेदेनं कपित्थ}\ A}}\skt{ }\edtext{\skt{पिपल्यादि}}{\lemma{\skt{पिपल्यादि}}\Afootnote{H; \skt{पत्र}\ A}}\skt{ }\edtext{\edtext{\edtext{\skt{कषा\-य\-म\-नु\-पा\-न\-म्पा\-य\-येत्}}{\lemma{\skt{कषायम्}}\Afootnote{H; \skt{कषा}\ A}}}{\lemma{\skt{अनुपानम्}}\Afootnote{H; \emph{om.}\ A}}}{\lemma{\skt{पाययेत्}}\Afootnote{H; \skt{येण}\ A}}\skt{~। }\edtext{\edtext{\skt{वल\-व\-ती\-म\-वलां}}{\lemma{\skt{वलवतीम्}}\Afootnote{H; \skt{वा कोष्णेन यथा}\ A}}}{\lemma{\skt{अवलां}}\Afootnote{H; \skt{कालं}\ A}}\skt{ }\edtext{\edtext{\edtext{\skt{यवा\-गू\-न्त्रि\-रा\-त्र\-म्प\-ञ्च\-रा\-त्रम्वा}}{\lemma{\skt{यवागून्}}\Afootnote{H; \skt{यथादो}\ A}}}{\lemma{\skt{त्रिरात्रम्}}\Afootnote{H; \skt{षं यथा}\ A}}}{\lemma{\skt{पञ्चरात्रम्वा}}\Afootnote{H; \skt{विभवं च}\ A}}\skt{~॥ }\edtext{\skt{अत}}{\lemma{\skt{अत}}\Afootnote{H; \emph{om.}\ A}}\skt{ }\edtext{\skt{ऊर्द्ध्वं}}{\lemma{\skt{ऊर्द्ध्वं}}\Afootnote{H; \emph{om.}\ A}}\skt{ }\edtext{\edtext{\skt{स्निग्धे\-ना\-न्न\-सं\-स\-र्ग्गे\-णो\-प\-च\-रे\-त्स्ने\-ह\-नित्या}}{\lemma{\skt{स्निग्धे\appbrk{}ना\appbrk{}न्न\appbrk{}सं\appbrk{}स\appbrk{}र्ग्गे\appbrk{}णो\appbrk{}प\appbrk{}च\appbrk{}रेत्}}\Afootnote{H; \emph{om.}\ A}}\skt{ च}}{\lemma{\skt{स्नेहनित्या च}}\Afootnote{H; \emph{om.}\ A}}\skt{ }\edtext{\skt{स्यात्}}{\lemma{\skt{स्यात्}}\Afootnote{H; \emph{om.}\ A}}\skt{~।}\pend

\pstart \ledsidenote{\passref{3.10.18}}\edtext{\skt{प्राय\-श्चै\-ना\-म्प्र\-भू\-ते\-नो\-ष्णो\-द\-केन}}{\lemma{\skt{प्रायश्चैनाम्}}\Afootnote{\skt{प्रायशश्चैनां}\ A; \skt{प्रायश्चैनाम्}\ H}}\skt{ परि\-षि\-ञ्चे\-त्क्रो\-धा\-या\-स\-मै\-थु\-ना\-दींश्च परिहरेत्~॥}\pend

\pstart\versepar \ledsidenote{\passref{3.10.19}}\edtext{\skt{भवति}}{\lemma{\skt{भवति}}\Afootnote{H; \emph{om.}\ A}}\skt{ }\edtext{\skt{चात्र}}{\lemma{\skt{चात्र}}\Afootnote{H; \emph{om.}\ A}}\skt{~॥ }\pend
\pstart\versepar \skt{}\edtext{\skt{मिथ्या\-चा\-रा\-त्सू\-ति\-काया यो}}{\lemma{\skt{सूतिकाया यो}}\Afootnote{A; \skt{सूतिकाया यो}\ H}}\skt{ }\edtext{\skt{व्याधि\-रु\-प\-जा\-यते}}{\lemma{\skt{व्याधिर्}}\Afootnote{A; \skt{व्याधि}\ H}}\skt{~। }\pend
\pstart\versepar \skt{स कृच्छ्रसाध्योऽसाध्यो वा }\edtext{\skt{भवे\-द्धा\-त्व\-प\-क\-र्ष\-र्णात्}}{\lemma{\skt{भवे\appbrk{}द्धा\appbrk{}त्व\appbrk{}प\appbrk{}क\appbrk{}र्ष\appbrk{}र्णात्}}\Afootnote{\skt{भवे\appbrk{}द\appbrk{}त्य\appbrk{}प\appbrk{}त\appbrk{}र्प\appbrk{}णात्}\ A; \skt{भवे\appbrk{}द्धा\appbrk{}त्व\appbrk{}प\appbrk{}क\appbrk{}र्ष\appbrk{}र्णात्}\ H}}\skt{~॥}\edtext{}{\lemma{\skt{उपजायते~। स}}\Afootnote{A; \skt{रुपजायते~। स}\ H}}\pend

\pstart\versepar \ledsidenote{\passref{3.10.20}}\skt{तस्मात्तां }\edtext{\skt{देशकालौ तु}}{\lemma{\skt{देशकालौ तु}}\Afootnote{H; \skt{देशकालौ च}\ A}}\skt{ व्याधिसात्म्येन कर्म्मणा~। }\pend
\pstart\versepar \skt{परी\-क्ष्यो\-प\-च\-रे\-न्नि\-त्य\-मेवं नात्य\-य\-मा\-प्नु\-यात्~॥}\pend

\pstart \ledsidenote{\passref{3.10.21}}\edtext{\edtext{\edtext{\edtext{\skt{अथा\-म\-रा\-म\-प\-त\-न्ती\-म्पा\-त\-ये\-त्क\-ण्ठ\-म\-स्याः\-के\-श\-वे\-ष्टि\-त\-या\-ङ्गुल्या}}{\lemma{\skt{अथा\appbrk{}म\appbrk{}रा\appbrk{}म\appbrk{}प\appbrk{}त\appbrk{}न्तीम्}}\Afootnote{H; \skt{अथा\appbrk{}प\appbrk{}राऽ\appbrk{}प\appbrk{}तन्त्य्}\ A}}}{\lemma{\skt{पातयेत्}}\Afootnote{H; \skt{आनाहाध्मानौ कुरुते, तस्मात्}\ A}}}{\lemma{\skt{कण्ठम्}}\Afootnote{H; \skt{कण्ठ}\ A}}}{\lemma{\skt{अस्याः\appbrk{}के\appbrk{}श\appbrk{}वे\appbrk{}ष्टि\appbrk{}त\appbrk{}या\appbrk{}ङ्गुल्या}}\Afootnote{H; \skt{मस्याः केश\appbrk{}वे\appbrk{}ष्टि\appbrk{}त\appbrk{}या\appbrk{}ङ्गुल्या}\ A}}\skt{ }\edtext{\skt{प्रमृशेत्कटुका}}{\lemma{\skt{प्रमृशेत्}}\Afootnote{H; \skt{प्रमृजेत्}\ A}}\skt{ }\edtext{\skt{लावु\-कृ\-त\-वे\-ध\-न\-स\-र्ष\-प\-स\-र्प्प\-नि\-र्म्मो\-च\-कैर्व्वा}}{\lemma{\skt{लावु\appbrk{}कृ\appbrk{}त\appbrk{}वे\appbrk{}ध\appbrk{}न\appbrk{}स\appbrk{}र्ष\appbrk{}प\appbrk{}स\appbrk{}र्प्प\appbrk{}नि\appbrk{}र्म्मो\appbrk{}च\appbrk{}कैर्व्वा}}\Afootnote{H; \skt{लाबु\appbrk{}कृ\appbrk{}त\appbrk{}वे\appbrk{}ध\appbrk{}न\appbrk{}स\appbrk{}र्ष\appbrk{}प\appbrk{}स\appbrk{}र्प\appbrk{}नि\appbrk{}र्मो\appbrk{}कैर्वा}\ A}}\skt{ }\edtext{\edtext{\edtext{\skt{कुट\-तै\-लै\-र्यो\-नि\-मु\-द्रू\-प\-येत्}}{\lemma{\skt{कुटतैलैर्}}\Afootnote{H; \skt{कटु\appbrk{}तै\appbrk{}ल\appbrk{}वि\appbrk{}मिश्रै}\ A}}}{\lemma{\skt{योनिम्}}\Afootnote{H; \skt{र्योनि}\ A}}}{\lemma{\skt{उद्रूपयेत्}}\Afootnote{H; \skt{मुखं धूपयेत्}\ A}}\skt{~। }\edtext{\skt{लाङ्ग\-ल\-की\-मू\-ल\-क\-ल्केन}}{\lemma{\skt{लाङ्ग\appbrk{}ल\appbrk{}की\appbrk{}मू\appbrk{}ल\appbrk{}क\appbrk{}ल्केन}}\Afootnote{H; \skt{लाङ्ग\appbrk{}ली\appbrk{}मू\appbrk{}ल\appbrk{}क\appbrk{}ल्केन}\ A}}\skt{ वास्याः }\edtext{\skt{पाणि\-पा\-द\-त\-ल\-मा\-लि\-म्पेत्}}{\lemma{\skt{आलिम्पेत्}}\Afootnote{H; \skt{आलिम्पेत्, मूर्ध्नि वास्या महा\appbrk{}वृ\appbrk{}क्ष\appbrk{}क्षी\appbrk{}र\appbrk{}म\appbrk{}नु\appbrk{}से\appbrk{}च\appbrk{}येत्}\ A}}\skt{~। }\edtext{\skt{कुष्ठ\-ला\-ङ्ग\-ल\-की\-मू\-ल\-क\-ल्कम्वा}}{\lemma{\skt{कुष्ठ\appbrk{}ला\appbrk{}ङ्ग\appbrk{}ल\appbrk{}की\appbrk{}मू\appbrk{}ल\appbrk{}क\appbrk{}ल्कम्वा}}\Afootnote{H; \skt{कुष्ठ\appbrk{}ला\appbrk{}ङ्ग\appbrk{}ली\appbrk{}मू\appbrk{}ल\appbrk{}कल्कं वा}\ A}}\skt{ मद्य\-मू\-त्र\-यो\-र\-न्य\-त\-रेण पाययेत्~। }\edtext{\edtext{\skt{शालि\-मू\-ल\-क\-ल्क\-म्वाऽपि}}{\lemma{\skt{शालि\appbrk{}मू\appbrk{}ल\appbrk{}क\appbrk{}ल्कम्वा}}\Afootnote{H; \skt{शालमूलकल्कं वा}\ A}}}{\lemma{\skt{ऽपि}}\Afootnote{\skt{पि}\ A H}}\skt{ }\edtext{\skt{प्पल्यादिकल्कम्वा}}{\lemma{\skt{प्पल्यादिकल्कम्वा}}\Afootnote{H; \skt{प्पल्यादिं वा}\ A}}\skt{ मद्येन~। }\edtext{\skt{सिद्धा\-र्थ\-क\-कु\-ष्ठै\-ला\-ला\-ङ्ग\-ल\-की\-म\-हा\-वृ\-क्ष\-क्षी\-र\-सु\-रा\-म\-ण्डै\-रा\-स्था\-प\-येत्~। एतैष्वेव च}}{\lemma{\skt{सिद्धा\appbrk{}र्थ\appbrk{}क\appbrk{}कु\appbrk{}ष्ठै\appbrk{}ला\appbrk{}ला\appbrk{}ङ्ग\appbrk{}ल\appbrk{}की\appbrk{}म\appbrk{}हा\appbrk{}वृ\appbrk{}क्ष\appbrk{}क्षी\appbrk{}र\appbrk{}सु\appbrk{}रा\appbrk{}म\appbrk{}ण्डै\appbrk{}रा\appbrk{}स्था\appbrk{}प\appbrk{}येत्~। एतैष्वेव च}}\Afootnote{H; \skt{सिद्धा\appbrk{}र्थ\appbrk{}क\appbrk{}कु\appbrk{}ष्ठ\appbrk{}ला\appbrk{}ङ्ग\appbrk{}ली\appbrk{}म\appbrk{}हा\appbrk{}वृ\appbrk{}क्ष\appbrk{}क्षी\appbrk{}र\appbrk{}मि\appbrk{}श्रेण सुरामण्डेन वास्थापयेत्, एतैरेव}\ A}}\skt{ सिद्धेन }\edtext{\edtext{\skt{सिद्धा\-र्थ\-क\-तै\-ले\-नो\-त्त\-र\-चू\-र्ण्ण\-ब\-स्ति\-श्चास्यै}}{\lemma{\skt{सिद्धा\appbrk{}र्थ\appbrk{}क\appbrk{}तै\appbrk{}ले\appbrk{}नो\appbrk{}त्त\appbrk{}र\appbrk{}चू\appbrk{}र्ण्ण\appbrk{}ब\appbrk{}स्तिश्}}\Afootnote{H; \skt{सिद्धा\appbrk{}र्थ\appbrk{}क\appbrk{}तै\appbrk{}ले\appbrk{}नो\appbrk{}त्त\appbrk{}र\appbrk{}बस्तिं}\ A}}}{\lemma{\skt{चास्यै}}\Afootnote{H; \emph{om.}\ A}}\skt{ दद्यात्~। }\edtext{\skt{हस्तेन}}{\lemma{\skt{हस्तेन}}\Afootnote{H; \skt{स्निग्धेन वा कृत्तनखेन हस्ते}\ A}}\skt{ }\edtext{\skt{चापहरेत्}}{\lemma{\skt{चापहरेत्}}\Afootnote{H; \skt{नापहरेत्}\ A}}\skt{~।}\pend

\pstart \ledsidenote{\passref{3.10.22}}\edtext{\edtext{\skt{प्रजातायास्तु}}{\lemma{\skt{प्रजातायास्}}\Afootnote{H; \skt{प्रजातायाश्च}\ A}}}{\lemma{\skt{तु}}\Afootnote{H; \skt{नार्या रूक्ष\appbrk{}श\appbrk{}री\appbrk{}रा\appbrk{}या\appbrk{}स्ती\appbrk{}क्ष्णै\appbrk{}र\appbrk{}वि\appbrk{}शो\appbrk{}धितं रक्तं वायुना तद्दे\appbrk{}श\appbrk{}ग\appbrk{}ते\appbrk{}ना\appbrk{}ति\appbrk{}सं\appbrk{}रुद्धं नाभेरधः पार्श्वयोर्बस्तौ}\ A}}\skt{ }\edtext{\skt{बस्ति\-शि\-रो\-द\-र\-यो\-नि\-शूलं}}{\lemma{\skt{बस्ति\appbrk{}शि\appbrk{}रो\appbrk{}द\appbrk{}र\appbrk{}यो\appbrk{}नि\appbrk{}शूलं}}\Afootnote{\skt{बस्तिशिरसि वा ग्रन्थिं करोति; ततश्च नाभि\appbrk{}ब\appbrk{}स्त्यु\appbrk{}द\appbrk{}र\appbrk{}शू\appbrk{}लानि भवन्ति, सूचीभिरिव निस्तुद्यते भिद्यते दीर्यत इव च पक्वाशयः}\ A; \skt{वस्तिशिरो}\ H}}\skt{ }\edtext{\skt{रक्तं}}{\lemma{\skt{रक्तं}}\Afootnote{\skt{सम\appbrk{}न्ता\appbrk{}दा\appbrk{}ध्मा\appbrk{}न\appbrk{}मुद}\ A; \skt{दरयोनिशूलं रक्तं}\ H}}\skt{ }\edtext{\skt{संरुद्ध्य}}{\lemma{\skt{संरुद्ध्य}}\Afootnote{H; \skt{रे}\ A}}\skt{ }\edtext{\skt{वायुः करोति}}{\lemma{\skt{वायुः करोति}}\Afootnote{H; \skt{मूत्रसङ्गश्च भवतीति मक्क}\ A}}\skt{ }\edtext{\skt{स}}{\lemma{\skt{स}}\Afootnote{H; \skt{ल्ल}\ A}}\skt{ }\edtext{\skt{मक्कलः}}{\lemma{\skt{मक्कलः}}\Afootnote{H; \skt{लक्षणम्}\ A}}\skt{~। }\edtext{\skt{तत्र}}{\lemma{\skt{तत्र}}\Afootnote{H; \skt{तत्र वीर\appbrk{}त\appbrk{}र्वा\appbrk{}दि\appbrk{}सिद्धं जल\appbrk{}मु\appbrk{}ष\appbrk{}का\appbrk{}दिप्र}\ A}}\skt{ }\edtext{\skt{यवक्षारचूर्ण्णं}}{\lemma{\skt{यवक्षारचूर्ण्णं}}\Afootnote{H; \skt{तीवापं पाययेत्, यवक्षारचूर्णं}\ A}}\skt{ }\edtext{\skt{सर्प्पिषा}}{\lemma{\skt{सर्प्पिषा}}\Afootnote{H; \skt{वा}\ A}}\skt{ }\edtext{\skt{सुखोदकेन}}{\lemma{\skt{सुखोदकेन}}\Afootnote{H; \skt{सुखोदकेन पिप्प\appbrk{}ल्या\appbrk{}दि\appbrk{}क्वा\appbrk{}थेन वा}\ A}}\skt{ }\edtext{\skt{पाययेत्}}{\lemma{\skt{पाययेत्}}\Afootnote{H; \skt{पिप्पल्या}\ A}}\skt{~॥ }\edtext{\skt{लवणचूर्ण्णम्वा}}{\lemma{\skt{लवणचूर्ण्णम्वा}}\Afootnote{H; \skt{दिचूर्णं वा}\ A}}\skt{ सुरामण्डेन~। }\edtext{\skt{पुराणगुडम्वा}}{\lemma{\skt{पुराणगुडम्वा}}\Afootnote{H; \skt{वरुणादिक्वाथं वा पञ्च\appbrk{}को\appbrk{}लै\appbrk{}ला\appbrk{}प्र\appbrk{}ती\appbrk{}वापं, पृथ\appbrk{}क्प\appbrk{}र्ण्या\appbrk{}दि\appbrk{}क्वाथं वा भद्र\appbrk{}दा\appbrk{}रु\appbrk{}म\appbrk{}रि\appbrk{}च\appbrk{}सं\appbrk{}सृष्टं, पुराणगुडं वा}\ A}}\skt{ त्रिकटुक }\edtext{\skt{त्रिजा\-त\-क\-कु\-स्तु\-म्बु\-रु\-मि\-श्रम्वा}}{\lemma{\skt{त्रिजा\appbrk{}त\appbrk{}क\appbrk{}कु\appbrk{}स्तु\appbrk{}म्बु\appbrk{}रु\appbrk{}मि\appbrk{}श्रम्वा}}\Afootnote{H; \skt{चतु\appbrk{}र्जा\appbrk{}त\appbrk{}क\appbrk{}कु\appbrk{}स्तु\appbrk{}म्बु\appbrk{}रु\appbrk{}मिश्रं}\ A}}\skt{ }\edtext{\skt{खादयेत्}}{\lemma{\skt{खादयेत्}}\Afootnote{H; \emph{om.}\ A}}\skt{~। }\edtext{\skt{पृथ\-क्प\-र्ण्ण्या\-दि\-क्वाथं}}{\lemma{\skt{पृथ\appbrk{}क्प\appbrk{}र्ण्ण्या\appbrk{}दि\appbrk{}क्वाथं}}\Afootnote{H; \emph{om.}\ A}}\skt{ }\edtext{\edtext{\skt{भद्र\-दा\-रु\-म\-रि\-च\-व\-चा\-सं\-सृ\-ष्ट\-म्पा\-य\-येत्}}{\lemma{\skt{भद्र\appbrk{}दा\appbrk{}रु\appbrk{}म\appbrk{}रि\appbrk{}च\appbrk{}व\appbrk{}चा\appbrk{}सं\appbrk{}सृ\appbrk{}ष्टम्}}\Afootnote{H; \emph{om.}\ A}}}{\lemma{\skt{पाययेत्}}\Afootnote{H; \skt{खादेत्, अच्छं वा पिबे\appbrk{}द\appbrk{}रि\appbrk{}ष्ट\appbrk{}मिति}\ A}}\skt{~।}\pend

\pstart \ledsidenote{\passref{3.10.23}}\skt{अथ बालं }\edtext{\skt{क्षौमपरिहितं}}{\lemma{\skt{क्षौमपरिहितं}}\Afootnote{H; \skt{क्षौमपरिवृतं}\ A}}\skt{ }\edtext{\edtext{\skt{क्षौम\-का\-र्प्पा\-सा\-स्तृ\-तायां}}{\lemma{\skt{क्षौमकार्प्पासास्}}\Afootnote{H; \skt{क्षौमवस्त्रास्}\ A}}}{\lemma{\skt{तृतायां}}\Afootnote{H; \skt{तृतायां शय्यायां}\ A}}\skt{ }\edtext{\skt{शाययीत}}{\lemma{\skt{शाययीत}}\Afootnote{H; \skt{शाययेत्}\ A}}\skt{~। }\edtext{\skt{निम्ब\-पी\-लु\-प\-रु\-ष\-क\-शा\-खा\-भि\-श्चैनं}}{\lemma{\skt{निम्ब\appbrk{}पी\appbrk{}लु\appbrk{}प\appbrk{}रु\appbrk{}ष\appbrk{}क\appbrk{}शा\appbrk{}खा\appbrk{}भिश्}}\Afootnote{H; \skt{पीलु\appbrk{}ब\appbrk{}द\appbrk{}री\appbrk{}नि\appbrk{}म्ब\appbrk{}प\appbrk{}रू\appbrk{}ष\appbrk{}क\appbrk{}शा\appbrk{}खा\appbrk{}भिश्}\ A}}\skt{ परिवीजयेत्~। मूर्द्ध्नि }\edtext{\skt{चास्य}}{\lemma{\skt{चास्य}}\Afootnote{H; \skt{चास्याहरह}\ A}}\skt{ }\edtext{\edtext{\skt{तैल\-पि\-चु\-म\-व\-ता\-र\-येत्}}{\lemma{\skt{तैलपिचुम्}}\Afootnote{H; \skt{स्तैलपिचु}\ A}}}{\lemma{\skt{अवतारयेत्}}\Afootnote{H; \skt{मवचारयेत्}\ A}}\skt{~। }\edtext{\skt{धूपयेच्चैनं}}{\lemma{\skt{धूपयेच्}}\Afootnote{H; \skt{धूपयेच्}\ A}}\skt{ रक्षोघ्नैर्धूपैः~। रक्षोघ्नानि चास्य }\edtext{\edtext{\skt{पाणि\-पा\-द\-शि\-रो\-ग्री\-वा\-स्व\-व\-सु\-जेत्}}{\lemma{\skt{पाणि\appbrk{}पा\appbrk{}द\appbrk{}शि\appbrk{}रो\appbrk{}ग्री\appbrk{}वास्व्}}\Afootnote{A; \skt{पाणि\appbrk{}पा\appbrk{}द\appbrk{}शि\appbrk{}रो\appbrk{}ग्री\appbrk{}वास्}\ H}}}{\lemma{\skt{अवसुजेत्}}\Afootnote{\skt{अवसृजेत्}\ A; \skt{ववसुजेत्}\ H}}\skt{~। }\edtext{\skt{तिला\-त\-सी\-स\-र्ष\-प\-क\-णा\-श्चात्र}}{\lemma{\skt{तिला\appbrk{}त\appbrk{}सी\appbrk{}स\appbrk{}र्ष\appbrk{}प\appbrk{}क\appbrk{}णाश्}}\Afootnote{H; \skt{तिला\appbrk{}त\appbrk{}सी\appbrk{}स\appbrk{}र्ष\appbrk{}प\appbrk{}क\appbrk{}णांश्}\ A}}\skt{ प्रकिरेत्~। अधिष्ठाने चाग्नि\-म्प्र\-ज्वा\-ल\-येत्~। }\edtext{\skt{व्रणि\-तो\-पा\-स\-नी\-य\-ञ्चा\-वे\-क्ष्येत}}{\lemma{\skt{चावेक्ष्येत}}\Afootnote{H; \skt{चावेक्षेत}\ A}}\pend

\pstart \ledsidenote{\passref{3.10.24}}\skt{ततो }\edtext{\skt{दशम्यां}}{\lemma{\skt{दशम्यां}}\Afootnote{H; \skt{दशमेऽ}\ A}}\skt{ }\edtext{\skt{मातापितरौ}}{\lemma{\skt{मातापितरौ}}\Afootnote{H; \skt{हनि मातापितरौ}\ A}}\skt{ कृत\-म\-ङ्ग\-ल\-कौ\-तुकौ }\edtext{\skt{स्वस्ति\-वा\-च\-क\-पू\-र्व्वकं}}{\lemma{\skt{स्वस्ति\appbrk{}वा\appbrk{}च\appbrk{}क\appbrk{}पू\appbrk{}र्व्वकं}}\Afootnote{H; \skt{स्वस्तिवाचनं कृत्वा}\ A}}\skt{ नाम कुर्यातां~। }\edtext{\skt{यथाभिप्रेतं}}{\lemma{\skt{यथाभिप्रेतं}}\Afootnote{H; \skt{यदभिप्रेतं}\ A}}\skt{ }\edtext{\skt{नक्षत्रन्नाम}}{\lemma{\skt{नक्षत्रन्}}\Afootnote{H; \skt{नक्षत्र}\ A}}\skt{ वा~।}\pend

\pstart \ledsidenote{\passref{3.10.25}}\skt{ततो यथावर्ण्णं }\edtext{\edtext{\skt{धात्री\-मु\-प\-स्था\-प\-येत्}}{\lemma{\skt{धात्रीम्}}\Afootnote{H; \skt{धात्री}\ A}}}{\lemma{\skt{उपस्थापयेत्}}\Afootnote{H; \skt{मुपे\appbrk{}या\appbrk{}न्म\appbrk{}ध्य\appbrk{}म\appbrk{}प्र\appbrk{}माणां मध्य\appbrk{}म\appbrk{}व\appbrk{}य\appbrk{}स्का(\appbrk{}सा)\appbrk{}म\appbrk{}रोगां शील\appbrk{}व\appbrk{}ती\appbrk{}म\appbrk{}च\appbrk{}प\appbrk{}ला\appbrk{}म\appbrk{}लो\appbrk{}लु\appbrk{}पाम}\ A}}\skt{~। }\edtext{\edtext{\edtext{\edtext{\edtext{\edtext{\skt{कृषा\-म\-स्थू\-ला\-म्व\-यः\-स्था\-म\-म\-व्या\-धि\-ता\-म्प्र\-स\-न्न\-क्षी\-रा\-म\-व\-ल\-म्वो\-र्द्ध\-स्त\-नी\-म\-व्य\-ङ्गा\-म\-व्य\-सनं}}{\lemma{\skt{कृषामस्थूलाम्}}\Afootnote{H; \skt{कृशामस्थूलां}\ A}}}{\lemma{\skt{वयःस्थामम्}}\Afootnote{H; \emph{om.}\ A}}}{\lemma{\skt{अव्याधिताम्}}\Afootnote{H; \emph{om.}\ A}}}{\lemma{\skt{प्रस\appbrk{}न्न\appbrk{}क्षी\appbrk{}रा\appbrk{}म\appbrk{}व\appbrk{}ल\appbrk{}म्वो\appbrk{}र्द्ध\appbrk{}स्त\appbrk{}नीम्}}\Afootnote{H; \skt{प्रस\appbrk{}न्न\appbrk{}क्षी\appbrk{}रा\appbrk{}म\appbrk{}ल\appbrk{}म्बौ\appbrk{}ष्ठी\appbrk{}म\appbrk{}ल\appbrk{}म्बो\appbrk{}र्ध्व\appbrk{}स्तनी}\ A}}}{\lemma{\skt{अव्यङ्गाम्}}\Afootnote{H; \skt{मव्यङ्गा}\ A}}}{\lemma{\skt{अव्यसनं}}\Afootnote{H; \skt{मव्यसनि}\ A}}\skt{ }\edtext{\skt{नीजी\-व\-व\-त्सां\-व\-त्सलां}}{\lemma{\skt{नीजी\appbrk{}व\appbrk{}व\appbrk{}त्सां\appbrk{}व\appbrk{}त्सलां}}\Afootnote{H; \skt{नीं जीवद्वत्सां दोग्ध्रीं वत्सला}\ A}}\skt{ }\edtext{\edtext{\skt{दोग्धी\-म\-क्षु\-द्रां\-कुले}}{\lemma{\skt{दोग्धीम्}}\Afootnote{H; \emph{om.}\ A}}}{\lemma{\skt{अक्षुद्रांकुले}}\Afootnote{H; \skt{मक्षुद्रकर्मिणीं कुले}\ A}}\skt{ }\edtext{\skt{जातमतो}}{\lemma{\skt{जातमतो}}\Afootnote{H; \skt{जातामतो}\ A}}\skt{ }\edtext{\skt{भूयिष्टैश्च}}{\lemma{\skt{भूयिष्टैश्च}}\Afootnote{H; \skt{भूयिष्ठैश्च}\ A}}\skt{ }\edtext{\edtext{\edtext{\skt{गुणै\-र\-न्वि\-ता\-मा\-रो\-ग्य\-व\-ल\-वृ\-द्धये}}{\lemma{\skt{गुणैर्}}\Afootnote{H; \skt{गुणै}\ A}}}{\lemma{\skt{अन्विताम्}}\Afootnote{H; \skt{रन्वितां}\ A}}}{\lemma{\skt{आरो\appbrk{}ग्य\appbrk{}व\appbrk{}ल\appbrk{}वृ\appbrk{}द्धये}}\Afootnote{H; \skt{श्मामा\appbrk{}मा\appbrk{}रो\appbrk{}ग्य\appbrk{}ब\appbrk{}ल\appbrk{}वृ\appbrk{}द्धये}\ A}}\skt{ }\edtext{\skt{कुमारस्य}}{\lemma{\skt{कुमारस्य}}\Afootnote{H; \skt{बालस्य}\ A}}\skt{~। }\edtext{\skt{तत्रा\-ति\-कृ\-षा\-ति\-स्थू\-ला\-नि\-न्दिता}}{\lemma{\skt{तत्रा\appbrk{}ति\appbrk{}कृ\appbrk{}षा\appbrk{}ति\appbrk{}स्थू\appbrk{}ला\appbrk{}नि\appbrk{}न्दिता}}\Afootnote{H; \skt{त}\ A}}\skt{ }\edtext{\skt{व्याधिताश्च}}{\lemma{\skt{व्याधिताश्च}}\Afootnote{H; \emph{om.}\ A}}\skt{ }\edtext{\edtext{\edtext{\skt{व्याख्या\-ता\-स्ता\-वि\-ष\-म\-धा\-तु\-त्वा\-द्रो\-गा\-ना\-पा\-द\-यन्ति}}{\lemma{\skt{व्याख्यातास्}}\Afootnote{H; \emph{om.}\ A}}}{\lemma{\skt{तावि\appbrk{}ष\appbrk{}म\appbrk{}धा\appbrk{}तु\appbrk{}त्वाद्}}\Afootnote{H; \skt{त्}\ A}}}{\lemma{\skt{रोगा\appbrk{}ना\appbrk{}पा\appbrk{}द\appbrk{}यन्ति}}\Afootnote{H; \skt{रो}\ A}}\skt{ }\edtext{\skt{बालस्य}}{\lemma{\skt{बालस्य}}\Afootnote{H; \emph{om.}\ A}}\skt{ }\edtext{\skt{विशेषतः}}{\lemma{\skt{विशेषतः}}\Afootnote{H; \emph{om.}\ A}}\skt{~। }\edtext{\skt{ऊर्ध्वस्तनी}}{\lemma{\skt{ऊर्ध्वस्तनी}}\Afootnote{H; \skt{र्ध्वस्तनी}\ A}}\skt{ करालं }\edtext{\skt{कुमारं}}{\lemma{\skt{कुमारं}}\Afootnote{H; \emph{om.}\ A}}\skt{ }\edtext{\skt{कुर्याल्लम्वस्तनी}}{\lemma{\skt{कुर्याल्}}\Afootnote{H; \skt{कुर्यात्}\ A}}\skt{ नासिकामुखं च्छादयित्वा }\edtext{\skt{मर\-ण\-मा\-पा\-द\-यन्ति}}{\lemma{\skt{आपादयन्ति}}\Afootnote{H; \skt{आपादयेत्}\ A}}\skt{~। }\edtext{\skt{अथ}}{\lemma{\skt{अथ}}\Afootnote{H; \skt{ततः}\ A}}\skt{ }\edtext{\skt{प्रशस्तां}}{\lemma{\skt{प्रशस्तां}}\Afootnote{H; \skt{प्रशस्तायां}\ A}}\skt{ }\edtext{\skt{शिरस्नाता}}{\lemma{\skt{शिरस्नाता}}\Afootnote{H; \skt{तिथौ शिरःस्नातम्}\ A}}\skt{ }\edtext{\edtext{\edtext{\skt{मह\-द्वा\-स\-सा\-मु\-द\-ङ्मु\-खी\-मु\-प\-वेश्य}}{\lemma{\skt{महद्वाससाम्}}\Afootnote{H; \skt{अह\appbrk{}त\appbrk{}वा\appbrk{}स\appbrk{}सम्}\ A}}}{\lemma{\skt{उदङ्मुखीम्}}\Afootnote{H; \skt{उदङ्मुखं शिशुमुपवेश्य धात्रीं प्राङ्मुखीं}\ A}}}{\lemma{\skt{उपवेश्य}}\Afootnote{H; \skt{चोपवेश्य}\ A}}\skt{ दक्षिणं स्तनं }\edtext{\skt{कुमारं}}{\lemma{\skt{कुमारं}}\Afootnote{H; \emph{om.}\ A}}\skt{ }\edtext{\edtext{\skt{पूर्व्वम्पाययेत्}}{\lemma{\skt{पूर्व्वम्}}\Afootnote{H; \emph{om.}\ A}}\skt{~॥ अनेन}}{\lemma{\skt{पाययेत्~॥ अनेन}}\Afootnote{H; \skt{धौतमीष}\ A}}\skt{ }\edtext{\skt{मन्त्रेणाभिमन्त्र्य}}{\lemma{\skt{मन्त्रेणाभिमन्त्र्य}}\Afootnote{H; \skt{त्परि\appbrk{}स्रु\appbrk{}त\appbrk{}म\appbrk{}भि\appbrk{}मन्त्र्य मन्त्रेणानेन पाययेत्}\ A}}\skt{~॥}\pend

\pstart\versepar \ledsidenote{\passref{3.10.26}}\skt{चत्वारः }\edtext{\skt{सागराः}}{\lemma{\skt{सागराः}}\Afootnote{H; \skt{सागरा}\ A}}\skt{ }\edtext{\skt{पुण्याः}}{\lemma{\skt{पुण्याः}}\Afootnote{H; \skt{स्तुभ्यं}\ A}}\skt{ }\edtext{\skt{स्तनानां}}{\lemma{\skt{स्तनानां}}\Afootnote{H; \skt{स्तनयोः}\ A}}\skt{ }\edtext{\skt{क्षीरवाहिनां}}{\lemma{\skt{क्षीरवाहिनां}}\Afootnote{H; \skt{क्षीरवाहिनः}\ A}}\skt{~। }\pend
\pstart\versepar \skt{}\edtext{\skt{सन्तु ते}}{\lemma{\skt{सन्तु ते}}\Afootnote{H; \skt{भवन्तु}\ A}}\skt{ }\edtext{\skt{शुभगे}}{\lemma{\skt{शुभगे}}\Afootnote{H; \skt{सुभगे}\ A}}\skt{ नित्यं }\edtext{\skt{बाल\-स्या\-नु\-व\-ल\-प्रदाः}}{\lemma{\skt{बाल\appbrk{}स्या\appbrk{}नु\appbrk{}व\appbrk{}ल\appbrk{}प्रदाः}}\Afootnote{H; \skt{बालस्य बलवृद्धये}\ A}}\skt{~।}\pend

\pstart\versepar \ledsidenote{\passref{3.10.27}}\edtext{\edtext{\edtext{\skt{पिबे\-न्बा\-लोऽ\-मृ\-त\-रसं}}{\lemma{\skt{पिबेन्}}\Afootnote{H; \emph{om.}\ A}}}{\lemma{\skt{बालो}}\Afootnote{\skt{पयो}\ A; \skt{वालो}\ H}}}{\lemma{\skt{ऽमृतरसं}}\Afootnote{A; \skt{मृतरसं}\ H}}\skt{ }\edtext{\edtext{\skt{पयस्तव}}{\lemma{\skt{पयस्}}\Afootnote{H; \skt{पीत्वा}\ A}}}{\lemma{\skt{तव}}\Afootnote{H; \skt{कुमारस्ते}\ A}}\skt{ शुभानने~। }\pend
\pstart\versepar \skt{}\edtext{\edtext{\edtext{\skt{दीर्घ\-मा\-यु\-र\-वा\-प्नोतु}}{\lemma{\skt{दीर्घम्}}\Afootnote{H; \skt{दीर्घ}\ A}}}{\lemma{\skt{आयुर्}}\Afootnote{H; \skt{मायु}\ A}}}{\lemma{\skt{अवाप्नोतु}}\Afootnote{H; \skt{रवाप्नोतु}\ A}}\skt{ देवाः प्राश्यामृतं यथा~॥}\pend

\pstart \ledsidenote{\passref{3.10.31}}\edtext{\skt{अथास्याः}}{\lemma{\skt{अथास्याः}}\Afootnote{H; \skt{अथा}\ A}}\skt{ }\edtext{\skt{क्षीर\-स\-ञ्ज\-न\-नार्थं}}{\lemma{\skt{क्षीर\appbrk{}स\appbrk{}ञ्ज\appbrk{}न\appbrk{}नार्थं}}\Afootnote{H; \emph{om.}\ A}}\skt{ }\edtext{\skt{यव\-गो\-धू\-म\-शा\-लि\-ष\-ष्टि\-क\-मां\-स\-क्षी\-र\-सु\-रा\-सौ\-वी\-र\-पि\-ण्याक}}{\lemma{\skt{यव\appbrk{}गो\appbrk{}धू\appbrk{}म\appbrk{}शा\appbrk{}लि\appbrk{}ष\appbrk{}ष्टि\appbrk{}क\appbrk{}मां\appbrk{}स\appbrk{}क्षी\appbrk{}र\appbrk{}सु\appbrk{}रा\appbrk{}सौ\appbrk{}वी\appbrk{}र\appbrk{}पि\appbrk{}ण्याक}}\Afootnote{H; \emph{om.}\ A}}\skt{ }\edtext{\skt{रसो\-न\-म\-त्स्य\-क\-शे\-रु\-का\-शृ\-ङ्गा\-ट\-क\-वि\-दा\-री\-क\-न्द\-श\-ता\-व\-री\-वि\-श\-मृ\-णा\-ल\-ना\-लि\-का\-ला\-वु\-का\-ल\-शा\-क\-प्र\-भृ\-तीनि}}{\lemma{\skt{रसो\appbrk{}न\appbrk{}म\appbrk{}त्स्य\appbrk{}क\appbrk{}शे\appbrk{}रु\appbrk{}का\appbrk{}शृ\appbrk{}ङ्गा\appbrk{}ट\appbrk{}क\appbrk{}वि\appbrk{}दा\appbrk{}री\appbrk{}क\appbrk{}न्द\appbrk{}श\appbrk{}ता\appbrk{}व\appbrk{}री\appbrk{}वि\appbrk{}श\appbrk{}मृ\appbrk{}णा\appbrk{}ल\appbrk{}ना\appbrk{}लि\appbrk{}का\appbrk{}ला\appbrk{}वु\appbrk{}का\appbrk{}ल\appbrk{}शा\appbrk{}क\appbrk{}प्र\appbrk{}भृ\appbrk{}तीनि}}\Afootnote{H; \emph{om.}\ A}}\skt{ }\edtext{\skt{विदध्यात्~। अथ}}{\lemma{\skt{विदध्यात्~। अथ}}\Afootnote{H; \emph{om.}\ A}}\skt{ }\edtext{\skt{तस्याः}}{\lemma{\skt{तस्याः}}\Afootnote{H; \skt{स्याः}\ A}}\skt{ स्तन्यमप्सु परीक्षेत }\edtext{\skt{तच्चे\-च्छी\-त\-ल\-म\-मलं}}{\lemma{\skt{तच्}}\Afootnote{A; \skt{तच्}\ H}}\skt{ }\edtext{\skt{तनुं}}{\lemma{\skt{तनुं}}\Afootnote{H; \skt{तनु}\ A}}\skt{ शंखावभासं अप्सु }\edtext{\skt{चैकीभावं}}{\lemma{\skt{चैकीभावं}}\Afootnote{H; \skt{न्यस्तमेकीभावं}\ A}}\skt{ }\edtext{\edtext{\skt{गच्छत्यफेनिलं}}{\lemma{\skt{गच्छत्य्}}\Afootnote{A; \skt{गच्छत्}\ H}}}{\lemma{\skt{अफेनिलं}}\Afootnote{A; \skt{यफेनिल}\ H}}\skt{ }\edtext{\skt{मत्रण्डुलं न}}{\lemma{\skt{मत्रण्डुलं न}}\Afootnote{H; \skt{अतन्तुमन्नोत्}\ A}}\skt{ प्लवते }\edtext{\skt{सीदति वा}}{\lemma{\skt{सीदति वा}}\Afootnote{H; \skt{ऽवसीदति वा}\ A}}\skt{ तच्छुद्धमिति }\edtext{\skt{विन्द्यात्}}{\lemma{\skt{विन्द्यात्}}\Afootnote{H; \skt{विद्यात्}\ A}}\skt{~। तेन कुमारस्यारोग्यं }\edtext{\skt{भवति}}{\lemma{\skt{भवति}}\Afootnote{H; \skt{शरीरोपचयो बलवृद्धिश्च भवति~। न च क्षुधि\appbrk{}त\appbrk{}शो\appbrk{}का\appbrk{}र्त\appbrk{}श्रा\appbrk{}न्त\appbrk{}प्र\appbrk{}दु\appbrk{}ष्ट\appbrk{}धा\appbrk{}तु\appbrk{}ग\appbrk{}र्भि\appbrk{}णी\appbrk{}ज्व\appbrk{}रि\appbrk{}ता\appbrk{}ति\appbrk{}क्षी\appbrk{}णा\appbrk{}ति\appbrk{}स्थू\appbrk{}ल\appbrk{}वि\appbrk{}द\appbrk{}ग्ध\appbrk{}भ\appbrk{}क्त\appbrk{}वि\appbrk{}रु\appbrk{}द्धा\appbrk{}हा\appbrk{}र\appbrk{}त\appbrk{}र्पि\appbrk{}तायाः स्तन्यं पाययेत्; नाजीर्णौषधं च बालं, दोषौ\appbrk{}ष\appbrk{}ध\appbrk{}म\appbrk{}लानां तीव्र\appbrk{}वे\appbrk{}गो\appbrk{}त्प\appbrk{}त्ति\appbrk{}भ\appbrk{}यात्}\ A}}\skt{~॥}\pend

\pstart\versepar \ledsidenote{\passref{3.10.32}}\edtext{\skt{भवति}}{\lemma{\skt{भवति}}\Afootnote{H; \skt{भवन्ति}\ A}}\skt{ चात्र~॥ }\pend
\pstart\versepar \skt{धात्र्यास्तु }\edtext{\edtext{\skt{गुरु\-भि\-र्भो\-ज्यै\-र्व्वि\-ष\-लै\-स्तथा}}{\lemma{\skt{भोज्यैर्}}\Afootnote{H; \skt{भोज्यै}\ A}}}{\lemma{\skt{व्विषलैस्}}\Afootnote{H; \skt{र्विष\appbrk{}मै\appbrk{}र्दो\appbrk{}ष\appbrk{}लैस्}\ A}}\skt{~। }\pend
\pstart\versepar \skt{दोषा देहे प्रकुप्यन्ति }\edtext{\skt{तत}}{\lemma{\skt{तत}}\Afootnote{H; \skt{ततः}\ A}}\skt{ स्तन्यम्प्रदुष्यति~॥}\pend

\pstart \ledsidenote{\passref{3.10.37}}\edtext{\skt{दुष्टं }\edtext{\skt{प्रागुक्तन्तेन}}{\lemma{\skt{तेन}}\Afootnote{H; \emph{om.}\ A}}}{\lemma{\skt{दुष्टं प्रागुक्तन्}}\Afootnote{H; \emph{om.}\ A}}\skt{ }\edtext{\skt{शरीराव्याधयः}}{\lemma{\skt{शरीराव्याधयः}}\Afootnote{H; \emph{om.}\ A}}\skt{ }\edtext{\skt{प्रादुर्भवन्ति}}{\lemma{\skt{प्रादुर्भवन्ति}}\Afootnote{\emph{om.}\ A; \skt{प्रादुर्ब्भवन्ति}\ H}}\skt{~। }\edtext{\skt{तेषु च}}{\lemma{\skt{तेषु च}}\Afootnote{H; \skt{तेषु}\ A}}\skt{ यथाभिहितं मृद्व\-च्छे\-द\-नी\-य\-मौ\-ष\-ध\-म्मा\-त्रया }\edtext{\skt{क्षीरस्य}}{\lemma{\skt{क्षीरस्य}}\Afootnote{H; \skt{क्षीरपस्य}\ A}}\skt{ }\edtext{\skt{क्षीरसर्प्पिः}}{\lemma{\skt{क्षीरसर्प्पिः}}\Afootnote{H; \skt{क्षीरसर्पिषा}\ A}}\skt{ संयुक्तं विदध्यात्~। धात्र्याश्च केवलं~। }\edtext{\skt{क्षीरान्नादस्य}}{\lemma{\skt{क्षीरान्नादस्य}}\Afootnote{H; \skt{क्षीरान्नादस्या}\ A}}\skt{ }\edtext{\skt{कल्की}}{\lemma{\skt{कल्की}}\Afootnote{H; \skt{त्म}\ A}}\skt{ }\edtext{\skt{कृतं}}{\lemma{\skt{कृतं}}\Afootnote{H; \skt{नि}\ A}}\skt{~। धात्र्याश्च पूर्व्ववत्~। अन्नादस्य }\edtext{\edtext{\edtext{\skt{कषा\-या\-दी\-ना\-मा\-त्म\-न्येव}}{\lemma{\skt{कषायादीनाम्}}\Afootnote{H; \skt{कषायादीना}\ A}}}{\lemma{\skt{आत्मन्य्}}\Afootnote{\skt{त्मन्य्}\ A; \skt{आत्मन्}\ H}}\skt{ न}}{\lemma{\skt{एव न}}\Afootnote{A; \skt{येव न}\ H}}\skt{ धात्र्याः~॥}\pend

\pstart\versepar \ledsidenote{\passref{3.10.44}}\edtext{\skt{भवति}}{\lemma{\skt{भवति}}\Afootnote{H; \skt{मारुतघ्नैः प्रशम}\ A}}\skt{ }\edtext{\skt{चात्र॥}}{\lemma{\skt{चात्र}}\Afootnote{H; \skt{येत्स्नेहस्वेदोप}\ A}}\skt{ }\pend
\pstart\versepar \skt{}\edtext{\skt{बालानां}}{\lemma{\skt{बालानां}}\Afootnote{H; \skt{नाहनैः}\ A}}\skt{ }\edtext{\skt{गुदपाकेषु}}{\lemma{\skt{गुदपाकेषु}}\Afootnote{H; \skt{गुदपाके तु}\ A}}\skt{ }\edtext{\skt{पित्तघ्नीं}}{\lemma{\skt{पित्तघ्नीं}}\Afootnote{H; \skt{बालानां पित्तघ्नीं}\ A}}\skt{ कारयेत्क्रियां~। }\pend
\pstart\versepar \skt{रसा\-ञ्ज\-न\-म्वि\-शे\-षेण }\edtext{\edtext{\skt{पाना\-ले\-प\-न\-यो\-र्हितं}}{\lemma{\skt{पाना\appbrk{}ले\appbrk{}प\appbrk{}न\appbrk{}योर्}}\Afootnote{A; \skt{पानालेपनयो}\ H}}}{\lemma{\skt{हितं}}\Afootnote{A; \skt{र्हितं}\ H}}\skt{~॥}\pend

\pstart \ledsidenote{\passref{3.10.45}}\edtext{\skt{क्षीराहारस्य सिद्धा\-र्थ\-क\-व\-चा\-प\-य\-स्या\-शा\-रि\-वा\-व्र\-ह्मी\-कु\-ष्ठ\-सै\-न्ध\-व\-सिद्धं}}{\lemma{\skt{क्षीराहारस्य सिद्धा\appbrk{}र्थ\appbrk{}क\appbrk{}व\appbrk{}चा\appbrk{}प\appbrk{}य\appbrk{}स्या\appbrk{}शा\appbrk{}रि\appbrk{}वा\appbrk{}व्र\appbrk{}ह्मी\appbrk{}कु\appbrk{}ष्ठ\appbrk{}सै\appbrk{}न्ध\appbrk{}व\appbrk{}सिद्धं}}\Afootnote{H; \skt{क्षीराहाराय सर्पिः पाय\appbrk{}ये\appbrk{}त्सि\appbrk{}द्धा\appbrk{}र्थ\appbrk{}क\appbrk{}व\appbrk{}चा\appbrk{}मां\appbrk{}सी\appbrk{}प\appbrk{}य\appbrk{}स्या\appbrk{}पा\appbrk{}मा\appbrk{}र्ग\appbrk{}श\appbrk{}ता\appbrk{}व\appbrk{}री\appbrk{}सा\appbrk{}रि\appbrk{}वा\appbrk{}ब्रा\appbrk{}ह्मी\appbrk{}पि\appbrk{}प्प\appbrk{}ली\appbrk{}ह\appbrk{}रि\appbrk{}द्रा\appbrk{}कु\appbrk{}ष्ठ\appbrk{}सै\appbrk{}न्ध\appbrk{}व\appbrk{}सिद्धं}\ A}}\skt{ }\edtext{\skt{सर्प्पिः}}{\lemma{\skt{सर्प्पिः}}\Afootnote{H; \emph{om.}\ A}}\skt{ }\edtext{\skt{पातुं}}{\lemma{\skt{पातुं}}\Afootnote{H; \emph{om.}\ A}}\skt{ }\edtext{\skt{प्रयच्छेत्}}{\lemma{\skt{प्रयच्छेत्}}\Afootnote{H; \emph{om.}\ A}}\skt{~। }\edtext{\skt{क्षीरान्नादस्य}}{\lemma{\skt{क्षीरान्नादस्य}}\Afootnote{H; \skt{क्षीरान्नादाय}\ A}}\skt{ }\edtext{\skt{मधु\-क\-व\-चा\-त्रि\-फ\-ला\-सिद्धं}}{\lemma{\skt{मधु\appbrk{}क\appbrk{}व\appbrk{}चा\appbrk{}त्रि\appbrk{}फ\appbrk{}ला\appbrk{}सिद्धं}}\Afootnote{H; \skt{मधु\appbrk{}क\appbrk{}व\appbrk{}चा\appbrk{}पि\appbrk{}प्प\appbrk{}ली\appbrk{}चि\appbrk{}त्र\appbrk{}क\appbrk{}त्रि\appbrk{}फ\appbrk{}ला\appbrk{}सि\appbrk{}द्धम्}\ A}}\skt{~। }\edtext{\skt{अन्नादस्य}}{\lemma{\skt{अन्नादस्य}}\Afootnote{H; \skt{अन्नादाय}\ A}}\skt{ }\edtext{\skt{द्विप\-ञ्च\-मू\-ली\-क्षी\-र\-त\-ग\-र\-भ\-द्र\-दा\-रु\-वि\-डङ्ग-}}{\lemma{\skt{द्विप\appbrk{}ञ्च\appbrk{}मू\appbrk{}ली\appbrk{}क्षी\appbrk{}र\appbrk{}त\appbrk{}ग\appbrk{}र\appbrk{}भ\appbrk{}द्र\appbrk{}दा\appbrk{}रु\appbrk{}वि\appbrk{}डङ्ग}}\Afootnote{H; \skt{द्विप\appbrk{}ञ्च\appbrk{}मू\appbrk{}ली\appbrk{}क्षी\appbrk{}र\appbrk{}त\appbrk{}ग\appbrk{}र\appbrk{}भ\appbrk{}द्र\appbrk{}दा\appbrk{}रु\appbrk{}म\appbrk{}रि\appbrk{}च\appbrk{}म\appbrk{}धु\appbrk{}क\appbrk{}वि\appbrk{}डङ्ग}\ A}}\skt{ }\edtext{\skt{द्राक्षादिसिद्धं}}{\lemma{\skt{द्राक्षादिसिद्धं}}\Afootnote{H; \skt{द्राक्ष\appbrk{}र्धि\appbrk{}ब्रा\appbrk{}ह्मी\appbrk{}सिद्धं}\ A}}\skt{ }\edtext{\skt{तेन}}{\lemma{\skt{तेन}}\Afootnote{H; \skt{ते}\ A}}\skt{ }\edtext{\skt{कुमा\-र\-स्या\-रो\-ग्य\-व\-ल\-मे\-धा\-यूंषि}}{\lemma{\skt{कुमा\appbrk{}र\appbrk{}स्या\appbrk{}रो\appbrk{}ग्य\appbrk{}व\appbrk{}ल\appbrk{}मे\appbrk{}धा\appbrk{}यूंषि}}\Afootnote{H; \skt{नारो\appbrk{}ग्य\appbrk{}ब\appbrk{}ल\appbrk{}मे\appbrk{}धा\appbrk{}यूंषि}\ A}}\skt{ }\edtext{\skt{भवन्ति}}{\lemma{\skt{भवन्ति}}\Afootnote{H; \skt{शिशोर्भवन्ति}\ A}}\skt{~।}\pend

\pstart \ledsidenote{\passref{3.10.46}}\skt{बालं }\edtext{\edtext{\skt{पुनर्ग्गात्रेषु}}{\lemma{\skt{पुनर्}}\Afootnote{H; \skt{पुन}\ A}}\skt{ न}}{\lemma{\skt{ग्गात्रेषु न}}\Afootnote{H; \skt{र्गात्रसुखं}\ A}}\skt{ }\edtext{\skt{गृङ्णीयात्~। न}}{\lemma{\skt{गृङ्णीयात्~। न}}\Afootnote{H; \skt{गृह्णीयात्, न}\ A}}\skt{ चैनं }\edtext{\skt{तर्ज्जयेत्सहसा वा न}}{\lemma{\skt{सहसा वा न}}\Afootnote{H; \skt{सहसा न}\ A}}\skt{ }\edtext{\skt{प्रबो\-ध\-ये\-द्वि\-त्रा\-स\-भ\-यात्}}{\lemma{\skt{प्रबोधयेद्}}\Afootnote{H; \skt{प्रतिबोधयेद्}\ A}}\skt{~। सहसा }\edtext{\edtext{\edtext{\skt{नाह\-रे\-दु\-त्क्षि\-पे\-द्वा\-ता\-भि\-घा\-त\-भ\-या\-न्नो\-प\-वे\-श\-ये\-त्कौ\-व्ज\-भ\-या\-न्नित्यं}}{\lemma{\skt{नाह\appbrk{}रे\appbrk{}दु\appbrk{}त्क्षि\appbrk{}पेद्}}\Afootnote{H; \skt{नाप\appbrk{}ह\appbrk{}रे\appbrk{}दु\appbrk{}त्क्षि\appbrk{}पेद्वा}\ A}}}{\lemma{\skt{वाता\appbrk{}भि\appbrk{}घा\appbrk{}त\appbrk{}भ\appbrk{}यान्}}\Afootnote{H; \skt{वाता\appbrk{}दि\appbrk{}वि\appbrk{}घा\appbrk{}त\appbrk{}भ\appbrk{}यात्}\ A}}}{\lemma{\skt{कौव्जभयान्}}\Afootnote{H; \skt{कौब्ज्यभयात्}\ A}}\skt{ चैन\-म\-नु\-व\-र्त्तेत }\edtext{\skt{प्रियहितैरेव}}{\lemma{\skt{प्रियहितैर्}}\Afootnote{H; \skt{प्रिय\appbrk{}श\appbrk{}तै\appbrk{}र\appbrk{}जि\appbrk{}घांसुः}\ A}}\skt{ }\edtext{\skt{मभि\-हि\-त\-म\-ना\-द्य\-म\-प्य\-भि\-व\-र्द्धते}}{\lemma{\skt{अभि\appbrk{}हि\appbrk{}त\appbrk{}म\appbrk{}ना\appbrk{}द्य\appbrk{}म\appbrk{}प्य\appbrk{}भि\appbrk{}व\appbrk{}र्द्धते}}\Afootnote{H; \skt{अविहतमना ह्यभिवर्धते}\ A}}\skt{ }\edtext{\skt{नित्य\-मु\-द\-ग्र\-स\-त्त्व\-स\-म्प\-न्न\-स्त्व\-रोगो}}{\lemma{\skt{नित्य\appbrk{}मु\appbrk{}द\appbrk{}ग्र\appbrk{}स\appbrk{}त्त्व\appbrk{}स\appbrk{}म्प\appbrk{}न्न\appbrk{}स्त्व\appbrk{}रोगो}}\Afootnote{H; \skt{नित्य\appbrk{}मु\appbrk{}द\appbrk{}ग्र\appbrk{}स\appbrk{}त्त्व\appbrk{}सं\appbrk{}पन्नो नीरोगः}\ A}}\skt{ }\edtext{\skt{भवति}}{\lemma{\skt{भवति}}\Afootnote{H; \skt{सुप्र\appbrk{}स\appbrk{}न्न\appbrk{}म\appbrk{}नाश्च भवति~। वाता\appbrk{}त\appbrk{}प\appbrk{}वि\appbrk{}द्यु\appbrk{}त्प्र\appbrk{}भा\appbrk{}पा\appbrk{}द\appbrk{}प\appbrk{}ल\appbrk{}ता\appbrk{}शू\appbrk{}न्या\appbrk{}गा\appbrk{}र\appbrk{}नि\appbrk{}म्न\appbrk{}स्था\appbrk{}न\appbrk{}ग्र\appbrk{}ह\appbrk{}च्छा\appbrk{}या\appbrk{}दिभ्यो दुर्ग्र\appbrk{}हो\appbrk{}प\appbrk{}स\appbrk{}र्ग\appbrk{}तश्च बालं रक्षेत्}\ A}}\skt{~॥}\pend

\pstart\versepar \ledsidenote{\passref{3.10.47}}\edtext{\skt{भवति}}{\lemma{\skt{भवति}}\Afootnote{H; \emph{om.}\ A}}\skt{ }\edtext{\skt{चात्र}}{\lemma{\skt{चात्र}}\Afootnote{H; \emph{om.}\ A}}\skt{~॥ }\pend
\pstart\versepar \skt{नाशुचौ विसृजेद्बालं }\edtext{\skt{नाकासे}}{\lemma{\skt{नाकासे}}\Afootnote{H; \skt{नाकाशे}\ A}}\skt{ }\edtext{\skt{विषमेण च}}{\lemma{\skt{विषमेण च}}\Afootnote{H; \skt{विषमे न च}\ A}}\skt{~। }\pend
\pstart\versepar \skt{}\edtext{\skt{नोष्ण\-मा\-रु\-त\-व\-र्षेषु}}{\lemma{\skt{नोष्ण\appbrk{}मा\appbrk{}रु\appbrk{}त\appbrk{}व\appbrk{}र्षेषु}}\Afootnote{H; \skt{नोष्म\appbrk{}मा\appbrk{}रु\appbrk{}त\appbrk{}व\appbrk{}र्षेषु}\ A}}\skt{ रजो\-धू\-मो\-द\-केषु }\edtext{\skt{वा}}{\lemma{\skt{वा}}\Afootnote{H; \skt{च}\ A}}\skt{~॥}\pend

\pstart\versepar \ledsidenote{\passref{3.10.48}}\skt{क्षीर\-सा\-त्म्य\-त\-या\-क्षी\-र\-माजं }\edtext{\edtext{\skt{गव्यमथापि}}{\lemma{\skt{गव्यम्}}\Afootnote{H; \skt{गव्य}\ A}}\skt{ वा}}{\lemma{\skt{अथापि वा}}\Afootnote{H; \skt{मथापि वा}\ A}}\skt{~। }\pend
\pstart\versepar \skt{}\edtext{\edtext{\skt{दद्या\-द\-स्त\-न्य\-म\-प्राप्तौ}}{\lemma{\skt{दद्यादस्तन्यम्}}\Afootnote{H; \skt{दद्यादास्तन्य}\ A}}}{\lemma{\skt{अप्राप्तौ}}\Afootnote{H; \skt{पर्याप्तेर्}\ A}}\skt{ }\edtext{\skt{बालेभ्यो}}{\lemma{\skt{बालेभ्यो}}\Afootnote{H; \skt{बालानां}\ A}}\skt{ वीक्ष्य मात्रया~॥}\pend

\pstart \ledsidenote{\passref{3.10.49}}\edtext{\edtext{\edtext{\edtext{\skt{षण्मा\-सा\-च्चै\-न\-म\-न्न\-म्प्रा\-श\-ये\-ल्ल\-घु\-हितं}}{\lemma{\skt{षण्मासाच्}}\Afootnote{H; \skt{षण्मासं}\ A}}}{\lemma{\skt{चैनम्}}\Afootnote{H; \skt{चैन}\ A}}}{\lemma{\skt{अन्नम्}}\Afootnote{H; \skt{मन्नं}\ A}}}{\lemma{\skt{लघुहितं}}\Afootnote{H; \skt{लघु हितं च}\ A}}\pend

\pstart \ledsidenote{\passref{3.10.50}}\edtext{\skt{नित्य\-म\-व\-रो\-धश्च}}{\lemma{\skt{अवरोधश्च}}\Afootnote{H; \skt{अव\appbrk{}रो\appbrk{}ध\appbrk{}र\appbrk{}तश्च}\ A}}\skt{~। स्यात्कृतरक्षः~। }\edtext{\skt{उपसर्ग्गभया}}{\lemma{\skt{उपसर्ग्गभया}}\Afootnote{H; \skt{उप\appbrk{}स\appbrk{}र्ग\appbrk{}भ\appbrk{}यात्}\ A}}\skt{ }\edtext{\skt{ग्रहोपसर्ग्गेभ्यो}}{\lemma{\skt{ग्रहोपसर्ग्गेभ्यो}}\Afootnote{H; \skt{प्रयत्नतश्च ग्रहोपसर्गेभ्यो}\ A}}\skt{ रक्ष्या बाला भवन्ति~॥}\pend

\pstart \ledsidenote{\passref{3.10.51}}\skt{अथ कुमार उद्विजते त्रस्यति रोदिति नष्टसंज्ञो भवति~॥ }\edtext{\edtext{\skt{नख\-द\-श\-नै\-र्द्धा\-त्री\-मा\-त्मा\-नञ्च}}{\lemma{\skt{द्धात्रीम्}}\Afootnote{H; \skt{धात्री}\ A}}}{\lemma{\skt{आत्मानञ्च}}\Afootnote{H; \skt{मात्मानं च}\ A}}\skt{ }\edtext{\skt{परितुदति}}{\lemma{\skt{परितुदति}}\Afootnote{H; \skt{परिणुदति}\ A}}\skt{ दन्तान्खादति कूजति जृम्भते }\edtext{\edtext{\skt{भ्रूवा\-वु\-त्क्षि\-प\-त्यू\-र्द्ध्व\-न्नि\-री\-क्षते}}{\lemma{\skt{भ्रूवावुत्क्षिपत्य्}}\Afootnote{H; \skt{भ्रुवौ विक्षिपत्}\ A}}}{\lemma{\skt{ऊर्द्ध्वन्}}\Afootnote{H; \skt{यूर्ध्वं}\ A}}\skt{ }\edtext{\edtext{\skt{फेनमुद्वमति}}{\lemma{\skt{फेनम्}}\Afootnote{H; \skt{फेन}\ A}}}{\lemma{\skt{उद्वमति}}\Afootnote{H; \skt{मुद्वमति}\ A}}\skt{ सन्दष्टौष्ठः }\edtext{\skt{क्रूर\-भि\-न्ना\-र्त्त\-स्वरो}}{\lemma{\skt{क्रूर\appbrk{}भि\appbrk{}न्ना\appbrk{}र्त्त\appbrk{}स्वरो}}\Afootnote{H; \skt{क्रूरो भिन्नामवर्चा दीनार्तस्वरो}\ A}}\skt{ निशि जागर्त्ति }\edtext{\skt{दुर्व्वलोम्नाङ्गो}}{\lemma{\skt{दुर्व्वलोम्नाङ्गो}}\Afootnote{H; \skt{दुर्बलो म्लानाङ्गो मत्स्य\appbrk{}च्छु\appbrk{}च्छ्रु\appbrk{}न्द\appbrk{}रि\appbrk{}म\appbrk{}त्कु\appbrk{}ण\appbrk{}गन्धो}\ A}}\skt{ }\edtext{\skt{यथापूर्व्वं}}{\lemma{\skt{यथापूर्व्वं}}\Afootnote{H; \skt{यथा पुरा}\ A}}\skt{ धात्र्याः }\edtext{\edtext{\skt{स्तन्य\-न्ना\-भि\-ल\-ष\-तीति}}{\lemma{\skt{स्तन्यन्}}\Afootnote{H; \skt{स्तन्य}\ A}}}{\lemma{\skt{नाभिलषतीति}}\Afootnote{H; \skt{मभिलषति तथा नाभिलषतीति}\ A}}\skt{ }\edtext{\skt{सामान्यं}}{\lemma{\skt{सामान्यं}}\Afootnote{H; \skt{सामान्येन}\ A}}\skt{ }\edtext{\skt{ग्रहो\-प\-सृ\-ष्ट\-ल\-क्ष\-णुक्तं}}{\lemma{\skt{ग्रहो\appbrk{}प\appbrk{}सृ\appbrk{}ष्ट\appbrk{}ल\appbrk{}क्ष\appbrk{}णुक्तं}}\Afootnote{H; \skt{ग्रहो\appbrk{}प\appbrk{}सृ\appbrk{}ष्ट\appbrk{}ल\appbrk{}क्ष\appbrk{}ण\appbrk{}मुक्तं}\ A}}\skt{~। विस्तरेणोत्त रे वक्ष्यामः~॥}\pend

\pstart \ledsidenote{\passref{3.10.52}}\skt{शक्तिमन्तं चैनं ज्ञात्वा यथा }\edtext{\skt{स्वम्विद्यां}}{\lemma{\skt{स्वम्}}\Afootnote{H; \skt{वर्णं}\ A}}\skt{ ग्राहयेत्~।}\pend

\pstart \ledsidenote{\passref{3.10.53}}\skt{अथास्मै पञ्च\-विं\-श\-ति\-व\-र्षाय }\edtext{\edtext{\edtext{\skt{षोड\-श\-व\-र्षा\-म्प\-त्नी\-मा\-व\-हे\-न्मि\-त्र\-ध\-र्म्मा\-र्थ\-का\-म\-प्रजाः}}{\lemma{\skt{षोडशवर्षाम्}}\Afootnote{H; \skt{षोड\appbrk{}श\appbrk{}षो\appbrk{}ड\appbrk{}श\appbrk{}द\appbrk{}श\appbrk{}वर्षां}\ A}}}{\lemma{\skt{आवहेन्}}\Afootnote{H; \skt{आवहेत्}\ A}}}{\lemma{\skt{मित्र\appbrk{}ध\appbrk{}र्म्मा\appbrk{}र्थ\appbrk{}का\appbrk{}म\appbrk{}प्रजाः}}\Afootnote{H; \skt{पित्र्य\appbrk{}ध\appbrk{}र्मा\appbrk{}र्थ\appbrk{}का\appbrk{}म\appbrk{}प्रजाः}\ A}}\skt{ }\edtext{\skt{प्राप्स्यामीति}}{\lemma{\skt{प्राप्स्यामीति}}\Afootnote{H; \skt{प्राप्स्यतीति}\ A}}\skt{~॥}\pend

\pstart\versepar \ledsidenote{\passref{3.10.54}}\edtext{\skt{भवति}}{\lemma{\skt{भवति}}\Afootnote{H; \emph{om.}\ A}}\skt{ }\edtext{\skt{चात्र}}{\lemma{\skt{चात्र}}\Afootnote{H; \emph{om.}\ A}}\skt{~॥ }\pend
\pstart\versepar \skt{ऊन\-षो\-ड\-श\-व\-र्षा\-या\-म\-प्राप्तः }\edtext{\skt{पञ्चविंशतिः}}{\lemma{\skt{पञ्चविंशतिः}}\Afootnote{H; \skt{पञ्चविंशतिम्}\ A}}\skt{~। }\pend
\pstart\versepar \skt{यद्याधत्ते }\edtext{\skt{पुमान्गर्भः}}{\lemma{\skt{गर्भः}}\Afootnote{\skt{गर्भं}\ A; \skt{गर्ब्भः}\ H}}\skt{ कुक्षिस्थः स विपद्यते~।}\pend

\pstart \ledsidenote{\passref{3.10.55}}\edtext{\skt{जातो वा न}}{\lemma{\skt{जातो वा न}}\Afootnote{A; \skt{जातोवा नि}\ H}}\skt{ चिरं }\edtext{\skt{जीवेद्वा}}{\lemma{\skt{जीवेद्वा}}\Afootnote{H; \skt{जीवेज्जीवेद्वा}\ A}}\skt{ }\edtext{\skt{दुर्बलेन्द्रियः}}{\lemma{\skt{दुर्बलेन्द्रियः}}\Afootnote{H; \skt{दुर्बलेन्द्रियः~। तस्मा\appbrk{}द\appbrk{}त्य\appbrk{}न्त\appbrk{}बा\appbrk{}लायां गर्भाधानं न कारयेत्}\ A}}\skt{~।}\pend

\pstart \ledsidenote{\passref{3.10.57.1}}\skt{तत्र पूर्वोक्तैः कारणैः }\edtext{\skt{प्रपतिष्यति}}{\lemma{\skt{प्रपतिष्यति}}\Afootnote{H; \skt{पतिष्यति}\ A}}\skt{ }\edtext{\skt{गर्भः}}{\lemma{\skt{गर्भः}}\Afootnote{\skt{गर्भे}\ A; \skt{गर्ब्भ}\ H}}\skt{~। }\edtext{\skt{गर्ब्भा\-श\-य\-श्रो\-णी\-वं\-क्ष\-ण\-ब\-स्ति\-शू\-लानि}}{\lemma{\skt{गर्ब्भा\appbrk{}श\appbrk{}य\appbrk{}श्रो\appbrk{}णी\appbrk{}वं\appbrk{}क्ष\appbrk{}ण\appbrk{}ब\appbrk{}स्ति\appbrk{}शू\appbrk{}लानि}}\Afootnote{H; \skt{गर्भा\appbrk{}श\appbrk{}य\appbrk{}क\appbrk{}टी\appbrk{}व\appbrk{}ङ्क्ष\appbrk{}ण\appbrk{}ब\appbrk{}स्ति\appbrk{}शू\appbrk{}लानि}\ A}}\skt{ }\edtext{\skt{भवन्ति}}{\lemma{\skt{भवन्ति}}\Afootnote{H; \emph{om.}\ A}}\skt{~। रक्तदर्शनञ्च~। }\edtext{\edtext{\edtext{\skt{तासा\-म्प\-रि\-षे\-का\-व\-गा\-हा\-ले\-प\-नै\-रु\-प\-च\-रे\-त्क्षी\-र\-पा\-नैश्च}}{\lemma{\skt{तासाम्}}\Afootnote{H; \skt{तत्र शीतैः}\ A}}}{\lemma{\skt{परि\appbrk{}षे\appbrk{}का\appbrk{}व\appbrk{}गा\appbrk{}हा\appbrk{}ले\appbrk{}प\appbrk{}नैर्}}\Afootnote{H; \skt{परि\appbrk{}षे\appbrk{}का\appbrk{}व\appbrk{}गा\appbrk{}ह\appbrk{}प्र\appbrk{}दे\appbrk{}हा\appbrk{}दि\appbrk{}भिर्}\ A}}}{\lemma{\skt{उपचरेत्}}\Afootnote{H; \skt{उप\appbrk{}च\appbrk{}रे\appbrk{}ज्जी\appbrk{}व\appbrk{}नी\appbrk{}य\appbrk{}शृत}\ A}}\skt{~। }\edtext{\skt{अल्पाल्पस्यन्दते}}{\lemma{\skt{अल्पाल्पस्यन्दते}}\Afootnote{H; \skt{गर्भस्फुरणे मुहुर्मुहुस्त}\ A}}\skt{ }\edtext{\skt{क्षीरमुत्पलादि}}{\lemma{\skt{क्षीरमुत्पलादि}}\Afootnote{H; \skt{त्सन्धारणार्थं क्षीरमुत्पलादि}\ A}}\skt{ }\edtext{\skt{सिद्धम्वा}}{\lemma{\skt{सिद्धम्वा}}\Afootnote{H; \skt{सिद्धं}\ A}}\skt{ पाययेत्~। }\edtext{\skt{संस्रमाने}}{\lemma{\skt{संस्रमाने}}\Afootnote{H; \skt{प्रस्रंसमाने}\ A}}\skt{ }\edtext{\skt{सवा\-त\-पा\-र्श्व\-शू\-ला\-नाह}}{\lemma{\skt{सवा\appbrk{}त\appbrk{}पा\appbrk{}र्श्व\appbrk{}शू\appbrk{}ला\appbrk{}नाह}}\Afootnote{H; \skt{सदा\appbrk{}ह\appbrk{}पा\appbrk{}र्श्व\appbrk{}पृ\appbrk{}ष्ठ\appbrk{}शू\appbrk{}ला\appbrk{}सृ\appbrk{}ग्द\appbrk{}रा\appbrk{}नाह}\ A}}\skt{ मूत्रसङ्गाः स्थानात्स्थानं }\edtext{\skt{चोत्क्रामति}}{\lemma{\skt{चोत्क्रामति}}\Afootnote{H; \skt{च प्रक्रामति}\ A}}\skt{ }\edtext{\skt{गर्भः}}{\lemma{\skt{गर्भः}}\Afootnote{\skt{गर्भे कोष्ठे संरम्भः}\ A; \skt{गर्ब्भः}\ H}}\skt{~। तत्र }\edtext{\skt{स्निग्धं}}{\lemma{\skt{स्निग्धं}}\Afootnote{H; \skt{स्निग्ध}\ A}}\skt{ शीताः क्रियाः~। वेदनायां }\edtext{\skt{महा\-स\-ह\-म\-धु\-क\-स्व\-दं\-ष्ट्र\-क\-ण्ट\-का\-रि\-का\-सिद्धं}}{\lemma{\skt{महा\appbrk{}स\appbrk{}ह\appbrk{}म\appbrk{}धु\appbrk{}क\appbrk{}स्व\appbrk{}दं\appbrk{}ष्ट्र\appbrk{}क\appbrk{}ण्ट\appbrk{}का\appbrk{}रि\appbrk{}का\appbrk{}सिद्धं}}\Afootnote{\skt{महा\appbrk{}स\appbrk{}हा\appbrk{}क्षु\appbrk{}द्र\appbrk{}स\appbrk{}हा\appbrk{}म\appbrk{}धु\appbrk{}क\appbrk{}श्व\appbrk{}दं\appbrk{}ष्ट्रा\appbrk{}क\appbrk{}ण्ट\appbrk{}का\appbrk{}रि\appbrk{}का\appbrk{}सिद्धं}\ A; \skt{महा\appbrk{}स\appbrk{}ह\appbrk{}म\appbrk{}धु\appbrk{}क\appbrk{}स्व\appbrk{}दं\appbrk{}ष्ट्र\appbrk{}क\appbrk{}ण्ट\appbrk{}का\appbrk{}रि\appbrk{}का\appbrk{}सिद्ध}\ H}}\skt{ पयः~। }\edtext{\skt{क्षौद्र\-श\-र्क\-रा\-मि\-श्र\-म्पा\-य\-येत्}}{\lemma{\skt{क्षौद्र\appbrk{}श\appbrk{}र्क\appbrk{}रा\appbrk{}मि\appbrk{}श्रम्}}\Afootnote{H; \skt{शर्क\appbrk{}रा\appbrk{}क्षौ\appbrk{}द्र\appbrk{}मिश्रं}\ A}}\skt{~। मूत्रसङ्गे }\edtext{\skt{दर्ब्भादिसिद्धं}}{\lemma{\skt{दर्ब्भादिसिद्धं}}\Afootnote{H; \skt{दर्भादिसिद्धम्}\ A}}\skt{~। आनाहे }\edtext{\skt{हिङ्गु\-सौ\-व\-र्च्च\-ल\-रसो न}}{\lemma{\skt{हिङ्गु\appbrk{}सौ\appbrk{}व\appbrk{}र्च्च\appbrk{}ल\appbrk{}रसो न}}\Afootnote{H; \skt{हिङ्गु\appbrk{}सौ\appbrk{}व\appbrk{}र्च\appbrk{}ल\appbrk{}ल\appbrk{}शुन}\ A}}\skt{ }\edtext{\skt{सिद्धं}}{\lemma{\skt{सिद्धं}}\Afootnote{\skt{वचासिद्धम्}\ A; \emph{om.}\ H}}\skt{ }\edtext{\skt{स्रवति}}{\lemma{\skt{स्रवति}}\Afootnote{\skt{अत्यर्थं स्रवति}\ A; \skt{सिद्धं स्रवति}\ H}}\skt{ रक्ते }\edtext{\edtext{\skt{कोष्ठा\-त्क\-ण्ट\-का\-रि\-का\-मृ\-त्पि\-ण्ड\-गै\-रि\-क\-स\-र्ज्ज\-र\-सा\-ञ्ज\-न\-चूर्ण्ण}}{\lemma{\skt{कोष्ठात्}}\Afootnote{H; \skt{कोष्ठा}\ A}}}{\lemma{\skt{कण्ट\appbrk{}का\appbrk{}रि\appbrk{}का\appbrk{}मृ\appbrk{}त्पि\appbrk{}ण्ड\appbrk{}गै\appbrk{}रि\appbrk{}क\appbrk{}स\appbrk{}र्ज्ज\appbrk{}र\appbrk{}सा\appbrk{}ञ्ज\appbrk{}न\appbrk{}चूर्ण्ण}}\Afootnote{H; \skt{गारि\appbrk{}का\appbrk{}गा\appbrk{}र\appbrk{}मृ\appbrk{}त्पि\appbrk{}ण्ड\appbrk{}स\appbrk{}म\appbrk{}ङ्गा\appbrk{}धा\appbrk{}त\appbrk{}की\appbrk{}कु\appbrk{}सु\appbrk{}म\appbrk{}न\appbrk{}व\appbrk{}मा\appbrk{}लि\appbrk{}का\appbrk{}गै\appbrk{}रि\appbrk{}क\appbrk{}स\appbrk{}र्ज\appbrk{}र\appbrk{}स\appbrk{}र\appbrk{}सा\appbrk{}ञ्ज\appbrk{}न\appbrk{}चूर्णं मधुनाव}\ A}}\skt{ }\edtext{\edtext{\skt{यथा\-ला\-भं\-न्य\-ग्रो\-धा\-दि\-त्व\-क्प्र\-बा\-ल\-क\-ल्कन्वा}}{\lemma{\skt{यथा\appbrk{}ला\appbrk{}भं\appbrk{}न्य\appbrk{}ग्रो\appbrk{}धा\appbrk{}दि\appbrk{}त्व\appbrk{}क्प्र\appbrk{}बा\appbrk{}ल\appbrk{}क\appbrk{}ल्कन्}}\Afootnote{H; \skt{लिह्यात्, यथालाभं न्यग्रो\appbrk{}धा\appbrk{}दि\appbrk{}त्व\appbrk{}क्प्र\appbrk{}वा\appbrk{}ल\appbrk{}कल्कं वा}\ A}}\skt{ पाययेत्~। उत्पलादिकल्कं वा}}{\lemma{\skt{वा पाययेत्~। उत्पलादिकल्कं वा}}\Afootnote{\skt{पयसा पाययेत्, उत्पलादिकल्कं वा}\ A; \skt{वा पाययेत्~। उत्पला}\ H}}\skt{ }\edtext{\skt{कशे\-रु\-शृ\-ङ्गा\-ट\-क\-शा\-लू\-क\-कल्कं वा}}{\lemma{\skt{कशे\appbrk{}रु\appbrk{}शृ\appbrk{}ङ्गा\appbrk{}ट\appbrk{}क\appbrk{}शा\appbrk{}लू\appbrk{}क\appbrk{}कल्कं वा}}\Afootnote{A; \emph{om.}\ H}}\skt{ }\edtext{\skt{शृतेन}}{\lemma{\skt{शृतेन}}\Afootnote{A; \emph{om.}\ H}}\skt{ }\edtext{\skt{पयसा,}}{\lemma{\skt{पयसा}}\Afootnote{A; \emph{om.}\ H}}\skt{ }\edtext{\skt{उदु\-म्ब\-र\-फ\-लौ\-द\-क\-क\-न्द\-क्वा\-थेन वा}}{\lemma{\skt{उदु\appbrk{}म्ब\appbrk{}र\appbrk{}फ\appbrk{}लौ\appbrk{}द\appbrk{}क\appbrk{}क\appbrk{}न्द\appbrk{}क्वा\appbrk{}थेन वा}}\Afootnote{A; \emph{om.}\ H}}\skt{ }\edtext{\skt{शर्क\-रा\-म\-धु\-म\-धु\-रेण}}{\lemma{\skt{शर्क\appbrk{}रा\appbrk{}म\appbrk{}धु\appbrk{}म\appbrk{}धु\appbrk{}रेण}}\Afootnote{A; \emph{om.}\ H}}\skt{ }\edtext{\skt{शालिपिष्टं;}}{\lemma{\skt{शालिपिष्टं}}\Afootnote{A; \emph{om.}\ H}}\skt{ }\edtext{\skt{न्यग्रो\-धा\-दि\-स्व\-र\-स\-प\-रि\-पीतं वा}}{\lemma{\skt{न्यग्रो\appbrk{}धा\appbrk{}दि\appbrk{}स्व\appbrk{}र\appbrk{}स\appbrk{}प\appbrk{}रि\appbrk{}पीतं वा}}\Afootnote{A; \skt{दि}\ H}}\skt{ }\edtext{\skt{वस्त्रावयवं}}{\lemma{\skt{वस्त्रावयवं}}\Afootnote{A; \emph{om.}\ H}}\skt{ }\edtext{\skt{योन्यां}}{\lemma{\skt{योन्यां}}\Afootnote{A; \emph{om.}\ H}}\skt{ }\edtext{\skt{धारयेत्~। अथा\-दृ\-ष्ट\-शो\-णि\-त\-वे\-द\-नायां}}{\lemma{\skt{धारयेत्~। अथा\appbrk{}दृ\appbrk{}ष्ट\appbrk{}शो\appbrk{}णि\appbrk{}त\appbrk{}वे\appbrk{}द\appbrk{}नायां}}\Afootnote{A; \emph{om.}\ H}}\skt{ }\edtext{\skt{मधु\-क\-दे\-व\-दा\-रु\-म\-ञ्जि\-ष्ठा\-प\-य\-स्या\-सिद्धं}}{\lemma{\skt{मधु\appbrk{}क\appbrk{}दे\appbrk{}व\appbrk{}दा\appbrk{}रु\appbrk{}म\appbrk{}ञ्जि\appbrk{}ष्ठा\appbrk{}प\appbrk{}य\appbrk{}स्या\appbrk{}सिद्धं}}\Afootnote{A; \emph{om.}\ H}}\skt{ }\edtext{\skt{पयः}}{\lemma{\skt{पयः}}\Afootnote{A; \emph{om.}\ H}}\skt{ }\edtext{\skt{पाययेत्,}}{\lemma{\skt{पाययेत्}}\Afootnote{A; \emph{om.}\ H}}\skt{ }\edtext{\skt{तदे\-वा\-श्म\-न्त\-क\-श\-ता\-व\-री\-प\-य\-स्या\-सिद्धं}}{\lemma{\skt{तदे\appbrk{}वा\appbrk{}श्म\appbrk{}न्त\appbrk{}क\appbrk{}श\appbrk{}ता\appbrk{}व\appbrk{}री\appbrk{}प\appbrk{}य\appbrk{}स्या\appbrk{}सिद्धं}}\Afootnote{A; \skt{क}\ H}}\skt{ }\edtext{\skt{विदा\-रि\-ग\-न्धा\-दि\-सिद्धं वा}}{\lemma{\skt{विदा\appbrk{}रि\appbrk{}ग\appbrk{}न्धा\appbrk{}दि\appbrk{}सिद्धं वा}}\Afootnote{A; \skt{सेदि सिद्धम्वा}\ H}}\skt{~। }\edtext{\skt{बृह\-ती\-द्व\-यो\-त्प\-ल\-शा\-रि\-वा\-प\-य\-स्या\-म\-धु\-क\-सि\-द्धम्वा}}{\lemma{\skt{बृह\appbrk{}ती\appbrk{}द्व\appbrk{}यो\appbrk{}त्प\appbrk{}ल\appbrk{}शा\appbrk{}रि\appbrk{}वा\appbrk{}प\appbrk{}य\appbrk{}स्या\appbrk{}म\appbrk{}धु\appbrk{}क\appbrk{}सि\appbrk{}द्धम्वा}}\Afootnote{H; \skt{बृह\appbrk{}ती\appbrk{}द्व\appbrk{}यो\appbrk{}त्प\appbrk{}ल\appbrk{}श\appbrk{}ता\appbrk{}व\appbrk{}री\appbrk{}सा\appbrk{}रि\appbrk{}वा\appbrk{}प\appbrk{}य\appbrk{}स्या\appbrk{}म\appbrk{}धु\appbrk{}क\appbrk{}सिद्धं वा}\ A}}\skt{~। }\edtext{\skt{एते\-नो\-प\-व\-र्त्तते}}{\lemma{\skt{एते\appbrk{}नो\appbrk{}प\appbrk{}व\appbrk{}र्त्तते}}\Afootnote{H; \skt{एव\appbrk{}मु\appbrk{}प\appbrk{}क्रा\appbrk{}न्ताया उपावर्तन्ते}\ A}}\skt{ }\edtext{\skt{रुजसु}}{\lemma{\skt{रुजसु}}\Afootnote{H; \skt{रुजो}\ A}}\skt{ }\edtext{\edtext{\skt{गर्भश्चाप्याय्यते}}{\lemma{\skt{गर्भश्}}\Afootnote{A; \skt{गर्ब्भश्}\ H}}}{\lemma{\skt{आय्यते}}\Afootnote{H; \skt{आयते}\ A}}\skt{~॥}\footnoteA{Passage 3.10.57 in the vulgate is split in the Nepalese version into two parts. In witness H, the second part of this passage appears after 3.10.65. Witness H is partially illegible here and the text has been emended based on the vulgate.}\pend

\pstart \ledsidenote{\passref{3.10.58}}\edtext{\skt{मासा\-नु\-मा\-सि\-क\-मत}}{\lemma{\skt{मासा\appbrk{}नु\appbrk{}मा\appbrk{}सि\appbrk{}कम्}}\Afootnote{H; \emph{om.}\ A}}\skt{ ऊर्ध्वं }\edtext{\skt{वक्ष्यामः}}{\lemma{\skt{वक्ष्यामः}}\Afootnote{H; \skt{मासानुमासिकं वक्ष्यामः}\ A}}\pend

\pstart\versepar \ledsidenote{\passref{3.10.59}}\skt{मधुकं }\edtext{\skt{शालिबीजञ्च}}{\lemma{\skt{शालिबीजञ्च}}\Afootnote{H; \skt{शाकबीजं च}\ A}}\skt{ पयस्या सुरदारु च~। }\pend
\pstart\versepar \skt{अश्मन्तकः }\edtext{\edtext{\skt{कृष्ण\-ति\-ला\-स्ता\-म्र\-वल्ली}}{\lemma{\skt{कृष्णतिलास्}}\Afootnote{H; \skt{तिलाः कृष्णा}\ A}}}{\lemma{\skt{ताम्रवल्ली}}\Afootnote{H; \skt{स्ताम्रवल्ली}\ A}}\skt{ शतावरी~॥}\pend

\pstart\versepar \ledsidenote{\passref{3.10.60}}\skt{वृक्षादनी पयस्या च लता }\edtext{\skt{सोत्पलशारिवा}}{\lemma{\skt{सोत्पलशारिवा}}\Afootnote{H; \skt{सोत्पलसारिवा}\ A}}\skt{~। }\pend
\pstart\versepar \skt{अनन्ता }\edtext{\skt{शारिवा}}{\lemma{\skt{शारिवा}}\Afootnote{H; \skt{सारिवा}\ A}}\skt{ रास्ना }\edtext{\skt{पद्म\-म्म\-धु\-क\-मेव}}{\lemma{\skt{पद्मम्}}\Afootnote{H; \skt{पद्मा}\ A}}\skt{ च~॥}\pend

\pstart\versepar \ledsidenote{\passref{3.10.61}}\edtext{\skt{वृह\-ती\-द्व\-य\-का\-श्म\-र्य\-क्षी\-री\-शु\-ङ्गा\-त्वचो}}{\lemma{\skt{वृह\appbrk{}ती\appbrk{}द्व\appbrk{}य\appbrk{}का\appbrk{}श्म\appbrk{}र्य\appbrk{}क्षी\appbrk{}री\appbrk{}शु\appbrk{}ङ्गा\appbrk{}त्वचो}}\Afootnote{H; \skt{बृहत्यौ काश्मरी चापि क्षीरिशुङ्गास्त्वचो}\ A}}\skt{ घृतं~। }\pend
\pstart\versepar \skt{}\edtext{\skt{पृथक्पर्ण्णी}}{\lemma{\skt{पृथक्पर्ण्णी}}\Afootnote{H; \skt{पृश्निपर्णी}\ A}}\skt{ वला }\edtext{\skt{शिग्र}}{\lemma{\skt{शिग्र}}\Afootnote{H; \skt{शिग्रुः}\ A}}\skt{ श्वदंष्ट्रा }\edtext{\skt{मधुयष्टिका}}{\lemma{\skt{मधुयष्टिका}}\Afootnote{H; \skt{मधुपर्णिका}\ A}}\skt{~॥}\pend

\pstart\versepar \ledsidenote{\passref{3.10.62}}\skt{शृङ्गाटकम्विसं द्राक्षा }\edtext{\skt{कसेरु}}{\lemma{\skt{कसेरु}}\Afootnote{H; \skt{कशेरु}\ A}}\skt{ मधुकं सिता~। }\pend
\pstart\versepar \skt{वत्सैते सप्तयोगाः }\edtext{\edtext{\skt{स्युर\-र्द्ध\-श्लो\-क\-स\-मा\-पनाः}}{\lemma{\skt{स्युर्}}\Afootnote{H; \skt{स्यु}\ A}}}{\lemma{\skt{अर्द्ध\appbrk{}श्लो\appbrk{}क\appbrk{}स\appbrk{}मा\appbrk{}पनाः}}\Afootnote{H; \skt{रर्ध\appbrk{}श्लो\appbrk{}क\appbrk{}स\appbrk{}मा\appbrk{}पनाः}\ A}}\skt{~॥ }\pend
\pstart\versepar \skt{}\edtext{\skt{यथाक्रमं}}{\lemma{\skt{यथाक्रमं}}\Afootnote{H; \skt{यथासंख्यं}\ A}}\skt{ प्रयोक्तव्या }\edtext{\skt{गर्भस्रावे}}{\lemma{\skt{गर्भस्रावे}}\Afootnote{A; \skt{गर्ब्भस्रावे}\ H}}\skt{ पयोयुताः~।}\pend

\pstart\versepar \ledsidenote{\passref{3.10.63}}\edtext{\skt{कपि\-त्थ\-वि\-ल्व\-वृ\-ह\-ती\-प\-टो\-ले\-क्षु\-नि\-दि\-ग्धिका}}{\lemma{\skt{कपि\appbrk{}त्थ\appbrk{}वि\appbrk{}ल्व\appbrk{}वृ\appbrk{}ह\appbrk{}ती\appbrk{}प\appbrk{}टो\appbrk{}ले\appbrk{}क्षु\appbrk{}नि\appbrk{}दि\appbrk{}ग्धिका}}\Afootnote{H; \skt{कपि\appbrk{}त्थ\appbrk{}बृ\appbrk{}ह\appbrk{}ती\appbrk{}बि\appbrk{}ल्व\appbrk{}प\appbrk{}टो\appbrk{}ले\appbrk{}क्षु\appbrk{}नि\appbrk{}दि\appbrk{}ग्धिका}\ A}}\skt{~॥ }\pend
\pstart\versepar \skt{मूलानि क्षीरसिद्धानि पाय\-ये\-द्भि\-ष\-ग\-ष्टमे~।}\pend

\pstart\versepar \ledsidenote{\passref{3.10.64}}\skt{नवमे मधुकानन्ता }\edtext{\skt{पय\-स्या\-शा\-रि\-वा\-म्पि\-वेत्}}{\lemma{\skt{पय\appbrk{}स्या\appbrk{}शा\appbrk{}रि\appbrk{}वाम्}}\Afootnote{H; \skt{पयस्यासारिवाः}\ A}}\skt{~॥ }\pend
\pstart\versepar \skt{}\edtext{\skt{पयस्तु}}{\lemma{\skt{पयस्तु}}\Afootnote{H; \skt{क्षीरं शुण्ठीप}\ A}}\skt{ }\edtext{\edtext{\skt{श्रित\-शी\-त\-म्प्र\-श\-स्यते}}{\lemma{\skt{श्रितशीतम्}}\Afootnote{H; \emph{om.}\ A}}}{\lemma{\skt{प्रशस्यते}}\Afootnote{H; \skt{यस्याभ्यां सिद्धं स्याद्दशमे हितम्}\ A}}\skt{~॥}\pend

\pstart\versepar \ledsidenote{\passref{3.10.65}}\skt{सक्षीरा वा हिता शुण्ठी मधुकं }\edtext{\skt{देवदारु च}}{\lemma{\skt{देवदारु च}}\Afootnote{H; \skt{सुरदारु च}\ A}}\skt{~। }\pend
\pstart\versepar \skt{}\edtext{\skt{एवमप्याय्यते}}{\lemma{\skt{अप्याय्यते}}\Afootnote{H; \skt{आप्यायते}\ A}}\skt{ }\edtext{\skt{गर्भस्तीव्रा}}{\lemma{\skt{गर्भस्}}\Afootnote{A; \skt{गर्ब्भस्}\ H}}\skt{ रुक्चोपशाम्यति~॥}\pend

\pstart \ledsidenote{\passref{3.10.57.2}}\edtext{\skt{व्यवस्थिते च}}{\lemma{\skt{व्यवस्थिते च}}\Afootnote{A; \skt{व्यवस्थि}\ H}}\skt{ }\edtext{\skt{गर्भे}}{\lemma{\skt{गर्भे}}\Afootnote{A; \emph{om.}\ H}}\skt{ गव्ये\-नो\-दु\-म्ब\-र\-श\-ला\-टु\-सि\-द्धेन पयसा भोजयेत्~। अतीते }\edtext{\skt{लवणं}}{\lemma{\skt{लवणं}}\Afootnote{H; \skt{लवण}\ A}}\skt{ }\edtext{\edtext{\skt{वर्ज्या\-भि\-रु\-द्दा\-लु\-का\-दीनां}}{\lemma{\skt{वर्ज्याभिर्}}\Afootnote{H; \skt{स्नेह\appbrk{}व\appbrk{}र्ज्या\appbrk{}भि\appbrk{}र्य\appbrk{}वा\appbrk{}गू\appbrk{}भिर्}\ A}}}{\lemma{\skt{उद्दालुकादीनां}}\Afootnote{H; \skt{उद्दालकादीनां}\ A}}\skt{ }\edtext{\edtext{\skt{पाच\-नी\-यो\-प\-सं\-भृ\-ता\-भि\-रु\-प\-क्रा\-मे\-त्या\-वन्तो}}{\lemma{\skt{पाच\appbrk{}नी\appbrk{}यो\appbrk{}प\appbrk{}सं\appbrk{}भृ\appbrk{}ता\appbrk{}भिर्}}\Afootnote{H; \skt{पाच\appbrk{}नी\appbrk{}यो\appbrk{}प\appbrk{}सं\appbrk{}स्कृ\appbrk{}ता\appbrk{}भिर्}\ A}}}{\lemma{\skt{उपक्रामेत्}}\Afootnote{H; \skt{उपक्रमेत}\ A}}\skt{ मासा गर्भस्य }\edtext{\skt{तावन्त्यहानि}}{\lemma{\skt{तावन्त्यहानि}}\Afootnote{\skt{तावन्त्यहानि}\ A H}}\skt{ }\edtext{\skt{वस्त्युदरशूले}}{\lemma{\skt{वस्त्युदरशूले}}\Afootnote{H; \skt{बस्त्यु\appbrk{}द\appbrk{}र\appbrk{}शू\appbrk{}लेषु}\ A}}\skt{ पुरा\-ण\-गु\-ड\-न्दी\-प\-नी\-य\-सं\-यु\-क्त\-म्पा\-य\-ये\-द\-रि\-ष्टम्वा~। }\edtext{\skt{वाता\-त\-प\-श्र\-म\-से\-विन्या}}{\lemma{\skt{वाता\appbrk{}त\appbrk{}प\appbrk{}श्र\appbrk{}म\appbrk{}से\appbrk{}विन्या}}\Afootnote{H; \skt{वातोपद्र}\ A}}\skt{ }\edtext{\skt{वातो\-प\-ही\-ही\-त\-त्वा\-त्स्रो\-तसां}}{\lemma{\skt{वातो\appbrk{}प\appbrk{}ही\appbrk{}ही\appbrk{}त\appbrk{}त्वात्}}\Afootnote{H; \skt{वगृहीतत्वात्}\ A}}\skt{ लीयते गर्भः }\edtext{\skt{सोति\-का\-ल\-म\-व\-ति\-ष्ठ\-मानो}}{\lemma{\skt{सोतिकालम्}}\Afootnote{H; \skt{सोऽतिकालम्}\ A}}\skt{ व्यापद्यते~। तां मृदुना }\edtext{\skt{स्नेहादिना}}{\lemma{\skt{स्नेहादिना}}\Afootnote{H; \skt{स्नेहादि}\ A}}\skt{ }\edtext{\skt{क्रमे\-णो\-प\-च\-रेत्~। उत्क्रोसरसं}}{\lemma{\skt{क्रमे\appbrk{}णो\appbrk{}प\appbrk{}च\appbrk{}रेत्~। उत्क्रोसरसं}}\Afootnote{H; \skt{क्रमे\appbrk{}णो\appbrk{}प\appbrk{}च\appbrk{}रेत्, उत्क्रोशरससं}\ A}}\skt{ }\edtext{\edtext{\edtext{\skt{यवा\-गू\-ञ्चा\-ल्प\-स्ने\-हा\-म्प्र\-य\-च्छेत्}}{\lemma{\skt{यवागूञ्}}\Afootnote{H; \skt{सिद्धाम}\ A}}}{\lemma{\skt{चाल्पस्नेहाम्}}\Afootnote{H; \skt{नल्पस्नेहां}\ A}}}{\lemma{\skt{प्रयच्छेत्}}\Afootnote{H; \skt{यवागूं पाययेत्, माष\appbrk{}ति\appbrk{}ल\appbrk{}बि\appbrk{}ल्व\appbrk{}श\appbrk{}ला\appbrk{}टु\appbrk{}सि\appbrk{}द्धान्वा कुल्मा\appbrk{}षा\appbrk{}न्भ\appbrk{}क्ष\appbrk{}ये\appbrk{}न्म\appbrk{}धु\appbrk{}मा\appbrk{}ध्वीकं चानु\appbrk{}पि\appbrk{}बे\appbrk{}त्स\appbrk{}प्त\appbrk{}रा\appbrk{}त्रम्}\ A}}\skt{~। काला\-ती\-त\-स्था\-यिनि }\edtext{\skt{गर्भविशेषतो}}{\lemma{\skt{गर्भविशेषतो}}\Afootnote{H; \skt{गर्भे विशेषतः}\ A}}\skt{ }\edtext{\skt{धान्य\-मु\-दू\-ख\-ल\-मु\-स\-ले\-ना\-भि\-ह\-न्यात्}}{\lemma{\skt{धान्यम्}}\Afootnote{H; \skt{सधान्यम्}\ A}}\skt{~। }\edtext{\skt{विषमम्वा}}{\lemma{\skt{विषमम्वा}}\Afootnote{H; \skt{विषमे वा}\ A}}\skt{ }\edtext{\edtext{\skt{याना\-स\-न\-मु\-प\-से\-वेत्}}{\lemma{\skt{यानासनम्}}\Afootnote{H; \skt{यानासने}\ A}}}{\lemma{\skt{उपसेवेत्}}\Afootnote{H; \skt{सेवेत}\ A}}\skt{~। वाताभिपन्न एव शुष्यति गर्भः समातुः कुक्षिन्न पूरयति~। मन्दं स्पन्दते च तं }\edtext{\skt{वृंहणीयः}}{\lemma{\skt{वृंहणीयः}}\Afootnote{H; \skt{बृंहणीयैः}\ A}}\skt{ }\edtext{\skt{पयो\-भि\-र्म्मां\-स\-र\-सै\-श्चो\-प\-च\-रेत्}}{\lemma{\skt{पयोभिर्}}\Afootnote{H; \skt{पयो भिर्}\ A}}\skt{~। शुक्रशोणितं }\edtext{\edtext{\edtext{\skt{वायु\-नो\-प\-रु\-द्ध\-म\-न\-व\-क्त्रा\-न्त\-जी\-व\-मा\-ध्मा\-प\-य\-त्यु\-द\-र\-न्त\-त्क\-दा\-चि\-द्य\-दृ\-च्छ\-यो\-प\-शा\-न्त\-न्नै\-ग\-मेषा}}{\lemma{\skt{वायु\appbrk{}नो\appbrk{}प\appbrk{}रु\appbrk{}द्धम्}}\Afootnote{H; \skt{वायु\appbrk{}ना\appbrk{}भि\appbrk{}प्र\appbrk{}पन्}\ A}}}{\lemma{\skt{अन\appbrk{}व\appbrk{}क्त्रा\appbrk{}न्त\appbrk{}जी\appbrk{}वम्}}\Afootnote{H; \skt{नम\appbrk{}व\appbrk{}क्रा\appbrk{}न्त\appbrk{}जी\appbrk{}वम्}\ A}}}{\lemma{\skt{तत्}}\Afootnote{H; \skt{तं}\ A}}\skt{ पहृतमिति }\edtext{\skt{वदन्ति~। तमेव}}{\lemma{\skt{वदन्ति~। तम्}}\Afootnote{H; \skt{भाषन्ते, तम्}\ A}}\skt{ }\edtext{\edtext{\skt{कदा\-चि\-त्प्र\-वि\-ली\-य\-मा\-न\-मा\-पू\-र्य\-मा\-णम्वा}}{\lemma{\skt{प्रवि\appbrk{}ली\appbrk{}य\appbrk{}मा\appbrk{}नम्}}\Afootnote{H; \skt{प्रलीयमानं}\ A}}}{\lemma{\skt{आपूर्यमाणम्वा}}\Afootnote{H; \emph{om.}\ A}}\skt{ नागो\-द\-र\-मि\-त्या\-हु\-स्त\-त्रापि }\edtext{\skt{लीन\-व\-त्प्र\-ति\-कारः}}{\lemma{\skt{प्रतिकारः}}\Afootnote{H; \skt{प्रतीकारः}\ A}}\skt{~। }\edtext{\skt{निवृ\-त्त\-प्र\-स\-वायाः}}{\lemma{\skt{निवृ\appbrk{}त्त\appbrk{}प्र\appbrk{}स\appbrk{}वायाः}}\Afootnote{H; \emph{om.}\ A}}\skt{ }\edtext{\skt{पुनः}}{\lemma{\skt{पुनः}}\Afootnote{H; \emph{om.}\ A}}\skt{ }\edtext{\skt{षड्भ्यो}}{\lemma{\skt{षड्भ्यो}}\Afootnote{H; \emph{om.}\ A}}\skt{ }\edtext{\skt{वर्षेभ्य}}{\lemma{\skt{वर्षेभ्य}}\Afootnote{H; \emph{om.}\ A}}\skt{ }\edtext{\edtext{\skt{ऊर्द्ध्वम्प्रसवत्याः}}{\lemma{\skt{ऊर्द्ध्वम्}}\Afootnote{H; \emph{om.}\ A}}}{\lemma{\skt{प्रसवत्याः}}\Afootnote{H; \emph{om.}\ A}}\skt{ }\edtext{\edtext{\edtext{\skt{कुमा\-रोऽ\-ल्पा\-यु\-र्भ\-वति}}{\lemma{\skt{कुमारो}}\Afootnote{H; \emph{om.}\ A}}}{\lemma{\skt{ऽल्पायुर्}}\Afootnote{H; \emph{om.}\ A}}}{\lemma{\skt{भवति}}\Afootnote{H; \emph{om.}\ A}}\skt{~।}\pend

\pstart \ledsidenote{\passref{3.10.67}}\skt{अथ }\edtext{\skt{गर्ब्भिणी}}{\lemma{\skt{गर्ब्भिणी}}\Afootnote{H; \skt{गर्भिणीं}\ A}}\skt{ }\edtext{\edtext{\skt{व्याध्यु\-त्प\-न्ना\-वा\-त्य\-ये\-च्छ\-र्द्द\-ये\-न्म\-धु\-रा\-म्ले\-ना\-न्नो\-प\-हि\-तेन}}{\lemma{\skt{व्याध्यु\appbrk{}त्प\appbrk{}न्ना\appbrk{}वा\appbrk{}त्य\appbrk{}ये\appbrk{}च्छ\appbrk{}र्द्द\appbrk{}येन्}}\Afootnote{\skt{व्याध्यु\appbrk{}त्प\appbrk{}त्ता\appbrk{}व\appbrk{}त्यये छर्दयेन्}\ A; \skt{व्याध्युत्पन्ना वात्यये च्छर्द्दयेन्}\ H}}}{\lemma{\skt{मधु\appbrk{}रा\appbrk{}म्ले\appbrk{}ना\appbrk{}न्नो\appbrk{}प\appbrk{}हि\appbrk{}तेन}}\Afootnote{H; \skt{मधु\appbrk{}रा\appbrk{}म्ले\appbrk{}ना\appbrk{}न्नो\appbrk{}प\appbrk{}हि\appbrk{}तेना}\ A}}\skt{~। }\edtext{\skt{अनुलोमयेच्च}}{\lemma{\skt{अनुलोमयेच्च}}\Afootnote{H; \skt{नुलोमयेच्च}\ A}}\skt{ }\edtext{\skt{संश\-म\-नी\-य\-म्मृदु}}{\lemma{\skt{संशमनीयम्}}\Afootnote{H; \skt{संशमनीयं च}\ A}}\skt{ विदध्यात्~। }\edtext{\skt{अन्न\-पा\-न\-यो\-रिति}}{\lemma{\skt{इति}}\Afootnote{H; \emph{om.}\ A}}\skt{~। अश्नीयाच्च }\edtext{\edtext{\skt{मृदु\-वी\-र्य\-ङ्ग\-र्ब्भा\-वि\-रु\-द्धांश्च}}{\lemma{\skt{मृदुवीर्यङ्}}\Afootnote{H; \skt{मृदुवीर्यं मधुरप्रायं}\ A}}}{\lemma{\skt{गर्ब्भाविरुद्धांश्च}}\Afootnote{H; \skt{गर्भाविरुद्धं च, गर्भाविरुद्धाश्च}\ A}}\skt{ }\edtext{\skt{क्रिया\-यो\-गा\-न्वि\-द\-धीत}}{\lemma{\skt{क्रियायोगान्}}\Afootnote{H; \skt{क्रिया यथायोगं}\ A}}\skt{ }\edtext{\skt{इति}}{\lemma{\skt{इति}}\Afootnote{H; \skt{मृदुप्रायाः}\ A}}\skt{~॥}\pend

\pstart\versepar \ledsidenote{\passref{3.10.68}}\edtext{\skt{सुवर्ण्णं}}{\lemma{\skt{सुवर्ण्णं}}\Afootnote{H; \skt{सौवर्णं}\ A}}\skt{ }\edtext{\skt{र्ण्णचूसुकृदं}}{\lemma{\skt{र्ण्णचूसुकृदं}}\Afootnote{H\textsuperscript{pc}; \skt{सुकृतं चूर्णं}\ A; \skt{सुकृदं}\ H\textsuperscript{ac}}}\skt{ }\edtext{\skt{कुष्ठम्मधुवचा}}{\lemma{\skt{मधुवचा}}\Afootnote{H; \skt{मधु}\ A}}\skt{ घृतं~। }\pend
\pstart\versepar \skt{}\edtext{\skt{मत्स्या\-क्ष\-कः\-शं\-ख\-पुष्पी}}{\lemma{\skt{मत्स्या\appbrk{}क्ष\appbrk{}कः\appbrk{}शं\appbrk{}ख\appbrk{}पुष्पी}}\Afootnote{H; \skt{वचा~। मत्स्याक्षकः शङ्खपुष्पी}\ A}}\skt{ }\edtext{\skt{सर्प्पिः}}{\lemma{\skt{सर्प्पिः}}\Afootnote{H; \skt{मधु सर्पिः}\ A}}\skt{ }\edtext{\edtext{\skt{काञ्चनमेव}}{\lemma{\skt{काञ्चनम्}}\Afootnote{H; \skt{सकाञ्चनम्}\ A}}}{\lemma{\skt{एव}}\Afootnote{H; \emph{om.}\ A}}\skt{ }\edtext{\skt{च}}{\lemma{\skt{च}}\Afootnote{H; \emph{om.}\ A}}\skt{~॥}\pend

\pstart\versepar \ledsidenote{\passref{3.10.69}}\skt{अर्क्कपुष्पी मधुघृतं चूर्ण्णितं }\edtext{\skt{कनकन्तथा}}{\lemma{\skt{तथा}}\Afootnote{H; \skt{वचा}\ A}}\skt{~। }\pend
\pstart\versepar \skt{}\edtext{\skt{सुव\-र्ण्ण\-चू\-र्ण्ण\-वै\-डूर्यं}}{\lemma{\skt{सुव\appbrk{}र्ण्ण\appbrk{}चू\appbrk{}र्ण्ण\appbrk{}वै\appbrk{}डूर्यं}}\Afootnote{H; \skt{हेमचूर्णानि कैडर्यः}\ A}}\skt{ }\edtext{\skt{स्वेता}}{\lemma{\skt{स्वेता}}\Afootnote{H; \skt{श्वेता}\ A}}\skt{ दूर्व्वा }\edtext{\skt{घृतम्मधुः}}{\lemma{\skt{मधुः}}\Afootnote{H; \skt{मधु}\ A}}\skt{~॥}\pend

\pstart\versepar \ledsidenote{\passref{3.10.70}}\skt{चत्वारोऽभिहिताः }\edtext{\skt{प्रास्या}}{\lemma{\skt{प्रास्या}}\Afootnote{H; \skt{प्राशाः}\ A}}\skt{ }\edtext{\skt{अर्द्ध\-श्लो\-क\-स\-मा\-पनाः}}{\lemma{\skt{अर्द्ध\appbrk{}श्लो\appbrk{}क\appbrk{}स\appbrk{}मा\appbrk{}पनाः}}\Afootnote{H; \skt{श्लोकार्धेषु चतुर्ष्वपि}\ A}}\skt{~। }\pend
\pstart\versepar \skt{कुमा\-रा\-णा\-म्व\-पु\-र्मेधा }\edtext{\skt{बलारोग्यकराः}}{\lemma{\skt{बलारोग्यकराः}}\Afootnote{H; \skt{बल}\ A}}\skt{ }\edtext{\skt{शुभा}}{\lemma{\skt{शुभा}}\Afootnote{H; \skt{बु}\ A}}\skt{ }\edtext{\skt{इति}}{\lemma{\skt{इति}}\Afootnote{H; \skt{द्धिविवर्धनाः}\ A}}\skt{~॥}\pend

\pstart \ledsidenote{\passref{3.10.trailer}}\skt{इति }\edtext{\skt{गर्ब्भि\-णी\-व्या\-क\-रणं}}{\lemma{\skt{गर्ब्भि\appbrk{}णी\appbrk{}व्या\appbrk{}क\appbrk{}रणं}}\Afootnote{\skt{सुश्रु\appbrk{}त\appbrk{}सं\appbrk{}हि\appbrk{}तायां शारीरस्थाने गर्भि\appbrk{}णी\appbrk{}व्या\appbrk{}क\appbrk{}रणं}\ A}}\skt{ }\edtext{\skt{शारीरे}}{\lemma{\skt{शारीरे}}\Afootnote{\skt{शारीरं}\ A}}\skt{ }\edtext{\edtext{\skt{दशमोऽध्यायः}}{\lemma{\skt{दशमो}}\Afootnote{\skt{नाम दशमो}\ A}}}{\lemma{\skt{ऽध्यायः}}\Afootnote{\skt{ऽध्यायः}\ A}}\skt{~॥ }\edtext{\skt{ओ}}{\lemma{\skt{ओ}}\Afootnote{\emph{om.}\ A}}\skt{~॥}\pend

\pstart\versepar \ledsidenote{\passref{3.10.70.1}}\skt{भूतचिन्ता शुक्रशुद्धिं }\edtext{\skt{गर्भ\-स\-ङ्क्र\-म\-ण\-न्तथा}}{\lemma{\skt{गर्भ\appbrk{}स\appbrk{}ङ्क्र\appbrk{}म\appbrk{}णन्}}\Afootnote{\skt{गर्ब्भ\appbrk{}स\appbrk{}ङ्क्र\appbrk{}म\appbrk{}णन्}\ H}}\skt{~। }\pend
\pstart\versepar \skt{}\edtext{\skt{गर्भ\-व्या\-क\-र\-ण\-ञ्चैव}}{\lemma{\skt{गर्भ\appbrk{}व्या\appbrk{}क\appbrk{}र\appbrk{}णञ्}}\Afootnote{\skt{गर्ब्भ\appbrk{}व्या\appbrk{}क\appbrk{}र\appbrk{}णञ्}\ H}}\skt{ काय\-व्या\-क\-र\-णञ्च यत्~॥}\pend

\pstart\versepar \ledsidenote{\passref{3.10.70.2}}\skt{मर्म्म\-सि\-रा\-वि\-भ\-क्तिश्च सिरा\-ता\-ड\-न\-मेव च~। }\pend
\pstart\versepar \skt{धम\-नी\-व्या\-क\-रण चैव गर्ब्भिण्या पूर्यते दश~॥}\pend

\pstart \ledsidenote{\passref{3.trailer}}\edtext{\skt{॥ इति शारीरस्थानं}}{\lemma{\skt{॥ इति शारीरस्थानं}}\Afootnote{H; \skt{इति भगवता श्रीध\appbrk{}न्व\appbrk{}न्त\appbrk{}रि\appbrk{}णो\appbrk{}प\appbrk{}दि\appbrk{}ष्टायां तच्छिष्येण महर्षिणा सुश्रुतेन विरचितायां सुश्रु\appbrk{}त\appbrk{}सं\appbrk{}हि\appbrk{}तायां तृतीयं शारीरस्थानं}\ A}}\skt{ }\edtext{\skt{समाप्तं}}{\lemma{\skt{समाप्तं}}\Afootnote{A; \skt{समाप्तं~॥ {\symfont ❈}}\ H}}\skt{~॥}\pend

\endnumbering
